\begin{bt}%[2D1N1-1]
% nguon: Archive/SP_SauLuoi_2024-2025/SP_HKI/2-TN-DS-TL-SGD-HUE-HKI-N24-25.tex #43
%[Lớp 12- Kiểm tra học kỳ I- SGD-HUẾ]%[Nguyễn Văn Sang]
    Tìm các khoảng đơn điệu của hàm số $y = x^3 - 3x + 1$.
    \loigiai{
        Ta có $y = x^3 - 3x + 1$.\\
        Tập xác định $D = \mathbb{R}$.\\
        Đạo hàm $y' = 3x^2 - 3$.\\
        Cho $y' = 0 \Leftrightarrow 3x^2 - 3 = 0 \Leftrightarrow x^2 = 1 \Leftrightarrow x = \pm 1$.\\
        Lập bảng xét dấu $y'$
        \begin{center}
            \begin{tikzpicture}[scale=1, font=\footnotesize, line join=round, line cap=round, >=stealth]
                \tkzTabInit[nocadre=false,lgt=1,espcl=2,deltacl=0.5]
                {$x$/.7 ,$y'$/.7,$y$/2}
                {$-\infty$ , $-1$ , $1$ , $+\infty$}
                \tkzTabLine{ , + , $0$ , - , $0$ , + , }
                \tkzTabVar{-/$-\infty$ , +/$ $ , -/$ $ , +/$+\infty$}
            \end{tikzpicture}
        \end{center}
        Từ bảng biến thiên, ta có
        \begin{itemize}
            \item Hàm số đồng biến trên các khoảng $(-\infty; -1)$ và $(1; +\infty)$.
            \item Hàm số nghịch biến trên khoảng $(-1; 1)$.
        \end{itemize}
    }
\end{bt}

\begin{bt}%[2D1V1-5]
% nguon: Archive/SP_SauLuoi_2024-2025/SP_HKI/2-TN-DS-TL-SGD-HUE-HKI-N24-25.tex #46
%[Lớp 12- Kiểm tra học kỳ I- SGD-HUẾ]%[Nguyễn Văn Sang]
    Một chất điểm đang chuyển động trên một đường thẳng và gặp một con dốc cao muốn vượt qua. Xem chân dốc là điểm $A$, khi chất điểm qua $A$ là thời điểm bắt đầu leo dốc với phương trình chuyển động là $s(t) = at - 2t^2$ cho đến khi lên đỉnh dốc ($a$ là tham số thực dương, $t \geq 0$ tính bằng giây và $s$ tính bằng mét). Biết dốc có độ dài bằng $50m$ và xe muốn leo dốc thành công thì phải tồn tại $t > 0$ sao cho $s(t) = 50$. Tính vận tốc ban đầu nhỏ nhất để xe leo dốc thành công.
    \loigiai{
        Ta có phương trình chuyển động của chất điểm $s(t) = at - 2t^2$.\\
        Để xe leo dốc thành công thì phải tồn tại $t > 0$ sao cho $s(t) = 50$, hay $at - 2t^2 = 50$ hay $2t^2 - at + 50 = 0$.\\
        Để phương trình trên có nghiệm $t > 0$ thì trước hết $\Delta \geq 0$, tức là $$(-a)^2 - 4.2.50 \geq 0 \Leftrightarrow a^2 \geq 400 \Leftrightarrow a \geq 20 \left(\text{ do } a>0\right).$$
       Gọi $t_1, t_2$ là hai nghiệm của phương trình $2t^2 - at + 50 = 0$. Theo định lý Viète, ta có
	   $$\heva{&t_1 + t_2 = \dfrac{a}{2}\\&t_1 t_2 = \dfrac{50}{2} = 25.}$$
       Để tồn tại nghiệm $t>0$, ta có $2$ trường hợp xảy ra
       \begin{itemize}
           \item  Cả $2$ nghiệm đều dương $\Leftrightarrow t_1 + t_2 > 0$ và $t_1t_2 >0$ (đã thoả mãn).
           \item Có $1$ nghiệm dương và 1 nghiệm âm. Lúc này ta có $t_1t_2<0$. Điều này mâu thuẫn với $t_1t_2=25>0$, vậy trường hợp này không xảy ra.
       \end{itemize}
        Vậy điều kiện để phương trình có nghiệm dương là  $\Delta \geq 0 \Leftrightarrow a\geq 20$.\\
        Vận tốc của chất điểm là $v(t) = s'(t) = a - 4t$. Vận tốc ban đầu ứng với $t=0$ là $v(0) = a$.\\
        Để xe leo dốc thành công thì phương trình $at - 2t^2 = 50$ phải có nghiệm dương, tức là $a \geq 20$.\\
        Vậy vận tốc ban đầu nhỏ nhất để xe leo dốc thành công là $20$ m/s.
    }
\end{bt}

\begin{bt}%[2D1H5-7]
% nguon: Archive/SP_SauLuoi_2024-2025/SP_HKI/2-TN-DS-TLN-TL-THPT-TranPhu-VungTau-HKI-NH24-25.tex #54
%[Dự án đề kiểm tra môn Toán khối 12 HKI NH24-25 - Đợt 4 - Thành Đức Trung]%[THPT Trần Phú - Bà Rịa Vũng Tàu]
Cho hàm số $y=x^3-3x^2+5$ có đồ thị $(C)$, gọi $A(x;y)$ là điểm cực đại của đồ thị hàm số. Tính $x+2y$.
\loigiai
{
Ta có $y'=3x^2-6x$. \\
Xét $y'=0 \Leftrightarrow \hoac{ & x=0 \\ & x=2.}$ \\
Bảng biến thiên
\begin{center}
\begin{tikzpicture}[scale=1]
\tkzTabInit[nocadre=false, lgt=1.2, espcl=2.5, deltacl=0.6]{$x$/0.6, $y'$/0.6, $y$/2}{$-\infty$, $0$, $2$, $+\infty$}
\tkzTabLine{,+,0,-,0,+,}
\tkzTabVar{-/$-\infty$, +/$5$, -/$1$, +/$+\infty$}
\end{tikzpicture}
\end{center}
Từ bảng biến thiên ta có $A(0;5)$. \\
Vậy $x+2y=0+2\cdot5=10$.
}
\end{bt}

\begin{bt}%[2D1H5-7]
% nguon: Archive/SP_SauLuoi_2024-2025/SP_HKI/2-TN-DS-TLN-TL-THPT-TranPhu-VungTau-HKI-NH24-25.tex #55
%[Dự án đề kiểm tra môn Toán khối 12 HKI NH24-25 - Đợt 4 - Thành Đức Trung]%[THPT Trần Phú - Bà Rịa Vũng Tàu]
Cho hàm số $y=f(x)=\dfrac{x^2-3x+7}{x-1}$. Tiệm cận xiên có dạng $y=ax+b$. Tính $b^2-a^2$.
\loigiai
{
Ta có $y=f(x)=\dfrac{x^2-3x+7}{x-1}=x-2+\dfrac{5}{x-1}$. \\
Vì $\lim\limits_{x\to-\infty} [f(x)-(x-2)]=\lim\limits_{x\to-\infty} \dfrac{5}{x-1}=0$ nên đồ thị hàm số $y=f(x)$ có đường tiệm cận xiên là $y=x-5$. \\
Vậy $b^2-a^2=\left(-5\right)^2-1^2=24$.
}
\end{bt}

\begin{bt}%[2D1C2-1]
% nguon: Archive/SP_SauLuoi_2024-2025/SP_HKI/2-TN-DS-TLN-THPT-DinhTienHoang-NinhBinh-HKI-NH24-25.tex #101
%[Dự án đề kiểm tra Toán 12 CHKI NH24-25- Nguyễn Trần Anh Tuấn]%[THPT Đinh Tiên Hoàng - Tỉnh Ninh Bình]
	Cho hàm số $y=f(x)$ có đạo hàm $f'(x)=\left(x^2-x\right)\left(x^2-4x+3\right)$, $\forall x \in \mathbb{R}$. Tính tổng tất cả các giá trị nguyên của tham số $m$ để hàm số $g(x)=f\left(x^2+m\right)$ có $3$ điểm cực trị.
	\loigiai{
		Ta có $f'(x)=x(x-3)(x-1)^2$.\\
		Ta có $g'(x)=2xf'(x^2+m)=2x(x^2+m)(x^2+m-3)(x^2+m-1)^2$.\\
		Do đó $g'(x)=0$ có ba điểm cực trị khi phương trình $g'(x)=0$ có $3$ nghiệm bội lẻ phân biệt.\\
		Hay $\hoac{&x^2=-m\\&x^2=-m+3}$ có $2$ nghiệm thực phân biệt $ \Leftrightarrow 0 \le m < 3 $.\\
		Vậy tổng tất cả các giá trị nguyên của tham số $m$ để hàm số $g(x)$ có $3$ điểm cực trị là $0+1+2=3$.
	}
\end{bt}

\begin{bt}%[2D1H3-1]
% nguon: Archive/SP_SauLuoi_2024-2025/SP_GHKI/2-TL-TN-TLN-THPT-PhamVanSang-HCM-GHKI-NH24-25€.tex #206
%[Dự án đề kiểm tra Toán 12 GHKI NH24-25- Dương Công Tạo]%[Phạm Văn Sáng - Tp HCM]
	Người ta giới thiệu một loại thuốc để kích thích sự sinh sản của một loại vi khuẩn. Sau $t$ phút, số lượng vi khuẩn được xác định theo công thức là $f(t)=-t^3+30t^2+1\,000$ với $0\leq t\leq 30$. Hỏi sau bao nhiêu phút thí số lượng vi khuẩn đạt giá trị cực đại
	\loigiai{
		Xét hàm số $f(t)=-t^3+30t^2+1\,000$ trên đoạn $[0;30]$.\\
		Ta có $f'(t)=-3t^2+60t$; $f'(t)=0\Leftrightarrow -3t^2+60t=0\Leftrightarrow\hoac{&t=0\in[0;30]\\&t=20\in[0;30].}$\\
		Mà $f(0)=1\,000$; $f(20)=5\,000$; $f(30)=1\,000$.\\
		Vậy $\max\limits_{[0;30]}f(t)=f(20)=5\,000$ hay số lượng vi khuẩn đạt giá trị cực đại là $5\,000$ sau $20$ phút.
	}
\end{bt}

\begin{bt}%[2D1V2-7]
% nguon: Archive/SP_SauLuoi_2024-2025/SP_GHKI/2-TL-TN-DS-THPT-PhanChuTrinh-Binhthuan-GHKI-NH24-25.tex #221
%[Dự án đề kiểm tra Toán 12 GHKI NH24-25- Hieu Phan]%[THPT Phan Chu Trinh- Bình Thuận]
    Một doanh nghiệp dự kiến lợi nhuận khi sản xuất $x$ sản phẩm $\left(0\le x\le 500\right)$ được cho bởi hàm số $y=-x^3+525x^2$ (đơn vị: đồng). Doanh nghiệp đó cần sản suất bao nhiêu sản phẩm để thu được lợi nhuận cao nhất?
    \loigiai{
        Xét $y=-x^3+525x^2$ với $0\le x\le 500$, $x\in\mathbb{N}$.\\
        $ y'=-3x^2+1050x $; 
        $  y'=0\Leftrightarrow-3x^2+1050x=0\Leftrightarrow \hoac{& x=0 \\ &x=350. }$\\
        Bảng biến thiên
        \begin{center}
            % Cần khai báo \usepackage{tkz-tab}
            \begin{tikzpicture}[>=stealth]
                \tkzTabInit[nocadre=false,lgt=1,espcl=2,deltacl=0.5]{$x$/.7 ,$y'$/.7,$y$/2}
                {$0$ , $350$ , $500$}
                \tkzTabLine{ , + , $0$ , - , }
                \tkzTabVar{-/$ $ , +/$y_{\text{CĐ}}$ , -/$ $ }
            \end{tikzpicture}
            
        \end{center}
        Vậy lợi nhuận cao nhất khi sản xuất $ 350 $ sản phẩm.
    }
    
\end{bt}

\begin{bt}%[2D1V1-3]
% nguon: Archive/SP_SauLuoi_2024-2025/SP_GHKI/2-TL-TN-DS-THPT-PhanChuTrinh-Binhthuan-GHKI-NH24-25.tex #222
%[Dự án đề kiểm tra Toán 12 GHKI NH24-25- Hieu Phan]%[THPT Phan Chu Trinh- Bình Thuận]
    Có bao nhiêu giá trị nguyên của tham số $m$ thuộc khoảng $(-10;20)$ sao cho ứng với mỗi $m$, hàm số $y=x^3-x^2+3mx+1$ đồng biến trên $\mathbb{R}$?
    \loigiai{
        $y'=3x^2-2x+3m$.\\
        Hàm số đồng biến trên $ \mathbb{R} $
        \begin{eqnarray*}
            &\Leftrightarrow &y'\ge 0,\,\forall x\in\mathbb{R}\\
            &\Leftrightarrow& 3x^2-2x+3m\ge 0,\,\forall x\in\mathbb{R}\\ 
            &\Leftrightarrow& 1-9m\le 0 \Leftrightarrow m\ge\dfrac{1}{9}.
        \end{eqnarray*}
        Kết hợp với điều kiện $m\in\left(-10;20\right)$ và $m$ là số nguyên nên $m\in \{ 1;\,2;\ldots;19\}$.\\
        Vậy có $ 19 $ giá trị nguyên của tham số $m$ thỏa yêu cầu đề bài.}
\end{bt}

\begin{bt}%[2D1V3-1]
% nguon: Archive/SP_SauLuoi_2024-2025/SP_GHKI/2-TL-TN-DS-THPT-PhanChuTrinh-Binhthuan-GHKI-NH24-25.tex #223
%[Dự án đề kiểm tra Toán 12 GHKI NH24-25- Hieu Phan]%[THPT Phan Chu Trinh- Bình Thuận]
    Gọi $M$, $m$ lần lượt là giá trị lớn nhất và giá trị nhỏ nhất của hàm số \break $f(x)=3x^3-9x^2+5x+2$ trên đoạn $\left[\dfrac{1}{2};\dfrac{5}{2}\right]$. Giá trị $M+m$ bằng bao nhiêu? (kết quả viết dưới dạng số thập phân và làm tròn đến hàng phần trăm).
    \loigiai{
        $f'(x)=9x^2-18x+5$.\\
        $f'(x)=0\Leftrightarrow 9x^2-18x+5=0\Leftrightarrow\hoac{& x=\dfrac{5}{3} \\ & x=\dfrac{1}{3} \text{ (loại)}.}$\\
        Ta có $f\left(\dfrac{5}{3}\right)=-\dfrac{7}{9}$,  $f\left(\dfrac{1}{2}\right)=\dfrac{21}{8}$, $f\left(\dfrac{5}{2}\right)=\dfrac{41}{8}$.\\
        Vậy $M=\dfrac{41}{8},\,\,m=-\dfrac{7}{9}$ nên $M+m=\dfrac{41}{8}+\left(-\dfrac{7}{9}\right)=\dfrac{313}{72}\approx 4{,}35$.}
\end{bt}

\begin{bt}%[2D1V5-8]
% nguon: Archive/SP_SauLuoi_2024-2025/SP_GHKI/2-TL-TN-DS-THPT-PhanChuTrinh-Binhthuan-GHKI-NH24-25.tex #224
%[Dự án đề kiểm tra Toán 12 GHKI NH24-25- Hieu Phan]%[THPT Phan Chu Trinh- Bình Thuận]
    \immini
    {
        Một cổng vòm có dạng nửa hình tròn trên mặt đất với bán kính $R=4m$.
    Người ta muốn đặt một khung hình chữ nhật $ABCD$ để thiết kế trang trí, với hai điểm $A$, $B$ đính trên vòm và $CD$ đặt trên mặt đất (tham khảo hình vẽ).
    Tìm khoảng cách giữa hai điểm $A$, $B$ (tính bằng mét, làm tròn kết quả đến hàng phần mười) để diện tích hình chữ nhật $ABCD$ lớn nhất.
    }
    { 
       
         \begin{tikzpicture}[scale=0.7, font=\footnotesize, line join=round, line cap=round, >=stealth]
        \draw[line width=0.4pt,black] (0,0) arc (0:180:3cm);
        \draw[line width=1.4pt](-6,0)--(0,0); 
         \draw (-5.22,0)node[above left]{$  D$}--(-5.22,2)node[ left]{$ 
             A$}--(-0.77,2)node[right]{$ B$}--(-0.77,0)node[above right]{$ C$}; 
    \end{tikzpicture}
    }
    \loigiai{
        Gọi khoảng cách giữa hai điểm $A$, $B$ là $2x$ với $0<x<4$ $\Rightarrow AD=\sqrt{16-x^2}$.\\
        Diện tích hình chữ nhật là $S=2x\sqrt{16-x^2}$.\\
        Xét hàm số $S(x)=2x\sqrt{16-x^2}$ trên khoảng $\left(0;4\right)$.\\
        Ta có $S'(x)=2\sqrt{16-x^2}-\dfrac{2x^2}{\sqrt{16-x^2}}$.\\
        $S'(x)=0\Rightarrow \hoac{&x=2\sqrt{2}  \\ & x=-2\sqrt{2}\text{ (loại)}.}$\\
         Bảng biến thiên
        \begin{center}
            \begin{tikzpicture}[>=stealth]
                \tkzTabInit[nocadre=false,lgt=1.3,espcl=2,deltacl=0.5]{$x$/.7 ,$S'(x)$/.7,$S(x)$/2}
                {$0$ , $2\sqrt{2}$ , $4$}
                \tkzTabLine{ , + , $0$ , - , }
                \tkzTabVar{-/$ $ , +/$y_{\text{CĐ}}$ , -/$ $}
            \end{tikzpicture}
        \end{center}
        Vậy diện tích hình chữ nhật $ABCD$ lớn nhất khi khoảng cách giữa hai điểm $A$, $B$ bằng $2x=2\cdot 2\sqrt{2}\approx 5{,}7$ m.
    }
\end{bt}

\begin{bt}%[2D1V5-1]
% nguon: Archive/SP_SauLuoi_2024-2025/SP_GHKI/2-TL-TN-DS-THPT-PhanChuTrinh-Binhthuan-GHKI-NH24-25.tex #225
%[Dự án đề kiểm tra Toán 12 GHKI NH24-25- Hieu Phan]%[THPT Phan Chu Trinh- Bình Thuận]
    
    \immini
    {
        Cho hàm số $y=\dfrac{ax-b}{x+c}$ với $a$, $b$, $c \in\mathbb{Z},\,\,ac+b\ne 0$ và có đồ thị như hình bên. Giá trị của $a+b-c$ bằng bao nhiêu?
    }
    {
        \begin{tikzpicture}[>=stealth,x=1cm,y=1cm,scale=0.4]
            \def\a{1}
            \def\b{-1}
            \def\c{1}
            \def\d{-2}
            \draw[->] (-4,0) -- (6,0) node[below] {\scriptsize $x$};
            \draw[->] (0,-3) -- (0,6) node[left] {\scriptsize $y$};
            \draw (0,0)node[below left]{\scriptsize $O$};
            \draw[blue] (2,-3)--(2,6) (-3,1)--(6,1); % Vẽ TCĐ và TCN
            \clip (-5,-5)rectangle(6,6);
            \pgfmathsetmacro{\can}{-(\d)/(\c)}
            \draw[samples=150,smooth,domain=-4:{\can-.25}] plot(\x,{(\a*\x+(\b))/(\c*\x+(\d))}); % Vẽ nhánh bên trái TCĐ
            \draw[samples=150,smooth,domain={\can+.1}:6] plot(\x,{(\a*\x+(\b))/(\c*\x+(\d))}); % Vẽ nhánh bên phải TCĐ
            \foreach \x in {1} \fill (\x,0) node[below]{\x} circle (1.2pt);
             \fill (0,1) node[above left]{$ 1 $} circle (1.2pt) (2,0) node[below right]{$ 2 $} circle (1.2pt);
        \end{tikzpicture}
        
    }
    \loigiai{
        Theo đề bài ta có\\
        Tiệm cận ngang $y=a=1$, tiệm cận đứng $x=-c=2\Rightarrow c=-2$.\\
        Giao điểm của đồ thị với trục $Ox$ là $\left(1;0\right)$ nên $\dfrac{b}{a}=1\Rightarrow b=1$.\\
        Vậy $a+b-c=1+1-(-2)=4$.}
\end{bt}

\begin{bt}%[2D1H4-1]
% nguon: Archive/SP_SauLuoi_2024-2025/SP_GHKI/2-TL-TN-DS-THPT-HungVuong-TpHCM-GHKI-NH24-25.tex #242
%[Dự án đề kiểm tra Toán 12 GHKI NH24-25- Phạm Hoài]%[THPT - Hùng Vương]
 Cho hàm số $y=\dfrac{3x^2-2x+6}{x+1}$ có đồ thị $(C)$. Phương trình đường tiệm cận xiên của đồ thị hàm số $(C)$ là $y=ax+b$. Biểu thức $P=a^2+b^2$ có giá trị là bao nhiêu?
\loigiai{
 hàm số $y=\dfrac{3x^2-2x+6}{x+1}$ có tập xác định $\mathscr{D}=\mathbb{R}\setminus\{-1\}$.\\
 Khi đó $y=3x-5-\dfrac{11}{x+1}$.\\
 Ta có 
 \begin{itemize}
 	\item $\lim\limits_{x\to +\infty}[y-(3x-5)]=\lim\limits_{x\to +\infty}\left(-\dfrac{11}{x+1}\right)=0$.
 	\item $\lim\limits_{x\to -\infty}[y-(3x-5)]=\lim\limits_{x\to -\infty}\left(-\dfrac{11}{x+1}\right)=0$.
 \end{itemize}
 Do đó phương trình đường tiệm cận xiên là $y=3x-5$, suy ra $a=3$, $b=-5$.\\
 Vậy $P=a^2+b^2=9+25=34$.	
	
}
\end{bt}

\begin{bt}%[2D1H4-4]
% nguon: Archive/SP_SauLuoi_2024-2025/SP_GHKI/2-TN-DS-TLN-TL-THPT-ChuyenLeKhiet-QuangNgai-GHKI-NH24-25.tex #342
%[KNTT - Lớp 12 - Kiểm tra giữa học kì 1 - THPT Liên Khiết-Quãng Ngãi]%[Ngô Tất Thành]
	Một công ty sản xuất đồ chơi $A$ phải chi $30\,000$ USD để thiết lập dây chuyền sản xuất ban đầu. Sau đó, cứ sản xuất được một sản phẩm đồ chơi $A$, công ty phải trả $4$ USD cho nguyên liệu thô và nhân công. Gọi $x \ (x \geq 1)$ là số đồ chơi mà công ty đã sản xuất và $T(x)$ (đơn vị USD) là tổng số tiền bao gồm cả chi phí ban đầu mà công ty phải trả khi sản xuất $x$ đồ chơi $A$. Người ta xác định chi phí trung bình cho mỗi sản phẩm đồ chơi $A$ là $M(x) = \dfrac{T(x)}{x}$. Khi $x$ đủ lớn $(x \to +\infty)$ thì chi phí trung bình cho mỗi sản phẩm đồ chơi $A$ gần nhất là bao nhiêu USD?
	\loigiai{ 
		Chi phí trung bình để sản xuất một sản phẩm $A$ là $M(x) = \dfrac{4x + 30\,000}{x}$. \\
		Ta có $\lim\limits_{x\to +\infty}M(x)=\lim\limits_{x\to +\infty}\dfrac{4x + 30\,000}{x}=4$. \\
		Vậy chi phí trung bình cho mỗi sản phẩm đồ chơi $A$ là gần $4$ USD.			
	}
\end{bt}

\begin{bt}%[2D1V3-6]
% nguon: Archive/SP_SauLuoi_2024-2025/SP_GHKI/2-TN-DS-TLN-TL-THPT-ChuyenLeKhiet-QuangNgai-GHKI-NH24-25.tex #343
%[KNTT - Lớp 12 - Kiểm tra giữa học kì 1 - THPT Liên Khiết-Quãng Ngãi]%[Ngô Tất Thành]
	Một người đang ở tại điểm $A$ trên sa mạc. Ông ta muốn đến điểm $B$ và cách $A$ một đoạn là $70$ km. Trong sa mạc thì xe ông ta chỉ có thể di chuyển với vận tốc $30$ km/h. Ông ấy phải đến được điểm $B$ trong $2$ giờ. Biết rằng có một con đường nhựa $HK$ song song với $AB$ và cách $AB$ một đoạn $10$ km. Trên đường nhựa này thì xe ông ấy có thể di chuyển với vận tốc $50$ km/h. Để đến $B$ sớm nhất (đảm bảo trong khung giờ cho phép) thì ông phải đi theo con đường nào?
	\begin{center}
		\begin{tikzpicture}
			\fill[yellow!50] (0,0) rectangle (8,2.5);
			\draw (0,0) rectangle (8,2.5);
			\node[above left] at (0,2.5) {$A$};
			\node[above right] at (8,2.5) {$B$};
			\node[below left] at (0,0) {$H$};
			\node[below right] at (8,0) {$K$};
			\draw[<->] (0,2.7) -- node[above] {70 km} (8,2.7);
			\draw[<->] (8.2,0) -- node[right] {10 km} (8.2,2.5);
			\draw[green, line width=2pt] (-1,0) -- (9,0); 
			\node at (4,1.25) {\textbf{Sa mạc}};
			\node[below,text=black] at (4,-0.6) {\textbf{Đường nhựa}};
			\draw[-latex,black] (2.7,-0.8) to[out=180, in=-45] (1,0);
		\end{tikzpicture}
	\end{center}
	\loigiai{
		\begin{center}
			\begin{tikzpicture}
				\fill[yellow!50] (0,0) rectangle (8,2.5); % Field with light yellow color
				\draw (0,0) rectangle (8,2.5);
				\node[above left] at (0,2.5) {$A$};
				\node[above right] at (8,2.5) {$B$};
				\node[below left] at (0,0) {$H$};
				\node[below right] at (8,0) {$K$};
				\draw[green, line width=2pt] (-1,0) -- (2,0)(6,0) -- (9,0); 
				\draw[red, line width=2pt] (0,2.5) -- (2,0) (8,2.5) -- (6,0) (2,0) -- (6,0);
				\node[above] at (0.7,0) {$x$};
				\node [below] at (2,0) {C};
				\node[above] at (7,0) {$y$};
				\node [below] at (6,0) {D};
				\draw[<->] (0,2.7) -- node[above] {70 km} (8,2.7);
				\draw[<->] (8.2,0) -- node[right] {10 km} (8.2,2.5);
				\node at (4,1.25) {Sa mạc};
				\node[below,text=black] at (4,-0.6) {\textbf{Đường nhựa}};
				\draw[-latex,black] (2.7,-0.8) to[out=180, in=-45] (1,0);
			\end{tikzpicture}
		\end{center}
		Thời gian nếu đi trực tiếp từ $A$ đến $B$ trên sa mạc là $\dfrac{70}{30} = \dfrac{7}{3} > 2$. \\ 
		Do đó, nhà địa chất học không thể đến đúng thời gian quy định. \\
		Vì vậy cần thiết phải chia quãng đường đi được thành 3 giai đoạn:
		$A \to C \to D \to B$.  \\
		Đặt $HC = x (0 < x < 70);\  DK = y (0 < y < 70);$ \\ 
		Thời gian đi từ $A \to C$ là $\dfrac{\sqrt{10^2 + x^2}}{30}$. \\ 
		Thời gian đi từ $C \to D$ là $\dfrac{70 - (x + y)}{50}$. \\ 
		Thời gian đi từ $D \to B$ là $\dfrac{\sqrt{10^2 + y^2}}{30}$. \\
		Tổng thời gian đi từ $A \to B$ theo cách này là 
		\begin{eqnarray*}
			& & \dfrac{\sqrt{10^2 + x^2}}{30} + \dfrac{70-(x+y)}{50} + \dfrac{\sqrt{10^2 + y^2}}{30}\\
			&= & \frac{\sqrt{10^2 + x^2}}{30} + \frac{35-x}{50} + \frac{\sqrt{10^2 + y^2}}{30} + \frac{35-y}{50} \\
			&= & f(x) + f(y).
		\end{eqnarray*}
		Xét $f(u) = \dfrac{\sqrt{10^2 + u^2}}{30} + \dfrac{35-u}{50}$, $0 < u < 70$. \\
		$\Rightarrow f'(u) = \dfrac{u}{30\sqrt{10^2 + u^2}} - \dfrac{1}{50}$; $f'(u) = 0 \Leftrightarrow u = \dfrac{15}{2}$. \\
		Lập bảng biến thiên ta được $\min\limits_{u\in(0;70)} f(u) = f\left( \dfrac{15}{2}\right)  = \dfrac{29}{30}$. \\
		Khi đó $f(x) + f(y) \geq \dfrac{29}{30} + \dfrac{29}{30} \approx 1,93$. \\
		Dấu \lq\lq$=$\rq\rq \, xảy ra khi $x = y = \dfrac{15}{2}$.\\
		Vậy để đến $B$ sớm nhất thì ông ta phải đi trên đoạn $AC=12{,}5$ km, đoạn $CD = 55$ km; $DB = 12{,}5$ km.}
\end{bt}

\begin{bt}%[2D1V3-6]
% nguon: Archive/2024/TOAN_10_XE/TOAN_10_XE/data4/10D6/DS10-C6-B18.tex #366

	Từ một tấm gỗ hình vuông cạnh $200$ cm. Cắt một tấm gỗ có hình tam giác vuông, có tổng của một cạnh góc vuông và cạnh huyền bằng $120$ cm và cạnh góc vuông còn lại bằng $\dfrac{\sqrt3}{5}$ cạnh của tấm gỗ hình vuông. Hỏi cạnh huyền của tấm gỗ hình tam giác bằng bao nhiêu?
	\loigiai{
		Kí hiệu cạnh góc vuông $AB = x,\, 0 < x < 60$.\\ 
		Khi đó cạnh huyền $BC = 120 - x$.\\
		Cạnh góc vuông kia là $AC = \sqrt {BC^2 - AB^2} = \sqrt{120^2 - 240x} $.\\ 
		Theo giả thiết ta có 	
		$$	\sqrt{120^2 - 240x}=\dfrac{\sqrt3}{5}\cdot 200 \Leftrightarrow x=40.$$
		Vậy cạnh huyền  $BC = 80$.
	}
\end{bt}

\begin{bt}%[2D1V3-6]
% nguon: Archive/2024/TOAN_10_XE/TOAN_10_XE/data4/10D6/DS10-C6-B18.tex #367
%[Phan Minh Quốc Vinh]
	\immini{Một ngọn hải đăng đạt tại vị trí $ A $ có khoảng cách đến bờ biển $ AB=5  $ km. Trên bờ biển có một cái kho ở vị trí $ C $ cách $ B $ một khoảng bằng $ BC= 7 $ km. Người canh hải đăng có thể chèo đò từ $ A $ đến vị trí $ M $ trên bờ biển với vận tốc $ 4 $ km/h rồi đi bộ đến $ C $ với vận tốc $ 6 $ km/h. Vị trí của điểm $ M $ cách $ B $ một khoảng bao nhiêu để thời gian người đó đi đến $C$ là $2$ giờ?(làm tròn đến hàng phần mười)
	}{
		\begin{tikzpicture}
			[scale=1,font=\footnotesize,line join=round,line cap=round,>=stealth]
			\tkzDefPoints{0/0/A, 0/-3/B, 5/-3/C, 2/-3/M}
			\tkzDrawSegments(A,B B,C A,M)
			\tkzLabelPoints[below left](B)
			\tkzLabelPoints[above left](A)
			\tkzLabelPoints[below](M,C)
			\tkzMarkRightAngle(A,B,M)
			\tkzDrawPoints(A,B,C,M)
		\end{tikzpicture}
	}
	\loigiai
	{Gọi khoảng cách từ $ M $ đến $ B $ là $ x $ km $ (0\le x\le 7) $.\\
		Khi đó, $ MC=7-x $ và $ AM=\sqrt{x^2+25} $.\\
		Người đó đi từ $ A $ đến $ C $ hết khoảng thời gian là $ \dfrac{7-x}{6}+\dfrac{\sqrt{x^2+25}}{4} $ giờ.\\
		Theo giả thiết ta có \allowdisplaybreaks
		\begin{eqnarray*} &&\dfrac{7-x}{6}+\dfrac{\sqrt{x^2+25}}{4}=2\Leftrightarrow 3\sqrt{x^2+5}=10+2x\\
			&\Leftrightarrow& 5x^2-40x-55=0\Leftrightarrow\hoac{&x=\dfrac{20+\sqrt{895}}{9}\\&x=\dfrac{20-\sqrt{895}}{9}\text{ (loại)}}
		\end{eqnarray*}
		Vậy $M$ cách $B$ một khoảng $5,5$ km.
	}
\end{bt}

\begin{bt}%[2D1V1-3]
% nguon: Archive/2024/San pham - To Latex Strong - Tex hoa Lop 12 moi - Dot 2 (Chuong 1 + 2 + 3)/San pham - To Latex Strong - Tex hoa Lop 12 moi - Dot 2 (Chuong 1 + 2 + 3)/Data/Chuong 1/Chan troi sang tao/To 27 - Dot 16 - Chuong 1 - Bai 4 - CTST - De 2.tex #839
%[Tổ 27 - Đợt 16 - Chương 1 - Bài 4 - CTST - Đề 2]%[Trần Thiên Đức]
Cho hàm số $y=-x^3-mx^2+(4m+9)x+5$, với $m$ là tham số. Có bao nhiêu giá trị nguyên của $m$ để hàm số nghịch biến trên $(-\infty;+\infty)$.
\loigiai{
Ta có\\
$y=-x^3-mx^2+(4m+9)x+5$.\\
$y'=-3x^2-2mx+(4m+9)$.\\
Hàm số nghịch biến trên $(-\infty;+\infty)$ khi $y'\leq0$, $\forall x\in(-\infty;+\infty)$.\\
$\Leftrightarrow\heva{
&a=-1<0\\
&\Delta'=m^2+12m+27\leq0.\\}$\\
$\Leftrightarrow-9\leq m\leq-3$.\\
Vì $m$ nguyên nên $m\in\{-9;-8;-7;-6;-5;-4;-3\}$.\\
Vậy có $7$ giá trị.
}
\end{bt}

\begin{bt}%[2D1V5-4]
% nguon: Archive/2024/San pham - To Latex Strong - Tex hoa Lop 12 moi - Dot 2 (Chuong 1 + 2 + 3)/San pham - To Latex Strong - Tex hoa Lop 12 moi - Dot 2 (Chuong 1 + 2 + 3)/Data/Chuong 1/Chan troi sang tao/To 27 - Dot 16 - Chuong 1 - Bai 4 - CTST - De 2.tex #840
%[Tổ 27 - Đợt 16 - Chương 1 - Bài 4 - CTST - Đề 2]%[Trần Thiên Đức]
Có bao nhiêu giá trị nguyên $m$ thuộc đoạn $[-5;5]$, để đường thẳng $d:y=mx+1$ cắt đồ thị hàm số $(C)\colon y=\dfrac{-2x+1}{x-1}$ tại hai điểm phân biệt?
\loigiai{
Điều kiện $x\neq1$. \\
Phương trình hoành độ giao điểm của $d$ và $(C)$ là \[\dfrac{-2x+1}{x-1}=mx+1\Leftrightarrow(mx+1)(x-1)=-2x+1
\Leftrightarrow mx^2+(3-m)x-2=0.\]
Đặt $g(x)=mx^2+(3-m)x-2=0$.\\
$d$ cắt $(C)$ tại hai điểm phân biệt khi phương trình $g(x)=0$ có hai nghiệm phân biệt $x_1$, $x_2$ khác $1$.\\
$\heva{
&a\neq0\\
&\Delta>0\\
&g(1)\neq0\\}\Leftrightarrow\heva{
&m\neq0\\
&(3-m)^2+8m>0\\
&m+3-m-2\neq0\\}\Leftrightarrow\heva{
&m\neq0\\
&m^2+2m+9>0\\}\Leftrightarrow m\neq0$.\\
Vì $\heva{
&m\in\mathbb{Z}\\
&m\in[-5;5]\\}$ nên $m\in\{-5;-4;-3;-2;-1;1;2;3;4;5\}$.\\
Vậy có $10$ giá trị.
}
\end{bt}

\begin{bt}%[2D1V3-2]
% nguon: Archive/2024/San pham - To Latex Strong - Tex hoa Lop 12 moi - Dot 2 (Chuong 1 + 2 + 3)/San pham - To Latex Strong - Tex hoa Lop 12 moi - Dot 2 (Chuong 1 + 2 + 3)/Data/Chuong 1/Chan troi sang tao/To 27 - Dot 16 - Chuong 1 - Bai 4 - CTST - De 2.tex #841
%[Tổ 27 - Đợt 16 - Chương 1 - Bài 4 - CTST - Đề 2]%[Trần Thiên Đức]
Gọi $m$ là giá trị nhỏ nhất của hàm số $y=\dfrac{x^2-2x+5}{x-1}$ trên khoảng $(1;+\infty)$. Giá trị của $m$ bằng bao nhiêu?
\loigiai{
Xét hàm số $y=\dfrac{x^2-2x+5}{x-1}$ trên khoảng $(1;\infty)$.\\
Ta có $y'=\dfrac{x^2-2x-3}{(x-1)^2}\Rightarrow y'=0\Leftrightarrow\hoac{
&x=-1\\
&x=3.\\}$\\
Bảng biến thiên\\
\begin{center}
\begin{tikzpicture}[>=stealth]
%Khai báo thông số các dòng
\tkzTabInit[nocadre=true,lgt=1.5,espcl=3,deltacl=0.5]{$x$/0.7 ,$h'(x)$/.7,$h(x)$/2.5}
%Dòng 2
{$-\infty$,$-1$,$1$,$3$,$+\infty$}
%Dòng 2 dấu
\tkzTabLine{ ,  ,$0$,,d,-,$0$,+, }
% Các vị trí
\path
%(N13)node[above](1){$-\infty$}
%(N22)node[below=.3cm](2){$\sqrt{2}$}
(N32)node[below,xshift=3mm](1){$+\infty$}
(N43)node[above](2){$4$}
(N52)node[below](3){$+\infty$}

;
\draw[double] (N32)--(N33);
\fill[pattern = north east lines,opacity=.6] (N33) rectangle (N11);
%Vẽ mũi tên
\foreach \x/\y in {1/2,2/3}
\draw[->] (\x)--(\y);
\end{tikzpicture}
\end{center}
Suy ra $m=\min\limits_{(1;+\infty)}y=4$ khi $x=3$.
}
\end{bt}

\begin{bt}%[2D1V5-8]
% nguon: Archive/2024/San pham - To Latex Strong - Tex hoa Lop 12 moi - Dot 2 (Chuong 1 + 2 + 3)/San pham - To Latex Strong - Tex hoa Lop 12 moi - Dot 2 (Chuong 1 + 2 + 3)/Data/Chuong 1/Chan troi sang tao/To 27 - Dot 16 - Chuong 1 - Bai 4 - CTST - De 2.tex #842
%[Tổ 27 - Đợt 16 - Chương 1 - Bài 4 - CTST - Đề 2]%[Trần Thiên Đức]
Một công ty chuyên sản xuất dụng cụ thể thao nhận được đơn đặt hàng sản xuất $8\,000$ quả bóng rổ. Công ty có một số máy, mỗi máy có khả năng sản xuất $30$ bóng rổ trong một giờ. Chi phí thiết lập mỗi máy là $200$ nghìn đồng. Sau khi thiết lập, quá trình sản xuất sẽ diễn ra hoàn toàn tự động và chỉ cần có người giám sát. Chi phí trả cho người giám sát là $192$ nghìn đồng mỗi giờ. Công ty cần sử dụng bao nhiêu máy để chi phí hoạt động đạt mức thấp nhất?
\loigiai{
Gọi số máy công ty sử dụng để sản xuất là $x$ ($x\in\mathbb{N}$, $x>0$).\\
Thời gian cần để sản xuất hết $8\,000$ quả bóng là: $\dfrac{8\,000}{30x}$.\\
Tổng chi phí để sản xuất là: $P(x)=200x+\dfrac{8\,000}{30x}\cdot192=200x+\dfrac{51\,200}{x}$.\\
Ta có $P'(x)=200-\dfrac{51\,200}{x^2}=0\Leftrightarrow x^2=256\Leftrightarrow\hoac{
&x=16\\
&x=-16~(\text{loại}).\\}$\\
\begin{center}
\begin{tikzpicture}
\tkzTabInit[nocadre=true,lgt=1.2,espcl=3]
{$x$/0.7,$P'(x)$/0.7,$P(x)$/2.5}{$0$,$16$,$+\infty$}
\tkzTabLine{,-,$0$,+,}
\tkzTabVar{+/{},-/$6\,400$,+/{}}
\end{tikzpicture}
\end{center}
Vậy công ty nên sử dụng $16$ máy để chi phí hoạt động là thấp nhất.
}
\end{bt}

\begin{bt}%[2D1H5-8]
% nguon: Archive/2024/San pham - To Latex Strong - Tex hoa Lop 12 moi - Dot 2 (Chuong 1 + 2 + 3)/San pham - To Latex Strong - Tex hoa Lop 12 moi - Dot 2 (Chuong 1 + 2 + 3)/Data/Chuong 1/Chan troi sang tao/To 27 - Dot 16 - Chuong 1 - Bai 4 - CTST - De 2.tex #843
%[Tổ 27 - Đợt 16 - Chương 1 - Bài 4 - CTST - Đề 2]%[Trần Thiên Đức]
Một vật chuyển động theo quy luật $s(t)=t^3-3t^2-9t+2$. Hỏi tại thời điểm nào thì vận tốc của vật đạt giá trị nhỏ nhất?
\loigiai{
\textbf{Cách 1:} Lấy đạo hàm $s'(t)$ ta có phương trình $v(t)=3t^2-6t-9$ có đồ thị là Parabol có đỉnh thấp nhất $I\heva{
&x_I=\dfrac{-b}{2a}=1\\
&y_I=v(1)=-12.\\}$\\
Vậy vận tốc đạt giá trị nhỏ nhất khi $t_{\min}=x_I=1$.\\
\textbf{Cách 2:} Nhập phương trình bậc hai vào máy tính 580VNX, ta được
$\heva{
&v(t)=y_{\min}=-12\\
&t_{\min}=x_{\min}=1.\\}$
}
\end{bt}

\begin{bt}%[2D1V5-8]
% nguon: Archive/2024/San pham - To Latex Strong - Tex hoa Lop 12 moi - Dot 2 (Chuong 1 + 2 + 3)/San pham - To Latex Strong - Tex hoa Lop 12 moi - Dot 2 (Chuong 1 + 2 + 3)/Data/Chuong 1/Chan troi sang tao/To 27 - Dot 16 - Chuong 1 - Bai 4 - CTST - De 2.tex #844
%[Tổ 27 - Đợt 16 - Chương 1 - Bài 4 - CTST - Đề 2]%[Trần Thiên Đức]
Trong phòng thí nghiệm người ta xác định được số lượng vi khuẩn được nuôi cấy tính theo công thức $N(t)=1\,000+\dfrac{100t}{100+t^2}$ (con vi khuẩn). Tính số lượng vi khuẩn lớn nhất kể từ khi nuôi cấy.
\loigiai{
Xét hàm số $N(t)=1\,000+\dfrac{100t}{100+t^2}$, $(t>0)$.\\
$N'(t)=\dfrac{100\left(100+t^2\right)-2t\cdot100t}{\left(100+t^2\right)^2}=\dfrac{100\left(100-t^2\right)}{\left(100+t^2\right)^2}$.\\
$N'(t)=0\Leftrightarrow100-t^2=0\Leftrightarrow\hoac{
&t=10~(\text{thỏa mãn})\\
&t=-10~(\text{loại}).\\}$\\
Ta có bảng biến thiên
\begin{center}
\begin{tikzpicture}
\tkzTabInit[nocadre=true,lgt=1.2,espcl=3]
{$t$/0.7,$N'(t)$/0.7,$N(t)$/2.5}{$0$,$10$,$+\infty$}
\tkzTabLine{,+,$0$,-,}
\tkzTabVar{-/$1\,000$,+/$1\,005$,-/$1\,000$}
\end{tikzpicture}
\end{center}
Vậy số lượng vi khuẩn lớn nhất nuôi cấy được là $1\,005$ con.
}
\end{bt}

\begin{bt}%[2D1V3-6]
% nguon: Archive/2024/0. [LVH] BG - 12 New - 4in1-20240625T224301Z-001/0. [LVH] BG - 12 New - 4in1/2. [LVH] TEX SGK 12 NEW-20240418T074833Z-001/2. [LVH] TEX SGK 12 NEW/[LVH] CTST-12new/CTST - CD12new/data/12-SGK-CTST-CD1-2.tex #1051

	Tại một xí nghiệp chuyên sản xuất vật liệu xây dựng, nếu trong một ngày xí nghiệp sản xuất $x$ (m$^3$) sản phẩm thì phải bỏ ra các khoản chi phí bao gồm: $4$ triệu đồng chi phí cố định; $0{,}2$ triệu đồng chi phí cho mỗi mét khối sản phẩm và $0{,}001 x^2$ triệu đồng chi phí bảo dưỡng máy móc. Biết rằng, mỗi ngày xí nghiệp sản xuất được tối đa $100$ m$^3$ sản phẩm.
	\begin{enumerate}
		\item Tính tổng chi phí (kí kiệu là $C$, đơn vị: triệu đồng) để xí nghiệp sản xuất $x$ (m$^3$) sản phẩm trong một ngày.
		\item Tính chi phí trung bình (kí hiệu là $\overline{C}$ ) trên mỗi mét khối sản phẩm.
		\item Tìm giá trị của $x$ sao cho chi phí trung bình $\overline{C}$ thấp nhất. Tìm giá trị thấp nhất đó.
	\end{enumerate}
	
	\loigiai{
		\begin{enumerate}
			\item Tổng chi phí (triệu đồng) để xí nghiệp sản xuất $x$ (m$^3$) sản phẩm trong một ngày là
			$$
			C=C(x)=4+0{,}2 x+0{,}001 x^2 \text { với } 0 \leq x \leq 100.
			$$
			\item Chi phí trung bình (triệu đồng) trên mỗi mét khối sản phẩm là
			$$
			\overline{C}=\overline{C}(x)=\dfrac{C(x)}{x}=\dfrac{4+0{,}2 x+0{,}001 x^2}{x}=0{,}001 x+\dfrac{4}{x}+0{,}2 \text { với } 0<x \leq 100.
			$$
			\item Ta có $\bar{C}'(x)=0{,}001-\dfrac{4}{x^2}$;
			$$
			\overline{C}'(x)=0 \Leftrightarrow 0{,}001-\dfrac{4}{x^2}=0 \Leftrightarrow x^2=4\,000 \Leftrightarrow x=20 \sqrt{10} \in(0 ; 100].
			$$
			
			Ta có $\overline{C}(20 \sqrt{10})=\dfrac{\sqrt{10}}{25}+\dfrac{1}{5} \approx 0,326$.\\
			Bảng biến thiên
			\begin{center}
				\begin{tikzpicture}
					\tikzset{double style/.append style={double distance=2pt}}
					\tkzTabInit[lgt=1.2, espcl=2]
					{$x$/0.6,$\overline{C'}(x)$/0.6,$\overline{C}(x)$/2.5}{$0$,$20\sqrt{10}$,$100$}
					\tkzTabLine{,-,0,+,}
					\tkzTabVar{+/$+\infty$,-/$\dfrac{\sqrt{10}}{25}+\dfrac{1}{5}$,+/$0{,}34$}
				\end{tikzpicture}
			\end{center}
			Từ bảng biến thiên, ta thấy chi phí trung bình thấp nhất là $\bar{C}(20 \sqrt{10}) \approx 0{,}326$ (triệu đồng/m$^3$ sản phẩm), đạt được khi $x=20 \sqrt{10} \approx 63$ (m$^3$).
		\end{enumerate}
	}
\end{bt}

\begin{bt}%[2D1V3-6]
% nguon: Archive/2024/0. [LVH] BG - 12 New - 4in1-20240625T224301Z-001/0. [LVH] BG - 12 New - 4in1/2. [LVH] TEX SGK 12 NEW-20240418T074833Z-001/2. [LVH] TEX SGK 12 NEW/[LVH] CTST-12new/CTST - CD12new/data/12-SGK-CTST-CD1-2.tex #1052

	  Nhà máy $A$ chuyên sản xuất một loại sản phẩm cung cấp cho nhà máy $B$. Hai nhà máy thoả thuận rằng, hằng tháng $A$ cung cấp cho $B$ số lượng sản phẩm theo đơn đặt hàng của $B$ (tối đa $100$ tấn sản phẩm). Nếu số lượng đặt hàng là $x$ tấn sản phẩm thì giá bán cho mỗi tấn sản phẩm là $P(x)=45-0{,}001 x^2$ (triệu đồng). Chi phí để $A$ sản xuất $x$ tấn sản phẩm trong một tháng là $C(x)=100+30 x$ (triệu đồng) (gồm $100$ triệu đồng chi phí cố định và $30$ triệu đồng cho mỗi tấn sản phẩm).
	\begin{enumerate}[a)]
		\item Nếu mỗi tháng $A$ bán $x$ tấn sản phẩm cho $B$ thì $A$ thu được lợi nhuận bao nhiêu?
		\item A bán cho $B$ bao nhiêu sản phẩm mỗi tháng thì thu được lợi nhuận lớn nhất? Tìm giá trị lợi nhuận lớn nhất đó.
	\end{enumerate}
	
	\loigiai{
		\begin{enumerate}
			\item Số tiền mà $A$ thu được (gọi là doanh thu) từ việc bán $x$ tấn sản phẩm $(0 \leq x \leq 100)$ cho $B$ là
			$$
			R(x)=x \cdot P(x)=x\left(45-0{,}001 x^2\right)=45 x-0{,}001 x^3 \text { (triệu đồng). }
			$$
			
			Lợi nhuận (triệu đồng) mà $A$ thu được là
			$$
			P(x)=R(x)-C(x)=x\left(45-0{,}001 x^2\right)-(100+30 x)=-0{,}001 x^3+15 x-100.
			$$
			\item Xét hàm số $P(x)=-0{,}001 x^3+15 x-100$ với $0 \leq x \leq 100$, ta có
			$$
			\begin{aligned}
				& P'(x)=-0{,}003 x^2+15; \\
				& P'(x)=0 \Leftrightarrow-0{,}003 x^2+15=0 \Leftrightarrow x^2=5\,000 \Leftrightarrow x=50 \sqrt{2} \in[0 ; 100].
			\end{aligned}
			$$
			
			Ta có $P(0)=-100$; $P(50 \sqrt{2})=500 \sqrt{2}-100 \approx 607$; $P(100)=400$.\\
			Bảng biến thiên
			\begin{center}
				\begin{tikzpicture}
					\tkzTabInit[lgt=1, espcl=4]
					{$x$/1,$y'$/0.6,$y$/3}{$0$,$50\sqrt{2}$,$100$}
					\tkzTabLine{,+,0,-,}
					\tkzTabVar{-/$100$,+/$500\sqrt{2}-100$,-/$400$}
				\end{tikzpicture}
			\end{center}
			
			Từ bảng biến thiên, ta có $\max \limits_{[0 ; 100]} P=P(50 \sqrt{2})=500 \sqrt{2}-100 \approx 607$.\\
			Vậy $A$ thu được lợi nhuận lớn nhất khi bán $50 \sqrt{2} \approx 70{,}7$ tấn sản phẩm cho $B$ mỗi tháng và lợi nhuận lớn nhất thu được khoảng $607$ triệu đồng.
		\end{enumerate}
	}
\end{bt}

\begin{bt}%[2D1V3-6]
% nguon: Archive/2024/0. [LVH] BG - 12 New - 4in1-20240625T224301Z-001/0. [LVH] BG - 12 New - 4in1/2. [LVH] TEX SGK 12 NEW-20240418T074833Z-001/2. [LVH] TEX SGK 12 NEW/[LVH] CTST-12new/CTST - CD12new/data/12-SGK-CTST-CD1-2.tex #1053

\immini{	Người ta muốn xây một đường cống thoát nước có mặt cắt ngang là hình tạo bởi một nửa hình tròn ghép với một hình chữ nhật (Hình 6). Biết rằng mặt cắt ngang có diện tích $2$ m$^2$. Các kích thước $x$, $y$ (đơn vị m) bằng bao nhiêu để chu vi của mặt cắt ngang là nhỏ nhất? Tính chu vi nhỏ nhất đó.}
{\begin{tikzpicture}[>=stealth,line join=round,line cap=round,font=\footnotesize,scale=0.65]
		\path 
		(0,0) coordinate (A)
		(4,0) coordinate (B)
		($(A)!1.2!90:(B)$) coordinate (D)
		($(B)!1.2!-90:(A)$) coordinate (C)
		($(C)!1/2!(D)$) coordinate (M)
		($(B)+(0.5,0)$) coordinate (u)
		($(C)+(0.5,0)$) coordinate (v)
		
		
		;
		\draw (D) arc(180:0:2);
		\draw(A)--(B)--(C)--(D)--(A);
		\draw[<->] (u)--(v)node[pos=0.5,right]{$y$};
		\draw (A)--(B)node[pos=0.5,below]{$x$};
		\draw (1.7,-0.5)node[below]{$\text{Hình 6}$ };
\end{tikzpicture}}
	\loigiai{
		Ta có diện tích mặt cắt ngang là \begin{eqnarray*}
			S&=&S_{hcn}+S_{\text{hình tròn}}\\
			&=&x\cdot y+\dfrac{1}{2}\cdot \pi \cdot R^2\\
			&=&x\cdot y+\dfrac{x^2}{8}\pi\\
			&=&2\, \text{m}^2.
		\end{eqnarray*} Suy ra $y=\dfrac{2}{x}-\dfrac{\pi}{8}x$.\\
		Chu vi của mặt cắt ngang là $C=2x+2y+\pi\dfrac{x}{2}=2x+2\left(\dfrac{2}{x}-\dfrac{\pi}{8}x\right)+\pi\dfrac{x}{2}=2x+\dfrac{4}{x}+\dfrac{\pi}{4}x.$\\
		$C'=2+\dfrac{\pi}{4}-\dfrac{4}{x^2}$.\\
		$C'=0\Leftrightarrow x=4\dfrac{\sqrt{\pi+8}}{\pi+8}$.\\
		
		Bảng biến thiên của hàm số đã cho như sau
		\begin{center}
			\begin{tikzpicture}
				\tkzTabInit[]{$x$/1,$C'$/1,$C$/2}{$0$, $4\dfrac{\sqrt{\pi+8}}{\pi+8}$, $+\infty$}
				\tikzset{double style/.append style={double distance=2.25pt}}
				\tkzTabLine{,-,0,+,}
				\tkzTabVar{+/$+\infty$,-/$6{,}68$,+/$+\infty$}
			\end{tikzpicture}
		\end{center}
		Chu vi của mặt cắt ngang nhỏ nhất bẳng $6{,}68$ khi $x=4\dfrac{\sqrt{\pi+8}}{\pi+8}\approx 1{,}2$ m;  $y\approx 6{,}68$ m.\\
	}
\end{bt}

\begin{bt}%[2D1V3-6]
% nguon: Archive/2024/0. [LVH] BG - 12 New - 4in1-20240625T224301Z-001/0. [LVH] BG - 12 New - 4in1/2. [LVH] TEX SGK 12 NEW-20240418T074833Z-001/2. [LVH] TEX SGK 12 NEW/[LVH] CTST-12new/CTST - CD12new/data/12-SGK-CTST-CD1-2.tex #1054

	\immini{Người ta muốn xây một bể bơi có dạng hình hộp chữ nhật, thể tích $1\,800$ m$^3$ và chiều sâu $2$ m (Hình 7). Biết rằng chi phí xây mỗi đơn vị diện tích của đáy bể gấp hai lần so với thành bể. Cần chọn chiều dài và chiều rộng của bể bằng bao nhiêu để tiết kiệm chi phí xây dựng bể nhất?}
	{\begin{tikzpicture}[>=stealth,line join=round,line cap=round,font=\footnotesize,scale=0.65]
			\path 
			(0,0) coordinate (A)
			(-1.3,-2) coordinate (B)
			(5,0) coordinate (D)
			($(B)+(D)-(A)$) coordinate (C)
			($(A)+(0,0.7)$) coordinate (A')
			($(A')+(B)-(A)$) coordinate (B')
			($(A')+(C)-(A)$) coordinate (C')
			($(A')+(D)-(A)$) coordinate (D')
			
			
			;
			\draw (B')--(A')--(D')--(C')--(B')--(B)--(C)--(C')(D')--(D)--(C) ;
			\draw[dashed](B)--(A)--(A')(D)--(A);
			\draw (B)--(C)node[pos=0.5,below]{$x$};
			\draw (C)--(D)node[pos=0.5,right]{$y$};
			\draw (D')--(D)node[pos=0.5,right]{$2$};
			\draw (2.2,-2.7)node[below]{$\text{Hình 7}$ };
	\end{tikzpicture}}
	\loigiai{
		Thể tích của bể là $V=B\cdot h$ nên $2xy=1\,800$, do đó $xy=900$.\\
		Khi đó diện tích thành bể cần xây là $2\cdot (2x)x+2\cdot (2y)$; diện tích đáy bể là $xy$.\\
		Tổng diện tích của bể là $4x+4y+xy=4x+4\cdot \dfrac{900}{x}+900$.\\
		Vì chi phí xây mỗi đơn vị diện tích của đáy bể gấp hai lần so với thành bể nên chi phí cần có là $4x+4\cdot \dfrac{900}{x}+2\cdot 900$.\\
		Đặt $f(x)=4x+4\cdot \dfrac{900}{x}+1\,800$.\\
		$f'(x)=4-\dfrac{3\,600}{x^2}$;\\
		 $f'(x)=0\Leftrightarrow x=30$ (do $x>0$).\\
		 
		Bảng biến thiên 
		\begin{center}
			\begin{tikzpicture}
				\tkzTabInit[nocadre=false, lgt=1.2, espcl=2.4]{$x$ /0.7,$f'(x)$ /0.7,$f(x)$ /2}{$0$,$30$,$+\infty$}
				\tkzTabLine{,-,$0$,+,}
				\tkzTabVar{+/  ,-/$2\,040$,+/$+\infty$}
			\end{tikzpicture}
		\end{center}
		Với $x=30$m thì $y=30$ m và thì chi phí xây dựng bể là thấp nhất.\\ 
		
		
	}
\end{bt}

\begin{bt}%[2D1C3-6]
% nguon: Archive/2024/0. [LVH] BG - 12 New - 4in1-20240625T224301Z-001/0. [LVH] BG - 12 New - 4in1/2. [LVH] TEX SGK 12 NEW-20240418T074833Z-001/2. [LVH] TEX SGK 12 NEW/[LVH] CTST-12new/CTST - CD12new/data/12-SGK-CTST-CD1-2.tex #1055

	\immini{Người ta muốn thiết kế một lồng nuôi cá có bề mặt hình chữ nhật bao gồm phần mặt nước có diện tích bằng $54$ m$^2$ và phần đường đi xung quanh với kích thước (đơn vị: m) như Hình $8$. Bề mặt của lồng có chiều dài và chiều rộng bằng bao nhiêu để diện tích phần đường đi là bé nhất?}
{\begin{tikzpicture}[>=stealth,line join=round,line cap=round,font=\footnotesize,scale=0.75]
		\path 
		(0,0) coordinate (A)
		(5,0) coordinate (B)
		(0,4) coordinate (D)
		($(D)+(B)-(A)$) coordinate (C)
		($(B)+(-0.5,0)$) coordinate (M)
		($(B)+(0,0.5)$) coordinate (N)
		($(A)+(1,0)$) coordinate (H)
		($(A)+(0,0.5)$) coordinate (T)
		($(C)+(-0.5,0)$) coordinate (Q)
		($(C)+(0,-0.5)$) coordinate (P)
		($(D)+(1,0)$) coordinate (R)
		($(D)+(0,-0.5)$) coordinate (S)
		($(T)+(H)-(A)$) coordinate (A')
		($(M)+(N)-(B)$) coordinate (B')
		($(Q)+(P)-(C)$) coordinate (C')
		($(R)+(S)-(D)$) coordinate (D')
		($(D)+(0,0.5)$) coordinate (x)
		($(C)+(0,0.5)$) coordinate (y)
		($(D)+(-0.5,0)$) coordinate (u)
		($(A)+(-0.5,0)$) coordinate (v)
		($(A)+(0,-0.5)$) coordinate (x')
		($(H)+(0,-0.5)$) coordinate (y')
		($(M)+(0,-0.5)$) coordinate (u')
		($(B)+(0,-0.5)$) coordinate (v')
		($(C)+(0.5,0)$) coordinate (x'')
		($(P)+(0.5,0)$) coordinate (y'')
		($(N)+(0.5,0)$) coordinate (u'')
		($(B)+(0.5,0)$) coordinate (v'')
		;
		\draw[fill=cyan!20!brown](A)--(B)--(C)--(D)--(A);
		\draw[fill=cyan!90!blue](A')--(B')--(C')--(D')--(A');
		\draw(D)--(x) (C)--(y)(D)--(u)(A)--(v)(A)--(x')(H)--(y')(M)--(u')(B)--(v')(C)--(x'') (P)--(y'') (N)--(u'') (B)--(v'');
		\draw[<->] (x)--(y)node[pos=0.5,above]{$a$};
		\draw[<->] (u)--(v)node[pos=0.5,left]{$b$};
		\draw[<->] (x')--(y')node[pos=0.5,below]{$2$};
		\draw[<->] (u')--(v')node[pos=0.5,below]{$1$};
		\draw[<->] (u'')--(v'')node[pos=0.5,right]{$1$};
		\draw[<->] (x'')--(y'')node[pos=0.5,right]{$1$};
		\draw (2.2,-0.5)node[below]{$\text{Hình 8}$ };
\end{tikzpicture}}
	\loigiai{
		Gọi $x$, $y$ lần lượt là độ dài, rộng của mặt nước. Điều kiện $x$, $y>0$.\\
		Phần mặt nước có diện tích bằng $54$ m$^2$ nên ta có $$x\cdot y=54.\, \hfill(*)$$
		Theo đề bài ta có $x=a-3$, $y=b-2$.\\
		Từ $(*)$ suy ra $$(a-3)(b-2)=54\Rightarrow b=\dfrac{54}{a-3}+2=\dfrac{2a+48}{a-3}$$
		Diện tích lối đi là \begin{eqnarray*}
			S(a)&=&a\cdot b-x\cdot y\\
			&=&ab-54\\
			&=&a\cdot \dfrac{2a+48}{a-3}-54\\
			&=&\dfrac{2a^2+48a}{a-3}-54.
		\end{eqnarray*}
		$S'(a)=\dfrac{2a^2-12a-144}{\left(a-3\right)^2}$;\\
		$S'(a)=0\Leftrightarrow \hoac{&a=-6\\&a=12.}$\\
		Bảng biến thiên
		\begin{center}
			\begin{tikzpicture}
				\tkzTabInit[nocadre=false, lgt=1.2, espcl=2.4]{$a$ /0.7,$S'(a)$ /0.7,$S(a)$ /2.5}{$0$,$3$,$12$,$+\infty$}
				\tkzTabLine{,-,d,-,$0$,+,}
				\tkzTabVar{+/$-54$  ,-D+/$-\infty$/$+\infty$,-/$42$,+/$+\infty$}
			\end{tikzpicture}
		\end{center}
		Vậy diện tích đường đi bé nhất khi $a=12$ m, $b=8$ m.
	}
\end{bt}

\begin{bt}%[2D1V3-6]
% nguon: Archive/2024/0. [LVH] BG - 12 New - 4in1-20240625T224301Z-001/0. [LVH] BG - 12 New - 4in1/2. [LVH] TEX SGK 12 NEW-20240418T074833Z-001/2. [LVH] TEX SGK 12 NEW/[LVH] CTST-12new/CTST - CD12new/data/12-SGK-CTST-CD1-2.tex #1056

	\immini{Một giếng dầu ngoài khơi được đặt ở vị trí $A$ cách bờ biển $3$ km, $B$ là vị trí trên bờ biển gần giếng dầu nhất. Nhà máy lọc dầu được đặt ở vị trí $C$ trên bờ biển, cách vị trí $B$ một khoảng $4$ km (Hình 9). Người ta dự định lắp đặt đường ống dẫn dầu gồm hai đoạn thẳng $AD$ và $DC$ ($D$ là một vị trí nằm giữa $B$ và $C$). Biết rằng mỗi mét đường ống đặt dưới biển có chi phí lắp đặt cao gấp đôi so với mỗi mét đường ống đặt trên bờ. Vị trí của $D$ như thế nào để giảm thiểu chi phí lắp đặt nhất?}
		{\begin{tikzpicture}[>=stealth,line join=round,line cap=round,font=\footnotesize,scale=0.75]
			\path 
			(0,0) coordinate (M)
			(7,0) coordinate (N)
			(0,1) coordinate (Q)
			($(N)+(Q)-(M)$) coordinate (P)
			(0,6) coordinate (S)
			($(S)+(P)-(Q)$) coordinate (R)
			($(Q)!0.3!(P)$) coordinate (B)
			($(Q)!0.45!(P)$) coordinate (D)
			($(Q)!0.9!(P)$) coordinate (C)
			($(B)!0.8!90:(C)$) coordinate (A)
			($(A)+(-0.5,0)$) coordinate (u)
			($(B)+(-0.5,0)$) coordinate (v)
			($(B)+(0,-0.4)$) coordinate (x)
			($(C)+(0,-0.4)$) coordinate (y)
			;
			\draw[fill=cyan!20!brown](M)--(N)--(P)--(Q)--(M);
			\draw[fill=cyan!90!blue](Q)--(P)--(R)--(S)--(Q);
			\draw[dashed](A)--(B)--(D);
			\draw (A)--(D)--(C) (A)--(u) (B)--(v) (B)--(x) (C)--(y); 
			\draw[<->] (x)--(y)node[pos=0.5,below]{$4$ km};
			\draw[<->] (u)--(v)node[pos=0.5,left]{$3$ km};
			\foreach \x/\y in {A/90,B/210,D/60,C/-30}
			\draw[fill=black] (\x) circle (1.1pt) + (\y:0.5cm) node{$\x$};
			\draw (3.2,-0.5)node[below]{$\text{Hình 9}$ };
	\end{tikzpicture}}
	\loigiai{
		Gọi $x$ là độ dài đoạn $BD$, điều kiện $x>0$.\\
		Ta có $AD=\sqrt{x^2+3^2}=\sqrt{x^2+9}$ (km);
		$  DC=4-x$ (km) với $0 \leq x \leq 4$ (km).\\
		Chi phí tiêu tốn là $2AD+DC=2\sqrt{x^2+9}+4-x$.\\ 
		
		Xét hàm số $y=2 \sqrt{x^2+9}+4-x$ với $0 \leq x \leq 4$, ta có
		$$
		\begin{aligned}
			& y'=\dfrac{2 x}{\sqrt{x^2+9}}-1;\\
			& y'=0 \Leftrightarrow \dfrac{2 x}{\sqrt{x^2+9}}-1=0 \Leftrightarrow 2 x=\sqrt{x^2+9} \Leftrightarrow 4 x^2=x^2+9 \\
			& \Leftrightarrow x^2=3 \Leftrightarrow x= \sqrt{3} \in[0 ; 4].
		\end{aligned}
		$$
		
		Ta có $y(0)=10$; $y(\sqrt{3})=3\sqrt{3}+4 \approx 9{,}20$; $y(4)=10$.\\
		Vậy $\min \limits_{[0,4]} y=y(\sqrt{3}) \approx 9{,}20$.
	}
\end{bt}

\begin{bt}%[2D1V3-6]
% nguon: Archive/2024/0. [LVH] BG - 12 New - 4in1-20240625T224301Z-001/0. [LVH] BG - 12 New - 4in1/2. [LVH] TEX SGK 12 NEW-20240418T074833Z-001/2. [LVH] TEX SGK 12 NEW/[LVH] CTST-12new/CTST - CD12new/data/12-SGK-CTST-CD1-2.tex #1057

	Tại một xí nghiệp, nếu trong một tuần xí nghiệp sản xuất $x$ nghìn sản phẩm thì chi phí sản xuất gồm: $10$ triệu đồng chi phí cố định, $3$ triệu đồng cho mỗi nghìn sản phẩm và $0{,}001 x^2$ triệu đồng chi phí bảo dưỡng thiết bị.
	\begin{enumerate}[a)]
		\item Tính chi phí trung bình trên mỗi nghìn sản phẩm theo $x$.
		\item Mỗi tuần xí nghiệp cần sản xuất bao nhiêu sản phẩm để chi phí trung bình thấp nhất?
	\end{enumerate}
	\loigiai{
		\begin{enumerate}[a)]
			\item 
			Chi phí sản xuất $x$ nghìn sản phẩm là 	$$
			C(x)=10+3 x+0{,}001 x^2 \text { (triệu đồng). }
			$$
			Chi phí sản xuất trung bình là $\overline{C}(x)=\dfrac{C(x)}{x}=\dfrac{10}{x}+\dfrac{3x}{x}+\dfrac{0{,}001x^2}{x}=\dfrac{10}{x}+3+0{,}001x$.\\
				Chi phí trung bình để sản xuất một nghìn sản phẩm là
				$$
				\overline{C}(1)=\dfrac{10}{1}+3+0{,}001\cdot 1^2=13{,}001\text { (triệu đồng). }
				$$
				\item $$
				\overline{C}(x)=\dfrac{10}{x}+3+0{,}001x \text { triệu đồng). }
				$$ Ta có 
				$$
				\overline{C}'(x)=-\dfrac{10}{x^2} +0{,}001
				$$
				$\overline{C}'(x)=0\Leftrightarrow -\dfrac{10}{x^2} +0{,}001 =0\Leftrightarrow x=100$.\\
				Bảng biến thiên
				\begin{center}
					\begin{tikzpicture}
						\tikzset{double style/.append style={double distance=2pt}}
						\tkzTabInit[lgt=1.3, espcl=2]
						{$x$/0.6,$\overline{C}'{x}$/0.6,$\overline{C}(x)$/1.5}{$0$,$100$,$+\infty$}
						\tkzTabLine{,-,0,+,}
						\tkzTabVar{+/,-/$3{,}2$,+/$+\infty$}
					\end{tikzpicture}
				\end{center}
				Vậy mỗi tháng xưởng sản xuất $3{,2}$ nghìn sản phẩm thì chi phí trung bình  sản xuất một sản phẩm thấp nhất.
			\end{enumerate}
		}
	\end{bt}

\begin{bt}%[2D1V3-6]
% nguon: Archive/2024/0. [LVH] BG - 12 New - 4in1-20240625T224301Z-001/0. [LVH] BG - 12 New - 4in1/2. [LVH] TEX SGK 12 NEW-20240418T074833Z-001/2. [LVH] TEX SGK 12 NEW/[LVH] CD-12new/SCD-12-CD/data/12-SCD-CD-2-2.tex #1058
%[ĐẠI DỰ ÁN TEX HOÁ SGK + BÀI GIẢNG LỚP 12]%[12-SCD-CD-2-2 - Thành Đức Trung]
	Độ giảm huyết áp của một bệnh nhân được cho bởi công thức $G(x)=0{,}025x^2(30-x)$, trong đó $x$ là liều lượng thuốc được tiêm cho bệnh nhân ($x$ được tính bằng miligam). Liều lượng thuốc cần tiêm cho bệnh nhân là bao nhiêu để huyết áp được giảm nhanh nhất?
	\begin{flushright}
		\textit{(Nguồn: Giải tích $12$ nâng cao, NXBGD Việt Nam, $2020$).}
	\end{flushright}
	\loigiai
	{ 
		Bài toán trở thành: Tìm $x\in[0;30]$ để hàm số $G(x)=0{,}025x^2(30-x)$ đạt giá trị lớn nhất. \\
		Ta có $G(x)=0{,}025\left(30x^2-x^3\right) \Rightarrow G'(x)=0{,}025\left(60x-3x^2\right)$. \\
		Xét $G'(x)=0 \Leftrightarrow \hoac{ & x=0 \\ & x=20.}$ \\
		Bảng biến thiên hàm số $G(x)$
		\begin{center}
			\begin{tikzpicture}[scale=1]
				\tkzTabInit[nocadre=false, lgt=1.2, espcl=3.5, deltacl=0.6]{$x$/0.6, $G'(x)$/0.6, $G(x)$/2}{$0$, $20$, $30$}
				\tkzTabLine{0,+,0,-,}
				\tkzTabVar{-/$0$, +/$100$, -/$0$}
			\end{tikzpicture}
		\end{center}
		Từ bảng biến thiên ta có $\max\limits_{[0;30]} G(x)=G(20)=100$. \\
		Vậy liều lượng thuốc cần tiêm cho bệnh nhân để huyết áp giảm nhanh nhất là $20$ mg.
	}
\end{bt}

\begin{bt}%[2D1V3-6]
% nguon: Archive/2024/0. [LVH] BG - 12 New - 4in1-20240625T224301Z-001/0. [LVH] BG - 12 New - 4in1/2. [LVH] TEX SGK 12 NEW-20240418T074833Z-001/2. [LVH] TEX SGK 12 NEW/[LVH] CD-12new/SCD-12-CD/data/12-SCD-CD-2-2.tex #1059
%[ĐẠI DỰ ÁN TEX HOÁ SGK + BÀI GIẢNG LỚP 12]%[12-SCD-CD-2-2 - Thành Đức Trung]
	Trong $25$ ngày đầu tiên sau khi phát hiện một bệnh dịch, các chuyên gia y tế ước tính số người nhiễm bệnh (kể từ ngày xuất hiện bệnh nhân thứ nhất) đến ngày thứ $t$ là $f(t)=45t^2-t^3$ với $t\in\{0;1;2;\ldots;25\}$.
	\begin{enumerate}
		\item Số người nhiễm bệnh (kể từ ngày xuất hiện bệnh nhân đầu tiên) lớn nhất là bao nhiêu?
		\item Nếu coi $f(t)$ là hàm số xác định trên đoạn $[0;25]$ thì $f'(t)$ được xem là tốc độ truyền bệnh (người/ngày) tại thời điểm $t$. Xác định ngày mà tốc độ truyền bệnh là lớn nhất và tính tốc độ đó.
	\end{enumerate}
	\begin{flushright}
		\textit{(Nguồn: Giải tích $12$ nâng cao, NXBGD Việt Nam, $2020$).}
	\end{flushright}
	\loigiai
	{
		\begin{enumerate}
			\item
			Bài toán trở thành: Tìm giá trị lớn nhất của hàm số $f(t)=45t^2-t^3$ trên đoạn $[0;25]$. \\
			Ta có $f'(t)=90t-3t^2$. \\
			Xét $f'(t)=0 \Leftrightarrow \hoac{ & t=0\in[0;25] \\ & t=30\notin[0;25].}$ \\
			%Suy ra $f'(t)\geqslant 0$ trên đoạn $[0;25]$.\\
			Bảng biến thiên của hàm số $f(t)$
			\begin{center}
				\begin{tikzpicture}[scale=1]
					\tkzTabInit[nocadre=false, lgt=1.2, espcl=5, deltacl=0.6]{$t$/0.6, $f'(t)$/0.6, $f(t)$/2}{$0$, $25$}
					\tkzTabLine{0,+,}
					\tkzTabVar{-/$0$, +/$12500$}
				\end{tikzpicture}
			\end{center}
			Từ bảng biến thiên, ta có $\max\limits_{[0;25]} f(t)=f(25)=12500$.
			Vậy sau $25$ ngày kể từ ngày xuất hiện bệnh nhân đầu tiên, số người nhiễm bệnh lớn nhất là $12500$ người.
			\item
			Bài toán trở thành: Tìm giá trị lớn nhất của hàm số $g(t)=f'(t)=90t-3t^2$ trên đoạn $[0;25]$. \\
			Ta có $g'(t)=90-6t$. \\
			Xét $g'(t)=0 \Leftrightarrow t=15$. \\
			Bảng biến thiên của hàm số $g(t)$
			\begin{center}
				\begin{tikzpicture}[scale=1]
					\tkzTabInit[nocadre=false, lgt=1.2, espcl=3.5, deltacl=0.8]{$t$/0.6, $g'(t)$/0.6, $g(t)$/2}{$0$, $15$, $25$}
					\tkzTabLine{,+,0,-,}
					\tkzTabVar{-/$0$, +/$675$, -/$375$}
				\end{tikzpicture}
			\end{center}
			Từ bảng biến thiên, ta có $\max\limits_{[0;25]} g(t)=g(15)=675$.
			Vậy tốc độ truyền bệnh lớn nhất là $675$ người/ngày và xảy ra tại ngày thứ $15$.
		\end{enumerate}
	}
\end{bt}

\begin{bt}%[2D1V3-6]
% nguon: Archive/2024/0. [LVH] BG - 12 New - 4in1-20240625T224301Z-001/0. [LVH] BG - 12 New - 4in1/2. [LVH] TEX SGK 12 NEW-20240418T074833Z-001/2. [LVH] TEX SGK 12 NEW/[LVH] CD-12new/SCD-12-CD/data/12-SCD-CD-2-2.tex #1060
%[ĐẠI DỰ ÁN TEX HOÁ SGK + BÀI GIẢNG LỚP 12]%[12-SCD-CD-2-2 - Thành Đức Trung]
	\immini
	{
		Người ta thiết kế một đường dây tải điện nối từ trạm biến áp ở vị trí $A$ đến trạm biến áp ở vị trí $C$. Do điều kiện thực tế nên phương án thiết kế được đưa ra như sau
		\begin{itemize}
			\item Chọn điểm $B$ sao cho $\widehat{ABC}=90^\circ$.
			\item Chọn một ví trí là điểm $S$ nằm giữa $A$ và $B$ để mắc đường dây tải điện đi từ $A$ đến $S$, rồi từ $S$ đến $C$ như hình vẽ.
		\end{itemize}
	}
	{
		\begin{tikzpicture}[scale=1, font=\footnotesize, line join=round, line cap=round, >=stealth]
			\path
			(0,0)coordinate(B)
			(0,1.5)coordinate(C)
			(1.5,0)coordinate(S)
			(4.5,0)coordinate(A)
			($(B)+(-0.5,0)$)coordinate(B')
			($(C)+(-0.5,0)$)coordinate(C')
			($(A)+(0,-0.5)$)coordinate(A'')
			($(B)+(0,-0.5)$)coordinate(B'')
			($(B')!0.5!(C')$)coordinate(X)node[above, rotate=90]{$1$ km}
			($(B'')!0.5!(A'')$)coordinate(Y)node[below]{$4$ km}
			;
			\draw(S)--(C)--(B)--(A);
			\draw[<->](B'')--(A'');
			\draw[<->](B')--(C');
			\foreach \p/\q in {S/-90, B/-90, A/-90, C/90} \fill[black] (\p) circle (1pt) ($(\p)+(\q:2.5mm)$) node{$\p$};
		\end{tikzpicture}
	}
	Biết rằng $AB=4$ km, $BC=1$ km, chi phí lắp đặt mỗi kilomet dây điện trên quãng đường $AS$, $SC$ lần lượt là $30$ triệu đồng và $50$ triệu đồng. Tính độ dài đoạn thẳng $AS$ để chi phí mắc đường dây điện là nhỏ nhất.
	\loigiai
	{
		Giả sử $AS=x$ (km) ($0<x<4$). \\
		Suy ra $BS=4-x$, $SC=\sqrt{1+\left(4-x\right)^2}$. \\
		Tổng chi phí để mắc đường dây tải điện $ASC$ là
		\begin{center}
			$T=30x+50\sqrt{1+\left(4-x\right)^2}$ (triệu đồng).\\
		\end{center}
		Xét hàm số $f(x)=30x+50\sqrt{1+\left(4-x\right)^2}$ với $x\in(0;4)$. \\
		Ta có $f'(x)=30-\dfrac{50(4-x)}{\sqrt{1+(4-x)^2}}$. \\
		Xét $f'(x)=0\Leftrightarrow 3\sqrt{1+(4-x)^2}=5(4-x)\Leftrightarrow (4-x)^2=\dfrac{9}{16} \Leftrightarrow x=\dfrac{13}{4}\text { (do } 0<x<4)$. \\
		Bảng biến thiên của hàm số $f(x)$
		\begin{center}
			\begin{tikzpicture}[scale=1]
				\tkzTabInit[nocadre=false, lgt=1.2, espcl=3.5, deltacl=0.9]{$x$/1.2, $f'(x)$/0.6, $f(x)$/2}{$0$, $\dfrac{13}{4}$, $4$}
				\tkzTabLine{,-,0,+,}
				\tkzTabVar{+/$50\sqrt{17}$, -/$160$, +/$170$}
			\end{tikzpicture}
		\end{center}
		Từ bảng biến thiên, ta có $\min\limits_{(0;4)} f(x)=f\left(\dfrac{13}{4}\right)=160$. \\
		Vậy độ dài đoạn thẳng $AS$ là $\dfrac{13}{4}$ km thì chi phí mắc đường dây điện là nhỏ nhất.
	}
\end{bt}

\begin{bt}%[2D1V3-6]
% nguon: Archive/2024/0. [LVH] BG - 12 New - 4in1-20240625T224301Z-001/0. [LVH] BG - 12 New - 4in1/2. [LVH] TEX SGK 12 NEW-20240418T074833Z-001/2. [LVH] TEX SGK 12 NEW/[LVH] CD-12new/SCD-12-CD/data/12-SCD-CD-2-2.tex #1061
%[ĐẠI DỰ ÁN TEX HOÁ SGK + BÀI GIẢNG LỚP 12]%[12-SCD-CD-2-2 - Thành Đức Trung]
	Một nhà máy cần sản xuất một bể nước không nắp bằng tôn có dạng hình hộp chữ nhật với đáy có chiều dài gấp hai lần chiều rộng và thể tích là $\dfrac{4}{3}$ m$^3$. Tính chiều rộng của đáy hình hộp đứng đó sao cho số tôn cần sử dụng là nhỏ nhất.
	\loigiai
	{
		\immini
		{
			Gọi chiều rộng của bể nước là $x$ (m) $(x>0)$. \\
			Khi đó, chiều dài của bể nước là $2x$ (m). \\
			Gọi chiều cao của bể nước là $h$ (m). \\
			Ta có thể tích bể nước là $V=x\cdot2x\cdot h=2x^2h=\dfrac{4}{3} \Rightarrow h=\dfrac{2}{3x^2}$. \\
			Số tôn cần sử dụng để sản xuất bể nước là 
			$$\begin{aligned}
				S
				&=S_{ABCD}+2S_{CDD'C'}+2S_{ADD'A'} \\
				&=2x^2+2xh+2\cdot2x\cdot h=2x^2+6xh=2x^2+\dfrac{4}{x} \left(\text{m}^2\right). 
			\end{aligned}$$
		}
		{
			\begin{tikzpicture}[scale=1, font=\footnotesize, line join=round, line cap=round, >=stealth]
				\def\a{1};
				\def\h{3};
				\path
				(0,0)coordinate(A)
				(-\a,-\a)coordinate(B)
				(\h,0)coordinate(D)
				($(B)+(D)-(A)$)coordinate(C)
				($(A)+(0,2)$)coordinate(A')
				($(A')+(B)-(A)$)coordinate(B')
				($(A')+(C)-(A)$)coordinate(C')
				($(A')+(D)-(A)$)coordinate(D')
				($(B)!0.5!(C)$)node[below]{$2x$}
				($(D)!0.5!(C)$)node[right]{$x$}
				($(D)!0.5!(D')$)node[right]{$h$}
				;
				\draw
				(B)--(B')
				(C)--(C')
				(D)--(D')
				(B)--(C)
				(C)--(D)
				(A')--(B')
				(B')--(C')
				(C')--(D')
				(D')--(A')
				;
				\draw[dashed]
				(B)--(A)
				(A)--(D)
				(A)--(A')
				;
				\foreach \p/\q in {A/150, B/-90, C/-90, D/0, A'/90, B'/-180, C'/90, D'/90} \fill[black] (\p) circle (1pt) ($(\p)+(\q:2.5mm)$) node{$\p$};
			\end{tikzpicture}
		}
		\noindent
		Xét hàm số $f(x)=2x^2+\dfrac{4}{x} \Rightarrow f'(x)=4x-\dfrac{4}{x^2}$. \\
		Xét $f'(x)=0 \Leftrightarrow x=1$.\\
		Bảng biến thiên của hàm số $f(x)$
		\begin{center}
			\begin{tikzpicture}[scale=1]
				\tkzTabInit[nocadre=false, lgt=1.2, espcl=3.5, deltacl=0.6]{$x$/0.6, $f'(x)$/0.6, $f(x)$/2}{$0$, $1$, $+\infty$}
				\tkzTabLine{,-,0,+,}
				\tkzTabVar{+/$+\infty$, -/$14$, +/$+\infty$}
			\end{tikzpicture}
		\end{center}
		Từ bảng biến thiên, ta có $\min\limits_{(0;+\infty)} f(x)=f(1)=14$. \\
		Vậy chiều rộng của đáy hình hộp đứng đó là $1$ m thì số tôn cần sử dụng là nhỏ nhất.
	}
\end{bt}

\begin{bt}%[2D1V3-6]
% nguon: Archive/2024/0. [LVH] BG - 12 New - 4in1-20240625T224301Z-001/0. [LVH] BG - 12 New - 4in1/2. [LVH] TEX SGK 12 NEW-20240418T074833Z-001/2. [LVH] TEX SGK 12 NEW/[LVH] CD-12new/SCD-12-CD/data/12-SCD-CD-2-2.tex #1062
%[ĐẠI DỰ ÁN TEX HOÁ SGK + BÀI GIẢNG LỚP 12]%[12-SCD-CD-2-2 - Thành Đức Trung]
	Giả sử lợi nhuận của doanh nghiệp khi sản xuất $x$ sản phẩm là $f(x)$ với $f(x)$ có đạo hàm tại mọi điểm $x$ trên khoảng $(0;+\infty)$. \\
	Khi đó, lợi nhuận gia tăng khi sản xuất thểm $1$ sản phẩm (từ $x$ sản phẩm lên $x+1$ sản phẩm) là hiệu $f(x+1)-f(x)$, hiệu này còn được gọi là giá trị $f$-cận biên tại $x$. \\
	Lợi nhuận cận biên khi sản xuất $x$ sản phẩm, kí hiệu là $Mf(x)$, được xác định bởi $f'(x)$.
	\begin{enumerate}
		\item Nêu một giá trị xấp xỉ của giá trị $f$-cận biên tại $x=64$, biết rằng $Mf(x)=\dfrac{8}{\sqrt[3]{x^2}}$.
		\item Giả sử hàm lợi nhuận cận biên $Mf(x)$ giảm dần khi $x$ càng lớn. Nêu nhận xét về hàm giá trị $f$-cận biên khi $x$ càng lớn.
	\end{enumerate}
	\loigiai
	{
		\begin{enumerate}
			\item
			Vì $f'(x)=\lim\limits_{\Delta x \to 0}\dfrac{f(x+\Delta x)-f(x)}{\Delta x}$ nên ta có thể lấy
			\begin{center}
				$f'(x)\approx\dfrac{f(x+1)-f(x)}{1}$ hay $f(x+1)-f(x)\approx Mf(x)$.
			\end{center}
			%Do đó, ta có thể lấy $Mf(x)$ là giá trị xấp xỉ của giá trị $f$-cận biên tại $x$. \\
			Vậy ta có thể lấy một giá trị xấp xỉ của giá trị $f$-cận biên tại $x=64$ là
			$$Mf(64)=\dfrac{8}{\sqrt[3]{64^2}}=0{,}5.$$
			\item
			Do $f(x+1)-f(x)\approx Mf(x)$ và $Mf(x)$ giảm dần khi $x$ càng lớn nên ta có thể thấy rằng hàm giá trị $f$-cận biên cũng giảm khi $x$ càng lớn.
		\end{enumerate}
	}
\end{bt}

\begin{bt}%[2D1V3-6]
% nguon: Archive/2024/0. [LVH] BG - 12 New - 4in1-20240625T224301Z-001/0. [LVH] BG - 12 New - 4in1/2. [LVH] TEX SGK 12 NEW-20240418T074833Z-001/2. [LVH] TEX SGK 12 NEW/[LVH] CD-12new/SCD-12-CD/data/12-SCD-CD-2-2.tex #1063
%[ĐẠI DỰ ÁN TEX HOÁ SGK + BÀI GIẢNG LỚP 12]%[12-SCD-CD-2-2 - Thành Đức Trung]
	Một doanh nghiệp sản xuất một loại sản phẩm. Giả sử tổng chi phí (đơn vị: triệu đồng) để sản xuất và bán hết $x$ sản phẩm đó được cho bởi
	$$f(x)=0{,}0001x^2+0{,}2x+10000.$$
	Giả sử tổng số tiền doanh nghiệp thu được (đơn vị: triệu đồng) khi bán hết $x$ sản phẩm đó là
	$$g(x)=2x+9000.$$
	\begin{flushright}
		\textit{(Nguồn: R.Larson and B.Edwards, Calculus $10$e, Cengage $2014$).}
	\end{flushright}
	\begin{enumerate}
		\item Tính số tiền lãi doanh nghiệp thu được khi bán hết $x$ sản phẩm đó. Từ đó, hãy cho biết doanh nghiệp cần sản xuất bao nhiêu sản phẩm để số tiền lãi thu được là lớn nhất.
		\item Tỉ số $M(x)=\dfrac{f(x)}{x}$ $(x\geqslant1)$ được gọi là chi phí trung bình cho một sản phẩm khi bán ra. Hãy cho biết doanh nghiệp cần sản xuất bao nhiêu sản phẩm để chi phí trung bình là nhỏ nhất.
	\end{enumerate}
	\loigiai
	{
		\begin{enumerate}
			\item
			Số tiền lãi doanh nghiệp thu được khi bán hết $x$ sản phẩm là
			$$g(x)-f(x)=(2x+9000)-\left(0{,}0001x^2+0{,}2x+10000\right)=-0{,}0001x^2+1{,}8x-1000.$$
			Xét hàm số $h(x)=-0{,}0001x^2+1{,}8x-1000$ với $x\geqslant1$. \\
			Ta có $h'(x)=-0{,}0002x+1{,}8$. \\
			Xét $h'(x)=0 \Leftrightarrow x=9000$. \\
			Bảng biến thiên của hàm số $h(x)$
			\begin{center}
				\begin{tikzpicture}[scale=1]
					\tkzTabInit[nocadre=false, lgt=1.2, espcl=3.5, deltacl=1.6]{$x$/0.6, $h'(x)$/0.6, $h(x)$/2}{$1$, $9000$, $+\infty$}
					\tkzTabLine{,+,0,-,}
					\tkzTabVar{-/$h(1)\approx-998{,}2$, +/$7100$, -/$-\infty$}
				\end{tikzpicture}
			\end{center}
			Từ bảng biến thiên, ta có $\max\limits_{[1;+\infty)} h(x)=h(9000)=7100$. \\
			Vậy doanh nghiệp cần sản xuất $9000$ sản phẩm để số tiền lãi thu được là lớn nhất.
			\item
			Ta có $M(x)=\dfrac{0{,}0001x^2+0{,}2x+10000}{x}=0{,}0001x+\dfrac{10000}{x}+0{,}2$. \\
			Xét hàm số $M(x)=0{,}0001x+\dfrac{10000}{x}=0{,}2$ với $x\geqslant 1$ \\
			Ta có $M'(x)=0{,}0001-\dfrac{10000}{x^2}$. \\
			Xét $M'(x)=0 \Leftrightarrow x=10000$ (do $x\geqslant1$). \\
			Bảng biến thiên của hàm số $M(x)$
			\begin{center}
				\begin{tikzpicture}[scale=1]
					\tkzTabInit[nocadre=false, lgt=1.5, espcl=3.5, deltacl=1.6]{$x$/0.6, $M'(x)$/0.6, $M(x)$/2}{$1$, $10000$, $+\infty$}
					\tkzTabLine{,-,0,+,}
					\tkzTabVar{+/$M(1)\approx10000{,}2$, -/$0{,}2$, +/$+\infty$}
				\end{tikzpicture}
			\end{center}
			Từ bảng biến thiên, ta có $\min\limits_{[1;+\infty)} M(x)=M(10000)=0{,}2$. \\
			Vậy doanh nghiệp cần sản xuất $10000$ sản phẩm để chi phí trung bình là nhỏ nhất.
		\end{enumerate}
	}
\end{bt}

\begin{bt}%[2D1V3-6]
% nguon: Archive/2024/0. [LVH] BG - 12 New - 4in1-20240625T224301Z-001/0. [LVH] BG - 12 New - 4in1/2. [LVH] TEX SGK 12 NEW-20240418T074833Z-001/2. [LVH] TEX SGK 12 NEW/[LVH] CD-12new/SCD-12-CD/data/12-SCD-CD-2-2.tex #1064
%[ĐẠI DỰ ÁN TEX HOÁ SGK + BÀI GIẢNG LỚP 12]%[12-SCD-CD-2-2 - Thành Đức Trung]
	Một cửa hàng nhập bưởi da xanh với giá nhập là $150000$ đồng/$1$ quả và bán bưởi đó với giá bán là $200000$ đồng/$1$ quả. Với giá bán này thì cửa hàng dự kiến chỉ bán được $50$ quả. Cửa hàng này quyết định giảm giá bán. Biết rằng nếu cửa hàng cứ giảm mỗi quả $10000$ đồng thì số bưởi bán được tăng thêm là $50$ quả. Xác định giá bán bưởi để cửa hàng đó thu được lợi nhuận lớn nhất.
	\loigiai
	{
		Gọi $x$ là số lần giảm giá bán bưởi $(x\in\mathbb{N}^{\ast})$. \\
		Số tiền giảm giá cho mỗi quả bưởi là $10000x$. \\
		Khi đó, giá bán của mỗi quả bưởi là $200000-10000x$ và số bưởi bán được là $50+50x$. \\
		Lợi nhuận cửa hàng thu được là
		$$\begin{aligned}
			(200000-10000x)(50+50x)-150000(50+50x)
			&=(50000-10000x)(50+50x) \\
			&=500000(5-x)(1+x) \\
			&=500000\left(5+4x-x^2\right).
		\end{aligned}$$
		Xét hàm số $f(x)=500000\left(5+4x-x^2\right)$ với $1\leqslant x\leqslant20$. \\
		Ta có $f'(x)=500000(4-2x)$. \\
		Xét $f'(x)=0 \Leftrightarrow x=2$. \\
		Bảng biến thiên của hàm số $f(x)$
		\begin{center}
			\begin{tikzpicture}[scale=1]
				\tkzTabInit[nocadre=false, lgt=1.2, espcl=3.5, deltacl=1.4]{$x$/0.6, $f'(x)$/0.6, $f(x)$/2}{$1$, $2$, $20$}
				\tkzTabLine{,+,0,-,}
				\tkzTabVar{-/$4000000$, +/$4500000$, -/$-157500000$}
			\end{tikzpicture}
		\end{center}
		Từ bảng biến thiên, ta có $\max\limits_{[1;20]} f(x)=f(2)=4500000$. \\
		Vậy cửa hàng cần bán bưởi với giá $180000$ đồng/$1$ quả để số tiền lãi thu được là lớn nhất.
	}
\end{bt}

\begin{bt}%[2D1V3-6]
% nguon: Archive/2024/0. [LVH] BG - 12 New - 4in1-20240625T224301Z-001/0. [LVH] BG - 12 New - 4in1/2. [LVH] TEX SGK 12 NEW-20240418T074833Z-001/2. [LVH] TEX SGK 12 NEW/[LVH] CD-12new/SCD-12-CD/data/12-SCD-CD-2-2.tex #1065
%[ĐẠI DỰ ÁN TEX HOÁ SGK + BÀI GIẢNG LỚP 12]%[12-SCD-CD-2-2 - Thành Đức Trung]
	Một công ty có $50$ căn phòng cho thuê. Biết rằng nếu công ty cho thuê mỗi căn phòng với giá $2$ triệu đồng/$1$ tháng thì mọi căn phòng đều có người ở, nhưng cứ mỗi lần tăng giá cho thuê mỗi căn phòng $100000$ đồng/$1$ tháng thì có thêm hai căn phòng bị bỏ trống. Công ty phải cho thuê mỗi căn phòng với giá là bao nhiêu để tổng số tiền thu được là lớn nhất.
	\loigiai
	{
		Gọi $x$ là số lần tăng giá cho thuê phòng $(x\in\mathbb{N}^{\ast})$. \\
		Số tiền tăng giá cho thuê mỗi căn phòng là $100x$ (nghìn đồng). \\
		Khi đó, giá cho thuê của mỗi phòng là $2000+100x$ và số phòng cho thuê là $50-2x$. \\
		Tổng số tiền công ty thu được là
		$$(2000+100x)(50-2x)=100000+1000x-200x^2.$$
		Xét hàm số $f(x)=100000+1000x-200x^2$ với $0\leqslant x\leqslant25$. \\
		Ta có $f'(x)=1000-400x$. \\
		Xét $f'(x)=0 \Leftrightarrow x=\dfrac{10}{4}=2{,}5$. \\
		Bảng biến thiên của hàm số $f(x)$
		\begin{center}
			\begin{tikzpicture}[scale=1]
				\tkzTabInit[nocadre=false, lgt=1.2, espcl=4.5, deltacl=1]{$x$/0.6, $f'(x)$/0.6, $f(x)$/2}{$0$, $2{,}5$, $25$}
				\tkzTabLine{,+,0,-,}
				\tkzTabVar{-/$100000$, +/$101250$, -/$0$}
			\end{tikzpicture}
		\end{center}
		Từ bảng biến thiên, ta có $\max\limits_{[0;25]} f(x)=f\left(2{,}5\right)$. \\
		Do $x\in\mathbb{N}^{\ast}$ nên ta xét $f(2)=101200$ và $f(3)=101200$. \\
		Vậy để tổng số tiền thu được là lớn nhất thì công ty có thể cho thuê phòng với giá $2{,}2$ triệu đồng hoặc $2{,}4$ triệu đồng.
	}
\end{bt}

\begin{bt}%[2D1V3-6]
% nguon: Archive/2024/0. [LVH] BG - 12 New - 4in1-20240625T224301Z-001/0. [LVH] BG - 12 New - 4in1/2. [LVH] TEX SGK 12 NEW-20240418T074833Z-001/2. [LVH] TEX SGK 12 NEW/[LVH] CD-12new/SCD-12-CD/data/12-SCD-CD-2-2.tex #1066
%[ĐẠI DỰ ÁN TEX HOÁ SGK + BÀI GIẢNG LỚP 12]%[12-SCD-CD-2-2 - Thành Đức Trung]
	Bạn Hà có một tấm bia hình vuông cạnh $60$ cm. Bạn muốn làm một cái hộp đựng đồ có dạng hình hộp chữ nhật mà có thể để được vào một ngăn sách có dạng hình hộp chữ nhật, đáy là hình vuông cạnh bằng $37$ cm, chiều cao bằng $28$ cm. Bạn cắt bốn góc của tấm bìa đó bốn hình vuông bằng nhau, mỗi hình vuông có cạnh bằng $x$ (cm), rồi gập lại thành một cái hộp không nắp. Tìm số nguyên dương $x$ để làm được cái hộp đựng đồ có thể tích lớn nhất.
	\begin{center}
		\begin{tikzpicture}[scale=1, font=\footnotesize, line join=round, line cap=round, >=stealth]
			\path
			(0,0)coordinate(A)
			(0,4)coordinate(B)
			(4,4)coordinate(C)
			($(C)+(A)-(B)$)coordinate(D)
			
			(0,3)coordinate(M)
			(4,3)coordinate(N)
			
			(0,1)coordinate(E)
			(4,1)coordinate(F)
			
			(1,0)coordinate(G)
			(1,4)coordinate(H)
			
			(3,0)coordinate(I)
			(3,4)coordinate(K)
			;
			\draw
			(A)--(B)--(C)--(D)--(A)
			;
			\draw[dashed]
			(M)--(N) (E)--(F) (I)--(K) (G)--(H)
			;
		\end{tikzpicture}
		\quad\quad\quad\quad\quad\quad
		\begin{tikzpicture}[scale=1, font=\footnotesize, line join=round, line cap=round, >=stealth]
			\path
			(0,0)coordinate(A)
			(-2,-1.5)coordinate(B)
			(5,0)coordinate(D)
			($(B)+(D)-(A)$)coordinate(C)
			($(A)+(0,2)$)coordinate(A')
			($(A')+(B)-(A)$)coordinate(B')
			($(A')+(C)-(A)$)coordinate(C')
			($(A')+(D)-(A)$)coordinate(D')
			(intersection of A--A' and B'--C')coordinate(E)
			;
			\draw
			(B)--(B')
			(C)--(C')
			(D)--(D')
			(B)--(C)
			(C)--(D)
			(A')--(B')
			(B')--(C')
			(C')--(D')
			(D')--(A')--(E)
			;
			\draw[dashed]
			(B)--(A)
			(A)--(D)
			(A)--(E)
			;
		\end{tikzpicture}
	\end{center}
	\loigiai
	{
		Hộp đựng đồ có đáy là hình vuông cạnh $60-2x$ (cm), chiều cao bằng $x$ (cm). \\
		Điều kiện $\heva{ & x\leqslant28 \\ & 60-2x\leqslant37} \Leftrightarrow \heva{ & x<28 \\ & x\geqslant11{,}5} \Leftrightarrow 11{,}5\leqslant x<28$. \\
		Thể tích hộp đựng đồ là $V=x\left(60-2x\right)^2=4\left(-x^3+60x^2-900x\right)$ với $x\in\left(11{,}5;28\right)$. \\
		Xét hàm số $f(x)=4\left(-x^3+60x^2-900x\right)$. \\
		Ta có $f'(x)=4\left(-3x^2+120x-900\right)$. \\
		Xét $f'(x)=0 \Leftrightarrow \hoac{ & x=30 \\ & x=10.}$ \\
		Bảng biến thiên của hàm số $f(x)$
		\begin{center}
			\begin{tikzpicture}[scale=1]
				\tkzTabInit[nocadre=false, lgt=1.2, espcl=4.5, deltacl=0.6]{$x$/0.6, $f'(x)$/0.6, $f(x)$/2}{$11{,}5$, $28$}
				\tkzTabLine{,-,}
				\tkzTabVar{+/, -/}
			\end{tikzpicture}
		\end{center}
		Từ bảng biến thiên và $x$ nguyên dương, ta có thể tích hộp đựng lớn nhất khi $x=12$. \\
		Vậy phải cắt mỗi hình vuông cạnh cạnh $12$ cm để được hộp hình chữ nhật có thể tích lớn nhất.
	}
\end{bt}

\begin{bt}%[2D1V3-6]
% nguon: Archive/2024/0. [LVH] BG - 12 New - 4in1-20240625T224301Z-001/0. [LVH] BG - 12 New - 4in1/2. [LVH] TEX SGK 12 NEW-20240418T074833Z-001/2. [LVH] TEX SGK 12 NEW/[LVH] CD-12new/SCD-12-CD/data/12-SCD-CD-2-2.tex #1067
%[ĐẠI DỰ ÁN TEX HOÁ SGK + BÀI GIẢNG LỚP 12]%[12-SCD-CD-2-2 - Thành Đức Trung]
	\immini
	{
		Hình vẽ minh họa một màn hình $BC$ có chiều cao $1{,}4$ m được đặt thẳng đứng và mép dưới của màn hình cách mặt đất một khoảng $BA=1{,}8$ m. Một chiếc đèn quan sát màn hình được đặt ở vị trí $O$ trên mặt đất. Hãy tính khoảng cách $AO$ sao cho góc quan sát $\widehat{BOC}$ là lớn nhất.
	}
	{
		\begin{tikzpicture}[scale=1, font=\footnotesize, line join=round, line cap=round, >=stealth]
			\path
			(0,0)coordinate(A)
			(0,1.8)coordinate(B)
			(0,3.2)coordinate(C)
			(5,0)coordinate(O)
			($(B)+(-0.5,0)$)coordinate(B')
			($(C)+(-0.5,0)$)coordinate(C')
			($(A)+(-0.5,0)$)coordinate(A')
			($(B')!0.5!(C')$)node[left]{$1{,}4$ m}
			($(B')!0.5!(A')$)node[left]{$1{,}8$ m}
			;
			\draw
			(O)--(C)--(B)--(O)--(A)
			;
			\draw[dashed]
			(A)--(B)
			(B)--(B')
			(C)--(C')
			(A)--(A')
			;
			\draw[dashed, <->]
			(B')--(C')
			;
			\draw[dashed, <->]
			(B')--(A')
			;
			\foreach \p/\q in {O/-90, A/-90, B/-135, C/90} \fill[black] (\p) circle (1pt) ($(\p)+(\q:2.5mm)$) node{$\p$};
		\end{tikzpicture}
	}
	\loigiai
	{
		Để góc $\widehat{BOC}$ lớn nhất thì $\tan \widehat{BOC}$ lớn nhất.\\
		Đặt $AO=x$ (m) với $x>0$. \\
		Ta có
		$$\begin{aligned}
			\tan \widehat{BOC}
			&=\tan \left(\widehat{AOC}-\widehat{AOB}\right)=\dfrac{\tan \widehat{AOC}-\tan\widehat {AOB}}{1+\tan \widehat{AOC}\cdot\tan\widehat{AOB}}\\
			&=\dfrac{\dfrac{AC}{OA}-\dfrac{AB}{OA}}{1+\dfrac{AC}{OA}\cdot\dfrac{AB}{OA}}=\dfrac{\dfrac{1{,}4}{x}}{1+\dfrac{3{,}2\cdot1{,}8}{x^2}}=\dfrac{1{,}4x}{x^2+5{,}76}.
		\end{aligned}$$
		Xét hàm số $f(x)=\dfrac{1{,}4x}{x^2+5{,}76}$ trên $(0;+\infty)$. \\
		Ta có $f'(x)=\dfrac{-1{,}4x^2+1{,}4\cdot5{,}76}{\left(x^2+5{,}76\right)^2}$. \\
		Xét $f'(x)=0 \Leftrightarrow x^2=5{,}76 \Leftrightarrow x=2{,}4$ (do $x>0$). \\
		Bảng biến thiên của hàm số $f(x)$
		\begin{center}
			\begin{tikzpicture}[scale=1]
				\tkzTabInit[nocadre=false, lgt=1.2, espcl=4.5, deltacl=0.6]{$x$/0.6, $f'(x)$/0.6, $f(x)$/2}{$0$, $2{,}4$, $+\infty$}
				\tkzTabLine{,+,0,-,}
				\tkzTabVar{-/$0$, +/$\dfrac{7}{24}$, -/$0$}
			\end{tikzpicture}
		\end{center}
		Từ bảng biến thiên, ta có $\max\limits_{(0;+\infty)} f(x)=f\left(2{,}4\right)$. \\
		Vậy $AO=2{,}4 $m thì góc quan sát $\widehat{BOC}$ là lớn nhất.
	}
\end{bt}

\begin{bt}%[2D1V3-6]
% nguon: Archive/2024/0. [LVH] BG - 12 New - 4in1-20240625T224301Z-001/0. [LVH] BG - 12 New - 4in1/2. [LVH] TEX SGK 12 NEW-20240418T074833Z-001/2. [LVH] TEX SGK 12 NEW/[LVH] CD-12new/SCD-12-CD/data/12-SCD-CD-2-2.tex #1068
%[ĐẠI DỰ ÁN TEX HOÁ SGK + BÀI GIẢNG LỚP 12]%[12-SCD-CD-2-2 - Thành Đức Trung]
	Khi nuôi cá thí nghiệm trong hồ, một nhà sinh vật học thấy rằng: Nếu trên mỗi đơn vị diện tích của mặt hồ có $n$ con cá thì trung bình mỗi con cá sau một vụ cân nặng
	\begin{center}
		$P(n)=480-20n$ (gam).
	\end{center}
	Hỏi phải thả bao nhiêu cá trên một đơn vị diện tích của mặt hồ để sau một vụ thu hoạch được nhiều cá nhất?
	\begin{flushright}
		\textit{(Nguồn: Giải tích $12$ - Nâng cao, NXBGD Việt Nam, $2020$).}
	\end{flushright}
	\loigiai
	{
		Khối lượng cá sau một vụ thu hoạch trên một đơn vị diện tích là 
		$$nP(n)=n(480-20n)=480n-20n^2 \ (\text{gam}).$$
		Xét hàm số $f(n)=480n-20n^2$ với $n\in\mathbb{N}^{\ast}$. \\
		Ta có $f'(n)=480-40n$. \\
		Xét $f'(n)=0 \Leftrightarrow n=12$. \\
		Bảng biến thiên của hàm số $f(n)$
		\begin{center}
			\begin{tikzpicture}[scale=1]
				\tkzTabInit[nocadre=false, lgt=1.2, espcl=4.5, deltacl=0.6]{$n$/0.6, $f'(n)$/0.6, $f(n)$/2}{$1$, $12$, $+\infty$}
				\tkzTabLine{,+,0,-,}
				\tkzTabVar{-/$460$, +/$2880$, -/$-\infty$}
			\end{tikzpicture}
		\end{center}
		Từ bảng biến thiên, ta có $\max\limits_{[1;+\infty)} f(n)=f(12)=2880$. \\
		Vậy phải thả $12$ con cá trên một đơn vị diện tích của mặt hồ để một vụ thu hoạch được nhiều cá nhất.
	}
\end{bt}

\begin{bt}%[2D1V3-6]
% nguon: Archive/2024/0. [LVH] BG - 12 New - 4in1-20240625T224301Z-001/0. [LVH] BG - 12 New - 4in1/2. [LVH] TEX SGK 12 NEW-20240418T074833Z-001/2. [LVH] TEX SGK 12 NEW/[LVH] CD-12new/SCD-12-CD/data/12-SCD-CD-2-2.tex #1069
%[ĐẠI DỰ ÁN TEX HOÁ SGK + BÀI GIẢNG LỚP 12]%[12-SCD-CD-2-2 - Thành Đức Trung]
	Một con cá hồi bơi ngược dòng để vượt một khoảng cách là $300$ km. Vận tốc dòng nước là $6$ km/h. Nếu vận tốc bơi của cá khi nước đứng yên là $v$ (km/h) thì năng lượng tiêu hao của cá trong $t$ giờ được cho bởi công thức
	$$E(v)=cv^3t,$$
	trong đó $c$ là một hằng số, $E$ được tính bằng jun. Tìm vận tốc bơi của cá khi nước đứng yên để năng lượng tiêu hao là ít nhất.
	\begin{flushright}
		\textit{(Nguồn: Giải tích $12$ - Nâng cao, NXBGD Việt Nam, $2020$).}
	\end{flushright}
	\loigiai
	{
		Vận tốc ngược dòng của cá khi bơi là $v-6$ (km/h). \\
		Thời gian để bơi $300$ km là $t=\dfrac{300}{v-6}$ (giờ). \\
		Năng lượng tiêu hao là $E=cv^3\dfrac{300}{v-6}$. \\
		Để năng lượng tiêu hao của cá là ít nhất thì $\left(v^3\dfrac{300}{v-6}\right)$ phải bé nhất do $c$ là hằng số. \\
		Xét hàm số $f(v)=v^3\dfrac{300}{v-6}$ với $v>6$. \\
		Ta có $f'(v)=300\cdot\dfrac{2v^3-18v^2}{\left(v-6\right)^2}$. \\
		Xét $f'(v)=0 \Leftrightarrow v=9$ (vì $v>6$). \\
		Bảng biến thiên của hàm số $f(v)$
		\begin{center}
			\begin{tikzpicture}[scale=1]
				\tkzTabInit[nocadre=false, lgt=1.2, espcl=4.5, deltacl=0.6]{$v$/0.6, $f'(v)$/0.6, $f(v)$/2}{$6$, $9$, $+\infty$}
				\tkzTabLine{,-,0,+,}
				\tkzTabVar{+/$+\infty$, -/$21600$, +/$+\infty$}
			\end{tikzpicture}
		\end{center}
		Từ bảng biến thiên, ta có $\min\limits_{(6;+\infty)} f(v)=f(9)=21600$. \\
		Vậy để tiêu hao ít năng lượng nhất thì vận tốc bơi của cá khi nước đứng yên là $9$ km/h.
	}
\end{bt}

\begin{bt}%[2D1V3-6]
% nguon: Archive/2024/0. [LVH] BG - 12 New - 4in1-20240625T224301Z-001/0. [LVH] BG - 12 New - 4in1/2. [LVH] TEX SGK 12 NEW-20240418T074833Z-001/2. [LVH] TEX SGK 12 NEW/[LVH] CD-12new/SCD-12-CD/data/12-SCD-CD-2-2.tex #1070
%[ĐẠI DỰ ÁN TEX HOÁ SGK + BÀI GIẢNG LỚP 12]%[12-SCD-CD-2-2 - Thành Đức Trung]
	Một nhà máy sản xuất xe đạp cho thị trường châu Âu theo đơn giá $120$ euro. Chi phí mỗi ngày của nhà máy được cho bởi hàm số
	$$K(x)=0{,}02x^3-3x^2+172x+2400,$$
	trong đó $x$ là số lượng xe đạp sản xuất được trong ngày hôm đó. Mỗi ngày có thể sản xuất tối đa $130$ xe đạp. Giả sử số xe đạp sản xuất được trong mỗi ngày đều được bán hết vào cuối ngày đó. Gọi $G(x)$ là hàm số biểu diễn lợi nhuận hàng ngày của nhà máy.
	\begin{enumerate}
		\item Vẽ đồ thị hàm số $G(x)$ trên đoạn $[0;140]$.
		\item Số lượng xe mỗi ngày cần sản xuất là bao nhiêu để nhà máy có lãi?
		\item Số lượng xe mỗi ngày cần sản xuất là bao nhiêu để nhà máy có lợi nhuận lớn nhất? 
		\item Giả sử nhà máy quyết định tận dụng tối đa công suất sản xuất $130$ xe đạp mỗi ngày. Nhà máy phải chọn đơn giá là bao nhiêu để có lãi?
	\end{enumerate}
	\begin{flushright}
		\textit{(Nguồn: A. Bigalke et al., Mathematik, Grundkurs ma-$1$, Cornelsen $2016$).}
	\end{flushright}
	\loigiai
	{
		\begin{enumerate}
			\item
			Ta có $G(x)=120x-K(x)=-0{,}02x^3+3x^2-52x-2400$. \\
			Ta có $G'(x)=-0{,}06x^2+6x-52=-0{,}02\left(3x^2-300x+2600\right)$. \\
			Xét $G'(x)=0 \Leftrightarrow \hoac{ & x=\dfrac{150+70\sqrt{3}}{3} \\ & x=\dfrac{150-70\sqrt{3}}{3}.}$ \\
			Bảng biến thiên của $G(x)$
			\begin{center}
				\begin{tikzpicture}[scale=1]
					\tkzTabInit[nocadre=false, lgt=1.2, espcl=4.5, deltacl=0.6]{$x$/1.2, $G'(x)$/0.6, $G(x)$/2}{$0$, $\dfrac{150-70\sqrt{3}}{3}$, $\dfrac{150+70\sqrt{3}}{3}$, $140$}
					\tkzTabLine{,-,0,+,0,-,}
					\tkzTabVar{+/, -/, +/, -/}
				\end{tikzpicture}
			\end{center}
			Đồ thị hàm số của $G(x)$
			\begin{center}
				\begin{tikzpicture}[scale=1, font=\footnotesize, line join=round, line cap=round, >=stealth, xscale=3, yscale=3]
					\draw[->](-0.5,0)--(1.7,0)node[below]{$x$};
					\draw[->](0,-0.5)--(0,0.5)node[left]{$y$};
					\draw[smooth, samples=100, domain=0:1.4]plot(\x,{-(\x+0.2)*(\x-0.5)*(\x-1.2)});
					\path
					(0,0)node[above left]{$O$}
					(1.4,0)node[above]{$140$}
					(0,-3/25)node[left]{$-2400$}
					(0,-36/125)node[left]{$-5760$}
					;
					\draw[dashed]
					(1.4,0)--(1.4,-36/125)--(0,-36/125)
					;
				\end{tikzpicture}
			\end{center}
			\item
			Để nhà máy có lãi thì lợi nhuận phải lớn hơn $0$, tức là
			$$G(x)>0 \Leftrightarrow -0{,}02x^3+3x^2-52x-2400>0 \Leftrightarrow 50<x<120 \ (\text{vì} \ x>0).$$
			Vậy số lượng xe cần sản xuất mỗi ngày để nhà máy có lãi nằm trong khoảng $(50;120)$.
			\item
			Bảng biến thiên của $G(x)$
			\begin{center}
				\begin{tikzpicture}[scale=1]
					\tkzTabInit[nocadre=false, lgt=1.2, espcl=2.5, deltacl=0.6]{$x$/1.2, $G(x)$/2}{$0$, $\dfrac{150-70\sqrt{3}}{3}$, $50$, $\dfrac{150+70\sqrt{3}}{3}$, $120$, $140$}
					\tkzTabVar{+/, -/, R, +/, R, -/}
				\end{tikzpicture}
			\end{center}
			Từ vào bảng biến thiên, ta có $\max\limits_{(50;120)} G(x)=G\left(\dfrac{150+70\sqrt{3}}{3}\right)$. \\
			Vì $x\in\mathbb{N}$ nên xét $G(90)=2640$ và $G(91)=2639{,}58$. \\
			Vậy để thu được lợi nhuận lớn nhất thì số lượng xe đạp cần sản xuất là $90$ cái.
			\item
			Gọi đơn giá là $y$ (euro) với $y>0$. \\
			Nếu mỗi ngày sản xuất được $130$ xe đạp thì chi phí là $K(130)=18000$. \\
			Khi đó lợi nhuận của nhà máy là $G(y)=130y-18000$. \\
			Để thu được lợi nhuận thì $G(y)>0 \Leftrightarrow y>\dfrac{18000}{130}$. \\
			Vậy nhà máy phải chọn đơn giá lớn hơn $\dfrac{18000}{130}$ euro để có lãi.
		\end{enumerate}
	}
\end{bt}

\begin{bt}%[2D1V3-6]
% nguon: Archive/2024/0. [LVH] BG - 12 New - 4in1-20240625T224301Z-001/0. [LVH] BG - 12 New - 4in1/2. [LVH] TEX SGK 12 NEW-20240418T074833Z-001/2. [LVH] TEX SGK 12 NEW/[LVH] CD-12new/SCD-12-CD/data/12-SCD-CD-2-2.tex #1071
%[ĐẠI DỰ ÁN TEX HOÁ SGK + BÀI GIẢNG LỚP 12]%[12-SCD-CD-2-2 - Thành Đức Trung]
	Một nhà máy sản xuất một loại sản phẩm cho thị trường Mỹ. Biết rằng
	\begin{itemize}
		\item Chi phí cho các công việc hành chính chung của nhà máy là $90$ đô la Mỹ (USD)/$1$ ngày.
		\item Chi phí sản xuất là $0{,}09$ USD/$1$ sản phẩm.
		\item Các loại chi phí khác trong mỗi một ngày là $\dfrac{x^2}{10000}$ (USD), trong đó $x$ là số sản phẩm nhà máy sản xuất được trong ngày hôm đó.
	\end{itemize}
	\begin{enumerate}
		\item Tính tổng chi phí $U(x)$ của mỗi một sản phẩm.
		\item Tìm $x$ sao cho $U(x)$ nhận giá trị nhỏ nhất.
	\end{enumerate}
	\loigiai
	{
		\begin{enumerate}
			\item
			Giả sử nhà máy sản xuất $x$ sản phẩm trong một ngày $\left(x\in\mathbb{N}^{\ast}\right)$.
			\begin{itemize}
				\item Chi phí cho các công việc hành chính chung của nhà máy cho mỗi một sản phẩm là $\dfrac{90}{x}$.
				\item Chi phí sản xuất là $0{,}09$ USD/$1$ sản phẩm.
				\item Các loại chi phí khác cho mỗi một sản phẩm là $\dfrac{x}{10000}$ (USD).
			\end{itemize}
			Tổng chi phí sản xuất mỗi một sản phẩm là $U(x)=\dfrac{90}{x}+0{,}09+\dfrac{x}{10000}$.
			\item
			Xét hàm số $U(x)=\dfrac{90}{x}+0{,}09+\dfrac{x}{10000}$ với $x>0$. \\
			Ta có $U'(x)=-\dfrac{90}{x^2}+\dfrac{1}{10000}$. \\
			Xét $U'(x)=0 \Leftrightarrow x=300\sqrt{10}$ (vì $x>0$). \\
			Bảng biến thiên của hàm số $U(x)$
			\begin{center}
				\begin{tikzpicture}[scale=1]
					\tkzTabInit[nocadre=false, lgt=1.2, espcl=4.5, deltacl=0.6]{$x$/0.6, $U'(x)$/0.6, $U(x)$/2}{$0$, $300\sqrt{10}$, $+\infty$}
					\tkzTabLine{,-,0,+,}
					\tkzTabVar{+/$+\infty$, -/$U\left(300\sqrt{10}\right)$, +/$+\infty$}
				\end{tikzpicture}
			\end{center}
			Từ bảng biến thiên, ta có $\min\limits_{(0;+\infty)} U(x)=U\left(300\sqrt{10}\right)$. \\
			Vì $x\in\mathbb{N}^{\ast}$ nên ta xét $U(948)\approx0{,}2797367089$ và $U(949)\approx0{,}2797366702$. \\
			Suy ra $U(949)<U(948)$. \\
			Vậy cần sản xuất $949$ sản phẩm để chi phí của mỗi một sản phẩm là nhỏ nhất.
		\end{enumerate}
	}
\end{bt}

\begin{bt}%[2D1H1-1]
% nguon: Archive/2024/0. [LVH] BG - 12 New - 4in1-20240625T224301Z-001/0. [LVH] BG - 12 New - 4in1/2. [LVH] TEX SGK 12 NEW-20240418T074833Z-001/2. [LVH] TEX SGK 12 NEW/[LVH] CD-12new/SGK-12-CD/data/12-SGK-CD-1-1.tex #1072

	Tìm các khoảng đơn điệu của mỗi hàm số sau:
	\begin{enumerate}[a)]	
		\item $y=-x^3+2x^2-3$;
		\item $y=x^4-2x^2+5$;
		\item $y=\dfrac{3x+1}{2-x}$;
		\item $y=\dfrac{x^2-2x}{x+1}$.
	\end{enumerate}
	\loigiai{
	\begin{enumerate}[a)]	
		\item $y=-x^3+2x^2-3$.\\
		Hàm số đã cho có tập xác định là $\mathbb{R}$.\\
		Ta có: $y'=-3x^2+4x$;\\
		$y'=0\Leftrightarrow -3x^2+4x=0\Leftrightarrow x=0$ hoặc $x=\dfrac{4}{3}$.\\
		Bảng biến thiên của hàm số như sau:
		\begin{center}
		\begin{tikzpicture}[font=\footnotesize,thick,>=stealth]
		\tikzset{double style/.append style={double distance=1.5pt}}
		\tkzTabInit[nocadre=false,lgt=1.2,espcl=2.5,deltacl=0.6,lw=.75pt,color,colorL=green!50,colorV=green!50]
		{$x$ /0.8,$y'$ /0.6,$y$ /2}
		{$-\infty$,$0$,$\dfrac{4}{3}$,$+\infty$}
		\tkzTabLine{ ,+,z,-,z,+, }
		\tkzTabVar{-/$-\infty$,+/$-3$,-/$-\dfrac{49}{27}$,+/$+\infty$}
		\end{tikzpicture}
		\end{center}
		Vậy hàm số đồng biến trên mỗi khoảng $(-\infty;0)$ và $\left(\dfrac{4}{3};+\infty\right)$; nghịch biến trên khoảng $\left(0;\dfrac{4}{3}\right)$.
		\item $y=x^4-2 x^2+5$.\\
		Hàm số đã cho có tập xác định là $\mathbb{R}$.\\
		Ta có: $y'=4x^3-4x$;\\
		$y'=0\Leftrightarrow 4x^3-4x=0\Leftrightarrow x=0$ hoặc $x=\pm 1$.\\
		Bảng biến thiên của hàm số như sau:
		\begin{center}
		\begin{tikzpicture}[font=\footnotesize,thick,>=stealth]
			\tikzset{double style/.append style={double distance=1.5pt}}
			\tkzTabInit[nocadre=false,lgt=1.2,espcl=2.5,deltacl=0.6,lw=.75pt,color,colorL=green!50,colorV=green!50]
			{$x$ /0.7, $f'(x)$ /0.8, $f(x)$ /2}
			{$-\infty$,$-1$,$0$,$1$,$+\infty$}
			\tkzTabLine{ ,-,z,+,z,-,z,+, }
			\tkzTabVar{+/$+\infty$,-/$4$,+/$5$,-/$4$,+/$+\infty$}
		\end{tikzpicture}
		\end{center}
		Vậy hàm số đồng biến trên mỗi khoảng $(-1;0)$ và $\left(1;+\infty\right)$; nghịch biến trên mỗi khoảng $\left(-\infty;-1\right)$ và $(0;1)$.
		\item $y=\dfrac{3x+1}{2-x}$.\\
		Hàm số đã cho có tập xác định là $\mathbb{R}\setminus\{2\}$.\\
		Ta có: $y'=\dfrac{7}{(2-x)^2}>0,\,\forall x\neq 2$.\\
		Vậy hàm số đồng biến trên mỗi khoảng $(-\infty;2)$ và $(2;+\infty)$.
		\item $y=\dfrac{x^2-2x}{x+1}$.\\
		Hàm số đã cho có tập xác định là $\mathbb{R}\setminus\{-1\}$.\\
		Ta có: $y'=\dfrac{x^2+2x-2}{(x+1)^2}$ với $x\neq -1$;\\
		$y'= 0\Leftrightarrow x^2+2x-2=0\Leftrightarrow x=-1-\sqrt{3}$ hoặc $x=-1+\sqrt{3}$.\\
		Bảng biến thiên của hàm số như sau:
		\begin{center}
			\begin{tikzpicture}[font=\footnotesize,thick,>=stealth]
				\tikzset{double style/.append style={double distance=1.5pt}}
				\tkzTabInit[nocadre=false,lgt=1.2,espcl=2.5,deltacl=0.6,lw=.75pt,color,colorL=green!50,colorV=green!50]
				{$x$ /0.6,$y'$ /0.6,$y$ /2}
				{$-\infty$,$-1-\sqrt{3}$,$-1$,$-1+\sqrt{3}$,$+\infty$}
				\tkzTabLine{,+,$0$,-,d,-,$0$,+,}
				\tkzTabVar{-/$-\infty$,+/$-4-2\sqrt{3}$,-D+/$-\infty$/$+\infty$,-/$-4+2\sqrt{3}$,+/$+\infty$}
			\end{tikzpicture}
		\end{center}
		Vậy hàm số đồng biến trên mỗi khoảng $(-\infty;-1-\sqrt{3})$ và $(-1+\sqrt{3};+\infty)$; nghịch biến trên mỗi khoảng $(-1-\sqrt{3};-1)$ và $(-1;-1+\sqrt{3})$.
	\end{enumerate} 
	}
\end{bt}

\begin{bt}%[2D1H2-1]
% nguon: Archive/2024/0. [LVH] BG - 12 New - 4in1-20240625T224301Z-001/0. [LVH] BG - 12 New - 4in1/2. [LVH] TEX SGK 12 NEW-20240418T074833Z-001/2. [LVH] TEX SGK 12 NEW/[LVH] CD-12new/SGK-12-CD/data/12-SGK-CD-1-1.tex #1073

	Tìm cực trị của mỗi hàm số sau:
	\begin{enumerate}[a)]
		\item $y=2x^3+3x^2-36x-10$;
		\item $y=x^4+2x^2-3$;
		\item $y=x-\dfrac{1}{x}$.
	\end{enumerate}
	\loigiai{
	\begin{enumerate}[a)]
		\item $y=2x^3+3x^2-36x-10$.\\
		Hàm số đã cho có tập xác định là $\mathbb{R}$.\\
		Ta có: $y'=6x^2+6x-36$;\\
		$y'=0\Leftrightarrow 6x^2+6x-36=0\Leftrightarrow x=-3$ hoặc $x=2$.\\
		Bảng biến thiên của hàm số như sau:
		\begin{center}
			\begin{tikzpicture}[font=\footnotesize,thick,>=stealth]
				\tikzset{double style/.append style={double distance=1.5pt}}
				\tkzTabInit[nocadre=false,lgt=1.2,espcl=2.5,deltacl=0.6,lw=.75pt,color,colorL=green!50,colorV=green!50]
				{$x$ /0.6,$y'$ /0.6,$y$ /2}
				{$-\infty$,$-3$,$2$,$+\infty$}
				\tkzTabLine{ ,+,z,-,z,+, }
				\tkzTabVar{-/$-\infty$,+/$71$,-/$-54$,+/$+\infty$}
			\end{tikzpicture}
		\end{center}
		Vậy hàm số đạt cực đại tại $x=-3$ và đạt cực tiểu tại $x=2$.
		\item $y=x^4+2x^2-3$.\\
		Hàm số đã cho có tập xác định là $\mathbb{R}$.\\
		Ta có: $y'=4x^3+4x$;\\
		$y'=0\Leftrightarrow =4x^3+4x=0\Leftrightarrow x=0$.\\
		Bảng biến thiên của hàm số như sau:
		\begin{center}
			\begin{tikzpicture}[font=\footnotesize,thick,>=stealth]
				\tikzset{double style/.append style={double distance=1.5pt}}
				\tkzTabInit[nocadre=false,lgt=1.2,espcl=2.5,deltacl=0.6,lw=.75pt,color,colorL=green!50,colorV=green!50]
				{$x$ /0.7, $y'$ /0.8, $y$ /2}
				{$-\infty$,$0$,$+\infty$}
				\tkzTabLine{ ,-,$0$,+, }
				\tkzTabVar{+/$+\infty$,-/$-3$,+/$+\infty$}
			\end{tikzpicture}
		\end{center}
		Vậy hàm số đạt cực tiểu tại $x=0$.
		\item $y=x-\dfrac{1}{x}$.\\
			Hàm số đã cho có tập xác định là $\mathbb{R}\setminus\{0\}$.\\
			Ta có: $y'=1+\dfrac{1}{x^2}>0,\,\forall x\neq 0$.\\
			Bảng biến thiên của hàm số như sau:
			\begin{center}
				\begin{tikzpicture}[font=\footnotesize,thick,>=stealth]
					\tikzset{double style/.append style={double distance=1.5pt}}
					\tkzTabInit[nocadre=false,lgt=1.2,espcl=2.5,deltacl=0.6,lw=.75pt,color,colorL=green!50,colorV=green!50]
					{$x$ /0.6,$y'$ /0.6,$y$ /2}
					{$-\infty$,$0$,$+\infty$}
					\tkzTabLine{,+,d,+,}
					\tkzTabVar{-/$-\infty$,+D-/$+\infty$/$-\infty$,+/$+\infty$}
				\end{tikzpicture}
			\end{center}
			Vậy hàm số không có cực trị.
	\end{enumerate}
}
\end{bt}

\begin{bt}%[2D1N1-2]
% nguon: Archive/2024/0. [LVH] BG - 12 New - 4in1-20240625T224301Z-001/0. [LVH] BG - 12 New - 4in1/2. [LVH] TEX SGK 12 NEW-20240418T074833Z-001/2. [LVH] TEX SGK 12 NEW/[LVH] CD-12new/SGK-12-CD/data/12-SGK-CD-1-1.tex #1074

	Cho hai hàm số $y=f(x)$, $y=g(x)$ có đồ thị lần lượt được cho ở \textit{Hình} $6a$, \textit{Hình} $6b$. Nêu khoảng đồng biến, nghịch biến và điểm cực trị của mỗi hàm số đó.
	\begin{center}
		\begin{tikzpicture}[scale=1, font=\footnotesize, line join=round, line cap=round, >=stealth]
		\draw[gray!50,thin,opacity=.5] (-3,-2) grid (4,4);
		\draw[->] (-2.9,0)--(4.1,0) node[below] {$x$};
		\draw[->] (0,-2.1)--(0,4.1) node[left] {$y$};
		\draw (0,0) node [below left] {$O$};
		\foreach \x/\nx in {-1/-1,1/1,2/2}
		\draw[thin] (\x,1pt)--(\x,-1pt) node [below] {$\nx$};
		\foreach \y/\ny in {1/1,2/2}
		\draw[thin] (1pt,\y)--(-1pt,\y) node [left] {$\ny$};
		\begin{scope}
			\clip (-2.8,-2) rectangle (4,4);
			\draw[samples=200,domain=-2.5:-0.32,smooth,variable=\x] plot (\x,{-1.4*(\x)^2-2.92*(\x)-0.51});
			\draw[samples=200,domain=-0.32:0.68,smooth,variable=\x] plot (\x,{2.9*(\x)^2+0.05*(\x)});
			\draw[samples=200,domain=0.68:1.44,smooth,variable=\x] plot (\x,{-5.06*(\x)^2+10.44*(\x)-3.38});
			\draw[samples=200,domain=1.44:3.5,smooth,variable=\x] plot (\x,{4.34*(\x)^2-16.99*(\x)+16.64});
		\end{scope}
		\draw (3.4,2) node{$y=f(x)$};
		\draw (0,-4) node{\text{\textit{Hình }} $6a$};
	\end{tikzpicture}
	\hspace{0.5cm}
	\begin{tikzpicture}[scale=1, font=\footnotesize, line join=round, line cap=round, >=stealth]
		\draw[gray!50,thin,opacity=.5] (-3,-8.5) grid (3,1);
		\draw[->] (-3,0)--(3,0) node[below] {$x$};
		\draw[->] (0,-8.5)--(0,1) node[left] {$y$};
		\draw (0,0) node [above left] {$O$};
		\foreach \x/\nx in {-2/-2,1/1}
		\draw[thin] (\x,1pt)--(\x,-1pt) node [below] {$\nx$};
		\foreach \y/\ny in {-8/-8}
		\draw[thin] (1pt,\y)--(-1pt,\y) node [left] {$\ny$};
		\begin{scope}
			\clip (-3,-8.6) rectangle (2.9,1.2);
			\draw[samples=200,domain=-2.76:-1.75,smooth,variable=\x] plot (\x,{14.08*(\x)^2+55.22*(\x)+46.11});
			\draw[samples=200,domain=-1.75:0.54,smooth,variable=\x] plot (\x,{-2.4*(\x)^2+0.03*(\x)});
			\draw[samples=200,domain=0.54:1.9,smooth,variable=\x] plot (\x,{2.76*(\x)^2-5.48*(\x)+1.47});
		\end{scope}
	\draw[dashed,thin] (0,-1.25)node[left]{$-\tfrac{5}{4}$}--(1,-1.25);
	\draw (-1,-4) node[rotate=75]{$y=g(x)$};
	\draw (0,-9) node{\text{\textit{Hình }} $6b$};
	\end{tikzpicture}
	\end{center}
	\loigiai{
	\begin{itemize}
		\item Quan sát đồ thị \textit{Hình} $6a$, ta thấy hàm số $y=f(x)$ đồng biến trên các khoảng $(-\infty;-1)$, $(0;1)$, $(2;+\infty)$; nghịch biến trên các khoảng $(-1;0)$, $(1;2)$.\\
		Hàm số đạt cực đại tại $x=-1$, $x=1$ và đạt cực tiểu tại $x=0$, $x=2$.
		\item Quan sát đồ thị \textit{Hình} $6b$, ta thấy hàm số $y=g(x)$ đồng biến trên các khoảng $(-2;0)$, $(1;+\infty)$; nghịch biến trên các khoảng $(-\infty;-2)$, $(0;1)$.\\
		Hàm số đạt cực đại tại $x=0$ và đạt cực tiểu tại $x=-2$, $x=1$.
	\end{itemize} 
	}
\end{bt}

\begin{bt}%[2D1V1-5]
% nguon: Archive/2024/0. [LVH] BG - 12 New - 4in1-20240625T224301Z-001/0. [LVH] BG - 12 New - 4in1/2. [LVH] TEX SGK 12 NEW-20240418T074833Z-001/2. [LVH] TEX SGK 12 NEW/[LVH] CD-12new/SGK-12-CD/data/12-SGK-CD-1-1.tex #1075

	Thể tích $V$ (đơn vị: centimét khối) của $1$ kg nước tại nhiệt độ $T\left(0^{\circ}\mathrm{C}\leq T \leq 30^{\circ}\mathrm{C}\right)$ được tính bởi công thức sau:
	$$V(T)=999{,}87-0{,}06426T+0{,}0085043T^2-0{,}0000679T^3.$$
	\centerline{
	(Nguồn: J. Srewart, Calculux, Sevenuh Ediriom, Brooks/Cole, CENGAGE Leaming 2012)}\\
	Hỏi thể tích $V(T), 0^{\circ}\mathrm{C}\leq T \leq 30^{\circ}\mathrm{C}$, giảm trong khoảng nhiệt độ nào?
	\loigiai{
	Xét hàm số  $V(T)=999{,}87-0{,}06426T+0{,}0085043T^2-0{,}0000679T^3$, với $T\in [0;30]$.\\
	Ta có $V'(T)=-0{,}0002037T^2+0{,}0170086T-0{,}06426$.\\
	$V'(T)=0\Leftrightarrow T=3{,}966514624=T_1$ hoặc $T=79{,}53176716\not\in [0;30]$.\\
	Bảng biến thiên của hàm số $V(T)$ như sau
	\begin{center}
	\begin{tikzpicture}[font=\footnotesize,thick,>=stealth]
		\tikzset{double style/.append style={double distance=1.5pt}}
		\tkzTabInit[nocadre=false,lgt=1.2,espcl=2.5,deltacl=0.6,lw=.75pt,color,colorL=green!50,colorV=green!50]
		{$T$ /0.7, $V'(T)$ /0.8, $V(T)$ /2}
		{$0$,$T_1$,$30$}
		\tkzTabLine{ ,-,$0$,+, }
		\tkzTabVar{+/$V(0)$,-/$V(T_1)$,+/$V(30)$}
	\end{tikzpicture}
	\end{center}
	Từ bảng biến thiên suy ra, thể tích $V(T), 0^{\circ}\mathrm{C}\leq T \leq 30^{\circ}\mathrm{C}$, giảm trong khoảng nhiệt độ từ $0^\circ$C đến $3{,}966514624^\circ$C.
	}
\end{bt}

\begin{bt}%[2D1V5-1]
% nguon: Archive/2024/0. [LVH] BG - 12 New - 4in1-20240625T224301Z-001/0. [LVH] BG - 12 New - 4in1/2. [LVH] TEX SGK 12 NEW-20240418T074833Z-001/2. [LVH] TEX SGK 12 NEW/[LVH] SGK-12-KNTT/data/C1/12-SGK-KNTT-1-4.tex #1077

	Khảo sát sự biến thiên và vẽ đồ thị của hàm số $y=\dfrac{x+1}{x-2}$.
	\loigiai{
		\begin{enumerate}
			\item Tập xác định của hàm số: $\mathbb{R} \setminus\{2\}$.
			
			\item Sự biến thiên
			\begin{itemize}
				\item Ta có: $y^{\prime}=-\dfrac{3}{(x-2)^{2}}<0$ với mọi $x \neq 2$.
				\item Hàm số nghịch biến trên từng khoảng $(-\infty ; 2)$ và $(2 ;+\infty)$.
				\item Hàm số không có cực trị.
				\item Tiệm cận
				$\begin{aligned}[t]
					&\lim _{x \rightarrow 2^{-}} y=\lim _{x \rightarrow 2^{-}} \dfrac{x+1}{x-2}=-\infty ; \lim _{x \rightarrow 2^{+}} y=\lim _{x \rightarrow 2^{+}} \dfrac{x+1}{x-2}=+\infty;\\
					&\lim _{x \rightarrow+\infty} y=\lim _{x \rightarrow+\infty} \dfrac{x+1}{x-2}=1 ; \lim _{x \rightarrow-\infty} y=\lim _{x \rightarrow-\infty} \dfrac{x+1}{x-2}=1.
				\end{aligned}
				$
				Do đó, đồ thị của hàm số có tiệm cận đứng là đường thẳng $x=2$, tiệm cận ngang là đường thẳng $y=1$.
				\item  Bảng biến thiên:
				\begin{center}
					\begin{tikzpicture}[line join=round,font=\footnotesize,>=stealth, scale=1]
						\tkzTabInit[nocadre=false,lgt=1.2,espcl=2.5,deltacl=0.6]
						{$x$ /0.6,$y'$ /0.6,$y$ /2}
						{$-\infty$, $2$, $+\infty$}
						\tkzTabLine{,-,d,-}
						\tkzTabVar{+/$1$, -D+/$-\infty$/$+\infty$, -/$1$}
					\end{tikzpicture}
				\end{center}
			\end{itemize}
			\item Đồ thị
			\immini{
				\begin{itemize}
					\item Giao điểm của đồ thị hàm số với trục tung là điểm $\left(0 ;-\dfrac{1}{2}\right)$.
					\item Giao điểm của đồ thị hàm số với trục hoành là điểm $(-1 ; 0)$.
					\item Đồ thị hàm số nhận giao điểm $(2 ; 1)$ của hai đường tiệm cận làm tâm đối xứng và nhận hai đường phân giác của các góc tạo bởi hai đường tiệm cận này làm trục đối xứng.
				\end{itemize}
			}{
				\begin{tikzpicture}[line cap=round, line join=round,font=\footnotesize,>=stealth, scale=1,x=0.4cm, y=0.4cm]
					\tikzset{label style/.style={font=\footnotesize}}
					\draw[->] (-6.5,0)--(9,0) node[below]{$x$};
					\draw[->](0,-5)--(0,7) node[left]{$y$};
					\foreach\x in {-6,-4,-2,-1,4,6,8} \draw[fill=black] (\x,0) circle (1pt) node[below]{$\x$};
					\foreach\y in {-4,-2,2,4,6}  \draw[fill=black] (0,\y) circle (1pt) node[left]{$\y$};
					\draw[smooth, samples=100] plot[domain=-6:1.5] (\x, { ((\x)+1)/((\x)-2) });
					\draw[smooth, samples=100] plot[domain=2.5:8] (\x, { ((\x)+1)/((\x)-2) });
					\draw (2,-5)--(2,7) node[pos=0.25,rotate=90, below left]{$x=2$}
					(-6.5,1) node[above right]{$y=1$}--(9,1) 
					(2,1) node[above right]{$I$}
					(0,0)+(0.5,0.4) node{$\tiny O$};
					\draw[fill=black] (2,0) circle (1pt) node[below right]{$2$}
					(0,1) circle (1pt) node[above right]{$1$};
					\draw (1,-6) node{\textit{\textbf{Hình 1.30}}};
				\end{tikzpicture}
			}
		\end{enumerate}
	}
	\begin{note}
		Đồ thị của hàm số phân thức $y=\dfrac{a x+b}{c x+d}(c \neq 0, a d-b c \neq 0)$ :
		\begin{itemize}
			\item Nhận giao điểm của tiệm cận đứng và tiệm cận ngang làm tâm đối xứng;
			\item Nhận hai đường phân giác của các góc tạo bởi hai đường tiệm cận này làm các trục đối xứng.
		\end{itemize}
	\end{note}
\end{bt}

\begin{bt}%[2D1V5-1]
% nguon: Archive/2024/0. [LVH] BG - 12 New - 4in1-20240625T224301Z-001/0. [LVH] BG - 12 New - 4in1/2. [LVH] TEX SGK 12 NEW-20240418T074833Z-001/2. [LVH] TEX SGK 12 NEW/[LVH] SGK-12-KNTT/data/C1/12-SGK-KNTT-1-4.tex #1078

	Khảo sát và vẽ đồ thị của hàm số $y=\dfrac{x^{2}-x-1}{x-2}$.
	\loigiai{
		\begin{enumerate}
			\item Tập xác định của hàm số là $\mathbb{R}\setminus\{2\}$.
			\item Sự biết thiên: Viết $y=x+1+\dfrac{1}{x-2}$.
			\begin{itemize}
				\item Ta có: $y^{\prime}=1-\dfrac{1}{(x-2)^{2}}=\dfrac{x^{2}-4 x+3}{(x-2)^{2}}$. Vậy $y^{\prime}=0 \Leftrightarrow \dfrac{x^{2}-4 x+3}{(x-2)^{2}}=0 \Leftrightarrow x=1$ hoặc $x=3$.
				\item Trên các khoảng $(-\infty ; 1)$ và $(3 ;+\infty), y^{\prime}>0$ nên hàm số đồng biến trên từng khoảng này.
				
				Trên các khoảng $(1 ; 2)$ và $(2 ; 3), y^{\prime}<0$ nên hàm số nghịch biến trên từng khoảng này.
				\item Hàm số đạt cực đại tại $x=1$ với $y_{C D}=1$; hàm số đạt cực tiểu tại $x=3$ với $y_{C T}=5$.
				\item $\lim\limits _{x \rightarrow-\infty} y=\lim\limits_{x \rightarrow-\infty} \dfrac{x^{2}-x-1}{x-2}=\lim\limits _{x \rightarrow-\infty} \dfrac{x-1-\dfrac{1}{x}}{1-\dfrac{2}{x}}=-\infty$;\\
				$\lim \limits_{x \rightarrow+\infty} y=\lim \limits_{x \rightarrow+\infty} \dfrac{x^{2}-x-1}{x-2}=\lim\limits _{x \rightarrow+\infty} \dfrac{x-1-\dfrac{1}{x}}{1-\dfrac{2}{x}}=+\infty$.
				\item Tiệm cận
				\begin{align*}
					&\lim _{x \rightarrow 2^{-}} y=\lim _{x \rightarrow 2^{-}}\left(x+1+\dfrac{1}{x-2}\right)=-\infty ; \lim _{x \rightarrow 2^{+}} y=\lim _{x \rightarrow 2^{+}}\left(x+1+\dfrac{1}{x-2}\right)=+\infty;\\
					&\lim _{x \rightarrow+\infty}[y-(x+1)]=\lim _{x \rightarrow+\infty} \dfrac{1}{x-2}=0 ; \lim _{x \rightarrow-\infty}[y-(x+1)]=\lim _{x \rightarrow-\infty} \dfrac{1}{x-2}=0
				\end{align*}
				Do đó, đồ thị hàm số có tiệm cận đứng là đường thẳng $x=2$, tiệm cận xiên là đường thẳng $y=x+1$.
				\item Bảng biến thiên
				\immini{
					\begin{center}
						\begin{tikzpicture}[line join=round,font=\footnotesize,>=stealth, scale=1]
							\tkzTabSetup[doubledistance=2pt]
							\tkzTabInit[nocadre=false,lgt=0.6,espcl=1.5,deltacl=0.6]
							{$x$ /0.6,$y'$ /0.6,$y$ /2}
							{$-\infty$, $1$, $2$, $3$, $+\infty$}
							\tkzTabLine{,+,0,-,d,-,0,+,}
							\tkzTabVar{-/$-\infty$, +/$1$, -D+/$-\infty$/$+\infty$, -/$5$, +/$+\infty$}
						\end{tikzpicture}
					\end{center}		
				}{
					\begin{tikzpicture}[line join=round, line cap=round,>=stealth,thick, scale=0.7]
						\tikzset{every node/.style={scale=0.9}}
						\draw[->] (-2.6,0)--(6.1,0) node[below left] {$x$};
						\draw[->] (0,-1.6)--(0,7.1) node[below left] {$y$};
						\draw (0,0) node [below left] {$O$};
						\foreach \x/\nx in {-2/-2,-1/ ,1/1,2/ ,3/3,4/4,5/5}
						\draw[thin] (\x,1pt)--(\x,-1pt) node [below] {$\nx$};
						\draw (-1,0) node [above]{$-1$}
						(2,0) node[below right]{$2$};
						\foreach \y/\ny in {-1/-1,1/1,2/2,3/3,4/4,5/5,6/6}
						\draw[thin] (1pt,\y)--(-1pt,\y) node [left] {$\ny$};
						\draw[dashed,thin](1,0)--(1,1)--(0,1);
						\draw[dashed,thin](3,0)--(3,5)--(0,5);
						\draw[dashed,thin] (2,-1.5)--(2,7) node[pos=0.9,above,rotate=90]{$x=2$};
						\begin{scope}
							\clip (-2.5,-1.5) rectangle (6,7);
							\draw[samples=200,domain=-2.1:1.99,smooth,variable=\x] plot (\x,{(1*((\x)^2)+-1*(\x)+-1)/(1*(\x)+-2)});
							\draw[samples=200,domain=2.01:6,smooth,variable=\x] plot (\x,{(1*((\x)^2)+-1*(\x)+-1)/(1*(\x)+-2)});
							\draw[dashed,thin] (-2.2,-1.2)--(6.1,7.1) node[pos=0.8,below,rotate=45]{$y=x+1$};
						\end{scope}
						\draw (2,-2) node{\textbf{\textit{Hình 1.31}}};
					\end{tikzpicture}
				}
			\end{itemize}
			\item Đồ thị (H.1.31):
			\begin{itemize}
				\item Giao điểm của đồ thị hàm số với trục tung là điểm $\left(0 ; \dfrac{1}{2}\right)$.
				\item Ta có $y=0 \Leftrightarrow \dfrac{x^{2}-x-1}{x-2}=0 \Leftrightarrow x=\dfrac{1-\sqrt{5}}{2}$ hoặc $x=\dfrac{1+\sqrt{5}}{2}$. Do đó giao điểm của đồ thị hàm số với trục hoành là các điểm $\left(\dfrac{1-\sqrt{5}}{2} ; 0\right)$ và $\left(\dfrac{1+\sqrt{5}}{2} ; 0\right)$.
				\item Đồ thị hàm số nhận giao điểm $I(2 ; 3)$ của hai đường tiệm cận làm tâm đối xứng và nhận hai đường phân giác của các góc tạo bởi hai đường tiệm cận này làm các trục đối xứng.
			\end{itemize}
		\end{enumerate}
	}
\end{bt}

\begin{bt}%[2D1V5-1]
% nguon: Archive/2024/0. [LVH] BG - 12 New - 4in1-20240625T224301Z-001/0. [LVH] BG - 12 New - 4in1/2. [LVH] TEX SGK 12 NEW-20240418T074833Z-001/2. [LVH] TEX SGK 12 NEW/[LVH] SGK-12-KNTT/data/C1/12-SGK-KNTT-1-4.tex #1079

	Khảo sát sự biến thiên và vẽ đồ thị của hàm số $y=\dfrac{x^{2}+x-2}{x+1}$.
	\loigiai{
		\begin{enumerate}
			\item Tập xác định của hàm số: $\mathbb{R} \setminus\{-1\}$.
			
			\item Sự biến thiên:
			\begin{itemize}
				\item Viết $y=x-\dfrac{2}{x+1}$, ta có $y^{\prime}=1+\dfrac{2}{(x+1)^{2}}>0$ với mọi $x \neq-1$.
				\item Hàm số đồng biến trên từng khoảng $(-\infty ;-1)$ và $(-1 ;+\infty)$.
				\item Hàm số không có cực trị.
				\item $\lim\limits _{x \rightarrow-\infty} y=\lim\limits _{x \rightarrow-\infty} \dfrac{x^{2}+x-2}{x+1}=\lim\limits _{x \rightarrow-\infty} \dfrac{x+1-\dfrac{2}{x}}{1+\dfrac{1}{x}}=-\infty$;\\ $\lim\limits _{x \rightarrow+\infty} y=\lim\limits _{x \rightarrow+\infty} \dfrac{x^{2}+x-2}{x+1}=\lim\limits _{x \rightarrow+\infty} \dfrac{x+1-\dfrac{2}{x}}{1+\dfrac{1}{x}}=+\infty$.
				\item Tiệm cận:
				\begin{align*}
					&\lim\limits _{x \rightarrow-1^{-}} y=\lim\limits _{x \rightarrow-1^{-}}\left(x-\dfrac{2}{x+1}\right)=+\infty ; \lim\limits _{x \rightarrow-1^{+}} y=\lim\limits _{x \rightarrow-1^{+}}\left(x-\dfrac{2}{x+1}\right)=-\infty;\\
					&\lim\limits _{x \rightarrow+\infty}[y-x]=\lim\limits _{x \rightarrow+\infty}\left(-\dfrac{2}{x+1}\right)=0 ; \lim\limits _{x \rightarrow-\infty}[y-x]=\lim\limits _{x \rightarrow-\infty}\left(-\dfrac{2}{x+1}\right)=0.
				\end{align*}
				Do đó, đồ thị hàm số có tiệm cận đứng là đường thẳng $x=-1$, tiệm cận xiên là đường thẳng $y=x$.
				\item Bảng biến thiên:
				\immini{
					\begin{tikzpicture}[line join=round,font=\footnotesize,>=stealth, scale=1]
						\tkzTabInit[nocadre=false,lgt=1.2,espcl=2.5,deltacl=0.6]
						{$x$ /0.6,$y'$ /0.6,$y$ /2}
						{$-\infty$, $-1$, $+\infty$}
						\tkzTabLine{,+,d,+,}
						\tkzTabVar{-/$-\infty$, +D-/$+\infty$/$-\infty$, +/$+\infty$}
					\end{tikzpicture}
				}{
					\begin{tikzpicture}[line cap=round, line join=round,font=\footnotesize,>=stealth, scale=1,x=0.5cm, y=0.5cm]
						\tikzset{label style/.style={font=\footnotesize}}
						\draw[->] (-6.1,0)--(6.1,0) node[below] {$x$};
						\draw[->] (0,-5.1)--(0,5.1) node[left] {$y$};
						\draw (0,0) node [above left] {\scriptsize$O$};
						\foreach \x/\nx in {-6/-6,-5/-5,-4/-4,-3/-3,-2/-2,-1/ ,1/1,2/2,3/3,4/4,5/5}
						\draw[thin] (\x,1pt)--(\x,-1pt) node [below] {\scriptsize$\nx$};
						\foreach \y/\ny in {-4/right,-3/right,-2/right,-1/left,1/left,2/left,3/left,4/left}
						\draw[thin] (1pt,\y)--(-1pt,\y) node [\ny] {\scriptsize $\y$};
						\draw[dashed,thin] (-1,-5)--(-1,5) node[above,rotate=90,pos=0.1]{$x=-1$};
						\draw (-1,0) node[above left]{$-1$};
						\begin{scope}
							\clip (-6,-5) rectangle (6,5);
							\draw[samples=200,domain=-6:-1.01,smooth,variable=\x,thick,magenta] plot (\x,{(1*((\x)^2)+1*(\x)+-2)/(1*(\x)+1)});
							\draw[samples=200,domain=-0.99:6,smooth,variable=\x,thick,magenta] plot (\x,{(1*((\x)^2)+1*(\x)+-2)/(1*(\x)+1)});
							\draw[dashed,thin] (-6.1,-6.1)--(6.1,6.1) node[above, rotate=45,pos=0.8]{$y=x$};
						\end{scope}
						\draw (0,-5.5) node{\textbf{\textit{Hình 1.32}}};
					\end{tikzpicture}
				}
			\end{itemize}
			\begin{itemize}
				\item Giao điểm của đồ thị hàm số với trục tung là điểm $(0 ;-2)$.
				\item Ta có $y=0 \Leftrightarrow \dfrac{x^{2}+x-2}{x+1}=0 \Leftrightarrow x=-2$ hoặc $x=1$.
				
				Do đó giao điểm của đồ thị hàm số với trục hoành là các điểm $(-2 ; 0)$ và $(1 ; 0)$.
				\item Đồ thị hàm số nhận giao điểm $I(-1 ;-1)$ của hai đường tiệm cận làm tâm đối xứng và nhận hai đường phân giác của góc tạo bởi hai đường tiệm cận này làm các trục đối xứng.
			\end{itemize}
		\end{enumerate}
		
	}
\end{bt}

\begin{bt}%[2D1N3-1]
% nguon: Archive/2024/0. [LVH] BG - 12 New - 4in1-20240625T224301Z-001/0. [LVH] BG - 12 New - 4in1/2. [LVH] TEX SGK 12 NEW-20240418T074833Z-001/2. [LVH] TEX SGK 12 NEW/[LVH] SGK-12-KNTT/data/C1/12-SGK-KNTT-1-2.tex #1080

	Tìm giá trị lớn nhất và giá trị nhỏ nhất của hàm số $y=f(x)=\sqrt{1-x^2}$.
	\loigiai{
		Tập xác định của hàm số là $[-1; 1]$.
		\begin{itemize}
			\item Cách $1$. Sử dụng định nghĩa.\\
			Ta có
			\begin{itemize}
				\item $f(x)=\sqrt{1-x^2} \geq 0$; dấu bằng xảy ra khi $1-x^2=0$, tức là khi $x=-1$ hoặc $x=1$.
			\end{itemize}
			Do đó $\min\limits_{[-1; 1]} f(x)=f(-1)=f(1)=0$.
			\begin{itemize}
				\item $f(x)=\sqrt{1-x^2} \leq 1$; dấu bằng xảy ra khi $1-x^2=1$, tức là khi $x=0$. Do đó $\max\limits_{[-1; 1]} f(x)=f(0)=1$.
			\end{itemize}
			\item Cách $2$. Sử dụng bảng biến thiên.\\
			Với $x \in(-1; 1)$, ta có $y'=\dfrac{\left(1-x^2\right)'}{2 \sqrt{1-x^2}}=-\dfrac{x}{\sqrt{1-x^2}}$; $y'=0 \Leftrightarrow x=0$.\\
			Lập bảng biến thiên của hàm số trên đoạn $[-1; 1]$:
			\begin{center}
				\begin{tikzpicture}[>=stealth,line join=round,line cap=round,font=\footnotesize,scale=1]
					\tkzTabInit[nocadre=false,lgt=1.2,espcl=3,deltacl=.55]
					{$x$/0.7, $y'$/0.7, $y$/1.8}
					{$-1$,$0$,$1$}
					\tkzTabLine{d,+,$0$,-,d}
					\tkzTabVar{-/$0$ ,+/ $1$ ,-/$0$}	
				\end{tikzpicture}
			\end{center}
			Từ bảng biến thiên, ta được $\min\limits_{[-1; 1]} f(x)=f(-1)=f(1)=0$; $\max\limits_{[-1; 1]} f(x)=f(0)=1$.		
		\end{itemize}			
	}
\end{bt}

\begin{bt}%[2D1N3-2]
% nguon: Archive/2024/0. [LVH] BG - 12 New - 4in1-20240625T224301Z-001/0. [LVH] BG - 12 New - 4in1/2. [LVH] TEX SGK 12 NEW-20240418T074833Z-001/2. [LVH] TEX SGK 12 NEW/[LVH] SGK-12-KNTT/data/C1/12-SGK-KNTT-1-2.tex #1081

	Tìm giá trị lớn nhất và giá trị nhỏ nhất (nếu có) của hàm số $y=x-2+\dfrac{1}{x}$ trên khoảng $(0;+\infty)$.	
	\loigiai{
		Ta có $y'=1-\dfrac{1}{x^2}; y'=0 \Leftrightarrow x=1$ (vì $x>0$).\\
		Tính các giới hạn
		$
		\lim\limits_{x \to 0^+} y=\lim\limits_{x \to 0^+}\left(x-2+\dfrac{1}{x}\right)=+\infty$; $\lim\limits_{x \to+\infty} y=\lim\limits_{x \to+\infty}\left(x-2+\dfrac{1}{x}\right)=+\infty
		$.\\
		Lập bảng biến thiên của hàm số trên khoảng $(0;+\infty)$ :
		\begin{center}
			\begin{tikzpicture}[>=stealth,line join=round,line cap=round,font=\footnotesize,scale=1]
				\tkzTabInit[nocadre=false,lgt=1.2,espcl=3,deltacl=.55]
				{$x$/0.7, $y'$/0.7, $y$/1.8}
				{$-\infty$,$1$,$+\infty$}
				\tkzTabLine{,-,$0$,+,}
				\tkzTabVar{+/$+\infty$ ,-/ $0$ ,+/$+\infty$}	
			\end{tikzpicture}	
		\end{center}
		Từ bảng biến thiên, ta được $\min\limits_{(0;+\infty)}y=y(1)=0$; hàm số không có giá trị lớn nhất trên khoảng $(0;+\infty)$.		
	}
\end{bt}

\begin{bt}%[2D1H3-6]
% nguon: Archive/2024/0. [LVH] BG - 12 New - 4in1-20240625T224301Z-001/0. [LVH] BG - 12 New - 4in1/2. [LVH] TEX SGK 12 NEW-20240418T074833Z-001/2. [LVH] TEX SGK 12 NEW/[LVH] SGK-12-KNTT/data/C1/12-SGK-KNTT-1-2.tex #1082

	\immini{
		Từ một tấm bìa carton hình vuông có độ dài cạnh bằng $60$ cm, người ta cắt bốn hình vuông bằng nhau ở bốn góc rồi gập thành một chiếc hộp có dạng hình hộp chữ nhật không có nắp (H.1.14). Tính cạnh của các hình vuông bị cắt sao cho thể tích của chiếc hộp là lớn nhất.	
	}{\begin{tabular}{ccc}
			\rule[-1ex]{0pt}{2.5ex}	\begin{tikzpicture}[>=stealth,line join=round,line cap=round,font=\footnotesize,scale=.6]
				\filldraw[pattern=north east lines] 
				(0,0)rectangle(1,1)(0,4)rectangle(1,5)
				(4,0)rectangle(5,1)(4,4)rectangle(5,5)
				;
				\draw[dashed](1,1)rectangle(4,4);
				\draw(0,0)rectangle(5,5);
				\draw[dashed](5,5)arc(45:-45:{5*1.41/2}) node[midway,fill=white]{$60$ cm};
			\end{tikzpicture}& & \begin{tikzpicture}[>=stealth,line join=round,line cap=round,font=\footnotesize,scale=.4]	
				\tikzset{
					pics/hhChuNhat/.style n args={8}{
						code={
							\tikzset{
								declare function={a=2.5;b=1.8;goc=-135;h=0.75;}
							}
							\path 
							(0,0)coordinate(#1)+(0:a)coordinate(#2)+(goc:b)coordinate(#4)+(90:h)coordinate(#5)
							($(#2)+(#4)-(#1)$)coordinate(#3)
							;
							\foreach \pone/\pname in{#2/#6,#3/#7,#4/#8}{
								\path 
								($(\pone)+(#5)-(#1)$)coordinate(\pname)
								;
							}
							\foreach \pointo/\pointt in{#2/#3,#3/#4,#5/#6,#6/#7,#7/#8,#8/#5,#2/#6,#3/#7,#4/#8}{
								\draw[fill=black](\pointo)--(\pointt);
							}
						}
					}
				}
				\path 
				(0,0) pic{hhChuNhat={A}{B}{C}{D}{A'}{B'}{C'}{D'}}
				(intersection of A--D and C'--D')coordinate(M)
				(intersection of A--B and B'--C')coordinate(N)
				;
				\foreach \pointo/\pointt in{A/M,A/N,A/A'}{
					\draw[fill=black](\pointo)--(\pointt);
				}
				\foreach \pointo/\pointt in{D/M,B/N}{
					\draw[fill=black,dashed](\pointo)--(\pointt);
				}
			\end{tikzpicture}\\
			\multicolumn{3}{c}{\color{orange}Hình 1.14}\\
		\end{tabular}
	}	
	\loigiai{
		Gọi $x$ (cm) là độ dài cạnh của các hình vuông nhỏ được cắt ở bốn góc của tấm bìa. Điều kiện $0<x<30$.\\
		Khi cắt bỏ bốn hình vuông nhỏ có cạnh $x$ cm ở bốn góc và gập lên thì ta được một chiếc hộp chữ nhật không có nắp, có đáy là hình vuông với độ dài cạnh bằng $(60-2 x)$ (cm) và chiều cao bằng $x$ cm. Thể tích của chiếc hộp này là
		\[
		V(x)=(60-2 x)^2 \cdot x=4 x^3-240 x^2+3600 x\left(\text{cm}^3\right).
		\]
		Ta có $V'(x)=12 x^2-480 x+3600$;  $V'(x)=0 \Leftrightarrow x^2-40 x+300=0 \Leftrightarrow x=10$ (thoả mãn điều kiện) hoặc $x=30$ (loại).\\
		Lập bảng biến thiên:
		\begin{center}
			\begin{tikzpicture}[>=stealth,line join=round,line cap=round,font=\footnotesize,scale=1]
				\tkzTabInit[nocadre=false,lgt=1.2,espcl=3,deltacl=.55]
				{$x$/0.7, $V'(x)$/0.7, $V(x)$/1.8}
				{$0$,$10$,$30$}
				\tkzTabLine{,+,$0$,-,}
				\tkzTabVar{-/$0$ ,+/ $16\,000$ ,-/$0$}	
			\end{tikzpicture}
		\end{center}
		Vậy để thể tích của chiếc hộp là lớn nhất thì độ dài cạnh của các hình vuông nhỏ phải cắt là $10$ cm.		
	}
\end{bt}

\begin{bt}%[2D1N3-1]
% nguon: Archive/2024/0. [LVH] BG - 12 New - 4in1-20240625T224301Z-001/0. [LVH] BG - 12 New - 4in1/2. [LVH] TEX SGK 12 NEW-20240418T074833Z-001/2. [LVH] TEX SGK 12 NEW/[LVH] SGK-12-KNTT/data/C1/12-SGK-KNTT-1-2.tex #1083

	Tìm giá trị lớn nhất và giá trị nhỏ nhất của hàm số $y=x^4-4x^2+3$ trên đoạn $[0; 4]$.
	\loigiai{
		Ta có $y^{\prime}=4x^3-8x=4x\left(x^2-2\right); y^{\prime}=0\Leftrightarrow x=0$ hoặc $x=\sqrt{2}$ (vì $x \in[0; 4]$);
		\[
		y(0)=3;~y(4)=195;~y\left(\sqrt{2}\right)=-1.
		\]
		Do đó: $\max \limits_{[0; 4]} y=y(4)=195; \min \limits_{[0; 4]} y=y\left(\sqrt{2}\right)=-1$.
	}
\end{bt}

\begin{bt}%[2D1H3-1]
% nguon: Archive/2024/0. [LVH] BG - 12 New - 4in1-20240625T224301Z-001/0. [LVH] BG - 12 New - 4in1/2. [LVH] TEX SGK 12 NEW-20240418T074833Z-001/2. [LVH] TEX SGK 12 NEW/[LVH] SGK-12-KNTT/data/C1/12-SGK-KNTT-1-2.tex #1084

	Tìm giá trị lớn nhất và giá trị nhỏ nhất của hàm số $y=\sin x+\cos x$ trên đoạn $[0; 2\pi]$.
	\loigiai{
		Ta có $y^{\prime}=\cos x-\sin x; y^{\prime}=0\Leftrightarrow \cos x=\sin x \Leftrightarrow x=\dfrac{\pi}{4}$ hoặc $x=\dfrac{5\pi}{4}$ (vì $x \in[0; 2\pi]$);
		\[
		y(0)=1; y(2 \pi)=1; y\left(\dfrac{\pi}{4}\right)=\sqrt{2}; y\left(\dfrac{5 \pi}{4}\right)=-\sqrt{2}.
		\]
		Do đó: $\max \limits_{[0; 2\pi]} y=y\left(\dfrac{\pi}{4}\right)=\sqrt{2}; \min \limits_{[0; 2\pi]} y=y\left(\dfrac{5 \pi}{4}\right)=-\sqrt{2}$.
	}
\end{bt}

\begin{bt}%[2D1H3-6]
% nguon: Archive/2024/0. [LVH] BG - 12 New - 4in1-20240625T224301Z-001/0. [LVH] BG - 12 New - 4in1/2. [LVH] TEX SGK 12 NEW-20240418T074833Z-001/2. [LVH] TEX SGK 12 NEW/[LVH] SGK-12-KNTT/data/C1/12-SGK-KNTT-1-2.tex #1085

	Giả sử sự lây lan của một loại virus ở một địa phương có thể được mô hình hoá bằng hàm số $N(t)=-t^3+12t^2$, $0\leq t \leq 12$, trong đó $N$ là số người bị nhiễm bệnh (tính bằng trăm người) và $t$ là thời gian (tuần).
	\begin{enumerate}
		\item Hãy ước tính số người tối đa bị nhiễm bệnh ở địa phương đó.
		\item Đạo hàm $N^{\prime}(t)$ biểu thị tốc độ lây lan của virus (còn gọi là tốc độ truyền bệnh). Hỏi virus sẽ lây lan nhậnh nhất khi nào?
	\end{enumerate}
	\loigiai{
		\begin{enumerate}
			\item Với $0\leq t \leq 12$ ta có\\
			$N^{\prime}(t)=-3t^2+24t, N^{\prime}(t)=0\Leftrightarrow-3t^2+24t=0\Leftrightarrow\hoac{&t=0\\
				&t=8.
				&}$\\
			Ta có $N(0)=0$, $N(8)=256$, $N(12)=0$.\\
			Do đó, số người tối đa bị nhiễm bệnh ở địa phương là $256$ người trong $12$ tuần đầu.
			\item Hàm số biểu thị tốc độ độ lây lan của virus là $N^{\prime}(t)=-3t^2+24t$.\\
			Đặt $f(t)=-3t^2+24t$, với $0\leq t \leq 12$.\\
			Ta có $f'(t)=-6t+24, f'(t)=0\Leftrightarrow t=4$.
			\[
			f(0)=0, f(4)=48, f(12)=-144.
			\]
			Do đó, virus sẽ lây lan nhanh nhất khi $t=4$ (tuần thứ $4$).
		\end{enumerate}
	}
\end{bt}

\begin{bt}%[2D1N3-2]
% nguon: Archive/2024/0. [LVH] BG - 12 New - 4in1-20240625T224301Z-001/0. [LVH] BG - 12 New - 4in1/2. [LVH] TEX SGK 12 NEW-20240418T074833Z-001/2. [LVH] TEX SGK 12 NEW/[LVH] SGK-12-KNTT/data/C1/12-SGK-KNTT-1-2.tex #1086

	Tìm giá trị lớn nhất và giá trị nhỏ nhất (nếu có) của các hàm số sau:
	\begin{enumEX}{2}
		\item  $y=-x^2+4 x+3$;
		\item  $y=x^3-2 x^2+1$ trên $[0;+\infty)$;
		\item  $y=\dfrac{x^2-2 x+3}{x-1}$ trên $(1;+\infty)$;
		\item  $y=\sqrt{4 x-2 x^2}$.			
	\end{enumEX}
	\loigiai{
		\begin{enumerate}
			\item Tập xác định của hàm số là $\mathbb{R}$.\\
			Ta có $y=-x^2+4 x+3\Rightarrow y'=-2x+4$.\\
			Với $y'=0\Leftrightarrow x=2$.\\
			Lập bảng biến thiên của hàm số
			\begin{center}
				\begin{tikzpicture}[>=stealth,line join=round,line cap=round,font=\footnotesize,scale=1]
					\tkzTabInit[nocadre=false,lgt=1.2,espcl=3,deltacl=.55]
					{$x$/0.7, $y'$/0.7, $y$/1.8}
					{$-\infty$,$2$,$+\infty$}
					\tkzTabLine{,+,$0$,-,}
					\tkzTabVar{-/$-\infty$,+/$7$,-/$-\infty$}	
				\end{tikzpicture}
			\end{center}
			Từ bảng biến thiên, ta được $\max y(2)=7$.	
			\item Ta có $y=x^3-2 x^2+1\Rightarrow y'=3x^2-4x$.\\
			Với $y'=0\Leftrightarrow \hoac{&x=0\\&x=\dfrac{4}{3}.}$\\
			Lập bảng biến thiên của hàm số
			\begin{center}
				\begin{tikzpicture}[>=stealth,line join=round,line cap=round,font=\footnotesize,scale=1]
					\tkzTabInit[nocadre=false,lgt=1.2,espcl=3,deltacl=.55]
					{$x$/0.7, $y'$/0.7, $y$/1.8}
					{$0$,$\dfrac{4}{3}$,$+\infty$}
					\tkzTabLine{,-,$0$,+,}
					\path 
					($(N12)!0.25!(N13)$) coordinate (N12) 
					;
					\tkzTabVar{+/$1$,-/$-\dfrac{5}{27}$,+/$-\infty$}	
				\end{tikzpicture}
			\end{center}
			Từ bảng biến thiên, ta được $\min y=y\left(\dfrac{4}{3}\right)=-\dfrac{5}{27}$.	
			\item Ta có $y=\dfrac{x^2-2 x+3}{x-1}\Rightarrow y'=\dfrac{x^2-2x-1}{x^2-2x+1}$.\\
			Với $y'=0\Leftrightarrow \hoac{&x=1+\sqrt{2}\\&x=1-\sqrt{2}\ne (1;+\infty).}$\\
			Lại có $\lim\limits_{x \to 1^+} y=\lim\limits_{x \to 1^+}\left(\dfrac{x^2-2 x+3}{x-1}\right)=+\infty$; $\lim\limits_{x \to +\infty} y=\lim\limits_{x \to +\infty}\left(\dfrac{x^2-2 x+3}{x-1}\right)=+\infty$.\\
			Lập bảng biến thiên của hàm số
			\begin{center}
				\begin{tikzpicture}[>=stealth,line join=round,line cap=round,font=\footnotesize,scale=1]
					\tkzTabInit[nocadre=false,lgt=1.2,espcl=3,deltacl=.55]
					{$x$/0.7, $y'$/0.7, $y$/1.8}
					{$1$,$1+\sqrt{2}$,$+\infty$}
					\tkzTabLine{,-,$0$,+,}
					\tkzTabVar{+/$+\infty$,-/$2\sqrt{2}$,+/$-\infty$}	
				\end{tikzpicture}
			\end{center}
			Từ bảng biến thiên, ta được $\min y=y\left(1+\sqrt{2}\right)=2\sqrt{2}$.	
			\item Tập xác định của hàm số là $[0;2]$.\\
			Ta có $y=\sqrt{4 x-2 x^2}\Rightarrow y'=\dfrac{-2x+2}{\sqrt{4x-2x^2}}$.\\
			Với $y'=0\Leftrightarrow x=1$.\\
			Lập bảng biến thiên của hàm số
			\begin{center}
				\begin{tikzpicture}[>=stealth,line join=round,line cap=round,font=\footnotesize,scale=1]
					\tkzTabInit[nocadre=false,lgt=1.2,espcl=3,deltacl=.55]
					{$x$/0.7, $y'$/0.7, $y$/1.8}
					{$0$,$1$,$2$}
					\tkzTabLine{d,+,$0$,-,d}
					\tkzTabVar{-/$0$,+/$\sqrt{2}$,-/$0$}	
				\end{tikzpicture}
			\end{center}
			Từ bảng biến thiên, ta được $\max y=y(1)=\sqrt{2}$.			
		\end{enumerate}		
	}
\end{bt}

\begin{bt}%[2D1N3-2]
% nguon: Archive/2024/0. [LVH] BG - 12 New - 4in1-20240625T224301Z-001/0. [LVH] BG - 12 New - 4in1/2. [LVH] TEX SGK 12 NEW-20240418T074833Z-001/2. [LVH] TEX SGK 12 NEW/[LVH] SGK-12-KNTT/data/C1/12-SGK-KNTT-1-2.tex #1087

	Tìm giá trị lớn nhất và giá trị nhỏ nhất (nếu có) của các hàm số sau:
	\begin{enumEX}{2}
		\item  $y=x^4-2 x^2+3$;
		\item  $y=x \mathrm{e}^{-x}$;
		\item  $y=x \ln x$;
		\item  $y=\sqrt{x-1}+\sqrt{3-x}$.		
	\end{enumEX}
	\loigiai{
		\begin{enumerate}
			\item Tập xác định của hàm số là $\mathbb{R}$.\\
			Ta có $y=x^4-2 x^2+3\Rightarrow y'=4x^3-4x$.\\
			Với $y'=0\Leftrightarrow \hoac{&x=0\\&x=1\\&x=-1.}$\\
			Lập bảng biến thiên của hàm số
			\begin{center}
				\begin{tikzpicture}[>=stealth,line join=round,line cap=round,font=\footnotesize,scale=1]
					\tkzTabInit[nocadre=false,lgt=1.2,espcl=3,deltacl=.55]
					{$x$/0.7, $y'$/0.7, $y$/1.8}
					{$-\infty$,$-1$,$0$,$1$,$+\infty$}
					\tkzTabLine{,-,$0$,+,$0$,-,$0$,+}
					\tkzTabVar{+/$-\infty$,-/$2$,+/$3$,-/$2$,+/$+\infty$}	
				\end{tikzpicture}
			\end{center}
			Từ bảng biến thiên, ta được $\min y=y(-1)=y(1)=2$.	
			\item Tập xác định của hàm số là $\mathbb{R}$.\\
			Ta có $y=x\mathrm{e}^{-x}\Rightarrow y'=\mathrm{e}^{-x}-x\mathrm{e}^{-x}$.\\
			Với $y'=0\Leftrightarrow x=1$.\\
			Lập bảng biến thiên của hàm số
			\begin{center}
				\begin{tikzpicture}[>=stealth,line join=round,line cap=round,font=\footnotesize,scale=1]
					\tkzTabInit[nocadre=false,lgt=1.2,espcl=3,deltacl=.55]
					{$x$/0.7, $y'$/0.7, $y$/2}
					{$-\infty$,$1$,$+\infty$}
					\tkzTabLine{,+,$0$,-,}
					\path 
					($(N32)!0.75!(N33)$) coordinate (N33) 
					;
					\tkzTabVar{-/$-\infty$,+/$\dfrac{1}{\mathrm{e}}$,-/$0$}	
				\end{tikzpicture}
			\end{center}
			Từ bảng biến thiên, ta được $\max y=(1)=\dfrac{1}{\mathrm{e}}$.	
			\item  Tập xác định của hàm số là $(0;+\infty)$.\\
			Ta có $y=x\ln x\Rightarrow y'=\ln x +1$.\\
			Với $y'=0\Leftrightarrow x=\dfrac{1}{e}$.\\
			Lập bảng biến thiên của hàm số
			\begin{center}
				\begin{tikzpicture}[>=stealth,line join=round,line cap=round,font=\footnotesize,scale=1]
					\tkzTabInit[nocadre=false,lgt=1.2,espcl=3,deltacl=.55]
					{$x$/0.7, $y'$/0.7, $y$/2}
					{$0$,$\dfrac{1}{e}$,$+\infty$}
					\tkzTabLine{,-,$0$,+}
					\path 
					($(N12)!0.25!(N13)$) coordinate (N12) 
					;
					\tkzTabVar{+/$0$,-/$-\dfrac{1}{\mathrm{e}}$,+/$+\infty$}	
				\end{tikzpicture}
			\end{center}
			Từ bảng biến thiên, ta được $\min y=y\left(\dfrac{1}{\mathrm{e}}\right)=-\dfrac{1}{\mathrm{e}}$.	
			\item  Tập xác định của hàm số là $[1;3]$.\\
			Ta có $y=\sqrt{x-1}+\sqrt{3-x}\Rightarrow y'=\dfrac{1}{2\sqrt{x-1}}-\dfrac{1}{2\sqrt{3-x}}$.\\
			Với $y'=0\Leftrightarrow x=2$.\\
			Lập bảng biến thiên của hàm số
			\begin{center}
				\begin{tikzpicture}[>=stealth,line join=round,line cap=round,font=\footnotesize,scale=1]
					\tkzTabInit[nocadre=false,lgt=1.2,espcl=3,deltacl=.55]
					{$x$/0.7, $y'$/0.7, $y$/2}
					{$1$,$2$,$3$}
					\tkzTabLine{d,+,$0$,-,d}
					\tkzTabVar{-/$\sqrt{2}$,+/$2$,-/$\sqrt{2}$}	
				\end{tikzpicture}
			\end{center}
			Từ bảng biến thiên, ta được $\min y=y(1)=\sqrt{2}$; $\max y=y(2)=2$.		
		\end{enumerate}		
	}
\end{bt}

\begin{bt}%[2D1H3-1]
% nguon: Archive/2024/0. [LVH] BG - 12 New - 4in1-20240625T224301Z-001/0. [LVH] BG - 12 New - 4in1/2. [LVH] TEX SGK 12 NEW-20240418T074833Z-001/2. [LVH] TEX SGK 12 NEW/[LVH] SGK-12-KNTT/data/C1/12-SGK-KNTT-1-2.tex #1088

	Tìm giá trị lớn nhất và giá trị nhỏ nhất của các hàm số sau:
	\begin{enumerate}
		\item $y=2x^3-6x+3$ trên đoạn $[-1; 2]$;
		\item $y=x^4-3x^2+2$ trên đoạn $[0; 3]$;
		\item $y=x-\sin 2x$ trên đoạn $[0; \pi]$;
		\item $y=\left(x^2-x\right) \mathrm{e}^x$ trên đoạn $[0; 1]$.
	\end{enumerate}
	\loigiai{
		\begin{enumerate}
			\item Ta có $y'=6x^2-6$; $y'=0\Leftrightarrow\hoac{&x=1 \\&x=-1.}$.\\
			\[
			y(-1)=7; y(1)=-1; y(2)=7.
			\]
			Do đó: $\max \limits_{[-1; 2]} y=y\left(-1\right)=y(2)=7; \min \limits_{[-1; 2]} y=y\left(1\right)=-1$.
			\item Ta có $y'=4x^3-6x$; $y'=0\Leftrightarrow\hoac{&x=0 \\&x=\dfrac{\sqrt{6}}{2}.}$ (vì $x\in[0; 3]$).\\
			\[
			y(0)=2; y\left(\dfrac{\sqrt{6}}{2}\right)=-\dfrac{1}{4}; y(3)=56.
			\]
			Do đó: $\max \limits_{[0; 3]} y=y\left(3\right)=56; \min \limits_{[0; 3]} y=y\left(\dfrac{\sqrt{6}}{2}\right)=-\dfrac{1}{4}$.
			\item Ta có $y'=1-2\cos2x$; $y'=0\Leftrightarrow\cos2x=\dfrac{1}{2}\Leftrightarrow x=\pm\dfrac{\pi}{6}+k\pi\Leftrightarrow\hoac{&x=\dfrac{\pi}{6} \\&x=\dfrac{5\pi}{6}.}$ (vì $x\in[0; \pi]$).\\
			\[
			y(0)=0;y\left(\dfrac{\pi}{6}\right)=\dfrac{\pi}{6}-\dfrac{\sqrt{3}}{2}; y\left(\dfrac{\pi}{6}\right)=\dfrac{5\pi}{6}+\dfrac{\sqrt{3}}{2}; y(\pi)=\pi.
			\]
			Do đó: $\max \limits_{[0; \pi]} y=y\left(\dfrac{\pi}{6}\right)=\dfrac{5\pi}{6}+\dfrac{\sqrt{3}}{2}; \min \limits_{[0; \pi]} y=y\left(\dfrac{\pi}{6}\right)=\dfrac{\pi}{6}-\dfrac{\sqrt{3}}{2}$.
			\item Ta có $y'=(2x-1)\mathrm{e}^x+\left(x^2-x\right) \mathrm{e}^x=\left(x^2+x-1\right) \mathrm{e}^x$; $y'=0\Leftrightarrow x=\dfrac{-1+\sqrt{5}}{2}$ (vì $x\in[0; 1]$).\\
			\[
			y(0)=0; y\left(\dfrac{-1+\sqrt{5}}{2}\right)=\left(2-\sqrt{5}\right)\mathrm{e}^{\frac{-1+\sqrt{5}}{2}}; y(1)=0.
			\]
			Do đó: $\max \limits_{[0; 1]} y=y\left(1\right)=y(0)=0; \min \limits_{[0; 1]} y\left(\dfrac{-1+\sqrt{5}}{2}\right)=\left(2-\sqrt{5}\right)\mathrm{e}^{\frac{-1+\sqrt{5}}{2}}$.
		\end{enumerate}
	}
\end{bt}

\begin{bt}%[2D1H3-6]
% nguon: Archive/2024/0. [LVH] BG - 12 New - 4in1-20240625T224301Z-001/0. [LVH] BG - 12 New - 4in1/2. [LVH] TEX SGK 12 NEW-20240418T074833Z-001/2. [LVH] TEX SGK 12 NEW/[LVH] SGK-12-KNTT/data/C1/12-SGK-KNTT-1-2.tex #1089

	Trong các hình chữ nhật có chu vi là $24\mathrm{~cm}$, hãy tìm hình chữ nhật có diện tích lớn nhất.
	\loigiai{
		Gọi $x$ (cm) là chiều rộng của hình chữ nhật.\\
		Điều kiện $0<x<12$.\\
		Chiều dài của hình chữ nhật là $12-x$ (cm).\\
		Diện tích của hình chữ nhật là $x(12-x)=-x^2+12x$ (cm$^2$).\\
		Xét hàm số $y=-x^2+12x$ với $0<x<12$.\\
		Ta có $y'=-2x+12$; $y'=0\Leftrightarrow x=6$.\\
		\[
		y(0)=0; y(6)=36; y(12)=0.
		\]
		Do đó: $\max \limits_{(0;12)} y=y(6)=36$.\\
		Trong các hình chữ nhật có chu vi là $24\mathrm{~cm}$ thì hình vuông có diện tích lớn nhất.
	}
\end{bt}

\begin{bt}%[2D1H3-6]
% nguon: Archive/2024/0. [LVH] BG - 12 New - 4in1-20240625T224301Z-001/0. [LVH] BG - 12 New - 4in1/2. [LVH] TEX SGK 12 NEW-20240418T074833Z-001/2. [LVH] TEX SGK 12 NEW/[LVH] SGK-12-KNTT/data/C1/12-SGK-KNTT-1-2.tex #1090

	\immini{
		Một nhà sản xuất muốn thiết kế một chiếc hộp có dạng hình hộp chữ nhật không có nắp, có đáy là hình vuông và diện tích bề mặt bằng $108\mathrm{~cm}^2$ như Hình 1.17. Tìm các kích thước của chiếc hộp sao cho thể tích của hộp là lớn nhất.
	}{
		\begin{tikzpicture}[>=stealth,line join=round,line cap=round,font=\footnotesize,scale=.4]	
			\tikzset{
				pics/hhChuNhat/.style n args={8}{
					code={
						\tikzset{
							declare function={a=2.5;b=1.8;goc=-135;h=0.75;}
						}
						\path 
						(0,0)coordinate(#1)+(0:a)coordinate(#2)+(goc:b)coordinate(#4)+(90:h)coordinate(#5)
						($(#2)+(#4)-(#1)$)coordinate(#3)
						;
						\foreach \pone/\pname in{#2/#6,#3/#7,#4/#8}{
							\path 
							($(\pone)+(#5)-(#1)$)coordinate(\pname)
							;
						}
						\foreach \pointo/\pointt in{#2/#3,#3/#4,#5/#6,#6/#7,#7/#8,#8/#5,#2/#6,#3/#7,#4/#8}{
							\draw[fill=black](\pointo)--(\pointt);
						}
					}
				}
			}
			\path 
			(0,0) pic{hhChuNhat={A}{B}{C}{D}{A'}{B'}{C'}{D'}}
			(intersection of A--D and C'--D')coordinate(M)
			(intersection of A--B and B'--C')coordinate(N)
			;
			\foreach \pointo/\pointt in{A/M,A/N,A/A'}{
				\draw[fill=black](\pointo)--(\pointt);
			}
			\foreach \pointo/\pointt in{D/M,B/N}{
				\draw[fill=black,dashed](\pointo)--(\pointt);
			}
			\path (C)--(D)node[midway, below]{$x$};
			\path (C)--(B)node[midway, right]{$x$};
			\path (B)--(B')node[midway, right]{$h$};
			\draw (current bounding box.south) node[below]{Hình 1.17};
		\end{tikzpicture}
	}
	\loigiai{
		Thể tích của chiếc hộp là $V=x^2\cdot h$ (cm$^3$).\\
		Vì diện tích bề mặt bằng $108$ cm$^2$ nên ta có 
		\begin{center}
			$x^2+4xh=108\Leftrightarrow h=\dfrac{108-x^2}{4x}$ (điều kiện $0< x < \sqrt{108}$).
		\end{center}
		Khi đó thể tích chiếc hộp là $V=\dfrac{x\left(108-x^2\right)}{4}=-\dfrac{x^3}{4}+27x$.\\
		Ta tìm giá trị lớn nhất của hàm số $V=-\dfrac{x^3}{4}+27 x~(0< x < \sqrt{108})$.\\
		Có $V^{\prime}=-\dfrac{3}{4} x^2+27; V^{\prime}=0\Leftrightarrow-\dfrac{3}{4} x^2+27=0\Leftrightarrow x=6$ (vì $0< x < \sqrt{108}$).\\
		Bảng biến thiên:
		\begin{center}
			\begin{tikzpicture}[scale=1, font=\footnotesize]
				\tkzTabInit[nocadre=false, lgt=1.2, espcl=2, deltacl=0.9]
				{$x$/0.8,$V'(x)$/0.6,$V(x)$/2}
				{$0$,$6$,$\sqrt{108}$};
				\tkzTabLine{,+,$0$,-,};
				\tkzTabVar{-/$0$,+/$108$,-/$V\left(\sqrt{108}\right)$};
			\end{tikzpicture}
		\end{center}
		Dựa vào bảng biến thiên, ta có thể tích lớn nhất của chiếc hộp là $108$ cm$^3$ khi $x=6$ cm và $h=3$ cm.
	}
\end{bt}

\begin{bt}%[2D1N4-1]
% nguon: Archive/2024/0. [LVH] BG - 12 New - 4in1-20240625T224301Z-001/0. [LVH] BG - 12 New - 4in1/2. [LVH] TEX SGK 12 NEW-20240418T074833Z-001/2. [LVH] TEX SGK 12 NEW/[LVH] SGK-12-KNTT/data/C1/12-SGK-KNTT-1-3.tex #1092

	Tìm tiệm cận ngang của đồ thị hàm số $y=f(x)=\dfrac{3 x-2}{x+1}$.
	\loigiai{
		Ta có: $\displaystyle\lim _{x \rightarrow+\infty} f(x)=\displaystyle\lim _{x \rightarrow+\infty} \dfrac{3 x-2}{x+1}=\displaystyle\lim _{x \rightarrow+\infty} \dfrac{3-\dfrac{2}{x}}{1+\dfrac{1}{x}}=3$. \\
		Tương tự, $\displaystyle\lim _{x \rightarrow-\infty} f(x)=3$.
		
		Vậy đồ thị hàm số $f(x)$ có tiệm cận ngang là đường thẳng $y=3$.
	}
\end{bt}

\begin{bt}%[2D1N4-1]
% nguon: Archive/2024/0. [LVH] BG - 12 New - 4in1-20240625T224301Z-001/0. [LVH] BG - 12 New - 4in1/2. [LVH] TEX SGK 12 NEW-20240418T074833Z-001/2. [LVH] TEX SGK 12 NEW/[LVH] SGK-12-KNTT/data/C1/12-SGK-KNTT-1-3.tex #1093

	Tìm các tiệm cận ngang của đồ thị hàm số $y=f(x)=\dfrac{\sqrt{x^{2}+1}}{x}$.
	\loigiai{
		\immini[thm]{
			Ta có:\\
			$\displaystyle\lim _{x \rightarrow+\infty} f(x)=\displaystyle\lim _{x \rightarrow+\infty} \dfrac{\sqrt{x^{2}+1}}{x}=\displaystyle\lim _{x \rightarrow+\infty} \sqrt{\dfrac{x^{2}+1}{x^{2}}}=\displaystyle\lim _{x \rightarrow+\infty} \sqrt{1+\dfrac{1}{x^{2}}}=1 ;$
			
			$\displaystyle\lim _{x \rightarrow-\infty} f(x)=\displaystyle\lim _{x \rightarrow-\infty} \dfrac{\sqrt{x^{2}+1}}{x}=\displaystyle\lim _{x \rightarrow-\infty}\left(-\sqrt{\dfrac{x^{2}+1}{x^{2}}}\right)=\displaystyle\lim _{x \rightarrow-\infty}\left(-\sqrt{1+\dfrac{1}{x^{2}}}\right)=-1$.\\
			Vậy đồ thị hàm số $f(x)$ có hai tiệm cận ngang là $y=1$ và $y=-1$. 	
		}{
			\begin{tikzpicture}[>=stealth,scale=.5,every node/.style={scale=0.8}]
				\draw[->] (-4.2,0) -- (4,0) node[below] {$x$};
				\draw[->] (0,-4) -- (0,3.8) node[below left] {$y$};
				\draw[thick,color=red] plot[domain=-4.2:-.25,samples=100] (\x,{sqrt(\x*\x+1)/\x});
				\draw[thick,color=red] plot[domain=.3:3.5,samples=100] (\x,{sqrt(\x*\x+1)/\x});
				\path 
				(0,-2)node[right]{$-2$}
				(0,-1)node[right]{$-1$}
				(0,1)node[below left]{$1$}
				(0,2)node[right]{$2$}
				;
				\draw[color=blue,thick] (-4.2,1)--(3.5,1);
				\foreach \x in {-3,-2,-1,1,2,3}{
					\path (\x,0)node[below]{$\x$};
					\draw[draw=black,fill=red] (\x,-2pt)--(\x,2pt)
					;}
				\foreach \y in {-3,-2,-1,1,2,3}{
					\draw[draw=black,fill=red] (-2pt,\y)--(2pt,\y)
					;}
			\end{tikzpicture}	
		}
		\noindent Nhận xét. Đồ thị hàm số $f(x)$ như hình vẽ.
	}
\end{bt}

\begin{bt}%[2D1N4-1]
% nguon: Archive/2024/0. [LVH] BG - 12 New - 4in1-20240625T224301Z-001/0. [LVH] BG - 12 New - 4in1/2. [LVH] TEX SGK 12 NEW-20240418T074833Z-001/2. [LVH] TEX SGK 12 NEW/[LVH] SGK-12-KNTT/data/C1/12-SGK-KNTT-1-3.tex #1094

	Tìm tiệm cận đứng của đồ thị hàm số $y=f(x)=\dfrac{3-x}{x+2}$.
	\loigiai{
		Ta có: $\displaystyle\lim _{x \rightarrow-2^{+}} f(x)=\displaystyle\lim _{x \rightarrow-2^{+}} \dfrac{3-x}{x+2}=+\infty$. \\
		Tương tự, $\displaystyle\lim _{x \rightarrow -2^{-}} f(x)=-\infty$. \\
		Vậy đồ thị hàm số $f(x)$ có tiệm cận đứng là đường thẳng $x=-2$.
	}
\end{bt}

\begin{bt}%[2D1N4-1]
% nguon: Archive/2024/0. [LVH] BG - 12 New - 4in1-20240625T224301Z-001/0. [LVH] BG - 12 New - 4in1/2. [LVH] TEX SGK 12 NEW-20240418T074833Z-001/2. [LVH] TEX SGK 12 NEW/[LVH] SGK-12-KNTT/data/C1/12-SGK-KNTT-1-3.tex #1095

	Tìm tiệm cận đứng của đồ thị hàm số $y=f(x)=\dfrac{x^{2}+2}{x}$.
	\loigiai{
		Ta có: $\displaystyle\lim _{x \rightarrow 0^{+}} f(x)=\displaystyle\lim _{x \rightarrow 0^{+}} \dfrac{x^{2}+2}{x}=+\infty$.\\
		Tương tự, $\displaystyle\lim _{x \rightarrow 0^{-}} f(x)=-\infty$. \\
		Vậy đồ thị hàm số $f(x)$ có tiệm cận đứng là đường thẳng $x=0$.
	}
\end{bt}

\begin{bt}%[2D1N4-4]
% nguon: Archive/2024/0. [LVH] BG - 12 New - 4in1-20240625T224301Z-001/0. [LVH] BG - 12 New - 4in1/2. [LVH] TEX SGK 12 NEW-20240418T074833Z-001/2. [LVH] TEX SGK 12 NEW/[LVH] SGK-12-KNTT/data/C1/12-SGK-KNTT-1-3.tex #1096

	Để loại bỏ $p\%$ một loài tảo độc khỏi một hồ nước, người ta ước tính chi phí bỏ ra là
	$$
	C(p)=\dfrac{45p}{100-p} \text { (triệu đồng), với } 0 \leq p<100.
	$$
	Tìm tiệm cận đứng của đồ thị hàm số $C(p)$ và nêu ý nghĩa thực tiễn của đường tiệm cận này.
	\loigiai{
		Ta có $\displaystyle\lim_{p\to 100^-}C(p)=\displaystyle\lim_{p\to 100^-}\dfrac{45p}{100-p}=+\infty$. \\
		Suy ra đồ thị hàm số có tiệm cận đứng là đường thẳng $p=100$.\\
		Ý nghĩa thực tiễn: Để loại bỏ hoàn toàn loài tảo độc trên thì chi phí vô cùng lớn (không thể thực hiện được).
	}
\end{bt}

\begin{bt}%[2D1H4-1]
% nguon: Archive/2024/0. [LVH] BG - 12 New - 4in1-20240625T224301Z-001/0. [LVH] BG - 12 New - 4in1/2. [LVH] TEX SGK 12 NEW-20240418T074833Z-001/2. [LVH] TEX SGK 12 NEW/[LVH] SGK-12-KNTT/data/C1/12-SGK-KNTT-1-3.tex #1097

	Cho hàm số $y=f(x)=x+\dfrac{1}{x+2}$. Tìm tiệm cận xiên của đồ thị hàm số $f(x)$
	\loigiai{
		Ta có $\lim\limits_{x \rightarrow+\infty}[f(x)-x)]=\lim\limits_{x \rightarrow+\infty} \dfrac{1}{x+2}=0$. Tương tự $\lim\limits_{x \rightarrow-\infty}[f(x)-x)]=0$.
		\\
		Do đó, đồ thị hàm số có tiệm cận xiên là đường thẳng $y=x$.	
	}
\end{bt}

\begin{bt}%[2D1H4-1]
% nguon: Archive/2024/0. [LVH] BG - 12 New - 4in1-20240625T224301Z-001/0. [LVH] BG - 12 New - 4in1/2. [LVH] TEX SGK 12 NEW-20240418T074833Z-001/2. [LVH] TEX SGK 12 NEW/[LVH] SGK-12-KNTT/data/C1/12-SGK-KNTT-1-3.tex #1098

	Tìm tiệm cận xiên của đồ thị hàm số $y=f(x)=\dfrac{x^2-x+2}{x+1}$.
	\loigiai{
		Ta có:
		$$
		\begin{aligned}
			& a=\lim _{x \rightarrow+\infty} \dfrac{f(x)}{x}=\lim _{x \rightarrow+\infty} \dfrac{x^2-x+2}{x^2+x}=1 ; \\
			& b=\lim _{x \rightarrow+\infty}[f(x)-x]=\lim _{x \rightarrow+\infty} \dfrac{-2 x+2}{x+1}=-2 .
		\end{aligned}
		$$
		(Tương tự, $\lim _{x \rightarrow-\infty} \dfrac{f(x)}{x}=1, \lim _{x \rightarrow-\infty}[f(x)-x]=-2$.)\\
		Vậy đồ thị hàm số $f(x)$ có tiệm cận xiên là đường thẳng $y=x-2$.
	}
\end{bt}

\begin{bt}%[2D1N1-2]
% nguon: Archive/2024/0. [LVH] BG - 12 New - 4in1-20240625T224301Z-001/0. [LVH] BG - 12 New - 4in1/2. [LVH] TEX SGK 12 NEW-20240418T074833Z-001/2. [LVH] TEX SGK 12 NEW/[LVH] SGK-12-KNTT/data/C1/12-SGK-KNTT-1-3.tex #1099

	Hình 1.26 là đồ thị của hàm số $y=f(x)=\dfrac{2 x^2}{x^2-1}$.
	\begin{center}
		\begin{tikzpicture}[scale=.7,>=stealth, font=\footnotesize, line join=round, line cap=round]
			\def\xmin{-6} \def\xmax{6}
			\def\ymin{-5} \def\ymax{7}
			%\draw[color=gray!50,dashed] (\xmin,\ymin) grid (\xmax,\ymax);
			\draw[->] (\xmin,0)--(\xmax,0) node [below]{$x$};
			\draw[->] (0,\ymin)--(0,\ymax) node [left]{$y$};
			\fill (0,0) circle (1pt) node[shift={(135:2.5mm)}]{$O$};
			\node at (current bounding box.south) [below=-2pt] {Hình 3};
			\clip (\xmin+0.1,\ymin+0.1) rectangle (\xmax-0.1,\ymax-0.1);
			\draw[smooth,red,samples=300,domain=(\xmin:-1.01)] plot(\x,{(2*(\x)^2)/((\x)^2-1)});		
			\draw[smooth,red,samples=300,domain=(-0.9:0.9)] plot(\x,{(2*(\x)^2)/((\x)^2-1)});
			\draw[smooth,red,samples=300,domain=(1.1:\xmax)] plot(\x,{(2*(\x)^2)/((\x)^2-1)});
			\draw[blue] (\xmin,2)--(\xmax,2);	
			\draw[blue] (-1,\ymin)--(-1,\ymax);	
			\draw[blue] (1,\ymin)--(1,\ymax);		
			\foreach \x in {\xmin,...,\xmax}
			\draw (\x,-0.1)--(\x,0.1);
			\foreach \y in {\ymin,...,\ymax}
			\draw (-0.1,\y)--(0.1,\y);
			\node at (-5,2)[below]{$y=2$};
			\node at (-1.2,-4)[left]{$x=-1$};
			\node at (1.2,-4)[right]{$x=1$};
			%\node at (-1,0)[shift={(-135:2.5mm)}]{$-1$};
			%\node at (.5,0)[shift={(-75:2.5mm)}]{$\dfrac{1}{2}$};
			%\node at (0,-1)[left]{$-1$};
			%\node at (0,2)[shift={(135:2.5mm)}]{$2$};		
		\end{tikzpicture}
	\end{center}
	Sử dụng đồ thị này, hãy:
	\begin{enumerate}[a)]
		\item Viết kết quả của các giới hạn sau: $\lim _{x \rightarrow-\infty} f(x)$ ; $\lim _{x \rightarrow+\infty} f(x)$ ; $\lim _{x \rightarrow 1^{-}} f(x)$ ; $\lim _{x \rightarrow-1^{+}} f(x)$.
		\item  Chỉ ra các tiệm cận của đồ thị hàm số đã cho.
	\end{enumerate}
	\loigiai{
		\begin{enumerate}[a)]
			\item  $\lim _{x \rightarrow-\infty} f(x)=2$ ; $\lim _{x \rightarrow+\infty} f(x)=2$ ; $\lim _{x \rightarrow 1^{-}} f(x)=-\infty$ ; $\lim _{x \rightarrow-1^{+}} f(x)-\infty$.
			\item  Đồ thị hàm số có tiệm cận đứng là đường thẳng $x=-1$ và $x=1$; tiệm cận ngang là đường thẳng $y=2$.
		\end{enumerate}
	}
\end{bt}

\begin{bt}%[2D1H1-1]
% nguon: Archive/2024/0. [LVH] BG - 12 New - 4in1-20240625T224301Z-001/0. [LVH] BG - 12 New - 4in1/2. [LVH] TEX SGK 12 NEW-20240418T074833Z-001/2. [LVH] TEX SGK 12 NEW/[LVH] SGK-12-KNTT/data/C1/12-SGK-KNTT-1-3.tex #1100

	Đường thẳng $x=1$ có phải là tiệm cận đứng của đồ thị hàm số $y=\dfrac{x^2+2 x-3}{x-1}$ không?
	
	\loigiai{Ta có:\\
		$\begin{aligned} 
			&\lim _{x \rightarrow 1^{+}}f(x)=\lim _{x \rightarrow 1^{+}}\dfrac{x^2+2 x-3}{x-1}=\lim _{x \rightarrow 1^{+}}\dfrac{(x-1)(x+3)}{x-1}=\lim _{x \rightarrow 1^{+}}(x+3)=4;\\
			&  \lim _{x \rightarrow 1^{-}}f(x)=\lim _{x \rightarrow 1^{-}}\dfrac{x^2+2 x-3}{x-1}=\lim _{x \rightarrow 1^{-}}\dfrac{(x-1)(x+3)}{x-1}=\lim _{x \rightarrow 1^{-}}(x+3)=4.
		\end{aligned}$\\
		Suy ra đường thẳng $x=1$ không phải là tiệm cận đứng của đồ thị hàm số $y=\dfrac{x^2+2 x-3}{x-1}$.
	}
\end{bt}

\begin{bt}%[2D1H1-1]
% nguon: Archive/2024/0. [LVH] BG - 12 New - 4in1-20240625T224301Z-001/0. [LVH] BG - 12 New - 4in1/2. [LVH] TEX SGK 12 NEW-20240418T074833Z-001/2. [LVH] TEX SGK 12 NEW/[LVH] SGK-12-KNTT/data/C1/12-SGK-KNTT-1-3.tex #1101

	Tìm các tiệm cận của đồ thị các hàm số sau:
	\begin{enumerate}[a)]
		\item $y=\dfrac{3-x}{2 x+1}$
		\item $y=\dfrac{2 x^2+x-1}{x+2}$.
	\end{enumerate}
	\loigiai{\begin{enumerate}[a)]
			\item \begin{itemize}
				\item 
				Ta có: $\lim _{x \rightarrow+\infty} f(x)=\lim _{x \rightarrow+\infty} \dfrac{3- x}{2x+1}=\lim _{x \rightarrow+\infty} \dfrac{\dfrac{3}{x}-1}{2+\dfrac{1}{x}}=-\dfrac{1}{2}$. Tương tự, $\lim _{x \rightarrow-\infty} f(x)=-\dfrac{1}{2}$.\\
				Vậy đồ thị hàm số $f(x)$ có tiệm cận ngang là đường thẳng $y=3$.
				\item Ta có: $\lim _{x \rightarrow-\dfrac{1}{2}^{+}} f(x)=\lim _{x \rightarrow-\dfrac{1}{2}^{+}} \dfrac{3-x}{2x+1}=+\infty$. Tương tự, $\lim _{x \rightarrow-\dfrac{1}{2}^{-}} f(x)=-\infty$. Vậy đồ thị hàm số $f(x)$ có tiệm cận đứng là đường thẳng $x=-\dfrac{1}{2}$.
			\end{itemize}
			\item \begin{itemize}
				\item 	Ta có: $\lim _{x \rightarrow+\infty} f(x)=\lim _{x \rightarrow+\infty} \dfrac{2 x^2+x-1}{x+2}=+\infty$. Tương tự, $\lim _{x \rightarrow+\infty} f(x)=\lim _{x \rightarrow-\infty} \dfrac{2 x^2+x-1}{x+2}=-\infty$.\\
				Vậy đồ thị hàm số $f(x)$ không có tiệm cận ngang.
				\item Ta có: $\lim _{x \rightarrow -2^{+}} f(x)=\lim _{x \rightarrow -2^{+}} \dfrac{2 x^2+x-1}{x+2}=+\infty$. Tương tự, $\lim _{x \rightarrow -2^{-}} f(x)=-\infty$.\\
				Vậy đồ thị hàm số $f(x)$ có tiệm cận đứng là đường thẳng $x=-2$.
				\item 	Ta có:
				$$
				\begin{aligned}
					& a=\lim _{x \rightarrow+\infty} \dfrac{f(x)}{x}=\lim _{x \rightarrow+\infty} \dfrac{2 x^2+x-1}{x^2+2x}=2 ; \\
					& b=\lim _{x \rightarrow+\infty}[f(x)-2x]=\lim _{x \rightarrow+\infty} \dfrac{-3 x-1}{x+2}=-3 .
				\end{aligned}
				$$
				(Tương tự, $\lim _{x \rightarrow-\infty} \dfrac{f(x)}{x}=2, \lim _{x \rightarrow-\infty}[f(x)+x]=-3$.)\\
				Vậy đồ thị hàm số $f(x)$ có tiệm cận xiên là đường thẳng $y=2x-3$.
			\end{itemize}
	\end{enumerate}}
\end{bt}

\begin{bt}%[2D1H1-1]
% nguon: Archive/2024/0. [LVH] BG - 12 New - 4in1-20240625T224301Z-001/0. [LVH] BG - 12 New - 4in1/2. [LVH] TEX SGK 12 NEW-20240418T074833Z-001/2. [LVH] TEX SGK 12 NEW/[LVH] SGK-12-KNTT/data/C1/12-SGK-KNTT-1-3.tex #1102

	Một công ty sản xuất đồ gia dụng ước tính chi phí để sản xuât $x$ (sản phẩm) là $C(x)=2 x+50$ (triệu đồng).\\
	Khi đó $f(x)=\dfrac{C(x)}{x}$ là chi phí sản xuất trung bình cho mỗi sản phẩm. Chứng tỏ rằng hàm số $f(x)$ giảm và $\lim _{x \rightarrow+\infty} f(x)=2$. Tính chất này nói lên điều gì?
	\loigiai{Ta có $f(x)=\dfrac{C(x)}{x}=\dfrac{2x+50}{x}$.\\
		$f^\prime(x)=\dfrac{-50}{x^2}$ suy ra hàm số giảm.\\
		$\lim _{x \rightarrow+\infty} f(x)=\lim _{x \rightarrow+\infty} \dfrac{2x+50}{x}=2$.
		\begin{center}
			\begin{tikzpicture}[font=\footnotesize,thick,>=stealth]
				\tikzset{double style/.append style={double distance=1.5pt}}
				\tkzTabInit[nocadre=false,lgt=1.2,espcl=2.5,deltacl=0.6,lw=.75pt,color,colorL=green!50,colorV=green!50]
				{$x$ /0.6,$f'(x)$ /0.6,$f(x)$ /2}
				{$0$,$+\infty$}
				\tkzTabLine{,-,}
				\tkzTabVar{D+/$+\infty$,-/$2$}
			\end{tikzpicture}
		\end{center}
		Từ đó suy ra số sản phẩm sản xuất càng nhiều thì chi phí sản xuất trung bình cho mỗi sản phẩm càng giảm nhưng luôn lớn hơn $2$ triệu đồng. 
	}
\end{bt}

\begin{bt}%[2D1V1-5]
% nguon: Archive/2024/0. [LVH] BG - 12 New - 4in1-20240625T224301Z-001/0. [LVH] BG - 12 New - 4in1/2. [LVH] TEX SGK 12 NEW-20240418T074833Z-001/2. [LVH] TEX SGK 12 NEW/[LVH] SGK-12-KNTT/data/C1/12-SGK-KNTT-1-3.tex #1103

	Một mảnh vườn hình chữ nhật có diện tích bằng $144 \mathrm{~m}^2$. Biết độ dài một cạnh của mảnh vườn là $x(m)$.
	\begin{enumerate}[a)]
		\item Viết biểu thức tính chu vi $P(x)$ (mét) của mảnh vườn.
		\item Tìm các tiệm cận của đồ thị hàm số $P(x)$.
	\end{enumerate}
	\loigiai{
		\begin{enumerate}[a)]
			\item Ta có độ dài một cạnh của mảnh vườn là $x(m)$ nên độ dài cạnh còn lại của mảnh vườn là $\dfrac{144}{x}(m)$.\\
			Vậy chu vi của mảnh vườn là $P(x)=2\left(x+\dfrac{144}{x}\right)=2x+\dfrac{72}{x}$.
			\item \begin{itemize}
				\item 	Ta có: $\lim _{x \rightarrow+\infty} P(x)=\lim _{x \rightarrow+\infty} 2x+\dfrac{72}{x}=+\infty$. Tương tự, $\lim _{x \rightarrow+\infty} P(x)=\lim _{x \rightarrow-\infty} \left(2x+\dfrac{72}{x}\right)=-\infty$.\\
				Vậy đồ thị hàm số $P(x)$ không có tiệm cận ngang.
				\item Ta có: $\lim _{x \rightarrow 0^{+}} f(x)=\lim _{x \rightarrow 0^{+}} \left(2x+\dfrac{72}{x}\right)=+\infty$. Tương tự, $\lim _{x \rightarrow 0^{-}} P(x)=-\infty$.\\
				Vậy đồ thị hàm số $f(x)$ có tiệm cận đứng là đường thẳng $x=0$.
				\item 	Ta có:
				$\begin{aligned} 
					&\lim _{x \rightarrow+\infty}[P(x)-2x]=\lim _{x \rightarrow+\infty} \dfrac{72}{x}=0;\\
					&  \lim _{x \rightarrow-\infty}[P(x)-2x]=\lim _{x \rightarrow-\infty} \dfrac{72}{x}=0 .
				\end{aligned}$\\
				Vậy đồ thị hàm số $P(x)$ có tiệm cận xiên là đường thẳng $y=2x$.
			\end{itemize}
		\end{enumerate}
	}
\end{bt}

\begin{bt}%[2D1N1-2]
% nguon: Archive/2024/0. [LVH] BG - 12 New - 4in1-20240625T224301Z-001/0. [LVH] BG - 12 New - 4in1/2. [LVH] TEX SGK 12 NEW-20240418T074833Z-001/2. [LVH] TEX SGK 12 NEW/[LVH] SGK-12-KNTT/data/C1/12-SGK-KNTT-1-1.tex #1104

	\immini{
		Hình 1.4 là đồ thị của hàm số $y=f(x)=|x|$. Hãy tìm các khoảng đồng biến, khoảng nghịch biến của hàm số.	
		
	}{
		\begin{tikzpicture}[>=stealth,line join=round,line cap=round,font=\footnotesize,thick,scale=.8]
			\tikzset{
				declare function={x2=2.5;y2=3.5;x1=-2.5;y1=-1;a=1;b=-3;c=-3;},
				declare function={f(\x)=a*abs(\x);}
			}
			
			%	\clip (x1,y1)rectangle (x2,y2);
			%	\draw[gray!20] (x1,y1)grid(x2,y2);
			\draw[fill=black](0,0)circle(1.2pt)node[below left,cyan]{$O$};
			%%%%%%%%
			\draw [smooth,samples=100,domain=x1:x2,orange]plot(\x,{f(\x)})node[below=9mm,cyan]{$y=f(x)=|x|$};	
			\draw[thick,->] (x1,0)--(x2+1,0)node[above left,cyan]{$x$};
			\draw[thick,->] (0,y1)--(0,y2)node[below right,cyan]{$y$};
			\path 
			(0,y1)node[below,orange]{Hình $1.4$}
			;
		\end{tikzpicture}		
		
	}
	\loigiai{
		Tập xác định của hàm số là $\mathbb{R}$.\\
		Từ đồ thị suy ra: Hàm số đồng biến trên khoảng $(0;+\infty)$, nghịch biến trên khoảng $(-\infty; 0)$.
	}
\end{bt}

\begin{bt}%[2D1N1-1]
% nguon: Archive/2024/0. [LVH] BG - 12 New - 4in1-20240625T224301Z-001/0. [LVH] BG - 12 New - 4in1/2. [LVH] TEX SGK 12 NEW-20240418T074833Z-001/2. [LVH] TEX SGK 12 NEW/[LVH] SGK-12-KNTT/data/C1/12-SGK-KNTT-1-1.tex #1105

	Tìm các khoảng đồng biến, khoảng nghịch biến của hàm số $y=x^2-4 x+2$.
	\loigiai{
		Tập xác định của hàm số là $\mathbb{R}$.\\
		Ta có $y'=2 x-4$; $y'>0$ với $x \in(2;+\infty)$; $y'<0$ với $x \in(-\infty; 2)$.\\
		Do đó, hàm số đồng biến trên khoảng $(2;+\infty)$, nghịch biến trên khoảng $(-\infty; 2)$.
	}
\end{bt}

\begin{bt}%[2D1H1-1]
% nguon: Archive/2024/0. [LVH] BG - 12 New - 4in1-20240625T224301Z-001/0. [LVH] BG - 12 New - 4in1/2. [LVH] TEX SGK 12 NEW-20240418T074833Z-001/2. [LVH] TEX SGK 12 NEW/[LVH] SGK-12-KNTT/data/C1/12-SGK-KNTT-1-1.tex #1106

	Tìm khoảng đơn điệu của hàm số $y = \dfrac{x^2 -2x + 5}{x - 1}$.
	\loigiai{
		Tập xác định của hàm số là $\mathbb{R}\setminus\{1\}$.\\
		Ta có $y' = \dfrac{\left(2x-2\right)\left(x-1\right) - \left(x^2 - 2x + 5\right)}{(x - 1)^{2}} = \dfrac{x^2 - 2x - 3}{(x - 1)^{2}}$.\\
		$y' = 0 \Leftrightarrow x = -1$ hoặc $x = 3$.\\
		Ta có bảng biến thiên:
		\begin{center}
			\begin{tikzpicture}[font=\footnotesize,thick,>=stealth]
				\tikzset{double style/.append style={double distance=1.5pt}}
				\tkzTabInit[nocadre=false,lgt=1.2,espcl=2.5,deltacl=0.6,lw=.75pt,color,colorL=green!50,colorV=green!50]
				{$x$ /0.6,$y'$ /0.6,$y$ /2}
				{$-\infty$,$-1$,$1$,$3$,$+\infty$}
				\tkzTabLine{,+,$0$,-,d,-,$0$,+,}
				\tkzTabVar{-/$-\infty$,+/$-4$,-D+/$-\infty$/$+\infty$,-/$4$,+/$+\infty$}
			\end{tikzpicture}
		\end{center}
		Từ bảng biến thiên, ta có\\
		Hàm số đồng biến trên các khoảng $\left(-\infty; -1\right)$ và $\left(3; +\infty\right)$.\\
		Hàm số nghịch biến trên các khoảng $\left(-1; 1\right)$ và $\left(1; 3\right)$.
	}
\end{bt}

\begin{bt}%[2D1H1-1]
% nguon: Archive/2024/0. [LVH] BG - 12 New - 4in1-20240625T224301Z-001/0. [LVH] BG - 12 New - 4in1/2. [LVH] TEX SGK 12 NEW-20240418T074833Z-001/2. [LVH] TEX SGK 12 NEW/[LVH] SGK-12-KNTT/data/C1/12-SGK-KNTT-1-1.tex #1107

	Xét chiều biến thiên của hàm số $y = \dfrac{x -2}{x + 1}$.
	\loigiai{
		Tập xác định của hàm số là $\mathbb{R} \setminus\{1\}$.\\
		Ta có $y' = \dfrac{\left(x - 1\right) - \left(x - 2\right)}{\left(x + 1\right)^{2}} = \dfrac{3}{\left(x + 1\right)^2} > 0$, với mọi $x \neq 1$.\\
		Lập bảng biến thiên của hàm số:	
		\begin{center}
			\begin{tikzpicture}[scale=1, font=\footnotesize]
				\tkzTabInit[nocadre=false, lgt=1.2, espcl=2, deltacl=0.6]
				{$x$/0.8,$f'(x)$/0.6,$f(x)$/2}
				{$-\infty$,$-1$,$+\infty$};
				\tkzTabLine{,+,d,+,};
				\tkzTabVar{-/$1$,+D-/$+\infty$/$-\infty$,+/$1$};
			\end{tikzpicture}
		\end{center}
		Từ bảng biến thiên, ta có\\
		Hàm số đồng biến trên các khoảng $\left(-\infty; -1\right)$.
	}
\end{bt}

\begin{bt}%[2D1N2-2]
% nguon: Archive/2024/0. [LVH] BG - 12 New - 4in1-20240625T224301Z-001/0. [LVH] BG - 12 New - 4in1/2. [LVH] TEX SGK 12 NEW-20240418T074833Z-001/2. [LVH] TEX SGK 12 NEW/[LVH] SGK-12-KNTT/data/C1/12-SGK-KNTT-1-1.tex #1109

		Hình vẽ bên là đồ thị của hàm số $y=f(x)$. Hãy tìm các cực trị của hàm số.
		\loigiai{
			Từ đồ thị hàm số, ta có:\\
			Hàm số đạt cực tiểu tại $x=-1$ và $y_{\text{CT}}=y(-1)=2$.\\
			Hàm số đạt cực đại tại $x=0$ và $y_{\text{CĐ}}=y(0)=3$.\\
			Hàm số đạt cực tiểu tại $x=1$ và $y_{\text{CT}}=y(1)=2$.	
		}
	\end{bt}

\begin{bt}%[2D1H2-1]
% nguon: Archive/2024/0. [LVH] BG - 12 New - 4in1-20240625T224301Z-001/0. [LVH] BG - 12 New - 4in1/2. [LVH] TEX SGK 12 NEW-20240418T074833Z-001/2. [LVH] TEX SGK 12 NEW/[LVH] SGK-12-KNTT/data/C1/12-SGK-KNTT-1-1.tex #1110

	Tìm cực trị của hàm số $y=x^{3}-6 x^{2}+9 x+30$.
	\loigiai{
		Tập xác định của hàm số là $\mathbb{R}$.
		
		Ta có: $y^{\prime}=3 x^{2}-12 x+9 ; y^{\prime}=0 \Leftrightarrow x=1$ hoặc $x=3$.
		
		Lập bảng biến thiên của hàm số:
		
		\begin{center}
			\begin{tikzpicture}
				\tkzTabInit[nocadre,lgt=1.2,espcl=2.5,deltacl=0.6]
				{$x$ /0.6,$y'$ /0.6,$y$ /1.5}
				{$-\infty$,$1$,$3$,$+\infty$}
				\tkzTabLine{,+,$0$,-,0,+}
				\tkzTabVar{-/$-\infty$,+/$34$,-/$30$,+/$+\infty$}
			\end{tikzpicture}
		\end{center}
		Từ bảng biến thiên, ta có:\\
		Hàm số đạt cực đại tại $x=1$ và $y_{\text{CĐ}}=y(1)=34$.\\
		Hàm số đạt cực tiểu tại $x=3$ và $y_{\text{CT}}=y(3)=30$.
	}
\end{bt}

\begin{bt}%[2D1H2-1]
% nguon: Archive/2024/0. [LVH] BG - 12 New - 4in1-20240625T224301Z-001/0. [LVH] BG - 12 New - 4in1/2. [LVH] TEX SGK 12 NEW-20240418T074833Z-001/2. [LVH] TEX SGK 12 NEW/[LVH] SGK-12-KNTT/data/C1/12-SGK-KNTT-1-1.tex #1111

	Tìm cực trị của hàm số $y=\dfrac{x^{2}-2 x+9}{x-2}$.
	\loigiai{
		Tập xác định của hàm số là $\mathbb{R}\setminus\{2\}$.\\
		Ta có: $y^{\prime}=\dfrac{(2 x-2)(x-2)-\left(x^{2}-2 x+9\right)}{(x-2)^{2}}=\dfrac{x^{2}-4 x-5}{(x-2)^{2}}$;\\
		$y^{\prime}=0 \Leftrightarrow x=-1$ hoặc $x=5$.\\
		Lập bảng biến thiên của hàm số:
		\begin{center}
			\begin{tikzpicture}
				\tikzset{double style/.append style={double distance=1.5pt}}
				\tkzTabInit[nocadre=false,lgt=1.2,espcl=2.5,deltacl=0.6]
				{$x$ /0.6,$y'$ /0.6,$y$ /2}
				{$-\infty$,$-1$,$2$,$5$,$+\infty$}
				\tkzTabLine{,+,$0$,-,d,-,$0$,+,}
				\tkzTabVar{-/$-\infty$,+/$-4$,-D+/$-\infty$/$+\infty$,-/$8$,+/$+\infty$}
			\end{tikzpicture}
		\end{center}
		Từ bảng biến thiên, ta có:\\
		Hàm số đạt cực đại tại $x=-1$ và $y_{\text{CĐ}}=y(-1)=-4$.\\
		Hàm số đạt cực tiểu tại $x=5$ và $y_{\text{CT}}=y(5)=8$.		
	}
\end{bt}

\begin{bt}%[2D1H2-1]
% nguon: Archive/2024/0. [LVH] BG - 12 New - 4in1-20240625T224301Z-001/0. [LVH] BG - 12 New - 4in1/2. [LVH] TEX SGK 12 NEW-20240418T074833Z-001/2. [LVH] TEX SGK 12 NEW/[LVH] SGK-12-KNTT/data/C1/12-SGK-KNTT-1-1.tex #1112

	Tìm cực trị của hàm số $y=\dfrac{x+1}{x-1}$.
	\loigiai{
		Tâp xác định của hàm số là $\mathbb{R} \setminus\{1\}$.\\
		Ta có: $y^{\prime}=\dfrac{(x-1)-(x+1)}{(x-1)^{2}}=\dfrac{-2}{(x-1)^{2}}<0$, vơi moi $x \neq 1$.\\
		Lập bảng biến thiên của hàm số:
		\begin{center}
			\begin{tikzpicture}
				\tkzTabInit[nocadre, lgt=2, espcl=4, deltacl=0.6]
				{$x$/0.8,$y'$/0.6,$y$/1.5}
				{$-\infty$,$1$,$+\infty$};
				\tkzTabLine{,-,d,-,};
				\tkzTabVar{+/$1$,-D+/$-\infty$/$+\infty$,-/$1$};
			\end{tikzpicture}
		\end{center}
		Từ bảng biên thiên suy ra hàm số không có cực trị.
	}
\end{bt}

\begin{bt}%[2D1H2-7]
% nguon: Archive/2024/0. [LVH] BG - 12 New - 4in1-20240625T224301Z-001/0. [LVH] BG - 12 New - 4in1/2. [LVH] TEX SGK 12 NEW-20240418T074833Z-001/2. [LVH] TEX SGK 12 NEW/[LVH] SGK-12-KNTT/data/C1/12-SGK-KNTT-1-1.tex #1113
]
	Một vật được phóng thẳng đứng lên trên từ độ cao $2 \mathrm{~m}$ với vận tốc ban đầu là $24{,}5 \mathrm{~m} / \mathrm{s}$. Trong Vật lí, ta biết rằng khi bỏ qua sức cản của không khí thì độ cao $h$ (mét) của vật sau $t$ (giây) được cho bởi công thức
	$$
	h(t)=2+24{,}5 t-4{,}9t^{2}.
	$$
	Hỏi tại thời điểm nào thì vật đạt độ cao lớn nhất?
	\loigiai{
		Xét hàm số $h(t)=2+24{,}5 t-4{,}9t^{2}$ với $t>0$.\\
		Ta có $h'(t)=24{,}5-9{,}8t$. \\
		$h'(t)=0\Leftrightarrow t=2{,}5$. \\
		Lập bảng biến thiên của $h(t)$: 
		\begin{center}
			\begin{tikzpicture}
				\tkzTabInit[nocadre,lgt=1.2,espcl=3.5,deltacl=0.6]
				{$t$ /0.6,$h'(t)$ /0.6,$h(t)$ /1.4}
				{$0$,$2{,}5$,$+\infty$}
				\tkzTabLine{,+,0,-,}
				\tkzTabVar{-/,+/$32{,}625$,-/}
			\end{tikzpicture}
		\end{center}
		Từ bảng biến thiên ta suy ra vật đạt độ cao lớn nhất tại thời điểm $t=2{,}5$ giây (Khi đó độ cao lớn nhất mà vật đạt được là $32{,}625$ mét).
	}
\end{bt}

\begin{bt}%[2D1N1-2]
% nguon: Archive/2024/0. [LVH] BG - 12 New - 4in1-20240625T224301Z-001/0. [LVH] BG - 12 New - 4in1/2. [LVH] TEX SGK 12 NEW-20240418T074833Z-001/2. [LVH] TEX SGK 12 NEW/[LVH] SGK-12-KNTT/data/C1/12-SGK-KNTT-1-1.tex #1114

	Tìm các khoảng đồng biến, khoảng nghịch biến của các hàm số có đồ thị như sau:
	\begin{center}
		\begin{tikzpicture}[scale=1,xscale=1.3, font=\footnotesize, line join=round, line cap=round, >=stealth]
			\draw[gray!50,thin,opacity=.5] (-2,-2) grid (3,3);
			\draw[->] (-1.6,0)--(2.6,0) node[below] {$x$};
			\draw[->] (0,-1.5)--(0,2.6) node[left] {$y$};
			\draw (0,0) node [above left] {$O$};
			\foreach \x/\nx in {-1/-1,1/1,2/2}
			\draw[thin] (\x,1pt)--(\x,-1pt) node [above] {$\nx$};
			\foreach \y/\ny in {1/1,2/2}
			\draw[thin] (1pt,\y)--(-1pt,\y) node [left] {$\ny$};
			\begin{scope}
				\clip (-2.8,-2) rectangle (4,4);
				\draw[samples=200,domain=-0.8:2,smooth,variable=\x] plot (\x,{(\x)^3-1.5*(\x)^2});
			\end{scope}
			\draw[dashed,thin] (0,-0.5)node[left]{$-\tfrac{1}{2}$}--(1,-0.5)--(1,0);
			\draw (1.2,0) node[below right]{$1,5$};
			\draw (0,3.4) node{\text{a) Đồ thị hàm số } $y=x^3-\dfrac{3}{2}x^2 \left(
				H.1.11\right)$ };
		\end{tikzpicture}
		\hspace{0.5cm}
		\begin{tikzpicture}[scale=1, font=\footnotesize, line join=round, line cap=round, >=stealth]
			\draw[gray!50,thin,opacity=.5] (-3,-1) grid (4,4);
			\draw[->] (-3,0)--(3.5,0) node[below] {$x$};
			\draw[->] (0,-1)--(0,3.5) node[left] {$y$};
			\draw (0,0) node [above left] {$O$};
			\foreach \x/\nx in {-2/-2,-1/-1,1/1,2/2,3/3}
			\draw[thin] (\x,1pt)--(\x,-1pt) node [below] {$\nx$};
			\foreach \y/\ny in {1/1,2/2,3/3}
			\draw[thin] (1pt,\y)--(-1pt,\y) node [left] {$\ny$};
			\begin{scope}
				%\clip (-3,-1) rectangle (2.9,1.2);
				\draw[samples=200,domain=-3:3,smooth,variable=\x] plot (\x,{(((\x)^2-4)^2)^(1/3)});
				%\draw[samples=200,domain=-1.999:1.999,smooth,variable=\x] plot (\x,{(((\x)^2-4)^2)^(1/3)});
				%\draw[samples=200,domain=2:3,smooth,variable=\x] plot (\x,{(((\x)^2-4)^2)^(1/3)});
			\end{scope}
			\draw[dashed,thin] (0,2.5)node[left]{$2,5$};
			\draw (0,3.85) node[above]{b) Đồ thị hàm số $y=\sqrt[3]{\left(x^2-4\right)^2} \left(H.1.12\right)$ };
			%\draw (0,-1) node{\text{\textit{Hình }} $6b$};
		\end{tikzpicture}
	\end{center}
	\loigiai{
		\begin{enumerate}[a)]
			\item Quan sát đồ thị \textit{Hình} $H.1.11$, ta thấy hàm số $y=x^3-\dfrac{3}{2}x^2$ đồng biến trên các khoảng $(-\infty;0)$,  $(1;+\infty)$; nghịch biến trên khoảng $(0;1)$.
			\item Quan sát đồ thị \textit{Hình} $H.1.12$, ta thấy hàm số $y=\sqrt[3]{\left(x^2-4\right)^2}$ đồng biến trên các khoảng $(-2;0)$, $(2;+\infty)$; nghịch biến trên các khoảng $(-\infty;-2)$, $(0;2)$.
		\end{enumerate} 
	}
\end{bt}

\begin{bt}%[2D1H1-1]
% nguon: Archive/2024/0. [LVH] BG - 12 New - 4in1-20240625T224301Z-001/0. [LVH] BG - 12 New - 4in1/2. [LVH] TEX SGK 12 NEW-20240418T074833Z-001/2. [LVH] TEX SGK 12 NEW/[LVH] SGK-12-KNTT/data/C1/12-SGK-KNTT-1-1.tex #1115

	Xét sự đồng biến, nghịch biến của các hàm số sau:
	\begin{enumerate}[a)]	
		\item $y=\frac{1}{3}x^3-2x^2+3x+1$;
		\item $y=-x^3+2x^2-5x+3$.
	\end{enumerate}
	\loigiai{
		\begin{enumerate}[a)]	
			\item $y=\frac{1}{3}x^3-2x^2+3x+1$.\\
			Hàm số đã cho có tập xác định là $\mathbb{R}$.\\
			Ta có: $y'=x^2-4x+3$;\\
			$y'=0\Leftrightarrow x^2-4x+3=0\Leftrightarrow x=1$ hoặc $x=3$.\\
			Bảng biến thiên của hàm số như sau:
			\begin{center}
				\begin{tikzpicture}[font=\footnotesize,thick,>=stealth]
					\tikzset{double style/.append style={double distance=1.5pt}}
					\tkzTabInit[nocadre=false,lgt=1.2,espcl=2.5,deltacl=0.6,lw=.75pt,color,colorL=green!50,colorV=green!50]
					{$x$ /0.8,$y'$ /0.6,$y$ /2}
					{$-\infty$,$1$,$3$,$+\infty$}
					\tkzTabLine{ ,+,z,-,z,+, }
					\tkzTabVar{-/$-\infty$,+/$\dfrac{7}{3}$,-/$1$,+/$+\infty$}
				\end{tikzpicture}
			\end{center}
			Vậy hàm số đồng biến trên mỗi khoảng $(-\infty;1)$ và $\left(3;+\infty\right)$; nghịch biến trên khoảng $\left(1;3\right)$.
			\item $y=-x^3+2x^2-5x+3$.\\
			Hàm số đã cho có tập xác định là $\mathbb{R}$.\\
			Ta có: $y'=-3x^2+4x-5$;\\
			$y'=0\Leftrightarrow -3x^2+4x-5=0\left(\text{phương trình vô nghiệm}\right) $ \\
			Bảng biến thiên của hàm số như sau:
			\begin{center}
				\begin{tikzpicture}[font=\footnotesize,thick,>=stealth]
					\tikzset{double style/.append style={double distance=1.5pt}}
					\tkzTabInit[nocadre=false,lgt=1.2,espcl=2.5,deltacl=0.6,lw=.75pt,color,colorL=green!50,colorV=green!50]
					{$x$ /0.7, $f'(x)$ /0.8, $f(x)$ /2}
					{$-\infty$,$+\infty$}
					\tkzTabLine{ ,-,}
					\tkzTabVar{+/$+\infty$,-/$-\infty$}
				\end{tikzpicture}
			\end{center}
			Vậy hàm số  nghịch biến trên  khoảng $\left(-\infty;+\infty\right)$.
		\end{enumerate} 
	}
\end{bt}

\begin{bt}%[2D1H1-1]
% nguon: Archive/2024/0. [LVH] BG - 12 New - 4in1-20240625T224301Z-001/0. [LVH] BG - 12 New - 4in1/2. [LVH] TEX SGK 12 NEW-20240418T074833Z-001/2. [LVH] TEX SGK 12 NEW/[LVH] SGK-12-KNTT/data/C1/12-SGK-KNTT-1-1.tex #1116

	Tìm các khoảng đơn điệu của các hàm số sau:
	\begin{multicols}{2}
		\begin{enumerate}[a)]	
			\item $y=\frac{2x-1}{x+2}$;
			\item $y=\frac{x^2+x+4}{x-3}$.
		\end{enumerate}
	\end{multicols}
	\loigiai{
		\begin{enumerate}[a)]	
			\item $y=\frac{2x-1}{x+2}$.\\
			Hàm số đã cho có tập xác định là $\mathbb{R}\setminus\{-2\}$.\\
			Ta có: $y'=\frac{5}{\left(x+2\right)^2}>0,\,\forall x\neq -2$.\\
			Bảng biến thiên của hàm số như sau:
			\begin{center}
				\begin{tikzpicture}[font=\footnotesize,thick,>=stealth]
					\tikzset{double style/.append style={double distance=1.5pt}}
					\tkzTabInit[nocadre=false,lgt=1.2,espcl=2.5,deltacl=0.6,lw=.75pt,color,colorL=green!50,colorV=green!50]
					{$x$ /0.6,$y'$ /0.6,$y$ /2}
					{$-\infty$,$-2$,$+\infty$}
					\tkzTabLine{,+,d,+,}
					\tkzTabVar{-/$2$,+D-/$+\infty$/$-\infty$,+/$2$}
				\end{tikzpicture}
			\end{center}
			Vậy hàm số đồng biến trên mỗi khoảng $(-\infty;-2)$ và $\left(-2;+\infty\right)$.
			\item $y=\frac{x^2+x+4}{x-3}$.\\
			Hàm số đã cho có tập xác định là $\mathbb{R}\setminus\{3\}$.\\
			Ta có: $y'=\frac{x^2-6x-7}{\left(x-3\right)^2}$.\\
			$y'=0\Leftrightarrow x^2-6x-7=0\Leftrightarrow x=-1$ hoặc $x=7$ \\
			Bảng biến thiên của hàm số như sau:
			\begin{center}
				\begin{tikzpicture}[font=\footnotesize,thick,>=stealth]
					\tikzset{double style/.append style={double distance=1.5pt}}
					\tkzTabInit[nocadre=false,lgt=1.2,espcl=2.5,deltacl=0.6,lw=.75pt,color,colorL=green!50,colorV=green!50]
					{$x$ /0.8,$y'$ /0.6,$y$ /2}
					{$-\infty$,$-1$,$3$,$7$,$+\infty$}
					\tkzTabLine{ ,+,z,-,d,-,z,+, }
					\tkzTabVar{-/ $-\infty$, +/ $-1$, -D+/ $-\infty$ / $+\infty$, -/ $15$, +/ $+\infty$}
				\end{tikzpicture}
			\end{center}
			Vậy hàm số đồng biến trên các khoảng $\left(-\infty;-1\right)$ và $\left(7;+\infty\right)$ ; nghịch biến trên mỗi khoảng $\left(-1;3\right)$ và $\left(3;7\right)$.
		\end{enumerate} 
	}
\end{bt}

\begin{bt}%[2D1H1-1]
% nguon: Archive/2024/0. [LVH] BG - 12 New - 4in1-20240625T224301Z-001/0. [LVH] BG - 12 New - 4in1/2. [LVH] TEX SGK 12 NEW-20240418T074833Z-001/2. [LVH] TEX SGK 12 NEW/[LVH] SGK-12-KNTT/data/C1/12-SGK-KNTT-1-1.tex #1117

	Xét chiều biến thiên của các hàm số sau:
	\begin{multicols}{2}	
		\begin{enumerate}[a)]	
			\item $y=\sqrt{4-x^2}$;
			\item $y=\frac{x}{x^2+1}$.
		\end{enumerate}
	\end{multicols}
	\loigiai{
		\begin{enumerate}[a)]	
			\item $y=\sqrt{4-x^2}$.\\
			Hàm số đã cho có tập xác định là $\left[-2; 2\right]$.\\
			Ta có: $y'=\frac{-x}{\sqrt{4-x^2}}$.\\
			$y'=0\Leftrightarrow x=0$. \\
			Bảng biến thiên của hàm số như sau:
			\begin{center}
				\begin{tikzpicture}[font=\footnotesize,thick,>=stealth]
					\tikzset{double style/.append style={double distance=1.5pt}}
					\tkzTabInit[nocadre=false,lgt=1.2,espcl=2.5,deltacl=0.6,lw=.75pt,color,colorL=green!50,colorV=green!50]
					{$x$ /0.8,$y'$ /0.6,$y$ /2}
					{$-2$,$0$,$2$}
					\tkzTabLine{ ,+,z,-, }
					\tkzTabVar{-/$0$,+/$2$,-/$0$}
				\end{tikzpicture}
			\end{center}	
			Vậy hàm số đồng biến trên  khoảng $\left(-2;0\right)$ ; nghịch biến trên khoảng  $\left(0;2\right)$.
			\item $y=\frac{x}{x^2+1}$.\\
			Hàm số đã cho có tập xác định là $\mathbb{R}$.\\
			Ta có: $y'=\frac{-x+1}{\left(x^2+1\right)^2}$.\\
			$y'=0\Leftrightarrow x=1$. \\
			Bảng biến thiên của hàm số như sau:
			\begin{center}
				\begin{tikzpicture}[font=\footnotesize,thick,>=stealth]
					\tikzset{double style/.append style={double distance=1.5pt}}
					\tkzTabInit[nocadre=false,lgt=1.2,espcl=2.5,deltacl=0.6,lw=.75pt,color,colorL=green!50,colorV=green!50]
					{$x$ /0.8,$y'$ /0.6,$y$ /2}
					{$-\infty$,$1$,$+\infty$}
					\tkzTabLine{ ,+,z,-, }
					\tkzTabVar{-/$0$,+/$\frac{1}{2}$,-/$0$}
				\end{tikzpicture}
			\end{center}	
			Vậy hàm số đồng biến trên  khoảng $\left(-\infty;1\right)$ ; nghịch biến trên khoảng  $\left(1;+\infty\right)$.    
		\end{enumerate} 
	}
\end{bt}

\begin{bt}%[2D1H2-2]
% nguon: Archive/2024/0. [LVH] BG - 12 New - 4in1-20240625T224301Z-001/0. [LVH] BG - 12 New - 4in1/2. [LVH] TEX SGK 12 NEW-20240418T074833Z-001/2. [LVH] TEX SGK 12 NEW/[LVH] SGK-12-KNTT/data/C1/12-SGK-KNTT-1-1.tex #1119
%[Dự án Tex hóa SKG KNTT - Nguyễn Mạnh Trung Anh]
	Đồ thị của đạo hàm bậc nhất $y=f'(x)$ của hàm số $f(x)$ được cho trong hình $1.13$.
	\begin{enumerate}
		\item Hàm số $f(x)$ đồng biến trên những khoảng nào? Giải thích.
		\item Tại giá trị nào của $x$ thì $f(x)$ có cực đại và cực tiểu? Giải thích.
	\end{enumerate}
	\begin{center}
		\begin{tikzpicture}[line join=round, line cap=round,scale=2,transform shape,scale=0.6]
			
			\def\xmin{-1} \def\xmax{9}
			\def\ymin{-2.5} \def\ymax{3} 
			\draw[->] (\xmin,0)--(\xmax,0) node [above left]{$x$};
			\draw[->] (0,\ymin)--(0,\ymax) node [below right]{$y$};
			\node at (0,0) [below left]{$O$};
			\node at (3,2.2) {$y=f'(x)$};
			\foreach \x in{1,2,3,4,5,6}\draw[fill=black] (\x,0) circle (.5pt) node[below]{\footnotesize $\x$};
			\clip (\xmin+0.1,\ymin+0.1) rectangle (\xmax-0.1,\ymax-0.1);
			\draw[line width=.5mm, red] (0,0)
			..controls +(-60:.2) and +(180:.5) ..(1,-1.2)
			..controls +(0:1) and +(180:.8) ..(3,1.7)
			..controls +(0:1.1) and +(180:1) ..(5,-2)
			..controls +(0:1) and +(180:.8) ..(6.8,1.8)
			..controls +(0:.3) and +(180:.3) ..(7.5,1.5)
			..controls +(0:.5) and +(-100:.5) ..(8.5,2.7)
			;
		\end{tikzpicture}
	\end{center}
	\loigiai{
		\begin{enumerate}
			\item Từ đồ thị hàm số, ta có $f'(x)>0\Leftrightarrow\heva{&x\in(2;4)\\&x\in(6;+\infty).}$ và liên tục trên hai khoảng đó.\\
			Do đó hàm số $f(x)$ đồng biến trên $(2;4)$ và $(6;+\infty)$.
			\item Từ đồ thị hàm số, ta có $f'(x)=0\Leftrightarrow\heva{&x=2\\&x=4\\&x=6}$ và qua các điểm đó thì $f'(x)$ đổi dấu.\\
			Vậy tại các điểm $x=2$, $x=4$, $x=6$ thì hàm số $f(x)$ có cực đại hoặc cực tiểu.
		\end{enumerate}
	}
\end{bt}

\begin{bt}%[2D1H2-1]
% nguon: Archive/2024/0. [LVH] BG - 12 New - 4in1-20240625T224301Z-001/0. [LVH] BG - 12 New - 4in1/2. [LVH] TEX SGK 12 NEW-20240418T074833Z-001/2. [LVH] TEX SGK 12 NEW/[LVH] SGK-12-KNTT/data/C1/12-SGK-KNTT-1-1.tex #1120
%[Dự án Tex hóa SKG KNTT - Nguyễn Mạnh Trung Anh]
	Tìm cực trị của các hàm số sau
	\begin{multicols}{2}
		\begin{enumerate}
			\item $y=2x^3-9x^2+12x-5$;
			\item $y=x^4-4x^2+2$;
			\item $y=\dfrac{x^2-2x+3}{x-1}$;
			\item $y=\sqrt{4x-2x^2}$.
		\end{enumerate}
	\end{multicols}
	\loigiai{
		\begin{enumerate}
			\item $y=2x^3-9x^2+12x-5$.\\
			Tập xác định của hàm số $\mathscr{D}=\mathbb{R}$.\\
			Ta có $y'=6x^2-18x+12$; $y'=0\Leftrightarrow x=1$ hoặc $x=2$.\\
			Lập bảng biến thiên của hàm số
			\begin{center}
				\begin{tikzpicture}
					\tkzTabInit[nocadre=false,lgt=1.2,espcl=2.5,deltacl=0.6]
					{$x$ /0.6, $y'$ /0.6, $y$ /2.5}
					{$-\infty$,$1$, $2$,$+\infty$}
					\tkzTabLine{,+,$0$,-,$0$,+,}
					\tkzTabVar{-/$-\infty$,+/$0$,-/$-1$,+/$+\infty$}
				\end{tikzpicture}
			\end{center}
			Hàm số đạt cực đại tại $x=1$ và $y_{\text{CĐ}}=0$.\\
			Hàm số đạt cực tiểu tại $x=2$ và $y_{CT}=-1$.
			\item $y=x^4-4x^2+2$.\\
			Tập xác định của hàm số $\mathscr{D}=\mathbb{R}$.\\
			Ta có $y'=4x^3-8x$; $y'=0\Leftrightarrow x=0$ hoặc $x=\pm\sqrt{2}$.\\
			Lập bảng biến thiên của hàm số
			\begin{center}
				\begin{tikzpicture}
					\tkzTabInit[nocadre=false,lgt=1.2,espcl=2.5,deltacl=0.6]
					{$x$ /0.6, $y'$ /0.6, $y$ /2.5}
					{$-\infty$,$-\sqrt{2}$,$0$, $\sqrt{2}$,$+\infty$}
					\tkzTabLine{,-,0,+,0,-,0,+,}
					\tkzTabVar{+/$+\infty$,-/$-2$,+/$2$,-/$-2$,+/$+\infty$}
				\end{tikzpicture}
			\end{center}
			Hàm số đạt cực tiểu tại $x=-\sqrt{2}$ và $x=\sqrt{2}$, $y_{CT}=-2$.\\
			Hàm số đạt cực đại tại $x=0$ và $y_{\text{CĐ}}=2$.
			\item $y=\dfrac{x^2-2x+3}{x-1}$.\\
			Tập xác định $\mathscr{D}=\mathbb{R}\setminus\{1\}$.\\
			Ta có $y'=\dfrac{x^2-2x-1}{(x-1)^2}$; $y'=0\Leftrightarrow x=1+\sqrt{2}$ hoặc $x=1-\sqrt{2}$.\\
			Lập bảng biến thiên của hàm số 
			\begin{center}
				\begin{tikzpicture}
					\tkzTabInit[nocadre=false,lgt=1.3,espcl=2.5,deltacl=0.6]
					{$x$ /0.6, $y'$ /0.6, $y$ /2.5}
					{$-\infty$,$1-\sqrt{2}$,$1$, $1+\sqrt{2}$,$+\infty$}
					\tkzTabLine{,+,0,-,d,-,0,+,}
					\tkzTabVar{-/$-\infty$,+/$-2\sqrt{2}$, -D+ /$-\infty$ /$+\infty$, -/$2\sqrt{2}$,+/$+\infty$}
				\end{tikzpicture}
			\end{center}
			Hàm số đặt cực đại tại $x=1-\sqrt{2}$ và $y_{\text{CĐ}}=-2\sqrt{2}$ trên khoảng $(-\infty;1)$.\\
			Hàm số đạt cực tiểu tại $x=1+\sqrt{2}$ và $y_{\text{CT}}=2\sqrt{2}$ trên khoảng $(1;+\infty)$.
			\item $y=\sqrt{4x-2x^2}$\\
			Tập xác định $\mathscr{D}=[0;2]$.\\
			Ta có $y'=\left(\sqrt{4x-2x^2}\right)'=\dfrac{4-4x}{2\sqrt{4x-2x^2}}$; $y'=0\Leftrightarrow x=1$.\\
			Lập bảng biến thiên của hàm số
			\begin{center}
				\begin{tikzpicture}
					\tkzTabInit[nocadre=false,lgt=1.3,espcl=2.5,deltacl=0.6]
					{$x$ /0.6, $y'$ /0.6, $y$ /2.5}
					{$0$, $1$, $2$}
					\tkzTabLine{,+,0,-,}
					\tkzTabVar{-/ $0$,+/$\sqrt{2}$,-/$0$}
				\end{tikzpicture}
			\end{center}
			Hàm số đạt cực đại tại $x=1$ và $y_{\text{CĐ}}=\sqrt{2}$.\\
			Hàm số đạt cực tiểu tại $x=0$ và $x=2$; $y_{CT}=0$.
		\end{enumerate}
	}
\end{bt}

\begin{bt}%[2D1V2-1]
% nguon: Archive/2024/0. [LVH] BG - 12 New - 4in1-20240625T224301Z-001/0. [LVH] BG - 12 New - 4in1/2. [LVH] TEX SGK 12 NEW-20240418T074833Z-001/2. [LVH] TEX SGK 12 NEW/[LVH] SGK-12-KNTT/data/C1/12-SGK-KNTT-1-1.tex #1121
%[Dự án Tex hóa SKG KNTT - Nguyễn Mạnh Trung Anh]
	Cho hàm số $y=f(x)=|x|$.
	\begin{enumerate}
		\item Tính các giới hạn $\lim\limits_{x\rightarrow0^+}\dfrac{f(x)-f(0)}{x-0}$ và $\lim\limits_{x\rightarrow0^-}\dfrac{f(x)-f(0)}{x-0}$.\\
		Từ đó suy ra hàm số không có đạo hàm tại $x=0$.
		\item Sử dụng định nghĩa, chứng minh hàm số có cực tiểu tại $x=0$.
	\end{enumerate}
	\loigiai{
		\begin{enumerate}
			\item Ta có
			\begin{itemize}
				\item $\lim\limits_{x\rightarrow0^+}\dfrac{f(x)-f(0)}{x-0}=\lim\limits_{x\rightarrow0^+}\dfrac{|x|-0}{x-0}=\lim\limits_{x\rightarrow0^+}\dfrac{x-0}{x-0}=1$.
				\item $\lim\limits_{x\rightarrow0^-}\dfrac{f(x)-f(0)}{x-0}=\lim\limits_{x\rightarrow0^-}\dfrac{|x|-0}{x-0}=\lim\limits_{x\rightarrow0^-}\dfrac{-x-0}{x-0}=-1$.
			\end{itemize}
			Suy ra $\lim\limits_{x\rightarrow0^+}\dfrac{f(x)-f(0)}{x-0}\neq\lim\limits_{x\rightarrow0^-}\dfrac{f(x)-f(0)}{x-0}$.\\
			Do đó hàm số không tồn tại giới hạn $\lim\limits_{x\rightarrow0}\dfrac{f(x)-f(0)}{x-0}$.\\
			Vậy hàm số không có đạo hàm tại điểm $x=0$.
			\item Do $f(0)=0$ và $f(x)>0$ với mọi $x\neq0$ nên hàm số đạt cực tiểu tại $x=0$.
			\begin{note}
				Một hàm số chỉ có thể đạt cực trị tại một điểm mà tại đó đạo hàm của hàm số bằng $0$, hoặc tại đó hàm số không có đạo hàm.
			\end{note}
		\end{enumerate}
	}
\end{bt}

\begin{bt}%[2D1V2-7]
% nguon: Archive/2024/0. [LVH] BG - 12 New - 4in1-20240625T224301Z-001/0. [LVH] BG - 12 New - 4in1/2. [LVH] TEX SGK 12 NEW-20240418T074833Z-001/2. [LVH] TEX SGK 12 NEW/[LVH] SGK-12-KNTT/data/C1/12-SGK-KNTT-1-1.tex #1122
%[Dự án Tex hóa SKG KNTT - Nguyễn Mạnh Trung Anh]
	Giả sử doanh số (tính bằng sản phẩm) của một sản phẩm mới (trong một năm nhất định) tuân theo quy luật logistic được mô hình hóa bằng hàm số \[f(t)=\dfrac{5000}{1+5e^{-t}},\, t\ge0,\]
	trong đó thời gian $t$ được tính bằng năm, kể từ khi phát hành sản phẩm mới. Khi đó đạo hàm $f'(t)$ biểu thị tốc độ bán hàng. Hỏi sau khi phát hành bao nhiêu năm thì tốc độ bán hàng là lớn nhất?
	\loigiai{
		Gọi $g(t)$ là hàm tốc độ bán hàng.\\
		Khi đó $g(t)=f'(t)=\dfrac{25\,000e^{-t}}{(1+5e^{-t})^2}$, $t\ge0$.\\
		Ta có $g'(t)=\dfrac{25\,000e^{-t}(1+5e^{-t})(5e^{-t}-1)}{(1+5e^{-t})^4}$; $g'(t)=0\Leftrightarrow t=-\ln{\dfrac{1}{5}}$.\\
		Bảng biến thiên hàm số 
		\begin{center}
			\begin{tikzpicture}
				\tkzTabInit[nocadre=false,lgt=1.3,espcl=2.5,deltacl=0.6]
				{$t$ /0.6, $y'$ /0.6, $y$ /2.5}
				{$0$, $-\ln{\frac{1}{5}}$, $+\infty$}
				\tkzTabLine{,+,0,-,}
				\tkzTabVar{-/ $694{,}4$,+/$1250$,-/$0$}
			\end{tikzpicture}
		\end{center}
		Hàm số đạt cực đại tại $t=-\ln{\dfrac{1}{5}}\approx1{,}6$.\\
		Vậy sau khi phát hành $1{,}6$ năm thì tốc độ bán hàng là lớn nhất.
	}
\end{bt}

\begin{bt}%[2D1H1-1]
% nguon: Archive/2024/0. [LVH] BG - 12 New - 4in1-20240625T224301Z-001/0. [LVH] BG - 12 New - 4in1/2. [LVH] TEX SGK 12 NEW-20240418T074833Z-001/2. [LVH] TEX SGK 12 NEW/[LVH] SGK-12-KNTT/data/C1/12-SGK-KNTT-1-BaiTapCuoiChuong.tex #1123

	Xét chiều biến thiên và tìm các cực trị (nếu có) của các hàm số sau 
	\begin{enumerate}[a) ]
		\item $y=x^3-3x^2+3x-1$. 
		\item $y=x^4-2x^2-1$.
		\item $y=\dfrac{2x-1}{3x+1}$. 
		\item $y=\dfrac{x^2+2x+2}{x+1}$.
	\end{enumerate}
	\loigiai{
		\begin{enumerate}[a) ]
			\item Tập xác định $D=\mathbb{R}$.\\
			Ta có $y'=3x^2-6x+3=3(x-1)^2 \ge 0$ với mọi $x \in \mathbb{R}$. \\ 
			Vậy, hàm số $y=x^3-3x^2+3x-1$ đồng biến trên $\mathbb{R}$. \\
			Hàm số $y=x^3-3x^2+3x-1$ không có cực trị.
			\item Tập xác định $D=\mathbb{R}$.\\
			Ta có $y'=4x^3-4x$. \\ 
			$y'=0\Leftrightarrow 4x^3-4x=0\Leftrightarrow \hoac{&x=0\\&x=\pm 1}$.\\
			Bảng biến thiên 
			\begin{center}
				\begin{tikzpicture}
					\tkzTabInit[nocadre=false,lgt=1,espcl=2.5]
					{$x$ /.7,$y'$ /.7,$y$ /2.5}{$-\infty$,$-1$,$0$,$1$,$+\infty$}
					\tkzTabLine{,-,$0$,+,$0$,-,$0$,+,}
					\tkzTabVar{+/ $+\infty$ ,-/$-2$,+/$-1$,-/$-2$,+/$+\infty$}
				\end{tikzpicture}
			\end{center}
			Hàm số đồng biến trên khoảng $\left(-1;0 \right)$ và $\left(1;+\infty \right)$.\\
			Hàm số nghịch biến trên khoảng $\left(-\infty;-1 \right)$ và $\left(0;1 \right)$.\\
			Hàm số đạt cực đại tại $x=0$ và $y_{CD}=-1$.\\
			Hàm số đạt cực tiểu tại $x=\pm 1$ và $y_{CT}=-2$.\\
			\item Tập xác định $D=\mathbb{R}\setminus \left\{-\dfrac{1}{3}\right\}$.\\
			Ta có $y'=\dfrac{2(3x+1)-3(2x-1)}{(3x+1)^2}=\dfrac{5}{(3x+1)^2}>0$ với mọi $x$ khác $- \dfrac{1}{3}$. \\ 
			Vậy, hàm số $y=\dfrac{2x-1}{3x+1}$ đồng biến trên $\left(-\infty; - \dfrac{1}{3}\right)$ và $\left(- \dfrac{1}{3};+\infty\right)$.\\
			Hàm số không có cực trị.
			\item Tập xác định $D=\mathbb{R}\setminus \{-1\}$.\\
			Ta có $y'=\dfrac{(2x+2)(x+1)-(x^2+2x+2)}{(x+1)^2}=\dfrac{x^2+2x}{(x+1)^2}$. \\ 
			$y'=0\Leftrightarrow \hoac{&x=0\\&x=-2}$ (thỏa mãn). \\
			Bảng biến thiên 
			\begin{center}
				\begin{tikzpicture}
					\tikzset{double style/.append style = {draw=\tkzTabDefaultWritingColor,double=\tkzTabDefaultBackgroundColor,double distance=2pt}}
					\tkzTabInit[nocadre=false,espcl=2.75]{$x$/.5,$f'(x)$/.7,$f(x)$/2.5}{$-\infty$,$-2$,$-1$,$0$,$+\infty$}
					\tkzTabLine{,+,z,-,d,-,z,+,} 
					\tkzTabVar{-/$-\infty$,+/$-2$,-D+/$-\infty$/$+\infty$,-/$2$,+/$+\infty$}
				\end{tikzpicture} 
			\end{center}
			Hàm số đồng biến trên khoảng $\left(-\infty;-2 \right)$ và $\left(0;+\infty \right)$. \\ 
			Hàm số nghịch biến trên khoảng $\left(-2;-1 \right)$ và $\left(-1;0 \right)$.\\
			Hàm số đạt cực đại tại $x=-2$ và $y_{CD}=-2$. Hàm số đạt cực tiểu tại $x=0$ và $y_{CT}=2$.
		\end{enumerate}
	}
\end{bt}

\begin{bt}%[2D1H3-2]
% nguon: Archive/2024/0. [LVH] BG - 12 New - 4in1-20240625T224301Z-001/0. [LVH] BG - 12 New - 4in1/2. [LVH] TEX SGK 12 NEW-20240418T074833Z-001/2. [LVH] TEX SGK 12 NEW/[LVH] SGK-12-KNTT/data/C1/12-SGK-KNTT-1-BaiTapCuoiChuong.tex #1124

	Tìm giá trị lớn nhất và giá trị nhỏ nhất (nếu có) của các hàm số sau
	\begin{enumerate}[a) ]
		\item $y=\dfrac{2x+1}{3x-2}$ trên nửa khoảng $\left[2;+\infty\right)$;
		\item $y=\sqrt{2-x^2}$.
	\end{enumerate}
	\loigiai{
		\begin{enumerate}[a) ]
			\item Ta có $y'=- \dfrac{7}{\left(3x-2 \right)^2}<0$ với mọi $x \in \left[2;+\infty\right)$.\\
			Suy ra $\underset{_{[2;+\infty)}}{\mathop{\max}}y=y(2)=\dfrac{2\cdot ^2+1}{3\cdot ^2-2}=\dfrac{5}{4}$, hàm số không có giá trị nhỏ nhất trên $\left[2;+\infty\right)$.
			\item Tập xác định $D=\left[-\sqrt{2};\sqrt{2}\right]$.\\
			Ta có $y'=\dfrac{-2x}{2\sqrt{2-x^2}}=\dfrac{-x}{\sqrt{2-x^2}}$. \\ 
			$y'=0\Leftrightarrow x=0$ (thỏa mãn).
			$y\left(-\sqrt{2} \right)=y \left(\sqrt{2} \right)=0$, $y \left(0 \right)=\sqrt{2}$.\\
			Vậy, $\underset{\left[_{-\sqrt{2};\sqrt{2}}\right]}{\mathop{\min}}y=y\left(-\sqrt{2} \right)=y \left(\sqrt{2} \right)=0$, $\underset{\left[_{-\sqrt{2};\sqrt{2}}\right]}{\mathop{\max}}y=y \left(0 \right)=\sqrt{2}$.
		\end{enumerate}
	}
\end{bt}

\begin{bt}%[2D1H4-1]
% nguon: Archive/2024/0. [LVH] BG - 12 New - 4in1-20240625T224301Z-001/0. [LVH] BG - 12 New - 4in1/2. [LVH] TEX SGK 12 NEW-20240418T074833Z-001/2. [LVH] TEX SGK 12 NEW/[LVH] SGK-12-KNTT/data/C1/12-SGK-KNTT-1-BaiTapCuoiChuong.tex #1125

	Tìm các tiệm cận của mỗi đồ thị hàm số sau
	\begin{enumerate}[a) ]
		\item $y=\dfrac{3x-2}{x+1}$.
		\item $y=\dfrac{x^2+2x-1}{2x-1}$.
	\end{enumerate}
	\loigiai{
		\begin{enumerate}[a) ]
			\item Ta có $\underset{x\to -1^{+}}{\mathop{\lim}}y=\underset{x\to -1^{+}}{\mathop{\lim}}\dfrac{3x-2}{x+1}=-\infty$; $\underset{x\to -1^{-}}{\mathop{\lim}}y=\underset{x\to -1^{-}}{\mathop{\lim}}\dfrac{3x-2}{x+1}=+\infty$. \\ 
			Vậy tiệm cận đứng của đồ thị hàm số $y=\dfrac{3x-2}{x+1}$ là đường thẳng $x=-1$. \\ 
			Ta có $\underset{x\to -\infty}{\mathop{\lim}}y=\underset{x\to -\infty}{\mathop{\lim}}\dfrac{3x-2}{x+1}=3;\underset{x\to +\infty}{\mathop{\lim}}y=\underset{x\to +\infty}{\mathop{\lim}}\dfrac{3x-2}{x+1}=3$.\\
			Vậy tiệm cận ngang của đồ thị hàm số $y=\dfrac{3x-2}{x+1}$ đường thẳng $y=3$.\\
			\item Ta có $\underset{x\to{\left(\dfrac{1}{2}\right)^{+}}}{\mathop{\lim}}y=\underset{x\to{\left(\dfrac{1}{2}\right)^{+}}}{\mathop{\lim}}\dfrac{x^2+2x-1}{2x-1}=+\infty;\underset{x\to{\left(\dfrac{1}{2}\right)^{-}}}{\mathop{\lim}}y=\underset{x\to{\left(\dfrac{1}{2}\right)^{-}}}{\mathop{\lim}}\dfrac{x^2+2x-1}{2x-1}=-\infty$. \\ 
			Vậy tiệm cận đứng của đồ thị hàm số $y=\dfrac{x^2+2x-1}{2x-1}$ là đường thẳng $x=\dfrac{1}{2}$.\\
			Ta có $\underset{x\to -\infty}{\mathop{\lim}}y=\underset{x\to -\infty}{\mathop{\lim}}\dfrac{x^2+2x-1}{2x-1}=-\infty;\underset{x\to +\infty}{\mathop{\lim}}y=\underset{x\to +\infty}{\mathop{\lim}}\dfrac{x^2+2x-1}{2x-1}=+\infty$
			Vậy đồ thị hàm số $y=\dfrac{x^2+2x-1}{2x-1}$ không có tiệm cận ngang.\\
			Ta có $y=\dfrac{x^2+2x-1}{2x-1}=\dfrac{x}{2}+\dfrac{5}{4}+\dfrac{1}{4(2x-1)}$. \\ 
			$\underset{x\to +\infty}{\mathop{\lim}}\left[y-\left(\dfrac{x}{2}+\dfrac{5}{4}\right)\right]=\underset{x\to +\infty}{\mathop{\lim}}\dfrac{1}{4(2x-1)}=0$, $\underset{x\to -\infty}{\mathop{\lim}}\left[y-\left(\dfrac{x}{2}+\dfrac{5}{4}\right)\right]=\underset{x\to -\infty}{\mathop{\lim}}\dfrac{1}{4(2x-1)}=0$. \\ 
			Vậy tiệm cận xiên của đồ thị hàm số $y=\dfrac{x^2+2x-1}{2x-1}$ là đường thẳng $y=\dfrac{x}{2}+\dfrac{5}{4}$.
		\end{enumerate}
	}
\end{bt}

\begin{bt}%[2D1H1-1]
% nguon: Archive/2024/0. [LVH] BG - 12 New - 4in1-20240625T224301Z-001/0. [LVH] BG - 12 New - 4in1/2. [LVH] TEX SGK 12 NEW-20240418T074833Z-001/2. [LVH] TEX SGK 12 NEW/[LVH] SGK-12-KNTT/data/C1/12-SGK-KNTT-1-BaiTapCuoiChuong.tex #1126

	Khảo sát sự biến thiên và vẽ đồ thị của các hàm số sau
	\begin{enumerate}[a) ]
		\item $y=-x^3+6x^2-9x+12$.
		\item $y=\dfrac{2x-1}{x+1}$.
		\item $y=\dfrac{x^2-2x}{x-1}$. 
	\end{enumerate}
	\loigiai{
		\begin{enumerate}[a)]
			\item Tập xác định $D=\mathbb{R}$. \\ 
			$\bullet$ Sự biến thiên. \\
			Ta có $y'=-3x^2+12x-9;y'=0\Leftrightarrow -3x^2+12x-9=0\Leftrightarrow \hoac{&x=1\\&x=3}$.\\
			$\underset{x\to -\infty}{\mathop{\lim}}y=\underset{x\to -\infty}{\mathop{\lim}}(-x^3+6x^2-9x+12)=\underset{x\to -\infty}{\mathop{\lim}}\left[x^3\left(-1+\dfrac{6}{x}-\dfrac{9}{x^2}+\dfrac{12}{x^3}\right)\right]=+\infty$. \\
			$\underset{x\to +\infty}{\mathop{\lim}}y=\underset{x\to +\infty}{\mathop{\lim}}(-x^3+6x^2-9x+12)=\underset{x\to +\infty}{\mathop{\lim}}\left[x^3\left(-1+\dfrac{6}{x}-\dfrac{9}{x^2}+\dfrac{12}{x^3}\right)\right]=-\infty$. \\ 
			Bảng biến thiên
			\begin{center}
				\begin{tikzpicture}
					\tkzTabInit[nocadre=false,lgt=1.5,espcl=2.5]
					{$x$/1,$y '(x)$ /1, $y (x)$ /2}
					{$-\infty$ , $1$ , $3$ , $+\infty$}
					\tkzTabLine{ ,-, 0 ,+, 0 ,-, }
					\tkzTabVar {+/$+\infty$,-/$8$ ,+/$12$,-/$-\infty$}
				\end{tikzpicture}
			\end{center}
			Hàm số đồng biến trên $\left(1;3\right)$, hàm số nghịch biến trên khoảng $\left(-\infty;1\right)$ và $\left(3;+\infty\right)$.\\
			Hàm số đạt cực đại tại $x=3$, giá trị cực đại $y=12$.\\
			Hàm số đạt cực tiểu tại $x=1$, giá trị cực tiểu $y_{CT}=8$.\\ 
			$\bullet$ Đồ thị. \\
			Giao điểm của đồ thị hàm số $y=-x^3+6x^2-9x+12$ với trục tung là $\left(0; 12 \right)$.\\
			Đồ thị hàm số $y=-x^3+6x^2-9x+12$ đi qua các điểm $\left(1;8\right)$, $\left(3;12\right)$, $\left(4;8 \right)$.\\
			Đồ thị hàm số có tâm đối xứng là điểm $\left(2; 10 \right)$.
			\begin{center}
				\begin{tikzpicture}[scale=0.7, line join=round, line cap=round, >=stealth]
					\tikzset{every node/.style={scale=1}}
					\def\xmin{-1}\def\xmax{5}\def\ymin{-1}\def\ymax{13}
					\draw[->] (\xmin-0.2,0)--(\xmax+0.2,0) node[below]{$x$};
					\draw[->] (0,\ymin-0.2)--(0,\ymax+0.2) node[right]{$y$};
					\draw (0,0) node[below left]{$O$};
					\foreach \x in {,1,2,3,4,}\draw (\x,0.1)--(\x,-0.1) node[below]{$\x$};
					\foreach \y in {,1,2,3,4,5,6,7,8,9,10,11,12,13}\draw (0.1,\y)--(-0.1,\y) node[left]{$\y$};
					\clip (\xmin,\ymin) rectangle (\xmax,\ymax);
					\draw[smooth,samples=200,domain=\xmin:\xmax] plot (\x,{-1*((\x)^3)+6*((\x)^2)+-9*(\x)+12});
					\draw[dashed] (2,0)--(2,10)--(0,10);\fill (2,10)circle(1.5pt);
					\draw[dashed] (1,0)--(1,8)--(0,8);\fill (1,8)circle(1.5pt);
					\draw[dashed] (3,0)--(3,12)--(0,12);\fill (3,12)circle(1.5pt);
				\end{tikzpicture}
			\end{center}
			\item Tập xác định $D=\mathbb{R}\setminus \{-1\}$.\\
			$\bullet$ Sự biến thiên.
			$y'=\dfrac{3}{(x+1)^2}>0$ với mọi $x$ khác $-1$. \\ 
			Hàm số đồng biến trên khoảng $\left(-\infty;-1\right)$ và $\left(-1;+\infty\right)$. \\ 
			Hàm số không có cực trị.\\
			Giới hạn \\ 
			$\underset{x\to +\infty}{\mathop{\lim}}y=\underset{x\to +\infty}{\mathop{\lim}}\dfrac{2x-1}{x+1}=2;\underset{x\to -\infty}{\mathop{\lim}}y=\underset{x\to -\infty}{\mathop{\lim}}\dfrac{2x-1}{x+1}=2$. \\ 
			$\underset{x\to -1^{-}}{\mathop{\lim}}y=\underset{x\to -1^{-}}{\mathop{\lim}}\dfrac{2x-1}{x+1}=+\infty;\underset{x\to -1^{+}}{\mathop{\lim}}y=\underset{x\to -1^{+}}{\mathop{\lim}}\dfrac{2x-1}{x+1}=-\infty$. \\
			Do đó, đồ thị hàm số nhận đường thẳng $x=-1$ làm tiệm cận đứng và đường thẳng $y=2$ làm tiệm cận ngang.\\
			Bảng biến thiên
			\begin{center}
				\begin{tikzpicture}[scale=.8]
					\tikzset{double style/.append style = {draw=\tkzTabDefaultWritingColor,double=\tkzTabDefaultBackgroundColor,double distance=2pt}}
					\tikzset{h style/.style = {pattern=north west lines}}
					\tkzTabInit[nocadre=false,lgt=1.5,espcl=3.5]
					{$x$ /1, $f'(x)$ /1,$f(x)$ /2}
					{$-\infty$,$-1$,$+\infty$}
					\tkzTabLine{  ,+,d ,+ , }
					\tkzTabVar{-/$2$,+D-/$+\infty$/$-\infty$,+/$2$ }
				\end{tikzpicture}
			\end{center}
			$\bullet$ Đồ thị. \\ 
			Giao điểm của đồ thị hàm số với trục tung là $\left(0;-1 \right)$.\\
			$y=0 \Leftrightarrow \dfrac{2x-1}{x+1}=0\Leftrightarrow x=\dfrac{1}{2}$. \\ 
			Giao điểm của đồ thị hàm số với trục hoành là điểm $\left(\dfrac{1}{2};0\right)$.\\
			Đồ thị hàm số nhận $I\left(-1; 2 \right)$ làm tâm đối xứng.
			\begin{center}
				\begin{tikzpicture}[scale=0.7, line join=round, line cap=round, >=stealth]
					\tikzset{every node/.style={scale=1}}
					\def\xmin{-5}\def\xmax{3}\def\ymin{-2}\def\ymax{5}
					\draw[->] (\xmin-0.2,0)--(\xmax+0.2,0) node[below]{$x$};
					\draw[->] (0,\ymin-0.2)--(0,\ymax+0.2) node[right]{$y$};
					\draw (0,0) node[below left]{$O$};
					\foreach \x in {-5, -4,-3,-2,-1,1,2}\draw (\x,0.1)--(\x,-0.1) node[below]{$\x$};
					\foreach \y in {-1,1,2,3,4,5}\draw (0.1,\y)--(-0.1,\y) node[left]{$\y$};
					\clip (\xmin,\ymin) rectangle (\xmax,\ymax);
					\draw[dashed] (\xmin,2.0)--(\xmax,2.0);
					\draw[dashed] (-1.0,\ymin)--(-1.0,\ymax);
					\draw[smooth,samples=200,domain=\xmin:-1.1] plot (\x,{(2*(\x)+-1)/(1*(\x)+1)});
					\draw[smooth,samples=200,domain=-0.9:\xmax] plot (\x,{(2*(\x)+-1)/(1*(\x)+1)});
					\draw[dashed] (-1,0)--(-1,2)--(0,2);\fill (-1,2)circle(1.5pt);
				\end{tikzpicture}
			\end{center}
			\item Tập xác định $D=\mathbb{R}\setminus \{-1\}$. \\ 
			$\bullet$ Sự biến thiên. \\ 
			Ta có $y=\dfrac{x^2-2x}{x-1}=x-1-\dfrac{1}{x-1}$. \\ 
			$y'=\dfrac{(2x-2)(x-1)-(x^2-2x)}{(x-1)^2}=\dfrac{x^2-2x+2}{(x-1)^2}=\dfrac{(x-1)^2+1}{(x-1)^2}>0$ với mọi $x$ khác $1$. \\ 
			Do đó, hàm số đồng biến trong khoảng $\left(-\infty;1 \right)$ và $\left(1;+\infty \right)$. \\
			Hàm số không có cực trị.\\
			Giới hạn \\ 
			$\underset{x\to +\infty}{\mathop{\lim}}y=\underset{x\to +\infty}{\mathop{\lim}}\dfrac{x^2-2x}{x-1}=+\infty;\underset{x\to -\infty}{\mathop{\lim}}y=\underset{x\to -\infty}{\mathop{\lim}}\dfrac{x^2-2x}{x-1}=-\infty$. \\ 
			$\underset{x\to{1^{-}}}{\mathop{\lim}}y=\underset{x\to{1^{-}}}{\mathop{\lim}}\dfrac{x^2-2x}{x-1}=+\infty;\underset{x\to{1^{+}}}{\mathop{\lim}}y=\underset{x\to{1^{+}}}{\mathop{\lim}}\dfrac{x^2-2x}{x-1}=-\infty$. \\ 
			$\underset{x\to +\infty}{\mathop{\lim}}[y-(x-1)]=\underset{x\to +\infty}{\mathop{\lim}}\left(x-1-\dfrac{1}{x-1}-(x-1)\right)=\underset{x\to +\infty}{\mathop{\lim}}-\dfrac{1}{x-1}=0$. \\ 
			$\underset{x\to -\infty}{\mathop{\lim}}[y-(x-1)]=\underset{x\to -\infty}{\mathop{\lim}}\left(x-1-\dfrac{1}{x-1}-(x-1)\right)=\underset{x\to -\infty}{\mathop{\lim}}-\dfrac{1}{x-1}=0$. \\ 
			Do đó, đồ thị hàm số nhận đường thẳng $x=1$ làm tiệm cận đứng và đường thẳng $y=x-1$ làm tiệm cận xiên.\\
			Bảng biến thiên
			\begin{center}
				\begin{tikzpicture}[scale=.8]
					\tikzset{double style/.append style = {draw=\tkzTabDefaultWritingColor,double=\tkzTabDefaultBackgroundColor,double distance=2pt}}
					\tikzset{h style/.style = {pattern=north west lines}}
					\tkzTabInit[nocadre=false,lgt=1.5,espcl=3.5]
					{$x$ /1, $f'(x)$ /1,$f(x)$ /2}
					{$-\infty$,$1$,$+\infty$}
					\tkzTabLine{  ,+,d ,+ , }
					\tkzTabVar{-/$-\infty$,+D-/$+\infty$/$-\infty$,+/$+\infty$ }
				\end{tikzpicture}
			\end{center}
			$\bullet$ Đồ thị. \\ 
			Giao điểm của đồ thị hàm số với trục tung là $\left(0;0 \right)$.\\
			$y=0\Leftrightarrow \dfrac{x^2-2x}{x-1}=0\Leftrightarrow \hoac{&x=0 \\&x=2}$.\\
			Đồ thị hàm số giao với $Ox$ tại các điểm $\left(0;0 \right)$ và $\left(2;0\right)$.\\
			Đồ thị hàm số nhận $I\left(1; 0 \right)$ làm tâm đối xứng.
			\begin{center}
				\begin{tikzpicture}[scale=0.7, line join=round, line cap=round, >=stealth]
					\tikzset{every node/.style={scale=1}}
					\def\xmin{-3}\def\xmax{5}\def\ymin{-4}\def\ymax{4}
					\draw[->] (\xmin-0.2,0)--(\xmax+0.2,0) node[below]{$x$};
					\draw[->] (0,\ymin-0.2)--(0,\ymax+0.2) node[right]{$y$};
					\draw (0,0) node[below left]{$O$};
					\foreach \x in {-3,-2,-1,1,2,3,4,5}\draw (\x,0.1)--(\x,-0.1) node[below]{$\x$};
					\foreach \y in {-4,-3,-2,-1,1,2,3,4}\draw (0.1,\y)--(-0.1,\y) node[left]{$\y$};
					\clip (\xmin,\ymin) rectangle (\xmax,\ymax);
					\draw[dashed] (1.0,\ymin)--(1.0,\ymax);
					\draw[dashed,domain=\xmin:\xmax] plot (\x,{1.0*(\x)+-1.0});
					\draw[smooth,samples=200,domain=\xmin:0.9] plot (\x,{(1*((\x)^2)+-2*(\x)+0)/(1*(\x)+-1)});
					\draw[smooth,samples=200,domain=1.1:\xmax] plot (\x,{(1*((\x)^2)+-2*(\x)+0)/(1*(\x)+-1)});
				\end{tikzpicture}
			\end{center}
		\end{enumerate}
	}
\end{bt}

\begin{bt}%[2D1H1-1]
% nguon: Archive/2024/0. [LVH] BG - 12 New - 4in1-20240625T224301Z-001/0. [LVH] BG - 12 New - 4in1/2. [LVH] TEX SGK 12 NEW-20240418T074833Z-001/2. [LVH] TEX SGK 12 NEW/[LVH] SGK-12-KNTT/data/C1/12-SGK-KNTT-1-BaiTapCuoiChuong.tex #1127

	Xét một thấu kính hội tụ có tiêu cự $f$ (H.1.39). Khoảng cách $p$ từ vật đến thấu kính liên hệ với khoảng cách $q$ từ ảnh đến thấu kính bởi hệ thức: $\dfrac{1}{p}+\dfrac{1}{q}=\dfrac{1}{f}$.
	\begin{center}
		\begin{tikzpicture}[line join=round, line cap=round,>=stealth]
			\tikzset{every node/.style={scale=0.9}}
			\def\f{3}
			\def\p{6}
			\def\h{1.5}
			\pgfmathsetmacro{\q}{1/(1/\f-1/\p)};
			\pgfmathsetmacro{\goc}{atan(\h/\p)};
			\path (-\p,0) coordinate (M)
			(\q,0) coordinate (M')
			(0,\h ) coordinate (A)
			(0,-\h) coordinate (B);
			\foreach\x/\y in {1/0.1, 2/0.25, 3/0.5, 4/0.75, 5/0.9}
			{
				\path ($(A)!\y!(B)$) coordinate (a\x);
				\draw[->] (M)--($(M)!0.8!(a\x)$);
				\draw ($(M)!0.8!(a\x)$)--(a\x);
				\draw[->] (M')--($(M')!0.8!(a\x)$);
				\draw ($(M')!0.8!(a\x)$)--(a\x);
			}
			\draw[dashed] (M) node[left]{$M$}--($(M)+(0,{-\h-0.5})$) (0,-\h)--(0,{-\h-0.5}) 	(M') node[right]{$M'$}--($(M')+(0,{-\h-0.5})$) (-\f,0) node[right,circle,inner sep=1pt,fill=white]{$F$}--(-\f,{\h+0.5}) (\f,0) node[left,circle, inner sep=1pt,fill=white]{$F'$}--(\f,{\h+0.5}) (0,\h)--(0,{\h+0.5});
			\draw[<->] ($(M)+(0,{-\h-0.5})$) --(0,{-\h-0.5}) node[below,midway]{$p$} ;
			\draw[<->] ($(M')+(0,{-\h-0.5})$) --(0,{-\h-0.5}) node[below,midway]{$q$} ;
			\draw[<->] (-\f,{\h+0.5})--(0,{\h+0.5}) node[above,midway]{$f$} ;
			\draw[<->] (\f,{\h+0.5})--(0,{\h+0.5}) node[above,midway]{$f$} ;
			\draw[fill=cyan] let \p1=($(M)-(0,\h)$),
			\p2=($(M')-(0,\h)$)
			in (0,\h) arc({\goc}:{-\goc}:{veclen(\p1)}) arc ({-180+\goc}:{-180-\goc}:{veclen(\p2)});
			
		\end{tikzpicture} \\
		Hình 1.39
	\end{center}
	\begin{enumerate}[a) ]
		\item Viết công thức tính $q=g\left(p \right)$ như một hàm số của biến $p\in \left(f;+\infty \right)$.
		\item Tính các giới hạn $\underset{p\to +\infty}{\mathop{\lim}}g(p)$, $\underset{p\to{f^{+}}}{\mathop{\lim}}g(p)$ và giải thích ý nghĩa các kết quả này.
		\item Lập bảng biến thiên của hàm số $q=g \left(p \right)$ trên khoảng $\left(f;+\infty \right)$. 
	\end{enumerate}
	\loigiai{
		\begin{enumerate}[a) ]
			\item Ta có $\dfrac{1}{p}+\dfrac{1}{q}=\dfrac{1}{f} \Rightarrow q=\dfrac{pf}{p-f}$. Do đó, $q=g \left(p \right)=\dfrac{pf}{p-f}$ với $p\in (f;+\infty)$. 
			\item $\underset{p\to +\infty}{\mathop{\lim}}g(p)=\underset{p\to +\infty}{\mathop{\lim}}\dfrac{pf}{p-f}=\underset{p\to +\infty}{\mathop{\lim}}\dfrac{f}{1-\dfrac{f}{p}}=f$, $\underset{p\to{f^{+}}}{\mathop{\lim}}g(p)=\underset{p\to{f^{+}}}{\mathop{\lim}}\dfrac{pf}{p-f}=+\infty$. \\ 
			Ý nghĩa của $\underset{p\to +\infty}{\mathop{\lim}}g(p)=f$: Khoảng cách từ vật đến thấu kính tiến ra vô cùng thì khoảng cách từ ảnh đến thấu kính xấp xỉ tiêu cự.\\
			Ý nghĩa của $\underset{p\to{f^{+}}}{\mathop{\lim}}g(p)=+\infty$: Khoảng cách từ vật đến thấu kính tiến gần về tiêu cự $f$ thì khoảng cách từ ảnh đến thấu kính là càng lớn.
			\item Ta có $q'=g' \left(p \right)=- \dfrac{f^2}{(p-f)^2}<0$ với mọi $p\in \left(f;+\infty \right)$ nên hàm số nghịch biến trên $ \left(f;+\infty \right)$. \\
			Bảng biến thiên
			\begin{center}
				\begin{tikzpicture}
					\tkzTabInit[nocadre=false,espcl=6]
					{$x$/1,$g'(p)$/1,$g (p)$/2}{$f$, $+\infty$}
					\tkzTabLine{,-,}
					\tkzTabVar{+/  $+\infty$,-/  $f$ } 
				\end{tikzpicture}
			\end{center}
		\end{enumerate}
	}
\end{bt}

\begin{bt}%[2D1H1-1]
% nguon: Archive/2024/0. [LVH] BG - 12 New - 4in1-20240625T224301Z-001/0. [LVH] BG - 12 New - 4in1/2. [LVH] TEX SGK 12 NEW-20240418T074833Z-001/2. [LVH] TEX SGK 12 NEW/[LVH] SGK-12-KNTT/data/C1/12-SGK-KNTT-1-BaiTapCuoiChuong.tex #1128

	Dân số của một quốc gia sau $t$ (năm) kể từ năm $2023$ được ước tính bởi công thức $N(t)=100\mathrm{e}^{0, 012t}$ ($N(t)$ được tính bằng triệu người, $0\le t\le 50$).
	\begin{enumerate}[a) ]
		\item Ước tính dân số của quốc gia này vào các năm $2030$ và $2035$ (kết quả tính bằng triệu người, làm tròn kết quả đến chữ số thập phân thứ ba). 
		\item Xem $N(t)$ là hàm số của biến số $t$ xác định trên đoạn $\left[0;50\right]$. Xét chiều biến thiên của hàm số $N(t)$ trên đoạn $\left[0;50\right]$. 
		\item Đạo hàm của hàm số $N(t)$ biểu thị tốc độ tăng dân số của quốc gia đó (tính bằng triệu người/ năm). Vào năm nào tốc độ tăng dân số của quốc gia đó là $1{,}6$ triệu người/ năm?
	\end{enumerate}
	\loigiai{
		\begin{enumerate}[a) ]
			\item Dân số của quốc gia vào năm $2030$ là $N(7)=100\mathrm{e}^{0{,}012\cdot ^7}=100\mathrm{e}^{0{,}084}=108{,}763$ (triệu người). \\ 
			Dân số của quốc gia vào năm $2035$ là $N(12)=100 \mathrm{e}^{0{,}012 \cdot 12}=100 \mathrm{e}^{0{,}144}=115,488$ (triệu người). 
			\item Trên đoạn $\left[0;50\right]$ ta có $N'(t)=0{,}012\cdot 100 \mathrm{e}^{0{,}012t}=1,2 \mathrm{e}^{0{,}012t}>0$ với mọi $t \in \left[0;50\right]$. \\ 
			Do đó, hàm số $N(t)$ đồng biến trên đoạn $\left[0;50\right]$.
			\item Ta có $N'(t)=1{,}2 \mathrm{e}^{0{,}012t}$. \\ 
			Với tốc độ tăng dân số của quốc gia đó là $1,6$ triệu người/ năm ta có \\ 
			$1{,}6=1{,}2\mathrm{e}^{0{,}012t}\Leftrightarrow{\mathrm{e}^{0{,}012t}}=\dfrac{4}{3}\Leftrightarrow t=\dfrac{250\ln \dfrac{4}{3}}{3}\approx 23{,}97$. \\ 
			Vậy vào năm $2046$ thì tốc độ tăng dân số của quốc gia đó là $1{,}6$ triệu người/ năm.
		\end{enumerate}
	}
\end{bt}

\begin{bt}%[2D1H1-1]
% nguon: Archive/2024/0. [LVH] BG - 12 New - 4in1-20240625T224301Z-001/0. [LVH] BG - 12 New - 4in1/2. [LVH] TEX SGK 12 NEW-20240418T074833Z-001/2. [LVH] TEX SGK 12 NEW/[LVH] SGK-12-KNTT/data/C1/12-SGK-KNTT-1-BaiTapCuoiChuong.tex #1129

	Một đường dây điện được nối từ một nhà máy điện ở $A$ đến một hòn đảo ở $C$ như Hình 1.40. Khoảng cách từ $C$ đến $B$ là $4$ $\mathrm{km}$. Bờ biển chạy thẳng từ $A$ đến $B$ với khoảng cách là $10$ $\mathrm{km}$. Tổng chi phí lắp đặt cho $1$ $\mathrm{km}$ dây điện trên biển là $50$ triệu đồng, còn trên đất liền là $30$ triệu đồng. Xác định vị trí điểm $M$ trên đoạn $AB$ (điểm nối dây từ đất liền ra đảo) để tổng chi phí lắp đặt là nhỏ nhất.
	\begin{center}
		\begin{tikzpicture}[scale=1, font=\footnotesize, line join=round, line cap=round, >=stealth]
			%%%%%%%%%%%%%%%%%%%%%%%%%%%%%%%%%
			\fill (10,0) circle(1pt) node[right]{$A$};
			\fill (0,0) circle(1pt) node[left]{$B$};
			\fill (3,0) circle(1pt) node[below]{$M$};
			\fill (0,4) circle(1pt) node[above]{$C$};
			\fill (0,2) circle(0pt) node[left]{$4 \mathrm{km}$};
			%%%%%%%%%%%%%%%%%%%%%%%%%%%%%%%%%
			\draw (0, 0) -- (0, 4) -- (3, 0) -- (10, 0) -- (0, 0); 
			%%%%%%%%%%%%%%%%%%%%%%%%%%%%%%%%%
		\end{tikzpicture} \\
		Hình 1.40
	\end{center}
	\loigiai{
		Đặt $MB=x(\mathrm{km}, 0\le x\le 10)$ $\Rightarrow AM=10-x$ (km), $MC=\sqrt{MB^2+CB^2}=\sqrt{x^2+16}(\mathrm{km})$.\\
		Khi đó, chi phí nối điện từ $A$ đến $C$ là $f(x)=30(10-x)+50\sqrt{x^2+16}$ (triệu đồng). \\ 
		Ta có $f'(x)=-30+\dfrac{50x}{\sqrt{x^2+16}}$. \\ 
		$f'(x)=0\Leftrightarrow \dfrac{x}{\sqrt{x^2+16}}=\dfrac{3}{5}\Leftrightarrow 25x^2=9x^2+144\Leftrightarrow x=3$ (do $0 \le x \le 10$).\\
		Ta có $f(0)=500$, $f(3)=460$, $f(10)=100 \sqrt{29}$ nên chi phí nhỏ nhất là $460$ triệu đồng khi $x=3$.\\
		Vậy $M$ cách $B$ một khoảng $3$ $\mathrm{km}$ trên đoạn $AB$ (điểm nối dây từ đất liền ra đảo) thì tổng chi phí lắp đặt là nhỏ nhất.
	}
\end{bt}

\begin{bt}%[2D1V1-5]
% nguon: Archive/2024/0. [LVH] BG - 12 New - 4in1-20240625T224301Z-001/0. [LVH] BG - 12 New - 4in1/2. [LVH] TEX SGK 12 NEW-20240418T074833Z-001/2. [LVH] TEX SGK 12 NEW/[LVH] SGK-12-KNTT/data/SCD/12-SCD-KNTT-1-2.tex #1170

	Tại một nhà máy sản xuất linh kiện điện tử, các linh kiện được xếp vào từng hộp một cách độc lập, mỗi hộp $10$ linh kiện. Hộp được xếp loại $I$ nếu hộp đó có nhiều nhất một linh kiện không đạt tiêu chuẩn. Biết rằng xác suất để nhà máy sản xuất ra một linh kiện điện tử không đạt tiêu chuẩn là $0{,}01$. Hỏi tỉ lệ những hộp linh kiện điện tử loại $I$ là bao nhiêu?
	\loigiai{
		Gọi $X$ là số linh kiện không đạt tiêu chuẩn thì $X$ là một biến ngẫu nhiên có phân bố nhị thức với tham số $n=10$, $p=0{,}01$.\\
		Hộp được xếp loại $I$ nếu hộp đó có nhiều nhất một linh kiện không đạt tiêu chuẩn tức là $X\le 1$.\\
		Theo chú ý về phân bố nhị thức ta có:
		$$\mathrm{P}(X\le 1)=\mathrm{C}^0_{10}\cdot \left(0{,}01\right)^0\cdot \left(0{,}99\right)^{10}+\mathrm{C}^1_{10}\cdot \left(0{,}01\right)^1\cdot \left(0{,}99\right)^{9}\approx 0{,}996.$$
		Vậy tỉ lệ những hộp linh kiện điện tử loại $I$ là $99{,}6\%$.
	}
\end{bt}

\begin{bt}%[2D1V1-5]
% nguon: Archive/2024/0. [LVH] BG - 12 New - 4in1-20240625T224301Z-001/0. [LVH] BG - 12 New - 4in1/2. [LVH] TEX SGK 12 NEW-20240418T074833Z-001/2. [LVH] TEX SGK 12 NEW/[LVH] SGK-12-KNTT/data/SCD/12-SCD-KNTT-1-2.tex #1171

	Một bài thi trắc nghiệm gồm $10$ câu hỏi, mỗi câu có $4$ phương án trả lời, trong đó chỉ có một phương án đúng. Mỗi câu trả lời đúng được $4$ điểm, mỗi câu trả lời sai trừ $1$ điểm. Một thí sinh làm bài bằng cách ở mỗi câu hỏi chọn ngẫu nhiên một phương án trả lời. Tính xác suất để thí sinh đó sau khi hoàn thành hết $10$ câu trong bài thi, có kết quả:
	\begin{enumerate}[a)]
		\item $15$ điểm;
		\item Bị điểm âm.
	\end{enumerate}
	\loigiai{
		Gọi $X$ là số câu trả lời đúng của thí sinh. $X$ là một biến ngẫu nhiên có phân bố nhị thức với tham số $n=10$, $p=\dfrac{1}{4}$.
		\begin{enumerate}[a)]
			\item Thí sinh đạt $15$ điểm thì có $5$ câu trả lời đúng và $5$ câu trả lời sai, tức là $X=5$. Khi đó, xác suất là:
			$$\mathrm{P}=\mathrm{C}^5_{10}\cdot \left(\dfrac{1}{4}\right)^5\cdot \left(\dfrac{3}{4}\right)^5\approx 0{,}058.$$
			\item Thí sinh bị điểm âm tức là thí sinh trả lời đúng nhiều nhất $1$ câu, tức là $X\le 1$.\\
			Theo chú ý về phân bố nhị thức ta có:
			$$\mathrm{P}(X\le 1)=\mathrm{C}^0_{10}\cdot \left(\dfrac{1}{4}\right)^0\cdot \left(\dfrac{3}{4}\right)^{10}+\mathrm{C}^1_{10}\cdot \left(\dfrac{1}{4}\right)^1\cdot \left(\dfrac{3}{4}\right)^{9}\approx 0{,}24.$$
		\end{enumerate}
	}
\end{bt}

\begin{bt}%[2D1V1-5]
% nguon: Archive/2024/0. [LVH] BG - 12 New - 4in1-20240625T224301Z-001/0. [LVH] BG - 12 New - 4in1/2. [LVH] TEX SGK 12 NEW-20240418T074833Z-001/2. [LVH] TEX SGK 12 NEW/[LVH] SGK-12-KNTT/data/SCD/12-SCD-KNTT-1-2.tex #1172

	Trong một trò chơi, mỗi ván người chơi gieo đồng thời $3$ xúc xắc cân đối, đồng chất. Nếu có ít nhất $2$ xúc xắc xuất hiện mặt $6$ chấm thì người chơi giành chiến thắng ván chơi đó. Bác Hưng tham gia chơi $3$ ván. Tính xác suất để bác Hưng thẳng ít nhất $2$ ván.
	\loigiai{
		Xác suất để một con xúc xắc xuất hiện mặt $6$ chấm là $\dfrac{1}{6}$.\\
		Gọ $X$ là số con xúc xắc xuất hiện mặt $6$ chấm.\\
		Bác Hưng thắng cuộc $1$ ván khi $X\ge 2$.\\
		Xác suất để bác Hưng thắng cuộc một ván là
		$$\mathrm{P}(X\ge 2)=\mathrm{C}^2_3\cdot \left(\dfrac{1}{6}\right)^2\cdot \dfrac{3}{4}+\left(\dfrac{1}{6}\right)^3=\dfrac{2}{27}.$$
		Gọi $Y$ là số ván thắng của bác Hưng.\\
		Xác suất để bác Hưng thắng ít nhất $2$ ván là
		$$\mathrm{P}(Y\ge 2)=\mathrm{C}^2_3\cdot \left(\dfrac{2}{27}\right)^2\cdot \dfrac{25}{27}+\left(\dfrac{2}{27}\right)^3=\dfrac{308}{19683}.$$
	}
\end{bt}

\begin{bt}%[2D1V1-5]
% nguon: Archive/2024/0. [LVH] BG - 12 New - 4in1-20240625T224301Z-001/0. [LVH] BG - 12 New - 4in1/2. [LVH] TEX SGK 12 NEW-20240418T074833Z-001/2. [LVH] TEX SGK 12 NEW/[LVH] SGK-12-KNTT/data/SCD/12-SCD-KNTT-1-2.tex #1173

	Màu hạt của đậu Hà Lan có hai kiểu hình: màu vàng và màu xanh. Có hai gene ứng với hai kiểu hình này là allele trội $A$ và allele lặn $a$. Khi cho lai hai cây đậu Hà Lan, cây con lấy ngẫu nhiên một gene từ cây bố và một gene từ cây mẹ để hình thành một cặp gene.\\
	Bốn bạn An, Bình, Sơn và Dương, mỗi bạn độc lập với nhau, thực hiện phép thử là lai hai cây đậu Hà Lan, trong đó cây bố có kiểu gene là $Aa$, cây mẹ có kiểu gene là $Aa$.\\
	Gọi $X$ là số cây con có hạt màu vàng trong số $4$ cây con.
	\begin{enumerate}[a)]
		\item Lập bảng phân bố xác suất của $X$.
		\item Hỏi trung bình có bao nhiêu cây con có hạt màu xanh?
	\end{enumerate}
	
	\loigiai{
		Ta vẽ sơ đồ hình cây để mô tả các kết quả có thể của kiểu gene ứng với màu hạt của cây con.
		\begin{center}
			\begin{tikzpicture}
				\def\gocm{20}
				\def\gocn{10}
				\def\r{4}
				\tikzset{s/.style={outer sep=0.5 mm,draw=magenta,rectangle,minimum width=2.75cm,text width=2.75cm,align=center,rounded corners=1mm,minimum height=1cm}}
				\path(0,0)node(O){}++(\gocm:\r)node[s](A1){A}++(\gocn:\r)node[s](A2){A};
				\path(A1)++({-\gocn}:\r)node[s](a2){a};
				\path(O)++(-\gocm:\r)node[s](B1){a}++(\gocn:\r)node[s](B2){A};
				\path(B1)++({-\gocn}:\r)node[s](b2){a};
				\foreach \x/\y in {
					O/A1,A1/A2,O/B1,B1/B2,A1/a2,B1/b2}
				\draw[-stealth](\x.east)--(\y.west);
				\path(O)--(A1.west)node[pos=0.5,above,sloped]{}(O)--(B1.west)node[pos=0.5,below,sloped]{}(B1.east)--(B2.west)node[pos=0.5,above,sloped]{}(A1.east)--(A2.west)node[pos=0.5,above,sloped]{}
				(A1.east)--(a2.west)node[pos=0.5,below,sloped]{}
				(B1.east)--(b2.west)node[pos=0.5,below,sloped]{};
				%%Node dòng trên
				\path(A2)++(0,1)node{\textbf{Cây mẹ}}++(180:4)node{\textbf{Cây bố}};
				\node[left] at (O) {\textbf{GỐC}};
			\end{tikzpicture}
		\end{center}
		\begin{enumerate}[a)]
			\item $X$ là số cây con có hạt màu vàng trong số $4$ cây con. $X$ là một biến ngẫu nhiên có phân bố nhị thức với $n=4$, $p=\dfrac{3}{4}$. Giá trị của $X$ thuộc tập $\{0; 1; 2; 3; 4\}$. Ta có
			\begin{itemize}
				\item $\mathrm{P}(X=0)=\left(\dfrac{1}{4}\right)^4=\dfrac{1}{256}$;
				\item $\mathrm{P}(X=1)=\mathrm{C}^1_4\cdot \dfrac{3}{4}\cdot \left(\dfrac{1}{4}\right)^3=\dfrac{12}{256}$;
				\item $\mathrm{P}(X=2)=\mathrm{C}^2_4\cdot \left( \dfrac{3}{4}\right) ^2\cdot \left(\dfrac{1}{4}\right)^2=\dfrac{54}{256}$;
				\item $\mathrm{P}(X=3)=\mathrm{C}^3_4\cdot \left( \dfrac{3}{4}\right)^3\cdot \left(\dfrac{1}{4}\right)^1=\dfrac{108}{256}$;
				\item $\mathrm{P}(X=4)=\mathrm{C}^4_4\cdot \left(\dfrac{3}{4}\right)^4=\dfrac{81}{256}$.
			\end{itemize}
			Bảng phân bố xác suất của $X$ như sau.
			\begin{center}
				\begin{tabular}{|c|c|c|c|c|c|}
					\hline
					$X$ & $0$ & $1$ & $2$ & $3$ & $4$ \\
					\hline
					$\mathrm{P}$ & $\dfrac{1}{256}$ & $\dfrac{12}{256}$ & $\dfrac{54}{256}$ & $\dfrac{108}{256}$ & $\dfrac{81}{256}$ \\
					\hline
				\end{tabular}
			\end{center}
			\item Gọi $Y$ là số cây con có hạt màu xanh. Khi đó $Y$ là biến ngẫu nhiên có phân bố nhị thức với tham số $n=4$, $p=\dfrac{1}{4}$.\\
			Trung bình có $\mathrm{E}(Y)=4\cdot\dfrac{1}{4}=1$ cây có hạt màu xanh.
		\end{enumerate}
	}
\end{bt}

\begin{bt}%[2D1V1-5]
% nguon: Archive/2024/0. [LVH] BG - 12 New - 4in1-20240625T224301Z-001/0. [LVH] BG - 12 New - 4in1/2. [LVH] TEX SGK 12 NEW-20240418T074833Z-001/2. [LVH] TEX SGK 12 NEW/[LVH] SGK-12-KNTT/data/SCD/12-SCD-KNTT-1-2.tex #1174

	Trong một lớp học có $6$ bóng đèn hoạt động độc lập với nhau. Mỗi bóng có xác suất bị hỏng là $0{,}25$. Gọi $X$ là số bóng sáng.
	\begin{enumerate}[a)]
		\item Gọi tên phân bố xác suất biến ngẫu nhiên $X$.
		\item Biết rằng lớp học có đủ ánh sáng nếu có ít nhất $4$ bóng sáng. Tính xác suất để lớp học đủ ánh sáng.
		\item Tính kì vọng, phương sai và độ lệch chuẩn của $X$.
	\end{enumerate}
	
	\loigiai{
		\begin{enumerate}[a)]
			\item $X$ là số bóng sáng. $X$ là một biến ngẫu nhiên có phân bố nhị thức với tham số $n=6$, $p=\dfrac{3}{4}$.
			\item Lớp học đủ ánh sáng nếu có ít nhất $4$ bóng sáng tức là $X\ge 4$. Khi đó
			$$\mathrm{P}(X\ge 4)=\mathrm{C}^4_6\cdot \left(\dfrac{3}{4}\right)^4\cdot \left(\dfrac{1}{4}\right)^2+\mathrm{C}^5_6\cdot \left(\dfrac{3}{4}\right)^5\cdot \dfrac{1}{4}+\mathrm{C}^6_6\cdot \left(\dfrac{3}{4}\right)^6=\dfrac{1701}{2048}.$$
			\item Kì vọng của $X$ là $\mathrm{E}(X)=6\cdot\dfrac{3}{4}=4{,}5$.\\
			Phương sai của $X$ là $\mathrm{V}(X)=6\cdot\dfrac{3}{4}\cdot\dfrac{1}{4}=1{,}125$.\\
			Độ lệch chuẩn của $X$ là $\sigma(X)=\sqrt{1{,}125}\approx 1{,}061$.
		\end{enumerate}
	}
\end{bt}

\begin{bt}%[2D1V1-5]
% nguon: Archive/2024/0. [LVH] BG - 12 New - 4in1-20240625T224301Z-001/0. [LVH] BG - 12 New - 4in1/2. [LVH] TEX SGK 12 NEW-20240418T074833Z-001/2. [LVH] TEX SGK 12 NEW/[LVH] SGK-12-KNTT/data/SCD/12-SCD-KNTT-1-2.tex #1175

	Sơn và Tùng thi đấu bóng bàn với nhau. Trận đấu gồm $5$ ván độc lập. Xác suất thắng của Sơn trong mỗi ván là $\dfrac{1}{4}$. Biết rằng mỗi ván không có kết quả hoà. Người thắng trận đấu nếu thắng ít nhất $3$ ván đấu.
	\begin{enumerate}[a)]
		\item Gọi $X$ là số trận thắng của Sơn. Hỏi $X$ là biến ngẫu nhiên có phân bố xác suất gì?
		\item Tính xác suất để Sơn thắng Tùng trong trận đấu.
	\end{enumerate}	
	\loigiai{
		\begin{enumerate}[a)]
			\item Gọi $X$ là số trận thắng của Sơn. $X$ là biến ngẫu nhiên có phân bố xác suất nhị thức với tham số $n=5$ và $p=\dfrac{1}{4}$.
			\item Sơn thắng Tùng trong trận đấu tức là $X\ge 3$. Khi đó ta có:
			$$\mathrm{P}(X\ge 3)=\mathrm{C}^3_5\cdot \left(\dfrac{1}{4}\right)^3\cdot \left(\dfrac{3}{4}\right)^2+\mathrm{C}^4_5\cdot \left(\dfrac{1}{4}\right)^4\cdot \dfrac{3}{4}+\mathrm{C}^5_5\cdot \left(\dfrac{1}{4}\right)^5=\dfrac{53}{512}.$$
		\end{enumerate} 
	}
\end{bt}

\begin{bt}%[2D1V1-5]
% nguon: Archive/2024/0. [LVH] BG - 12 New - 4in1-20240625T224301Z-001/0. [LVH] BG - 12 New - 4in1/2. [LVH] TEX SGK 12 NEW-20240418T074833Z-001/2. [LVH] TEX SGK 12 NEW/[LVH] SGK-12-KNTT/data/SCD/12-SCD-KNTT-1-2.tex #1176

	Cam xuất khẩu được đóng thành từng thùng. Xác suất để một quả cam không đạt chất lượng là $0{,}03$. Vì số lượng cam trong mỗi thùng rất lớn nên không thể kiểm tra toàn bộ số cam trong thùng, người ta lấy ngẫu nhiên từ thùng cam $20$ lần một cách độc lập, mỗi lần lấy $1$ quả để kiểm tra rồi trả lại nó vào thùng. Gọi $X$ là số quả cam không đạt chất lượng.
	\begin{enumerate}[a)]
		\item Gọi tên phân bố xác suất biến ngẫu nhiên $X$.
		\item Các thùng cam được phân thành ba loại theo cách sau:\\
		Trong 20 lần lấy đó:
		\begin{itemize}
			\item Nếu tất cả các quả cam lấy ra đều đạt chất lượng thì thùng được xếp loại $I$;
			\item Nếu có $1$ hoặc $2$ quả cam không đạt chất lượng thì thùng được xếp loại $II$;
			\item Nếu có ít nhất $3$ quả cam không đạt chất lượng thì thùng được xếp loại $III$.
		\end{itemize}
		Tính tỉ lệ các thùng cam được xếp loại $I$, $II$, $III$.
	\end{enumerate}
	\loigiai{
		\begin{enumerate}[a)]
			\item $X$ là số quả cam không đạt chất lượng. $X$ là biến ngẫu nhiên có phân bố xác suất nhị thức với tham số $n=20$, $p=0{,}03$.
			\item 
			\begin{itemize}
				\item Thùng cam được xếp loại $I$ nếu $X=0$. Khi đó xác suất là $$\mathrm{P}_1=\mathrm{P}(X=0)=0{,}97^{20}\approx 0{,}5438.$$
				\item Thùng cam được xếp loại $II$ nếu $X=1$ hoặc $X=2$. Khi đó xác suất là $$\mathrm{P}_2=\mathrm{P}(X=1)+\mathrm{P}(X=2)=\mathrm{C}^1_{20}\cdot 0{,}03\cdot 0{,}97^{19}+\mathrm{C}^2_{20}\cdot 0{,}03^2\cdot 0{,}97^{18} \approx 0{,}4352.$$
				\item Thùng cam được xếp loại $III$ nếu $X\ge 3$. Khi đó xác suất là $$\mathrm{P}_3=\mathrm{P}(X\ge 3)=1-\mathrm{P}_1-\mathrm{P}_2=1-0{,}5438-0{,}4352=0{,}021.$$
			\end{itemize}
			Vậy tỉ lệ các thùng cam được xếp loại $I$, $II$, $III$ lần lượt là $54{,}38\%$, $43{,}52\%$ và $2{,}1\%$.
		\end{enumerate}
		
	}
\end{bt}

\begin{bt}%[2D1H5-1]
% nguon: Archive/2024/TaiLieuGiangDay/1.HamSo/1. K12 - 4in1 - Hamso/ThubaiGoc/Bai5-Dang1.tex #1999
%[Dự án Bài Giảng 12 Nhóm Toán & LaTex, Lê Minh Thiện Anh]
	Khảo sát sự biến thiên và vẽ đồ thị của hàm số $y=-x^3+3x^2-4$.
	\loigiai{
		\begin{enumerate}
			\item Tập xác định $\mathscr{D}=\mathbb{R}$.
			\item Sự biến thiên
			\begin{itemize}
				\item Ta có $y'=-3x^2+6x$. Vậy $y'=0$ khi $x=0$ hoặc $x=2$.
				\item Trên khoảng $(0;2)$, $y'>0$ nên hàm số đồng biến. Trên các khoảng $(-\infty;0)$ và $(2;+\infty)$, $y'<0$ nên hàm số nghịch biến trên mỗi khoảng đó.
				\item Hàm số đạt cực tiểu tại $x=0$, giá trị cực tiểu $y_{CT}=-4$. Hàm số đạt cực đại tại $x=2$, giá trị cực đại $y_{\text{CĐ}}=0$.
				\item Giới hạn vô cực\\ 
				$\lim\limits_{x \to -\infty}y=\lim\limits_{x \to -\infty}x^3\left(-1+\dfrac{3}{x}-\dfrac{4}{x^3}\right)=+\infty$;\\ 
				$\lim\limits_{x \to +\infty}y=\lim\limits_{x \to +\infty}x^3\left(-1+\dfrac{3}{x}-\dfrac{4}{x^3}\right)=-\infty$.
				\item Bảng biến thiên
				\begin{center}
					\begin{tikzpicture}
						\tikzset{double style/.append style={double distance=2pt}}
						\tkzTabInit[nocadre=false,lgt=1.2,espcl=2.5,deltacl=0.6]
						{$x$/0.6,$y'$/0.6,$y$/2}
						{$-\infty$,$0$,$2$,$+\infty$}
						\tkzTabLine{,-,0,+,0,-,}
						\tkzTabVar{+/$+\infty$ ,-/$-4$,+/ $0$, -/ $-\infty$}
					\end{tikzpicture}
				\end{center}
			\end{itemize}
			\item Đồ thị
			\immini{
				\begin{itemize}
					\item Giao điểm của đồ thị với trục tung là điểm $(0;-4)$.
					\item Ta có $y=0\Leftrightarrow-x^3+3x^2-4=0\Leftrightarrow-(x-2)^2(x+1)=0$ $\Leftrightarrow x=-1$ hoặc $x=2$. Do đó giao điểm của đồ thị hàm số với trục hoành là các điểm $(-1;0)$ và $(2;0)$.
					\item Đồ thị của hàm số có tâm đối xứng là điểm $(1;-2)$.
				\end{itemize}
			}{\begin{tikzpicture}[>=stealth, scale=0.8]
					%\draw[step=1,gray,very thin] (-3,-3) grid (3,3);
					\draw[->] (-2,0) --(0,0) node[below right]{$O$} -- (4,0) node[below]{$x$};
					\draw[->](0,-5)--(0,2) node[right]{$y$};
					\draw[smooth,blue, thick] plot[domain=-1.1:3.1] (\x,{(-1)*(\x)^(3)+3*(\x)^(2)-4});
					\draw (-1,0) node [below left]{$-1$} (0,-4) node [below left]{$-4$} (2,0) node [above]{$2$};
					\draw [dashed] (0,-2) node [left]{\tiny$-2$}--(1,-2)--(1,0) node [above]{$1$};
					\draw [fill=black] (1,-2) circle (1.5pt);
			\end{tikzpicture}}
		\end{enumerate}
	}
\end{bt}

\begin{bt}%[2D1H5-1]
% nguon: Archive/2024/TaiLieuGiangDay/1.HamSo/1. K12 - 4in1 - Hamso/ThubaiGoc/Bai5-Dang1.tex #2000
%[Dự án Bài Giảng 12 Nhóm Toán & LaTex, Lê Minh Thiện Anh]
	Khảo sát sự biến thiên và vẽ đồ thị hàm số $y=x^3-2x^2+2x-1$.
	\loigiai{
		\begin{enumerate}
			\item Tập xác định $\mathscr{D}=\mathbb{R}$.
			\item Sự biến thiên
			\begin{itemize}
				\item Ta có $y'=3x^2-4x+2$. Vậy $y'>0, \forall x\in\mathbb{R}$.
				\item Hàm số đồng biến trên khoảng $(-\infty; +\infty)$.
				\item Hàm số không có cực trị.
				\item Giới hạn vô cực:\\
				$\lim\limits_{x \to -\infty}y=\lim\limits_{x \to -\infty}x^3\left(1-\dfrac{2}{x}+\dfrac{2}{x^2}-\dfrac{1}{x^3}\right)=-\infty$,\\
				$\lim\limits_{x \to +\infty}y=\lim\limits_{x \to +\infty}x^3\left(1-\dfrac{2}{x}+\dfrac{2}{x^2}-\dfrac{1}{x^3}\right)=+\infty$.
				\item Bảng biến thiên
				\begin{center}
					\begin{tikzpicture}
						\tikzset{double style/.append style={double distance=2pt}}
						\tkzTabInit[nocadre=false,lgt=1.2,espcl=4.5,deltacl=0.6]
						{$x$/0.6,$y'$/0.6,$y$/2}
						{$-\infty$,$+\infty$}
						\tkzTabLine{,+,}
						\tkzTabVar{-/$-\infty$ , +/ $+\infty$}
					\end{tikzpicture}
				\end{center}
			\end{itemize}
			\item Đồ thị
			\immini{
				\begin{itemize}
					\item Giao điểm của đồ thị hàm số với trục tung là điểm $(0;-1)$.
					\item Ta có $y=0\Leftrightarrow x^3-2x^2+2x-1=0\Leftrightarrow x=1$.\\
					Do đó giao điểm của đồ thị hàm số với trục hoành là điểm $(1;0)$.
					\item Đồ thị hàm số có tâm đối xứng là điểm $\left(\dfrac{2}{3};\dfrac{-7}{27}\right)$.
				\end{itemize}
			}{\begin{tikzpicture}[>=stealth, scale=1.2]
					%\draw[step=1,gray,very thin] (-3,-3) grid (3,3);
					\draw[->] (-2,0) --(0,0) node[above left]{$O$} -- (3,0) node[below]{$x$};
					\draw[->](0,-2)--(0,3) node[right]{$y$};
					\draw[smooth,blue, thick] plot[domain=-0.4:2] (\x,{(\x)^(3)-2*(\x)^(2)+2*(\x)-1});
					\draw (0,-1) node [right]{$-1$} (1,0) node [below right]{$1$};
					\draw [dashed] (0,-7/27) node [left]{$-\frac{7}{27}$}--(2/3,-7/27)--(2/3,0) node [above]{$\frac{2}{3}$};
					\draw [fill=black] (2/3,-7/27) circle (1.5pt);
			\end{tikzpicture}}
		\end{enumerate}
	}
\end{bt}

\begin{bt}%[2D1H5-1]
% nguon: Archive/2024/TaiLieuGiangDay/1.HamSo/1. K12 - 4in1 - Hamso/ThubaiGoc/Bai5-Dang1.tex #2001
%[Dự án Bài Giảng 12 Nhóm Toán & LaTex, Lê Minh Thiện Anh]
	Khảo sát sự biến thiên và vẽ đồ thị hàm số $y=-2x^3+3x^2-5x$.
	\loigiai{
		\begin{enumerate}
			\item Tập xác định $\mathscr{D}=\mathbb{R}$.
			\item Sự biến thiên
			\begin{itemize}
				\item Ta có $y'=-6x^2+6x-5$. Vậy $y'<0, \forall x\in\mathbb{R}$.
				\item Hàm số không có cực trị.
				\item Giới hạn vô cực\\
				$\lim\limits_{x \to -\infty}y=\lim\limits_{x \to -\infty}(-2x^3+3x^2-5x)=\lim\limits_{x \to -\infty}x^3\left(-2+\dfrac{3}{x}-\dfrac{5}{x^2}\right)=+\infty$;\\
				$\lim\limits_{x \to +\infty}y=\lim\limits_{x \to +\infty}(-2x^3+3x^2-5x)=\lim\limits_{x \to +\infty}x^3\left(-2+\dfrac{3}{x}-\dfrac{5}{x^2}\right)=-\infty$.
				\item Bảng biến thiên
				\begin{center}
					\begin{tikzpicture}
						\tikzset{double style/.append style={double distance=2pt}}
						\tkzTabInit[nocadre=false,lgt=1.2,espcl=4.5,deltacl=0.6]
						{$x$/0.6,$y'$/0.6,$y$/2}
						{$-\infty$,$+\infty$}
						\tkzTabLine{,-,}
						\tkzTabVar{+/$+\infty$ , -/ $-\infty$}
					\end{tikzpicture}
				\end{center}
			\end{itemize}
			\item Đồ thị
			\immini{
				\begin{itemize}
					\item Giao điểm của đồ thị hàm số với trục tung là điểm $(0,0)$.
					\item Đồ thị của hàm số có tâm đối xứng là điểm $\left(\dfrac{1}{2};-2\right)$.
				\end{itemize}
			}{
				\begin{tikzpicture}[>=stealth, scale=0.8]
					%\draw[step=1,gray,very thin] (-3,-3) grid (3,3);
					\draw[->] (-2,0) --(0,0) node[above left]{$O$} -- (3,0) node[below]{$x$};
					\draw[->](0,-4)--(0,3) node[right]{$y$};
					\draw[smooth,blue, thick] plot[domain=-0.4:1] (\x,{(-2)*(\x)^(3)+3*(\x)^(2)-5*(\x)});
					\draw[dashed] (0,-2)node [above left]{$-2$}--(0.5,-2)--(0.5,0) node [above]{$\frac{1}{2}$};
					\draw [fill=black] (0.5,-2) circle (1.5pt);
				\end{tikzpicture}
			}
		\end{enumerate}
	}
\end{bt}

\begin{bt}%[2D1H5-1]
% nguon: Archive/2024/TaiLieuGiangDay/1.HamSo/1. K12 - 4in1 - Hamso/ThubaiGoc/Bai5-Dang1.tex #2005
%[Dự án Bài Giảng 12 Nhóm Toán & LaTex, Lê Minh Thiện Anh]
	Khảo sát sự biến thiên và vẽ đồ thị các hàm số sau:
	\begin{listEX}[2]
		\item $y=x^3+3x^2-4$;
		\item $y=x^3+4x^2+4x$;
		\item $y=-2x^3+2$;
		\item $y=-x^3-x^2-x+1$.
	\end{listEX}
	\loigiai{
	\begin{enumerate}
		\item $y=x^3+3x^2-4$
		\begin{itemize}
			\item Tập xác định $\mathscr{D}=\mathbb{R}$.
			\item Sự biến thiên
			\begin{itemize}
				\item Giới hạn tại vô cực:
				$\lim\limits_{x\to -\infty}y=-\infty$; 
				$\lim\limits_{x\to +\infty}y=+\infty$.
				\item Ta có $y'=3x^2+6x$; $y'=0\Leftrightarrow3x^2+6x=0\Leftrightarrow x=0$ hoặc $x=-2$.
				\item Bảng biến thiên
				\begin{center}
					\begin{tikzpicture}
						\tkzTabInit[lgt=1.2,espcl=2.5]{$x$/.6,$y'$/.6,$y$/2}{$-\infty$,$-2$,$0$,$+\infty$}
						\tkzTabLine{,+,0,-,0,+,}
						\tkzTabVar{-/$-\infty$ , +/$0$,-/$-4$, +/$+\infty$}%
					\end{tikzpicture}
				\end{center}
				\item Chiều biến thiên: Hàm số đồng biến trên các khoảng $(-\infty;-2)$ và $(0;+\infty)$, nghịch biến trên khoảng $(-2;0)$.
				\item Cực trị:	Hàm số đạt cực đại tại $x=-2$, $y_{\text{CĐ}}=0$;\\
				Hàm số đạt cực tiểu tại $x=0$, $y_{\text{CT}}=-4$.
			\end{itemize}
			\item Đồ thị:\\
			\immini{
			Đồ thị hàm số giao với trục $Oy$ tại điểm $(0;-4)$.\\
			Ta có $x^3+3x^2-4=0\Leftrightarrow x=1 $ hoặc $x=-2$, do đó đồ thị hàm số giao với trục $Ox$ tại hai điểm là $(1;0)$ và $(-2;0)$.\\
			Đồ thị hàm số đi qua điểm $(-3;-4)$ và có tâm đối xứng là điểm $I(-1;-2)$.}{\begin{tikzpicture}[scale=0.8, font=\footnotesize,line join=round, line cap=round, >=stealth]
				\draw[->](-4,0)--(2,0);
				\foreach \x in {-3,-2}
				\draw[shift={(\x,0)},color=black] (0pt,2pt)--(0pt,-2pt) node[above]{$\x$};
				\draw[->,color=black](0,-5)--(0,2.5);
				\foreach \y in {}
				\draw[shift={(0,\y)},color=black](2pt,0pt)--(-2pt,0pt) node[left]{$\y$};
				\draw[color=black](2,0)node[above]{$x$};
				\draw[color=black](0,2.5)node[right]{$y$};
				\draw[color=black](0pt,-8pt)node[right]{$O$};
				\clip(-4,-5) rectangle (3.5,2.4);
				\draw[smooth,samples=100,domain=-4:3.4]plot(\x,{(\x)^3+3*(\x)^2-4});
				\draw[dashed](-3,0)--(-3,-4)--(0,-4);
				\draw[dashed](-1,0)--(-1,-2)--(0,-2);
				\node[below left] at (-1,-2){$I$};
				\node[below right] at (0,-4){$-4$};
				\node[above] at (-1,0) {$-1$};
				\node[below left] at (0,-2) {$-2$};
				\node[above right] at (1,0) {$1$};
				\end{tikzpicture}}
		\end{itemize}
		\item $y=x^3+4x^2+4x$
		\begin{itemize}
			\item Tập xác định $\mathscr{D}=\mathbb{R}$.
			\item Sự biến thiên
			\begin{itemize}
				\item Giới hạn tại vô cực:
				$\lim\limits_{x\to -\infty}y=-\infty$; 
				$\lim\limits_{x\to +\infty}y=+\infty$.
				\item Ta có $y'=3x^2+8x+4$; $y'=0\Leftrightarrow3x^2+8x+4=0\Leftrightarrow x=-\dfrac{2}{3}$ hoặc $x=-2$.
				\item Bảng biến thiên
				\begin{center}
					\begin{tikzpicture}
						\tkzTabInit[lgt=1.2,espcl=2.5]{$x$/1,$y'$/.6,$y$/2}{$-\infty$,$-2$,$-\dfrac{2}{3}$,$+\infty$}
						\tkzTabLine{,+,0,-,0,+,}
						\tkzTabVar{-/$-\infty$ , +/$0$,-/$-\dfrac{32}{27}$, +/$+\infty$}%
					\end{tikzpicture}
				\end{center}
				\item Chiều biến thiên: Hàm số đồng biến trên các khoảng $(-\infty;-2)$ và $\left(-\dfrac{2}{3};+\infty\right)$, nghịch biến trên khoảng $\left(-2;-\dfrac{2}{3}\right)$.
				\item Cực trị:	Hàm số đạt cực đại tại $x=-2$, $y_{\text{CĐ}}=0$;\\
				Hàm số đạt cực tiểu tại $x=-\dfrac{2}{3}$, $y_{\text{CT}}=-\dfrac{32}{27}$.
			\end{itemize}
			\item Đồ thị:\\
			\immini{
				Đồ thị hàm số giao với trục $Oy$ tại điểm $(0;0)$.\\
				Ta có $x^3+4x^2+4x=0\Leftrightarrow x=0 $ hoặc $x=-2$, do đó đồ thị hàm số giao với trục $Ox$ tại hai điểm là $(0;0)$ và $(-2;0)$.\\
				Đồ thị hàm số đi qua điểm $(-3;-3)$ và có tâm đối xứng là điểm $I\left (-\dfrac{4}{3};-\dfrac{16}{27}\right )$.}{\begin{tikzpicture}[scale=0.8, font=\footnotesize,line join=round, line cap=round, >=stealth]
					\draw[->](-4,0)--(2,0);
					\foreach \x in {-3,-2}
					\draw[shift={(\x,0)},color=black] (0pt,2pt)--(0pt,-2pt) node[below]{$\x$};
					\draw[->,color=black](0,-5)--(0,2.5);
					\foreach \y in {}
					\draw[shift={(0,\y)},color=black](2pt,0pt)--(-2pt,0pt) node[left]{$\y$};
					\draw[color=black](2,0)node[above]{$x$};
					\draw[color=black](0,2.5)node[right]{$y$};
					\draw[color=black](0pt,-8pt)node[right]{$O$};
					\clip(-4,-5) rectangle (3.5,2.4);
					\draw[smooth,samples=100,domain=-4:3.4]plot(\x,{(\x)^3+4*(\x)^2+4*(\x)});
					\draw[dashed](-3,0)--(-3,-3)--(0,-3);
					\draw[dashed](-0.666,0)--(-0.666,-1.185)--(0,-1.185);
					\draw[dashed] (-1.333,0)--(-1.333,-0.592)--(0,-0.592) ;
					\node[below left] at (-1.333,-0.592){$I$};
					\node[below right] at (0,-3){$-3$};
					\node[above] at (-0.666,0) {$-\dfrac{2}{3}$};
					\node[right] at (0,-1.185) {$-\dfrac{32}{27}$};
			\end{tikzpicture}}
		\end{itemize}
		\item $y=-2x^3+2$
		\begin{itemize}
			\item Tập xác định $\mathscr{D}=\mathbb{R}$.
			\item Sự biến thiên
			\begin{itemize}
				\item Giới hạn tại vô cực:
				$\lim\limits_{x\to -\infty}y=+\infty$; 
				$\lim\limits_{x\to +\infty}y=-\infty$.
				\item Ta có $y'=-6x^2$; $y'=0\Leftrightarrow -6x^2=0\Leftrightarrow x=0$.
				\item Bảng biến thiên
				\begin{center}
					\begin{tikzpicture}[font=\normalsize,t style/.style={style=solid}]
						%dòng khai báo
						\tkzTabInit[lgt=1.2,espcl=5,deltacl=0.5]
						{$x$ /0.6, $y'$/0.6, $y$/2}
						{$ -\infty $,$ +\infty $}
						%dòng xét dấu
						\tkzTabLine{ , - , } % z, t, d;
						%dòng biến thiên
						\path ($(N12)!0.1!(N13)$) node (A1){$ +\infty $}
						($(N22)!0.9!(N23)$) node (A2){$ -\infty $};
						\foreach \x/\y in {A1/A2}{
							\draw[-stealth] (\x)--(\y);
						}
					\end{tikzpicture}
				\end{center}
				\item Chiều biến thiên: Hàm số nghịch biến trên $\mathbb{R}$.
				\item Cực trị:	Hàm số không có cực trị.
			\end{itemize}
			\item Đồ thị:\\
			\immini{
				Đồ thị hàm số giao với trục $Oy$ tại điểm $(0;2)$.\\
				Ta có $-2x^3+2=0\Leftrightarrow x=1 $, do đó đồ thị hàm số giao với trục $Ox$ tại điểm $(1;0)$\\
				Đồ thị hàm số đi qua điểm $(-1;4)$ và có tâm đối xứng là điểm $I(0;2)$.}{\begin{tikzpicture}[scale=0.8, font=\footnotesize,line join=round, line cap=round, >=stealth]
					\draw[->](-2,0)--(2,0);
					\foreach \x in {}
					\draw[shift={(\x,0)},color=black] (0pt,2pt)--(0pt,-2pt) node[below]{$\x$};
					\draw[->,color=black](0,-2)--(0,4.5);
					\foreach \y in {}
					\draw[shift={(0,\y)},color=black](2pt,0pt)--(-2pt,0pt) node[left]{$\y$};
					\draw[color=black](2,0)node[above]{$x$};
					\draw[color=black](0,4.5)node[right]{$y$};
					\draw[color=black](0pt,-8pt)node[left]{$O$};
					\clip(-2,-2) rectangle (2,4.5);
					\draw[smooth,samples=100,domain=-2:2]plot(\x,{-2*(\x)^3+2});
					\draw[dashed] (-1,0)--(-1,4)--(0,4) ;
					\node[above right] at (0,2){$I$};
					\node[below left] at (0,2) {$2$};
					\node[below left] at (1,0) {$1$};
					\node[below] at (-1,0) {$-1$};
					\node[right] at (0,4) {$4$};
			\end{tikzpicture}}
		\end{itemize}
		\item $y=-x^3-x^2-x+1$
		\begin{itemize}
			\item Tập xác định $\mathscr{D}=\mathbb{R}$.
			\item Sự biến thiên
			\begin{itemize}
				\item Giới hạn tại vô cực:
				$\lim\limits_{x\to -\infty}y=+\infty$; 
				$\lim\limits_{x\to +\infty}y=-\infty$.
				\item Ta có $y'=-3x^2-2x-1$; $y'=0\Leftrightarrow -3x^2-2x-1=0$ vô nghiệm.
				\item Bảng biến thiên
				\begin{center}
					\begin{tikzpicture}[font=\normalsize,t style/.style={style=solid}]
						%dòng khai báo
						\tkzTabInit[lgt=1.2,espcl=5,deltacl=0.5]
						{$x$ /0.6, $y'$/0.6, $y$/2}
						{$ -\infty $,$ +\infty $}
						%dòng xét dấu
						\tkzTabLine{ , - , } % z, t, d;
						%dòng biến thiên
						\path ($(N12)!0.1!(N13)$) node (A1){$ +\infty $}
						($(N22)!0.9!(N23)$) node (A2){$ -\infty $};
						\foreach \x/\y in {A1/A2}{
							\draw[-stealth] (\x)--(\y);
						}
					\end{tikzpicture}
				\end{center}
				\item Chiều biến thiên: Hàm số nghịch biến trên $\mathbb{R}$.
				\item Cực trị:	Hàm số không có cực trị.
			\end{itemize}
			\item Đồ thị:\\
			\immini{
				Đồ thị hàm số giao với trục $Oy$ tại điểm $(0;1)$.\\
				Ta có $-x^3-x^2-x+1=0\Leftrightarrow x=1 $, do đó đồ thị hàm số giao với trục $Ox$ tại điểm $(1;0)$.\\
				Đồ thị hàm số đi qua điểm $(-1;2)$ và có tâm đối xứng là điểm $I\left (-\dfrac{1}{3};\dfrac{34}{27}\right )$.}{\begin{tikzpicture}[scale=0.8, font=\footnotesize,line join=round, line cap=round, >=stealth]
					\draw[->](-3,0)--(2,0);
					\foreach \x in {-1}
					\draw[shift={(\x,0)},color=black] (0pt,2pt)--(0pt,-2pt) node[below]{$\x$};
					\draw[->,color=black](0,-3)--(0,5);
					\foreach \y in {}
					\draw[shift={(0,\y)},color=black](2pt,0pt)--(-2pt,0pt) node[left]{$\y$};
					\draw[color=black](2,0)node[above]{$x$};
					\draw[color=black](0,5)node[right]{$y$};
					\draw[color=black](0pt,-8pt)node[right]{$O$};
					\clip(-3,-3) rectangle (2,5);
					\draw[smooth,samples=100,domain=-4:3.4]plot(\x,{-(\x)^3-(\x)^2-(\x)+1});
					\draw[dashed](-1,0)--(-1,2)--(0,2);
					\draw[dashed](-0.333,0)--(-0.333,1.259)--(0,1.259);
					\node[below left] at (-0.333,1.259){$I$};
					\node[right] at (0,1.259){$\frac{34}{27}$};
					\node[below] at (-0.333,0) {$-\frac{1}{3}$};
					\node[right] at (0,2) {$2$};
					\node[below right] at (1,0) {$1$};
			\end{tikzpicture}}
		\end{itemize}
	\end{enumerate}
	}
\end{bt}

\begin{bt}%[2D1H5-1]
% nguon: Archive/2024/TaiLieuGiangDay/1.HamSo/1. K12 - 4in1 - Hamso/ThubaiGoc/Bai5-Dang1.tex #2006
%[Dự án Bài Giảng 12 Nhóm Toán & LaTex, Lê Minh Thiện Anh]
	Khảo sát sự biến thiên và vẽ đồ thị của hàm số $y=-x^3-3x^2+2$.
	\loigiai{
		\begin{itemize}
			\item Tập xác định: $\mathscr{D}=\mathbb{R}$.
			\item Sự biến thiên\\
			$y'=-3x^2-6x$, 
			$y'=0\Leftrightarrow \hoac{&x=-2\\&x=0.}$\\
			$y'<0\Leftrightarrow x\in (-\infty;-2)\cup (0;+\infty)$;\\
			$y'>0\Leftrightarrow x\in (-2;0)$.\\
			Suy ra hàm số nghịch biến trên các khoảng $(-\infty;-2)$ và $(0;+\infty)$; đồng biến trên $(-2;0)$.\\
			Hàm số đạt cực đại tại $x=0$; $y_{\text{CĐ}}=y(0)=2$.\\
			Hàm số đạt cực tiểu tại $x=-2$; $y_{\text{CT}}=y(-2)=-2$.\\
			$\lim\limits_{x\to -\infty}y=+\infty$.\\
			$\lim\limits_{x\to +\infty}y=-\infty$.\\
			Bảng biến thiên
			\begin{center}
				\begin{tikzpicture}
				\tkzTabInit[lgt=1.2,espcl=2.5]{$x$/.6,$y'$/.6,$y$/2}{$-\infty$,$-2$,$0$,$+\infty$}
				\tkzTabLine{,-,0,+,0,-,}
				\tkzTabVar{+/$+\infty$ , -/$-2$,+/$2$, -/$-\infty$}%
				\end{tikzpicture}
			\end{center}
			\item Đồ thị\\
			\immini{
				Đồ thị cắt trục $Oy$ tại $(0;2)$.\\
				Đồ thị đi qua các điểm $(-2;-2)$, $(1;-2)$, $(-3;2)$.\\
				Đồ thị nhận điểm $I(-1;0)$ làm tâm đối xứng.
			}
			{
				\begin{tikzpicture}[>=stealth,line cap=round,line join=round,x=1cm,y=1cm]
				\draw[->](-4,0)--(3,0);
				\foreach \x in {-3,2}
				\draw[shift={(\x,0)},color=black] (0pt,2pt)--(0pt,-2pt) node[below]{$\x$};
				\draw[->,color=black](0,-3)--(0,3);
				\foreach \y in {-1,1}
				\draw[shift={(0,\y)},color=black](2pt,0pt)--(-2pt,0pt) node[left]{\normalsize $\y$};
				\draw[color=black](3,.2)node[right]{$x$};
				\draw[color=black](.2,3)node[right]{$y$};
				\draw[color=black](0pt,-8pt)node[right]{\normalsize $O$};
				\clip(-4,-3) rectangle (3,3);
				%Vẽ đồ thị
				\draw[smooth,samples=100,domain=-4:3]plot(\x,{-(\x)^3-3*(\x)^2+2});
				%Vẽ râu ria
				\draw[dashed](-2,0)--(-2,-2)--(0,-2);
				\draw[dashed](-3,0)--(-3,2)--(0,2);
				\draw[dashed](1,0)--(1,-2)--(0,-2);
				\node[above left] at (-1,0){$I$};
				\node[above right] at (0,2){$2$};
				\node[below right] at (0,-2){$-2$};
				\node[above] at (-2,0){$-2$};
				\node[below right] at (-1,0) {$-1$};
				\node[above] at (1,0) {$1$};
				\end{tikzpicture}
			}
		\end{itemize}
	}
\end{bt}

\begin{bt}%[2D1H5-1]
% nguon: Archive/2024/TaiLieuGiangDay/1.HamSo/1. K12 - 4in1 - Hamso/ThubaiGoc/Bai5-Dang1.tex #2007
%[Dự án Bài Giảng 12 Nhóm Toán & LaTex, Lê Minh Thiện Anh]
	\immini{Cho hàm số $y=ax^3+bx^2+cx+d$ có đồ thị như hình vẽ bên. Tính $S=a+b$.}
	{	\begin{tikzpicture}[scale=0.7, font=\footnotesize,line join=round, line cap=round,>=stealth]
	\draw[->] (-1.5,0)--(3.5,0) node[below] {$x$};
	\draw[->] (0,-2.5)--(0,2.5) node[right] {$y$};
	\draw (0,0) node [below left] {$O$};
	\foreach \x in {-1,1,2,3}\draw (\x,-0.05)--(\x,0.05) node [above] {\footnotesize $\x$};
	\foreach \y in {-2,-1,1,2}\draw (0.05,\y)--(-0.05,\y) node [left] {\footnotesize $\y$};
	\clip (-1.5,-2.5) rectangle (3.5,2.5);
	\draw[smooth,samples=200,domain=-1.5:3.5] plot (\x,{(\x)^3+-3*((\x)^2)+2});
	\fill (0,0) circle (1pt);
	\fill (0,2) circle (1pt);
	\draw[dashed] (2,0)--(2,-2)--(0,-2);\fill (2,-2) circle (1pt);
	\end{tikzpicture}}
	\loigiai{
	Ta có $y'=3ax^2+2bx+c$. Từ đồ thị hàm số ta suy ra $y'$ có hai nghiệm $x_1=0$ và $x_2=2$ nên 
	$$\heva{&c=0\\&12a+4b+c=0}\Leftrightarrow \heva{&c=0\\&b=-3a.}$$
	Ta cũng có đồ thị đi qua điểm $(0;2)$ nên suy ra $d=2$.\\
	Do đó, hàm số có dạng $y=ax^3-3ax^2+2$.\\
	Mặt khác, đồ thị hàm số có điểm cực tiểu $(2;-2)$ nên $a\cdot2^3-3a\cdot 2^2+2=-2\Leftrightarrow a=1\Rightarrow b=-3$.\\
	Vậy $S=a+b=1+\left(-3\right)=-2$.
	}
\end{bt}

\begin{bt}%[2D1H5-1]
% nguon: Archive/2024/TaiLieuGiangDay/1.HamSo/1. K12 - 4in1 - Hamso/ThubaiGoc/Bai5-Dang1.tex #2008
%[Dự án Bài Giảng 12 Nhóm Toán & LaTex, Lê Minh Thiện Anh]
	Cho hàm số $y=x^3-3x^2+2$
	\begin{enumerate}
		\item Tìm điểm $I$ thuộc đồ thị hàm số biết hoành độ của $I$ là nghiệm của phương trình $y''=0$.
		\item Chứng minh rằng $I$ là trung điểm nối hai điểm cực trị của đồ thị hàm số.
	\end{enumerate}
	\loigiai{
		\begin{enumerate}
			\item Ta có $y'=3x^2-6x$.\\
			Suy ra $y''=6x-6$. $y''=0 \Rightarrow 6x-6=0 \Leftrightarrow x=1$.\\
			Do điểm $I$ thuộc đồ thị hàm số biết hoành độ của $I$ là nghiệm của phương trình $y''=0$ nên $I(1;0)$
			\item Do  $y'=3x^2-6x$. $y'=0 \Rightarrow x=0$ và $x=2$.\\
			Mặc khác $y''(0)=-6<0$ nên hàm số đạt cực đại tại $(0;2)$ và $y''(2)=6>0$ nên hàm số đạt cực tiểu tại $(2;-2)$.\\
			Do đó $I$ là là trung điểm nối hai điểm cực trị của đồ thị hàm số.		
		\end{enumerate}
	}
\end{bt}

\begin{bt}%[2D1V5-8]
% nguon: Archive/2024/TaiLieuGiangDay/1.HamSo/1. K12 - 4in1 - Hamso/ThubaiGoc/Bai5-Dang1.tex #2009
%[Dự án Bài Giảng 12 Nhóm Toán & LaTex, Lê Minh Thiện Anh]
	Bạn Việt muốn dùng tấm bìa hình vuông cạnh $6$ dm làm một chiếc hộp không nắp, có đáy là hình vuông bằng cách cắt bỏ đi $4$ hình vuông nhỏ ở bốn góc của tấm bìa (minh họa như hình bên dưới).
	\begin{center}
		\begin{tikzpicture}[>=stealth,line join=round,line cap=round,font=\footnotesize,scale=1]
			\fill[blue!20] (0,0)--(3,0)--(3,3)--(0,3);
			\fill[white] (0,0)--(0,0.4)--(0.4,0.4)--(0.4,0) (2.6,0)--(3,0)--(3,0.4)--(2.6,0.4)
			(2.6,2.6)--(2.6,3)--(3,3)--(3,2.6) (0,2.6)--(0.4,2.6)--(0.4,3)--(0,3);
			\draw (0,0)--(3,0)--(3,3)--(0,3)--cycle;
			\draw[line width=0.2pt] (0.4,0.4)--(2.6,0.4)--(2.6,2.6)--(0.4,2.6)--cycle;
			\draw [line width=0.05pt,<->] (3.05,2.6)--(3.05,3);
			\path (3,2.6)--(3,3)node[pos=0.5,sloped,black,below]{$x$};					
		\end{tikzpicture}~\begin{tikzpicture}[>=stealth,line join=round,line cap=round,font=\footnotesize,scale=1]
			\fill[blue!20] (0,0)--(1.5,-1)--(3,0)--(1.5,1);		
			\fill[black!20] (0,0)--(-0.1,0.4)--(1.4,1.4)--(1.5,1);
			\draw (0,0)--(-0.1,0.4)--(1.4,1.4)--(1.5,1)--cycle;
			\fill[black!20] (1.5,1)--(1.6,1.4)--(3.1,0.4)--(3,0);
			\draw (1.5,1)--(1.6,1.4)--(3.1,0.4)--(3,0)--cycle;
			\fill[black!20] (1.5,-1)--(1.5,-0.6)--(3.1,0.4)--(3,0);
			\draw (1.5,-1)--(1.5,-0.6)--(3.1,0.4)--(3,0)--cycle;
			\fill[black!20] (1.5,-1)--(1.3,-0.7)--(-0.3,0.3)--(0,0);
			\draw (1.5,-1)--(1.3,-0.7)--(-0.3,0.3)--(0,0)--cycle;
%			\node [above=-1mm, fill=white,font=\footnotesize] at (0,-1.5) {\it Hình $11$};
		\end{tikzpicture}
	\end{center}	
	Bạn Việt muốn tìm độ dài cạnh hình vuông cần cắt bỏ để chiếc hộp đạt thể tích lớn nhất.
	\begin{enumerate}
		\item Hãy thiết lập hàm số biểu thị thể tích hộp theo $x$ với $x$ là độ dài cạnh hình vuông cần cắt đi.
		\item Khảo sát và vẽ đồ thị hàm số tìm được.\\		
		Từ đó, hãy tư vấn cho bạn Việt cách giải quyết vấn đề và giải thích vì sao cần chọn giá trị này. (Làm tròn kết quả đến hàng phần mười.)
	\end{enumerate}
	\loigiai{
		\begin{enumerate}
			\item Hãy thiết lập hàm số biểu thị thể tích hộp theo $x$ với $x$ là độ dài cạnh hình vuông cần cắt đi.
			Mặt đáy của hộp là hình vuông có cạnh bằng $6-2x$ (cm), với $0<x<3$. Vậy diện tích của đáy hộp là $S=(6-2x)^2$.\\
			Khối hộp có chiều cao $h=x$ (cm).\\
			Vậy thể tích hộp là $V=S\cdot h=(6-2x)^2 \cdot x=4x^3-24x^2+36x$ (cm$^3$).
			\item Khảo sát và vẽ đồ thị hàm số tìm được.\\
			Xét hàm $f(x)=4x^3-24x^2+36x,\,\,0<x<3$.
			\begin{enumerate}
				\item Tập xác định: $\mathscr{D}=(0;3)$.
				\item Sự biến thiên.
				\begin{itemize}
					\item Giới hạn: $\displaystyle\displaystyle\lim _{x \rightarrow 0^{+}} f(x)=0, \displaystyle\displaystyle\lim _{x \rightarrow 3^{-}} f(x)=0$.
					\item Ta có $f'(x)=12x^2-48x+36\Rightarrow f'(x)=0\Leftrightarrow x^2-4x+3=0\Leftrightarrow \hoac{& x=1\\ & x=3.}$\\
					Ta có bảng biến thiên:
					\begin{center}
						\begin{tikzpicture}[>=stealth]
							\tkzTabInit[nocadre=false,lgt=1,espcl=2,deltacl=0.5]{$x$/.7 ,$y'$/.7,$y$/2}
							{$0$, $1$ , $3$ }
							\tkzTabLine{ , + , $0$ , - , $0$ }
							\tkzTabVar{-/$0$ , +/$16$ , -/$0$ }
						\end{tikzpicture}
					\end{center}
					Hàm số đồng biến trên $(0;1)$ và nghịch biến trên khoảng $(1;3)$.\\
					Hàm số đạt cực đại tại $x=1$.
					\item Đồ thị hàm số đi qua các điểm $(0 ; 0),(1 ; 16),(3 ; 0)$.
					\begin{center}
						\begin{tikzpicture}[line cap=butt,line join=miter,>=stealth,scale=0.8,font=\footnotesize,y=0.5cm]
							\tikzset{declare function={xmin=-1;xmax=4;ymin=-1;ymax=17;},
								smooth,samples=450}
							\draw[->] (xmin,0)--(xmax,0) node[shift={(0:7pt)},]{$ x $};
							\draw[->] (0,ymin)--(0,ymax) node[shift={(90:7pt)}]{$ y $};
							\fill (0,0) node[shift={(-150:7pt)}]{$ O $};
							\clip (xmin,ymin) rectangle (xmax,ymax);
							\foreach \i in {1}{
								\draw(\i,1.5pt)--(\i,-1.5pt)node[below]{$\i$};}
							\foreach \j in {16}{
								\draw(-1.5pt,\j)--(1.5pt,\j) node[left]{$\j$};}
							%	\draw(-1.5pt,-1)--(1.5pt,-1)node[shift={(7pt,0pt)}]{$-1$};	
							\def\f(#1){4*(#1)^3-24*(#1)^2+36*(#1)} % Đồ thị hàm số y=x^3+3x^2+3x+1
							\def\c{-1}
							\def\d{0}
							\def\e{-2}	
							\pgfmathsetmacro\fc{\f(\c)}
							\pgfmathsetmacro\fd{\f(\d)}
							\pgfmathsetmacro\fe{\f(\e)}
							\draw[samples=100] plot[domain=0:3] (\x,{\f(\x)});
							\foreach \x/\y in {\c/\fc,\d/\fd,\e/\fe}{
								\draw[dashed] (1,0)--(1,16)--(0,16);
								\fill[white,draw=black] (\x,\y) circle (1pt);}	
							%\node at (-2,-3.2) [right,fill=white,font=\footnotesize]{\it Hình LT $2b$};
						\end{tikzpicture}
					\end{center}
				\end{itemize}
			\end{enumerate}
			Vậy hình vuông mà bạn Việt cần cắt bỏ pải có độ dài cạnh $x=1$ dm thì chiếc hộp đạt thể tích lớn nhất.
		\end{enumerate}
	}
\end{bt}

\begin{bt}%[2D1N3-4]
% nguon: Archive/2024/TaiLieuGiangDay/1.HamSo/1. K12 - 4in1 - Hamso/ThubaiGoc/Bai3-Dang6.tex #2020
[Mức độ 1]%[BG12-4IN1, Nguyễn Khánh Trọng]
	Cho hàm số $y=f(x)$ có bảng biến thiên như hình vẽ.
	\begin{center}
		\begin{tikzpicture}
			\tkzTabInit[lgt=1.2,espcl=2]
			{$x$ /.7, $f’(x)$ /.7, $f(x)$ /2}
			{$-\infty$,$-3$,$4$,$+\infty$}
			\tkzTabLine{ ,+,z,-,z,+, }
			\tkzTabVar{-/$-\infty$,+/$2$,-/$-2$,+/$+\infty$}
		\end{tikzpicture}
	\end{center}
	Tìm tất cả giá trị của tham số $m$ để bất phương trình $f(x)\le m$ có nghiệm $x\in(-3;+\infty)$.
	\loigiai{
		Ta có $f(x)\le m$ có nghiệm $x\in(-3;+\infty)$ khi và chỉ khi
		$\min\limits_{(-3;+\infty)}f(x)\le m\Leftrightarrow m\ge -2$.
	}	
\end{bt}

\begin{bt}%[2D1N3-4]
% nguon: Archive/2024/TaiLieuGiangDay/1.HamSo/1. K12 - 4in1 - Hamso/ThubaiGoc/Bai3-Dang6.tex #2021
[Mức độ 1]%[BG12-4IN1, Nguyễn Khánh Trọng]
	\immini{Cho hàm số $y=f(x)=\log_2x$ có đồ thị như hình vẽ bên. Tìm tất cả giá trị của tham số $m$ để bất phương trình $f(x)\ge m$ có nghiệm $x\in(0;4)$.
	}{\begin{tikzpicture}[line join=round,line cap=round,>=stealth,scale=0.5,font=\footnotesize]
		\def \xmin{-1} 		\def \xmax{6}
		\def \ymin{-1} 		\def \ymax{3.5}
		\def \fx{ln(\x)/ln(2)}
		\draw[->] (\xmin,0)--(\xmax,0) node[below left] {\normalsize$x$};
		\draw[->] (0,\ymin)--(0,\ymax) node[below left] {\normalsize$y$};
		\draw (0,0) node[below left] {$O$};
		\foreach \x in {1,4} \fill (\x,0) circle (1.25pt) node[below] {$\x$};
		\foreach \y in {2} \fill (0,\y) circle (1.25pt) node[left] {$\y$};
		\draw[dashed] (4,0)--(4,2)--(0,2);
		\begin{scope}
			\clip (\xmin+0.01,\ymin+0.01) rectangle (\xmax-0.01,\ymax-0.01);
			\draw[thick,samples=150,domain=0.1:\xmax-0.01,smooth,variable=\x] plot (\x,{\fx});
		\end{scope}
		\end{tikzpicture}}
	\loigiai{
		Ta có bất phương trình $f(x)\ge m$ có nghiệm $x\in(0;4)$ khi và chỉ khi
		$$f(4)> m\Leftrightarrow m<2.$$
	}	
\end{bt}

\begin{bt}%[2D1H3-4]
% nguon: Archive/2024/TaiLieuGiangDay/1.HamSo/1. K12 - 4in1 - Hamso/ThubaiGoc/Bai3-Dang6.tex #2022
[Mức độ 2]%[BG12-4IN1, Nguyễn Khánh Trọng]
	Tìm tất cả giá trị của tham số $m$ để bất phương trình $\mathrm{e}^x\ge m$ có nghiệm $x\in[0;1]$.
	\loigiai{
		Xét hàm số $f(x)=\mathrm{e}^x$ trên đoạn $[0;1]$.\\
		Ta có $f(x)$ luôn đồng biến trên $(0;1)\Rightarrow \max\limits_{[0;1]}f(x)=f(1)=\mathrm{e}$.\\
		Khi đó, bất phương trình $\mathrm{e}^x\ge m$ có nghiệm $x\in[0;1]$ khi và chỉ khi
		$$\max\limits_{[0;1]}f(x)\ge m\Leftrightarrow m\le \mathrm{e}.$$
		Vậy $m\in(-\infty;\mathrm{e}]$.
	}	
\end{bt}

\begin{bt}%[2D1H3-4]
% nguon: Archive/2024/TaiLieuGiangDay/1.HamSo/1. K12 - 4in1 - Hamso/ThubaiGoc/Bai3-Dang6.tex #2023
[Mức độ 2]%[BG12-4IN1, Nguyễn Khánh Trọng]
	Tìm tất cả giá trị của tham số $m$ để bất phương trình $x^3+3x+1\le m$ có nghiệm $x\in[-1;1]$.
	\loigiai{
		Xét hàm số $f(x)=x^3+3x+1$ trên đoạn $[-1;1]$.\\
		Ta có $f'(x)=3x^2+3>0$, $\forall x\in[-1;1]$ nên hàm số $f(x)$ luôn đồng biến trên $(-1;1)$.\\
		Suy ra $\min\limits_{[-1;1]}f(x)=f(-1)=-3$.\\
		Khi đó, bất phương trình $x^3+3x+1\le m$ có nghiệm $x\in[-1;1]$ khi và chỉ khi
		$$\min\limits_{[-1;1]}f(x)\le m\Leftrightarrow m\ge -3.$$
		Vậy $m\in[-3;+\infty)$.
	}	
\end{bt}

\begin{bt}%[2D1H3-4]
% nguon: Archive/2024/TaiLieuGiangDay/1.HamSo/1. K12 - 4in1 - Hamso/ThubaiGoc/Bai3-Dang6.tex #2024
[Mức độ 2]%[BG12-4IN1, Nguyễn Khánh Trọng]
	Tìm tất cả giá trị của tham số $m$ để bất phương trình $x-m\le mx-3$ có nghiệm $x\in[0;3]$.
	\loigiai{
		Với $x\in[0;3]$, ta có 
		$$x-m\le mx-3\Leftrightarrow x+3\le mx+m\Leftrightarrow x+3\le m(x+1)\Leftrightarrow \dfrac{x+3}{x+1}\le m.$$
		Xét hàm số $f(x)=\dfrac{x+3}{x+1}$ trên đoạn $[0;3]$.\\
		Ta có $f'(x)=-\dfrac{2}{(x+1)^2}<0$, $\forall x\in[0;3]$ nên hàm số $f(x)$ luôn nghịch biến trên $(0;3)$.\\
		Suy ra $\min\limits_{[0;3]}f(x)=f(3)=\dfrac{3}{2}$.\\
		Khi đó, bất phương trình $x-m\le mx-3$ có nghiệm $x\in[0;3]$ khi và chỉ khi
		$$\min\limits_{[0;3]}f(x)\le m\Leftrightarrow m\ge \dfrac{3}{2}.$$
		Vậy $m\in\left[\dfrac{3}{2};+\infty\right)$.
	}	
\end{bt}

\begin{bt}%[2D1H3-4]
% nguon: Archive/2024/TaiLieuGiangDay/1.HamSo/1. K12 - 4in1 - Hamso/ThubaiGoc/Bai3-Dang6.tex #2025
[Mức độ 2]%[BG12-4IN1, Nguyễn Khánh Trọng]
	Tìm tất cả giá trị của tham số $m$ để bất phương trình $\dfrac{(x-2)^2}{x-1}>m$ có nghiệm $x\in[2;3]$.
	\loigiai{
		Xét hàm số $f(x)=\dfrac{(x-2)^2}{x-1}=\dfrac{x^2-4x+4}{x-1}$ trên đoạn $[2;3]$.\\
		Khi đó ta có $f'(x)=\dfrac{x^2-2x}{(x-1)^2}$; $f'(x)=0\Rightarrow \hoac{&x=0\notin[2;3]\\&x=2\in[2;3].}$\\
		Ta có $f(2)=0$; $f(3)=1$.\\
		Suy ra $\max\limits_{[2;3]}f(x)=f(3)=1$.\\
		Do đó, bất phương trình $\dfrac{(x-2)^2}{x-1}>m$ có nghiệm $x\in[2;3]$ khi và chỉ khi
		$$\max\limits_{[2;3]}f(x)>m\Leftrightarrow m<1.$$
		Vậy $m\in\left(-\infty;1\right)$.
	}	
\end{bt}

\begin{bt}%[2D1H3-4]
% nguon: Archive/2024/TaiLieuGiangDay/1.HamSo/1. K12 - 4in1 - Hamso/ThubaiGoc/Bai3-Dang6.tex #2026
[Mức độ 2]%[BG12-4IN1, Nguyễn Khánh Trọng]
	Tìm tất cả giá trị của tham số $m$ để bất phương trình $\ln x-x\le m$ có nghiệm $x\in\left(\dfrac{1}{\mathrm{e}};\mathrm{e}\right)$.
	\loigiai{
		Xét hàm số $f(x)=\ln x-x$ trên đoạn $\left[\dfrac{1}{\mathrm{e}};\mathrm{e}\right]$.\\
		Khi đó ta có $f'(x)=\dfrac{1}{x}-1=\dfrac{1-x}{x}$; $f'(x)=0\Rightarrow x=1$.\\
		Ta có $f\left(\dfrac{1}{\mathrm{e}}\right)=-1-\dfrac{1}{\mathrm{e}}$; $f(1)=-1$; $f(\mathrm{e})=1-\mathrm{e}$.\\
		Suy ra $\min\limits_{\left[\tfrac{1}{\mathrm{e}};\mathrm{e}\right]}f(x)=f(\mathrm{e})=1-\mathrm{e}$.\\
		Do đó, bất phương trình $\ln x-x\le m$ có nghiệm $x\in\left(\dfrac{1}{\mathrm{e}};\mathrm{e}\right)$ khi và chỉ khi
		$$\min\limits_{\left[\tfrac{1}{\mathrm{e}};\mathrm{e}\right]}f(x)<m\Leftrightarrow m>1-\mathrm{e}.$$
		Vậy $m\in\left(1-\mathrm{e};+\infty\right)$.
	}	
\end{bt}

\begin{bt}%[2D1H3-4]
% nguon: Archive/2024/TaiLieuGiangDay/1.HamSo/1. K12 - 4in1 - Hamso/ThubaiGoc/Bai3-Dang6.tex #2030
[Mức độ 2]%[BG12-4IN1, Nguyễn Khánh Trọng]
	Tìm $m$ để bất phương trình $x^4-4x^2+9>m$ có nghiệm trên đoạn $\left[-\sqrt{3};2\right]$.
	\dapso{$m<9$}
	\loigiai{
		Hàm số $f(x)=x^4-4x^2+9$ liên tục trên $\left[-\sqrt{3};2\right]$.\\
		Ta có $f'(x)=4x^3-8x$;
		$f'(x)=0\Leftrightarrow\hoac{&x=0\\&x=\pm\sqrt{2}.}$\\
		$f(0)=9$, $f\left(\pm\sqrt{2}\right)=5$, $f\left(-\sqrt{3}\right)=6$, $f(2)=9$.\\
		Do đó $\max\limits_{\left[-\sqrt{3};2\right]} f(x)=9$.\\
		Bất phương trình $x^4-4x^2+9>m$ có nghiệm trên đoạn $\left[-\sqrt{3};2\right]$ khi và chỉ khi
		$$\max\limits_{\left[-\sqrt{3};2\right]} f(x)>m\Leftrightarrow m<9.$$
	}
\end{bt}

\begin{bt}%[2D1H3-4]
% nguon: Archive/2024/TaiLieuGiangDay/1.HamSo/1. K12 - 4in1 - Hamso/ThubaiGoc/Bai3-Dang6.tex #2031
[Mức độ 2]%[BG12-4IN1, Nguyễn Khánh Trọng]
	Tìm giá trị của tham số $m$ để  bất phương trình $\dfrac{x^2+3x+3}{x+1}\geq m$ có nghiệm $x\in [0;1]$.
	\dapso{$m\le \dfrac{7}{2}$}
	\loigiai{Để bất phương trình $\dfrac{x^2+3x+3}{x+1}\geq m$ có nghiệm $x\in [0;1]$ tương đương với
		$$\max\limits_{[0;1]}f(x)\ge m\,\, \text{với} \,\, f(x)=\dfrac{x^2+3x+3}{x+1}.$$
		Ta có $f'(x)=\dfrac{x^2+2x}{(x+1)^2}=0\Leftrightarrow x=0$ (nhận) hoặc $x=-2$ (loại).\\
		$f(0)=3$, $f(1)=\dfrac{7}{2}$. \\
		Suy ra $\max\limits_{[0;1]}f(x)=\dfrac{7}{2}$. Do đó $\dfrac{7}{2}\geq m$.
	}
\end{bt}

\begin{bt}%[2D1H3-4]
% nguon: Archive/2024/TaiLieuGiangDay/1.HamSo/1. K12 - 4in1 - Hamso/ThubaiGoc/Bai3-Dang6.tex #2032
[Mức độ 2]%[BG12-4IN1, Nguyễn Khánh Trọng]
	Tìm $m$ để bất phương trình $x-\sqrt{16-x^2}-m\le0$ có nghiệm.
	\dapso{$m\ge -4\sqrt{2}$}
	\loigiai{
		Ta có $x-\sqrt{16-x^2}-m\le0\Leftrightarrow x-\sqrt{16-x^2}\le m$.\\
		Gọi $f(x)=x-\sqrt{16-x^2}$ có tập xác định $\mathscr{D}=[-4;4]$.\\
		Ta có $f'(x)= 1 + \dfrac{x}{\sqrt {16-x^2}} = \dfrac{\sqrt{16-x^2}+x}{\sqrt {16-x^2}}$.\\
		$f'(x)=0 \Leftrightarrow \sqrt {16-x^2}+x= 0 \Leftrightarrow \sqrt {16-x^2}=-x \Leftrightarrow \heva{
			&16-x^2=(-x)^2\\
			&-x\ge 0
		} \Leftrightarrow \heva{
			&2x^2 = 16\\
			&x\le 0
		} \Leftrightarrow x=-2\sqrt2$.\\
		Ta có $f(-4)=-4$; $f(4)=4$; $f(-2\sqrt2)=-4\sqrt{2}$.\\
		Suy ra $\min\limits_{\mathscr{D}}f(x)=-4\sqrt{2}$.\\
		Bất phương trình $f(x)=x-\sqrt{16-x^2}-m\le0$ có nghiệm khi và chỉ khi $$\min\limits_{\mathscr{D}}f(x)\le m\Leftrightarrow m\ge -4\sqrt{2}.$$
	}
\end{bt}

\begin{bt}%[2D1H3-4]
% nguon: Archive/2024/TaiLieuGiangDay/1.HamSo/1. K12 - 4in1 - Hamso/ThubaiGoc/Bai3-Dang6.tex #2033
[Mức độ 2]%[BG12-4IN1, Nguyễn Khánh Trọng]
	Tìm $m$ để bất phương trình $\mathrm{e}^{x^3-3x+3}\ge m$ có nghiệm trên đoạn $[0;2]$.
	\dapso{$m\le \mathrm{e}^5$}
	\loigiai{
		Xét hàm số $f(x)=\mathrm{e}^{x^3-3x+3}$, có $f'(x)=(3x^2-3)\mathrm{e}^{x^3-3x+3}$.\\
		$f'(x)=0\Leftrightarrow 3x^2-3=0\Leftrightarrow\hoac{&x=1\in[0;2]\\&x=-1\notin[0;2].}$\\
		Có $f(0)=\mathrm{e}^3$, $f(2)=\mathrm{e}^5$, $f(1)=\mathrm{e}^3$.\\
		Suy ra $\max\limits_{[0;2]}f(x)=\mathrm{e}^5$.\\
		Bất phương trình $\mathrm{e}^{x^3-3x+3}\ge m$ có nghiệm trên đoạn $[0;2]$ khi và chỉ khi
		$$\max\limits_{[0;2]}f(x)\ge m\Leftrightarrow \mathrm{e}^5\ge m.$$
	}
\end{bt}

\begin{bt}%[2D1V3-4]
% nguon: Archive/2024/TaiLieuGiangDay/1.HamSo/1. K12 - 4in1 - Hamso/ThubaiGoc/Bai3-Dang6.tex #2034
[Mức độ 3]%[BG12-4IN1, Nguyễn Khánh Trọng]
	Tìm tất cả các giá trị của tham số $m$ để bất phương trình $m\left( x-1 \right)-\sqrt{x-4}\le m+1$ có nghiệm.
	\dapso{$m\in \left( -\infty ;\dfrac{1+\sqrt{3}}{4} \right]$}
	\loigiai{
		Xét bất phương trình $m\left( x-1 \right)-\sqrt{x-4}\le m+1$\quad (1).\\
		Điều kiện xác định của bất phương trình là $x\ge 4$.\\
		Đặt $t=\sqrt{x-4}\Rightarrow t\ge 0$ và $x=t^2+4$, khi đó (1) trở thành
		$m\left(t^2+3 \right)-t\le m+1\Leftrightarrow m\le \dfrac{t+1}{t^2+2}$\quad (2).\\
		Xét hàm số $f(t)=\dfrac{t+1}{t^2+2}$ với $t\ge 0$. Ta có
		$f'(t)=\dfrac{-t^2-2t+2}{(t^2+2)^2}$.\\
		Cho $f'(t)=0\Leftrightarrow -t^2-2t+2=0\Leftrightarrow \hoac{& t=-1+\sqrt{3} \\ & t=-1-\sqrt{3}\,\,(\text{loại}).}$\\
		Bảng biến thiên của hàm số $f(t)$ trên nửa khoảng $[0;+\infty)$.
		\begin{center}
			\begin{tikzpicture}
				\tkzTabInit[lgt=1.2,espcl=3]
				{$t$ /.7, $f’(t)$ /.7, $f(t)$ /2.5}
				{$0$,$-1+\sqrt{3}$,$+\infty$}
				\tkzTabLine{ ,+,z,- }
				\tkzTabVar{-/$\dfrac{1}{2}$,+/$\dfrac{1+\sqrt{3}}{4}$,-/$0$}
			\end{tikzpicture}
		\end{center}
		Bất phương trình (1) có nghiệm $\Leftrightarrow $ (2) có nghiệm $t\in [0;+\infty)\Leftrightarrow m\le \dfrac{1+\sqrt{3}}{4}$.
	}
\end{bt}

\begin{bt}%[2D1H5-7]
% nguon: Archive/2024/TaiLieuGiangDay/1.HamSo/1. K12 - 4in1 - Hamso/ThubaiGoc/Bai5-Dang10.tex #2045
%[Dự Án Bài Giảng 12 4 in 1, Lê Văn Toàn]
	Tìm tâm đối xứng của đồ thị hàm số $y=\dfrac{x-2}{-x+2}$.
	\loigiai{
		Tâm đối xứng của đồ thị hàm số $y=\dfrac{x-2}{-x+2}$ chính là giao điểm của $2$ đường tiệm cận $x=2$ và $y=-1$. \\
		Vậy $I(2;-1)$.
	}
\end{bt}

\begin{bt}%[2D1H5-7]
% nguon: Archive/2024/TaiLieuGiangDay/1.HamSo/1. K12 - 4in1 - Hamso/ThubaiGoc/Bai5-Dang10.tex #2046
%[Dự Án Bài Giảng 12 4 in 1, Lê Văn Toàn]
	Trên đồ thị hàm số $y=\dfrac{x+1}{2x-1}$ có bao nhiêu điểm mà tung độ bằng $\dfrac{1}{2}$?
	\loigiai{
		Gọi $M\left(x;\dfrac{1}{2}\right)$ là điểm thuộc đồ thị hàm số, $x\neq \dfrac{1}{2}$.\\
		Ta có $\dfrac{x+1}{2x-1}=\dfrac{1}{2}\Leftrightarrow 2x+2=2x-1\Leftrightarrow 0x=-3$. Phương trình này vô nghiệm.\\
		Vậy không có điểm nào thuộc đồ thị hàm số đã cho mà có tung độ bằng $\dfrac{1}{2}$.
	}
\end{bt}

\begin{bt}%[2D1H5-7]
% nguon: Archive/2024/TaiLieuGiangDay/1.HamSo/1. K12 - 4in1 - Hamso/ThubaiGoc/Bai5-Dang10.tex #2047
%[Dự Án Bài Giảng 12 4 in 1, Lê Văn Toàn]
	Đồ thị hàm số $y=\dfrac{x-1}{x+1}$ có bao nhiêu điểm mà tọa độ của nó đều là các số nguyên?
	\loigiai{
		Ta có $y=\dfrac{x-1}{x+1}=1-\dfrac{2}{x+1}$.\\
		Suy ra tọa độ các điểm thuộc đồ thị hàm số đã cho có dạng $\left(x;1-\dfrac{2}{x+1}\right), x\neq -1$.\\
		Do đó điểm thuộc đồ thị có tọa độ là các số nguyên khi và chỉ khi $x\in\mathbb{Z}$ và $x+1$ là ước của $2$. Suy ra $x+1=\pm 1$ hoặc $x+1=\pm 2\Leftrightarrow\hoac{&x=0\\&x=-2\\&x=1\\&x=-3}$. Vậy có $4$ điểm.
	}
\end{bt}

\begin{bt}%[2D1H5-7]
% nguon: Archive/2024/TaiLieuGiangDay/1.HamSo/1. K12 - 4in1 - Hamso/ThubaiGoc/Bai5-Dang10.tex #2048
%[Dự Án Bài Giảng 12 4 in 1, Lê Văn Toàn]
	Đồ thị của hàm số $y=\dfrac{4x-3}{x-2}$ nhận điểm $I(a;b)$ làm tâm đối xứng. Tính giá trị của $a+b$.
	\loigiai{
		Đồ thị hàm số nhận giao điểm của hai đường tiệm cận làm tâm đối xứng.\\
		Đồ thị hàm số có tiệm cận đứng $x=2$ và tiệm cận ngang $y=4$.\\
		Vậy $I(2;4)$ là tâm đối xứng của đồ thị, suy ra $a+b=2+4=6$.
	}
\end{bt}

\begin{bt}%[2D1N5-7]
% nguon: Archive/2024/TaiLieuGiangDay/1.HamSo/1. K12 - 4in1 - Hamso/ThubaiGoc/Bai5-Dang10.tex #2052
%[Dự Án Bài Giảng 12 4 in 1, Lê Văn Toàn]
	Cho hàm số $y=x^3-3x^2+2$ có đồ thị $(C)$. Tìm các cặp điểm $M$, $N$ nằm trên $(C)$ đối xứng nhau qua điểm $I(1; 0)$
	\loigiai{
		Ta có $y'=3x^2-6x$, $y''=6 x-6$, $y''=0 \Leftrightarrow x=1 \Rightarrow y=0$.
		Do đó $I(1; 0)$ là tâm đối xứng của $(C)$.
	}
\end{bt}

\begin{bt}%[2D1H5-7]
% nguon: Archive/2024/TaiLieuGiangDay/1.HamSo/1. K12 - 4in1 - Hamso/ThubaiGoc/Bai5-Dang10.tex #2053
%[Dự Án Bài Giảng 12 4 in 1, Lê Văn Toàn]
	Đồ thị hàm số $y=\dfrac{3}{2 x-1}$ có bao nhiêu điểm có tọa độ nguyên dương?
	\loigiai{
	Giả sử $M\left(x_0; \dfrac{3}{2 x_0-1}\right)$ điểm nằm trên đồ thị hàm số.\\
	Khi đó $M$ có tọa độ nguyên dương khi $x_0$ nguyên dương và $\dfrac{3}{2 x_0-1}$ nguyên dương.\\
	Suy ra $2 x_0-1 \in\{1; 3\}$ hay $x_0 \in\{1; 2\}$.\\
	Vậy có $2$ điểm thỏa mãn yêu cầu bài toán.
	}
\end{bt}

\begin{bt}%[2D1H5-7]
% nguon: Archive/2024/TaiLieuGiangDay/1.HamSo/1. K12 - 4in1 - Hamso/ThubaiGoc/Bai5-Dang10.tex #2054
%[Dự Án Bài Giảng 12 4 in 1, Lê Văn Toàn]
	Đồ thị của hàm số $y=x^3-3x^2+m x+m$ với $m$ là tham số luôn đi qua một điểm $M$ cố định. Tính độ dài của đoạn $O M$ với $O$ là gốc tọa độ.
	\loigiai{
		Ta có $y=x^3-3 x^2+m x+m \Leftrightarrow(x+1) m+x^3-3 x^2-y=0$.\\
		Phương trình nghiệm đúng với mọi $m$ khi $$\heva{&x+1=0\\&x^3-3x^2-y=0}\Leftrightarrow\heva{&x=-1\\&y=-4.}$$
		Suy ra $M(-1;-4)$ là điểm cố định của đồ thị hàm số.
		Khi đó $O M=\sqrt{17}$.
	}
\end{bt}

\begin{bt}%[2D1V5-7]
% nguon: Archive/2024/TaiLieuGiangDay/1.HamSo/1. K12 - 4in1 - Hamso/ThubaiGoc/Bai5-Dang10.tex #2055
%[Dự Án Bài Giảng 12 4 in 1, Lê Văn Toàn]
	Xác định tọa độ điểm $M$ có hoành độ dương nằm trên đồ thị hàm số $(C)\colon y=\dfrac{x+1}{x-2}$ sao cho khoảng cách từ $M$ đến tiệm cận ngang của $(C)$ bằng $1$.
	\loigiai{
	Giả sử $M\left(x_0; \dfrac{x_0+1}{x_0-2}\right) \in(C)\left(x_0>0; x_0 \neq 2\right)$.\\
	Đồ thị $(C)$ có đường tiệm cận ngang là $y=1$.\\
	Theo giả thiết 
	\begin{eqnarray*}
		& & \left|\dfrac{x_0+1}{x_0-2}-1\right|=1\\
		&\Leftrightarrow & \left|\dfrac{3}{x_0-2}\right|=1\\
		&\Leftrightarrow & \hoac{&x_0=5\Rightarrow y_0=2\\&x_0=-1\text{ (loại)}.}
	\end{eqnarray*}
	}
\end{bt}

\begin{bt}%[2D1V5-7]
% nguon: Archive/2024/TaiLieuGiangDay/1.HamSo/1. K12 - 4in1 - Hamso/ThubaiGoc/Bai5-Dang10.tex #2056
%[Dự Án Bài Giảng 12 4 in 1, Lê Văn Toàn]
	Đồ thị của hàm số $y=\dfrac{x-m}{x+m}$ với $m$ là tham số $(m \neq 0)$ luôn đi qua một điểm $M$ cố định. Tính khoảng cách từ $M$ đến đường tiệm cận ngang của đồ thị hàm số.
	\loigiai{
	Đồ thị hàm số có đường tiệm cận ngang $\Delta \colon y-1=0$.\\
	Ta có $y=\dfrac{x-m}{x+m} \Leftrightarrow(y+1) m+x y-x=0$.\\
	Phương trình nghiệm đúng với mọi $m$ khi $$\heva{&y+1=0\\&xy-x=0}\Leftrightarrow\heva{&x=0\\&x=-1.}$$
	Suy ra $M(0;-1)$ là điểm cố định của đồ thị hàm số.
	Khi đó $\mathrm{d}(M; \Delta)=|-1-1|=2$.
	}
\end{bt}

\begin{bt}%[2D1N5-8]
% nguon: Archive/2024/TaiLieuGiangDay/1.HamSo/1. K12 - 4in1 - Hamso/ThubaiGoc/Bai5-Dang11.tex #2067
 [SGK12-KNTT ]
	Giả sử một hạt chuyển động trên một trục thẳng đứng chiều dương hướng lên trên sao cho toạ độ của hạt (đơn vị: mét) tại thời điểm $t$ (giây) là $y=t^3-12t+3, t \geq 0$. Tìm vận tốc của hạt sau $3$ giây.
	\loigiai{
		Vận tốc của hạt chuyển động là $v(t)=y'=3t^2-12$.\\
		Suy ra vận tốc của hạt sau $3$ giây là $v(3)=y'=3\cdot 3^2-12=15$ m/s.
	}
\end{bt}

\begin{bt}%[2D1N5-8]
% nguon: Archive/2024/TaiLieuGiangDay/1.HamSo/1. K12 - 4in1 - Hamso/ThubaiGoc/Bai5-Dang11.tex #2068
[SGK12-KNTT ]
	Để loại bỏ $x\%$ chất gây ô nhiễm không khí từ khí thải của một nhà máy, người ta ước tính chi phí cần bỏ ra là
	$$
	C(x)=\dfrac{300 x}{100-x} \text { (triệu đồng), } 0 \leq x < 100.
	$$		
	Tính chi phí cần bỏ ra để loại bỏ được $90 \%$ chất gây ô nhiễm không khí.
	\loigiai{
		Chi phí cần bỏ ra để loại bỏ được $90 \%$ chất gây ô nhiễm không khí là 
		$$
		C(90)=\dfrac{300 \cdot 90}{100-90} \text { (triệu đồng)
		}=2700 \text { (triệu đồng) }.
		$$
	}
\end{bt}

\begin{bt}%[2D1H5-8]
% nguon: Archive/2024/TaiLieuGiangDay/1.HamSo/1. K12 - 4in1 - Hamso/ThubaiGoc/Bai5-Dang11.tex #2069
[SGK12-KNTT ]
	Giả sử chi phí $C(x)$ (nghìn đồng) để sản xuất $x$ đơn vị của một loại hàng hoá nào đó được cho bởi hàm số $C(x)=30\,000+300x-2{,}5x^2+0{,}125x^3$. Tính chi phí sản xuất đơn vị hàng hoá thứ 201.
	\loigiai{
		Chi phí sản xuất đơn vị hàng hoá thứ $201$ là
		$$
		C(201)-C(200)=1\,004\,372{,}625- 990\,000=14\,372{,}625 \text { (nghìn đồng).}
		$$				
	}
\end{bt}

\begin{bt}%[2D1H5-8]
% nguon: Archive/2024/TaiLieuGiangDay/1.HamSo/1. K12 - 4in1 - Hamso/ThubaiGoc/Bai5-Dang11.tex #2070
[SGK12-KNTT ]
	Khi máu di chuyển từ tim qua các động mạch chính rồi đến các mao mạch và quay trở lại qua các tĩnh mạch, huyết áp tâm thu (tức là áp lực của máu lên động mạch khi tim co bóp) liên tục giảm xuống. Giả sử một người có huyết áp tâm thu $P$ (tính bằng mmHg) được cho bởi hàm số
	$$
	P(t)=\dfrac{25 t^2+125}{t^2+1}, 0 \leq t \leq 10,
	$$
	trong đó thời gian $t$ được tính bằng giây. Tính tốc độ thay đổi của huyết áp sau $5$ giây kể từ khi máu rời tim.
	\loigiai{
		Ta có tốc độ thay đổi của huyết áp là $P'(t)=\dfrac{-200t}{(t^2+1)^2}$.\\
		Do đó tốc độ thay đổi huyết áp sau $5$ giây là $P'(5)=-\dfrac{250}{169}$.
	}
\end{bt}

\begin{bt}%[2D1H5-8]
% nguon: Archive/2024/TaiLieuGiangDay/1.HamSo/1. K12 - 4in1 - Hamso/ThubaiGoc/Bai5-Dang11.tex #2071
 [SGK12-KNTT ]
	Một cốc chứa $30 \mathrm{ml}$ dung dịch $\mathrm{KOH}$ (potassium hydroxide) với nồng độ $100 \mathrm{mg} / \mathrm{ml}$.
	Một bình chứa dung dịch $\mathrm{KOH}$ khác với nồng độ $8 \mathrm{mg} / \mathrm{ml}$ được trộn vào cốc.
	Tính nồng độ $\mathrm{KOH}$ trong cốc sau khi trộn $10 (\mathrm{ml})$ từ bình chứa.
	\loigiai{
		Giả sử trộn $x(\mathrm{ml})$ từ bình chứa vào cốc, ta được  nồng độ $\mathrm{KOH}$ trong cốc sau khi trộn là $C(x)=\dfrac{8x+3000}{x+100}$.\\
		Khi đó, với $x=10$ ta được $C(10)=\dfrac{8\cdot 10 +3000}{10+100}=28 \mathrm{mg} / \mathrm{ml}$. 
	}
\end{bt}

\begin{bt}%[2D1H5-8]
% nguon: Archive/2024/TaiLieuGiangDay/1.HamSo/1. K12 - 4in1 - Hamso/ThubaiGoc/Bai5-Dang11.tex #2072
 [SGK12-Cùng Khám Phá]
	Một chất điểm chuyển động theo quy luật $s(t)=-t^3+2t^2-t$, với $t$ (đơn vị: giây) là khoảng thời gian tính từ khi vật bắt đầu chuyển động và $s$ (đơn vị: mét) là quãng đường chất điểm di chuyển được trong khoảng thời gian đó. Trong khoảng thời gian 2 giây kể từ khi bắt đầu chuyển động, chất điểm đạt được vận tốc lớn nhất là bao nhiêu?
	\loigiai{
		Vận tốc của chuyển động chính là đạo hàm cấp một của quãng đường. \\
		Do đó $v(t)=s'(t)=-3t^2+4t-1=-3\left(t-\dfrac{2}{3}\right)^2+\dfrac{1}{3}\le\dfrac{1}{3},\; \forall t\in \mathbb{R}$.\\
		Suy ra trong khoảng thời gian 2 giây kể từ khi bắt đầu chuyển động, chất điểm đạt được vận tốc lớn nhất là $\dfrac{2}{3}\mathrm{~m/s}$.
	}
\end{bt}

\begin{bt}%[2D1H5-8]
% nguon: Archive/2024/TaiLieuGiangDay/1.HamSo/1. K12 - 4in1 - Hamso/ThubaiGoc/Bai5-Dang11.tex #2074
 
	Một chất điểm chuyển động theo quy luật $s(t)=-t^3+2t^2-t$, với $t$ (đơn vị: giây) là khoảng thời gian tính từ khi vật bắt đầu chuyển động và $s$ (đơn vị: mét) là quãng đường chất điểm di chuyển được trong khoảng thời gian đó.
	\begin{listEX}[1]
		\item Khảo sát sự biến thiên và vẽ đồ thị của hàm số $s=s(t)$ trên hệ trục toạ độ $tOs$.
		\item Trong khoảng thời gian $2$ giây kể từ khi bắt đầu chuyển động, chất điểm đạt được vận tốc lớn nhất là bao nhiêu?
	\end{listEX}
	\loigiai{
		\begin{enumerate}
			\item Khảo sát sự biến thiên và vẽ đồ thị của hàm số $s=s(t)$ trên hệ trục toạ độ $tOs$.
			\begin{itemize}
				\item $s'(t)=-3t^2+4t-1$, $s'(t)=0\Leftrightarrow \hoac{& t_1=1\\ & t_2=\dfrac{1}{3}.}$
				\item Bảng biến thiên
				\begin{center}
					\begin{tikzpicture}
						\tkzTab[nocadre=false,lgt=1.2,espcl=2.5,deltacl=0.6] 
						{$t$/0.6, $s'$/0.6, $s$/2}  
						{$-\infty$, $\tfrac{1}{3}$, $1$, $+\infty$}
						{,-,z,+,z,-,}  
						{+/ $+\infty$, -/ $-\tfrac{4}{27}$, +/ $0$ , -/ $-\infty$} 
					\end{tikzpicture}
				\end{center}
				\item  Đồ thị hàm số
				\begin{center}
					\begin{tikzpicture}[scale=0.7, font=\footnotesize, line join=round, line cap=round, >=stealth]
						\def\a{-1} % Hệ số a phải khác 0
						\def\b{2}
						\def\c{-1}
						\def\d{0}
						\draw[->] (-1.5,0) -- (3.5,0) node[above] {$t$};
						\draw[->] (0,-3.2) -- (0,5) node[left] {$s$};
						\draw (0,0) circle (1pt) node[above right]{$O$};
						\draw[samples=150,smooth,domain=-1.1:2.2] plot(\x,{\a*(\x)^3+(\b)*(\x)^2+(\c)*(\x)+(\d)});
						%\draw[dashed] (-1,0) node[below]{$ -1 $}|-(0,4)node[above right]{$ 4 $}-|(2,0)node[below]{$ 2 $};
						\foreach \x in {-1,0,1,2,3} 
						\draw[shift={(\x,0)}] (0,1pt)--(0,-1pt);
						\foreach \y in {-1,1,2,3,4} 
						\draw[shift={(0,\y)}] (1pt,0)--(-1pt,0);
					\end{tikzpicture}
				\end{center}
			\end{itemize}
			\item Ta có $v(t)=s'(t)=-3t^2+4t-1$, suy ra $v'(t)=-6t+4=0\Leftrightarrow t=\dfrac{2}{3}<2$.\\
			Ta có bảng biến thiên
			\begin{center}
				\begin{tikzpicture}
					\tkzTabInit[nocadre=false,lgt=1.2,espcl=2.5,deltacl=0.6] 
					{$t$/0.6, $v’(t)$/0.6, $v(t)$/2} % cột đầu tiên
					{$0$, $\tfrac{2}{3}$, $2$} % hàng 1 cột 2
					\tkzTabLine{,+,0,-,} % hàng 2 cột 2
					\tkzTabVar{-/ , +/ $\tfrac{1}{3}$ , -/ } % hàng 3 cột 2
				\end{tikzpicture}
			\end{center}
			Dựa vào bảng biến thiên, vận tốc đạt giá trị lớn nhất là $\dfrac{1}{3}$ khi $t=\dfrac{2}{3}$.
		\end{enumerate}
	}
\end{bt}

\begin{bt}%[2D1H5-8]
% nguon: Archive/2024/TaiLieuGiangDay/1.HamSo/1. K12 - 4in1 - Hamso/ThubaiGoc/Bai5-Dang11.tex #2075
 
	Một trang sách có dạng hình chữ nhật với diện tích là $384$ cm$^2$. Sau khi để lề trên và lề dưới đều là $3$ cm, để lề trái và lề phải đều là $2$ cm. Phần còn lại của trang sách được in chữ. Kích thước tối ưu của trang sách là bao nhiêu để phần in chữ trên trang sách có diện tích lớn nhất?
	\loigiai{
		Gọi chiều rộng, chiều dài của trang sách hình chữ nhật lần lượt là $a$, $b>0$.\\
		Khi đó $ab=384\Rightarrow b=\dfrac{384}{a}$.\\
		Sau khi để lề trên, lề dưới và lề trái, lề phải diện tích còn lại của trang sách là:\\
		$(b-6)(a-4)=\left(\dfrac{384}{a}-6\right)(a-4)=-\dfrac{1536}{a}-6a+408$,\\
		suy ra $f(a)=-\dfrac{1536}{a}-6a+408$.\\
		$f'(a)=\dfrac{1536}{a^2}-6=0\Rightarrow a=16$.\\
		Bảng biến thiên
		\begin{center}
			\begin{tikzpicture}
				\tkzTabInit[nocadre=false,lgt=1.2,espcl=2.5,deltacl=0.6] 
				{$a$/0.6, $f’(a)$/0.6, $f(a)$/2} % cột đầu tiên
				{$0$, $16$, $384$} % hàng 1 cột 2
				\tkzTabLine{,+,0,-,} % hàng 2 cột 2
				\tkzTabVar{-/ , +/ $24$ , -/ } % hàng 3 cột 2
			\end{tikzpicture}
		\end{center}
		Dựa vào bảng biến thiên, để phần in chữ trên trang sách có diện tích lớn nhất thì kích thước dài, rộng là $24$, $16$ cm.
	}
\end{bt}

\begin{bt}%[2D1H5-8]
% nguon: Archive/2024/TaiLieuGiangDay/1.HamSo/1. K12 - 4in1 - Hamso/ThubaiGoc/Bai5-Dang11.tex #2076
 
	Trong một nhà hàng, mỗi tuần để chế biến $x$ phần ăn ($x$ lấy giá trị trong khoảng từ $30$ đến $120$) thì chi phí trung bình (đơn vị: nghìn đồng) của một phần ăn được cho bởi công thức:	
	$$C(x)=2 x-230+\dfrac{7200}{x}.$$
	\begin{enumerate}
		\item Khảo sát và vẽ đồ thị hàm số $C(x)$ trên $[30; 120]$.
		\item Từ kết quả trên, tìm số phần ăn sao cho chi phí trung bình của một phần ăn là thấp nhất.
	\end{enumerate}
	\loigiai{
		\begin{enumerate}
			\item 
			\begin{itemize}
				\item Ta có $C'(x)=2-\dfrac{7200}{x^2}=0\Rightarrow \hoac{& x=60& (\text{nhận})\\ & x=-60&(\text{loại}).}$
				\item Bảng biến thiên
				\begin{center}
					\begin{tikzpicture}
						\tkzTabInit[nocadre=false,lgt=1.2,espcl=2.5,deltacl=0.6] 
						{$x$/0.6, $C’(x)$/0.6, $C(x)$/2} 
						{$30$, $60$, $120$}  
						\tkzTabLine{,-,0,+,} 
						\tkzTabVar{+/ , -/ $10$ , +/ }  
					\end{tikzpicture}
				\end{center}
				\item Đồ thị của hàm số
				\begin{center}
					\begin{tikzpicture}[scale=0.2, font=\footnotesize, line join=round, line cap=round, >=stealth, xscale=0.4,yscale=0.4]
						\draw[->] (-12,0) -- (125,0) node[above] {$x$}; 
						\draw[->] (0,-8) -- (0,75) node[below right] {$y$}; 
						\draw[smooth,domain=30:120,blue,thick] plot(\x,{7200/(\x)+2*(\x)-230});
						\draw[dashed] (30,0)node[below]{$30$}--(30,70) (120,0)node[below]{$120$}--(120,70) (60,0)node[below]{$60$}|-(0,10)node[left]{$10$};
						\fill (0,0) node[below left]{$O$}circle(14pt);
					\end{tikzpicture} 
				\end{center}
			\end{itemize}
			\item Dựa vào đồ thị, số phần ăn $60$ thì chi phí trung bình của một phần ăn thấp nhất $10$ (nghìn đồng).
		\end{enumerate}
	}
\end{bt}

\begin{bt}%[2D1H5-8]
% nguon: Archive/2024/TaiLieuGiangDay/1.HamSo/1. K12 - 4in1 - Hamso/ThubaiGoc/Bai5-Dang11.tex #2077
 
	Một người nông dân có $15\;000\;000$ đồng để làm một hàng rào hình chữ $E$ dọc theo một con sông bao quanh hai khu đất trồng rau có dạng hai hình chữ nhật bằng nhau. Đối với mặt hàng rào song song với bờ sông thì chi phí nguyên vật liệu là $60\;000$ đồng/mét, còn đối vối ba mặt hàng rào song song nhau thì chi phí nguyên vật liệu là $50\;000$ đồng/mét, mặt giáp với bờ sông không phải rào. Tìm diện tích lớn nhất của hai khu đất thu được sau khi làm hàng rào.
	\begin{center}
		\begin{tikzpicture}[line join = round, line cap=round,>=stealth,font=\footnotesize,scale=1]
			\draw[fill=green] 
			(-3,-2.1) rectangle (11,0);		
			\draw[fill=cyan!50] 
			(-3,0) rectangle (11,2);	
			
			\draw[fill=green!20] (0,0)coordinate (A)--(10,0)coordinate (B)--(8,-2)coordinate(C)--(-2,-2)coordinate (D)--cycle
			(5,0)coordinate (E)--(3,-2)coordinate (F)
			;
			\foreach  \t in {0,0.1,0.2,...,1}
			\draw[brown,line width=2pt] ($(E)!\t!(F)$)--
			($(E)!\t!(F)+(0,1)$);
			
			\foreach  \t in {0,0.1,0.2,...,1}
			\draw[brown,line width=2pt] ($(A)!\t!(D)$)--
			($(A)!\t!(D)+(0,1)$);
			
			\foreach  \r in {0,0.1,0.2,...,1}
			\draw[brown,line width=2pt] ($(B)!\r!(C)$)--
			($(B)!\r!(C)+(0,1)$);
			
			\foreach  \i in {0,0.025,0.05,...,1}
			\draw[brown,line width=2pt] ($(C)!\i!(D)$)--
			($(C)!\i!(D)+(0,1)$);
			\draw[brown,line width=2pt] ($(C)+(0,0.8)$)--($(D)+(0,0.8)$);
			\draw[brown,line width=2pt] ($(C)+(0,0.8)$)--($(B)+(0,0.8)$);
			\draw[brown,line width=2pt] ($(A)+(0,0.8)$)--($(D)+(0,0.8)$);
		\end{tikzpicture}
	\end{center}
	\loigiai{
		Gọi chiều dài, chiều rộng của khu đất trồng rau lần lượt là $a$ và $b$, điều kiện $a$, $b>0$.\\
		Theo đề bài $2\cdot 60\cdot b+3\cdot 50\cdot a=15\;000\Rightarrow b=\dfrac{1500-15a}{12}$.\\
		Diện tích của khu đất trồng rau là $S=2ab=2a\left(\dfrac{1500-15a}{12}\right)=-\dfrac{5a^2}{2}+250a$.\\
		$S'(a)=-5a+250=0\Rightarrow a=50$.\\
		Bảng biến thiên
		\begin{center}
			\begin{tikzpicture}
				\tkzTabInit[nocadre=false,lgt=1.2,espcl=2.5,deltacl=0.6] 
				{$a$/0.6, $S’(a)$/0.6, $S(a)$/2} % cột đầu tiên
				{$0$, $50$, } % hàng 1 cột 2
				\tkzTabLine{,+,0,-,} % hàng 2 cột 2
				\tkzTabVar{-/ , +/ $6250$ , -/ } % hàng 3 cột 2
			\end{tikzpicture}
		\end{center}
		Dựa vào bảng biến thiên, suy ra diện tích lớn nhất của hai khu đất là $6250$ khi chiều rộng là $a=50$ m, chiều dài là $b=62{,}5$ m.
		
	}
\end{bt}

\begin{bt}%[2D1V5-8]
% nguon: Archive/2024/TaiLieuGiangDay/1.HamSo/1. K12 - 4in1 - Hamso/ThubaiGoc/Bai5-Dang11.tex #2078
 
	Một công ty kinh doanh bất động sản có $20$ căn hộ cho thuê. Biết rằng nếu cho thuê mỗi căn hộ với giá $2$ triệu đồng/$1$ tháng thì tất cả các căn hộ đều có người thuê. Nhưng cứ mỗi lần tăng giá cho thuê mỗi căn hộ thêm $200$ nghìn đồng/$1$ tháng thì có thêm một căn hộ bị bỏ trống. Hỏi công ty nên cho thuê mỗi căn hộ bao nhiêu tiền một tháng để tổng số tiền thu được là lớn nhất?
	\loigiai{
		Gọi $p$ là số căn hộ cho thuê; $x$ (nghìn đồng) là giá thuê mỗi căn hộ trên tháng. Ta cần xác định hàm số $p=p(x)$. Theo giả thiết tốc độ của $x$ tỉ lệ với tốc độ thay đổi của $p$ nên hàm số $p=p(x)$ là hàm số bậc nhất.\\
		Giá thuê phòng $x_1=2000$ ứng với $p_1=20$ và $x_2=2200$ ứng với $p_2=19$.\\
		Do đó, phương trình đường thẳng $p=ax+b$ đi qua hai điểm $(2000;20)$ và $(2200;19)$.\\ Khi đó $p=-\dfrac{1}{200}x+30\Rightarrow x=-200p+6000$.\\
		Hàm doanh thu từ tiền cho thuê phòng là $R(p)=px=p(-200p+6000)=-200p^2+6000p$.\\
		Ta cần tìm $p$ sao cho $R$ đạt giá trị lớn nhất:\\
		$R'(p)=-400p+6000=0\Rightarrow p=15$.\\
		Bảng biến thiên
		\begin{center}
			\begin{tikzpicture}
				\tkzTabInit[nocadre=false,lgt=1.2,espcl=2.5,deltacl=0.6] 
				{$p$/0.6, $R’(p)$/0.6, $R(p)$/2}
				{$0$, $15$,$40$ }
				\tkzTabLine{,+,0,-,}
				\tkzTabVar{-/ , +/ $45000$ , -/ } 
			\end{tikzpicture}
		\end{center}	
		Vậy với $15$ căn hộ cho thuê với mức giá $3$ triệu đồng/$1$ tháng thì công ty thu về lợi nhuận lớn nhất là $45$ triệu đồng/$1$ tháng.
	}
\end{bt}

\begin{bt}%[2D1H3-6]
% nguon: Archive/2024/TaiLieuGiangDay/1.HamSo/1. K12 - 4in1 - Hamso/ThubaiGoc/Bai3-Dang7.tex #2088
[SGK12 - Cánh Diều, Mức độ 2]%[BG12-4IN1, Nguyễn Khánh Trọng]
	Trong $5$ giây đầu tiên, một chất điểm chuyển động theo phương trình $$s(t)=-t^3+6t^2+t+5,$$ trong đó $t$ tính bằng giây và $s$ tính bằng mét. Chất điểm có vận tốc tức thời lớn nhất bằng bao nhiêu trong $5$ giây đầu tiên đó?
	\loigiai{
		Xét hàm số $s(t)=-t^3+6t^2+t+5$, trên đoạn $\left[0;5\right]$.\\
		Vận tốc tức thời $v(t)=s'(t)=-3t^2+12t+1$.\\
		Đạo hàm $v'(t)=-6t+12$. Cho $v'(t)=0 \Leftrightarrow -6t+12=0 \Leftrightarrow t=2\in[0;5]$.\\
		Các giá trị $v(0)=1$, $v(2)=9$, $v(5)=-24$.\\
		So sánh các giá trị, ta có $\max\limits_{[0;5]}v(t)=9$, $\min\limits_{[0;5]}v(t)=-24$.
	}
\end{bt}

\begin{bt}%[2D1H3-6]
% nguon: Archive/2024/TaiLieuGiangDay/1.HamSo/1. K12 - 4in1 - Hamso/ThubaiGoc/Bai3-Dang7.tex #2089
[Mức độ 2]%[BG12-4IN1, Nguyễn Khánh Trọng]
	Sau khi phát hiện một bệnh dịch, các chuyên gia y tế ước tính số người nhiễm bệnh kể từ ngày xuất hiện bệnh nhân đầu tiên đến ngày thứ $t$ là $f(t)=45t^2-t^3$. Nếu xem $f'(t)$ là tốc độ truyền bệnh (người/ngày) tại thời điểm $t$. Tốc độ truyền bệnh sẽ lớn nhất vào ngày thứ mấy?
	\loigiai{
		Ta có $f'(t)=90t-3t^2$. Cần tính giá trị lớn nhất của hàm số $g(t)=f'(t)$.\\
		Khi đó: $g'(t)=f''(t)=90-6t$; $g'(t)=0\Leftrightarrow t=15$.\\
		Bảng biến thiên.
		\begin{center}
			\begin{tikzpicture}
				\tkzTabInit[nocadre=false,lgt=1.2,espcl=2.5,deltacl=0.6]
				{$t$ /0.6,$g'(t)$ /0.6,$g(t)$ /2}
				{$0$,$15$,$+\infty$}
				\tkzTabLine{,+,$0$,-,}
				\tkzTabVar{-/$0$, +/$675$,-/$0$}
			\end{tikzpicture}
		\end{center}
		Vậy tốc độ truyền bệnh lớn nhất vào ngày thứ $15$.
	}
\end{bt}

\begin{bt}%[2D1H3-6]
% nguon: Archive/2024/TaiLieuGiangDay/1.HamSo/1. K12 - 4in1 - Hamso/ThubaiGoc/Bai3-Dang7.tex #2090
[SGK 12 - Cùng Khám Phá, Mức độ 2]%[BG12-4IN1, Nguyễn Khánh Trọng]
	Trong các hình chữ nhật có cùng chu vi bằng $18 \mathrm{~cm}$, hãy tìm hình chữ nhật có diện tích lớn nhất.
	\loigiai{
		Gọi chiều rộng của hình chữ nhật là $x$ cm, với $x \in \left(0;\dfrac{9}{2}\right]$.\\
		Khi đó chiều dài của hình chữ nhật là $9-x$ cm.\\
		Diện tích của hình chữ nhật là $f(x)=x(9-x)=-x^2+9x$.\\
		Ta có $f^{\prime}(x)=-2x+9$ và $f^{\prime}(x)=0  \Leftrightarrow x = \dfrac{9}{2}$.\\
		Bảng biến thiên
		\begin{center}
			\begin{tikzpicture}
				\tkzTabInit[nocadre=false,lgt=1.2,espcl=3.5,deltacl=0.6]
				{$x$/.7,$f'(x)$/.7,$f(x)$/2}
				{$0$,	$4{,}5$}
				\tkzTabLine{	,+,		}
				\tkzTabVar{-/$0$,	+/$\tfrac{81}{4}$}
			\end{tikzpicture}
		\end{center}
		Vậy hình chữ nhật có diện tích lớn nhất là hình vuông có cạnh $x = \dfrac{9}{2}$.
	}
\end{bt}

\begin{bt}%[2D1H3-6]
% nguon: Archive/2024/TaiLieuGiangDay/1.HamSo/1. K12 - 4in1 - Hamso/ThubaiGoc/Bai3-Dang7.tex #2091
[SGK 12 - KNTT, Mức độ 2]%[BG12-4IN1, Nguyễn Khánh Trọng]
	\immini{
		Từ một tấm bìa carton hình vuông có độ dài cạnh bằng $60$ cm, người ta cắt bốn hình vuông bằng nhau ở bốn góc rồi gập thành một chiếc hộp có dạng hình hộp chữ nhật không có nắp (H.1.14). Tính cạnh của các hình vuông bị cắt sao cho thể tích của chiếc hộp là lớn nhất.	
	}{\begin{tabular}{ccc}
			\rule[-1ex]{0pt}{2.5ex}	\begin{tikzpicture}[>=stealth,line join=round,line cap=round,font=\footnotesize,scale=.6]
				\filldraw[pattern=north east lines] 
				(0,0)rectangle(1,1)(0,4)rectangle(1,5)
				(4,0)rectangle(5,1)(4,4)rectangle(5,5)
				;
				\draw[dashed](1,1)rectangle(4,4);
				\draw(0,0)rectangle(5,5);
				\draw[dashed](5,5)arc(45:-45:{5*1.41/2}) node[midway,fill=white]{$60$ cm};
			\end{tikzpicture}& & \begin{tikzpicture}[>=stealth,line join=round,line cap=round,font=\footnotesize,scale=.4]	
				\tikzset{
					pics/hhChuNhat/.style n args={8}{
						code={
							\tikzset{
								declare function={a=2.5;b=1.8;goc=-135;h=0.75;}
							}
							\path 
							(0,0)coordinate(#1)+(0:a)coordinate(#2)+(goc:b)coordinate(#4)+(90:h)coordinate(#5)
							($(#2)+(#4)-(#1)$)coordinate(#3)
							;
							\foreach \pone/\pname in{#2/#6,#3/#7,#4/#8}{
								\path 
								($(\pone)+(#5)-(#1)$)coordinate(\pname)
								;
							}
							\foreach \pointo/\pointt in{#2/#3,#3/#4,#5/#6,#6/#7,#7/#8,#8/#5,#2/#6,#3/#7,#4/#8}{
								\draw[fill=black](\pointo)--(\pointt);
							}
						}
					}
				}
				\path 
				(0,0) pic{hhChuNhat={A}{B}{C}{D}{A'}{B'}{C'}{D'}}
				(intersection of A--D and C'--D')coordinate(M)
				(intersection of A--B and B'--C')coordinate(N)
				;
				\foreach \pointo/\pointt in{A/M,A/N,A/A'}{
					\draw[fill=black](\pointo)--(\pointt);
				}
				\foreach \pointo/\pointt in{D/M,B/N}{
					\draw[fill=black,dashed](\pointo)--(\pointt);
				}
			\end{tikzpicture}\\
			\multicolumn{3}{c}{\color{orange}Hình 1.14}\\
		\end{tabular}
	}	
	\loigiai{
		Gọi $x$ (cm) là độ dài cạnh của các hình vuông nhỏ được cắt ở bốn góc của tấm bìa. Điều kiện $0<x<30$.\\
		Khi cắt bỏ bốn hình vuông nhỏ có cạnh $x$ cm ở bốn góc và gập lên thì ta được một chiếc hộp chữ nhật không có nắp, có đáy là hình vuông với độ dài cạnh bằng $(60-2 x)$ (cm) và chiều cao bằng $x$ cm. Thể tích của chiếc hộp này là
		\[
		V(x)=(60-2 x)^2 \cdot x=4 x^3-240 x^2+3600 x\left(\text{cm}^3\right).
		\]
		Ta có $V'(x)=12 x^2-480 x+3600$;  $V'(x)=0 \Leftrightarrow x^2-40 x+300=0 \Leftrightarrow x=10$ (thoả mãn điều kiện) hoặc $x=30$ (loại).\\
		Lập bảng biến thiên:
		\begin{center}
			\begin{tikzpicture}[>=stealth,line join=round,line cap=round,font=\footnotesize,scale=1]
				\tkzTabInit[nocadre=false,lgt=1.2,espcl=3,deltacl=.55]
				{$x$/0.7, $V'(x)$/0.7, $V(x)$/1.8}
				{$0$,$10$,$30$}
				\tkzTabLine{,+,$0$,-,}
				\tkzTabVar{-/$0$ ,+/ $16\,000$ ,-/$0$}	
			\end{tikzpicture}
		\end{center}
		Vậy để thể tích của chiếc hộp là lớn nhất thì độ dài cạnh của các hình vuông nhỏ phải cắt là $10$ cm.		
	}
\end{bt}

\begin{bt}%[2D1V3-6]
% nguon: Archive/2024/TaiLieuGiangDay/1.HamSo/1. K12 - 4in1 - Hamso/ThubaiGoc/Bai3-Dang7.tex #2092
[Mức độ 3]%[BG12-4IN1, Nguyễn Khánh Trọng]
	Khi xây nhà, chủ nhà cần làm một bồn nước bằng gạch và xi măng có dạng hình hộp đứng đáy là hình chữ nhật có chiều rộng là $x$ (m), chiều dài gấp 2 lần chiều rộng và không nắp, có chiều cao là $h$ (m), có thể tích là $\dfrac{4}{3}$ m$^3$. Tìm chiều rộng của đáy hình chữ nhật để chi phí xây dựng là thấp nhất. 
	\loigiai{
		Thể tích của bồn nước là $\dfrac{4}{3}$ m$^3$ nên $2x^2 h=\dfrac{4}{3} \Leftrightarrow h=\dfrac{2}{3x^2}$.\\
		Diện tích xung quanh của bồn nước (không nắp) là $$S=2(xh+2xh)+2x^2=6xh+2x^2=\dfrac{4}{x}+2x^2.$$
		Ta có $S'(x)=-\dfrac{4}{x^2}+4x$; $S'(x)=0\Leftrightarrow x=1$.\\
		Bảng biến thiên
		\begin{center}
			\begin{tikzpicture}
				\tkzTabInit[nocadre=false,lgt=1.2,espcl=2.5,deltacl=0.6]
				{$x$ /0.6,$S'(x)$ /0.6,$S(x)$ /2}
				{$0$,$1$,$+\infty$}
				\tkzTabLine{,-,$0$,+,}
				\tkzTabVar{+/$+\infty$, -/$6$,+/$+\infty$}
			\end{tikzpicture}
		\end{center}
		Vậy với chiều rộng là $1$ (m) thì chi phí xây dựng là thấp nhất.
	}
\end{bt}

\begin{bt}%[2D1V3-6]
% nguon: Archive/2024/TaiLieuGiangDay/1.HamSo/1. K12 - 4in1 - Hamso/ThubaiGoc/Bai3-Dang7.tex #2093
[Mức độ 3]%[BG12-4IN1, Nguyễn Khánh Trọng]
	\immini{
		Một công ty muốn làm một đường ống dẫn từ một điểm $A$ trên bờ đến một điểm $B$ trên một hòn đảo. Hòn đảo cách bờ biển $6$ km. Giá để xây đường ống trên bờ là $50\,000$ USD mỗi km, và $130\,000$ USD mỗi km để xây dưới nước. Gọi $B'$ là điểm trên bờ biển sao cho $BB'$ vuông góc với bờ biển. Khoảng cách từ $A$ đến $B'$ là $9$ km. Vị trí $C$ trên đoạn $AB'$ sao cho khi nối ống theo $ACB$ thì số tiền ít nhất. Khi đó khoảng cách từ $C$ đến $A$ là bao nhiêu?
	}{
		\begin{tikzpicture}[scale=0.7, font=\footnotesize, line join=round, line cap=round, >=stealth]
			\tkzDefPoints{0/0/B',5/0/A,2/0/C,0/3/B}
			\tkzDrawSegments(A,C C,B)
			\tkzDrawSegments[dashed](C,B' B',B)
			\tkzDrawPoints[fill=black,size=2](A,B,B',C)
			\tkzLabelPoints[below](A,C,B')
			\tkzLabelPoints[above](B)
			\tkzLabelSegment[left](B',B){$6$ km}
			\tkzLabelSegment[below](A,C){$9$ km}
		\end{tikzpicture}
	}
	\loigiai{
		Đặt $B'C=x$ (km), $x\in[0;9]$.\\
		$BC=\sqrt{x^2+6^2}$; $AC=9-x$.\\
		Chi phí xây dựng đường ống là $f(x)=130
		,000\sqrt{x^2+6^2}+50\,000(9-x)$ (USD).\\
		Ta tìm $x\in[0;9]$ sao cho $f(x)$ nhỏ nhất.\\
		Ta có $f'(x)=10\,000\cdot\left(\dfrac{13x}{\sqrt{x^2+36}}-5\right)$.\\
		$f'(x)=0\Leftrightarrow 13x=5\sqrt{x^2+36}\Leftrightarrow 169x^2=25\left(x^2+36\right)\Leftrightarrow x^2=\dfrac{25}{4}\Leftrightarrow x=\dfrac{5}{2}$.\\
		$f(0)=1\,230\,000$; $f\left(\dfrac{5}{2}\right)=1\,170\,000$; $f(9)\approx 1\,406\,165$.\\
		Vậy chi phí thấp nhất khi $x=2{,}5$. Vậy $C$ cần cách $A$ một khoảng $6{,}5$ km.
	}
\end{bt}

\begin{bt}%[2D1V3-6]
% nguon: Archive/2024/TaiLieuGiangDay/1.HamSo/1. K12 - 4in1 - Hamso/ThubaiGoc/Bai3-Dang7.tex #2094
[Mức độ 3]%[BG12-4IN1, Nguyễn Khánh Trọng]
	Một công ty bất động sản có $50$ căn hộ cho thuê. Biết rằng nếu cho thuê mỗi căn hộ với giá $2000000$ đồng một tháng thì mọi căn hộ đều có người thuê và cứ mỗi lần tăng giá cho thuê mỗi căn hộ thêm $50000$ đồng một tháng thì có thêm một căn hộ bị bỏ trống. Công ty đã tìm ra phương án cho thuê đạt lợi nhuận lớn nhất. Hỏi thu nhập nhiều nhất công ty có thể đạt được trong một tháng là bao nhiêu?
	\loigiai{
		Gọi $n$ là số lần tăng giá ($n$ là số tự nhiên). \\
		Khi đó số căn hộ bị bỏ trống cũng là $n$. \\
		Do đó số tiền thu được khi cho thuê $50-n$ căn hộ là \begin{center}
			$A=\left({2\cdot 10}^6+{5\cdot 10}^4 n\right)(50-n)=-5\cdot 10^4n^2+5\cdot 10^5n+10^8$, với $n<50$.
		\end{center}
		Xét hàm số $f(x)=-5\cdot 10^4x^2+5\cdot 10^5x+10^8$, với $0\leq x<50$.\\
		Ta có $f'(x)=-10^5x+5\cdot 10^5$; $f'(x)=0\Leftrightarrow x=5$.\\
		Bảng biến thiên
		\begin{center}
			\begin{tikzpicture}
				\tkzTabInit[nocadre=false,lgt=1.2,espcl=2.5,deltacl=0.6]
				{$x$ /0.6,$f'(x)$ /0.6,$f(x)$ /2}
				{$0$,$5$,$50$}
				\tkzTabLine{,+,$0$,-,}
				\tkzTabVar{-/$f(0)$, +/$f(5)$,-/$f(50)$}
			\end{tikzpicture}
		\end{center}
		Vậy $\max\limits_{[0;50]} f(x)=f(5)=101250000$.\\
		Vậy thu nhập cao nhất công ty có thể đạt được trong một tháng là $101250000$ đồng.
	}
\end{bt}

\begin{bt}%[2D1H3-6]
% nguon: Archive/2024/TaiLieuGiangDay/1.HamSo/1. K12 - 4in1 - Hamso/ThubaiGoc/Bai3-Dang7.tex #2098
[Mức độ 2]%[BG12-4IN1, Nguyễn Khánh Trọng]
	Một vật chuyển động theo quy luật $s=-\dfrac{1}{2}t^3+9t^2$ với $t$ (giây) là khoảng thời gian tính từ lúc bắt đầu chuyển động và $s$ (mét) là quãng đường vật đi được trong khoảng thời gian đó. Hỏi trong khoảng thời gian $10$ giây, kể từ lúc bắt đầu chuyển động, vận tốc lớn nhất của vật đạt được bằng bao nhiêu?
	\dapso{$54$ m/s}
	\loigiai{
		Vận tốc tại thời điểm $t$ là $v(t)=s'(t)=-\dfrac{3}{2}t^2+18t$ với $t\in [0;10]$.\\
		Ta có $v(t)=-\dfrac{3}{2}(t-6)^2+54\le 54$.\\
		Vậy $\max\limits_{[0;10]} v(t)=v(6)=54$ m/s.
	}
\end{bt}

\begin{bt}%[2D1H3-6]
% nguon: Archive/2024/TaiLieuGiangDay/1.HamSo/1. K12 - 4in1 - Hamso/ThubaiGoc/Bai3-Dang7.tex #2099
[SGK 12 - KNTT, Mức độ 2]%[BG12-4IN1, Nguyễn Khánh Trọng]
	Trong các hình chữ nhật có chu vi là $24\mathrm{~cm}$, hãy tìm hình chữ nhật có diện tích lớn nhất.
	\loigiai{
		Gọi $x$ (cm) là chiều rộng của hình chữ nhật.\\
		Điều kiện $0<x<12$.\\
		Chiều dài của hình chữ nhật là $12-x$ (cm).\\
		Diện tích của hình chữ nhật là $x(12-x)=-x^2+12x$ (cm$^2$).\\
		Xét hàm số $y=-x^2+12x$ với $0<x<12$.\\
		Ta có $y'=-2x+12$; $y'=0\Leftrightarrow x=6$.\\
		\[
		y(0)=0; y(6)=36; y(12)=0.
		\]
		Do đó: $\max \limits_{(0;12)} y=y(6)=36$.\\
		Trong các hình chữ nhật có chu vi là $24\mathrm{~cm}$ thì hình vuông có diện tích lớn nhất.
	}
\end{bt}

\begin{bt}%[2D1H3-6]
% nguon: Archive/2024/TaiLieuGiangDay/1.HamSo/1. K12 - 4in1 - Hamso/ThubaiGoc/Bai3-Dang7.tex #2100
[SGK 12 - Cùng Khám Phá, Mức độ 2]%[BG12-4IN1, Nguyễn Khánh Trọng]
	Trong một trò chơi, mỗi đội được phát một tấm bìa hình vuông có cạnh bằng $30$~cm. Nhiệm vụ của mỗi đội chơi là cắt ở bốn góc của tấm bìa này bốn hình vuông bằng nhau, rồi gấp tầm bìa lại (Hình 1.6) và dán keo để được một cái hộp không nắp có dạng hình hộp chữ nhật. Đội nào thiết kế được một cái hộp có thể tích lớn nhất sẽ dành chiến thắng. Hãy xác định cạnh của các hình vuông bị cắt để thu được hộp có thể tích lớn nhất.
	\dapso{$5$ cm}
	\loigiai{
		\begin{center}
			\def\kyhieu(#1,#2)(#3){
				\path (#1)--(#2)--([turn]-90:9pt) coordinate (X)
				($ (X)+(#1)-(#2) $)  coordinate (Y);
				\draw[dashed] (#1)--(Y) (#2)--(X);
				\draw[|<->|] (X)--(Y) node[midway,sloped,fill=white,inner sep=1pt]{\scriptsize #3};
			}
			\begin{tikzpicture}[scale=.6, line join = round, line cap = round]
				\tikzset{label style/.style={font=\footnotesize}}
				\def\a{4}
				\path 	(0:0) coordinate (A)
				++(0:\a) coordinate (B)
				++(90:\a) coordinate (C)
				($(A)-(B)+(C)$) coordinate (D)
				($(A)+(45:.35*\a)$) coordinate (A')
				($(B)+(135:.35*\a)$) coordinate (B')
				($(C)+(-135:.35*\a)$) coordinate (C')
				($(D)+(-45:.35*\a)$) coordinate (D')
				($(A)+(0:.35*\a)$) coordinate (E);
				\fill[pattern=dots] 	(A')--(B')--(C')--(D')--cycle;
				\draw[dashed] 	(A')--(B')--(C')--(D')--cycle;
				\draw[thick] 	(A)--(B)--(C)--(D)--cycle;
				\draw[pattern=north west lines] (A) rectangle (A')	(B) rectangle (B')	(C) rectangle (C')	(D) rectangle (D');
				\kyhieu(D,A)($30$ cm);
				\draw[draw=none] 	(A)--(0:.2*\a) node[pos=0.5,sloped,below]{$x$}; 
				\draw[draw=none] (A)--(B) node[pos=0.5,sloped,below=12pt]{Hình 1.6}; 
			\end{tikzpicture}\hspace{3cm}
			\begin{tikzpicture}[scale=.8]
				\def\a{3}	
				\def\b{1.65}
				\def\h{1.5}
				\path 	(0:0) coordinate (A)
				++(0:\a) coordinate (D)
				++(-150:\b) coordinate (C)
				($(A)+(C)-(D)$) coordinate (B)
				($(A)+(90:\h)$) coordinate (A')
				($(B)+(90:\h)$) coordinate (B')
				($(C)+(90:\h)$) coordinate (C')
				($(D)+(90:\h)$) coordinate (D');
				\fill[pattern=dots] 	(A)--(B)--(C)--(D)--cycle;
				\draw[dashed,thick] 	(B)--(A)--(D)	(A)--(A');
				\draw[thick] 	(C)--(C') 	(D)--(D') 	(B)--(B')	(C)--(C')
				(B)--(C)--(D) 
				(A')--(B')--(C')--(D')--cycle;
			\end{tikzpicture}
		\end{center}
		\loigiai{
			Gọi $x$ (cm) chiều dài cạnh hình vuông bị cắt.\\
			Khi đó, hình hộp được tao thành có các kích thước $x\times (30-2x)\times (30-2x)$, với $0<x<15$.\\
			Thể tích của khối hộp là $V(x)=x(30-2x)^2=4x^3-120x^2+900x$.\\
			Xét hàm số $V=4x^3-120x^2+900x$ trên khoảng $(0;\,15)$ ta có 
			\begin{itemize}
				\item $V'=12x^2-240x+900$.
				\item $V'=0\Leftrightarrow 12x^2-240x+900=0\Leftrightarrow x=5$ (do $x\in (0;\,15)$ ).
			\end{itemize}
			Bảng biến thiên
			\begin{center}
				\begin{tikzpicture}[>=stealth]
					\tkzTabInit[nocadre=false,lgt=1.5,espcl=3,deltacl=0.6]{$x$/.6 ,$V'(x)$/.7,$V(x)$/2}
					{$0$ , $5$ , $15$}
					\tkzTabLine{ ,+ , $0$ , - , }
					\tkzTabVar{-/$0$ , +/$2000$ , -/$0$}
				\end{tikzpicture}
			\end{center}
			Trên khoảng $(0;15)$, hàm số $V(x)$ đạt giá trị lớn nhất bằng $2\,000$ tại $x=5$.\\
			Vậy ta phải cắt các hình vuông có cạnh bằng $x=5$ (cm) thì thu được khối hộp có thể tích lớn nhất.
		}
	}
\end{bt}

\begin{bt}%[2D1H3-6]
% nguon: Archive/2024/TaiLieuGiangDay/1.HamSo/1. K12 - 4in1 - Hamso/ThubaiGoc/Bai3-Dang7.tex #2101
[SGK 12 - CTST, Mức độ 2]%[BG12-4IN1, Nguyễn Khánh Trọng]
	Tam giác vuông có cạnh huyền bằng $5 \mathrm{~cm}$ có thể có diện tích lớn nhất bằng bao nhiêu?
	\dapso{$\dfrac{625}{4}$}
	\loigiai{Gọi một cạnh góc vuông là $x$ ($0<x<5$) thì cạnh góc vuông còn lại là $\sqrt{25-x^2}$.\\ Như vậy, diện tích tam giác là $S=\dfrac{x\cdot\sqrt{25-x^2}}{2}$.
		Đặt $f(x)=25x^2-x^4$. 
		\\Ta có $f'(x)=50x-4x^3$. Khi đó
		$f'(x)=0 \Leftrightarrow x=\dfrac{5\sqrt{2}}{2}$.\\
		Vì vậy $\displaystyle\max _{(0;5)} f(x)=f\left( \dfrac{5\sqrt{2}}{2}\right) =\dfrac{625}{4}$.\\
		Vậy tam giác vuông có cạnh huyền bằng $5 \mathrm{~cm}$ có thể có diện tích lớn nhất bằng $\dfrac{625}{4}$.}
\end{bt}

\begin{bt}%[2D1V3-6]
% nguon: Archive/2024/TaiLieuGiangDay/1.HamSo/1. K12 - 4in1 - Hamso/ThubaiGoc/Bai3-Dang7.tex #2102
[SGK 12 - KNTT, Mức độ 3]%[BG12-4IN1, Nguyễn Khánh Trọng]
	\immini{
		Một nhà sản xuất muốn thiết kế một chiếc hộp có dạng hình hộp chữ nhật không có nắp, có đáy là hình vuông và diện tích bề mặt bằng $108\mathrm{~cm}^2$ như Hình 1.17. Tìm các kích thước của chiếc hộp sao cho thể tích của hộp là lớn nhất.
		\dapso{$x=6$ cm và $h=3$ cm}
	}{
		\begin{tikzpicture}[>=stealth,line join=round,line cap=round,font=\footnotesize,scale=.4]	
			\tikzset{
				pics/hhChuNhat/.style n args={8}{
					code={
						\tikzset{
							declare function={a=2.5;b=1.8;goc=-135;h=0.75;}
						}
						\path 
						(0,0)coordinate(#1)+(0:a)coordinate(#2)+(goc:b)coordinate(#4)+(90:h)coordinate(#5)
						($(#2)+(#4)-(#1)$)coordinate(#3)
						;
						\foreach \pone/\pname in{#2/#6,#3/#7,#4/#8}{
							\path 
							($(\pone)+(#5)-(#1)$)coordinate(\pname)
							;
						}
						\foreach \pointo/\pointt in{#2/#3,#3/#4,#5/#6,#6/#7,#7/#8,#8/#5,#2/#6,#3/#7,#4/#8}{
							\draw[fill=black](\pointo)--(\pointt);
						}
					}
				}
			}
			\path 
			(0,0) pic{hhChuNhat={A}{B}{C}{D}{A'}{B'}{C'}{D'}}
			(intersection of A--D and C'--D')coordinate(M)
			(intersection of A--B and B'--C')coordinate(N)
			;
			\foreach \pointo/\pointt in{A/M,A/N,A/A'}{
				\draw[fill=black](\pointo)--(\pointt);
			}
			\foreach \pointo/\pointt in{D/M,B/N}{
				\draw[fill=black,dashed](\pointo)--(\pointt);
			}
			\path (C)--(D)node[midway, below]{$x$};
			\path (C)--(B)node[midway, right]{$x$};
			\path (B)--(B')node[midway, right]{$h$};
			\draw (current bounding box.south) node[below]{Hình 1.17};
		\end{tikzpicture}
	}
	\loigiai{
		Thể tích của chiếc hộp là $V=x^2\cdot h$ (cm$^3$).\\
		Vì diện tích bề mặt bằng $108$ cm$^2$ nên ta có 
		\begin{center}
			$x^2+4xh=108\Leftrightarrow h=\dfrac{108-x^2}{4x}$ (điều kiện $0< x < \sqrt{108}$).
		\end{center}
		Khi đó thể tích chiếc hộp là $V=\dfrac{x\left(108-x^2\right)}{4}=-\dfrac{x^3}{4}+27x$.\\
		Ta tìm giá trị lớn nhất của hàm số $V=-\dfrac{x^3}{4}+27 x~(0< x < \sqrt{108})$.\\
		Có $V^{\prime}=-\dfrac{3}{4} x^2+27; V^{\prime}=0\Leftrightarrow-\dfrac{3}{4} x^2+27=0\Leftrightarrow x=6$ (vì $0< x < \sqrt{108}$).\\
		Bảng biến thiên:
		\begin{center}
			\begin{tikzpicture}[scale=1, font=\footnotesize]
				\tkzTabInit[nocadre=false, lgt=1.2, espcl=2, deltacl=0.9]
				{$x$/0.8,$V'(x)$/0.6,$V(x)$/2}
				{$0$,$6$,$\sqrt{108}$};
				\tkzTabLine{,+,$0$,-,};
				\tkzTabVar{-/$0$,+/$108$,-/$V\left(\sqrt{108}\right)$};
			\end{tikzpicture}
		\end{center}
		Dựa vào bảng biến thiên, ta có thể tích lớn nhất của chiếc hộp là $108$ cm$^3$ khi $x=6$ cm và $h=3$ cm.
	}
\end{bt}

\begin{bt}%[2D1H5-1]
% nguon: Archive/2024/TaiLieuGiangDay/1.HamSo/1. K12 - 4in1 - Hamso/ThubaiGoc/Bai5-Dang2.tex #2113
%[Dự án Bài Giảng 12 Nhóm Toán & LaTex, Lê Minh Thiện Anh]
	Khảo sát sự biến thiên và vẽ đồ thị các hàm số sau:
	\begin{listEX}[2]
		\item $y=\dfrac{x-2}{2x+1}$;
		\item $y=\dfrac{1-2x}{2x+4}$.
	\end{listEX}
	\loigiai{
		\begin{enumerate}
			\item $y=\dfrac{x-2}{2x+1}$
				\begin{itemize}
					\item Tập xác định $\mathscr{D}=\mathbb{R}\backslash \left\{ -\dfrac{1}{2} \right\}$.
					\item Đạo hàm $y'=\dfrac{5}{(2x+1)^2}>0, \forall x\ne -\dfrac{1}{2}$.\\ 
					Hàm số đồng biến trên các khoảng $\left(-\infty;-\dfrac{1}{2}\right)$ và $\left(-\dfrac{1}{2};+\infty\right)$.
					\item Giới hạn \& tiệm cận.\\
					+ $ \lim\limits_{x \to {-\tfrac{1}{2}^-}}y=+\infty; \lim\limits_{x \to {-\tfrac{1}{2}^+}}y=-\infty$.
					Đồ thị hàm số có tiệm cận đứng $x=-\dfrac{1}{2}$.\\
					+$\lim\limits_{x \to \pm \infty}y=\dfrac{1}{2}$. Đồ thị hàm số có tiệm cận ngang $y=\dfrac{1}{2}$.
					\item Bảng biến thiên
					\begin{center}
						\begin{tikzpicture}[>=stealth,line join=round,line cap=round,font=\footnotesize,scale=1.3]
							%========================
							\tkzTabInit[nocadre=false,lgt=1.2,espcl=2.5,deltacl=0.6]
							{$x$ /0.6,$y’$ /0.6,$y$ /2}
							{$-\infty$,$-\dfrac{1}{2}$,$+\infty$}
							\tkzTabLine{,+,d,+,}
							\tkzTabVar{-/$\dfrac{1}{2}$,+D-/$+\infty$/$-\infty$,+/$\dfrac{1}{2}$}
							%========================
						\end{tikzpicture}
					\end{center}
					\item Đồ thị\\
					\immini{
					\begin{itemize}
						\item Vẽ các đường tiệm cận $x=-\dfrac{1}{2}$ và $y=\dfrac{1}{2}$.\\
						\item Đồ thị hàm số giao với trục $Oy$ tại điểm $(0;-2)$ và giao với trục $Ox$ tại điểm $(2;0)$.\\
						Đồ thị hàm số nhận giao điểm $I\left(-\dfrac{1}{2};\dfrac{1}{2}\right)$ của hai đường tiệp cận của đồ thị làm tâm đối xứng và nhận hai đường phân giác của các góc tạo bởi hai đường tiệm cận đó làm trục đối xứng.
					\end{itemize}}
					 {
					 	\begin{tikzpicture}[>=stealth,line join=round,line cap=round,font=\footnotesize,scale=.7]
							%========================
							\def\hsf(#1){(1*(#1)-2)/(2*(#1)+1)}
							\draw[->](-5,0)--(5,0);
							\draw[->](0,-5)--(0,5);
							\draw (5,0) node[above]{$x$} (0,5) node[left]{$y$} (0,0) node[below right]{$O$}; 
							\draw[red] (-0.5,-5)--(-0.5,5) (-5,0.5)--(5,0.5); 
							\clip (-5,-5)rectangle(5,5);
							\draw[thick,samples= 100,smooth,domain=-5:-0.6,blue] plot(\x,{\hsf(\x)}); 
							\draw[thick,samples= 100,smooth,domain=-0.4:5,blue] plot(\x,{\hsf(\x)});
							
							\foreach \x in {}{
								\fill (\x,0) node[below]{$\x$} circle(1pt);}
							\foreach \y in {}{
								\fill (0,\y) node[left]{$\y$} circle(1pt);}
							\draw(-0.5,0) circle (.02) node[below left]{$-\tfrac{1}{2}$};
							\draw(0,0.5) circle (.02) node[above right]{$\tfrac{1}{2}$};
							\draw (.25,0)|-(0,.25);\fill (.25,.25) circle(0pt);
							\node[right] at (0,-2){$-2$};
							\node[below] at (2,0){$2$};
							%========================
					\end{tikzpicture}
				}
				\end{itemize}
				
				\item $y=\dfrac{1-2x}{2x+4}$
				\begin{itemize}
					\item Tập xác định $\mathscr{D}=\mathbb{R}\backslash \left\{ -2 \right\}$.
					\item Đạo hàm $y'=\dfrac{8}{(2x+4)^2}>0, \forall x\ne -2$.\\ 
					Hàm số đồng biến trên các khoảng $\left(-\infty;-2\right)$ và $\left(-2;+\infty\right)$.
					\item Giới hạn \& tiệm cận.\\
					+ $ \lim\limits_{x \to {-2^-}}y=+\infty; \lim\limits_{x \to {-2^+}}y=-\infty$.
					Đồ thị hàm số có tiệm cận đứng $x=-2$.\\
					+$\lim\limits_{x \to \pm \infty}y=-1$. Đồ thị hàm số có tiệm cận ngang $y=-1$.
					\item Bảng biến thiên
						\begin{center}
							\begin{tikzpicture}[>=stealth,line join=round,line cap=round,font=\footnotesize,scale=1.3]
							%========================
							\tkzTabInit[nocadre=false,lgt=1.2,espcl=2.5,deltacl=0.6]
							{$x$ /0.6,$y’$ /0.6,$y$ /2}
							{$-\infty$,$-2$,$+\infty$}
							\tkzTabLine{,+,d,+,}
							\tkzTabVar{-/$-1$,+D-/$+\infty$/$-\infty$,+/$-1$}
							%========================
						\end{tikzpicture}
						\end{center}
					\item Đồ thị\\
					\immini{
						\begin{itemize}
							\item Vẽ các đường tiệm cận $x=-2$ và $y=-1$.\\
							\item Đồ thị hàm số giao với trục $Oy$ tại điểm $\left(\dfrac{1}{2};0\right)$ và giao với trục $Ox$ tại điểm $\left(0;\dfrac{1}{4}\right)$.\\
							Đồ thị hàm số nhận giao điểm $I\left(-2;-1\right)$ của hai đường tiệp cận của đồ thị làm tâm đối xứng và nhận hai đường phân giác của các góc tạo bởi hai đường tiệm cận đó làm trục đối xứng.
					\end{itemize}}
					{
						\begin{tikzpicture}[>=stealth,line join=round,line cap=round,font=\footnotesize,scale=.7]
							%========================
							\def\hsf(#1){(-2*(#1)+1)/(2*(#1)+4)}
							\draw[->](-7,0)--(3,0);
							\draw[->](0,-6)--(0,4);
							\draw (3,0) node[above]{$x$} (0,4) node[left]{$y$} (0,0) node[below right]{$O$}; 
							\draw[red] (-2,-6)--(-2,4) (-7,-1)--(3,-1); 
							\clip (-7,-6)rectangle(3,4);
							\draw[samples= 100,smooth,domain=-7:-2.1,blue] plot(\x,{\hsf(\x)}); 
							\draw[samples= 100,smooth,domain=-1.9:3,blue] plot(\x,{\hsf(\x)});
							
							\foreach \x in {}{
								\fill (\x,0) node[below]{$\x$} circle(1pt);}
							\foreach \y in {}{
								\fill (0,\y) node[left]{$\y$} circle(1pt);}
							\draw(-2,0) circle (.02) node[below left]{$-2$};
							\draw(0,-1) circle (.02) node[below right]{$-1$};
							\draw (.25,0)|-(0,.25);\fill (.25,.25) circle(0pt);
							%========================
					\end{tikzpicture}
				}
				\end{itemize}
		\end{enumerate}
	}
\end{bt}

\begin{bt}%[2D1H5-1]
% nguon: Archive/2024/TaiLieuGiangDay/1.HamSo/1. K12 - 4in1 - Hamso/ThubaiGoc/Bai5-Dang2.tex #2114
%[Dự án Bài Giảng 12 Nhóm Toán & LaTex, Lê Minh Thiện Anh]
	Khảo sát sự biến thiên và vẽ đồ thị của các hàm số
	\begin{listEX}[2]
		\item $y=\dfrac{2x+1}{x+1}$;
		\item $y=\dfrac{x+3}{1-x}$.
	\end{listEX}
	\loigiai{
		\begin{enumerate}
			\item \begin{enumerate}
				\item Tập xác định $\mathscr{D}=\mathbb{R} \setminus\{-1\}$.		
				\item Sự biến thiên
				\begin{itemize}
					\item Ta có: $y^{\prime}=\dfrac{1}{(x+1)^{2}}>0$ với mọi $x \neq -1$.
					\item Hàm số đồng biến trên từng khoảng $(-\infty ; -1)$ và $(-1 ;+\infty)$.
					\item Hàm số không có cực trị.
					\item Tiệm cận
					$\begin{aligned}[t]
						&\lim _{x \rightarrow (-1)^{-}} y=\lim _{x \rightarrow (-1)^{-}} \dfrac{2x+1}{x+1}=+\infty ; \lim _{x \rightarrow )(-1)^{+}} y=\lim _{x \rightarrow (-1)^{+}} \dfrac{2x+1}{x+1}=-\infty;\\
						&\lim _{x \rightarrow+\infty} y=\lim _{x \rightarrow+\infty} \dfrac{2x+1}{x+1}=2 ; \lim _{x \rightarrow-\infty} y=\lim _{x \rightarrow-\infty} \dfrac{2x+1}{x+1}=2.
					\end{aligned}
					$
					
					Do đó, đồ thị của hàm số có tiệm cận đứng là đường thẳng $x=-1$, tiệm cận ngang là đường thẳng $y=2$.
					\item  Bảng biến thiên:
					\begin{center}
						\begin{tikzpicture}[line join=round,font=\footnotesize,>=stealth, scale=1]
							\tkzTabInit[nocadre=false,lgt=1.2,espcl=3.5,deltacl=0.6]
							{$x$ /0.6,$y'$ /0.6,$y$ /2}
							{$-\infty$, $2$, $+\infty$}
							\tkzTabLine{,-,d,-}
							\tkzTabVar{-/$2$, +D-/$+\infty$/$-\infty$, +/$2$}
						\end{tikzpicture}
					\end{center}
				\end{itemize}
				\item Đồ thị
				\immini{
					\begin{itemize}
						\item Giao điểm của đồ thị hàm số với trục tung là điểm $\left(0 ;1\right)$.
						\item Giao điểm của đồ thị hàm số với trục hoành là điểm $\left(0;-\dfrac{1}{2}\right)$.
						\item Đồ thị hàm số nhận giao điểm $I(-1 ; 2)$ của hai đường tiệm cận làm tâm đối xứng và nhận hai đường phân giác của các góc tạo bởi hai đường tiệm cận này làm trục đối xứng.
					\end{itemize}
				}{
					\begin{tikzpicture}[line cap=round, line join=round,font=\footnotesize,>=stealth, scale=1,x=0.5cm, y=0.5cm]
						\tikzset{label style/.style={font=\footnotesize}}
						\draw[->] (-5,0)--(0,0)node[below right]{$O$}--(3,0) node[below]{$x$};
						\draw[->](0,-2)--(0,6) node[left]{$y$};
						\draw[smooth, samples=100] plot[domain=-5:-1.25](\x, { (2*(\x)+1)/(\x+1) });
						\draw[smooth, samples=100] plot[domain=-0.75:3](\x, { (2*(\x)+1)/(\x+1) });
						\foreach\x in {-1} \draw[fill=black] (\x,0) circle (.5pt) node[above left]{$\x$};
						\foreach\y in {2}  \draw[fill=black] (0,\y) circle (.5pt) node[above right]{$\y$};
						\draw (-1,-2) node[below,rotate=0]{$x=-1$}--(-1,6)
						(-5,2)--(3,2) node[above]{$y=2$};
						
					\end{tikzpicture}
				}
			\end{enumerate}
			\item \begin{enumerate}
				\item Tập xác định của hàm số: $\mathbb{R} \setminus\{1\}$.	
				\item Sự biến thiên
				\begin{itemize}
					\item Ta có: $y'=\dfrac{4}{(x-1)^{2}}>0$ với mọi $x \neq -1$.
					\item Hàm số đồng biến trên từng khoảng $(-\infty ; 1)$ và $(1 ;+\infty)$.
					\item Hàm số không có cực trị.
					\item Tiệm cận
					$\begin{aligned}[t]
						&\lim _{x \rightarrow 1^{-}} y=\lim _{x \rightarrow (-1)^{-}} \dfrac{x+3}{-x+1}=+\infty ; \lim _{x \rightarrow 1^{+}} y=\lim _{x \rightarrow 1^{+}} \dfrac{x+3}{-x+1}=-\infty;\\
						&\lim _{x \rightarrow+\infty} y=\lim _{x \rightarrow+\infty} \dfrac{x+3}{-x+1}=-1 ; \lim _{x \rightarrow-\infty} y=\lim _{x \rightarrow-\infty} \dfrac{x+3}{-x+1}=-1.
					\end{aligned}
					$
					
					Do đó, đồ thị của hàm số có tiệm cận đứng là đường thẳng $x=1$, tiệm cận ngang là đường thẳng $y=-1$.
					\item  Bảng biến thiên:
					\begin{center}
						\begin{tikzpicture}[line join=round,font=\footnotesize,>=stealth, scale=1]
							\tkzTabInit[nocadre=false,lgt=1.2,espcl=3.5,deltacl=0.6]
							{$x$ /0.6,$y'$ /0.6,$y$ /2}
							{$-\infty$, $1$, $+\infty$}
							\tkzTabLine{,-,d,-}
							\tkzTabVar{-/$-1$, +D-/$+\infty$/$-\infty$, +/$-1$}
						\end{tikzpicture}
					\end{center}
				\end{itemize}
				\item Đồ thị
				\immini{
					\begin{itemize}
						\item Giao điểm của đồ thị hàm số với trục hoành là điểm $\left(-3 ;0\right)$.
						\item Giao điểm của đồ thị hàm số với trục tung là điểm $(0;3)$.
						\item Đồ thị hàm số nhận giao điểm $I(1 ; -1)$ của hai đường tiệm cận làm tâm đối xứng và nhận hai đường phân giác của các góc tạo bởi hai đường tiệm cận này làm trục đối xứng.
					\end{itemize}
				}{
					\begin{tikzpicture}[line cap=round, line join=round,font=\footnotesize,>=stealth, scale=1,x=0.5cm, y=0.5cm]
						\tikzset{label style/.style={font=\footnotesize}}
						\draw[->] (-5,0)--(0,0)node[below left]{$O$}--(7,0) node[below]{$x$};
						\draw[->](0,-7)--(0,6) node[left]{$y$};
						\draw[smooth, samples=100] plot[domain=-5:0.4](\x, { ((\x)+3)/(-\x+1) });
						\draw[smooth, samples=100] plot[domain=1.7:7](\x, { ((\x)+3)/(-\x+1) });
						\foreach\x in {-3,1} \draw[fill=black] (\x,0) circle (.5pt) node[above right]{$\x$};
						\foreach\y in {-1,3}  \draw[fill=black] (0,\y) circle (.5pt) node[below right]{$\y$};
						\draw (1,-7) node[above right,rotate=90]{$x=1$}--(1,6)
						(-5,-1)node[below]{$y=-1$}--(7,-1) ;
						
					\end{tikzpicture}
				}
			\end{enumerate}
		\end{enumerate} 
	}
\end{bt}

\begin{bt}%[2D1H5-1]
% nguon: Archive/2024/TaiLieuGiangDay/1.HamSo/1. K12 - 4in1 - Hamso/ThubaiGoc/Bai5-Dang2.tex #2115
%[Dự án Bài Giảng 12 Nhóm Toán & LaTex, Lê Minh Thiện Anh]
	Khảo sát sự biến thiên và vẽ đổ thị của hàm số $y=f(x)=\dfrac{2x + 4}{2x + 1}$.
	\loigiai{
		\begin{enumerate}
			\item Tập xác định: $\mathscr{D}=\mathbb{R}\setminus\{-\frac{1}{2}\}$.
			\item Sự biến thiên:
			\begin{itemize}
				\item Giới hạn, tiệm cận
				\begin{itemize}
					\item $\lim\limits_{x \rightarrow (-\frac{1}{2})^+} f(x)=\lim\limits_{x \rightarrow -(-\frac{1}{2})^+}\dfrac{2x + 4}{2x + 1}=+\infty$; $\lim\limits_{x \rightarrow (-\frac{1}{2})^-} f(x)=\lim\limits_{x \rightarrow (-\frac{1}{2})^-}\dfrac{2x + 4}{2x + 1}= -\infty$.\\
					Suy ra đường thẳng $x = -\frac{1}{2}$ là đường tiệm cận đứng của đồ thị hàm số đã cho.
					\item $\lim\limits_{x \rightarrow+\infty} f(x)=\lim\limits_{x \rightarrow+\infty}\dfrac{2x + 4}{2x + 1}=1; \lim\limits_{x \rightarrow-\infty} f(x)=\lim\limits_{x \rightarrow -\infty}\dfrac{2x + 4}{2x + 1}=1$.\\
					Suy ra đường thẳng $y = 1$ là đường tiệm cận ngang của đồ thị hàm số đã cho.
				\end{itemize}
				Ta có $y'=-\dfrac{6}{(2x + 1)^2}< 0, \forall x \in \mathscr{D}$; $y'$ không xác định tại $x = -\frac{1}{2}$.
				\item Bảng biến thiên
				\begin{center}
					\begin{tikzpicture}[>=stealth]
						\tkzTabInit[nocadre=false,lgt=1,espcl=3.5,deltacl=0.5]{$x$/.6 ,$y'$/.6,$y$/2}
						{$-\infty$ , $-\frac{1}{2}$ , $+\infty$}
						\tkzTabLine{ , - , d , - , }
						\tkzTabVar{+/$1$ , -D+/$-\infty$/$+\infty$ , -/$1$}
					\end{tikzpicture}
				\end{center}
				\item Chiều biến thiên:\\
				Hàm số nghịch biến trên các khoảng $(-\infty ; -\frac{1}{2})$ và $(-\frac{1}{2} ;+\infty)$.\\
				Hàm số không có cực trị.
			\end{itemize}
				\item Đồ thị
				\immini{
				Vẽ các đường tiệm cận $x = -\dfrac{1}{2}$ và $y = 1$.\\
				Đồ thị hàm số giao với trục $Oy$ tại điểm $(0 ; 4)$ và giao với trục $Ox$ tại điểm $(-2;0)$.\\
				Đồ thị hàm số đi qua điểm $(3 ;-2)$, nhận giao điểm $I\left(-\dfrac{1}{2} ;1\right)$ của hai đường tiệm cận của đồ thị làm tâm đối xứng và nhận hai đường phân giác của các góc tạo bởi hai đường tiệm cận đó làm trục đối xứng.}{
				\begin{tikzpicture}[>=stealth,scale=0.7]
					\tikzset{every node/.style={scale=0.7}}
					\draw[->] (-5,0)--(5,0)node[below]{$x$};
					\draw[->] (0,-4)--(0,5)node[right]{$y$};
					\fill (0,0)node[below right]{ $O$} (-0.5,0) (-0.5,1)node[above left]{ $I$};
					\draw[dashed,black,fill=black] (-2,0)circle(1.5pt);
					%\draw[dashed,black,fill=black] (3,-2)circle(1.5pt)edge(3,0)edge(0,-2);
					\draw[black,samples=150,smooth,domain=-5:-0.8] plot(\x,{(2*(\x) + 4)/(2*(\x) + 1)});
					\draw[black,samples=150,smooth,domain=-0.1:5] plot(\x,{(2*(\x) + 4)/(2*(\x) + 1)});
					\draw[red,fill=black] (-5,1)--(5,1)node[below left]{$y = 1$};
					\draw[red,fill=black] (-0.5,-4)--(-0.5,5)node[below left]{$x = -\frac{1}{2}$};
				\end{tikzpicture}
				}
			\end{enumerate}
		}
\end{bt}

\begin{bt}%[2D1H5-1]
% nguon: Archive/2024/TaiLieuGiangDay/1.HamSo/1. K12 - 4in1 - Hamso/ThubaiGoc/Bai5-Dang2.tex #2119
%[Dự án Bài Giảng 12 Nhóm Toán & LaTex, Lê Minh Thiện Anh]
	Khảo sát sự biến thiên và vẽ đồ thị của các hàm số $y=\dfrac{2x-1}{x+1}$.
	\loigiai{
		\begin{itemize}
			\item Tập xác định $D=\mathbb{R}\setminus \{-1\}$.\\
			$\bullet$ Sự biến thiên.
			$y'=\dfrac{3}{(x+1)^2}>0$ với mọi $x$ khác $-1$. \\ 
			Hàm số đồng biến trên khoảng $\left(-\infty;-1\right)$ và $\left(-1;+\infty\right)$. \\ 
			Hàm số không có cực trị.\\
			Giới hạn \\ 
			$\underset{x\to +\infty}{\mathop{\lim}}y=\underset{x\to +\infty}{\mathop{\lim}}\dfrac{2x-1}{x+1}=2;\underset{x\to -\infty}{\mathop{\lim}}y=\underset{x\to -\infty}{\mathop{\lim}}\dfrac{2x-1}{x+1}=2$. \\ 
			$\underset{x\to -1^{-}}{\mathop{\lim}}y=\underset{x\to -1^{-}}{\mathop{\lim}}\dfrac{2x-1}{x+1}=+\infty;\underset{x\to -1^{+}}{\mathop{\lim}}y=\underset{x\to -1^{+}}{\mathop{\lim}}\dfrac{2x-1}{x+1}=-\infty$. \\
			Do đó, đồ thị hàm số nhận đường thẳng $x=-1$ làm tiệm cận đứng và đường thẳng $y=2$ làm tiệm cận ngang.\\
			Bảng biến thiên
			\begin{center}
				\begin{tikzpicture}[scale=1]
					\tikzset{double style/.append style = {draw=\tkzTabDefaultWritingColor,double=\tkzTabDefaultBackgroundColor,double distance=2pt}}
					\tikzset{h style/.style = {pattern=north west lines}}
					\tkzTabInit[nocadre=false,lgt=1.5,espcl=3.5]
					{$x$ /.6, $f'(x)$ /.6,$f(x)$ /2}
					{$-\infty$,$-1$,$+\infty$}
					\tkzTabLine{  ,+,d ,+ , }
					\tkzTabVar{-/$2$,+D-/$+\infty$/$-\infty$,+/$2$ }
				\end{tikzpicture}
			\end{center}
			$\bullet$ Đồ thị. \\ 
			Giao điểm của đồ thị hàm số với trục tung là $\left(0;-1 \right)$.\\
			$y=0 \Leftrightarrow \dfrac{2x-1}{x+1}=0\Leftrightarrow x=\dfrac{1}{2}$. \\ 
			Giao điểm của đồ thị hàm số với trục hoành là điểm $\left(\dfrac{1}{2};0\right)$.\\
			Đồ thị hàm số nhận $I\left(-1; 2 \right)$ làm tâm đối xứng.
			\begin{center}
				\begin{tikzpicture}[scale=0.7, line join=round, line cap=round, >=stealth]
					\tikzset{every node/.style={scale=1}}
					\def\xmin{-5}\def\xmax{3}\def\ymin{-2}\def\ymax{5}
					\draw[->] (\xmin-0.2,0)--(\xmax+0.2,0) node[below]{$x$};
					\draw[->] (0,\ymin-0.2)--(0,\ymax+0.2) node[right]{$y$};
					\draw (0,0) node[below left]{$O$};
					\foreach \x in {-5, -4,-3,-2,-1,1,2}\draw (\x,0.1)--(\x,-0.1) node[below]{$\x$};
					\foreach \y in {-1,1,2,3,4,5}\draw (0.1,\y)--(-0.1,\y) node[left]{$\y$};
					\clip (\xmin,\ymin) rectangle (\xmax,\ymax);
					\draw[dashed] (\xmin,2.0)--(\xmax,2.0);
					\draw[dashed] (-1.0,\ymin)--(-1.0,\ymax);
					\draw[smooth,samples=200,domain=\xmin:-1.1] plot (\x,{(2*(\x)+-1)/(1*(\x)+1)});
					\draw[smooth,samples=200,domain=-0.9:\xmax] plot (\x,{(2*(\x)+-1)/(1*(\x)+1)});
					\draw[dashed] (-1,0)--(-1,2)--(0,2);\fill (-1,2)circle(1.5pt);
				\end{tikzpicture}
			\end{center}
		\end{itemize}
	}
\end{bt}

\begin{bt}%[2D1H5-1]
% nguon: Archive/2024/TaiLieuGiangDay/1.HamSo/1. K12 - 4in1 - Hamso/ThubaiGoc/Bai5-Dang2.tex #2120
%[Dự án Bài Giảng 12 Nhóm Toán & LaTex, Lê Minh Thiện Anh]
	Khảo sát và vẽ đồ thị của các hàm số sau
	\begin{multicols}{2}
		\begin{enumerate}
			\item $y=3+ \dfrac{1}{x}$.
			\item $y=\dfrac{x-3}{1-x}$.
		\end{enumerate}
	\end{multicols}
	\loigiai{
		\begin{enumerate}
			\item Khảo sát sự biến thiên và vẽ đồ thị của hàm số $y=3+ \dfrac{1}{x}$.
			\begin{itemize}
				\item Tập xác định: $\mathscr{D}=\mathbb{R}\setminus\{0\}$.
				\item Đạo hàm: $y'=-\dfrac{1}{x^2}<0$ với mọi $x\in\mathscr{D}$.\\ 
				Suy ra hàm số nghịch biến trên từng khoảng $(-\infty;0)$ và $(0;+\infty)$.
				\item Giới hạn: $\lim\limits_{x\to\pm\infty}y=3$, $\lim\limits_{x\to(0)^-}y=-\infty$, $\lim\limits_{x\to(0)^+}y=+\infty$.\\ 
				Suy ra đồ thị hàm số có tiệm cận ngang $y=3$ và tiệm cận đứng $x=0$.
				\item Bảng biến thiên:
				\begin{center}
					\begin{tikzpicture}
						\tkzTabInit[lgt=1.2, espcl=3, deltacl=0.6]
						{$x$/0.6, $y'$/0.6, $y$/2}
						{$-\infty$, $0$, $+\infty$}
						\tkzTabLine{, -, d, -, }
						\tkzTabVar{+/$3$, -D+/$-\infty$/$+\infty$, -/$3$}
					\end{tikzpicture}
				\end{center}
				\item Bảng giá trị:
				\begin{center}
					\begin{tikzpicture}
						\tkzTabInit[lgt=1.2, espcl=1.5, deltacl=1]
						{$x$/0.7, $y$/1}
						{ $-3$, $-2$, $0$, $1$, $2$}
						\tkzTabLine{ \dfrac{8}{3}, , \dfrac{5}{2}, , 2, , 4, , \dfrac{7}{2}}
					\end{tikzpicture}
				\end{center}
				\item Đồ thị:
				\begin{center}
					\begin{tikzpicture}[x=1cm,y=1cm,scale=1,>=stealth]
						\def\aa{1} \def\bb{1} \def\cc{1} \def\dd{0}
						\draw[->] (-4,0)--(4,0)node[below right]{$x$};
						\draw[->] (0,-2)--(0,5)node[left]{$y$};
						\fill (0,0)node[below left]{$O$};
						\draw[thin] (-4,1)--(4,1) (0,-2)--(0,5);
						\draw[black,samples=150,smooth,domain=-4:-0.34] plot(\x,{(\aa*\x+\bb)/(\cc*\x+\dd)});
						\draw[black,samples=150,smooth,domain=0.25:4] plot(\x,{(\aa*\x+\bb)/(\cc*\x+\dd)});
						\foreach \x/\g in {-3/-90,-2/-90,-1/-90,1/-90,2/-90,3/-90}
						\draw[thin] (\x,2pt)--(\x,-2pt) + (\g:3mm) node [scale=0.5] {$\x$};
						%Vẽ các điểm trên trục Oy
						\foreach \y/\g in {4/180,-1/0,1/180,2/180,3/180}
						\draw[thin] (2pt,\y)--(-2pt,\y) + (\g:3mm) node [scale=0.5] {$\y$};
					\end{tikzpicture}
				\end{center}
			\end{itemize}
			\item Khảo sát sự biến thiên và vẽ đồ thị của hàm số $y=\dfrac{x-3}{1-x}$.
			\begin{itemize}
				\item Tập xác định: $\mathscr{D}=\mathbb{R}\setminus\{1\}$.
				\item Đạo hàm: $y'=\dfrac{-2}{(1-x)^2}<0$ với mọi $x\in\mathscr{D}$.\\ 
				Suy ra hàm số nghịch biến trên từng khoảng $(-\infty;1)$ và $(1;+\infty)$.
				\item Giới hạn: $\lim\limits_{x\to\pm\infty}y=-1$, $\lim\limits_{x\to(1)^-}y=-\infty$, $\lim\limits_{x\to(1)^+}y=+\infty$.\\ 
				Suy ra đồ thị hàm số có tiệm cận ngang $y=-1$ và tiệm cận đứng $x=1$.
				\item Bảng biến thiên:
				\begin{center}
					\begin{tikzpicture}
						\tkzTabInit[lgt=1.2, espcl=3, deltacl=0.6]
						{$x$/0.6, $y'$/0.6, $y$/2}
						{$-\infty$, $1$, $+\infty$}
						\tkzTabLine{, -, d, -, }
						\tkzTabVar{+/$-1$, -D+/$+\infty$/$-\infty$, -/$-1$}
					\end{tikzpicture}
				\end{center}
				\item Bảng giá trị:
				\begin{center}
					\begin{tikzpicture}
						\tkzTabInit[lgt=1.2, espcl=1.5, deltacl=1]
						{$x$/0.7, $y$/1}
						{$-2$, $-1$, $0$, $2$, $3$}
						\tkzTabLine{-\dfrac{5}{3}, , -2, , -3, , 1, , 0}
					\end{tikzpicture}
				\end{center}
				\item Đồ thị:
				\begin{center}
					\begin{tikzpicture}[x=1cm,y=1cm,scale=1,>=stealth]
						\def\aa{1} \def\bb{-3} \def\cc{-1} \def\dd{1}
						\draw[->] (-4,0)--(5,0)node[below right]{$x$};
						\draw[->] (0,-5)--(0,3)node[left]{$y$};
						\fill (0,0)node[below left]{$O$};
						\draw[thin] (-4,-1)--(5,-1) (1,-5)--(1,3);
						\draw[black,samples=150,smooth,domain=-4:0.50] plot(\x,{(\aa*\x+\bb)/(\cc*\x+\dd)});
						\draw[black,samples=150,smooth,domain=1.50:5] plot(\x,{(\aa*\x+\bb)/(\cc*\x+\dd)});
						\foreach \x/\g in {-3/-90,-2/-90,-1/-90,1/-90,2/-90,3/-90}
						\draw[thin] (\x,2pt)--(\x,-2pt) + (\g:3mm) node [scale=0.5] {$\x$};
						%Vẽ các điểm trên trục Oy
						\foreach \y/\g in {-3/180,-1/0,1/180,2/180,-2/180}
						\draw[thin] (2pt,\y)--(-2pt,\y) + (\g:3mm) node [scale=0.5] {$\y$};
					\end{tikzpicture}				
				\end{center}
			\end{itemize}			
		\end{enumerate}
	}
\end{bt}

\begin{bt}%[2D1H5-1]
% nguon: Archive/2024/TaiLieuGiangDay/1.HamSo/1. K12 - 4in1 - Hamso/ThubaiGoc/Bai5-Dang2.tex #2121
%[Dự án Bài Giảng 12 Nhóm Toán & LaTex, Lê Minh Thiện Anh] 
	Cho hàm số $y=\dfrac{2x+1}{x-1}$.
	\begin{enumerate}
		\item Khảo sát và vẽ đồ thị của hàm số.
		\item Gọi $A$ là giao điểm của đồ thị hàm số với trục $Oy$, $I$ là giao điểm của hai đường tiệm cận của đồ thị hàm số. Tìm điểm $B$ đối xứng với $A$ qua $I$. Chứng minh rằng điểm $B$ cũng thuộc đồ thị hàm số này.
	\end{enumerate}
	\loigiai{
		\begin{enumerate}
			\item Khảo sát và vẽ đồ thị của hàm số
			\begin{itemize}
			\item Tập xác định $\mathscr{D}=\mathbb{R}\setminus \{1\}$.
			\item Sự biến thiên\\
				Đạo hàm $y'=-\dfrac{3}{(x-1)^2}$. Vì $y'<0$ với mọi $x\neq 1$ nên hàm số nghịch biến trên mỗi khoảng $(-\infty;1)$ và $(1;+\infty)$.\\
			Ta có $\lim\limits_{x\to -\infty} y=\lim\limits_{x\to -\infty} \dfrac{2x+1}{x-1}=2$; $\lim\limits_{x\to +\infty} y=\lim\limits_{x\to +\infty} \dfrac{2x+1}{x-1}=2$.\\
				$\Rightarrow$ Đường thẳng $y=2$ là tiệm cận ngang của đồ thị hàm số.\\
				Lại có $\lim\limits_{x\to 1^{-}} y=\lim\limits_{x\to 1^{-}} \dfrac{2x+1}{x-1}=-\infty$; $\lim\limits_{x\to 1^{+}} y=\lim\limits_{x\to 1^{+}} \dfrac{2x+1}{x-1}=+\infty$.\\
				$\Rightarrow$ Đường thẳng $x=1$ là tiệm cận đứng của đồ thị hàm số.\\
				Bảng biến thiên
				\begin{center}
			\begin{tikzpicture}
					\tkzTabInit[nocadre=false,lgt=1.2,espcl=4,deltacl=0.6]
					{$x$ /0.6,$y'$ /0.6,$y$ /2}
					{$-\infty$,$1$,$+\infty$}
					\tkzTabLine{,-,d,-,}
					\tkzTabVar{+/$2$,-D+/$-\infty$/$+\infty$,-/$2$}
				\end{tikzpicture}
			\end{center}
			\item Đồ thị
				\immini{
					Đồ thị của hàm số giao với trục $Ox$ tại điểm $\left(-\dfrac{1}{2};0\right)$, giao với trục $Oy$ tại điểm $(0;-1)$.\\
					Đồ thị của hàm số đã cho được biểu diễn như hình bên. Đồ thị của hàm số có tâm đối xứng là điểm $I(1;2)$.
					}{
				\begin{tikzpicture}[scale=0.7, font=\footnotesize, line join=round, line cap=round, >=stealth]
					\def\a{2} \def\b{1} \def\c{1} \def\d{-1} % Hệ số
					\def\xt{-3} \def\xp{5} \def\yt{6} \def\yd{-2} % x_trái, x_phải, y_trên, y_dưới (giới hạn)
					\draw[->] (\xt,0)--(\xp,0) node [below]{$x$};
					\draw[->] (0,\yd)--(0,\yt) node [left]{$y$};
					\node at (0,0) [below right]{$O$};
					\clip (\xt+0.1,\yd+0.1) rectangle (\xp-0.1,\yt-0.1);
					\draw[thick,red,smooth,samples=300,domain=\xt:(-\d/\c-0.1)] plot(\x,{(\a*(\x)+\b)/(\c*(\x)+\d)});
					\draw[thick,red,smooth,samples=300,domain=(-\d/\c+0.1:\xp)] plot(\x,{(\a*(\x)+\b)/(\c*(\x)+\d)});
					\draw[dashed] (-\d/\c,\yd)--(-\d/\c,\yt);
					\draw[dashed] (\xt,\a/\c)--(\xp,\a/\c);
					\fill (1,2) circle (1.5pt) node[below right]{$I$};
					\fill (0,-1) circle (1.5pt) node[left]{$-1$};
					\fill (-0.5,0) circle (1.5pt) node[above]{$-\frac{1}{2}$};
					\foreach \x in {-2,-1,1,2,3,4}
					\draw[shift={(\x,0)},color=black] (0pt,2pt) -- (0pt,-2pt);
				\foreach \y in {-1,1,2,3,4,5}
					\draw[shift={(0,\y)},color=black] (2pt,0pt) -- (-2pt,0pt);
					\node at (4.2,2) [below]{$y=2$};
					\node at (1,-1.5) [right]{$x=1$};
				\end{tikzpicture}
			}
			\end{itemize}
			\item Ta có $A(0;-1)$ và $I(1;2)$.\\
				Do điểm $B$ đối xứng với $A$ qua $I$ nên $I$ là trung điểm của đoạn $AB$ $\Rightarrow B(2;5)$.\\
				Thay tọa độ $B$ vào đồ thị hàm số $y=\dfrac{2x+1}{x-1}$ ta có $5=\dfrac{2\cdot 2+1}{2-1} \Leftrightarrow 5=5$ (đúng).\\
				Vậy điểm $B$ cũng thuộc đồ thị hàm số đã cho.
		\end{enumerate}
		}
\end{bt}

\begin{bt}%[2D1V5-8]
% nguon: Archive/2024/TaiLieuGiangDay/1.HamSo/1. K12 - 4in1 - Hamso/ThubaiGoc/Bai5-Dang2.tex #2122
%[Dự án Bài Giảng 12 Nhóm Toán & LaTex, Lê Minh Thiện Anh]
	Một trung tâm dạy nghề cần thiết kế mô hình đánh giá kĩ năng của một học viên theo học nghề đánh máy. Người ta có thể làm như sau:
	\begin{itemize}
		\item Để xây dựng mô hình toán học cho bài toán trên, ta sử dụng thống kê. Bằng cách khảo sát tốc độ đánh máy trung bình $S$ (tính bằng từ trên phút) của học viên đó sau $t$ tuần học ($5\leq t\leq 30$), ta thu thập các số liệu thống kê được cho trong bảng bên dưới. \textit{(Nguồn: R. Larson and B. Edwards, Calculus 10e, Cengage 2014).}
		\begin{center}
			\begin{tabular}{|c|c|c|c|c|c|c|c|}
				\hline
				$t$ & $5$ & $10$ & $15$ & $20$ & $25$ & $30$\\
				\hline
				$S$ & $38$ & $56$ & $79$ & $90$ & $93$ & $94$\\
				\hline
			\end{tabular}
		\end{center}
		\item Ta cần chọn hàm số $y=f(t)$ để biểu diễn các số liệu ở bảng trên, tức là ở hệ trục tọa độ $Oty$, đồ thị của hàm số đó trên khoảng $(0;+\infty)$ \lq\lq gần\rq\rq\ với các điểm $A(5;38)$, $B(10;56)$, $C(15;79)$, $D(20;90)$, $E(25;93)$, $G(30;94)$. Ngoài ra, do tốc độ đánh máy trung bình của học viên tăng theo thời gian $t$ và chỉ đến một giới hạn $M$ nào đó cho dù thời gian $t$ có kéo dài đến vô cùng nên hàm số $y=f(t)$ phải thỏa mãn thêm hai điều kiện: Hàm số đó đồng biến trên khoảng $(0;+\infty)$ và $\lim\limits_{t\to +\infty}f(t)=M\in\mathbb{R}$, $M>94$. Vì các hàm đa thức (với bậc lớn hơn hoặc bằng $1$) không thỏa mãn hai điều kiện đó nên ta chọn một hàm phân thức hữu tỉ để biểu diễn các số liệu ở bảng trên. Ta có thể chọn hàm số có dạng $f(t)=\dfrac{at+b}{ct+d}$ ($ac\neq 0$) cho mục đích đó. Dựa vào bảng trên, ta chọn hàm số $f(t)=\dfrac{110t-280}{t+2}$ $(t>0)$.
	\end{itemize}
	\begin{enumerate}
		\item Dựa theo mô hình đó, dự đoán tốc độ đánh máy trung bình của học viên đó sau $40$ tuần (làm tròn kết quả đến hàng đơn vị của từ/phút).
		\item Xem $y=f(t)$ là một hàm số xác định trên khoảng $(0;+\infty)$, hãy tìm tiệm cận ngang của đồ thị hàm số đó.
		\item Nêu nhận xét về tốc độ đánh máy trung bình của học viên đó sau thời gian $t$ ngày càng lớn.
	\end{enumerate}
	\loigiai{
		\begin{enumerate}
			\item Ta có $f(40)=\dfrac{110\cdot40-280}{40+2}\approx 98$. Vậy tốc độ đánh máy trung bình của học viên đó sau $40$ tuần là khoảng $98$ từ/phút.
			\item Ta có $\lim\limits_{t\to+\infty}f(t)=\lim\limits_{t\to+\infty}\dfrac{110t-280}{t+2}=\lim\limits_{t\to+\infty}\left(110-\dfrac{500}{t+2}\right)=110$.\\
			Vậy đường thẳng $y=110$ là tiệm cận ngang của đồ thị hàm số $y=f(t)$.
			\item Do đường thẳng $y=110$ là tiệm cận ngang của đồ thị hàm số $y=f(t)$ nên khi $t$ càng lớn thì tốc độ đánh máy trung bình của học viên đó sẽ tiến gần đến mức $110$ từ/phút.
		\end{enumerate}
		\immini{
			\begin{note}
				Đồ thị hàm số $y=f(t)=\dfrac{110t-280}{t+2}$ ($t>0$) được cho ở hình bên. Đồ thị đó giao với trục tung tại điểm có tọa độ là $(0;-140)$.
			\end{note}
		}{
			\begin{tikzpicture}[>=stealth, yscale=0.6]
				\draw[->] (0,0)--(7,0) node[below] {$t$};
				\draw[->] (0,0)--(0,5) node[left] {$y$};
				\draw[domain=0.166:6, smooth, samples=500] plot (\x, {(12*\x-2)/(3*\x-0)});
				\draw (0,4)--(6.5,4) node[above left=0 and 0.5] {$y=110$};
				\draw[fill=black] (0,0) node[below left=-0.1] {$O$} circle (1.2pt);
				\draw[fill=black] (0,4) node[left] {$110$} circle (1.2pt);
			\end{tikzpicture}
		}
	}
\end{bt}

\begin{bt}%[2D1V5-8]
% nguon: Archive/2024/TaiLieuGiangDay/1.HamSo/1. K12 - 4in1 - Hamso/ThubaiGoc/Bai5-Dang2.tex #2123
%[Dự án Bài Giảng 12 Nhóm Toán & LaTex, Lê Minh Thiện Anh]
	Ở một bể chứa nước có chứa 1000 lít nước ngọt. Người ta bơm nước biển có nồng độ muối là 30 gam/lít vào bể nước với tốc độ là 25 lít/phút.
	\begin{enumerate}
		\item Chứng minh rằng nồng độ muối của nước trong bể sau $t$ phút kể từ khi bắt đầu bơm là $C(t)=\dfrac{30 t}{40+t}$ (gam/lít).
		\item Khảo sát sự biến thiên của hàm số $y=C(t)$ sau 10 tiếng kể từ lúc bắt đầu bơm, từ đó nhận xét về nồng độ muối trong bể khi thời gian $t$ càng lớn.
	\end{enumerate}
	\loigiai{\begin{enumerate}
			\item Ta có
			\begin{itemize}
				\item Tại thời điểm $t$, tổng dung tích của bể sẽ là $1000 + 25t$ lít.
				\item Khối lượng muối bơm vào sau $t$ phút ta được $30 \cdot 25t$ gam.
				\end{itemize}
				Do đó nồng độ muối của nước trong bể sau $t$ phút kể từ khi bắt đầu bơm là $C(t)=\dfrac{30 t}{40+t}$ (gam/lít).
				\item Khảo sát hàm số $C(t)=\dfrac{30t}{40+t}$ trên cơ sở đó nhận xét về nồng độ muối trong bể khi $t$ càng lớn.\\
				Ta có $\lim\limits_{t \rightarrow+\infty} C(t)=\lim\limits_{t \rightarrow+\infty}\dfrac{30 t}{40+t}=30$.\\
				Suy ra đường thẳng $C(t) = 30$ là đường tiệm cận ngang của đồ thị hàm số đã cho.\\
				Ta có bẳng biến thiên
				\begin{center}
			\begin{tikzpicture}[>=stealth]
					\tkzTabInit[nocadre=false,lgt=1.5,espcl=4.5,deltacl=0.5]{$t$/.7 ,$C(t)'$/.7,$C(t)$/1.5}
					{$0$ , $+\infty$}
					\tkzTabLine{ , + , }
					\tkzTabVar{-/$0$ , +/$30$}
				\end{tikzpicture}
				\end{center}
				Ta thấy, khi $t$ càng lớn thì nồng độ muối trong bể sẽ tiến gần đến mực 30 (gam/lít). Lúc đó, nồng độ muối trong bể sẽ gần như bằng nồng độ muối trong nước muối được bơm vào bể.
			\end{enumerate}
		}
\end{bt}

\begin{bt}%[2D1N5-4]
% nguon: Archive/2024/TaiLieuGiangDay/1.HamSo/1. K12 - 4in1 - Hamso/ThubaiGoc/Bai5-Dang3.tex #2134
[Mức độ 1]%[Bài giảng 12 mới, Lê Hồng Phi]
Tìm tọa độ giao điểm giữa đồ thị hàm số $y=\dfrac{x^2-x+1}{x-1}$ với trục tung.
\loigiai
{Gọi $M(x_M;y_M)$ là giao điểm giữa đồ thị hàm số đã cho với trục tung.\\
Vì $M\in Oy$ nên $x_M=0$.\\
Do $M$ thuộc đồ thị hàm số đã cho nên $y_M=\dfrac{x_M^2-x_M+1}{x_M-1}=\dfrac{0^2-0+1}{0-1}=-1$.\\
Vậy tọa độ giao điểm cần tìm là $M(0;-1)$.
}
\end{bt}

\begin{bt}%[2D1N5-4]
% nguon: Archive/2024/TaiLieuGiangDay/1.HamSo/1. K12 - 4in1 - Hamso/ThubaiGoc/Bai5-Dang3.tex #2135
[Mức độ 1]
Tìm tọa độ giao điểm của đồ thị hàm số $y=\dfrac{x^2+x-2}{x+1}$ với trục hoành.
\loigiai
{Phương trình hoành độ giao điểm giữa đồ thị hàm số đã cho với trục hoành là \begin{eqnarray*}
&&\dfrac{x^2+x-2}{x+1}=0\qquad (x\neq -1)\\
&\Leftrightarrow & x^2+x-2=0\\
&\Leftrightarrow &\hoac{& x=1 \\ & x=-2.}
\end{eqnarray*}
Vậy tọa độ các giao điểm cần tìm là $(1;0)$ và $(-2;0)$.
}
\end{bt}

\begin{bt}%[2D1H5-3]
% nguon: Archive/2024/TaiLieuGiangDay/1.HamSo/1. K12 - 4in1 - Hamso/ThubaiGoc/Bai5-Dang3.tex #2136
[Cùng khám phá, Mức độ 2]
\immini{Đồ thị $(C)$ của hàm số $y=\dfrac{-x^2-3x-6}{x+1}$ là đường cong như hình vẽ bên. Tìm $m$ để đường thẳng $d\colon y=m$ cắt đồ thị $(C)$ tại hai điểm phân biệt.}{\begin{tikzpicture}[scale=1, font=\footnotesize, line join=round, line cap=round,>=stealth,x=0.4cm,y=0.4cm]
\def \xmin{-9.5};
\def \xmax{7};
\def \ymin{-7.5};
\def \ymax{7.5};
\draw[->] (\xmin, 0.) -- (\xmax,0.) node[anchor=north] {$x$};
\draw[->] (0.,\ymin) -- (0.,\ymax) node[anchor=west] {$y$};
\clip(\xmin+0.1,\ymin+0.1) rectangle (\xmax-0.1,\ymax-0.1);
\draw[smooth,samples=300,domain=\xmin:-1.01] plot(\x,{(-(\x)^2-3*(\x)-6)/((\x)+1)});
\draw[smooth,samples=300,domain=-0.99:\xmax] plot(\x,{(-(\x)^2-3*(\x)-6)/((\x)+1)});
\draw[smooth,samples=200,domain=\xmin:\xmax] plot(\x,{-(\x)-2});
\draw (-1,\ymin)--(-1,\ymax);
\draw[dashed] (-3,0)node[below]{$-3$}|-(0,3)node[right]{$3$} (1,0)node[above]{$1$}|-(0,-5)node[left]{$-5$};
\foreach \x/\y/\l/\g  in {0/0/O/135,0/3//0,1/0//0,-3/0//0,0/-5//0,-3/3//0,1/-5//0}{\fill[black] (\x,\y) circle(1pt)+(\g:0.6) node{$\l$};}
\end{tikzpicture}}
\loigiai
{\immini{Đường thẳng $d\colon y=m$ vuông góc và cắt trục $Oy$ tại điểm có tung độ bằng $m$.\\
Dựa vào đồ thị, đường thẳng $d$ cắt đồ thị $(C)$ tại hai điểm phân biệt khi $m>3$ hoặc $m<-5$.}{\begin{tikzpicture}[scale=1, font=\footnotesize, line join=round, line cap=round,>=stealth,x=0.4cm,y=0.4cm]
\def \xmin{-9.5};
\def \xmax{7};
\def \ymin{-7.5};
\def \ymax{7.5};
\draw[->] (\xmin, 0.) -- (\xmax,0.) node[anchor=north] {$x$};
\draw[->] (0.,\ymin) -- (0.,\ymax) node[anchor=west] {$y$};
\clip(\xmin+0.1,\ymin+0.1) rectangle (\xmax-0.1,\ymax-0.1);
\draw[smooth,samples=300,domain=\xmin:-1.01] plot(\x,{(-(\x)^2-3*(\x)-6)/((\x)+1)});
\draw[smooth,samples=300,domain=-0.99:\xmax] plot(\x,{(-(\x)^2-3*(\x)-6)/((\x)+1)});
\draw[smooth,samples=200,domain=\xmin:\xmax] plot(\x,{-(\x)-2});
\draw (-1,\ymin)--(-1,\ymax) (\xmin,4)--(\xmax,4)node[above,pos=0.9]{$y=m$};
\draw[dashed] (-3,0)node[below]{$-3$}|-(0,3)node[right]{$3$} (1,0)node[above]{$1$}|-(0,-5)node[left]{$-5$};
\foreach \x/\y/\l/\g  in {0/0/O/135,0/3//0,1/0//0,-3/0//0,0/-5//0,-3/3//0,1/-5//0,0/4/m/45}{\fill[black] (\x,\y) circle(1pt)+(\g:0.6) node{$\l$};}
\end{tikzpicture}}
}
\end{bt}

\begin{bt}%[2D1H5-1]
% nguon: Archive/2024/TaiLieuGiangDay/1.HamSo/1. K12 - 4in1 - Hamso/ThubaiGoc/Bai5-Dang3.tex #2137
[CTST, Mức độ 2]%[Bài giảng 12 mới, Lê Hồng Phi]
Khảo sát và vẽ đồ thị hàm số $y=\dfrac{x^2+2x-2}{x-1}$.
	\loigiai{
\begin{enumEX}{1}
\item Tập xác định $\mathscr{D}=\mathbb{R}\setminus\{1\}$.
\item Sự biến thiên
\begin{itemize}
\item Chiều biến thiên:\\
Ta có $y=x+3+\dfrac{1}{x-1}$ và $y'=\dfrac{x^2-2x}{(x-1)^2}$.\\
Khi đó $y'=0\Leftrightarrow\hoac{& x=0 \\ & x=2.}$\\
Trên các khoảng $(-\infty;0)$, $(2;+\infty)$ thì $y'>0$ nên hàm số đồng biến trên các khoảng đó.\\
Trên khoảng $(0;2)$ thì $y'<0$ nên hàm số nghịch biến trên khoảng đó.
\item Cực trị:\\
Hàm số đạt cực tiểu tại $x=2$ và $y_\text{CT}=6$.\\
Hàm số đạt cực đại tại $x=0$ và $y_\text{CĐ}=2$.
\item Các giới hạn và tiệm cận:\\
$\lim\limits_{x\to -\infty}y=-\infty$, $\lim\limits_{x\to +\infty}y=+\infty$.\\
$\lim\limits_{x\to \pm\infty}\left[y-(x+3)\right]=0$ nên đường thẳng $y=x+3$ là tiệm cận xiên của đồ thị hàm số.\\
$\lim\limits_{x\to 1^-}y=-\infty$, $\lim\limits_{x\to 1^+}y=+\infty$ nên đường thẳng $x=1$ là tiệm cận đứng của đồ thị hàm số.
\item Bảng biến thiên
\begin{center}
\begin{tikzpicture}
\tkzTabInit[nocadre=false, lgt=1.2, espcl=2.5, deltacl=0.6]{$x$/0.6,$y'$/0.6,$y$/2}
{$-\infty$, $0$, $1$, $2$, $+\infty$}
\tkzTabLine {,+,0,-,d,-,0,+,}
\tkzTabVar{-/$-\infty$, +/$2$, -D+/$-\infty$ /$+\infty$,-/$6$, +/$+\infty$}
\end{tikzpicture}
\end{center}
\end{itemize}
\item Đồ thị\\
\immini{Đồ thị cắt trục $Oy$ tại điểm $(0;2)$.\\
Đồ thị nhận giao điểm của hai tiệm cận là $I(1;4)$ làm tâm đối xứng.\\
Các trục đối xứng của đồ thị hàm số là hai đường phân giác của các góc tạo bởi hai đường tiệm cận $x=1$ và $y=x+3$.}{\begin{tikzpicture}[scale=1, font=\footnotesize, line join=round, line cap=round,>=stealth,x=0.6cm,y=0.6cm]
\def \xmin{-4.5};
\def \xmax{6.7};
\def \ymin{-1.5};
\def \ymax{8.5};
\draw[->] (\xmin, 0.) -- (\xmax,0.) node[anchor=north] {$x$};
\draw[->] (0.,\ymin) -- (0.,\ymax) node[anchor=west] {$y$};
\clip(\xmin+0.1,\ymin+0.1) rectangle (\xmax-0.1,\ymax-0.1);
\draw[smooth,samples=200,domain=\xmin:0.99] plot(\x,{((\x)^2+2*(\x)-2)/((\x)-1)});
\draw[smooth,samples=200,domain=1.01:\xmax] plot(\x,{((\x)^2+2*(\x)-2)/((\x)-1)});
\draw[smooth,samples=200,domain=\xmin:\xmax] plot(\x,{(\x)+3});
\draw (1,\ymin)--(1,\ymax);
\draw[dashed] (0,6) node[left]{$6$} -| (2,0)node[below]{$2$};
\foreach \x/\y/\l/\g  in {0/0/O/225,0/2/2/135,2/6//0,1/0/1/-45,0/6//0,2/0//0}{\fill[black] (\x,\y) circle(1pt)+(\g:0.5) node{$\l$};}
\end{tikzpicture}}
\end{enumEX}
	}	
\end{bt}

\begin{bt}%[2D1H5-1]
% nguon: Archive/2024/TaiLieuGiangDay/1.HamSo/1. K12 - 4in1 - Hamso/ThubaiGoc/Bai5-Dang3.tex #2138
[CTST, Mức độ 2]%[Bài giảng 12 mới, Lê Hồng Phi]
Khảo sát và vẽ đồ thị hàm số $y=\dfrac{-x^2-3x+4}{x+2}$.
	\loigiai{
\begin{enumEX}{1}
\item Tập xác định $\mathscr{D}=\mathbb{R}\setminus\{-2\}$.
\item Sự biến thiên
\begin{itemize}
\item Chiều biến thiên:\\
Ta có $y=-x-1+\dfrac{6}{x+2}$ và $y'=-1-\dfrac{6}{(x+2)^2}<0,\ \forall x\in\mathscr{D}$.\\
 Suy ra hàm số nghịch biến trên mỗi khoảng $(-\infty;-2)$, $(-2;+\infty)$.
\item Các giới hạn và tiệm cận:\\
$\lim\limits_{x\to -\infty}y=+\infty$, $\lim\limits_{x\to +\infty}y=-\infty$.\\
$\lim\limits_{x\to \pm\infty}\left[y-(-x-1)\right]=0$ nên đường thẳng $y=-x-1$ là tiệm cận xiên của đồ thị hàm số.\\
$\lim\limits_{x\to -2^-}y=-\infty$, $\lim\limits_{x\to -2^+}y=+\infty$ nên đường thẳng $x=1$ là tiệm cận đứng của đồ thị hàm số.
\item Bảng biến thiên
\begin{center}
\begin{tikzpicture}
\tkzTabInit[nocadre=false, lgt=1.2, espcl=2.5, deltacl=0.6]{$x$/0.6,$y'$/0.6,$y$/2}
{$-\infty$, $-2$, $+\infty$}
\tkzTabLine {,-,d,-,}
\tkzTabVar{+/$+\infty$, -D+/$-\infty$ /$+\infty$,-/$-\infty$}
\end{tikzpicture}
\end{center}
\end{itemize}
\item Đồ thị\\
\immini{Đồ thị cắt trục $Oy$ tại điểm $(0;2)$.\\
Ta có $y=0\Leftrightarrow -x^2-3x+4=0\Leftrightarrow\hoac{& x=-4 \\ & x=1.}$\\
Suy ra đồ thị cắt trục $Ox$ tại các điểm $(-4;0)$, $(1;0)$.\\
Đồ thị nhận giao điểm của hai tiệm cận là $I(-2;1)$ làm tâm đối xứng.\\
Các trục đối xứng của đồ thị hàm số là hai đường phân giác của các góc tạo bởi hai đường tiệm cận $x=-2$ và $y=-x-1$.}{\begin{tikzpicture}[scale=1, font=\footnotesize, line join=round, line cap=round,>=stealth,x=0.6cm,y=0.6cm]
\def \xmin{-8.0};
\def \xmax{3.7};
\def \ymin{-3.5};
\def \ymax{6.5};
\draw[->] (\xmin, 0.) -- (\xmax,0.) node[anchor=north] {$x$};
\draw[->] (0.,\ymin) -- (0.,\ymax) node[anchor=west] {$y$};
\clip(\xmin+0.1,\ymin+0.1) rectangle (\xmax-0.1,\ymax-0.1);
\draw[smooth,samples=200,domain=\xmin:-2.01] plot(\x,{(-(\x)^2-3*(\x)+4)/((\x)+2)});
\draw[smooth,samples=200,domain=-1.99:\xmax] plot(\x,{(-(\x)^2-3*(\x)+4)/((\x)+2)});
\draw[smooth,samples=200,domain=\xmin:\xmax] plot(\x,{-(\x)-1});
\draw (-2,\ymin)--(-2,\ymax);

\foreach \x/\y/\l/\g  in {0/0/O/135,0/2/2/180,1/0/1/-90,-4/0/-4/-135,-2/0/-2/210,-2/1/I/30}{\fill[black] (\x,\y) circle(1pt)+(\g:0.5) node{$\l$};}
\end{tikzpicture}}
\end{enumEX}
	}	
\end{bt}

\begin{bt}%[2D1H5-7]
% nguon: Archive/2024/TaiLieuGiangDay/1.HamSo/1. K12 - 4in1 - Hamso/ThubaiGoc/Bai5-Dang3.tex #2139
[Mức độ 2]%[Bài giảng 12 mới, Lê Hồng Phi]
Cho hàm số $y=\dfrac{2x^2-5x+5}{x-2}$ có đồ thị $(C)$. Tìm tất cả các điểm thuộc đồ thị $(C)$ có tọa độ là số nguyên.
\loigiai
{Tập xác định $\mathscr{D}=\mathbb{R}\setminus\{2\}$.\\
Ta có $y=2x-1+\dfrac{3}{x-2}$ nên điểm $M(x_M;y_M)$ thuộc $(C)$ có tọa độ nguyên khi và chỉ khi \[\heva{& x_M\in\mathbb{Z} \\ & 3\ \vdots\ (x_M-2)}\Leftrightarrow x_M-2\in\{-3;-1;1;3\}\Leftrightarrow x_M\in\{-1;1;3;5\}.
\]
Vậy có $4$ điểm thuộc $(C)$ có tọa độ nguyên là $(-1;-4)$, $(1;-2)$, $(3;8)$, $(5;10)$.
}
\end{bt}

\begin{bt}%[2D1H5-6]
% nguon: Archive/2024/TaiLieuGiangDay/1.HamSo/1. K12 - 4in1 - Hamso/ThubaiGoc/Bai5-Dang3.tex #2140
[Mức độ 2]%[Bài giảng 12 mới, Lê Hồng Phi]
Cho hàm số $y=\dfrac{-x^2+x+1}{x+1}$ có đồ thị $(C)$. Viết phương trình đường thẳng đi qua điểm $A(-1;5)$ và tiếp xúc với đồ thị $(C)$.
\loigiai
{Tập xác định $\mathscr{D}=\mathbb{R}\setminus\{-1\}$.\\
Ta có $y'=\dfrac{-x^2-2x}{(x+1)^2}$.\\
Gọi $M\left(x_0;\dfrac{-x_0^2+x_0+1}{x_0+1}\right)$  là tiếp điểm $(x_0\neq -1)$.\\
Phương trình tiếp tuyến $\Delta$  của đồ thị $(C)$ tại $M$ là $y=\dfrac{-x_0^2-2x_0}{(x_0+1)^2}\cdot (x-x_0)+\dfrac{-x_0^2+x_0+1}{x_0+1}$.\\
Vì $\Delta$ đi qua $A(-1;5)$ nên \begin{eqnarray*}
&& 5=\dfrac{-x_0^2-2x_0}{(x_0+1)^2}\cdot (-1-x_0)+\dfrac{-x_0^2+x_0+1}{x_0+1}\\
&\Leftrightarrow & 5=\dfrac{x_0^2+2x_0}{x_0+1}+\dfrac{-x_0^2+x_0+1}{x_0+1}\\
&\Leftrightarrow & 5(x_0+1)=3x_0+1\\
&\Leftrightarrow & x_0=-2.
\end{eqnarray*}
Với $x_0=2$ ta có phương trình tiếp tuyến $\Delta$ là $y=5$.\\
Vậy phương trình đường thẳng cần tìm là $y=5$.
}
\end{bt}

\begin{bt}%[2D1V5-4]
% nguon: Archive/2024/TaiLieuGiangDay/1.HamSo/1. K12 - 4in1 - Hamso/ThubaiGoc/Bai5-Dang3.tex #2143
%[Bài giảng 12 mới, Lê Hồng Phi]
Tìm tất cả các giá trị của $m$ để đường thẳng $y=-x+m$ cắt đồ thị hàm số $y=\dfrac{x^2-1}{x}$ tại hai điểm phân biệt $A$, $B$ sao cho $AB=\sqrt{6}$.
\loigiai{Phương trình hoành độ giao điểm là $\dfrac{x^2-1}{x}=-x+m\Leftrightarrow\heva{& x\neq 0 \\ & 2x^2-mx-1=0.\quad (*)}$\\
Đường thẳng cắt đồ thị tại hai điểm $A$, $B$ khi và chỉ khi phương trình $(*)$ có hai nghiệm phân biệt khác $0$. Điều kiện tương đương là \[\heva{& \Delta=(-m)^2+2>0 \\ & 2\cdot 0^2-m\cdot 0-1\neq 0}\Leftrightarrow\heva{& m^2+2>0 \\ & -1\neq 0}\Leftrightarrow\forall m\in\mathbb{R}.\]
Như vậy đường thẳng luôn cắt đồ thị hàm số tại hai điểm phân biệt $A\left(x_A;-x_A+m\right)$, $B\left(x_B;-x_B+m\right)$ với $x_A$, $x_B$ là nghiệm của phương trình $(*)$.\\
Theo Vi-et thì $x_A+x_B=\dfrac{m}{2}$ và $x_A x_B=-\dfrac{1}{2}$.\\
Ta có \begin{eqnarray*}
AB=\sqrt{6} &\Leftrightarrow & AB^2=6\\
&\Leftrightarrow & \left(x_B-x_A\right)^2+(x_A-x_B)^2=6\\
&\Leftrightarrow & 2\left[\left(x_A+x_B\right)^2-4 x_A x_B\right]=6\\
&\Leftrightarrow & \dfrac{m^2}{2}+4=6\\ &\Leftrightarrow & m=\pm 2.
\end{eqnarray*}
Vậy các giá trị $m$ cần tìm là $m=\pm 2$.
} 
\end{bt}

\begin{bt}%[2D1H5-1]
% nguon: Archive/2024/TaiLieuGiangDay/1.HamSo/1. K12 - 4in1 - Hamso/ThubaiGoc/Bai5-Dang3.tex #2145
%[Bài giảng 12 mới, Lê Hồng Phi]
Khảo sát và vẽ đồ thị hàm số $y=\dfrac{x^2+x-2}{x+1}$.
	\loigiai{
\begin{enumEX}{1}
\item Tập xác định $\mathscr{D}=\mathbb{R}\setminus\{-1\}$.
\item Sự biến thiên
\begin{itemize}
\item Chiều biến thiên:\\
Ta có $y=x-\dfrac{2}{x+1}$ và $y'=1+\dfrac{2}{(x+1)^2}>0,\ \forall x\in\mathscr{D}$.\\
 Suy ra hàm số đồng biến trên mỗi khoảng $(-\infty;-1)$, $(-1;+\infty)$.
\item Các giới hạn và tiệm cận:\\
$\lim\limits_{x\to -\infty}y=-\infty$, $\lim\limits_{x\to +\infty}y=+\infty$.\\
$\lim\limits_{x\to \pm\infty}\left[y-x\right]=0$ nên đường thẳng $y=x$ là tiệm cận xiên của đồ thị hàm số.\\
$\lim\limits_{x\to -1^-}y=+\infty$, $\lim\limits_{x\to -1^+}y=-\infty$ nên đường thẳng $x=1$ là tiệm cận đứng của đồ thị hàm số.
\item Bảng biến thiên
\begin{center}
\begin{tikzpicture}
\tkzTabInit[nocadre=false, lgt=1.2, espcl=2.5, deltacl=0.6]{$x$/0.6,$y'$/0.6,$y$/2}
{$-\infty$, $-1$, $+\infty$}
\tkzTabLine {,+,d,+,}
\tkzTabVar{-/$-\infty$, +D-/$+\infty$ /$-\infty$,+/$+\infty$}
\end{tikzpicture}
\end{center}
\end{itemize}
\item Đồ thị\\
\immini{Đồ thị cắt trục $Oy$ tại điểm $(0;-2)$.\\
Ta có $y=0\Leftrightarrow x^2+x-2=0\Leftrightarrow\hoac{& x=-2 \\ & x=1.}$\\
Suy ra đồ thị cắt trục $Ox$ tại các điểm $(-2;0)$, $(1;0)$.\\
Đồ thị nhận giao điểm của hai tiệm cận là $I(-1;-1)$ làm tâm đối xứng.\\
Các trục đối xứng của đồ thị hàm số là hai đường phân giác của các góc tạo bởi hai đường tiệm cận $x=-1$ và $y=x$.}{\begin{tikzpicture}[scale=1, font=\footnotesize, line join=round, line cap=round,>=stealth,x=0.6cm,y=0.6cm]
\def \xmin{-6.0};
\def \xmax{6.7};
\def \ymin{-5.5};
\def \ymax{5.5};
\draw[->] (\xmin, 0.) -- (\xmax,0.) node[anchor=north] {$x$};
\draw[->] (0.,\ymin) -- (0.,\ymax) node[anchor=west] {$y$};
\clip(\xmin+0.1,\ymin+0.1) rectangle (\xmax-0.1,\ymax-0.1);
\draw[smooth,samples=200,domain=\xmin:-1.01] plot(\x,{((\x)^2+(\x)-2)/((\x)+1)});
\draw[smooth,samples=200,domain=-0.99:\xmax] plot(\x,{((\x)^2+(\x)-2)/((\x)+1)});
\draw[smooth,samples=200,domain=\xmin:\xmax] plot(\x,{(\x)});
\draw (-1,\ymin)--(-1,\ymax);

\foreach \x/\y/\l/\g  in {0/0/O/135,0/-2/-2/0,1/0/1/-90,-2/0/-2/135,-1/-1/I/-30,-1/0/-1/135}{\fill[black] (\x,\y) circle(1pt)+(\g:0.5) node{$\l$};}
\end{tikzpicture}}
\end{enumEX}
}	
\end{bt}

\begin{bt}%[2D1H5-7]
% nguon: Archive/2024/TaiLieuGiangDay/1.HamSo/1. K12 - 4in1 - Hamso/ThubaiGoc/Bai5-Dang3.tex #2146
%[Bài giảng 12 mới, Lê Hồng Phi]
Cho hàm số $y=\dfrac{x^2-2x-3}{-x+2}$ có đồ thị $(C)$. Tìm tất cả các điểm thuộc đồ thị $(C)$ có tọa độ là số nguyên.
\loigiai
{Tập xác định $\mathscr{D}=\mathbb{R}\setminus\{2\}$.\\
Ta có $y=-x-\dfrac{3}{-x+2}$ nên điểm $M(x_M;y_M)$ thuộc $(C)$ có tọa độ nguyên khi và chỉ khi \[\heva{& x_M\in\mathbb{Z} \\ & 3\ \vdots\ (-x_M+2)}\Leftrightarrow -x_M+2\in\{-3;-1;1;3\}\Leftrightarrow x_M\in\{-1;1;3;5\}.
\]
Vậy có $4$ điểm thuộc $(C)$ có tọa độ nguyên là $(-1;0)$, $(1;-4)$, $(3;0)$, $(5;-4)$.
}
\end{bt}

\begin{bt}%[2D1H5-7]
% nguon: Archive/2024/TaiLieuGiangDay/1.HamSo/1. K12 - 4in1 - Hamso/ThubaiGoc/Bai5-Dang3.tex #2147
%[Bài giảng 12 mới, Lê Hồng Phi]
Tìm tọa độ điểm $M$ mà đồ thị hàm số $y=\dfrac{mx^2-mx+1}{x-1}$ luôn đi qua với mọi $m$.
\loigiai
{Tập xác định $\mathscr{D}=\mathbb{R}\setminus\{1\}$.\\
Gọi $M(x_0;y_0)$ là điểm mà đồ thị hàm số đã cho luôn đi qua với mọi $m$.\\
Khi đó \begin{eqnarray*}
&&y_0=\dfrac{mx_0^2-mx_0+1}{x_0-1},\ \forall m\\
&\Leftrightarrow & \heva{& x_0\neq 1 \\ & y_0(x_0-1)=m\left(x_0^2-x_0\right)+1,\ \forall m}\\
&\Leftrightarrow & \heva{& x_0\neq 1 \\ & m\left(x_0^2-x_0\right)+1-y_0(x_0-1)=0,\ \forall m}\\
&\Leftrightarrow &\heva{& x_0\neq 1 \\ & x_0^2-x_0=0\\ &1-y_0(x_0-1)=0}\\
&\Leftrightarrow & \heva{& x_0\neq 1 \\ & \hoac{& x_0=1 \\ & x_0=0}\\ &y_0(x_0-1)=1}\\
&\Leftrightarrow &\heva{& x_0=0 \\ & y_0=-1.}
\end{eqnarray*}
Vậy điểm mà đồ thị hàm số đã cho luôn đi qua với mọi $m$ là $M(0;-1)$.
}
\end{bt}

\begin{bt}%[2D1V5-3]
% nguon: Archive/2024/TaiLieuGiangDay/1.HamSo/1. K12 - 4in1 - Hamso/ThubaiGoc/Bai5-Dang3.tex #2148
%[Bài giảng 12 mới, Lê Hồng Phi]
Tìm $m$ để phương trình $x^2+2x+5=\left(m^2+2m+5\right)(x+1)$ có hai nghiệm dương phân biệt.
\loigiai
{Nhận thấy $x=-1$ không phải là nghiệm của phương trình đã cho. Do đó
\[x^2+2x+5=\left(m^2+2m+5\right)(x+1)\Leftrightarrow\dfrac{x^2+2x+5}{x+1}=m^2+2m+5.\quad (*)\]
Ta có $(*)$ là phương trình hoành độ giao điểm giữa đồ thị $(C)$ của hàm số $y=\dfrac{x^2+2x+5}{x+1}$ và đường thẳng $d\colon y=m^2+2m+5$.\\
Hàm số $y=\dfrac{x^2+2x+5}{x+1}$ có tập xác định $\mathscr{D}=\mathbb{R}\setminus\{-1\}$ và $y'=\dfrac{x^2+2x-3}{(x+1)^2}$.\\
Phương trình $y'=0\Leftrightarrow x^2+2x-3=0\Leftrightarrow\hoac{& x=1 \\ & x=-3.}$\\
Bảng biến thiên
\begin{center}
\begin{tikzpicture}
\tkzTabInit[nocadre=false, lgt=1.2, espcl=2.5, deltacl=0.6]{$x$/0.6,$y'$/0.6,$y$/2}
{$-\infty$, $-3$, $-1$, $0$, $1$, $+\infty$}
\tkzTabLine {,+,0,-,d,,-,,0,+,}
\tkzTabVar{-/$-\infty$, +/$-4$, -D/$-\infty$ }
\path (N32)node(C)[below right]{$+\infty$} (N53)node(E)[above]{$4$}($(N32)!0.5!(N53)$)node(D)[]{$5$} (N62)node(F)[below]{$+\infty$};
\draw[-stealth] (C)--(D);
\draw[-stealth] (D)--(E);
\draw[-stealth] (E)--(F);
\end{tikzpicture}
\end{center}
Dựa vào bảng biến thiên, đường thẳng $d$ cắt đồ thị $(C)$ tại hai điểm phân biệt có hoành độ dương khi và chỉ khi \begin{eqnarray*}
4<m^2+2m+5<5 &\Leftrightarrow &\heva{& m^2+2m+5>4 \\ & m^2+2m+5<5}\\
&\Leftrightarrow & \heva{& m^2+2m+1>0 \\ & m^2+2m<0}\\
&\Leftrightarrow & \heva{& m\neq -1 \\ & -2<m<0.}
\end{eqnarray*}
Vậy tập hợp các giá trị $m$ cần tìm là $m\in (-2;0)\setminus\{-1\}$.
}
\end{bt}

\begin{bt}%[2D1N5-4]
% nguon: Archive/2024/TaiLieuGiangDay/1.HamSo/1. K12 - 4in1 - Hamso/ThubaiGoc/Bai5-Dang7.tex #2159

	Tìm giao điểm của đồ thị hàm số $y=\dfrac{2 x^2+1}{x-1}$ với đường thẳng $y=5x-1$.
	
	\loigiai{
		Tập xác đinh $\mathscr{D}=\mathbb{R}\setminus \{1\}$.
		Phương trình hoành độ giao điểm của $y=\dfrac{2 x^2+1}{x-1}$ và $y=5x-1$ là
		$$\dfrac{2 x^2+1}{x-1}=5x-1\Leftrightarrow \hoac{&x=0\\&x=2.}$$
		\begin{itemize}
			\item $x=0\Rightarrow y=-1.$
			\item $x=2\Rightarrow y=9.$ 
		\end{itemize}
	Vậy giao điểm của đồ thị hàm số $y=\dfrac{2 x^2+1}{x+1}$ với đường thẳng $y=5x-1$ là $(0;-1)$ và $(2;9).$
	}
\end{bt}

\begin{bt}%[2D1N5-4]
% nguon: Archive/2024/TaiLieuGiangDay/1.HamSo/1. K12 - 4in1 - Hamso/ThubaiGoc/Bai5-Dang7.tex #2160

	Cho hàm số $y=\dfrac{x-1}{x-2}$ có đồ thị $(C)$. Tìm số giao điểm của $(C)$ và trục hoành.
	\loigiai{
		Tập xác đinh $\mathscr{D}=\mathbb{R}\setminus \{2\}$.\\
		Phương trình hoành độ giao điểm của $(C)$ và trục hoành $\dfrac{x-1}{x-2}=0 \Leftrightarrow x=1.$\\
		Vậy số giao điểm của $(C)$ và trục hoành là $1$.}
\end{bt}

\begin{bt}%[2D1N5-4]
% nguon: Archive/2024/TaiLieuGiangDay/1.HamSo/1. K12 - 4in1 - Hamso/ThubaiGoc/Bai5-Dang7.tex #2161

	Đồ thị của hàm số $y=x^4-2x^2+2$ và đồ thị của hàm số $y=-x^2+4$ có tất cả bao nhiêu điểm chung?
	\loigiai{
		Số giao điểm của hai đồ thị chính bằng số nghiệm của phương trình hoành độ giao điểm của hai đồ thị hàm số.\\
		Ta có phương trình hoành độ giao điểm: $x^4-2x^2+2=-x^2+4 \Leftrightarrow x^4-x^2-2=0 \Leftrightarrow \hoac{&x=\sqrt{2}\\&x=-\sqrt{2}.}$\\
		Vậy có tất cả $2$ giao điểm cần tìm.}
\end{bt}

\begin{bt}%[2D1N5-4]
% nguon: Archive/2024/TaiLieuGiangDay/1.HamSo/1. K12 - 4in1 - Hamso/ThubaiGoc/Bai5-Dang7.tex #2162

	Tìm tọa độ giao điểm (\textbf{nếu có}) của hai đồ thị hàm số $y=3x+4$ và $y=x^3+2x+4$.
	\loigiai{
		Tọa độ giao điểm của hai đồ thị là nghiệm của hệ $\heva{&y=3x+4\\&y=x^3+2x+4.}$\\
		Giải hệ ta được $\heva{&x=0\\&y=4.}$\\
		Vậy tọa độ cần tìm là $(0;4)$.
	}
\end{bt}

\begin{bt}%[2D1H5-4]
% nguon: Archive/2024/TaiLieuGiangDay/1.HamSo/1. K12 - 4in1 - Hamso/ThubaiGoc/Bai5-Dang7.tex #2163

	Gọi $A$, $B$ là giao điểm của đồ thị hàm số $y=\dfrac{2x-1}{x+2}$ với đường thẳng $y=x-2$. Tính độ dài $AB$.
	\loigiai
	{Tập xác đinh $\mathscr{D}=\mathbb{R}\setminus \{-2\}$.\\
		Hoành độ giao điểm của hai đồ thị là nghiệm của phương trình
		\begin{center}
			$\dfrac{2x-1}{x+2}=x-2\Leftrightarrow \heva{&2x-1=(x-2)(x+2)\\&x\ne 2}\Leftrightarrow \heva{&x^2-2x-3=0\\&x\ne 2}\Leftrightarrow \hoac{&x=-1\\&x=3.}$
		\end{center}
		Khi $x=3\Rightarrow y=1$, $x=-1\Rightarrow y=-3$.\\
		Tọa độ $A(3;1)$, $B(-1;-3)$, $\overrightarrow{AB}=(-4;-4)\Rightarrow AB=\sqrt{4^2+(-4)^2}=4\sqrt{2}$.
	}
\end{bt}

\begin{bt}%[2D1H5-4]
% nguon: Archive/2024/TaiLieuGiangDay/1.HamSo/1. K12 - 4in1 - Hamso/ThubaiGoc/Bai5-Dang7.tex #2164

	Gọi $M$, $N$ là giao điểm của đường thẳng $(d)\colon y=4x-4$ và đường cong $(C)\colon y=\dfrac{2x^2+4}{x-1}$. Tìm hoành độ trung điểm $I$ của đoạn thẳng $MN$.
	\loigiai{
		Tập xác đinh $\mathscr{D}=\mathbb{R}\setminus \{1\}$.\\
		Phương trình hoành độ giao điểm
		\begin{eqnarray*}
			& &4x-4=\dfrac{2x^2+4}{x-1}\\
			&\Leftrightarrow&\hoac{&x=0\\&x=4.} 
		\end{eqnarray*}
		$\Rightarrow M\left(0;y_M\right)$, $N\left(4;y_N\right)$.\\
		Suy ra hoành độ trung điểm $I$ của đoạn $MN$ bằng $\dfrac{x_M+x_N}{2}=2$.
	}
\end{bt}

\begin{bt}%[2D1V5-4]
% nguon: Archive/2024/TaiLieuGiangDay/1.HamSo/1. K12 - 4in1 - Hamso/ThubaiGoc/Bai5-Dang7.tex #2165

	Cho hàm số $y=\dfrac{2x+1}{x+1}$ có đồ thị $(C)$ và đường thẳng $d\colon y=x+m$. Tìm giá trị của tham số $m$ để $d$ cắt $(C)$ tại hai điểm phân biệt $A, B$ sao cho $AB=\sqrt{10}$.
	\loigiai
	{Tập xác đinh $\mathscr{D}=\mathbb{R}\setminus \{-1\}$.\\
		Phương trình hoành độ giao điểm của đồ thị $(C)$ và đường thẳng $d$ là
		$$\dfrac{2x+1}{x+1}=x+m\Leftrightarrow\heva{&x\neq -1\\&x^2+\left(m-1\right)x+m-1=0 \quad(1)}$$
		Khi đó $d$ cắt $(C)$ tại hai điểm phân biệt $A, B$ khi và chỉ khi phương trình $(1)$ có hai nghiệm phân biệt khác $-1$ khi
		$$\heva{&\left(m-1\right)^2-4\left(m-1\right)>0\\&\left(-1\right)^2-\left(m-1\right)+m-1\neq 0}\Leftrightarrow \heva{&m<1\\&m>5\quad(*)}$$
		Ta có
		$$A\left(x_1;x_1+m\right),B\left(x_2;x_2+m\right)\Rightarrow \overline{AB}=\left(x_2-x_1;x_2-x_1\right)\Rightarrow AB=\sqrt{2{\left(x_2-x_1\right)}^2}=\sqrt{2}\left|{x_2-x_1}\right|$$
		Và $\heva{&x_1+x_2=1-m\\&x_1x_2=m-1.}$ Từ đây ta có
		\begin{eqnarray*}
			AB=\sqrt{10}
			&\Leftrightarrow  & \left|{x_2-x_1}\right|=\sqrt{5}\\
			&\Leftrightarrow & \left(x_2+x_1\right)^2-4x_1x_2=5\\
			&\Leftrightarrow & \left(1-m\right)^2-4\left(m-1\right)=5\\
			&\Leftrightarrow & m^2-6m=0\\
			&\Leftrightarrow & \hoac{&m=0\\&m=6}\left(\text{thỏa mãn}\right) (*).
		\end{eqnarray*}
	Vậy $\hoac{&m=0\\&m=6}$.
	}
\end{bt}

\begin{bt}%[2D1H5-4]
% nguon: Archive/2024/TaiLieuGiangDay/1.HamSo/1. K12 - 4in1 - Hamso/ThubaiGoc/Bai5-Dang7.tex #2169

		Cho đồ thị hàm số $y=\dfrac{x^2-4}{x+2}$ $(C)$. Tìm tọa độ giao điểm của $(C)$ với các trục tọa độ.
		\loigiai{
			Tập xác đinh $\mathscr{D}=\mathbb{R}\setminus \{-2\}$.\\
			\begin{itemize}
				\item Thay $x=0$ vào $y=\dfrac{x^2-4}{x+2}$, ta được $y=-2.$
				\item Thay $y=0$ vào $y=\dfrac{x^2-4}{x+2}$, ta được $0=\dfrac{x^2-4}{x+2}\Leftrightarrow\hoac{&x=2\\&x=-2.}$
			\end{itemize}
		Vậy tọa độ giao điểm của $(C)$ với trục $Ox$ là $(2;0)$ và $(-2;0)$; với trục $Oy$ là $(0;-2)$.
		}
	\end{bt}

\begin{bt}%[2D1H5-4]
% nguon: Archive/2024/TaiLieuGiangDay/1.HamSo/1. K12 - 4in1 - Hamso/ThubaiGoc/Bai5-Dang7.tex #2170

	Tìm tung độ giao điểm của đồ thị hàm số $y=-3x+4$ và đồ thị hàm số $y=x^3+2x+4$.
	\loigiai{
		Phương trình hoành độ giao điểm $x^3+2x+4 = -3x+4 \Leftrightarrow x\left(x^2+5\right)=0\Leftrightarrow x=0 \Rightarrow y=4.$
	}
\end{bt}

\begin{bt}%[2D1H5-4]
% nguon: Archive/2024/TaiLieuGiangDay/1.HamSo/1. K12 - 4in1 - Hamso/ThubaiGoc/Bai5-Dang7.tex #2171

	Gọi $A,B$ lần lượt là giao điểm của đồ thị hàm số $y=\dfrac{2x-3}{x+1}$ với các trục $Ox,Oy$. Tính diện tích tam giác $OAB$.
	\loigiai{
		Tập xác đinh $\mathscr{D}=\mathbb{R}\setminus \{-1\}$.\\
		\begin{itemize}
			\item Thay $y=0$ vào $y=\dfrac{2x-3}{x+1}$, ta được $x=\dfrac{3}{2}.$
			\item Thay $x=0$ vào $y=\dfrac{2x-3}{x+1}$, ta được $y=-3.$
		\end{itemize}
		Do đó $A\left(\dfrac{3}{2};0\right)$ và $B(0;-3)$.\\
		Vậy $S_{\Delta OAB}=\dfrac{OA\cdot OB}{2}=\dfrac{9}{4}$.
	}
\end{bt}

\begin{bt}%[2D1H5-4]
% nguon: Archive/2024/TaiLieuGiangDay/1.HamSo/1. K12 - 4in1 - Hamso/ThubaiGoc/Bai5-Dang7.tex #2172

	Cho hàm số $y=\dfrac{x+1}{x-2}$ có đồ thị $(C)$ và đường thẳng $d\colon y=2x+m$. Tìm $m$ để đường thẳng $d$ cắt $(C)$ tại hai điểm phân biệt.
	\loigiai{
		Tập xác đinh $\mathscr{D}=\mathbb{R}\setminus \{2\}$.\\
		Phương trình hoành độ giao điểm của $(C)$ và $d$ là
		$$\dfrac{x+1}{x-2}=2x+m\Leftrightarrow x+1=(2x+m)\cdot (x-2)\Leftrightarrow 2x^2+(m-5)x-2m-1=0 \hfill(*)$$
		$(C)$ cắt $d$ tại hai điểm phần biệt $\Leftrightarrow$ $(*)$ có hai nghiệm phân biệt khác $2$ khi 
		$$\heva{&\Delta >0\\&g(2)\ne 0}
		\Leftrightarrow m^2+6m+33>0 \Leftrightarrow m\in \mathbb{R}.$$
	}
\end{bt}

\begin{bt}%[2D1H3-1]
% nguon: Archive/2024/TaiLieuGiangDay/1.HamSo/1. K12 - 4in1 - Hamso/ThubaiGoc/Bai3-Dang1.tex #2185
[CKP]%[Bài giảng 12-4in1, Nhật Thiện]
	Tìm giá trị lớn nhất, giá trị nhỏ nhất của hàm số $y=f(x)=x^3-3x^2-9x+5$ trên đoạn $[0;5]$.
	\loigiai{
		Ta có $f'(x)=3x^2-6x-9$; $f'(x)=0\Leftrightarrow \hoac{&x=-1\notin (0;5)\\&x=3\in (0;5).}$\\
		Mà $f(0)=5$; $f(5)=10$; $f(3)=-22$.\\
		Vậy $\max\limits_{[0;5]}f(x)=10$, $\min\limits_{[0;5]}f(x)=-22$.
	}
\end{bt}

\begin{bt}%[2D1H3-1]
% nguon: Archive/2024/TaiLieuGiangDay/1.HamSo/1. K12 - 4in1 - Hamso/ThubaiGoc/Bai3-Dang1.tex #2186
[CTST]%[Bài giảng 12-4in1, Nhật Thiện]
	Tìm giá trị lớn nhất, giá trị nhỏ nhất của hàm số $f(x)=x^4-8 x^2+9$ trên đoạn $[-1 ; 3]$.
	\loigiai {\\
		Ta có $f'(x)=4 x^3-16 x$;
		$
		f'(x)=0 \Leftrightarrow x=0$ hoặc  $x=2$ hoặc $x=-2$ (loại vì không thuộc $[-1 ; 3]$);
		$f(-1)=2 ; f(0)=9 ; f(2)=-7 ; f(3)=18$.\\
		Vậy $\max \limits_{[-1 ; 3]} f(x)=f(3)=18$ và $\min\limits _{[-1 ; 3]} f(x)=f(2)=-7$.
	}
\end{bt}

\begin{bt}%[2D1H3-1]
% nguon: Archive/2024/TaiLieuGiangDay/1.HamSo/1. K12 - 4in1 - Hamso/ThubaiGoc/Bai3-Dang1.tex #2187
[CD]%[Bài giảng 12-4in1, Nhật Thiện]
	Tìm giá trị lớn nhất và giá trị nhỏ nhất của hàm số $f(x)=\dfrac{x^3}{3}-2x^2+3x+1$ trên đoạn $\left[-3;2\right]$
	\loigiai{
			Ta có $f'(x)=x^2-4x+3$.\\
				Cho $f'(x)=0 \Leftrightarrow x^2-4x+3=0 \Leftrightarrow \hoac{&x=1\in\left[-3;2\right]\\ &x=3\notin\left[-3;2\right].}$\\
				Các giá trị $f(1)=\dfrac{7}{3}, f(-3)=-35, f(2)=\dfrac{5}{3}$.\\
				Vậy
				$\begin{aligned}[t]
					&\max\limits_{\left[-3;2\right]}f(x)=\dfrac{7}{3} \text{ tại } x=1,\\
					&\min\limits_{\left[-3;2\right]}f(x)=-35 \text{ tại } x=-3.
				\end{aligned}$
	}
\end{bt}

\begin{bt}%[2D1H3-1]
% nguon: Archive/2024/TaiLieuGiangDay/1.HamSo/1. K12 - 4in1 - Hamso/ThubaiGoc/Bai3-Dang1.tex #2188
[CTST]%[Bài giảng 12-4in1, Nhật Thiện]
	Tìm giá trị lớn nhất, giá trị nhỏ nhất của hàm số $g(x)=x+\dfrac{4}{x^2}$ trên đoạn $[1 ; 4]$.
	\loigiai{Ta có $f'(x)=1-\dfrac{8}{x^3}$. Khi đó
		$f'(x)=0 \Leftrightarrow x=2$. \\Ta có
		$
		f(1)=5$; $f(2)=3$; $f(4)=\dfrac{17}{4}=4{,}25
		$.\\
		Vậy $\displaystyle\max _{[1 ; 4]} f(x)=f(1)=5$ và $\displaystyle\min _{[1 ; 4]} f(x)=f(2)=3$.}
\end{bt}

\begin{bt}%[2D1H3-1]
% nguon: Archive/2024/TaiLieuGiangDay/1.HamSo/1. K12 - 4in1 - Hamso/ThubaiGoc/Bai3-Dang1.tex #2189
%[Bài giảng 12-4in1, Nhật Thiện]
	Tìm giá trị lớn nhất và giá trị nhỏ nhất của hàm số $f(x)=\sin2x-2x$ trên đoạn $\left[\dfrac{\pi}{2};\dfrac{3\pi}{2}\right]$.
	\loigiai{
		Ta có $f'(x)=2\cos 2x-2>0$, $\forall x\in \left(\dfrac{\pi}{2};\dfrac{3\pi}{2}\right)$.\\
		Suy ra $\max\limits_{\left[\dfrac{\pi}{2};\dfrac{3\pi}{2}\right]} f(x)=-\pi$, $\min\limits_{\left[\dfrac{\pi}{2};\dfrac{3\pi}{2}\right]} f(x)=-3\pi$.
	}
\end{bt}

\begin{bt}%[2D1H3-1]
% nguon: Archive/2024/TaiLieuGiangDay/1.HamSo/1. K12 - 4in1 - Hamso/ThubaiGoc/Bai3-Dang1.tex #2192
%[Bài giảng 12-4in1, Nhật Thiện]
	Tìm giá trị lớn nhất và giá trị nhỏ nhất của hàm số $f(x)=x^3-\dfrac{3}{2}x^2$ trên đoạn $\left[-1;2\right]$
	\loigiai{
		Xét hàm số $f(x)=x^3-\dfrac{3}{2}x^2$, trên đoạn $\left[-1;2\right]$.\\
			Đạo hàm $f'(x)=3x^2-3x$.\\
			Cho $f'(x)=0 \Leftrightarrow 3x^2-3x=0 \Leftrightarrow \hoac{&x=1\in[-1;2]\\&x=0\in[-1;2].}$\\
			Các giá trị $f(-1)=-\dfrac{5}{2}$, $f(0)=0$, $f(1)=-\dfrac{1}{2}$, $f(2)=2$.\\
			So sánh các giá trị, ta có $\max\limits_{[-1;2]}f(x)=2$, $\min\limits_{[-1;2]}f(x)=-\dfrac{5}{2}$.
	}
\end{bt}

\begin{bt}%[2D1H3-1]
% nguon: Archive/2024/TaiLieuGiangDay/1.HamSo/1. K12 - 4in1 - Hamso/ThubaiGoc/Bai3-Dang1.tex #2193
%[Bài giảng 12-4in1, Nhật Thiện]
	Tìm giá trị lớn nhất và giá trị nhỏ nhất của hàm số $f(x)=x^4-2x^3+x^2+1$ trên đoạn $\left[-1;1\right]$.
	\loigiai{
			Xét hàm số $f(x)=x^4-2x^3+x^2+1$, trên đoạn $\left[-1;1\right]$.\\
			Đạo hàm $f'(x)=4x^3-6x^2+2x$.\\
			Cho $f'(x)=0 \Leftrightarrow 4x^3-6x^2+2x \Leftrightarrow \hoac{&x=1\in[-1;1]\\&x=\dfrac{1}{2}\in[-1;1]\\&x=0\in[-1;1].}$\\
			Các giá trị $f(-1)=5$, $f(0)=1$, $f\left(\dfrac{1}{2}\right)=\dfrac{17}{16}$, $f(1)=1$.\\
			So sánh các giá trị, ta có $\max\limits_{[-1;1]}f(x)=5$, $\min\limits_{[-1;1]}f(x)=1$.
	}
\end{bt}

\begin{bt}%[2D1H3-1]
% nguon: Archive/2024/TaiLieuGiangDay/1.HamSo/1. K12 - 4in1 - Hamso/ThubaiGoc/Bai3-Dang1.tex #2194
[CKP]%[Bài giảng 12-4in1, Nhật Thiện]
	Tìm giá trị lớn nhất và giá trị nhỏ nhất của hàm số $y=f(x)=\mathrm{e}^x-x$ trên đoạn $[-1 ; 2]$.
	\loigiai{
		Xét hàm số $y=f(x)=\mathrm{e}^x-x$ trên đoạn $[-1 ; 2]$;\\
		Ta có 
		$f'(x)=\mathrm{e}^x-1$; $f'(x)=0 \Leftrightarrow x=0\in(-1 ; 2)$.\\
		Mà $f(-1)=\mathrm{e}^{-1}+1 ; f(0)=1 ; f(2)=\mathrm{e}^2-2$.\\
		Vậy $\max\limits_{x \in [-1;2]} f(x)=f(2)=\mathrm{e}^2-2; \min\limits_{x \in [-1;2]} f(x)=f(0)=1$.
	}
\end{bt}

\begin{bt}%[2D1H3-1]
% nguon: Archive/2024/TaiLieuGiangDay/1.HamSo/1. K12 - 4in1 - Hamso/ThubaiGoc/Bai3-Dang1.tex #2195
%[Bài giảng 12-4in1, Nhật Thiện]
	Tìm giá trị lớn nhất và giá trị nhỏ nhất của hàm số $f(x)=\cos2x+2x+1$ trên đoạn $\left[\dfrac{-\pi}{2};\pi\right]$.
	\loigiai{
	Xét hàm số $f(x)=\cos2x+2x+1$ trên đoạn $\left[\dfrac{-\pi}{2};\pi\right]$ .\\
			Đạo hàm $f'(x)=-2\sin2x+2$.\\
			Cho $f'(x)=0 \Leftrightarrow -2\sin2x+2=0 \Leftrightarrow x=\dfrac{\pi}{4}+k\pi,\,k\in\mathbb{Z}$.\\
			Vì $x\in \left[-\dfrac{\pi}{2};\pi\right] \Rightarrow x=\dfrac{\pi}{4}$.\\
			Các giá trị $f(-\dfrac{\pi}{2})=-\pi$, $f\left(\dfrac{\pi}{4}\right)=\dfrac{\pi}{2}+1$, $f(1)=-\dfrac{1}{2}$, $f(2)=2$.\\
			So sánh các giá trị, ta có $\max\limits_{\left[-\tfrac{\pi}{2};\pi\right]}f(x)=2\pi+2$, $\min\limits_{\left[-\tfrac{\pi}{2};\pi\right]}f(x)=-\pi$.
	}
\end{bt}

\begin{bt}%[2D1H3-1]
% nguon: Archive/2024/TaiLieuGiangDay/1.HamSo/1. K12 - 4in1 - Hamso/ThubaiGoc/Bai3-Dang1.tex #2196
[CTST]%[Bài giảng 12-4in1, Nhật Thiện]
	Tìm giá trị lớn nhất, giá trị nhỏ nhất của hàm số $y=\dfrac{2 x+1}{x-2}$ trên đoạn $[3 ; 7]$.
	\loigiai{
			Ta có $f'(x)=\dfrac{-5}{(x-2)^2}<0,\quad \forall x\neq 2$.\\
			Vậy $\displaystyle\max _{[3 ; 7]} f(x)=f(3)=7$ và $\displaystyle\min _{[3 ; 7]} f(x)=f(7)=3$.
		}
\end{bt}

\begin{bt}%[2D1H3-1]
% nguon: Archive/2024/TaiLieuGiangDay/1.HamSo/1. K12 - 4in1 - Hamso/ThubaiGoc/Bai3-Dang1.tex #2197
[CTST]%[Bài giảng 12-4in1, Nhật Thiện]
	Tìm giá trị lớn nhất, giá trị nhỏ nhất của hàm số $y=\sin 2 x$ trên đoạn $\left[0 ; \dfrac{7 \pi}{12}\right]$.
	\loigiai{
	Ta có $f'(x)=2\cos2x$. Khi đó
			$f'(x)=0 \Leftrightarrow x=\dfrac{\pi}{4}$.
			\\ Ta có
			$
			f(0)=0 ; f\left( \dfrac{\pi}{4}\right) =1 ; f\left( \dfrac{7\pi}{12}\right) =\dfrac{-1}{2}.
			$\\
			Vậy $\displaystyle\max _{\left[ 0 ; \tfrac{7\pi}{12}\right] } f(x)=f\left( \dfrac{\pi}{4}\right) =1$ và $\displaystyle\min _{\left[ 0 ; \tfrac{7\pi}{12}\right] } f(x)=f(\sqrt[3]{2})=\dfrac{-1}{2}$.
}
\end{bt}

\begin{bt}%[2D1H3-1]
% nguon: Archive/2024/TaiLieuGiangDay/1.HamSo/1. K12 - 4in1 - Hamso/ThubaiGoc/Bai3-Dang1.tex #2198
[CKP]%[Bài giảng 12-4in1, Nhật Thiện]
	Tìm giá trị lớn nhất và giá trị nhỏ nhất của hàm số $y=f(x)=x+\dfrac{1}{x}-2$ trên khoảng $(-\infty ; 0)$.
	\loigiai{
		Xét hàm số $y=f(x)=x+\dfrac{1}{x}-2$ trên khoảng $(-\infty ; 0)$;\\
			Trên khoảng $(-\infty ; 0)$, ta có $f'(x)=1-\dfrac{1}{x^2}=\dfrac{x^2-1}{x^2}$.
			\begin{itemize}
				\item $f'(x)=0 \Leftrightarrow x=-1$.
				\item $\lim\limits_{x\to -\infty}f(x)=\lim\limits_{x\to -\infty}\left(x+\dfrac{1}{x}-2\right)=-\infty$.
				\item $\lim\limits_{x\to 0^-}f(x)=\lim\limits_{x\to 0^-}\left(x+\dfrac{1}{x}-2\right)=-\infty$.
			\end{itemize}
			Bảng biến thiên
			\begin{center}
				\begin{tikzpicture}
					\tkzTabInit[nocadre=false,lgt=1.2,espcl=3.5,deltacl=0.6]
					{$x$/1,$f'(x)$/1,$f(x)$/3}
					{$-\infty$,$-1$,$0$}
					\tkzTabLine{,+,0,-,} 
					\tkzTabVar{-/$-\infty$,+/$-4$,-/$-\infty$}
				\end{tikzpicture}
			\end{center}	
			Từ bảng biến thiên, ta có $f(x)\geq -4$. Suy ra, trên khoảng $(-\infty;0)$
			\begin{itemize}
				\item $\max\limits_{x \in (-\infty;0)} f(x)=f(-1)=-4$;
				\item không tồn tại giá trị nhỏ nhất.
			\end{itemize}
	}
\end{bt}

\begin{bt}%[2D1H3-1]
% nguon: Archive/2024/TaiLieuGiangDay/1.HamSo/1. K12 - 4in1 - Hamso/ThubaiGoc/Bai3-Dang1.tex #2199
[CKP]%[Bài giảng 12-4in1, Nhật Thiện]
	Tìm giá trị lớn nhất và giá trị nhỏ nhất của hàm số $y=f(x)=\sqrt{4-x^2}$.
	\loigiai{
		 Xét hàm số $y=f(x)=\sqrt{4-x^2}$;\\
		 Tập xác định của hàm số là $ \mathscr{D}=[-2;2]$\\
		 Ta có $f'(x)=\dfrac{-x}{\sqrt{4-x^2}}$; $f'(x)=0 \Leftrightarrow x=0\in(-2 ; 2)$.\\
		 Mà $f(-2)=0 ; f(0)=2 ; f(2)=0$.\\
		 Vậy $\max\limits_{x \in [-2;2]} f(x)=f(0)=2$; $\min\limits_{x \in [-2;2]} f(x)=f(2)=0$.
	}
\end{bt}

\begin{bt}%[2D1H3-1]
% nguon: Archive/2024/TaiLieuGiangDay/1.HamSo/1. K12 - 4in1 - Hamso/ThubaiGoc/Bai3-Dang1.tex #2200
%[Bài giảng 12-4in1, Nhật Thiện]
	Tìm giá trị lớn nhất và giá trị nhỏ nhất của hàm số $f(x)=\mathrm{e}^x\left(x^2-5x+7\right)$ trên đoạn $\left[0;3\right]$.
	\loigiai{
		Xét hàm số $f(x)=\mathrm{e}^x\left(x^2-5x+7\right)$, trên đoạn $\left[0;3\right]$.\\
		Đạo hàm $f'(x)=\mathrm{e}^x(x^2-3x+2)$.\\
		Cho $f'(x)=0 \Leftrightarrow \mathrm{e}^x(x^2-3x+2)=0 \Leftrightarrow \hoac{&x=2\in[0;3]\\&x=1\in[0;3].}$\\
		Các giá trị $f(0)=7$, $f(1)=3e$, $f(2)=\mathrm{e}^2$, $f(3)=\mathrm{e}^3$.\\
		So sánh các giá trị, ta có $\max\limits_{[0;3]}f(x)=\mathrm{e}^3$, $\min\limits_{[0;3]}f(x)=7$.
	}
\end{bt}

\begin{bt}%[2D1H3-1]
% nguon: Archive/2024/TaiLieuGiangDay/1.HamSo/1. K12 - 4in1 - Hamso/ThubaiGoc/Bai3-Dang1.tex #2201
[CKP]%[Bài giảng 12-4in1, Nhật Thiện]
	Tìm giá trị lớn nhất và giá trị nhỏ nhất của hàm số $y=f(x)=x \ln x$ trên đoạn $\left[\mathrm{e}^{-2} ; \mathrm{e}\right]$.
	\loigiai{
		Xét hàm số $y=f(x)=x \ln x$ trên đoạn $\left[\mathrm{e}^{-2} ; \mathrm{e}\right]$.\\
			Ta có $f'(x)=\ln x -1$; $f'(x)=0 \Leftrightarrow x=\mathrm{e}\notin \left[\mathrm{e}^{-2} ; \mathrm{e}\right]$.\\
			Mà $f(\mathrm{e}^{-2})=-2\mathrm{e}^{-2}$; $f(\mathrm{e})=3$; $f(\mathrm{e})=\mathrm{e}$.\\
			Vậy $\max\limits_{x \in [\mathrm{e}^{-2};\mathrm{e}]} f(x)=f(\mathrm{e})=\mathrm{e}$; $\min\limits_{x \in [\mathrm{e}^{-2} ; \mathrm{e}]} f(x)=f(\mathrm{e}^{-2})=-2\mathrm{e}^{-2}$.
	}
\end{bt}

\begin{bt}%[2D1N5-3]
% nguon: Archive/2024/TaiLieuGiangDay/1.HamSo/1. K12 - 4in1 - Hamso/ThubaiGoc/Bai5-Dang6.tex #2210

	\immini{
		Cho hàm số $y=f(x)=ax^3+bx^2+cx+d$ có đồ thị như hình vẽ bên. Tìm số nghiệm thực phân biệt của phương trình $f(x)=1$.}{
		\begin{tikzpicture}[>=stealth,x=1cm,y=1cm,scale=0.7]
			\draw[->,line width = 1pt] (-4,0)--(0,0) node[below left]{$O$}--(3,0) node[above]{$x$};
			\draw[->,line width = 1pt] (0,-5) --(0,2) node[left]{$y$};
			\foreach \x in {-2,,,}{\draw[-] (\x,-0.1)--(\x,0.1) node[above]{$\x$};}
			\foreach \y in {-4,,,,,}{\draw[-] (-0.1,\y)--(0.1,\y) node[below right]{$\y$};}
			\draw [line width = 1pt, domain=-3.12:1.19, samples=100] plot (\x, {(\x)^3+3*(\x)^2-4});
			\draw (1,0) node[above right]{$1$};
		\end{tikzpicture}
	}
	\loigiai{
		Số nghiệm của phương trình $f(x)=1$ là số giao điểm của đường thẳng $y=1$ với đồ thị hàm số $y=f(x)$.\\
		Đường thẳng $y=1$ qua điểm $(0;1)$ song song với $Ox$ nên cắt đồ thị hàm số $y=f(x)$ tại đúng $1$ điểm.\\
		Do đó phương trình $f(x)=1$ có đúng $1$ nghiệm.
	}
\end{bt}

\begin{bt}%[2D1N5-3]
% nguon: Archive/2024/TaiLieuGiangDay/1.HamSo/1. K12 - 4in1 - Hamso/ThubaiGoc/Bai5-Dang6.tex #2211

	Cho hàm số $f(x)$ có bảng biến thiên như sau
	\begin{center}
		\begin{tikzpicture}[>=stealth,scale=1]
			\tkzTabInit[lgt=1.2,espcl=3]
			{$x$ /1.2, $f’(x)$ /1.2, $f(x)$ /2.5}
			{$-\infty$,$-1$,$1$,$+\infty$}
			\tkzTabLine{ ,+,0,-,0,+, }
			\tkzTabVar{-/$-\infty$,+/$4$,-/$0$,+/$+\infty$}
		\end{tikzpicture}
	\end{center}
	Xác định số nghiệm của phương trình $f(x) + 7 = 0$.
	\loigiai{
		Ta nhận thấy hàm số nhận giá trị $-7$ tại đúng một điểm trong mỗi khoảng $(-\infty;-1)$ và $(1;+\infty )$ nhưng không nhận giá trị $-7$ trong $(-1;1)$. Do đó, phương trình $f(x) + 7 = 0$ có $2$ nghiệm.
	}
\end{bt}

\begin{bt}%[2D1N5-3]
% nguon: Archive/2024/TaiLieuGiangDay/1.HamSo/1. K12 - 4in1 - Hamso/ThubaiGoc/Bai5-Dang6.tex #2212

	Hãy xác định tọa độ giao điểm của đồ thị $\left(C\right)\colon y=\dfrac{2x-1}{x+1}$ cắt đường thẳng $d \colon y=2x-3$.
	\loigiai{
		Phương trình hoành độ giao điểm: $\dfrac{2x-1}{x+1}=2x-3\Leftrightarrow\heva{&x\ne-1\\ &2x^2-3x-2=0}\Leftrightarrow\hoac{&x=2\\ &x=-\dfrac 12.}$\\
		Với $ x=2\Rightarrow y=1$.\\
		Với $ x=- \dfrac{1}{2}\Rightarrow y=-4$.\\
		Vậy đồ thị $\left(C\right)\colon y=\dfrac{2x-1}{x+1}$ cắt đường thẳng $d \colon y=2x-3$ tại các điểm là $\left(2;1\right)$, $\left(- \dfrac{1}{2}; - 4\right)$.
	}
\end{bt}

\begin{bt}%[2D1N5-3]
% nguon: Archive/2024/TaiLieuGiangDay/1.HamSo/1. K12 - 4in1 - Hamso/ThubaiGoc/Bai5-Dang6.tex #2213

	\immini{Cho hàm số $y=f(x)$ có đồ thị như đường cong hình bên. Phương trình $3f(x)+6=2f(x)+8$ có bao nhiêu nghiệm?
	}
	{\begin{tikzpicture}[line width=0.6pt,line cap=round,line join=round,>=stealth,smooth,samples=300,scale=0.8]
			\draw[->] (-2.5,0)--(0,0)%
			node[below right]{$O$}--(2.5,0) node[below]{\footnotesize $x$};
			\draw[->] (0,-0.8) --(0,3.5) node[right]{\footnotesize $y$};
			\foreach \x in {}{
				\draw (\x,0) node[below]{$\x$};%Ox
			}
			\foreach \y in {1,2}{
				\draw (0,\y) node[left]{\footnotesize $\y$};%Oy
			}
			\draw[dashed,line width=0.1pt,color=black!30!] (-2.5,-0.8) grid (2.5,3.5);
			\draw [black, thick, domain=-1.505:2.2, samples=100] %
			plot (\x, {(((\x)+1))^2*((\x)-1)*((\x)-2)+1});
			\fill[fill=black,draw=black]
			(0,1) circle (1pt)
			(0,0) circle (1pt)
			(0,2) circle (1pt)
			(0,3) circle (1pt);
	\end{tikzpicture}}
	\loigiai{
		Ta có $3f(x)+6=2f(x)+8\Leftrightarrow f(x)=2$.\\
		Số nghiệm phương trình $f(x)=2$ là số giao điểm của đồ thị hàm số $y=f(x)$ và đường thẳng $y=2$. \\
		Dựa vào đồ thị suy ra phương trình $f(x)=2$ có $4$ nghiệm phân biệt.}
\end{bt}

\begin{bt}%[2D1N5-3]
% nguon: Archive/2024/TaiLieuGiangDay/1.HamSo/1. K12 - 4in1 - Hamso/ThubaiGoc/Bai5-Dang6.tex #2214

	Cho hàm số $y=f(x)$ xác định trên $\mathbb{R}\setminus\{2\}$, liên tục trên mỗi khoảng xác định và có bảng biến thiên như hình vẽ sau. Tìm tập hợp tất cả các giá trị của tham số thực $m$ sao cho phương trình  $f(x)=m$ có ba nghiệm phân biệt.
	\begin{center}
		\begin{tikzpicture}[>=stealth]
			\tkzTabInit[lgt=1.2,espcl=3]
			{$x$ /1,$y'$ /1,$y$ /2.5}
			{$-\infty$,$2$,$3$,$+\infty$}
			\tkzTabLine{t,-,d,+,z,-,}%
			\tkzTabVar{+/ $+\infty$ , -D-/$2$ /$-\infty$, +/ $3$, -/ $-\infty$}
		\end{tikzpicture}
	\end{center}
	\loigiai{
		Phương trình $f(x)=m$ có $3$ nghiệm phân biệt khi và chỉ khi đồ thị $y=f(x)$ và đường thẳng $y=m$ cắt nhau tại $3$ điểm phân biệt.\\
		Dựa vào bảng biến thiên ta thấy điều này tương đương với $2<m<3$.}
\end{bt}

\begin{bt}%[2D1H5-3]
% nguon: Archive/2024/TaiLieuGiangDay/1.HamSo/1. K12 - 4in1 - Hamso/ThubaiGoc/Bai5-Dang6.tex #2215

	Tìm tất cả các giá trị thực của $m$ để phương trình $x+1=m\sqrt{2x^2+1}$ có hai nghiệm phân biệt.
	\loigiai{
		Phương trình $x+1=m\sqrt{2x^2+1}\Leftrightarrow m=\dfrac{x+1}{\sqrt{2x^2+1}}$.\\
		Xét hàm số $y=\dfrac{x+1}{\sqrt{2x^2+1}}\Rightarrow y'=\dfrac{1-2x}{\sqrt{2x^2+1}}\Rightarrow y'=0\Leftrightarrow x=\dfrac{1}{2}$.\\
		Bảng biến thiên
		\begin{center}
			\begin{tikzpicture}
				\tkzTabInit[nocadre,lgt=2,espcl=4]{$x$/0.7, $y'$/0.7, $y$/2.3}{$-\infty$,$\tfrac{1}{2}$,$+\infty$}
				\tkzTabLine{,+,0,-,}
				\tkzTabVar{-/$-\tfrac{\sqrt{2}}{2}$,+/$\tfrac{\sqrt{6}}{2}$,-/$\tfrac{\sqrt{2}}{2}$}
			\end{tikzpicture}
		\end{center}
		Dựa vào bảng biến thiên ta thấy phương trình có hai nghiệm phân biệt $\Leftrightarrow \dfrac{\sqrt{2}}{2}<m<\dfrac{\sqrt{6}}{2}$.
	}
\end{bt}

\begin{bt}%[2D1N5-3]
% nguon: Archive/2024/TaiLieuGiangDay/1.HamSo/1. K12 - 4in1 - Hamso/ThubaiGoc/Bai5-Dang6.tex #2218

	\immini{Cho hàm số $y=f(x)$ có đồ thị như đường cong hình bên. Phương trình $f(x)=2$ có bao nhiêu nghiệm?}
	{\begin{tikzpicture}[line width=0.6pt,line cap=round,line join=round,>=stealth,smooth,samples=300,scale=0.8]
			\draw[->] (-2.5,0)--(0,0)%
			node[below right]{$O$}--(2.5,0) node[below]{\footnotesize $x$};
			\draw[->] (0,-0.8) --(0,3.5) node[right]{\footnotesize $y$};
			\foreach \x in {}{
				\draw (\x,0) node[below]{$\x$};%Ox
			}
			\foreach \y in {1,2}{
				\draw (0,\y) node[left]{\footnotesize $\y$};%Oy
			}
			\draw[dashed,line width=0.1pt,color=black!30!] (-2.5,-0.8) grid (2.5,3.5);
			\draw [black, thick, domain=-1.505:2.2, samples=100] %
			plot (\x, {(((\x)+1))^2*((\x)-1)*((\x)-2)+1});
			\fill[fill=black,draw=black]
			(0,1) circle (1pt)
			(0,0) circle (1pt)
			(0,2) circle (1pt)
			(0,3) circle (1pt);
	\end{tikzpicture}}
	\loigiai{
		Số nghiệm phương trình $f(x)=2$ là số giao điểm của đồ thị hàm số $y=f(x)$ và đường thẳng $y=2$. \\
		Dựa vào đồ thị suy ra phương trình $f(x)=2$ có $4$ nghiệm phân biệt.}
\end{bt}

\begin{bt}%[2D1N5-3]
% nguon: Archive/2024/TaiLieuGiangDay/1.HamSo/1. K12 - 4in1 - Hamso/ThubaiGoc/Bai5-Dang6.tex #2219

	\immini{Cho hàm số $ y=f(x) $ có đồ thị như hình vẽ bên. Có bao nhiêu giá trị nguyên của tham số $ m $ để phương trình $ f(x)=m $ có $ 3 $ nghiệm phân biệt.
	}{
		\begin{tikzpicture}[>=stealth,scale=0.6]
			\draw[->,black] (-2.5,0) -- (3,0)node[above left] {$x$};
			\draw[dashed] (-1,0)--(-1,-2) (-1,-2)--(0,-2);
			\draw[dashed] (2,0)--(2,1) (2,1)--(0,1);
			\draw[->,black] (0,-2.5) -- (0,3.1)node[right] {$y$};
			\node at (1,0) [below] {\footnotesize $ 1 $};
			\node at (0,1) [left] {\footnotesize $ 1 $};
			\node at (0,-2) [right] {\footnotesize $ -2 $};
			\node at (0,1) [left] {\footnotesize $ 1 $};
			\node at (0,0) [below left] {\footnotesize $O$};
			\draw[smooth,samples=100,domain=-1.05:2.3] plot(\x,{(\x)^3-2*(\x)^2+1});
		\end{tikzpicture}
	}
	\loigiai{
		Số nghiệm của phương trình $ f(x)=m $ là số giao điểm của hai đồ thị hàm số $(C)\colon y=f(x) $ với đường thẳng $ y=m $.\\
		Dựa vào đồ thị hàm số $ y=f(x) $, phương trình $ f(x)=m $ có $ 3 $ nghiệm phân biệt khi chỉ khi $ m $ nhận giá trị nguyên bằng $ 0 $.
	}
\end{bt}

\begin{bt}%[2D1H5-3]
% nguon: Archive/2024/TaiLieuGiangDay/1.HamSo/1. K12 - 4in1 - Hamso/ThubaiGoc/Bai5-Dang6.tex #2220

	Cho đồ thị hàm số $y=f(x)$ như hình bên dưới. Tìm tất cả các giá trị của tham số $m$ để phương trình $f(x)-m=0$ có $4$ nghiệm phân biệt.
	\begin{center}
		\begin{tikzpicture}[line join=round, line cap = round, >=stealth, scale=1.2,font=\footnotesize,transform shape]
			\pgfmathsetmacro\a{sqrt(2)}
			\draw[->] (-2.2,0) -- (2.2,0)node[above]{$x$};
			\draw[->] (0,-4.5) -- (0,1.2)node[right]{$y$};
			%\draw[-] (-1.5,-3.5) -- (1.5,-3.5) node[right]{$y=m$};
			\draw[fill=black] 
			(0,0) circle(1pt) node[below left]{$O$} 
			(0,-3) circle(1pt) node[above right]{$-3$}
			(0,-4) circle(1pt) node[below right]{$-2$}
			(-1,0) circle(1pt) node[above]{$-1$}
			(1,0) circle(1pt) node[above]{$1$}
			;
			\draw[dashed] 
			(-1,0)--(-1,-4)--(1,-4)--(1,0)
			;
			\draw[smooth,samples=100,domain=-1.8:1.8] plot(\x,{(\x)^4-2*(\x)^2-3});
		\end{tikzpicture}
	\end{center}
	\loigiai{
		Ta có $f(x)-m=0 \Leftrightarrow f(x)=m$. \quad $(*)$\\
		Phương trình $(*)$ có đúng $4$ nghiệm phân biệt khi và chỉ khi đường thẳng $y=m$ cắt đồ thị hàm số $y=f(x)$ tại đúng $4$ điểm phân biệt.\\	
		Từ đồ thị ta thấy điều này xảy ra khi và chỉ khi $-4<m<-3$.\\
		Vậy phương trình đã cho có đúng $4$ nghiệm phân biệt khi và chỉ khi $-4<m<-3$.
	}
\end{bt}

\begin{bt}%[2D1V5-3]
% nguon: Archive/2024/TaiLieuGiangDay/1.HamSo/1. K12 - 4in1 - Hamso/ThubaiGoc/Bai5-Dang6.tex #2221

	Cho đồ thị hàm số $y=f(x)$ như hình bên dưới. Hỏi phương trình $f(x)=x$ có bao nhiêu nghiệm?
	\begin{center}
		\begin{tikzpicture}[>=stealth,line join=round,line cap=round,font=\footnotesize,scale=.6]
			\tikzset{declare function={
					x1=-4.5;x2=4.5;y1=-3.5;y2=3.5;% Miền hình vẽ
					a=0;b=1;c=0; m=1;n=-1;
					f(\x)=(a*(\x)^2+b*(\x)+c)/(m*(\x)^2+n);	
					a1=a/m;b1=(b-a*n/m)/m;
					g(\x)=a1*\x+b1;
					q1=-(n/m)-.1;q2=-n/m+.1;
					a2=1;b2=0;
					h(\x)=a2*\x+b2;
			}}
			\begin{scope}
				\clip (x1,y1)rectangle(x2,y2);
				\draw [smooth,samples=300,domain=x1:q1]plot(\x,{f(\x)});
				\draw [smooth,samples=300,domain=q2:x2]plot(\x,{f(\x)});
				\draw [smooth,samples=100,domain=x1:x2,dashed]plot(\x,{g(\x)});
				%				\draw[smooth,samples=100,domain=x1:x2]plot(\x,{h(\x)});
				\draw[dashed] (-n/m,y1)--(-n/m,y2);
			\end{scope}
			\draw[->] (0,y1)--(0,y2)node[left]{$y$};
			\draw[->] (x1,0)--(x2,0)node[above]{$x$};
			\filldraw(0,0)circle(1.2pt)node[below left]{$O$};
			\draw[dashed]
			;
		\end{tikzpicture}
		
	\end{center}
	\loigiai{
		Vẽ đồ thị của hàm số $y=f(x)$ và $y=x$ trên cùng một đồ thị như hình dưới đây.
		\begin{center}
			\begin{tikzpicture}[>=stealth,line join=round,line cap=round,font=\footnotesize,scale=.6]
				\tikzset{declare function={
						x1=-4.5;x2=4.5;y1=-3.5;y2=3.5;% Miền hình vẽ
						a=0;b=1;c=0; m=1;n=-1;
						f(\x)=(a*(\x)^2+b*(\x)+c)/(m*(\x)^2+n);	
						a1=a/m;b1=(b-a*n/m)/m;
						g(\x)=a1*\x+b1;
						q1=-(n/m)-.1;q2=-n/m+.1;
						a2=1;b2=0;
						h(\x)=a2*\x+b2;
				}}
				\begin{scope}
					\clip (x1,y1)rectangle(x2,y2);
					\draw [smooth,samples=300,domain=x1:q1]plot(\x,{f(\x)});
					\draw [smooth,samples=300,domain=q2:x2]plot(\x,{f(\x)});
					\draw [smooth,samples=100,domain=x1:x2,dashed]plot(\x,{g(\x)});
					\draw[smooth,samples=100,domain=x1:x2]plot(\x,{h(\x)});
					\draw[dashed] (-n/m,y1)--(-n/m,y2);
				\end{scope}
				\draw[->] (0,y1)--(0,y2)node[left]{$y$};
				\draw[->] (x1,0)--(x2,0)node[above]{$x$};
				\filldraw(0,0)circle(1.2pt)node[below left]{$O$};
				\draw[dashed]
				;
			\end{tikzpicture}
			
		\end{center}
		Nhìn vào đồ thị ta thấy đường thẳng $y=x$ cắt đồ thị hàm số $y=f(x)$ tại $3$ điểm.\\
		Vậy phương trình đã cho có $3$ nghiệm.
	}
\end{bt}

\begin{bt}%[2D1V5-3]
% nguon: Archive/2024/TaiLieuGiangDay/1.HamSo/1. K12 - 4in1 - Hamso/ThubaiGoc/Bai5-Dang6.tex #2222

	Cho bảng biến thiên của hàm số $y=f(x)=ax^3+bx^2+cx+d$. Tìm tất cả các giá trị $m$ để phương trình $|f(x)|=m$ có $4$ nghiệm phân biệt thỏa mãn $x_1<x_2<x_3<\dfrac{1}{2}<x_4$?
	\begin{center}
		\begin{tikzpicture}[>=stealth]
			\tkzTabInit[nocadre=false,lgt=1,espcl=2,deltacl=0.5]{$x$/.7 ,$y'$/.7,$y$/2}
			{$-\infty$ , $0$ , $1$ , $+\infty$}
			\tkzTabLine{ , + , $0$ , - , $0$ , + , }
			\tkzTabVar{-/$-\infty$ , +/$1$ , -/$0$ , +/$+\infty$}
		\end{tikzpicture}	
	\end{center}
	\loigiai{
		\immini{
			Từ bảng biến thiên ta có\\ $\heva{&d=1\\&a+b+c+d=0\\&c=0\\&3a+2b+c=0.}\Leftrightarrow \heva{&a=2\\&b=-3\\&c=0\\&d=1.}$\\
			Vậy hàm số $y=f(x)$ có phương trình $y=2x^3-3x^2+1$.
			Từ đó suy ra hàm số $y=|f(x)|$ có đồ thị như hình vẽ.\\
			Để phương trình $|f(x)|=m$ có $4$ nghiệm phân biệt thỏa mãn $x_1<x_2<x_3<\dfrac{1}{2}<x_4$ thì $\dfrac{1}{2}<m<1$.
		}
		{
			% Đồ thị hàm y=ax^3+bx^2+cx+d. Nếu hệ số lớn cần điều chỉnh hệ trục, vùng lưới, domain và lệnh \clip
			\begin{tikzpicture}[>=stealth,x=1cm,y=1cm,scale=1]
				\def\a{2} % Hệ số a phải khác 0
				\def\b{-3}
				\def\c{0}
				\def\d{1}
				\draw[->] (-2,0) -- (3,0)node[below]{\scriptsize $x$};
				\draw[->] (0,-1) -- (0,3.5) node[left] {\scriptsize $y$};
				\draw (0,0)node[below left]{\scriptsize $O$};
				\draw (0,1)node[above left]{\scriptsize $1$};
				\draw[dashed] (.5,0)node[below]{\scriptsize $\frac{1}{2}$}--(.5,.5)--(0,.5)node[left]{\scriptsize $\frac{1}{2}$};
				\clip (-3,0)rectangle(3,3);
				\draw[thick,samples=150,smooth,domain=-5:5] plot(\x,{\a*(\x)^3+(\b)*(\x)^2+(\c)*\x+(\d)});
				\draw[thick,samples=150,smooth,domain=-5:5] plot(\x,{-2*(\x)^3+3*(\x)^2-1});
				\fill (.5,.5) circle (1.5pt) (0,1) circle (1.5pt) ;
			\end{tikzpicture}
		}	
	}
\end{bt}

\begin{bt}%[2D1N5-4]
% nguon: Archive/2024/TaiLieuGiangDay/1.HamSo/1. K12 - 4in1 - Hamso/ThubaiGoc/Bai5-Dang4.tex #2233
[Mức độ 1]%[Bài giảng 12 mới, Lê Hồng Phi]
Tìm tọa độ giao điểm của đồ thị hàm số $y=3^{-x}$ với trục tung.
\loigiai
{Gọi $M(x_M;y_M)$ là giao điểm của đồ thị hàm số đã cho với trục tung.\\
Ta có $M\in Oy$ nên $x_M=0$.\\
Khi đó $y_M=3^{-0}=1$.\\
Vậy tọa độ giao điểm cần tìm là $M(0;1)$.
}
\end{bt}

\begin{bt}%[2D1H5-2]
% nguon: Archive/2024/TaiLieuGiangDay/1.HamSo/1. K12 - 4in1 - Hamso/ThubaiGoc/Bai5-Dang4.tex #2234
[Mức độ 1]%[Bài giảng 12 mới, Lê Hồng Phi]
Đồ thị hàm số $y=\log_2 x$ và đường thẳng $x=-1$ cắt nhau tại bao nhiêu điểm?
\loigiai
{Đồ thị hàm số $y=\log_2 x$  nằm phía bên phải trục $Oy$.\\
Đường thẳng $x=-1$ nằm bên trái trục $Oy$ nên đồ thị hàm số đã cho và đường thẳng $x=-1$ không cắt nhau.
}
\end{bt}

\begin{bt}%[2D1H5-2]
% nguon: Archive/2024/TaiLieuGiangDay/1.HamSo/1. K12 - 4in1 - Hamso/ThubaiGoc/Bai5-Dang4.tex #2235
[Mức độ 2]%[Bài giảng 12 mới, Lê Hồng Phi]
Tìm phương trình đường tiệm cận đứng của đồ thị hàm số $y=\log_2 (x-3)$.
\loigiai
{Đồ thị hàm số $y=f(x)=\log_2 x$ có tiệm cận đứng là $x=0$.\\
Đồ thị hàm số $y=\log_2(x-3)=f(x-3)$ có được bằng cách tịnh tiến đồ thị hàm số $y=\log_2 x$ sang phải $3$ đơn vị nên có phương trình đường tiệm cận đứng là $x=3$.
}
\end{bt}

\begin{bt}%[2D1H5-1]
% nguon: Archive/2024/TaiLieuGiangDay/1.HamSo/1. K12 - 4in1 - Hamso/ThubaiGoc/Bai5-Dang4.tex #2236
[Mức độ 2]%[Bài giảng 12 mới, Lê Hồng Phi]
Khảo sát và vẽ đồ thị hàm số $y=2^x$.
\loigiai
{\begin{enumEX}{1}
\item Tập xác định $\mathscr{D}=\mathbb{R}$.
\item Sự biến thiên
\begin{itemize}
\item Chiều biến thiên:\\
Ta có $y'=2^x\ln 2>0$, $\forall x\in\mathbb{R}$.\\
Hàm số đồng biến trên $\mathbb{R}$.
\item Các giới hạn và tiệm cận:\\
$\lim\limits_{x\to -\infty}y=0$ nên đường thẳng $y=0$ là tiệm cận ngang,\\
$\lim\limits_{x\to +\infty}y=+\infty$.
\item Bảng biến thiên
\begin{center}
\begin{tikzpicture}
\tkzTabInit[nocadre=false, lgt=1.2, espcl=2.5, deltacl=0.6]{$x$/0.6,$y'$/0.6,$y$/2}
{$-\infty$,  $+\infty$}
\tkzTabLine {,+,}
\tkzTabVar{-/$0$, +/$+\infty$}
\end{tikzpicture}
\end{center}
\end{itemize}
\item Đồ thị
\immini{Đồ thị hàm số nằm phía trên trục $Ox$.\\
Đồ thị cắt trục tung tại $(0;1)$ và đi qua điểm $(1;2)$.}{\begin{tikzpicture}[scale=1, font=\footnotesize, line join=round, line cap=round,>=stealth,x=0.6cm,y=0.6cm]
\def \xmin{-3.5};
\def \xmax{2.7};
\def \ymin{-1.0};
\def \ymax{4.0};
\draw[->] (\xmin, 0.) -- (\xmax,0.) node[anchor=north] {$x$};
\draw[->] (0.,\ymin) -- (0.,\ymax) node[anchor=west] {$y$};
\clip(\xmin,\ymin+0.1) rectangle (\xmax-0.1,\ymax-0.1);
\draw[smooth,samples=100,domain=\xmin:\xmax] plot(\x,{2^(\x)});
\draw[dashed] (1,0) node[below]{$1$} |- (0,2)node[left]{$2$};
\foreach \x/\y/\l/\g  in {0/0/O/225,1/2//0,0/1/1/150,1/0//0,0/2//0}{\fill[black] (\x,\y) circle(1pt)+(\g:0.5) node{$\l$};}
\end{tikzpicture}}
\end{enumEX}
}
\end{bt}

\begin{bt}%[2D1H5-1]
% nguon: Archive/2024/TaiLieuGiangDay/1.HamSo/1. K12 - 4in1 - Hamso/ThubaiGoc/Bai5-Dang4.tex #2237
[Mức độ 2]%[Bài giảng 12 mới, Lê Hồng Phi]
Khảo sát và vẽ đồ thị hàm số $y=\log_{\tfrac{1}{3}} x$.
\loigiai
{\begin{enumEX}{1}
\item Tập xác định $\mathscr{D}=(0;+\infty)$.
\item Sự biến thiên
\begin{itemize}
\item Chiều biến thiên:\\
Ta có $y'=\dfrac{1}{x\ln\tfrac{1}{3}}<0$, $\forall x\in\mathscr{D}$.\\
Hàm số nghịch biến trên $(0;+\infty)$.
\item Các giới hạn và tiệm cận:\\
$\lim\limits_{x\to 0^+}y=+\infty$ nên đường thẳng $x=0$ là tiệm cận đứng,\\
$\lim\limits_{x\to +\infty}y=-\infty$.
\item Bảng biến thiên
\begin{center}
\begin{tikzpicture}
\tkzTabInit[nocadre=false, lgt=1.2, espcl=2.5, deltacl=0.6]{$x$/0.6,$y'$/0.6,$y$/2}
{$0$,  $+\infty$}
\tkzTabLine {,-,}
\tkzTabVar{+/$+\infty$, -/$-\infty$}
\end{tikzpicture}
\end{center}
\end{itemize}
\item Đồ thị
\immini{Đồ thị hàm số nằm phía bên phải $Oy$.\\
Đồ thị cắt trục hoành tại $(1;0)$ và đi qua điểm $(3;-1)$.}{\begin{tikzpicture}[scale=1, font=\footnotesize, line join=round, line cap=round,>=stealth,x=0.6cm,y=0.6cm]
\def \xmin{-1.5};
\def \xmax{4.7};
\def \ymin{-2.0};
\def \ymax{3.2};
\draw[->] (\xmin, 0.) -- (\xmax,0.) node[anchor=north] {$x$};
\draw[->] (0.,\ymin) -- (0.,\ymax) node[anchor=east] {$y$};
\clip(\xmin,\ymin+0.1) rectangle (\xmax-0.1,\ymax-0.1);
\draw[smooth,samples=100,domain=0.01:\xmax] plot(\x,{ln(\x)/ln(1/3)});
\draw[dashed] (3,0) node[above]{$3$} |- (0,-1)node[left]{$-1$};
\foreach \x/\y/\l/\g  in {0/0/O/225,3/-1//0,1/0/1/90,3/0//0,0/-1//0}{\fill[black] (\x,\y) circle(1pt)+(\g:0.5) node{$\l$};}
\end{tikzpicture}}
\end{enumEX}
}
\end{bt}

\begin{bt}%[2D1H5-4]
% nguon: Archive/2024/TaiLieuGiangDay/1.HamSo/1. K12 - 4in1 - Hamso/ThubaiGoc/Bai5-Dang4.tex #2238
[Mức độ 2]%[Bài giảng 12 mới, Lê Hồng Phi]
\immini{Hàm số $y=\log_a x$ và $y=\log_b x$ có đồ thị như hình vẽ bên. Đường thẳng $y=3$ cắt hai đồ thị tại các điểm có hoành độ $x_1$, $x_2$. Biết rằng $x_2=2x_1$, tính giá trị của $\dfrac{a}{b}$.}
{				\begin{tikzpicture}[scale=0.8,font=\footnotesize, line join=round, line cap=round, >=stealth]
			\draw[->] (-1,0)--(5,0) node[below]{$x$};
			\draw[->] (0,-2.5)--(0,3.5) node[right]{$y$};
			\draw (0,0) node[below left]{$O$};
			\fill(0,0)circle (1pt);
			\draw[line width=0.8pt,color=violet,domain=0.15:4,smooth,samples=300,]plot(\x,{ln((\x))/ln(3)});
			\draw[line width=0.8pt,color=blue,domain=0.4:3.2,smooth,samples=300,]plot(\x,{ln((\x))/ln(1.5)});
			\draw(1,0) circle(1pt)(1.5,1) circle(1pt) (3,1) circle(1pt);
			\draw(0,1)circle(1pt)node[left]{\small$3$}--(4,1);  \draw[dashed](1.5,1)--(1.5,0)circle(1pt)node[below]{\small$x_1$};
				\fill(1.5,1)circle (1pt);
				\fill(3,1)circle (1pt);
			 \draw[dashed](3,1)--(3,0)circle(1pt)node[below]{\small$x_2$};
			\node[right] at (3,1.5) {\small$y=\log_ax$};
			\node[right] at (3,2.7) {\small$y=\log_bx$};
			%\node[right] at (0.2,2.4) {\small$y=\log_cx$};
			%\node[above] at (1,0) {\small$1$};
		\end{tikzpicture}
}
\loigiai{
Ta có $\heva{& \log_b x_1=3 \Leftrightarrow x_1=b^3 \\ &  \log_a x_2=3 \Leftrightarrow x_2=a^3.}$\\
Do $x_2=2x_1$ nên $a^3=2 b^3 \Leftrightarrow \dfrac{a}{b}=\sqrt[3]{2}$.\\
Vậy $\dfrac{a}{b}=\sqrt[3]{2}$.
}
\end{bt}

\begin{bt}%[2D1H5-4]
% nguon: Archive/2024/TaiLieuGiangDay/1.HamSo/1. K12 - 4in1 - Hamso/ThubaiGoc/Bai5-Dang4.tex #2239
[Mức độ 2]%[Bài giảng 12 mới, Lê Hồng Phi]
	\immini{
		Cho số thực dương $a\neq 1$. Biết rằng bất kì đường thẳng nào song song với trục $Ox$ mà cắt các đồ thị $y=4^x,\,y=a^x$, trục tung lần lượt tại $M$, $N$ và $A$ thì $AN=2AM$ (hình vẽ bên). Tính giá trị của $a$.}
	{
		\begin{tikzpicture}[scale=0.9,>=stealth, font=\footnotesize , line join=round, line cap=round]
		\def \xmin {-2.5}
		\def \xmax {4}
		\def \ymax {4}
		\draw [->] (\xmin ,0)--(\xmax ,0) node [below]{$x$};
		\draw [->] (0,-0.9)--(0,\ymax ) node [left]{$y$};
		\node at (0,0) [below left]{$O$};
		\clip (\xmin ,-0.9) rectangle (\xmax -0.1,\ymax -0.1);
		\draw[samples=200,domain=-2:4,smooth,variable=\x] plot (\x,{4^(\x)});
		\draw[samples=200,domain=-2:4,smooth,variable=\x] plot (\x,{0.5^(\x)});
		\coordinate [label=above left:$A$] (A) at (0,3);
		\coordinate [label=below right:$M$] (M) at (0.792,3);
		\coordinate [label=below left:$N$] (N) at (-1.585,3);
		\draw (M)--(N);
		\foreach \diem in {A,M,N}
		\fill (\diem )circle(1pt);
		\end{tikzpicture}
	}

	\loigiai{
Từ $AN=2AM$ suy ra $-x_N=2x_M$.\\
Ta có $MN$ song song với $Ox$ nên $$4^{x_M}=a^{x_N}\Leftrightarrow 4^{x_M}=a^{-2x_M}\Leftrightarrow 4^{x_M}=\left(\dfrac{1}{a^2}\right)^{x_M}\Leftrightarrow 4=\dfrac{1}{a^2}\Leftrightarrow a=\dfrac{1}{2}\quad (\text{do } a>0).$$
Vậy $a=\dfrac{1}{2}$.}
\end{bt}

\begin{bt}%[2D1H5-1]
% nguon: Archive/2024/TaiLieuGiangDay/1.HamSo/1. K12 - 4in1 - Hamso/ThubaiGoc/Bai5-Dang4.tex #2243
%[Bài giảng 12 mới, Lê Hồng Phi]
Khảo sát và vẽ đồ thị hàm số $y=3^{-x}$.
\loigiai
{\begin{enumEX}{1}
\item Ta có $y=3^{-x}=\left(\dfrac{1}{3}\right)^x$. Tập xác định $\mathscr{D}=\mathbb{R}$.
\item Sự biến thiên
\begin{itemize}
\item Chiều biến thiên:\\
Ta có $y'=\left(\dfrac{1}{3}\right)^x\ln \dfrac{1}{3}<0$, $\forall x\in\mathbb{R}$.\\
Hàm số nghịch biến trên $\mathbb{R}$.
\item Các giới hạn và tiệm cận:\\
$\lim\limits_{x\to +\infty}y=0$ nên đường thẳng $y=0$ là tiệm cận ngang,\\
$\lim\limits_{x\to -\infty}y=+\infty$.
\item Bảng biến thiên
\begin{center}
\begin{tikzpicture}
\tkzTabInit[nocadre=false, lgt=1.2, espcl=2.5, deltacl=0.6]{$x$/0.6,$y'$/0.6,$y$/2}
{$-\infty$,  $+\infty$}
\tkzTabLine {,-,}
\tkzTabVar{+/$+\infty$,-/$0$}
\end{tikzpicture}
\end{center}
\end{itemize}
\item Đồ thị
\immini{Đồ thị hàm số nằm phía trên trục $Ox$.\\
Đồ thị cắt trục tung tại $(0;1)$ và đi qua điểm $(-1;3)$.}{\begin{tikzpicture}[scale=1, font=\footnotesize, line join=round, line cap=round,>=stealth,x=0.6cm,y=0.6cm]
\def \xmin{-2.5};
\def \xmax{3.3};
\def \ymin{-1.0};
\def \ymax{4.5};
\draw[->] (\xmin, 0.) -- (\xmax,0.) node[anchor=north] {$x$};
\draw[->] (0.,\ymin) -- (0.,\ymax) node[anchor=west] {$y$};
\clip(\xmin,\ymin+0.1) rectangle (\xmax-0.1,\ymax-0.1);
\draw[smooth,samples=100,domain=\xmin:\xmax] plot(\x,{(1/3)^(\x)});
\draw[dashed] (-1,0) node[below]{$-1$} |- (0,3)node[right]{$3$};
\foreach \x/\y/\l/\g  in {0/0/O/-45,-1/3//0,0/1/1/0,-1/0//0,0/3//0}{\fill[black] (\x,\y) circle(1pt)+(\g:0.5) node{$\l$};}
\end{tikzpicture}}
\end{enumEX}
}
\end{bt}

\begin{bt}%[2D1H5-1]
% nguon: Archive/2024/TaiLieuGiangDay/1.HamSo/1. K12 - 4in1 - Hamso/ThubaiGoc/Bai5-Dang4.tex #2244
%[Bài giảng 12 mới, Lê Hồng Phi]
Khảo sát và vẽ đồ thị hàm số $y=\log_2 x$.
\loigiai
{\begin{enumEX}{1}
\item Tập xác định $\mathscr{D}=(0;+\infty)$.
\item Sự biến thiên
\begin{itemize}
\item Chiều biến thiên:\\
Ta có $y'=\dfrac{1}{x\ln2}>0$, $\forall x\in\mathscr{D}$.\\
Hàm số đồng biến trên $(0;+\infty)$.
\item Các giới hạn và tiệm cận:\\
$\lim\limits_{x\to 0^+}y=-\infty$ nên đường thẳng $x=0$ là tiệm cận đứng,\\
$\lim\limits_{x\to +\infty}y=+\infty$.
\item Bảng biến thiên
\begin{center}
\begin{tikzpicture}
\tkzTabInit[nocadre=false, lgt=1.2, espcl=2.5, deltacl=0.6]{$x$/0.6,$y'$/0.6,$y$/2}
{$0$,  $+\infty$}
\tkzTabLine {,+,}
\tkzTabVar{-/$-\infty$,+/$+\infty$}
\end{tikzpicture}
\end{center}
\end{itemize}
\item Đồ thị
\immini{Đồ thị hàm số nằm phía bên phải $Oy$.\\
Đồ thị cắt trục hoành tại $(1;0)$ và đi qua điểm $(2;1)$.}{\begin{tikzpicture}[scale=1, font=\footnotesize, line join=round, line cap=round,>=stealth,x=0.6cm,y=0.6cm]
\def \xmin{-1.5};
\def \xmax{4.7};
\def \ymin{-2.0};
\def \ymax{3.2};
\draw[->] (\xmin, 0.) -- (\xmax,0.) node[anchor=north] {$x$};
\draw[->] (0.,\ymin) -- (0.,\ymax) node[anchor=east] {$y$};
\clip(\xmin,\ymin+0.1) rectangle (\xmax-0.1,\ymax-0.1);
\draw[smooth,samples=100,domain=0.01:\xmax] plot(\x,{ln(\x)/ln(2)});
\draw[dashed] (2,0) node[below]{$2$} |- (0,1)node[left]{$1$};
\foreach \x/\y/\l/\g  in {0/0/O/225,2/1//0,1/0/1/-90,2/0//0,0/1//0}{\fill[black] (\x,\y) circle(1pt)+(\g:0.5) node{$\l$};}
\end{tikzpicture}}
\end{enumEX}
}
\end{bt}

\begin{bt}%[2D1H5-4]
% nguon: Archive/2024/TaiLieuGiangDay/1.HamSo/1. K12 - 4in1 - Hamso/ThubaiGoc/Bai5-Dang4.tex #2245
%[Bài giảng 12 mới, Lê Hồng Phi]
Tìm $m$ để đồ thị hàm số $y=2024^{mx^2-2x+m-2}$ cắt đường thẳng $y=1$ tại hai điểm có hoành độ trái dấu.
\loigiai
{Phương trình hoành độ giao điểm là \[2024^{mx^2-2x+m-2}=1\Leftrightarrow mx^2-2x+m-2=0.\quad (*)\]
Phương trình $(*)$ có hai nghiệm trái dấu khi và chỉ khi \[m(m-2)<0\Leftrightarrow 0<m<2.\]
Vậy giá trị $m$ cần tìm là $0<m<2$.
}
\end{bt}

\begin{bt}%[2D1H5-7]
% nguon: Archive/2024/TaiLieuGiangDay/1.HamSo/1. K12 - 4in1 - Hamso/ThubaiGoc/Bai5-Dang4.tex #2246
%[Bài giảng 12 mới, Lê Hồng Phi]
Có bao nhiêu điểm $M$ thuộc đồ thị hàm số $y=2^x$ có hoành độ là số nguyên dương và nằm dưới đường thẳng $y=2024$?
\loigiai
{Vì điểm $M$ thuộc đồ thị hàm số $y=2^x$ nên $M\left(x;2^x\right)$.\\
Ta có điểm $M$ nằm dưới đường thẳng $y=2024$ nên \[2^x<2024\Leftrightarrow x<\log_2 2024\approx 10{,}98.\]
Do $x$ nguyên dương nên $x\in\{1;2;\ldots;10\}$.\\
Vậy có $10$ điểm thỏa mãn bài toán.
}
\end{bt}

\begin{bt}%[2D1V5-3]
% nguon: Archive/2024/TaiLieuGiangDay/1.HamSo/1. K12 - 4in1 - Hamso/ThubaiGoc/Bai5-Dang4.tex #2247
%[Bài giảng 12 mới, Lê Hồng Phi]
Tìm $m$ để đường thẳng $y=m$ cắt đồ thị hàm số $y=\log_2\sqrt{1-x^2}$ tại hai điểm phân biệt.
\loigiai
{Phương trình hoành độ giao điểm giữa đồ thị và đường thẳng là $\log_2\sqrt{1-x^2}=m$.\quad $(*)$\\
Điều kiện $1-x^2>0\Leftrightarrow -1<x<1$.\\
Bài toán trở thành tìm $m$ để phương trình $(*)$ có hai nghiệm phân biệt thuộc $(-1;1)$.\\
Hàm số $y=\log_2\sqrt{1-x^2}$ có $y'=\dfrac{-x}{(1-x^2)\ln 2}$ và $y'=0\Leftrightarrow x=0$.\\
Bảng biến thiên của hàm số $y=\log_2\sqrt{1-x^2}$ trên $(-1;1)$ như sau
\begin{center}
\begin{tikzpicture}
\tkzTabInit[nocadre=false, lgt=1.2, espcl=2.5, deltacl=0.6]{$x$/0.6,$y'$/0.6,$y$/2}
{$-1$, $0$, $1$}
\tkzTabLine {,+,0,-,}
\tkzTabVar{-/$-\infty$, +/$0$, -/$-\infty$}
\end{tikzpicture}
\end{center}
Từ bảng biến thiên ta có giá trị $m$ cần tìm là $m<0$.
}
\end{bt}

\begin{bt}%[2D1H3-2]
% nguon: Archive/2024/TaiLieuGiangDay/1.HamSo/1. K12 - 4in1 - Hamso/ThubaiGoc/Bai3-Dang3.tex #2258
[CTST, Mức độ 2]%[Nguyen Huynh]
	Tìm giá trị lớn nhất và giá trị nhỏ nhất của hàm số có đồ thị được cho ở Hình 5.
	\begin{enumerate}[a.]
		\begin{multicols}{2}
			\item \begin{tikzpicture}[>=stealth,scale=0.6]
				%	\draw[line width=0.1pt,dashed] (0,0) grid (6,6); 
				\draw[-> ] (-1,0) -- (7,0)node[above left] {$x$};
				\draw[-> ] (0,-1) -- (0,7)node[below left] {$y$};
				\draw [thick,red](0,0) node[below right] {\small $O$};
				\draw [thick,red](1,6) parabola bend (3,2) (2,3);
				\draw [thick,red](3,2) parabola  (3.5,2.5); 
				\draw [thick,red](3.5,2.5) parabola [bend at end]  (4,3);
				\draw [thick,red](4,3) parabola  (4.75,1.3); 
				\draw [thick,red](4.75,1.3) parabola[bend at end]  (5,1);
				\draw [thick,red](5,1) parabola  (6,5); 
				
				\foreach \x in {1,2,3,4,5,6}\draw (\x,0.1)--(\x,-0.1) node [below] {\footnotesize $\x$};
				\foreach \y in {5,6,1,2,3}\draw (0.1,\y)--(-0.1,\y) node [left] {\footnotesize $\y$};
				\draw[dashed] (1,0) -- (1,6) -- (0,6);
				\draw[dashed] (2,0) -- (2,3);
				\draw[dashed] (3,0) -- (3,2) -- (0,2);
				\draw[dashed] (4,0) -- (4,3) -- (0,3);
				\draw[dashed] (5,0) -- (5,1) -- (0,1);
				\draw[dashed] (6,0) -- (6,5) -- (0,5);
			\end{tikzpicture}
			\item \begin{tikzpicture}[>=stealth,scale=0.6]
				%	\draw[line width=0.1pt,dashed] (0,0) grid (6,6); 
				\draw[-> ] (-4,0) -- (4,0)node[above left] {$x$};
				\draw[-> ] (0,-1) -- (0,8.5)node[below left] {$y$};
				\draw [thick,red](0,0) node[below right] {\small $O$};
				\draw [thick,red] (-3,1) parabola [bend at end] (-2,3) ;
				\draw[thick,red] (-2,3) parabola  (-1.1,1.2);
				\draw [thick,red](-1.1,1.2) parabola[bend at end]  (-1,1);
				\draw[thick,red] (-1,1) parabola bend (1,7)  (2,4); 
				\draw[thick,red] (2,4) to[bend left=10]  (3,5);
				
				\foreach \x in {1,2,3,-3,-2,-1}\draw (\x,0.1)--(\x,-0.1) node [below] {\footnotesize $\x$};
				\foreach \y in {1,3,4}\draw (0.1,\y)--(-0.1,\y) node [right] {\footnotesize $\y$};
				\foreach \y in {5,7}\draw (0.1,\y)--(-0.1,\y) node [left] {\footnotesize $\y$};
				\draw[dashed] (-3,0) -- (-3,1) -- (0,1);
				\draw[dashed] (-2,0) -- (-2,3)--(0,3);
				\draw[dashed] (-1,0) -- (-1,1) -- (0,1);
				\draw[dashed] (2,0) -- (2,4) -- (0,4);
				\draw[dashed] (3,0) -- (3,5) -- (0,5);
				\draw[dashed] (1,0) -- (1,7) -- (0,7);
			\end{tikzpicture}
		\end{multicols}
	\end{enumerate}
	\loigiai{
		Ở hình a, ta có $\displaystyle\min f(x)=1$ tại $x=5$, $\displaystyle\max f(x)=6$ tại $x=1$.\\
		Ở hình b, ta có $\displaystyle\min f(x)=1$ tại $x=-3$ và $x= -1$, $\displaystyle\max f(x)=7$ tại $x=1$.
	}
\end{bt}

\begin{bt}%[2D1H3-2]
% nguon: Archive/2024/TaiLieuGiangDay/1.HamSo/1. K12 - 4in1 - Hamso/ThubaiGoc/Bai3-Dang3.tex #2259
%[Dự án BG 4in1, Nguyễn Văn Nay]
	Cho hàm số $f(x)$ có đồ thị như hình vẽ
	\begin{center}
		\begin{tikzpicture}[>=stealth,x=1cm,y=1cm,scale=.8]
			\def\a{1} % Hệ số a phải khác 0
			\def\b{-3}
			\def\c{0}
			\def\d{2}
			%\draw[color=gray,dash pattern=on 1pt off 1pt,xstep=1.0cm,ystep=1.0cm] (-5.2,-5.2) grid (5.2,5.2);
			\draw[->] (-1.5,0) -- (3.2,0)node[below]{\scriptsize $x$};
			\draw[->] (0,-2.5) -- (0,3) node[left] {\scriptsize $y$};
			\draw (0,0)node[below left]{\scriptsize $O$};
			\clip (-1.5,-2.5)rectangle(3.5,3);
			\draw[thick,samples=150,smooth,domain=-5:5] plot(\x,{\a*(\x)^3+(\b)*(\x)^2+(\c)*\x+(\d)});
			\draw[dashed](-1,0)|-(0,-2)-|(2,0);
			\foreach \x in{-1,1,2} {\fill (\x,0) circle (1.2pt);\node[above] at (\x,0){$\x$};}
			\fill (0,-2) circle (1.2pt) node[below right]{$-2$};
			\fill (0,2) circle (1.2pt) node[above right]{$2$};
		\end{tikzpicture}
	\end{center}
	\begin{listEX}
		\item Tìm giá trị nhỏ nhất của hàm số trên $(0;+\infty)$;
		\item Hàm số đạt giá trị lớn nhất trên $(-\infty;2)$ tại $x$ bằng bao nhiêu?
		\item Tìm giá trị lớn nhất, giá trị nhỏ nhất của hàm số $g(x)=f(\sin x)+3$
	\end{listEX}
	\loigiai{
		\begin{listEX}
			\item Dựa vào đồ thị hàm số, giá trị nhỏ nhất của hàm số trên $(0;+\infty)$ bằng $-2$ khi $x=2$.
			\item Từ đồ thị hàm số, suy ra hàm số đạt giá trị lớn nhất trên $(-\infty;2)$ tại $x=0$.
			\item Đặt $t=\sin x$, $t\in [-1;1]$.\\
			Bài toán trở thành tìm giá trị lớn nhất, giá trị nhỏ nhất của hàm số $f(t)+3$ khi $t\in [-1;1]$.\\
			Từ đồ thị hàm số, suy ra $\max\limits_{[-1;1]} f(t)=f(0)=2$, $\min\limits_{[-1;1]} f(t)=f(-1)=-2$.\\
			Suy ra $\max g(x)=2+3=5$, $\min g(x)=-2+3=1$.
		\end{listEX}
	}
\end{bt}

\begin{bt}%[2D1H3-2]
% nguon: Archive/2024/TaiLieuGiangDay/1.HamSo/1. K12 - 4in1 - Hamso/ThubaiGoc/Bai3-Dang3.tex #2260
%[Dự án BG 4in1, Nguyễn Văn Nay]
	Cho hàm số $f(x)$ có bảng biến thiên như hình vẽ
	\begin{center}
		\begin{tikzpicture}[>=stealth,scale=1]
			\tkzTabInit[deltacl=.5,lgt=1.5,espcl=3]
			{$t$/1,$f’(t)$/1,$f(t)$/3}
			{$-\infty$,$-2$,$2$,$\dfrac{5}{2}$,$+\infty$}
			\tkzTabLine{,-,,h,,-,z,+,}
			\tkzTabVar{+/$+\infty$,-H/$22$,+/$2$,-/$\dfrac{7}{4}$,+/$+\infty$}
		\end{tikzpicture}
	\end{center}
	\begin{listEX}
		\item Tìm giá trị lớn nhất, giá trị nhỏ nhất của hàm số trên tập xác định.
		\item Tìm giá trị nhỏ nhất của hàm số trên $(2;10)$.
		\item Phương trình $f(x^2)=3$ có bao nhiêu nghiệm?
	\end{listEX}
	\loigiai{
		\begin{listEX}
			\item Dựa vào bảng biến thiên, suy ra hàm số không có giá trị lớn nhất, giá trị nhỏ nhất bằng $\dfrac{7}{4}$ khi $x=\dfrac{5}{2}$.
			\item Trên $(2;10)$, ta có $\min\limits_{(2;10)} f(x)=f(\dfrac{5}{2}=\dfrac{7}{4}$.
			\item Đặt $t=x^2\ge 0$.\\
			Phương trình $f(t)=3$,  có hai nghiệm $t_1<-2$, $t_2>\dfrac{5}{2}$.\\
			Từ điều kiện $t=x^2\ge 0$ suy ra $t_1$ không thỏa mãn.\\
			Phương trình đã cho tương đương $x^2=t_2 \Leftrightarrow x=\pm \sqrt{t_2}$.\\
			Vậy phương trình đã cho có $2$ nghiệm.
		\end{listEX}
	}
\end{bt}

\begin{bt}%[2D1H3-2]
% nguon: Archive/2024/TaiLieuGiangDay/1.HamSo/1. K12 - 4in1 - Hamso/ThubaiGoc/Bai3-Dang3.tex #2261
%[Dự án BG 4in1, Nguyễn Văn Nay]
	\immini
		{Cho hàm số $y=f(x)$ liên tục trên $\mathbb{R}$ và có đồ thị hàm số $y=f'(x)$ như hình vẽ. Hàm số $y=f(x)$ đạt giá trị nhỏ nhất trên đoạn $\left[0;\dfrac{7}{2}\right]$ tại điểm có hoành độ bằng bao nhiêu?
		}
		{ \begin{tikzpicture}[line join = round, line cap = round,>=stealth,font=\footnotesize,scale=0.7]
				\def\xt{-1} \def\xp{4.3} \def\yd{-3} \def\yt{3.5}
				\draw[->] (\xt,0)--(\xp,0) node[below]{$x$};
				\draw[->] (0,\yd)--(0,\yt) node[left]{$y$};
				\fill (0,0) circle (1.5pt) node[below left]{$O$};
				\begin{scope}
					\clip (\xt+0.1,\yd+0.1) rectangle (\xp-0.1,\yt-0.1);
					\draw[samples=100,domain=\xt:\xp,smooth] plot (\x, {1/1.5*(\x-1)^2*(\x-3)});
				\end{scope}
				\draw[dashed] (3.5,0)node[below]{$\dfrac{7}{2}$}--(3.5,2.083);
				\node at (1,0)[above]{$1$};
				\node at (3,0)[above left]{$3$};
			\end{tikzpicture}
		} 
		\loigiai{
			Dự vào đồ thị của hàm số $y=f'(x)$, ta có bảng biến thiên sau:
			\begin{center}
				\begin{tikzpicture}
					\tkzTabInit[nocadre=false,lgt=1.2,espcl=2.5,deltacl=0.6]{$x$/.6 ,$f'(x)$/.6,$f(x)$/2}
					{$0$ ,$1$, $3$ , $\tfrac{7}{2}$}
					\tkzTabLine{ , - , $0$ ,-,$0$, + , }
					\tkzTabVar{+/ ,R, -/$f(3)$ , +/}
				\end{tikzpicture}
			\end{center}
			Dựa vào bảng biến thiên trên, ta có giá trị nhỏ nhất trên đoạn $\left[0;\dfrac{7}{2}\right]$ là $f(3)$.
		} 
\end{bt}

\begin{bt}%[2D1H3-2]
% nguon: Archive/2024/TaiLieuGiangDay/1.HamSo/1. K12 - 4in1 - Hamso/ThubaiGoc/Bai3-Dang3.tex #2262
%[Dự án BG 4in1, Nguyễn Văn Nay]
	\immini
	{Cho hàm số bậc ba $y=f(x)$ có đạo hàm trên $\mathbb{R}$ và có đồ thị như hình vẽ. 
	\begin{listEX}
		\item  Tìm giá trị lớn nhất và giá trị nhỏ nhất của hàm số trên $[-2;2]$.
		\item Tìm tập hợp các giá trị của $m$ để $f(x)=m$ có nghiệm trên $(-1;2)$.
		\item Tìm tập hợp các giá trị $a$ để $f(x)>a$, với mọi $x\in (-1;1)$.
		\item Xét hàm số $g(x)=m^2-3m-4+f(2x^3+x-1)$. Tìm các giá trị của $m$ để $\max\limits_{x\in[0;1]}g(x)=3$.
	\end{listEX}
	}
	{
		\begin{tikzpicture}[line cap=round,line join=round,font=\footnotesize,>=stealth,scale=0.75]
			\draw[->] (-2.5,0)--(2.5,0) node[above] {$x$};
			\draw[->] (0,-2)--(0,3.5) node[left] {$y$};
			\fill (0,0) circle (1pt) node [below left] {$O$}
			(-2,0) circle (1pt) node [shift={(120:2ex)}] {$-2$}
			(0,3) circle (1pt) node [shift={(45:2ex)}] {$3$}
			(0,-1) circle (1pt) node [shift={(-135:2ex)}] {$-1$}
			(1,0) circle (1pt) node [above] {$1$}
			(-1,3) circle (1pt) (1,-1) circle (1pt) (2,3) circle (1pt) (-2,-1) circle (1pt);
			\foreach \x in {-1,2} \fill (\x,0) circle (1pt) node [below] {$\x$};
			\draw [samples=100, domain=-2.1:2.03] plot (\x, {(\x)^3-3*(\x)+1});
			\draw[dashed](1,0)|-(0,-1) (-1,0)|-(0,3) (2,0)|-(0,3) (-2,0)|-(0,-1);
	\end{tikzpicture}}
	\loigiai{
		\begin{listEX}
			\item  Dựa vào đồ thị, trên $[-2;2]$, giá trị nhỏ nhất của hàm số bằng $-1$ khi $x=1$, giá trị lớn nhất của hàm số bằng $3$ khi $x=-1$.
			\item Số nghiệm của phương trình $f(x)=m$ là số điểm chung của hai đồ thị $y=f(x)$ và $y=m$.\\
			Dựa vào hình vẽ, để hai đồ thị có điểm chung trên $(-1;2)$ thì $-1\le m <3$.
			\item Từ đồ thị, ta có $f(x)>-1$, với mọi $x \in (-1;1)$. Do đó để $f(x)>a$, với mọi $x\in (-1;1)$ thì $a\le -1$.
			\item Ta có $g'(x)=(6x^2+1)\cdot f'(2x^3+x-1)$ và $$g'(x)=0\Leftrightarrow f'(2x^3+x-1)=0\Leftrightarrow \hoac{&2x^3+x-1=-1\\&2x^3+x-1=1}\Leftrightarrow\hoac{&x=0\\&x=-1\\&x=1.}$$
			Ta có $g(0)=m^2-3m-4+f(-1)=m^2-3m-1$ và $g(1)=m^2-3m-4+f(2)=m^2-3m-1$.\\
			Do $\max\limits_{x\in[0;1]}g(x)=3\Leftrightarrow m^2-3m-1=3\Leftrightarrow m^2-3m-4=0\Leftrightarrow m=-1$ hoặc $m=4$.
		\end{listEX}
	}
\end{bt}

\begin{bt}%[2D1H3-2]
% nguon: Archive/2024/TaiLieuGiangDay/1.HamSo/1. K12 - 4in1 - Hamso/ThubaiGoc/Bai3-Dang3.tex #2263
%[Dự án BG 4in1, Nguyễn Văn Nay]
	\immini{
		Cho hàm số $ y=f(x)$ trên $\left[ -2;4 \right]$ như hình vẽ. Gọi $S$ là tập chứa các giá trị của $ m$ để hàm số $ y=\left( f\left(2-x \right)+m \right)^{2}$ có giá trị lớn nhất trên đoạn $\left[ -2;4 \right]$ bằng $49$. Tính tổng các phần tử tập $S$.
	}{
		\begin{tikzpicture}[yscale=0.3]
			\def\xmax{5}
			\def\xmin{-3}
			\def\ymax{7}
			\def\ymin{-4.5}
			\path
			(-3,-4) coordinate (b)
			(-2,-2) coordinate (c)
			(-1,-3) coordinate (d)
			(0,-4) coordinate (e)
			(2,6) coordinate (f)
			(4,-2) coordinate (g)
			;
			\begin{scope}
				\clip (\xmin,\ymin) rectangle (\xmax,\ymax);
				\draw[name path=dothi]
				(b) .. controls +(80:0.5) and +(180:0.5) ..
				(c) .. controls +(0:0.5) and +(180:0.6) ..
				(e) .. controls +(0:1) and +(-91:0.1)..
				(f) .. controls +(-10:1) and +(100:1).. (g)
				;
				\draw[-stealth,name path=ox] (\xmin,0)--(\xmax,0) node[below left]{$x$};
				\draw[-stealth] (0,\ymin)--(0,\ymax) node[below right]{$y$};
				\path [name intersections = {of =dothi and ox,by={P}}];
			\end{scope}
			\draw
			(-2,0) node[above] {$-2$}
			(2,0) node[below] {$2$}
			(4,0) node[above] {$4$}
			(0,6) node[left]{$6$}
			(0,-4) node[below right]{$-4$}
			(0,0) node[below left] {$O$};
			\draw[dashed] (-2,0)--(c)--(4,-2)--(4,0) (2,0)--(f)--(0,6);
		\end{tikzpicture}
	}
	\loigiai{
		Do $ y=\left(f\left( 2-x \right)+m \right)^{2}$ có giá trị lớn nhất trên đoạn $[-2;4]$ bằng $49$ nên với $\forall x\in [-2;4]$ ta có 
		\[
		y=\left(f\left( 2-x \right)+m \right)^{2}\le 49
		\Leftrightarrow -7\le f\left( 2-x \right)+m\le 7
		\Leftrightarrow -7-m\le f\left( 2-x \right)\le 7-m.
		\]
		Mặt khác $-4\le f\left( 2-x \right)\le 6,\forall x\in \left[ -2;4 \right]$, suy ra $\hoac{& 7-m=6 \\ 
			& -7-m=-4}\Leftrightarrow \hoac{& m=1 \\ 
			& m=-3.}$\\
		Vậy $S=\left\{ 1;-3 \right\}$. Khi đó tổng giá trị của các phần tử của là $-2.$
	}
\end{bt}

\begin{bt}%[2D1H3-1]
% nguon: Archive/2024/TaiLieuGiangDay/1.HamSo/1. K12 - 4in1 - Hamso/ThubaiGoc/Bai3-Dang2.tex #2275
%[Dự án BG 4in1, Nguyễn Văn Nay][CKP, Mức độ 1]
	Tìm giá trị nhỏ nhất, giá trị lớn nhất của hàm số $y=f(x)=x^3-6x^2+9x+1$ trên nửa khoảng $[-1;4)$.
	\loigiai{
		Trên nửa khoảng $[-1;4)$, ta có $y'=f'(x)=3x^2-12x+9$.
		\begin{itemize}
			\item $y'=0\Leftrightarrow \hoac{&x=3\in [-1;4)\\&x=1\in [-1;4).}$
			\item $\lim\limits_{x\to 4^-}f(x)=\lim\limits_{x\to 4^-}\left(x^3-6x^2+9x+1\right)=5$.
		\end{itemize}
		Bảng biến thiên
		\begin{center}
			\begin{tikzpicture}[>=stealth]
				\tkzTabInit[nocadre=false,lgt=1.5,espcl=3,deltacl=0.6]{$x$/.6 ,$f'(x)$/.6,$f(x)$/2}
				{$-1$ , $1$ , $3$ , $4$}
				\tkzTabLine{ , + , $0$ , - ,  $0$ , +, }
				\tkzTabVar{-/$-15$ , +/$5$ , -/$1$,	+/$5$ ,}
			\end{tikzpicture}
		\end{center}
		Từ bảng biến thiên, ta có $-15\leq f(x)\leq 5,\forall x\in [-1;4)$.\\
		Suy ra trên khoảng $[-1;4)$, 
		\begin{itemize}
			\item Hàm số có giá trị nhỏ nhất là $-15$ tại $x=-1$.
			\item Hàm số có giá trị lớn nhất là $5$ tại $x=1$.
		\end{itemize}
	}
\end{bt}

\begin{bt}%[2D1H3-1]
% nguon: Archive/2024/TaiLieuGiangDay/1.HamSo/1. K12 - 4in1 - Hamso/ThubaiGoc/Bai3-Dang2.tex #2276
%[Dự án BG 4in1, Nguyễn Văn Nay][CTST, Mức độ 2]
	Tìm giá trị lớn nhất, giá trị nhỏ nhất của hàm  số $g(x)=x+\dfrac{1}{x}$ trên khoảng $\left(0 ; 5\right)$.
	\loigiai{
		Xét $g(x)=x+\dfrac{1}{x}$ trên khoảng $(0;5)$. Nhận thấy hàm số xác định và liên tục trên khoảng $(0;5)$.\\
		Ta có $g'(x)=1-\dfrac{1}{x^2}=\dfrac{x^2-1}{x^2}=0\Leftrightarrow x=1$.\\
		Bảng biến thiên
		\begin{center}
			\begin{tikzpicture}
				\tkzTabInit[espcl=3,lgt=1.5]
				{$x$/0.7,$g'(x)$/0.7,$g(x)$/2.1}
				{$0$,$1$,$5$}
				\tkzTabLine{,-,0,+,}
				\tkzTabVar{+/$+\infty$/$+\infty$,-/$2$,+/$\dfrac{26}{5}$}
			\end{tikzpicture}
		\end{center}
		Từ bảng biến thiên, ta thấy $\min\limits_{(0;5)}f(x)=2$ và hàm số không có giá trị lớn nhất trên khoảng $(0;5)$.
	}	
\end{bt}

\begin{bt}%[2D1H3-1]
% nguon: Archive/2024/TaiLieuGiangDay/1.HamSo/1. K12 - 4in1 - Hamso/ThubaiGoc/Bai3-Dang2.tex #2277
%[Dự án BG 4in1, Nguyễn Văn Nay][CD, Mức độ 2]%[Tên dự án, Tên GV]%[ID6]
	Tìm giá trị lớn nhất và giá trị nhỏ nhất (nếu có) của hàm số $f(x)=\dfrac{x^2+9}{x}$ trên khoảng $(0;+\infty)$.
	\loigiai{
		\begin{itemize}
			\item Xét hàm số $f(x)=\dfrac{x^2+9}{x}$ với $x\in(0;+\infty)$.
			\item Ta có $f'(x)=\dfrac{x^2-9}{x^2}$.\\
			Khi đó $f'(x)=0 \Leftrightarrow x^2-9=0 \Leftrightarrow \hoac{&x=3\in(0;+\infty)\\&x=-3\notin(0;+\infty).} $\\
			Các giới hạn $\lim\limits_{x\to0^+}f(x)=+\infty$, $\lim\limits_{x\to+\infty}f(x)=+\infty$.
		\end{itemize}
			Bảng biến thiên
			\begin{center}
				\begin{tikzpicture}[thick]
					\tkzTabInit[nocadre=false,lgt=1.2,espcl=2.5,deltacl=0.6]
					{$x$/0.8,$f'(x)$/0.8,$f(x)$/2}
					{$0$,$3$,$+\infty$}
					\tkzTabLine{d,-,0,+,}
					\tkzTabVar{D+/$+\infty$,-/$6$,+/$+\infty$}
				\end{tikzpicture}	
			\end{center}				
		Căn cứ bảng biến thiên, ta có $\min\limits_{(0;+\infty)}f(x)=6$ tại $x=3$ và hàm số $f(x)$ không có giá trị lớn nhất.
	}
\end{bt}

\begin{bt}%[2D1H3-1]
% nguon: Archive/2024/TaiLieuGiangDay/1.HamSo/1. K12 - 4in1 - Hamso/ThubaiGoc/Bai3-Dang2.tex #2286
%[Dự án BG 4in1, Nguyễn Văn Nay]
	Tìm giá trị lớn nhất của hàm số $f(x)=x^3-2x^2+x-6$ trên khoảng $(-1;1)$.
	\loigiai{
		Tập xác định $\mathscr{D}=\mathbb{R}$, ta chỉ xét trên khoảng $(-1;1)$.\\
		Đạo hàm $y'=3x^2-4x+1; y'=0\Leftrightarrow\hoac{&x=1\\&x=\frac{1}{3}.}$\\
		Bảng biến thiên:
		\begin{center}
			\begin{tikzpicture}
				\tkzTabInit[lgt=1.5,espcl=2.5,deltacl=.5]
				{$x$ /.7, $y’$ /.7,$y$ /2}
				{$-1$ , $\frac{1}{3}$, $1$}
				\tkzTabLine{,+,0,-,}
				\tkzTabVar{-/$-10$ ,+/ $-\frac{158}{27}$ ,-/$-6$}
			\end{tikzpicture}
		\end{center}
		Từ bảng biến thiên ta có $\max\limits_{(-1;1)}y=-\dfrac{158}{27}=y\left(\dfrac{1}{3}\right)$.}
\end{bt}

\begin{bt}%[2D1H3-1]
% nguon: Archive/2024/TaiLieuGiangDay/1.HamSo/1. K12 - 4in1 - Hamso/ThubaiGoc/Bai3-Dang2.tex #2287
%[Dự án BG 4in1, Nguyễn Văn Nay]
	Tìm giá trị lớn nhất, giá trị nhỏ nhất của hàm số $f(x)=x^3(1-x)^2$ trên nửa khoảng $\left[\dfrac{1}{2};+\infty\right)$.
	\loigiai{
		Tập xác định $\mathscr{D}=\mathbb{R}$, ta chỉ xét trên nửa khoảng $\left[\dfrac{1}{2};+\infty\right)$.\\
		Đạo hàm $f'(x)=x^2(1-x)(3-5x); f'(x)=0\Leftrightarrow x=0, x=1, x=\dfrac{3}{5}$.\\Bảng biến thiên:
		\begin{center}
			\begin{tikzpicture}
				\tkzTabInit[lgt=1.5,espcl=2.5,deltacl=.5]
				{$x$ /1, $f'(x)$ /.7,$f(x)$ /2}
				{$\frac{1}{2}$ , $\frac{3}{5}$ , $1$ , $+\infty$}
				\tkzTabLine{,+,0,-,0,+,}
				\tkzTabVar{-/$\frac{1}{32}$ ,+/ $\frac{108}{3215}$ ,-/$0$, +/$+\infty$}
			\end{tikzpicture}
		\end{center}
		Từ bảng biến thiên cho ta\\ $\min\limits_{\left[\frac{1}{2};+\infty\right)}f(x)=f(1)=0$ và $\max\limits_{\left[\frac{1}{2};+\infty\right)}f(x)$ không tồn tại vì $\lim\limits_{x\to +\infty}f(x)=+\infty$.
	}
\end{bt}

\begin{bt}%[2D1H3-1]
% nguon: Archive/2024/TaiLieuGiangDay/1.HamSo/1. K12 - 4in1 - Hamso/ThubaiGoc/Bai3-Dang2.tex #2288
%[Dự án BG 4in1, Nguyễn Văn Nay]
	Tìm giá trị lớn nhất và giá trị nhỏ nhất của hàm số $f(x)=\dfrac{x}{4+x^2}$ trên $\mathbb{R}$.
	\loigiai
	{Tập xác định $\mathscr{D}=\mathbb{R}$.\\
		Đạo hàm $f'(x)=\dfrac{4-x^2}{\left(4+x^2\right)^2}; f'(x)=0\Leftrightarrow x=\pm 2$.\\
		Bảng biến thiên:
		\begin{center}
			\begin{tikzpicture}
				\tkzTabInit[lgt=1.5,espcl=2.5,deltacl=.5]
				{$x$ /.7, $f'(x)$ /.7,$f(x)$ /2}
				{$-\infty$ , $-2$ , $2$ , $+\infty$}
				\tkzTabLine{,-,z,+,z,-,}
				\tkzTabVar{+/$0$ ,-/ $-\dfrac{1}{4}$ ,+/$\dfrac{1}{4}$, -/$0$}
			\end{tikzpicture}
		\end{center}
		Vậy $\min\limits_{\mathbb{R}}f(x)=f(-2)=-\dfrac{1}{4}$ và $\max\limits_{\mathbb{R}}f(x)=f(2)=\dfrac{1}{4}$.
	}
\end{bt}

\begin{bt}%[2D1H3-1]
% nguon: Archive/2024/TaiLieuGiangDay/1.HamSo/1. K12 - 4in1 - Hamso/ThubaiGoc/Bai3-Dang2.tex #2289
%[Dự án BG 4in1, Nguyễn Văn Nay]
	Tìm giá trị nhỏ nhất và giá trị lớn nhất (nếu có) của hàm số $f(x)=x^2+\dfrac{2}{x}$ trên khoảng $\left(0;+\infty\right)$.
	\loigiai{
	Điều kiện $x\ne 0$.\\
	Đạo hàm $f'(x)=2x-\dfrac{2}{x^2}=\dfrac{2x^3-2}{x^2}$.\\
	$f'(x)=0\Leftrightarrow x=1$.\\
	Bảng biến thiên
	\begin{center}
		\begin{tikzpicture}
			\tkzTabInit[nocadre=false,lgt=1.2,espcl=2.5,deltacl=0.6]
			{$x$ /0.6,$f'(x)$ /0.6,$f(x)$ /2}
			{$-\infty$,$0$,$1$,$+\infty$}
			\tkzTabLine{,-,d,-,$0$,+,}
			\tkzTabVar{-/$-\infty$,+D-/$+\infty$/$-\infty$, +/$3$,-/$-\infty$}
		\end{tikzpicture}
	\end{center}
	Dựa vào bảng biến thiên, suy ra giá trị lớn nhất của hàm số trên $(0;+\infty)$ là $3$ khi $x=1$, hàm số không có giá trị nhỏ nhất trên $(0;+\infty)$.
	}
\end{bt}

\begin{bt}%[2D1H3-1]
% nguon: Archive/2024/TaiLieuGiangDay/1.HamSo/1. K12 - 4in1 - Hamso/ThubaiGoc/Bai3-Dang2.tex #2290
%[Dự án BG 4in1, Nguyễn Văn Nay]
	Tìm giá trị nhỏ nhất của hàm số $f(x)=2\sqrt{x+1}-\sqrt{x-2}$ trên tập xác định của nó.
	\loigiai
	{
		Tập xác định $\mathscr{D}=\left[2;+\infty\right)$.\\
		Đạo hàm $f'(x)=\dfrac{1}{\sqrt{x+1}}-\dfrac{1}{2\sqrt{x-2}}=\dfrac{2\sqrt{x-2}-\sqrt{x-1}}{2\sqrt{x+1}\cdot\sqrt{x-2}}, x>2$.\\
		Ta có $f'(x)=0\Leftrightarrow 2\sqrt{x-2}-\sqrt{x-1}\Leftrightarrow 2\sqrt{x-2}=\sqrt{x+1}\Leftrightarrow 4(x-2)=x+1\Leftrightarrow x=3$ (nhận).\\
		Bảng biến thiên:
		\begin{center}
			\begin{tikzpicture}
				\tkzTabInit[lgt=1.5,espcl=2.5,deltacl=.5]
				{$x$ /.7, $y’$ /.7,$y$ /2}
				{$2$ , $3$, $+\infty$}
				\tkzTabLine{,-,0,+,}
				\tkzTabVar{+/ ,-/ $3$ ,+/}
			\end{tikzpicture}
		\end{center}
		Vậy $\min\limits_{\left[2;+\infty\right)}f(x)=f(3)=3$.
	}
\end{bt}

\begin{bt}%[2D1H3-1]
% nguon: Archive/2024/TaiLieuGiangDay/1.HamSo/1. K12 - 4in1 - Hamso/ThubaiGoc/Bai3-Dang2.tex #2291
%[Dự án BG 4in1, Nguyễn Văn Nay]
	Cho hàm số $f(x)=x+\sqrt{2x^2+1}$. Tìm tham số $m$ thỏa mãn
	\begin{enumerate}[a)]
		\item Phương trình $x+\sqrt{2x^2+1} = m$ có nghiệm.
		\item Bất phương trình $x + \sqrt{2x^2+1}>m$ có nghiệm với mọi $\forall x \in \mathbb{R}$.
	\end{enumerate}
	\loigiai{
		Đặt $f(x) = x+\sqrt{2x^2+1} \Rightarrow f'(x) = 1+\dfrac{2x}{\sqrt{2x^2+1}}$.\\
		Ta có $f'(x) = 0 
		\Leftrightarrow \sqrt{2x^2+1} = -2x \Leftrightarrow x = -\dfrac{1}{\sqrt{2}}$.\\
		Bảng biến thiên:
		\begin{center}
			\begin{tikzpicture}[>=stealth,scale=1]
				\tkzTabInit[deltacl=.5,lgt=1.2,espcl=3]
				{$x$/1,$f’(x)$/1,$f(x)$/2.5}
				{$-\infty$,$-\dfrac{1}{\sqrt{2}}$,$+\infty$}
				\tkzTabLine{ ,-,z,+, }
				\tkzTabVar{+/$+\infty$,-/$\dfrac{1}{\sqrt{2}}$,+/$+\infty$}
			\end{tikzpicture}
		\end{center}
		\begin{enumerate}[a)]
			\item Nhìn bảng biến thiên suy ra phương trình $f(x) = m$ có nghiệm $\Leftrightarrow m \ge f\left(-\dfrac{1}{\sqrt{2}}\right) = \dfrac{1}{\sqrt{2}}$.
			\item $f(x) = x+\sqrt{2x^2+1} > m, \forall x \in \mathbb{R} \Leftrightarrow m < \min{f(x)} \Leftrightarrow m < f\left(-\dfrac{1}{\sqrt{2}}\right) = \dfrac{1}{\sqrt{2}}$.
		\end{enumerate}
	}
\end{bt}

\begin{bt}%[2D1H5-5]
% nguon: Archive/2024/TaiLieuGiangDay/1.HamSo/1. K12 - 4in1 - Hamso/ThubaiGoc/Bai5-Dang5.tex #2304

	Cho hàm số $y=\dfrac{2x-2}{x+1}$ có đồ thị như hình vẽ sau. Hãy  suy ra đồ thị của hàm số $y=\left| \dfrac{2x-2}{x+1} \right|$ ?
	\begin{center}
		\begin{tikzpicture}[>=stealth,x=1cm,y=1cm,scale=0.4]
		\def\a{2}
		\def\b{-2}
		\def\c{1}
		\def\d{1}
		\draw[->] (-7.5,0) -- (5.5,0) node[below] {\scriptsize $x$};
		\draw[->] (0,-5.5) -- (0,8.5) node[right] {\scriptsize $y$};
		\fill (0,0)node[above right]{\scriptsize $O$} circle(2pt);
		\draw[dashed,blue] (-1,-5)--(-1,8) (-7,2)--(5,2); % Vẽ TCĐ và TCN
		\clip (-7,-5)rectangle(5,8);
		\pgfmathsetmacro{\can}{-(\d)/(\c)}
		\draw[thick,samples=150,smooth,domain=-7:{\can-.1}] plot(\x,{(\a*\x+(\b))/(\c*\x+(\d))}); % Vẽ nhánh bên trái TCĐ
		\draw[thick,samples=150,smooth,domain={\can+.1}:5] plot(\x,{(\a*\x+(\b))/(\c*\x+(\d))}); % Vẽ nhánh bên phải TCĐ
		\fill (0,2) node[above right] {\scriptsize $2$} circle(2pt)
		(0,-2) node[right] {\scriptsize $-2$} circle(2pt)
		(-1,0) node[below left] {\scriptsize $-1$} circle(2pt)
		(1,0) node[below right] {\scriptsize $1$} circle(2pt);
		\end{tikzpicture}
	\end{center}
	\loigiai{
		Tập xác định $\mathscr{D}=\mathbb{R} \setminus \{-1\}$.\\
		Đồ thị hàm số $y=\left| \dfrac{2x-2}{x+1} \right|$ có được bằng cách:
		\begin{itemize}
			\item Giữ nguyên phần đồ thị của hàm số $y= \dfrac{2x-2}{x+1}$ nằm trên trục hoành.
			\item Lấy đối xứng qua trục hoành phần đồ thị của $y= \dfrac{2x-2}{x+1}$ nằm dưới trục hoành.
		\end{itemize}
		Do đó đồ thị của hàm số $y=\left| \dfrac{2x-2}{x+1} \right|$ là
		\begin{center}
			\begin{tikzpicture}[>=stealth,x=1cm,y=1cm,scale=0.4]
				\def\a{2}
				\def\b{-2}
				\def\c{1}
				\def\d{1}
				\draw[->] (-7.5,0) -- (5.5,0) node[below] {\scriptsize $x$};
				\draw[->] (0,-5.5) -- (0,8.5) node[right] {\scriptsize $y$};
				\fill (0,0)node[below right]{\scriptsize $O$} circle(2pt);
				\draw[dashed,blue] (-1,-5)--(-1,8) (-7,2)--(5,2); % Vẽ TCĐ và TCN
				\clip (-7,-5)rectangle(5,8);
				\pgfmathsetmacro{\can}{-(\d)/(\c)}
				\draw[thick,samples=150,smooth,domain=-7:{\can-.1}] plot(\x,{(\a*\x+(\b))/(\c*\x+(\d))}); % Vẽ nhánh bên trái TCĐ
				\draw[thick,samples=150,smooth,domain={\can+.1}:5] plot(\x,{abs((\a*\x+(\b))/(\c*\x+(\d)))}); % Vẽ nhánh bên phải TCĐ
				\fill (0,2) node[above right] {\scriptsize $2$} circle(2pt)
				(0,-2) node[right] {\scriptsize $-2$} circle(2pt)
				(-1,0) node[below left] {\scriptsize $-1$} circle(2pt)
				(1,0) node[below right] {\scriptsize $1$} circle(2pt);
			\end{tikzpicture}
		\end{center} 
	}

\end{bt}

\begin{bt}%[2D1H5-2]
% nguon: Archive/2024/TaiLieuGiangDay/1.HamSo/1. K12 - 4in1 - Hamso/ThubaiGoc/Bai5-Dang5.tex #2305

	Cho đồ thị hàm số $y=x^3-6x^2+9x$ như hình dưới. Hãy suy ra đồ thị của hàm số $y=|x|^3-6x^2+9|x|$.
	
	\begin{center}
		\begin{tikzpicture}[>=stealth,x=1cm,y=1cm,scale=0.4]
		\def\a{1} % Hệ số a phải khác 0
		\def\b{-6}
		\def\c{9}
		\def\d{0}
		\draw[->] (-5,0) -- (5,0) node[below]{\scriptsize $x$};
		\foreach \x in {-4,-3,-2,-1,1,2,3,4}
		\draw[shift={(\x,0)}] (0,2pt)--(0,-2pt) node[below]{};
		\draw[->] (0,-1) -- (0,5)node[right]{\scriptsize $y$};
		\foreach \y in {1,2,3,4,5}
		\draw[shift={(0,\y)}] (2pt,0)--(-2pt,0) node[right]{};
		\fill (0,0)node[below left]{\scriptsize $O$} circle(2pt);
		\clip (-5,-2.5)rectangle(5,5);
		\draw[thick,samples=150,smooth,domain=-0.1:5] plot(\x,{\a*(\x)^3+(\b)*(\x)^2+(\c)*\x+(\d)});
		\draw (0,-2) node{$\text{Hình } 1$};
		\end{tikzpicture}
	\end{center}
	\loigiai{
		Tập xác định $\mathscr{D}=\mathbb{R}$.\\
		Ta thấy đồ thị hàm số $y=|x|^3-6|x|^2+9|x|$ có được bằng cách:
		\begin{itemize}
			\item Giữ nguyên phần đồ thị của hàm số $y=x^3-6x^2+9x$ nằm bên phải trục $Oy$.
			\item Lấy đối xứng qua trục tung phần đồ thị của $y=x^3-6x^2+9x$ nằm bên phải trục $Oy$.
		\end{itemize}
		Do đó đồ thị hàm số $y=|x|^3-6|x|^2+9|x|$ là
		\begin{center}
			
			\begin{tikzpicture}[>=stealth,x=1cm,y=1cm,scale=0.4]
				\def\a{1} % Hệ số a phải khác 0
				\def\b{-6}
				\def\c{9}
				\def\d{0}
				\draw[->] (-5,0) -- (5,0) node[below]{\scriptsize $x$};
				\foreach \x in {-4,-3,-2,-1,1,2,3,4}
				\draw[shift={(\x,0)}] (0,2pt)--(0,-2pt) node[below]{};
				\draw[->] (0,-1) -- (0,5)node[right]{\scriptsize $y$};
				\foreach \y in {1,2,3,4,5}
				\draw[shift={(0,\y)}] (2pt,0)--(-2pt,0) node[right]{};
				\fill (0,0)node[below left]{\scriptsize $O$} circle(2pt);
				\clip (-5,-2.5)rectangle(5,5);
				\draw[thick,samples=250,smooth,domain=-5:5] plot(\x,{\a*abs((\x))^3+(\b)*(\x)^2+(\c)*abs(\x)+(\d)});
				\draw (0,-2) node{$\text{Hình } 2$};
			\end{tikzpicture}			
		\end{center}
	}

\end{bt}

\begin{bt}%[2D1H5-2]
% nguon: Archive/2024/TaiLieuGiangDay/1.HamSo/1. K12 - 4in1 - Hamso/ThubaiGoc/Bai5-Dang5.tex #2306
 
	Cho hàm số $y=f(x)=\dfrac{x^2-2x+1}{x+1}$ có đồ thị như hình dưới. Hãy suy ra đồ thị của hàm số $y=f(|x|)$?
	\begin{center}
			\begin{tikzpicture}[>=stealth,line join=round,line cap=round,font=\footnotesize,scale=.4]
			\tikzset{declare function={
					x1=-6.5;x2=4.5;y1=-12.5;y2=5.5;% Miền hình vẽ
					a=1;b=-2;c=1; m=1;n=1;
					f(\x)=(a*(\x)^2+b*(\x)+c)/(m*(\x)+n);	
					a1=a/m;b1=(b-a*n/m)/m;
					g(\x)=a1*\x+b1;
					q1=-(n/m)-.1;q2=-n/m+.1;
			}}
			\begin{scope}
				\clip (x1,y1)rectangle(x2,y2);
				\draw [smooth,samples=300,domain=x1:q1]plot(\x,{f(\x)});
				\draw [smooth,samples=300,domain=q2:x2]plot(\x,{f(\x)});
				\draw [smooth,samples=100,domain=x1:x2,dashed]plot(\x,{g(\x)});
				\draw[dashed] (-n/m,y1)--(-n/m,y2);
			\end{scope}
			\draw[->] (0,y1)--(0,y2)node[left]{$y$};
			\draw[->] (x1,0)--(x2,0)node[above]{$x$};
			\filldraw[black](0,0)circle(1.2pt)node[above left]{$O$};
			\draw[dashed]
			(-3,0)|-(0,-8)
			;
			
			\foreach \x/\goc in {-3/-90,-1/-90,1/-70,2/-90}{
				\draw[fill=black] (\x,0)circle(1.2pt) node[shift={(\goc:2.8mm)}]{$\x$};
			}
			\foreach \y/\goc in {-3,-4,1,2,-8}{
				\draw[fill=black] (0,\y)circle(1.2pt)node[shift={(\goc:2.8mm)}]{$\y$};
			}
		\end{tikzpicture}
	\end{center}
	\loigiai{
		Tập xác định $\mathscr{D}=\mathbb{R}\setminus {-1}$.\\
		Ta thấy đồ thị hàm số $y=f(|x|)$ có được bằng cách:
		\begin{itemize}
			\item Giữ nguyên phần đồ thị của hàm số $y=f(x)=\dfrac{x^2-2x+1}{x+1}$ nằm bên phải trục $Oy$.
			\item Lấy đối xứng qua trục tung phần đồ thị của $y=f(x)=\dfrac{x^2-2x+1}{x+1}$ nằm bên phải trục $Oy$.
		\end{itemize}
		Do đó đồ thị hàm số $y=f(|x|)=\dfrac{|x|^2-2|x|+1}{|x|+1}$ là
		\begin{center}
		\begin{tikzpicture}[>=stealth,line join=round,line cap=round,font=\footnotesize,scale=.4]
			\tikzset{declare function={
					x1=-6.5;x2=4.5;y1=-4.5;y2=5.5;% Miền hình vẽ
					a=1;b=-2;c=1; m=1;n=1;
					f(\x)=(a*(\x)^2+b*(\x)+c)/(m*(\x)+n);	
					a1=a/m;b1=(b-a*n/m)/m;
					g(\x)=a1*\x+b1;
					q1=-(n/m)-.1;q2=-n/m+.1;
			}}
			\begin{scope}
				\clip (x1,y1)rectangle(x2,y2);
				\draw [smooth,samples=300,domain=x1:q1]plot(\x,{f(abs(\x))});
				\draw [smooth,samples=300,domain=q2:x2]plot(\x,{f(abs(\x))});
				
				\draw [smooth,samples=100,domain=x1:x2,dashed]plot(\x,{g(abs(\x))});
				%		\draw[dashed] (-n/m,y1)--(-n/m,y2);
			\end{scope}
			\draw[->] (0,y1)--(0,y2)node[left]{$y$};
			\draw[->] (x1,0)--(x2,0)node[above]{$x$};
			\filldraw[black](0,0)circle(1.2pt)node[below right]{$O$};
			\draw[dashed]
			%	(-3,0)|-(0,-8)
			;
			
			\foreach \x/\goc in {-3/-90,-1/-90,1/-70,2/-90}{
				\draw[fill=black] (\x,0)circle(1.2pt) node[shift={(\goc:2.8mm)}]{$\x$};
			}
			\foreach \y/\goc in {1,2,-3}{
				\draw[fill=black] (0,\y)circle(1.2pt)node[shift={(\goc:2.8mm)}]{$\y$};
			}
		\end{tikzpicture}
	\end{center}
	}

\end{bt}

\begin{bt}%[2D1H5-2]
% nguon: Archive/2024/TaiLieuGiangDay/1.HamSo/1. K12 - 4in1 - Hamso/ThubaiGoc/Bai5-Dang5.tex #2307

	Cho hàm số $y=log_4 x$ có đồ thị như hình sau. Hãy suy ra đồ thị của hàm số $y=|log_4 x|$.
	\begin{center}
		\begin{tikzpicture}[>=stealth,line join=round,line cap=round,font=\footnotesize,scale=.8]
			\tikzset{declare function={
					x1=-1.5;x2=3.5;y1=-2.5;y2=2.5;% Miền hình vẽ
					a=4;b=-2;c=1; m=1;n=1;
					f(\x)=ln(\x)/ln(a);	
					x0=0.02;
			}}
			\begin{scope}
				\clip (x1,y1)rectangle(x2,y2);
				\draw [smooth,samples=300,domain=x0:x2]plot(\x,{f(\x)});
			\end{scope}
			\draw[->] (0,y1)--(0,y2)node[left]{$y$};
			\draw[->] (x1,0)--(x2,0)node[above]{$x$};
			\filldraw[black](0,0)circle(1.2pt)node[above left]{$O$};
			\draw[dashed];
			
			\foreach \x/\goc in {-1/-90,1/-70,2/-90}{
				\draw[fill=black] (\x,0)circle(1.2pt) node[shift={(\goc:2.8mm)}]{$\x$};
			}
			\foreach \y/\goc in {1/180}{
				\draw[fill=black] (0,\y)circle(1.2pt)node[shift={(\goc:2.8mm)}]{$\y$};
			}
		\end{tikzpicture}
		
	\end{center}
	\loigiai{
	Tập xác định $\mathscr{D}=(0;+\infty)$.\\
	Đồ thị hàm số $y=|log_4 x|$ có được bằng cách:
	\begin{itemize}
		\item Giữ nguyên phần đồ thị của hàm số $y=log_4 x$ nằm trên trục hoành.
		\item Lấy đối xứng qua trục hoành phần đồ thị của $y=log_4 x$ nằm dưới trục hoành.
	\end{itemize}
	Do đó đồ thị của hàm số $y=|log_4 x|$ là	
	\begin{center}
		\begin{tikzpicture}[>=stealth,line join=round,line cap=round,font=\footnotesize,scale=.8]
			\tikzset{declare function={
					x1=-1.5;x2=3.5;y1=-2.5;y2=2.5;% Miền hình vẽ
					a=4;b=-2;c=1; m=1;n=1;
					f(\x)=ln(\x)/ln(a);	
					x0=0.02;
			}}
			\begin{scope}
				\clip (x1,y1)rectangle(x2,y2);
				\draw [smooth,samples=300,domain=x0:x2]plot(\x,{abs(f(\x))});
			\end{scope}
			\draw[->] (0,y1)--(0,y2)node[left]{$y$};
			\draw[->] (x1,0)--(x2,0)node[above]{$x$};
			\filldraw[black](0,0)circle(1.2pt)node[above left]{$O$};
			\draw[dashed];
			
			\foreach \x/\goc in {-1/-90,1/-70,2/-90}{
				\draw[fill=black] (\x,0)circle(1.2pt) node[shift={(\goc:2.8mm)}]{$\x$};
			}
			\foreach \y/\goc in {1/180}{
				\draw[fill=black] (0,\y)circle(1.2pt)node[shift={(\goc:2.8mm)}]{$\y$};
			}
		\end{tikzpicture}
	\end{center}
	}
\end{bt}

\begin{bt}%[2D1H5-2]
% nguon: Archive/2024/TaiLieuGiangDay/1.HamSo/1. K12 - 4in1 - Hamso/ThubaiGoc/Bai5-Dang5.tex #2308
 
	Cho hàm số $y=\left(\dfrac{1}{3}\right)^x$ có đồ thị như hình dưới. Hãy suy ra đồ thị của hàm số $y=\left(\dfrac{1}{3}\right)^{|x|}$?
	\begin{center}
		\begin{tikzpicture}[>=stealth,line join=round,line cap=round,font=\footnotesize,scale=.8]
			\tikzset{declare function={
					x1=-2.5;x2=3.5;y1=-1.5;y2=2.5;% Miền hình vẽ
					a=1/3;
					f(\x)=a^(\x);	
			}}
			\begin{scope}
				\clip (x1,y1)rectangle(x2,y2);
				\draw [smooth,samples=300,domain=x1:x2]plot(\x,{f(\x)});
			\end{scope}
			\draw[->] (0,y1)--(0,y2)node[left]{$y$};
			\draw[->] (x1,0)--(x2,0)node[above]{$x$};
			\filldraw[black](0,0)circle(1.2pt)node[above left]{$O$};
			\draw[dashed];
			
			\foreach \x/\goc in {-2/-90,-1/-90,1/-70,2/-90}{
				\draw[fill=black] (\x,0)circle(1.2pt) node[shift={(\goc:2.8mm)}]{$\x$};
			}
			\foreach \y/\goc in {1/180}{
				\draw[fill=black] (0,\y)circle(1.2pt)node[shift={(\goc:2.8mm)}]{$\y$};
			}
		\end{tikzpicture}
	\end{center}
	\loigiai{
		Tập xác định $\mathscr{D}=\mathbb{R}$.\\
		Ta thấy đồ thị hàm số $y=\left(\dfrac{1}{3}\right)^{|x|}$ có được bằng cách:
		\begin{itemize}
			\item Giữ nguyên phần đồ thị của hàm số $y=\left(\dfrac{1}{3}\right)^x$ nằm bên phải trục $Oy$.
			\item Lấy đối xứng qua trục tung phần đồ thị của $y=\left(\dfrac{1}{3}\right)^x$ nằm bên phải trục $Oy$.
		\end{itemize}
		Do đó đồ thị hàm số $y=\left(\dfrac{1}{3}\right)^{|x|}$ là
		\begin{center}
		\begin{tikzpicture}[>=stealth,line join=round,line cap=round,font=\footnotesize,scale=.8]
			\tikzset{declare function={
					x1=-2.5;x2=3.5;y1=-1.5;y2=2.5;% Miền hình vẽ
					a=1/3;
					f(\x)=a^(\x);	
			}}
			\begin{scope}
				\clip (x1,y1)rectangle(x2,y2);
				\draw [smooth,samples=300,domain=x1:x2]plot(\x,{f(abs(\x))});
			\end{scope}
			\draw[->] (0,y1)--(0,y2)node[left]{$y$};
			\draw[->] (x1,0)--(x2,0)node[above]{$x$};
			\filldraw[black](0,0)circle(1.2pt)node[above left]{$O$};
			\draw[dashed];
			
			\foreach \x/\goc in {-2/-90,-1/-90,1/-70,2/-90}{
				\draw[fill=black] (\x,0)circle(1.2pt) node[shift={(\goc:2.8mm)}]{$\x$};
			}
			\foreach \y/\goc in {1/180}{
				\draw[fill=black] (0,\y)circle(1.2pt)node[shift={(\goc:2.8mm)}]{$\y$};
			}
		\end{tikzpicture}
		\end{center}
	}
	
\end{bt}

\begin{bt}%[2D1V5-2]
% nguon: Archive/2024/TaiLieuGiangDay/1.HamSo/1. K12 - 4in1 - Hamso/ThubaiGoc/Bai5-Dang5.tex #2309
 
	Cho hàm số $y=-x^3+3x^2-2$ có đồ thị như hình sau. Từ đó suy ra đồ thị của hàm số $y=|-|x|^3+3|x|^2-2|$ ?
	\begin{center}
		\begin{tikzpicture}[>=stealth,line join=round,line cap=round,font=\footnotesize,scale=.6]
			\tikzset{declare function={
					x1=-4.5;x2=4.5;y1=-2.5;y2=3.5;
					a=-1;b=3;c=0;d=-2;
					f(\x)=a*(\x)^3+b*(\x)^2+c*(\x)+d;
			}}
			\begin{scope}
				\clip (x1,y1)rectangle(x2,y2);
				\draw [smooth,samples=100,domain=x1:x2]plot(\x,{f(\x)});
			\end{scope}
			\draw[->] (0,y1)--(0,y2)node[left]{$y$};
			\draw[->] (x1,0)--(x2,0)node[above]{$x$};
			\filldraw(0,0)circle(1.2pt)node[below left]{$O$};
			\draw[dashed] 
			%	(1,0)|-(0,-2) (3,0)|-(0,-4)
			;
			\foreach \x/\goc in {-2/-110,-3/-110,-1/-110,1/-70,2/-90,3/-90}{
				\draw[fill=black] (\x,0)circle(1.2pt) node[shift={(\goc:2.8mm)}]{$\x$};
			}
			\foreach \y/\goc in {-2/180,1/180,2/180,3/180}{
				\draw[fill=black] (0,\y)circle(1.2pt)node[shift={(\goc:2.8mm)}]{$\y$};
			}
		\end{tikzpicture}
	\end{center}
	\loigiai{
		Tập xác định $\mathscr{D}=\mathbb{R}$.\\
		Ta thấy đồ thị hàm số $y=|-|x|^3+3|x|^2-2|$ có được bằng cách:
		\begin{itemize}
			\item Giữ nguyên phần đồ thị của hàm số $y=-x^3+3x^2-2$ nằm bên phải trục $Oy$.
			\item Lấy đối xứng qua trục tung phần đồ thị của $y=-x^3+3x^2-2$ nằm bên phải trục $Oy$.
			\item Sau đó, giữ nguyên phần đồ thị nằm phía trên trục hoành.
			\item Lấy đối xứng qua trục hoành phần đồ thị nằm bên dưới trục hoành.
		\end{itemize}
		Do đó đồ thị hàm $y=|-|x|^3+3|x|^2-2|$ là 
		\begin{center}
			\begin{tikzpicture}[>=stealth,line join=round,line cap=round,font=\footnotesize,scale=.6]
				\tikzset{declare function={
						x1=-4.5;x2=4.5;y1=-2.5;y2=3.5;
						a=-1;b=3;c=0;d=-2;
						f(\x)=a*(\x)^3+b*(\x)^2+c*(\x)+d;
				}}
				\begin{scope}
					\clip (x1,y1)rectangle(x2,y2);
					\draw [smooth,samples=500,domain=x1:x2]plot(\x,{abs(f(abs(\x)))});
				\end{scope}
				\draw[->] (0,y1)--(0,y2)node[left]{$y$};
				\draw[->] (x1,0)--(x2,0)node[above]{$x$};
				\filldraw(0,0)circle(1.2pt)node[below left]{$O$};
				\draw[dashed] 
				%	(1,0)|-(0,-2) (3,0)|-(0,-4)
				;
				\foreach \x/\goc in {-2/-110,-3/-110,-1/-110,1/-70,2/-90,3/-90}{
					\draw[fill=black] (\x,0)circle(1.2pt) node[shift={(\goc:2.8mm)}]{$\x$};
				}
				\foreach \y/\goc in {-2/180,1/180,2/180,3/180}{
					\draw[fill=black] (0,\y)circle(1.2pt)node[shift={(\goc:2.8mm)}]{$\y$};
				}
			\end{tikzpicture}
		\end{center}
	}	
\end{bt}

\begin{bt}%[2D1V5-2]
% nguon: Archive/2024/TaiLieuGiangDay/1.HamSo/1. K12 - 4in1 - Hamso/ThubaiGoc/Bai5-Dang5.tex #2310

	Cho hàm số $y=\dfrac{x^2+x-4}{x-1}$ có đồ thị như hình sau. Từ đó suy ra đồ thị của hàm số $y=\left|\dfrac{|x|^2+|x|-4}{|x|-1}\right|$ ?
	\begin{center}
	\begin{tikzpicture}[>=stealth,line join=round,line cap=round,font=\footnotesize,scale=.6]
		\tikzset{declare function={
				x1=-3.5;x2=4.5;y1=-3.5;y2=6.5;
				a=1;b=1;c=-4;m=1;n=-1;
				f(\x)=(a*(\x)^2+b*\x+c)/(m*(\x)+n);
				a1=a/m;b1=b-(a*n)/m*m;
				g(\x)=a1*\x+b1;
				x01=-n/m-.02;x02=-n/m+.02;
		}}
		\begin{scope}
			\clip (x1,y1)rectangle(x2,y2);
			\draw [smooth,samples=100,domain=x1:x01]plot({\x},{f(\x)});
			\draw [smooth,samples=100,domain=x02:x2]plot({\x},{f((\x))});
			\draw [smooth,samples=100,domain=x1:x2,dashed]plot({\x},{g((\x))});
			\draw[dashed] (-n/m,y1)--(-n/m,y2);
		\end{scope}
		\draw[->] (0,y1)--(0,y2)node[left]{$y$};
		\draw[->] (x1,0)--(x2,0)node[above]{$x$}; % Corrected label
		\filldraw(0,0)circle(1.2pt)node[above left]{$O$};
		\foreach \x/\goc in {1/-90,-1/-90}{
			\draw[fill=black] (\x,0)circle(1.2pt) node[shift={(\goc:2.8mm)}]{$\x$};
		}
		\foreach \y/\goc in {1,2,3,4}{
			\draw[fill=black] (0,\y)circle(1.2pt)node[shift={(\goc:2.8mm)}]{$\y$};
		}
	\end{tikzpicture}
	
	\end{center}
	\loigiai{
		Tập xác định $\mathscr{D}=\mathbb{R}\setminus {1}$.\\
		Ta thấy đồ thị hàm số $y=\left|\dfrac{|x|^2+|x|-4}{|x|-1}\right|$ có được bằng cách:
		\begin{itemize}
			\item Giữ nguyên phần đồ thị của hàm số $y=\dfrac{x^2+x-4}{x-1}$ nằm bên phải trục $Oy$.
			\item Lấy đối xứng qua trục tung phần đồ thị của $y=\dfrac{x^2+x-4}{x-1}$ nằm bên phải trục $Oy$.
			\item Sau đó, giữ nguyên phần đồ thị nằm phía trên trục hoành.
			\item Lấy đối xứng qua trục hoành phần đồ thị nằm bên dưới trục hoành.
		\end{itemize}
		Do đó đồ thị hàm $y=\left|\dfrac{|x|^2+|x|-4}{|x|-1}\right|$ là 
		\begin{center}
			\begin{tikzpicture}[>=stealth,line join=round,line cap=round,font=\footnotesize,scale=.6]
				\tikzset{declare function={
						x1=-4.5;x2=4.5;y1=-3.5;y2=6.5;
						a=1;b=1;c=-4;m=1;n=-1;
						f(\x)=(a*(\x)^2+b*\x+c)/(m*(\x)+n);
						a1=a/m;b1=b-(a*n)/m*m;
						g(\x)=a1*\x+b1;
						x01=-n/m-.02;x02=-n/m+.02;
				}}
				\begin{scope}
					\clip (x1,y1)rectangle(x2,y2);
					\draw [smooth,samples=160,domain=x1:x01]plot({\x},{abs(f(abs(\x)))});
					\draw [smooth,samples=200,domain=x02:x2]plot({\x},{abs(f(abs(\x)))});
					\draw [smooth,samples=100,domain=x1:x2,dashed]plot({\x},{g(\x)});
					\draw [smooth,samples=100,domain=x1:x2,dashed]plot({\x},{g(-\x)});
					\draw[dashed] (-n/m,y1)--(-n/m,y2);
					\draw[dashed] (n/m,y1)--(n/m,y2);
				\end{scope}
				\draw[->] (0,y1)--(0,y2)node[left]{$y$};
				\draw[->] (x1,0)--(x2,0)node[above]{$x$}; % Corrected label
				\filldraw(0,0)circle(1.2pt)node[above left]{$O$};
				\foreach \x/\goc in {1/-90,-1/-90}{
					\draw[fill=black] (\x,0)circle(1.2pt) node[shift={(\goc:2.8mm)}]{$\x$};
				}
				\foreach \y/\goc in {1,2,3,4}{
					\draw[fill=black] (0,\y)circle(1.2pt)node[shift={(\goc:2.8mm)}]{$\y$};
				}
			\end{tikzpicture}
			
		\end{center}	
	}	
\end{bt}

\begin{bt}%[2D1N5-2]
% nguon: Archive/2024/TaiLieuGiangDay/1.HamSo/1. K12 - 4in1 - Hamso/ThubaiGoc/Bai5-Dang5.tex #2314

	\immini{Cho hàm số $y=f(x)$ có đồ thị như hình vẽ sau. Hãy xác định khoảng đồng biến của hàm số $y=f(|x|)$.}
	{
		\begin{tikzpicture}[>=stealth,line join=round,line cap=round,font=\footnotesize,scale=.6]
			\tikzset{declare function={
					x1=-3.5;x2=4.5;y1=-2.5;y2=4.5;% Miền hình vẽ
					a=1;b=-1;c=-1;d=1; m=1;n=1;
					f(\x)=a*(\x)^3+b*(\x)^2+c*(\x)+d;	
					a1=a/m;b1=(b-a*n/m)/m;
					g(\x)=a1*\x+b1;
					q1=-(n/m)-.1;q2=-n/m+.1;
			}}
			\begin{scope}
				\clip (x1,y1)rectangle(x2,y2);
				\draw [smooth,samples=100,domain=x1:x2]plot(\x,{f(\x)});
			\end{scope}
			\draw[->] (0,y1)--(0,y2)node[left]{$y$};
			\draw[->] (x1,0)--(x2,0)node[above]{$x$};
			\filldraw(0,0)circle(1.2pt)node[below left]{$O$};
			`	\draw[dashed]
			;
			% Đánh dấu tọa độ trên trục ox và oy
			\foreach \x/\goc in {-3/-90,-1/-90,1/-70,2/-90}{
				\draw[fill=black] (\x,0)circle(1.2pt) node[shift={(\goc:2.8mm)}]{$\x$};
			}
			\foreach \y/\goc in {-2,1,2}{
				\draw[fill=black] (0,\y)circle(1.2pt)node[shift={(\goc:2.8mm)}]{$\y$};
			}
		\end{tikzpicture}
	}
	\loigiai{	
		Ta thấy đồ thị hàm số $y=f(|x|)$ có được bằng cách:
		\begin{itemize}
			\item Giữ nguyên phần đồ thị của hàm số $y=f(x)$ nằm bên phải trục $Oy$.
			\item Lấy đối xứng qua trục tung phần đồ thị của $y=f(x)$ nằm bên phải trục $Oy$.
		\end{itemize}
		Do đó đồ thị hàm số $y=f(|x|)$ là
		\begin{center}
			\begin{tikzpicture}[>=stealth,line join=round,line cap=round,font=\footnotesize,scale=.6]
				\tikzset{declare function={
						x1=-3.5;x2=4.5;y1=-2.5;y2=4.5;% Miền hình vẽ
						a=1;b=-1;c=-1;d=1; m=1;n=1;
						f(\x)=a*(\x)^3+b*(\x)^2+c*(\x)+d;	
						a1=a/m;b1=(b-a*n/m)/m;
						g(\x)=a1*\x+b1;
						q1=-(n/m)-.1;q2=-n/m+.1;
				}}
				\begin{scope}
					\clip (x1,y1)rectangle(x2,y2);
					\draw [smooth,samples=100,domain=x1:x2]plot(\x,{f(abs(\x))});
				\end{scope}
				\draw[->] (0,y1)--(0,y2)node[left]{$y$};
				\draw[->] (x1,0)--(x2,0)node[above]{$x$};
				\filldraw(0,0)circle(1.2pt)node[below left]{$O$};
				`	\draw[dashed]
				;
				% Đánh dấu tọa độ trên trục ox và oy
				\foreach \x/\goc in {-3/-90,-1/-90,1/-70,2/-90}{
					\draw[fill=black] (\x,0)circle(1.2pt) node[shift={(\goc:2.8mm)}]{$\x$};
				}
				\foreach \y/\goc in {-2,1,2}{
					\draw[fill=black] (0,\y)circle(1.2pt)node[shift={(\goc:2.8mm)}]{$\y$};
				}
			\end{tikzpicture}
		\end{center}
		Từ đồ thị hàm số $y=f(\left|x\right|)$ ta thấy hàm số $y=f(\left|x\right|)$ đồng biến trên khoảng $(-1;0)$ và $(1;+\infty)$.	
		
	}
\end{bt}

\begin{bt}%[2D1H5-2]
% nguon: Archive/2024/TaiLieuGiangDay/1.HamSo/1. K12 - 4in1 - Hamso/ThubaiGoc/Bai5-Dang5.tex #2315

	Cho hàm số $y=f(x)=x^3+2x^2+x-1$. Hãy xác định khoảng biến thiên của hàm số $y=f(-x)$.
	
	\loigiai{
		Ta có, $y=f(x)=x^3+2x^2+x-1 \Rightarrow f(-x)=(-x)^3+2(-x)^2+(-x)-1=-x^3+2x^2-x-1$.\\
		$y'=f'(-x)=0\Leftrightarrow -3x^2+4x-1=0 \Leftrightarrow \hoac{&x=\dfrac{1}{3}\\&x=1.}$\\
		Khi đó ta có bảng biến thiên
		\begin{center}
			\begin{tikzpicture}[>=stealth,line join=round,line cap=round,font=\footnotesize,scale=1]
				\tkzTabInit[nocadre=false,lgt=1.2,espcl=3,deltacl=.55]
				{$x$/0.7, $y'$/0.7, $y$/1.8}
				{$-\infty$,$\dfrac{1}{3}$,$1$,$+\infty$}
				\tkzTabLine{,-,$0$,+,$0$,-,}
				\tkzTabVar{+/$\infty$,-/$f\left(\dfrac{1}{3}\right)$,+/$f(1)$,-/$-\infty$}	
			\end{tikzpicture}
			
		\end{center}
		
		Vậy hàm số đồng biến trên khoảng $\left(\dfrac{1}{3};1\right)$; nghịch biến trên khoảng $\left(-\infty;\dfrac{1}{3}\right)$ và $\left(1;+\infty\right)$.}
\end{bt}

\begin{bt}%[2D1H5-2]
% nguon: Archive/2024/TaiLieuGiangDay/1.HamSo/1. K12 - 4in1 - Hamso/ThubaiGoc/Bai5-Dang5.tex #2316

	\immini{
		Cho hàm số $ y=\dfrac{2x-1}{x-1}$ có đồ thị như hình vẽ. Tìm các khoảng đồng biến của hàm số $ y=\dfrac{2\left|x\right|-1}{\left|x\right|-1}$. 
	}
	{\begin{tikzpicture}[>=stealth,line join=round,line cap=round,font=\footnotesize,scale=.6]
			\tikzset{declare function={
					x1=-3.5;x2=4.5;y1=-1.5;y2=6.5;% Miền hình vẽ
					a=2;b=-1; m=1;n=-1;
					f(\x)=(a*(\x)+b)/(m*(\x)+n);	
					a1=a/m;b1=(b-a*n/m)/m;
					g(\x)=a/m;
					q1=-(n/m)-.1;q2=-n/m+.1;
			}}
			\begin{scope}
				\clip (x1,y1)rectangle(x2,y2);
				\draw [smooth,samples=300,domain=x1:q1]plot(\x,{f(\x)});
				\draw [smooth,samples=300,domain=q2:x2]plot(\x,{f(\x)});
				\draw [smooth,samples=100,domain=x1:x2]plot(\x,{g(\x)});
				\draw (-n/m,y1)--(-n/m,y2);
			\end{scope}
			\draw[->] (0,y1)--(0,y2)node[left]{$y$};
			\draw[->] (x1,0)--(x2,0)node[above]{$x$};
			\filldraw(0,0)circle(1.2pt)node[below left]{$O$};
			\draw[dashed]
			;
			\foreach \x/\goc in {1/-60}{
				\draw[fill=black] (\x,0)circle(1.2pt) node[shift={(\goc:2.8mm)}]{$\x$};
			}
			\foreach \y/\goc in {1/180,2/150}{
				\draw[fill=black] (0,\y)circle(1.2pt)node[shift={(\goc:2.8mm)}]{$\y$};
			}
		\end{tikzpicture}
	}
	\loigiai{
		\immini {
			Đồ thị hàm số $ y=\dfrac{2\left|x\right|-1}{\left|x\right|-1}$ có được bằng cách lấy đối xứng phần đồ thị phía bên phải trục tung qua trục $Oy$ và xóa phần đồ thị phía bên trái trục tung của đồ thị hàm số $ y=\dfrac{2x-1}{x-1}$.\\
			Dựa vào đồ thị ta thấy hàm số $ y=|x|^3-3|x|+1 $ đồng biến trên khoảng $(-\infty;-1)$ và $(-1;0)$.
		}
		{\begin{tikzpicture}[>=stealth,line join=round,line cap=round,font=\footnotesize,scale=.6]
				\tikzset{declare function={
						x1=-3.5;x2=4.5;y1=-1.5;y2=6.5;% Miền hình vẽ
						a=2;b=-1; m=1;n=-1;
						f(\x)=(a*(\x)+b)/(m*(\x)+n);	
						a1=a/m;b1=(b-a*n/m)/m;
						g(\x)=a/m;
						q1=-(n/m)-.1;q2=-n/m+.1;
				}}
				\begin{scope}
					\clip (x1,y1)rectangle(x2,y2);
					\draw [smooth,samples=300,domain=x1:x2]plot(\x,{f(abs(\x))});
					\draw [smooth,samples=300,domain=q2:x2]plot(\x,{f(abs(\x))});
					\draw [smooth,samples=100,domain=x1:x2]plot(\x,{g(\x)});
					\draw (-n/m,y1)--(-n/m,y2);
				\end{scope}
				\draw[->] (0,y1)--(0,y2)node[left]{$y$};
				\draw[->] (x1,0)--(x2,0)node[above]{$x$};
				\filldraw(0,0)circle(1.2pt)node[below left]{$O$};
				\draw[dashed]
				;
				\foreach \x/\goc in {1/-60,-1/-130}{
					\draw[fill=black] (\x,0)circle(1.2pt) node[shift={(\goc:2.8mm)}]{$\x$};
				}
				\foreach \y/\goc in {1/180,2/150}{
					\draw[fill=black] (0,\y)circle(1.2pt)node[shift={(\goc:2.8mm)}]{$\y$};
				}
		\end{tikzpicture}}
	}
\end{bt}

\begin{bt}%[2D1V5-2]
% nguon: Archive/2024/TaiLieuGiangDay/1.HamSo/1. K12 - 4in1 - Hamso/ThubaiGoc/Bai5-Dang5.tex #2317

	\immini{
		Cho đồ thị của hàm số $y=f(x)=\log_4 x$ có đồ thị như hình vẽ. Từ đó, hãy suy ra đồ thị của hàm số $y=\left|f(x-2)\right|$.
	}{
		\begin{tikzpicture}[>=stealth,line join=round,line cap=round,font=\footnotesize,scale=.8]
			\tikzset{declare function={
					x1=-2.2;x2=4.5;y1=-2.5;y2=2.5;
					a=4;b=0;
					f(\x)=ln(\x+b)/ln(a);
					x0=-b+.02;
			}}
			\begin{scope}
				\clip (x1,y1)rectangle(x2,y2);
				\draw [smooth,samples=200,domain=x0:x2]plot(\x,{f(\x)});
			\end{scope}
			\draw[->] (0,y1)--(0,y2)node[left]{$y$};
			\draw[->] (x1,0)--(x2,0)node[above]{$x$};
			\filldraw(0,0)circle(1.2pt)node[below left]{$O$};
			\foreach \x/\goc in {1/90}{
				\draw[fill=black] (\x,0)circle(1.2pt) node[shift={(\goc:2.8mm)}]{$\x$};
			}
			\foreach \y/\goc in {}{
				\draw[fill=black] (0,\y)circle(1.2pt)node[shift={(\goc:2.8mm)}]{$\y$};
			}
		\end{tikzpicture}	
	}	
	\loigiai{	
		Trước tiên ta vẽ đồ thị hàm số $y=\log_4 x$.\\	
		Ta thấy đồ thị hàm số $y=|\log_4x|$ có được bằng cách:
		\begin{itemize}
			\item Giữ nguyên phần đồ thị của hàm số $y=\log_4x$ nằm trên trục hoành.
			\item Lấy đối xứng qua trục hoành phần đồ thị của $y=\log_4x$ nằm dưới trục hoành.
		\end{itemize}
		Do đó đồ thị hàm số $y=|\log_4x|$ là
		\begin{center}
			\begin{tikzpicture}[>=stealth,line join=round,line cap=round,font=\footnotesize,scale=.8]
				\tikzset{declare function={
						x1=-2.2;x2=4.5;y1=-2.5;y2=2.5;
						a=4;b=0;
						f(\x)=ln(\x+b)/ln(a);
						x0=-b+.02;
				}}
				\begin{scope}
					\clip (x1,y1)rectangle(x2,y2);
					\draw [smooth,samples=200,domain=x0:x2]plot(\x,{abs(f(\x))});
				\end{scope}
				\draw[->] (0,y1)--(0,y2)node[left]{$y$};
				\draw[->] (x1,0)--(x2,0)node[above]{$x$};
				\filldraw(0,0)circle(1.2pt)node[below left]{$O$};
				\foreach \x/\goc in {1/90}{
					\draw[fill=black] (\x,0)circle(1.2pt) node[shift={(\goc:2.8mm)}]{$\x$};
				}
				\foreach \y/\goc in {}{
					\draw[fill=black] (0,\y)circle(1.2pt)node[shift={(\goc:2.8mm)}]{$\y$};
				}
			\end{tikzpicture}
		\end{center}
		Đồ thị hàm số $y=\left|\log_4(x-2)\right|$ có được bằng cách tịnh tiến đồ thị hàm số $y=\left|\log_4x\right|$ sang phải $2$ đơn vị theo phương của trục $Ox$.
		\begin{center}
			\begin{tikzpicture}[>=stealth,line join=round,line cap=round,font=\footnotesize,scale=.8]
				\tikzset{declare function={
						x1=-2.2;x2=4.5;y1=-2.5;y2=2.5;
						a=4;b=-2;
						f(\x)=ln(\x+b)/ln(a);
						x0=-b+.02;
				}}
				\begin{scope}
					\clip (x1,y1)rectangle(x2,y2);
					\draw [smooth,samples=200,domain=x0:x2]plot(\x,{abs(f(\x))});
				\end{scope}
				\draw[->] (0,y1)--(0,y2)node[left]{$y$};
				\draw[->] (x1,0)--(x2,0)node[above]{$x$};
				\filldraw(0,0)circle(1.2pt)node[below left]{$O$};
				\draw (2,y1)--(2,y2);
				\foreach \x/\goc in {2/-110,3/-90}{
					\draw[fill=black] (\x,0)circle(1.2pt) node[shift={(\goc:2.8mm)}]{$\x$};
				}
				\foreach \y/\goc in {}{
					\draw[fill=black] (0,\y)circle(1.2pt)node[shift={(\goc:2.8mm)}]{$\y$};
				}
			\end{tikzpicture}
		\end{center}	
	}
\end{bt}

\begin{bt}%[2D1H2-6]
% nguon: Archive/2024/TaiLieuGiangDay/1.HamSo/1. K12 - 4in1 - Hamso/ThubaiGoc/Bai2-Dang6.tex #2328
[MĐ2]
	Cho hàm số $f(x)=\dfrac{x^2-2mx+1}{x-1}$. Tìm tất cả các giá trị của $m$ để hàm số $f(x)$ có đúng hai điểm cực trị.
	\loigiai{
		Ta có $f'(x)=\dfrac{x^2-2x+2m-1}{(x-1)^2}$, $\forall x\ne 1$. Do đó $f(x)$ có hai điểm cực trị khi và chỉ khi phương trình
		$$x^2-2x+2m-1=0$$
		 có hai nghiệm phân biệt khác $1$ hay
		\[ \heva{&\Delta'=1-2m+1> 0 \\ &1^2-2\cdot 1+2m-1\ne 0}\Leftrightarrow \heva{&m< 1\\ &m\ne 1}\Leftrightarrow m<1. \] 
	}
\end{bt}

\begin{bt}%[2D1H2-6]
% nguon: Archive/2024/TaiLieuGiangDay/1.HamSo/1. K12 - 4in1 - Hamso/ThubaiGoc/Bai2-Dang6.tex #2329
[MĐ2]
	Tìm tất cả các giá trị của tham số $m$ để hàm số $f(x)=\dfrac{x^2-mx}{x+1}$ không có điểm cực trị nào.
	\loigiai{
		Ta có $f'(x)= \dfrac{x^2+2x-m}{(x+1)^2}$, $\forall x\ne -1$. Hàm số $f(x)$ không có điểm cực trị khi và chỉ khi phương trình
		\[  x^2+2x-m=0\]
		vô nghiệm hoặc có nghiệm kép hay
		\[\Delta'=1+m\le 0 \Leftrightarrow m\le -1. \]
	}
\end{bt}

\begin{bt}%[2D1V2-6]
% nguon: Archive/2024/TaiLieuGiangDay/1.HamSo/1. K12 - 4in1 - Hamso/ThubaiGoc/Bai2-Dang6.tex #2330
[MĐ3]
	Cho hàm số $f(x)=x+1+\dfrac{m}{x}$, với $m$ là tham số thực. Tìm tất cả các giá trị của $m$ sao cho $f(x)$ có hai điểm cực trị trái dấu.
	\loigiai{
		Ta có $f'(x)=\dfrac{x^2-m}{x^2}$, $\forall x\ne 0$. Hàm số $f(x)$ có hai điểm cực trị khi và chỉ khi phương trình $x^2-m=0$ có hai nghiệm phân biệt hay $m>0$. Khi đó, nghiệm nói trên là trái dấu. Vậy $m>0$ là tất cả các giá trị cần tìm.
	}
\end{bt}

\begin{bt}%[2D1V2-6]
% nguon: Archive/2024/TaiLieuGiangDay/1.HamSo/1. K12 - 4in1 - Hamso/ThubaiGoc/Bai2-Dang6.tex #2331
[MĐ3]
	Tìm tất cả các giá trị của tham số $m$ để hàm số $f(x)=\dfrac{x^2-3x+2m}{x-2}$ có hai điểm cực trị dương.
	\loigiai{
		Ta có $f'(x)=\dfrac{x^2-4x+6-2m}{(x-2)^2}$, $\forall x\ne 2$. Khi đó yêu cầu bài toán tương đương phương trình
		\[x^2-4x+6-2m=0 \]
		có hai nghiệm dương phân biệt khác $2$ hay
		\[ \heva{&\Delta'=4-6+2m>0\\ &6-2m>0\\ &4-4\cdot 2+6-2m\ne 0} \Leftrightarrow \heva{&m>1\\ &m<3\\ &m\ne 1} \Leftrightarrow 1<m<3. \] 
	}
\end{bt}

\begin{bt}%[2D1V2-6]
% nguon: Archive/2024/TaiLieuGiangDay/1.HamSo/1. K12 - 4in1 - Hamso/ThubaiGoc/Bai2-Dang6.tex #2332
[MĐ3]
	Tìm tất cả các giá trị của tham số $m$ để hàm số $f(x)=\dfrac{x^2-2mx+1}{x-1}$ có một điểm cực trị là $x_0=2$.
	\loigiai{
		Ta có $f'(x)=\dfrac{x^2-2x+2m-1}{(x-1)^2}$, với mọi $x\ne 1$. Do $x_0=2$ là điểm cực trị của hàm số $f(x)$ nên
		\[ f'(2)=0 \Rightarrow 4-2\cdot 2+2m-1=0 \Rightarrow m=\dfrac{1}{2}. \]
		Thử lại, $m=\dfrac{1}{2}$ thỏa mãn yêu cầu.
	}
\end{bt}

\begin{bt}%[2D1V2-6]
% nguon: Archive/2024/TaiLieuGiangDay/1.HamSo/1. K12 - 4in1 - Hamso/ThubaiGoc/Bai2-Dang6.tex #2333
[MĐ3]
	Tìm tất cả các giá trị của tham số $m$ để hàm số $f(x)=\dfrac{x^2-mx}{(x+1)^2}$ nhận $x_0=1$ là điểm cực tiểu.
	\loigiai{
		Ta có $f'(x)=\dfrac{x^2+2x-m}{(x+1)^2}$, $\forall x\ne -1$. Do $x_0=1$ là điểm cực trị của hàm số $f(x)$ nên
		\[ f'(1)=0 \Rightarrow 3-m=0\Rightarrow m=3. \]
		Với $m=3$, ta có $f'(x)=\dfrac{x^2+2x-3}{(x+1)^2}$. Suy ra bảng xét dấu của $f'(x)$
		\begin{center}
			\begin{tikzpicture}[font=\footnotesize,>=stealth, scale=1]
				\tkzTabInit[nocadre=false,lgt=1.2,espcl=2.5,deltacl=0.6]
				{$x$ /0.6,$f'(x)$ /0.6}
				{$-\infty$, $-3$, $-1$, $1$, $+\infty$}
				\tkzTabLine{,+,0,-,d,-,0,+,}
			\end{tikzpicture}
		\end{center}
		Từ bảng xét dấu, $m=3$ thỏa mãn yêu cầu. Vậy $m=3$ là giá trị cần tìm.
	}
\end{bt}

\begin{bt}%[2D1V2-6]
% nguon: Archive/2024/TaiLieuGiangDay/1.HamSo/1. K12 - 4in1 - Hamso/ThubaiGoc/Bai2-Dang6.tex #2334
[MĐ3]
	Tìm tất cả các giá trị của tham số $m$ để đồ thị hàm số $f(x)=x+\dfrac{m}{x}$ có hai điểm cực trị $A$ và $B$ sao cho $AB=2\sqrt{5}$.
	\loigiai{
		Ta có $f'(x)=\dfrac{x^2-m}{x^2}$, $\forall x\ne 0$. Khi đó, đồ thị hàm số $f(x)$ có hai điểm cực trị khi và chỉ khi phương trình $x^2-m=0$ có hai nghiệm phân biệt khác $0$ hay $m>0$.
		
		Khi đó, đồ thị hàm số $f(x)$ có hai điểm cực trị là $A\left(-\sqrt{m},-2\sqrt{m}\right)$ và $B\left(\sqrt{m},2\sqrt{m}\right)$. Do đó
		\[ 20=AB^2=\left(2\sqrt{m}\right)^2+(4\sqrt{m})^2=20m \Leftrightarrow m=1.  \]
		Vậy $m=1$ là giá trị cần tìm.
	}
\end{bt}

\begin{bt}%[2D1H2-6]
% nguon: Archive/2024/TaiLieuGiangDay/1.HamSo/1. K12 - 4in1 - Hamso/ThubaiGoc/Bai2-Dang6.tex #2338
[MĐ2]
	Cho hàm số $f(x)=\dfrac{2x^2+3x+m-2}{x+2}$. Tìm $m$ để hàm số có hai điểm cực trị.
		\loigiai{
			 Ta có  $f'(x)=\dfrac{2x^2+8x+8-m}{(x+2)^2}$, $\forall x\ne -2$. Hàm số có hai điểm cực trị khi và chỉ khi phương trình $2x^2+8x+8-m=0$ có hai nghiệm phân biệt hay
				\[ \Delta'=16-2(8-m)>0 \Leftrightarrow  m>0. \]
		}
\end{bt}

\begin{bt}%[2D1H2-6]
% nguon: Archive/2024/TaiLieuGiangDay/1.HamSo/1. K12 - 4in1 - Hamso/ThubaiGoc/Bai2-Dang6.tex #2339
[MĐ2]
	Cho hàm số $f(x)=\dfrac{x^2+(m+1)x+m+1}{x+1}$ có đồ thị $(C_m)$. Chứng minh rằng $(C_m)$ luôn có hai điểm cực trị với mọi $m$, đồng thời khoảng cách giữa hai điểm đó bằng $\sqrt{20}$.
	\loigiai{
		Ta có $f'(x)=\dfrac{x^2+2x}{(x+1)^2}$, suy ra $f'(x)=0$ có hai nghiệm phân biệt $x_1=0$, $x_2=-2$ nên đồ thị $(C_m)$ luôn có hai điểm cực trị $A(0;m+1)$ và $B(-2;m-3)$. Hơn nữa
		\[ AB=\sqrt{(0+2)^2+(m+1-(m-3)^2)}=\sqrt{20}. \]
	}
\end{bt}

\begin{bt}%[2D1H2-6]
% nguon: Archive/2024/TaiLieuGiangDay/1.HamSo/1. K12 - 4in1 - Hamso/ThubaiGoc/Bai2-Dang6.tex #2340
[MĐ2]
	Tìm tất cả các giá trị của $m$ để hàm số $f(x)=\dfrac{x^2+x-m}{x-1}$ có điểm cực trị là $x=2$.
	\loigiai{
		Ta có $f'(x)=\dfrac{x^2-2x+m-1}{(x-1)^2}$, với mọi $x\ne 1$. Do đó yêu cầu bài toán tương đương
		\[ f'(2)=0 \Leftrightarrow 4-4+m-1=0 \Leftrightarrow m=1. \] 
	}
\end{bt}

\begin{bt}%[2D1V2-6]
% nguon: Archive/2024/TaiLieuGiangDay/1.HamSo/1. K12 - 4in1 - Hamso/ThubaiGoc/Bai2-Dang6.tex #2341
[MĐ2]
	Cho hàm số $f'(x)=\dfrac{x^2-mx+m}{x-m}$. Tìm tất cả các giá trị của $m$ để hàm số có điểm cực tiểu là $x=3$.
	\loigiai{
		Ta có $f'(x)=\dfrac{x^2-2mx+m^2-m}{(x-m)^2}$, với mọi $x\ne m$. Do $x=2$ là điểm cực trị của hàm số $f(x)$ nên
		\[ f'(3)=0 \Leftrightarrow \frac{4-4m+m^2-m}{(2-m)^2} \Leftrightarrow \hoac{&m=1\\ &m=4.}  \]
		Với $m=1$ thì $f'(x)=\dfrac{x(x-2)}{(x-1)^2}$ nên $f'(x)$ đổi dấu từ âm sang dương khi $x$ đi qua $2$ hay $x=2$ là điểm cực tiểu của hàm số $f(x)$.
		
		Với $m=4$ thì $f'(x)=\dfrac{(x-2)(x-6)}{(x-4)^2}$ nên $f'(x)$ đổi dấu từ dương sang âm khi $x$ qua $2$ hay $x=2$ là điểm cực đại của hàm số $f(x)$.
		
		Vậy $m=1$ là giá trị cần tìm.
	}
\end{bt}

\begin{bt}%[2D1V2-6]
% nguon: Archive/2024/TaiLieuGiangDay/1.HamSo/1. K12 - 4in1 - Hamso/ThubaiGoc/Bai2-Dang6.tex #2342
[MĐ3]
	Cho hàm số $f(x)=mx+\dfrac{1}{x}$ có đồ thị $(C_m)$. Tìm $m$ để hàm số có cực trị và khoảng cách từ điểm cực tiểu đến đường thẳng $\Delta:y=mx$ bằng $\dfrac{1}{\sqrt{2}}$.
	\loigiai{
		Ta có $f'(x)=m-\dfrac{1}{x^2}$, với mọi $x\ne 0$. Suy ra đồ thị hàm số có hai điểm cực trị khi và chỉ khi phương trình
		\[ f'(x)=0 \Leftrightarrow m=\frac{1}{x^2} \]
		có hai nghiệm phân biệt hay $m>0$. Khi đó hai điểm cực trị của hàm số là $-\dfrac{1}{\sqrt{m}}$ và $\dfrac{1}{\sqrt{m}}$. Ta có bảng xét dấu của $f'(x)$
		\begin{center}
			\begin{tikzpicture}[line cap=round, line join=round,font=\footnotesize,>=stealth, scale=1]
				\tkzTabInit[nocadre=false,lgt=1.2,espcl=2.5,deltacl=0.6]
				{$x$ /1,$f'(x)$ /0.6}
				{$-\infty$, $-\dfrac{1}{\sqrt{m}}$, $0$, $\dfrac{1}{\sqrt{m}}$, $+\infty$}
				\tkzTabLine{,+,0,-,d,-,0,+,}
			\end{tikzpicture}
		\end{center}
		Suy ra điểm cực tiểu của $(C_m)$ là $A\left(\dfrac{1}{\sqrt{m}},2\sqrt{m}\right)$. Suy ra
		\[ \dfrac{1}{\sqrt{2}}=\mathrm{d}(A,\Delta)=\frac{\left| \sqrt{m}-2\sqrt{m} \right|}{\sqrt{m^2+1}}=\frac{\sqrt{m}}{\sqrt{m^2+1}}. \]
		Giải phương trình, kết hợp với $m>0$,  ta được $m=1$.		 
	}
\end{bt}

\begin{bt}%[2D1N4-1]
% nguon: Archive/2024/TaiLieuGiangDay/1.HamSo/1. K12 - 4in1 - Hamso/ThubaiGoc/Bai4-Dang1.tex #2353
%[BG-12NEW-4in1, Nguyen Huynh]
	Tìm phương trình đường tiệm cận ngang của đồ thị hàm số $y = \dfrac{x + 1}{2 - x}$.
	\loigiai{
		Ta có $\displaystyle\lim_{x\rightarrow + \infty}\dfrac{x + 1}{2 - x} = \lim\limits_{x\rightarrow + \infty}\dfrac{1 + \dfrac{1}{x}}{\dfrac{2}{x} - 1} = - 1$ và $\lim\limits_{x\rightarrow - \infty}\dfrac{x + 1}{2 - x} = \lim\limits_{x\rightarrow - \infty}\dfrac{1 + \dfrac{1}{x}}{\dfrac{2}{x} - 1} = - 1$  nên đường thẳng $y = -1$ là tiệm cận ngang của đồ thị hàm số.
	}
\end{bt}

\begin{bt}%[2D1N4-1]
% nguon: Archive/2024/TaiLieuGiangDay/1.HamSo/1. K12 - 4in1 - Hamso/ThubaiGoc/Bai4-Dang1.tex #2354
%[BG-12NEW-4in1, Nguyen Huynh]
	Tìm phương trình đường tiệm cận xiên của đồ thị hàm số $y=f(x)=\dfrac{-x^2-3x-3}{x+1}$.
	\loigiai{
		Ta có $y=f(x)=\dfrac{-x^2-3x-3}{x+1}=-x-2-\dfrac{1}{x+1}$.\\
		Mà $\lim\limits_{x\to+\infty}\left(-\dfrac{1}{x+1}\right)=0$ nên đường thẳng $y=-x-2$ là tiệm cận xiên của đồ thị hàm số đã cho.
	}
\end{bt}

\begin{bt}%[2D1H4-1]
% nguon: Archive/2024/TaiLieuGiangDay/1.HamSo/1. K12 - 4in1 - Hamso/ThubaiGoc/Bai4-Dang1.tex #2355
%[BG-12NEW-4in1, Nguyen Huynh]
	Tính khoảng cách giữa hai tiệm cận đứng của đồ thị hàm số $y = \dfrac{1}{x^2 - 2}$.
	\loigiai{
		Hàm số $y = \dfrac{1}{x^2 - 2}$ có hai tiệm cận đứng là $x = \sqrt{2}$ và $x = - \sqrt{2}$.
		\\ Ta có khoảng cách giữa hai đường thẳng này bằng $2\sqrt{2}$.
	}
\end{bt}

\begin{bt}%[2D1H4-1]
% nguon: Archive/2024/TaiLieuGiangDay/1.HamSo/1. K12 - 4in1 - Hamso/ThubaiGoc/Bai4-Dang1.tex #2356
%[BG-12NEW-4in1, Nguyen Huynh]
	Tổng số tiệm cận đứng và tiệm cận ngang của đồ thị hàm số $y=\dfrac{2x^2-5x+2}{x^2-4}$ bằng bao nhiêu?
	\loigiai{
		Vì $\lim\limits_{x  \rightarrow +\infty} \dfrac{2x^2-5x+2}{x^2-4}=2$ và $\lim\limits_{x  \rightarrow -\infty} \dfrac{2x^2-5x+2}{x^2-4}=2$ nên $y=2$ là tiệm cận ngang.\\
		Lại có $x^2-4=0 \Leftrightarrow x=\pm 2$ nên
		\begin{itemize}
			\item Vì $\lim\limits_{x  \rightarrow 2} \dfrac{2x^2-5x+2}{x^2-4}=\lim\limits_{x  \rightarrow 2} \dfrac{(2x-1)(x-2)}{(x-2)(x+2)}=\lim\limits_{x  \rightarrow 2} \dfrac{2x-1}{x+2}=\dfrac{3}{4}$ nên $x=2$ \textbf{không} là tiệm cận đứng.
			\item $\lim\limits_{x  \rightarrow -2^+} \dfrac{2x^2-5x+2}{x^2-4}=\lim\limits_{x  \rightarrow -2^+} \dfrac{(2x-1)(x-2)}{(x-2)(x+2)}=\lim\limits_{x  \rightarrow -2^+} \dfrac{2x-1}{x+2}=-\infty$.
			\item $\lim\limits_{x  \rightarrow -2^-} \dfrac{2x^2-5x+2}{x^2-4}=\lim\limits_{x  \rightarrow -2^-} \dfrac{(2x-1)(x-2)}{(x-2)(x+2)}=\lim\limits_{x  \rightarrow -2-} \dfrac{2x-1}{x+2}=+\infty$.
		\end{itemize}
		Do đó ta có $x=-2$ là tiệm cận đứng của đồ thị hàm số.\\
		Vậy đồ thị hàm số có tổng số tiệm cận đứng và tiệm cận ngang là $2$.}
\end{bt}

\begin{bt}%[2D1H4-1]
% nguon: Archive/2024/TaiLieuGiangDay/1.HamSo/1. K12 - 4in1 - Hamso/ThubaiGoc/Bai4-Dang1.tex #2357
%[BG-12NEW-4in1, Nguyen Huynh]
	Tổng số tiệm cận đứng và tiệm cận ngang của đồ thị hàm số $ y= \dfrac{5x^2-4x-1}{x^2-1} $ bằng bao nhiêu?
	\loigiai{
		Tập xác định của hàm số $ \mathscr{D}=\mathbb{R}\setminus \{-1;1\} $.\\
		Ta có
		\begin{itemize}
			\item Vì  $  \lim\limits_{x \rightarrow +\infty} \dfrac{5x^2-4x-1}{x^2-1} = 5 $ và $  \lim\limits_{x \rightarrow -\infty} \dfrac{5x^2-4x-1}{x^2-1} = 5 $ nên đường thẳng $  y = 5  $ là đường tiệm cận ngang của đồ thị hàm số.
			\item $  \lim\limits_{x \rightarrow -1^+} \dfrac{5x^2-4x-1}{x^2-1} =  \lim\limits_{x \rightarrow -1^+ } \dfrac{5x+1}{x+1} = -\infty $ suy ra đường thẳng $x =-1  $ là đường tiệm cận đứng của đồ thị hàm số.
		\end{itemize}
		Vậy đồ thị hàm số đã cho có hai đường tiệm cận.
	}
\end{bt}

\begin{bt}%[2D1H4-1]
% nguon: Archive/2024/TaiLieuGiangDay/1.HamSo/1. K12 - 4in1 - Hamso/ThubaiGoc/Bai4-Dang1.tex #2358
%[BG-12NEW-4in1, Nguyen Huynh]
	Tìm giao điểm $I$ giữa hai tiệm cận của đồ thị hàm số $y=\dfrac{2{{x}^{2}}-x+1}{x-1}$.
	\loigiai{
		Hàm số có $\mathscr{D}=\mathbb{R}\setminus\{1\}$.
		Ta xét
		\begin{itemize}
			\item Vì $\lim\limits_{x\rightarrow 1^+}\dfrac{2{{x}^{2}}-x+1}{x-1}=+\infty
			$ nên $x =1$ là tiệm cận đứng của hàm số.
			\item $\lim\limits_{x\rightarrow +\infty}\left[\dfrac{2{{x}^{2}}-x+1}{x-1} -(2x+1)\right]
			=\lim\limits_{x\rightarrow +\infty}{\dfrac{2}{x-1}}=0 \Rightarrow y=2x+1$ là tiệm cận xiên của hàm số.
		\end{itemize}	
		Mà hai tiệm cận này cắt nhau tại $I\ (1;3)$.
	}
\end{bt}

\begin{bt}%[2D1N4-1]
% nguon: Archive/2024/TaiLieuGiangDay/1.HamSo/1. K12 - 4in1 - Hamso/ThubaiGoc/Bai4-Dang1.tex #2359
%[BG-12NEW-4in1, Nguyen Huynh]
	Tìm các tiệm cận đứng và tiệm cận xiên của đồ thị hàm số $y=f(x)=\dfrac{x^2-4 x+2}{1-x}$.	
	\loigiai{
		Ta có $\lim\limits _{x \rightarrow1^+}f(x)=\lim\limits _{x \rightarrow1^+}\dfrac{x^2-4 x+2}{1-x}=-\infty$.\\ Tương tự $\lim\limits _{x \rightarrow1^-}f(x)=+\infty$.\\
		Ta có
		$$
		\begin{aligned}
			& a=\lim _{x \rightarrow+\infty} \dfrac{f(x)}{x}=\lim _{x \rightarrow+\infty} \dfrac{x^2-4 x+2}{x-x^2}=-1 ; \\
			& b=\lim _{x \rightarrow+\infty}[f(x)+x]=\lim _{x \rightarrow+\infty} \dfrac{-3 x+2}{1-x}=3 .
		\end{aligned}
		$$
		Tương tự, $\lim\limits _{x \rightarrow-\infty} \dfrac{f(x)}{x}=-1, \lim\limits _{x \rightarrow-\infty}[f(x)+x]=3$.\\
		Vậy đồ thị hàm số $f(x)$ có tiệm cận đứng là đường thẳng $x=1$ tiệm cận xiên là đường thẳng $y=-x+3$.
	}
\end{bt}

\begin{bt}%[2D1H4-1]
% nguon: Archive/2024/TaiLieuGiangDay/1.HamSo/1. K12 - 4in1 - Hamso/ThubaiGoc/Bai4-Dang1.tex #2363
%[BG-12NEW-4in1, Nguyen Huynh]
	Tìm các tiệm cận xiên của đồ thị hàm số $y=\dfrac{2x^{2}+9x+11}{2x+5}$.
	\loigiai{ Ta có $\begin{aligned}[t]
			a&=\lim\limits_{x \rightarrow+\infty} \dfrac{f(x)}{x}=\lim\limits_{x \rightarrow+\infty} \dfrac{2x+9+\dfrac{11}{x}}{2x+5}=1;\\
			b&=\lim\limits_{x \rightarrow+\infty}[f(x)-ax]=\lim\limits_{x \rightarrow+\infty}\left(\dfrac{2x^{2}+9x+11}{2x+5}-x\right)=\lim\limits_{x \rightarrow+\infty} \dfrac{4x+11}{2x+5}=2.
		\end{aligned}$\\
		Ta cũng có $\lim\limits_{x \rightarrow-\infty} \dfrac{f(x)}{x}=1$; $\lim\limits_{x \rightarrow-\infty}[f(x)-x]=2$.\\
		Do đó, đồ thị hàm số có tiệm cận xiên là đường thẳng $y=x+2$.
		
	}
\end{bt}

\begin{bt}%[2D1N4-1]
% nguon: Archive/2024/TaiLieuGiangDay/1.HamSo/1. K12 - 4in1 - Hamso/ThubaiGoc/Bai4-Dang1.tex #2364
%[BG-12NEW-4in1, Nguyen Huynh]
	Tìm đường tiệm cận đứng của đồ thị hàm số $y=\dfrac{\sqrt{9x^2+6x+4}}{x+2}$.
	\loigiai{
		Ta có $\displaystyle\lim_{x\rightarrow -2^{\pm}}y=\displaystyle\lim_{x\rightarrow -2^{\pm}}\dfrac{\sqrt{9x^2+6x+4}}{x+2}=
		\pm\infty$ nên đồ thị có một đường tiệm cận đứng $x=-2$.
	}
\end{bt}

\begin{bt}%[2D1V4-1]
% nguon: Archive/2024/TaiLieuGiangDay/1.HamSo/1. K12 - 4in1 - Hamso/ThubaiGoc/Bai4-Dang1.tex #2365
%[BG-12NEW-4in1, Nguyen Huynh]
	Tìm tất cả các giá trị thực của tham số $m$ để đồ thị hàm số $y=\dfrac{mx-8}{x+2}$ có tiệm cận đứng.
	\loigiai{
		Ta có $\lim\limits_{x \to (-2)^{-}} (x+2)=0$ và $\lim\limits_{x \to (-2)^{-}} (mx-8)=-2m-8$; $\lim\limits_{x \to (-2)^{+}} (x+2)=0$ và $\lim\limits_{x \to (-2)^{+}} (mx-8)=-2m-8$.
		\\Do đó, đồ thị hàm số $y=\dfrac{mx-8}{x+2}$ có tiệm cận đứng $ \Leftrightarrow -2m-8 \ne 0 \Leftrightarrow m \ne -4$.}
\end{bt}

\begin{bt}%[2D1V4-1]
% nguon: Archive/2024/TaiLieuGiangDay/1.HamSo/1. K12 - 4in1 - Hamso/ThubaiGoc/Bai4-Dang1.tex #2366
%[BG-12NEW-4in1, Nguyen Huynh]
	Tính tổng tất cả các giá trị của tham số $m$ để đồ thị hàm số $y=\dfrac{x+1}{x^2-2x+m}$ có đúng hai đường tiệm cận.
	\loigiai{
		Ta có $\displaystyle\lim_{x\rightarrow\infty} y=\displaystyle\lim_{x\rightarrow\infty} \dfrac{x+1}{x^2-2x+m}=0$ nên đồ thị hàm số có một tiệm cận ngang $y=0$.\\
		Do đó để đồ thị hàm số có đúng $2$ đường tiệm cận thì đồ thị hàm số có đúng một tiệm cận đứng, suy ra phương trình $x^2-2x+m=0$ có nghiệm duy nhất khác $-1$ hoặc có hai nghiệm phân biệt $x_1 \ne -1$; $x_2=-1$.
		\begin{center}
			$\Leftrightarrow \hoac{&m_1=1 \\& m_2=-3}\Rightarrow m_1+m_2=-2$.
	\end{center}}
\end{bt}

\begin{bt}%[2D1C4-1]
% nguon: Archive/2024/TaiLieuGiangDay/1.HamSo/1. K12 - 4in1 - Hamso/ThubaiGoc/Bai4-Dang1.tex #2367
%[BG-12NEW-4in1, Nguyen Huynh]
	Cho hàm số $y=\dfrac{2 x+1}{x-3}$ có đồ thị $(C)$. Gọi $M$ là một điểm thuộc đồ thị $(C)$. Tìm điểm $M$ để tổng các khoảng cách từ $M$ đến hai đường tiệm cận của $(C)$ đạt giá trị nhỏ nhất.
	\loigiai{
		Gọi khoảng cách từ $M(x_0;y_0)$ thuộc đồ thị hàm số $y=\dfrac{ax+b}{cx+d}$ đến tiệm cận đứng và tiệm cận ngang lần lượt là $p$ và $q$. Ta chứng minh được 
		$p\cdot q=\dfrac{|ad-bc|}{c^2}$.\\
		Thật vậy, 
		\begin{itemize}
			\item $p=\left|x_0+\dfrac{d}{c}\right|=\left|\dfrac{cx_0+d}{c}\right|$,
			\item $q=\left|y_0-\dfrac{a}{c}\right|=\left|\dfrac{ax_0+b}{cx_0+d}-\dfrac{a}{c}\right|=\left|\dfrac{c(ax_0+b)-a(cx_0+d)}{c(cx_0+d)}\right|=\left|\dfrac{ad-bc}{c(cx_0+d)}\right|$.
		\end{itemize}
		Do đó $pq=\dfrac{|ad-bc|}{c^2}$.\\
		Dễ thấy $I(3 ; 2)$ là giao của hai đường tiệm cận và cũng là tâm đối xứng của đồ thị $(C)$.\\		
		Ta có
		$$pq=\dfrac{|a d-b c|}{c^{2}}=7.$$
		Tổng khoảng cách từ $M$ đến hai đường tiệm cận là $$p+q \geq 2 \sqrt{p q}=2 \cdot \sqrt{7}.$$
		Dấu ``='' xảy ra khi $p=q=\sqrt{7}$. \\
		Suy ra tọa độ điểm $M$ là $M=(3 \pm p ; 2 \pm q)=(3 \pm \sqrt{7} ; 2 \pm \sqrt{7})$.
	}
\end{bt}

\begin{bt}%[2D1H2-1]
% nguon: Archive/2024/TaiLieuGiangDay/1.HamSo/1. K12 - 4in1 - Hamso/ThubaiGoc/Bai2-Dang7.tex #2378
[MĐ2]
	Tìm các điểm cực trị của hàm số $f(x)=x^3+\dfrac{3}{x}-2$.
	\loigiai{
		Tập xác định của hàm số là $\mathbb{R}\setminus \{2\}$. Ta có $f'(x)=3x^2-\dfrac{3}{x^2}=\dfrac{3x^4-3}{x^2}$, suy ra $f'(x)=0 \Leftrightarrow x\pm 1$. Ta có bảng xét dấu của $f'(x)$
		\begin{center}
			\begin{tikzpicture}[font=\footnotesize,>=stealth, scale=1]
				\tkzTabInit[nocadre=false,lgt=1.2,espcl=2.5,deltacl=0.6]
				{$x$ /0.6,$f'$ /0.6}
				{$-\infty$, $-1$, $1$, $+\infty$}
				\tkzTabLine{,+,0,-,0,+,}
			\end{tikzpicture}
		\end{center}
		Vậy hàm số đã cho có hai điểm cực trị là $x_1=-1$ và $x_2=1$.
	}
\end{bt}

\begin{bt}%[2D1H2-1]
% nguon: Archive/2024/TaiLieuGiangDay/1.HamSo/1. K12 - 4in1 - Hamso/ThubaiGoc/Bai2-Dang7.tex #2379
[MĐ2]
	Tìm các điểm cực trị của hàm số $f(x)=3^{x^3-3x+2}$.
	\loigiai{
		Ta có $f'(x)=3\left(x^2-1\right)\cdot 3^{x^3-3x+2}\ln 3$. Khi đó $f'(x)=0 \Leftrightarrow x=\pm 1$. Ta có bảng xét dấu của $f'(x)$
		\begin{center}
			\begin{tikzpicture}[font=\footnotesize,>=stealth, scale=1]
				\tkzTabInit[nocadre=false,lgt=1.2,espcl=2.5,deltacl=0.6]
				{$x$ /0.6,$f'$ /0.6}
				{$-\infty$,$-1$, $1$, $+\infty$}
				\tkzTabLine{,+,0,-,0,+,}
			\end{tikzpicture}
		\end{center} 
		Từ bảng xét dấu, ta có $f(x)$ có hai điểm cực trị là $x_1=-1$ và $x_2=1$.	
	} 
\end{bt}

\begin{bt}%[2D1H2-1]
% nguon: Archive/2024/TaiLieuGiangDay/1.HamSo/1. K12 - 4in1 - Hamso/ThubaiGoc/Bai2-Dang7.tex #2380
[MĐ2]
	Tìm các điểm cực đại của hàm số $f(x)=\cos^3x-3\cos^2x+2$ trên đoạn $\left[0;2\pi\right]$.
	\loigiai{
		Ta có $f'(x)=-3\sin x\cos x\left(\cos x-2\right)$. Suy ra
		\[  f'(x)=0 \Leftrightarrow \sin 2x=0 \Leftrightarrow 2x=k\pi \Leftrightarrow x=\frac{k\pi}{2}, k\in \mathbb{Z}. \]
		 Do $x\in [0;2\pi]$ nên $f'(x)=0 \Leftrightarrow x=\dfrac{k\pi}{2}$, với $k\in \{0;1;2;3;4\}$. Ta có bảng xét dấu
		 \begin{center}
		 	\begin{tikzpicture}[font=\footnotesize,>=stealth, scale=1]
		 		\tkzTabInit[nocadre=false,lgt=1.2,espcl=2.5,deltacl=0.6]
		 		{$x$ /1,$f'$ /0.6}
		 		{$0$, $\dfrac{\pi}{2}$, $\pi$, $\dfrac{3\pi}{2}$, $2\pi$}
		 		\tkzTabLine{0,+,0,-,0,+,0,-,0}
		 	\end{tikzpicture}
		 \end{center}
		Vậy hàm số $f(x)$ có hai điểm cực đại là $x_1=\dfrac{\pi}{2}$ và	 $x_2=\dfrac{3\pi}{2}$ trên đoạn $[0;2\pi]$.
	}
\end{bt}

\begin{bt}%[2D1H2-1]
% nguon: Archive/2024/TaiLieuGiangDay/1.HamSo/1. K12 - 4in1 - Hamso/ThubaiGoc/Bai2-Dang7.tex #2381
[MĐ2]
	Tìm các điểm cực trị của hàm số $f(x)=\dfrac{\ln \left(x^2+2x\right)+1}{2\ln \left(x^2+2x\right)-3}$.
	\loigiai{
		Điều kiện xác định của hàm số là 
		\begin{align*}
			\heva{&x^2+2x>0\\ &\ln\left(x^2+2x\right)\ne \dfrac{3}{2}} \Leftrightarrow \heva{&\hoac{&x>0\\ &x<-2}\\ &x^2+2x\ne \sqrt{\mathrm{e}^3}.}
		\end{align*}
		Với điều kiện đó, ta có $f'(x)=\dfrac{2x+2}{x^2+2x}\cdot \dfrac{-5}{\left(2\ln\left(x^2+2x\right)-3\right)^2}$. Suy ra $f'(x)$ trái dấu với $x+1$. Ta có bảng xét dấu của $f'(x)$
		\begin{center}
			\begin{tikzpicture}[font=\footnotesize,>=stealth, scale=1]
				\tikzset{h style/.style = {pattern=north west lines}}
				\tkzTabInit[nocadre=false,lgt=1.2,espcl=2.5,deltacl=0.6]
				{$x$ /0.6,$f'$ /0.6}
				{$-\infty$, $-2$,  $0$, $+\infty$}
				\tkzTabLine{,+,d,h,d,-,}
			\end{tikzpicture}
		\end{center}
		Vậy $f(x)$ có điểm cực trị nào.
	}
\end{bt}

\begin{bt}%[2D1V2-2]
% nguon: Archive/2024/TaiLieuGiangDay/1.HamSo/1. K12 - 4in1 - Hamso/ThubaiGoc/Bai2-Dang7.tex #2382
[MĐ3]
	\vspace*{-1em}	\immini{
		Cho hàm số $f(x)$ có bảng xét dấu của $f'(x)$ như hình vẽ. Số điểm cực trị của hàm số $y=f(2x-2)$ là bao nhiêu?
		
	}{\vspace*{-1em}
		\begin{tikzpicture}[font=\footnotesize,>=stealth, scale=1]
			\tkzTabInit[nocadre=false,lgt=1.2,espcl=1.5,deltacl=0.6]
			{$x$ /0.6,$f'$ /0.6}
			{$-\infty$, $0$, $2$, $+\infty$}
			\tkzTabLine{,+,0,-,0,+}
		\end{tikzpicture}
	}
	\loigiai{
		Ta có $y'=2f'(2x-2)$. Suy ra
		\[ y'=0 \Leftrightarrow f'(2x-2)=0 \Leftrightarrow \hoac{&2x-2=0\\ &2x-2=2} \Leftrightarrow \hoac{&x=1\\ &x=2.} \]
		Ta có bảng xét dấu
		\begin{center}
			\begin{tikzpicture}[font=\footnotesize,>=stealth, scale=1]
				\tkzTabInit[nocadre=false,lgt=1.2,espcl=2.5,deltacl=0.6]
				{$x$ /0.6,$y'$ /0.6}
				{$-\infty$, $1$, $2$, $+\infty$}
				\tkzTabLine{,+,0,-,0,+,}
			\end{tikzpicture}
		\end{center}
		Vậy $y=f(2x-2)$ có hai điểm cực trị.	 
	}
\end{bt}

\begin{bt}%[2D1V2-2]
% nguon: Archive/2024/TaiLieuGiangDay/1.HamSo/1. K12 - 4in1 - Hamso/ThubaiGoc/Bai2-Dang7.tex #2383
[MĐ3]
	\vspace*{-1em}
	\immini{
		Cho hàm số $f(x)$ có đạo hàm liên tục trên $\mathbb{R}$ và có bảng xét dấu đạo hàm như hình vẽ. Tìm số điểm cực đại của hàm số $g(x)=f(3-x)$.
	}{\vspace*{-0.5em}
		\begin{tikzpicture}[font=\footnotesize,>=stealth, scale=1]
			\tkzTabInit[nocadre=false,lgt=1.2,espcl=1.2,deltacl=0.6]
			{$x$ /0.6,$f'$ /0.6}
			{$-\infty$, $0$, $1$, $2$, $+\infty$}
			\tkzTabLine{,+,0,-,0,+,0,-,}
		\end{tikzpicture}
	}
	\loigiai{
		Ta có $g'(x)=-f'(3-x)$, suy ra
		$$g'(x)=0 \Leftrightarrow f'(3-x)=0 \Leftrightarrow \hoac{&3-x=0\\ &3-x=1\\ &3-x=2} \Leftrightarrow \hoac{&x=3\\ &x=2\\ &x=1.}$$
		Ta có bảng xét dấu
		\begin{center}
			\begin{tikzpicture}[font=\footnotesize,>=stealth, scale=1]
				\tkzTabInit[nocadre=false,lgt=1.2,espcl=2.5,deltacl=0.6]
				{$x$ /0.6,$g'$ /0.6}
				{$-\infty$, $1$, $2$, $3$, $+\infty$}
				\tkzTabLine{,+,0,-,0,+,0,-,}
			\end{tikzpicture}
		\end{center}
		Vậy $f(x)$ có hai điểm cực đại.
	}
\end{bt}

\begin{bt}%[2D1V2-1]
% nguon: Archive/2024/TaiLieuGiangDay/1.HamSo/1. K12 - 4in1 - Hamso/ThubaiGoc/Bai2-Dang7.tex #2384
[MĐ3]
	Cho hàm số $f(x)$ có đạo hàm liên tục trên $\mathbb{R}$ và có bảng biến thiên
	\begin{center}
		\begin{tikzpicture}[font=\footnotesize,>=stealth, scale=1]
			\tkzTabInit[nocadre=false,lgt=1.2,espcl=2.5,deltacl=0.6]
			{$x$ /0.6,$f'$ /0.6,$f$ /2}
			{$-\infty$, $-2$, $0$, $2$, $+\infty$}
			\tkzTabLine{,-,0,+,0,-,0,+}
			\tkzTabVar{+/$+\infty$, -/$1$, +/$5$, -/$3$, +/$+\infty$}
		\end{tikzpicture}
	\end{center}
	Hàm số $g(x)=\ln(f(x))+\dfrac{2}{f(x)}$ có bao nhiêu điểm cực trị?
	\loigiai{
		Tập xác định của $g(x)$ là $\mathbb{R}$. Ta có 
		\[g'(x)=f'(x)\left(\dfrac{1}{f(x)}-\frac{2}{\left[f(x)\right]^2}\right)=\frac{f'(x)\left(f(x)-2\right)}{\left[f(x)\right]^2}.\] 
		Suy ra 
		\[ g'(x)=0 \Leftrightarrow \hoac{&f'(x)=0\\ &f(x)=2} \Leftrightarrow \hoac{&x\in \{-2;0;2\}\\ &x=a<-2\\ &x=b\in (-2;0).} \]
		Từ đó suy ra $g(x)$ có $5$ điểm cực trị.
	}
\end{bt}

\begin{bt}%[2D1H2-6]
% nguon: Archive/2024/TaiLieuGiangDay/1.HamSo/1. K12 - 4in1 - Hamso/ThubaiGoc/Bai2-Dang7.tex #2388
[MĐ2]
	Tìm điểm cực trị của đồ thị hàm số $f(x)=x+\sqrt{9-x^2}$.
	\loigiai{
		Hàm số đã cho có tập xác đinh $[-3;3]$. Ta có $$f'(x)=1-\dfrac{x}{\sqrt{9-x^2}}=\dfrac{\sqrt{9-x^2}-x}{\sqrt{9-x^2}}, \forall -3<x<3.$$
		Suy ra
		\[ f'(x)=0 \Leftrightarrow \sqrt{9-x^2}=x \Leftrightarrow \heva{&x\ge 0 \\ &9-x^2=x^2} \Leftrightarrow x=\frac{3}{\sqrt{2}}.  \]
		Ta có bảng biến thiên  của $f(x)$
		\begin{center}
			\begin{tikzpicture}[font=\footnotesize,>=stealth, scale=1]
				\tkzTabInit[nocadre=false,lgt=1.2,espcl=2.5,deltacl=0.6]
				{$x$ /1,$f'$ /0.6,$f$ /2}
				{$-3$, $\dfrac{3}{\sqrt{2}}$, $3$}
				\tkzTabLine{d,+,0,-,d}
				\tkzTabVar{-/$-3$, +/$3\sqrt{2}$, -/$3$}
			\end{tikzpicture}
		\end{center} 
		Vậy đồ thị hàm số $f(x)$ có một điểm cực trị duy nhất là $\left(\dfrac{3}{\sqrt{2}};3\sqrt{2}\right)$.
	}
\end{bt}

\begin{bt}%[2D1H2-6]
% nguon: Archive/2024/TaiLieuGiangDay/1.HamSo/1. K12 - 4in1 - Hamso/ThubaiGoc/Bai2-Dang7.tex #2389
[MĐ2]
	Tìm các điểm cực trị của hàm số $f(x)=x^2\ln x$.
	\loigiai{
		\immini{
			Hàm số có tập xác định là $(0;+\infty)$. Ta có
			$$f'(x)=2x\ln x+x.$$
			Suy ra
		\[ f'(x)=0 \Leftrightarrow \ln x=-\dfrac{1}{2} \Leftrightarrow x=\frac{1}{\sqrt{\mathrm{e}}}. \]
		Ta có bảng biến thiên như hình bên. Vậy hàm số đã cho có một điểm cực trị $x=\dfrac{1}{\sqrt{\mathrm{e}}}$. 
		}{
			\begin{tikzpicture}[font=\footnotesize,>=stealth, scale=1]
				\tkzTabInit[nocadre=false,lgt=1.2,espcl=2.5,deltacl=0.6]
				{$x$ /1,$f'$ /0.6,$f$ /2}
				{$0$, $\dfrac{1}{\sqrt{\mathrm{e}}}$, $+\infty$}
				\tkzTabLine{d,-,0,+,}
				\tkzTabVar{D+/$0$, -/$\dfrac{-1}{2\mathrm{e}}$, +/$+\infty$}
			\end{tikzpicture}
		}
	}
\end{bt}

\begin{bt}%[2D1H2-6]
% nguon: Archive/2024/TaiLieuGiangDay/1.HamSo/1. K12 - 4in1 - Hamso/ThubaiGoc/Bai2-Dang7.tex #2390
[MĐ2]
	Tìm các điểm cực trị của hàm số $f(x)=\dfrac{x^2+x-1}{x^2-1}$.
	\loigiai{
		Hàm số đã cho có tập xác định $\mathbb{R}\setminus\{-1;1\}$. Ta có $f'(x)=\dfrac{-x^2-4x-1}{\left(x^2-1\right)^2}$, suy ra 
		\[ f'(x)=0\Leftrightarrow x=-2\pm \sqrt{3}. \]
		Ta có bảng xét dấu $f'(x)$
		\begin{center}
			\begin{tikzpicture}[font=\footnotesize,>=stealth, scale=1]
				\tkzTabInit[nocadre=false,lgt=1.2,espcl=2.5,deltacl=0.6]
				{$x$ /0.6,$f'$ /0.6}
				{$-\infty$, $-2-\sqrt{3}$, $-1$, $-2+\sqrt{3}$, $1$, $+\infty$}
				\tkzTabLine{,-,0,+,d,+,0,-,d,-,}
				\tkzTabVar{f}
			\end{tikzpicture}
		\end{center} 
		Từ đó hàm số đã cho có hai điểm cực trị $x=-2-\sqrt{3}$ và $x=-2+\sqrt{3}$.
	}
\end{bt}

\begin{bt}%[2D1H2-6]
% nguon: Archive/2024/TaiLieuGiangDay/1.HamSo/1. K12 - 4in1 - Hamso/ThubaiGoc/Bai2-Dang7.tex #2391
[MĐ2]
	Tìm các điểm cực trị của  hàm số $f(x)=x^{2025}(x-4)^2$.
	\loigiai{
		Ta có
		$$f'(x)=2025x^{2024}(x-4)^2+2x^{2025}(x-4)=x^{2024}(x-4)\left(2027x-4\cdot 2025\right).$$ 
		Do đó $f'(x)=0 \Leftrightarrow x=0$ hoặc $x=4$ hoặc $x=\dfrac{4\cdot 2025}{2027}$. Ta có bảng xét dấu
		\begin{center}
			\begin{tikzpicture}[font=\footnotesize,>=stealth, scale=1]
				\tkzTabInit[nocadre=false,lgt=1.2,espcl=2.5,deltacl=0.6]
				{$x$ /1,$f'$ /0.6}
				{$-\infty$, $0$, $\dfrac{4\cdot 2025}{2027}$, $4$, $+\infty$}
				\tkzTabLine{,+,0,+,0,-,0,+,}
			\end{tikzpicture}
		\end{center} 
		Do đó $f(x)$ có hai điểm cực trị là $\dfrac{4\cdot 2025}{2027}$ và $4$.
	}
\end{bt}

\begin{bt}%[2D1V2-6]
% nguon: Archive/2024/TaiLieuGiangDay/1.HamSo/1. K12 - 4in1 - Hamso/ThubaiGoc/Bai2-Dang7.tex #2392
[MĐ3]
	Tìm giá trị cực trị của hàm số $f(x)=x^x$ với $x>\dfrac{1}{\mathrm{e}^{3}}$.
	\loigiai{
		Ta có $\ln f(x)=x\ln x$. Đạo hàm hai vế, ta được
		\[ \frac{f'(x)}{f(x)}=1+\ln x \Leftrightarrow f'(x)=(1+\ln x)f(x). \]
		Do đó $f'(x)=0 \Leftrightarrow x=\mathrm{e}^{-1}$. Ta có bảng xét dấu
		\begin{center}
			\begin{tikzpicture}[font=\footnotesize,>=stealth, scale=1]
				\tkzTabInit[nocadre=false,lgt=1.2,espcl=2.5,deltacl=0.6]
				{$x$ /0.6,$y'$ /0.6}
				{$\mathrm{e}^{-3}$, $\mathrm{e}^{-1}$, $+\infty$}
				\tkzTabLine{,-,0,+,}
			\end{tikzpicture}
		\end{center} 
		Suy ra $x=\mathrm{e}^{-1}$ là điểm cực trị của $f(x)$ và giá trị cực trị tương ứng là $f\left(\mathrm{e}^{-1}\right)=\mathrm{e}^{\mathrm{e}^{-1}}$.	
	}
\end{bt}

\begin{bt}%[2D1N2-5]
% nguon: Archive/2024/TaiLieuGiangDay/1.HamSo/1. K12 - 4in1 - Hamso/ThubaiGoc/Bai2-Dang5.tex #2402
[THPT Kim Liên - Hà Nội - Lần 1 - Năm 2018, Mức độ N]%[Dự án Bài giảng 12 - Trung Anh]
	Cho hàm số $y=( m+1 ){x^4}-m{x^2}+3$. Tìm tất cả các giá trị thực của tham số $m$ để hàm số có ba điểm cực trị.
	\loigiai{
		Để hàm số có ba điểm cực trị thì $( m+1 )m>0\Leftrightarrow \left[ \begin{aligned}
			&m<-1 \\
			&m>0 .\\
		\end{aligned} \right. $\\ 
		Vậy $m\in ( -\infty ;-1 )\cup ( 0;+\infty )$.}
\end{bt}

\begin{bt}%[2D1N2-5]
% nguon: Archive/2024/TaiLieuGiangDay/1.HamSo/1. K12 - 4in1 - Hamso/ThubaiGoc/Bai2-Dang5.tex #2403
[Thi thử lần 1 - THPT Tiên Du số 1 - Bắc Ninh - 2021, Mức độ N]%[Dự án Bài giảng 12 - Trung Anh]
	Cho hàm số $y=-x^4+(m^2-m)x^2$. Tìm $m$ để hàm số có đúng một cực trị.
	\loigiai{
		Tập xác định $\mathscr{D}=\mathbb{R}$.\\
		Ta có $y'=-4x^3+2(m^2-m)x=2x(-2x^2+m^2-m)$.\\
		$y'=0\Leftrightarrow \hoac{&x=0\\&x^2=\dfrac{m^2-m}{2}.\quad (*)}$\\
		Hàm số có đúng một cực trị khi và chỉ khi phương trình $(*)$ vô nghiệm hoặc có nghiệm kép
		$$\Leftrightarrow \dfrac{m^2-m}{2}\le 0\Leftrightarrow m\in [0;1].$$
	}
\end{bt}

\begin{bt}%[2D1H2-5]
% nguon: Archive/2024/TaiLieuGiangDay/1.HamSo/1. K12 - 4in1 - Hamso/ThubaiGoc/Bai2-Dang5.tex #2404
[Thi thử THPTQG  Chuyên Thái Bình lần 1, năm 2018 - 2019, Mức độ H]%[Dự án Bài giảng 12 - Trung Anh]
	Tìm tất cả giá trị thực của tham số $m$ để đồ thị hàm số $y=x^4-2mx^2+2m-3$ có ba điểm cực trị là ba đỉnh của tam giác cân.
	\loigiai{
		\textbf{Nhận xét:}
		\begin{itemize}
			\item Hàm số $y=ax^4+bx^2+c$ có ba cực trị $\Leftrightarrow ab<0$.
			\item Nếu đồ thị hàm số $y=ax^4+bx^2+c$ có ba điểm cực trị thì các điểm cực trị này luôn là ba đỉnh của một tam giác cân.
		\end{itemize}
		Từ nhận xét trên ta suy ra yêu cầu bài toán thỏa mãn khi và chỉ khi hàm số $y=x^4-2mx^2+2m-3$ có ba cực trị khi và chỉ khi $-2m<0$ hay $m>0$.
	}
\end{bt}

\begin{bt}%[2D1H2-5]
% nguon: Archive/2024/TaiLieuGiangDay/1.HamSo/1. K12 - 4in1 - Hamso/ThubaiGoc/Bai2-Dang5.tex #2405
[Thi thử tốt nghiệp, Trung tâm toán học Pytago - Lào Cai, Mức độ H]%[Dự án Bài giảng 12 - Trung Anh]
	Biết hàm số $y=ax^4+bx^2+c$ có ba điểm cực trị là $A$, $B$, $C$. Có nhận xét gì về hình dạng của $\triangle{ABC}$?
	\loigiai{Ta có $\triangle ABC$ là tam giác cân.	 
	}
\end{bt}

\begin{bt}%[2D1H2-5]
% nguon: Archive/2024/TaiLieuGiangDay/1.HamSo/1. K12 - 4in1 - Hamso/ThubaiGoc/Bai2-Dang5.tex #2406
[Sách tham khảo, Mức độ H]%[Dự án Bài giảng 12 - Trung Anh]
	Có bao nhiêu giá trị nguyên của tham số $ m $ để hàm số $ y=(m+3)x^4 +(2m-13)x^2 +6m-5$ có $ 3 $ điểm cực trị?
	\loigiai{
		Ta có $ y'=4(m+3)x^3+2(2m-13)x=2x\left ( 2(m+3)x^2 +(2m-13)\right ) \,\forall x\in \mathbb{R}$.\\
		Hàm số đã cho có ba điểm cực trị khi và chỉ khi $ (m+3)(2m-13)<0\Leftrightarrow -3<m<\dfrac{13}{2} $.\\
		Vì $ m\in \mathbb{Z} $ nên có $ 9 $ giá trị nguyên của $ m $ thoả mãn.
	}
\end{bt}

\begin{bt}%[2D1H2-5]
% nguon: Archive/2024/TaiLieuGiangDay/1.HamSo/1. K12 - 4in1 - Hamso/ThubaiGoc/Bai2-Dang5.tex #2407
[Thi thử tốt nghiệp - Lương Thế Vinh - Hà Nội - 22, Mức độ H]%[Dự án Bài giảng 12 - Trung Anh]
	Gọi $S$ là tập hợp tất cả các giá trị nguyên âm của tham số $m$ để đồ thị hàm số $y=x^4-2(m+1)x^2+m^2$ có ba điểm cực trị. Số phần tử của tập hợp $S$ là bao nhiêu?
	\loigiai{
		Hàm số $y=x^4-2(m+1)x^2+m^2$ có $3$ điểm cực trị khi và chỉ khi $1\cdot \left[-2(m+1)\right]<0\Leftrightarrow m+1>0\Leftrightarrow m>-1$.\\
		$S$ là tập hợp các giá trị nguyên âm của $m$ nên $S=\varnothing$ hay số phần tử của $S$ là $0$.
	}
\end{bt}

\begin{bt}%[2D1H2-5]
% nguon: Archive/2024/TaiLieuGiangDay/1.HamSo/1. K12 - 4in1 - Hamso/ThubaiGoc/Bai2-Dang5.tex #2408
[Đề thi HK1, Lương Thế Vinh, Hà Nội, 2018-2019, Mức độ H]%[Dự án Bài giảng 12 - Trung Anh]
	Có bao nhiêu số nguyên $m$ để đồ thị của hàm số \ $y=(m-1)x^4+(6-m)x^2+m$ có đúng một cực trị?
\loigiai{
	Xét $m=1$. Khi đó hàm số là $y=5x^2+1$, hàm số này có một cực trị là $1$.\\
	Xét $m\ne 1$. Khi đó $y'=4(m-1)x^3+12(6-m)x, y'=0\Leftrightarrow 4x\left[(m-1)x^2+3(6-m)	\right]=0$.\\
	Hàm số luôn có một điểm cực trị là $x=0$.\\
	Vậy yêu cầu bài toán tương đương phương trình $(m-1)x^2+3(6-m)=0$ có nghiệm kép $x=0$ hoặc vô nghiệm. Điều này tương đương
	$$3(m-1)(6-m)\ge 0\Leftrightarrow 1<m\le 6.$$
	Kết hợp hai trường hợp ta có $1\le m\le 6$. Mà $m\in\mathbb{Z}$ nên $m\in\{1;2;3;4;5;6\}$.
}
\end{bt}

\begin{bt}%[2D1H2-5]
% nguon: Archive/2024/TaiLieuGiangDay/1.HamSo/1. K12 - 4in1 - Hamso/ThubaiGoc/Bai2-Dang5.tex #2412
[Chuyên Bắc Ninh, Mức H]%[Dự án bài giảng 12, Trung Anh]
	Tìm tất cả các giá trị thực của tham số $m$ sao cho đồ thị của hàm số $y={x^4}-2( m+1 ){x^2}+{{m}^2}$ có ba điểm cực trị tạo thành một tam giác vuông cân.
	\loigiai{
			Ta có $y'=4x( {x^2}-m-1 )$\\
			Xét $y'=0\Leftrightarrow \left[ \begin{aligned}
				&x=0 \\
				&{x^2}=m+1 \\
			\end{aligned} \right. $.\\ 
			Để đồ thị số có ba điểm cực trị thì $m>-1 ( * )$\\
			Tọa độ ba điểm cực trị là $A( 0;{{m}^2} )$, $B( \sqrt{m+1};-2m-1 )$, $C( -\sqrt{m+1};-2m-1 )$\\
			Gọi $H$ là trung điểm của đoạn thẳng $BC$ thì $H( 0;-2m-1 )$\\
			Khi đó ba điểm cực trị lập thành tam giác vuông cân khi $$AH=BH \Leftrightarrow \sqrt{{{( m+1 )}^4}}=\sqrt{m+1} \Leftrightarrow m=0\, (\text{thỏa mãn}).$$
		}
\end{bt}

\begin{bt}%[2D1H2-5]
% nguon: Archive/2024/TaiLieuGiangDay/1.HamSo/1. K12 - 4in1 - Hamso/ThubaiGoc/Bai2-Dang5.tex #2413
[NXBGD - Đoàn Quỳnh - Đề số 03 - Năm 2017 - 2018, Mức H]%[Dự án bài giảng 12, Trung Anh]
	Cho hàm số phụ thuộc tham số $m$ là $y={x^4}-8{x^2}+m$; kí hiệu đồ thị của nó là $( {{C}_m} )$. Tìm giá trị của tham số $m$ để $( {{C}_m} )$ có các điểm cực tiểu cách gốc tọa độ $O$ một khoảng bằng $2$.
	\loigiai{
		Ta có
		$y'=0\Leftrightarrow 4{x^3}-16x=0\Leftrightarrow \left[ \begin{aligned}
			&x=0\Rightarrow y=m \\
			&x=-2\Rightarrow y=m-16 \\
			&x=2\Rightarrow y=m-16 \\
		\end{aligned} \right. $\\
		Dựa vào hình dáng đồ thị ta có điểm cực tiểu có tọa độ $A( 2;m-16 )$.\\
		Khi đó $OA=2\Leftrightarrow {{2}^2}+{{( m-16 )}^2}=4\Leftrightarrow m-16=0\Leftrightarrow m=16$.
		
	}
\end{bt}

\begin{bt}%[2D1H2-5]
% nguon: Archive/2024/TaiLieuGiangDay/1.HamSo/1. K12 - 4in1 - Hamso/ThubaiGoc/Bai2-Dang5.tex #2414
[HSG THPT 2018-2019, Mức H]%[Dự án bài giảng 12, Trung Anh]
	Tìm các giá trị của $m$ để hàm số $y=mx^4+(m^2+m)x^2-2019$ có đúng một điểm cực trị.
	\loigiai{
		\begin{itemize}
			\item \textbf{Trường hợp 1.} Với $m=0$ ta có hàm số $y=-2019$, hàm số không có điểm cực trị nên $m=0$ không thỏa mãn.
			\item \textbf{Trường hợp 2.} Với $m\neq 0$ hàm số đã cho là hàm trùng phương. Khi đó hàm số đã cho có đúng một điểm cực trị khi và chỉ khi 
			$$m(m^2+m)\geq 0
			\Leftrightarrow {m^2}(m+1)\geq 0
			\Leftrightarrow m\geq -1.$$
			Đối chiếu điều kiện $m\neq 0$ ta được $m\in \left[-1;0\right)\cup (0;+\infty)$.
		\end{itemize}
		Vậy $m\in \left[-1;0\right)\cup (0;+\infty)$ là các giá trị cần tìm.
	}
\end{bt}

\begin{bt}%[2D1H2-5]
% nguon: Archive/2024/TaiLieuGiangDay/1.HamSo/1. K12 - 4in1 - Hamso/ThubaiGoc/Bai2-Dang5.tex #2415
[Thi thử - THPT Chuyên Hà Tĩnh - 2021, Mức H]%[Dự án bài giảng 12, Trung Anh]
	Cho đồ thị hàm số $y=ax^4+bx^2+c$ có điểm cực đại $A(0;-3)$ và điểm cực tiểu $B(-1;-5)$. Tính giá trị của $P=a+2b+3c$.
	\loigiai{
		Ta có $y'=4ax^3+2bx$, $a>0$.\\
		Vì đồ thị hàm số có điểm cực đại $A(0;-3)$ và điểm cực tiểu $B(-1;-5)$ nên ta có hệ phương trình
		\[\heva{&y(0)=-3\\&y(-1)=-5\\&y'(-1)=0}\Leftrightarrow \heva{&c=-3\\&a+b-3=-5\\&-4a-2b=0}\Leftrightarrow \heva{&c=-3\\&a=2\\&b=-4.}\]
		Suy ra $P=a+2b+3c=2+2\cdot(-4)+3\cdot(-3)=-15$.
	}
\end{bt}

\begin{bt}%[2D1V2-5]
% nguon: Archive/2024/TaiLieuGiangDay/1.HamSo/1. K12 - 4in1 - Hamso/ThubaiGoc/Bai2-Dang5.tex #2416
[Sách tham khảo, Mức V]%[Dự án bài giảng 12, Trung Anh]
	Tìm các giá trị của tham số $m$ để đồ thị hàm số $y=x^4-2mx^2+\dfrac{m}{2}$ có ba điểm cực trị và đường tròn đi qua ba điểm cực trị này có bán kính bằng 1.
	\loigiai{
		Ta có $y'=4x^3-4mx=0\Leftrightarrow \left[ \begin{aligned}
			&x=0 \\
			&x^2=m \\
		\end{aligned} \right. $.\\ 
		Hàm số có 3 điểm cực trị khi và chỉ khi $y=0$ có 3 nghiệm phân biệt và $y'$ đổi dấu khi đi qua các nghiệm đó $\Leftrightarrow m>0$.\\
		Khi đó $A\left(0;\dfrac{m}{2} \right)$, $B\left(\sqrt{m};\dfrac{m}{2}-m^2 \right)$, $C\left(-\sqrt{m};\dfrac{m}{2}-m^2 \right)$ \\
		Suy ra $AB=AC=\sqrt{m+m^4}$, $BC=2\sqrt{m}$ \\
		Và ${{S}_{\Delta ABC}}=\dfrac{1}{2}d[ A;(BC) ]\cdot BC=\dfrac{1}{2}| {{y}_B}-{{y}_A} |\cdot | {{x}_C}-{{x}_B} |=m^2\sqrt{m}$.\\
		Ta có
		\begin{align*}
			&R=\dfrac{AB\cdot BC\cdot AC}{4{{S}_{\Delta ABC}}}=1\\
			\Leftrightarrow&\, (m+m^4)\cdot 2\sqrt{m}=4m^2\sqrt{m} \Leftrightarrow m^3-2m+1=0\\
			 \Leftrightarrow&\, \left[ \begin{aligned}
				&m=1(n) \\
				&m=\dfrac{-1+\sqrt{5}}{2}(n) \\
				&m=\dfrac{-1-\sqrt{5}}{2}(l) \\
			\end{aligned} \right.
		\end{align*}.
		
	}
\end{bt}

\begin{bt}%[2D1N4-1]
% nguon: Archive/2024/TaiLieuGiangDay/1.HamSo/1. K12 - 4in1 - Hamso/ThubaiGoc/Bai4-Dang2.tex #2427
%[BG-12NEW-4in1, Nguyen Huynh]
	Tìm các đường tiệm cận của đồ thị hàm số $y=2^{x}$.
	\loigiai{
		Đồ thị của hàm số $y=2^x$.
		\begin{center}
			\begin{tikzpicture}[line join = round, line cap = round,>=stealth,font=\footnotesize,scale=1]
				\def\xt{-2} \def\xp{2} \def\yd{-1} \def\yt{4.5}
				\def\a{3.5} 
				\draw[->] (\xt,0)--(\xp,0) node[above]{$x$};
				\draw[->] (0,\yd)--(0,\yt) node[left]{$y$};
				%\draw[dashed,thin,black] (-1,0)--(-1,3)--(0,3);
				%\draw[dashed,thin,black] (1,0)--(1,-1)--(0,-1);
				\fill (0,0) circle (1.5pt) node[below left]{$O$} (0,1) circle (1.5pt) node[right]{$1$};
				\begin{scope}
					\clip (\xt+0.1,\yd-0.1) rectangle (\xp-0.3,\yt-0.1);
					\draw[samples=150,thick,smooth,domain=\xt:\xp] plot(\x,{\a^\x});
				\end{scope}
			\end{tikzpicture}
		\end{center}
		Căn cứ vào đồ thị của hàm số $y=2^{x}$, suy ra đồ thị hàm số đã cho có một tiệm cận duy nhất là tiệm cận ngang $y=0$.
	}	
\end{bt}

\begin{bt}%[2D1N4-1]
% nguon: Archive/2024/TaiLieuGiangDay/1.HamSo/1. K12 - 4in1 - Hamso/ThubaiGoc/Bai4-Dang2.tex #2428
%[BG-12NEW-4in1, Nguyen Huynh]
	Tìm các đường tiệm cận của đồ thị hàm số $y=\left( \dfrac{1}{2}\right) ^{x}$.
	\loigiai{Đồ thị đồ thị hàm số $y=\left( \dfrac{1}{2}\right) ^{x}$.
		\begin{center}
			\begin{tikzpicture}[line join = round, line cap = round,>=stealth,font=\footnotesize,scale=1]
				\def\xt{-2} \def\xp{2.5} \def\yd{-1} \def\yt{4.5}
				\def\a{0.35} 
				\draw[->] (\xt,0)--(\xp,0) node[above]{$x$};
				\draw[->] (0,\yd)--(0,\yt) node[left]{$y$};
				%\draw[dashed,thin,black] (-1,0)--(-1,3)--(0,3);
				%\draw[dashed,thin,black] (1,0)--(1,-1)--(0,-1);
				\fill (0,0) circle (1.5pt) node[below left]{$O$} (0,1) circle (1.5pt) node[right]{$1$};
				\begin{scope}
					\clip (\xt+0.1,\yd-0.1) rectangle (\xp-0.3,\yt-0.1);
					\draw[samples=150,thick,smooth,domain=\xt:\xp] plot(\x,{\a^\x});
				\end{scope}
			\end{tikzpicture}
		\end{center}
		Căn cứ vào đồ thị của hàm số $y=\left( \dfrac{1}{2}\right) ^{x}$, suy ra đồ thị hàm số đã cho có một tiệm cận duy nhất là tiệm cận ngang $y=0$.
	}	
\end{bt}

\begin{bt}%[2D1H4-1]
% nguon: Archive/2024/TaiLieuGiangDay/1.HamSo/1. K12 - 4in1 - Hamso/ThubaiGoc/Bai4-Dang2.tex #2431
%[BG-12NEW-4in1, Nguyen Huynh]
	Tìm các đường tiệm cận của đồ thị hàm số $y=2^{x}-1$.
	\loigiai{
		Tịnh tiến đồ thị $y=2^{x}$ xuống dưới 1 đơn vị, ta được đồ thị $y=2^{x}-1$
		\begin{center}
			\begin{tikzpicture}[line join = round, line cap = round,>=stealth,font=\footnotesize,scale=1]
				\def\xt{-2} \def\xp{2} \def\yd{-2} \def\yt{3}
				\def\a{3.5} 
				\draw[->] (\xt,0)--(\xp,0) node[above]{$x$};
				\draw[->] (0,\yd)--(0,\yt) node[left]{$y$};
				\draw[dashed,thick,black] (-2,-1)--(2,-1);
				%\draw[dashed,thin,black] (1,0)--(1,-1)--(0,-1);
				\fill (0,0) circle (1.5pt) node[below right]{$O$} (0.8,-0.8) node[right]{$y=-1$};
				\begin{scope}
					\clip (\xt+0.1,\yd-0.1) rectangle (\xp-0.3,\yt-0.1);
					\draw[samples=150,thick,smooth,domain=\xt:\xp] plot(\x,{\a^\x - 1});
				\end{scope}
			\end{tikzpicture}
		\end{center}
		Căn cứ vào đồ thị, suy ra đồ thị hàm số đã cho có một tiệm cận duy nhất là tiệm cận ngang $y=-1$.
	}	
\end{bt}

\begin{bt}%[2D1H4-1]
% nguon: Archive/2024/TaiLieuGiangDay/1.HamSo/1. K12 - 4in1 - Hamso/ThubaiGoc/Bai4-Dang2.tex #2433
%[BG-12NEW-4in1, Nguyen Huynh]
	Tìm các đường tiệm cận của đồ thị hàm số $y = \log_{2}x + 2^{x}.$
	\loigiai{
		Tập xác định $\mathscr D = \left( 0; +\infty \right)$.\\
		Ta có $\lim\limits_{x \to +\infty} y = +\infty$ và $\lim\limits_{x \to 0^{+}} y = -\infty$.\\
		Suy ra đồ thị hàm số đã cho không có tiệm cận ngang và có duy nhất 1 tiệm cận đứng là $x=0$.
	}	
\end{bt}

\begin{bt}%[2D1N4-1]
% nguon: Archive/2024/TaiLieuGiangDay/1.HamSo/1. K12 - 4in1 - Hamso/ThubaiGoc/Bai4-Dang2.tex #2437
%[BG-12NEW-4in1, Nguyen Huynh]
	Tìm các đường tiệm cận của đồ thị hàm số $y=\left( \dfrac{1}{5}\right) ^{x}$.
	\loigiai{Đồ thị đồ thị hàm số $y=\left( \dfrac{1}{5}\right) ^{x}$.
		\begin{center}
			\begin{tikzpicture}[line join = round, line cap = round,>=stealth,font=\footnotesize,scale=1]
				\def\xt{-2} \def\xp{2.5} \def\yd{-1} \def\yt{4.5}
				\def\a{0.35} 
				\draw[->] (\xt,0)--(\xp,0) node[above]{$x$};
				\draw[->] (0,\yd)--(0,\yt) node[left]{$y$};
				%\draw[dashed,thin,black] (-1,0)--(-1,3)--(0,3);
				%\draw[dashed,thin,black] (1,0)--(1,-1)--(0,-1);
				\fill (0,0) circle (1.5pt) node[below left]{$O$} (0,1) circle (1.5pt) node[right]{$1$};
				\begin{scope}
					\clip (\xt+0.1,\yd-0.1) rectangle (\xp-0.3,\yt-0.1);
					\draw[samples=150,thick,smooth,domain=\xt:\xp] plot(\x,{\a^\x});
				\end{scope}
			\end{tikzpicture}
		\end{center}
		Căn cứ vào đồ thị của hàm số $y=\left( \dfrac{1}{5}\right) ^{x}$, suy ra đồ thị hàm số đã cho có một tiệm cận duy nhất là tiệm cận ngang $y=0$.
	}	
\end{bt}

\begin{bt}%[2D1H4-1]
% nguon: Archive/2024/TaiLieuGiangDay/1.HamSo/1. K12 - 4in1 - Hamso/ThubaiGoc/Bai4-Dang2.tex #2438
%[BG-12NEW-4in1, Nguyen Huynh]
	Tìm các đường tiệm cận của đồ thị hàm số $y=\log_{3}(x-2)$.
	\loigiai{
		Tịnh tiến đồ thị $y=\log_{3}x$ qua phải $2$ đơn vị, ta được đồ thị $y=\log_{3}(x-2)$
		\begin{center}
			\begin{tikzpicture}
				\draw[->] (-1,0)--(0,0) node[below left]{O}--(5.2,0) node[below]{$x$};
				\draw[->] (0,-3.3)--(0,2) node[below left]{$y$};
				\draw[smooth, black, line width=1,samples=100] plot[domain=2.04:4.5] (\x,{{ln(\x - 2)/ln(3)}}) ;
				\draw[dashed,thick,black] (2,2)--(2,-3);
				\draw[dashed] (2,0) circle (1.5pt) node[below right]{$2$} (3,0) circle (1.5pt) node[below right]{$3$};
			\end{tikzpicture}
		\end{center}
		Căn cứ vào hình vẽ, suy ra đồ thị hàm số đã cho có một tiệm cận duy nhất là tiệm cận đứng $x=2$.
	}	
\end{bt}

\begin{bt}%[2D1H4-1]
% nguon: Archive/2024/TaiLieuGiangDay/1.HamSo/1. K12 - 4in1 - Hamso/ThubaiGoc/Bai4-Dang2.tex #2439
%[BG-12NEW-4in1, Nguyen Huynh]
	Tìm các đường tiệm cận của đồ thị hàm số $y=4^{x}+2$.
	\loigiai{
		Tịnh tiến đồ thị $y=4^{x}$ lên trên $2$ đơn vị, ta được đồ thị $y=4^{x}+2$
		\begin{center}
			\begin{tikzpicture}[line join = round, line cap = round,>=stealth,font=\footnotesize,scale=1]
				\def\xt{-2} \def\xp{2} \def\yd{-1} \def\yt{4}
				\def\a{3.5} 
				\draw[->] (\xt,0)--(\xp,0) node[above]{$x$};
				\draw[->] (0,\yd)--(0,\yt) node[left]{$y$};
				\draw[dashed,thick,black] (-2,2)--(2,2);
				%\draw[dashed,thin,black] (1,0)--(1,-1)--(0,-1);
				\fill (0,0) circle (1.5pt) node[below right]{$O$} (0.8,1.8) node[right]{$y=2$};
				\begin{scope}
					\clip (\xt+0.1,\yd-0.1) rectangle (\xp-0.3,\yt-0.1);
					\draw[samples=150,thick,smooth,domain=\xt:\xp] plot(\x,{\a^\x + 2});
				\end{scope}
			\end{tikzpicture}
		\end{center}
		Căn cứ vào đồ thị, suy ra đồ thị hàm số đã cho có một tiệm cận duy nhất là tiệm cận ngang $y=2$.
	}	
\end{bt}

\begin{bt}%[2D1C4-1]
% nguon: Archive/2024/TaiLieuGiangDay/1.HamSo/1. K12 - 4in1 - Hamso/ThubaiGoc/Bai4-Dang2.tex #2440
%[BG-12NEW-4in1, Nguyen Huynh]
	Cho đồ thị hàm số $ y=f(x) $ như hình vẽ. Đồ thị của hàm số $ y=(0{,}7)^{f(x)} $ có tất cả bao nhiêu đường tiệm cận ?
	\begin{center}
		\begin{tikzpicture}[>=latex,line join=round,line cap=round,line width=1pt,font=\footnotesize,scale=0.85]
			\draw[->](-4,0)--(0,0)node [above right=-2pt]{$ O $}--(4,0) node [above]{$ x $};
			\draw[->](0,-2)--(0,5)node [right]{$ y $};
			\draw [dashed] (-4,2)--(4,2);
			\draw [dashed] (-1,-2)--(-1,5);
			\draw [dashed] (1,-2)--(1,5);
			\coordinate (A) at (-1.1,-2);
			\coordinate (B) at (-1.5,1);
			\coordinate (C) at (-3,3);
			\coordinate (D) at (-3.9,-2);
			\coordinate (E) at (-0.9,5);	
			\coordinate (F) at (0.9,-2);
			\coordinate (G) at (1.1,-2);
			\coordinate (H) at (3.9,1.9);
			\foreach \diem  in {(-1,0),(1,0),(0,2)}	{\fill \diem circle (1.5pt);}
			\draw (A) .. controls +(95:0.5) and +(-30:0.5) .. (B) .. controls +(150:0.5) and + (0:0.5)..(C)..controls +(180:0.5) and +(80:0.5).. (D); 
			\draw (E).. controls + (-90:3) and +(90:3)..(F);
			\draw (G).. controls + (85:2) and +(-180:3)..(H);
			\node [below right] at (-1,0){$ a $};
			\node [above left] at (1,0){$ b $};
			\node [above right] at (0,2){$ c $};
		\end{tikzpicture}
	\end{center}
	\loigiai{
		Tập xác định của hàm số $ y=0{,}7^{f(x)} $ là $ \mathbb{R} \setminus \{a ; b\} $.\\
		Xét tiệm cận đứng \\
		$ \lim \limits_{x \to a^-} 0{,}7^{f(x)} =\lim \limits_{x \to -\infty} 0{,}7^{x} =+\infty$ $ \Rightarrow $ $ x=a $ là một tiệm cận đứng.\\
		$ \lim \limits_{x \to b^-} 0{,}7^{f(x)} = \lim \limits_{x \to -\infty} 0{,}7^{x}=+\infty$ $ \Rightarrow $ $ x=b $ là một tiệm cận đứng.\\
		Xét tiệm cận ngang \\
		$ \lim \limits_{x \to -\infty} 0{,}7^{f(x)} = \lim \limits_{x \to -\infty} 0{,}7^{x}=+\infty$.\\
		$ \lim \limits_{x \to +\infty} 0{,}7^{f(x)} = 0{,}7^{c}$ $ \Rightarrow $ $ y=0{,}7^c $ là tiệm cận ngang.\\
		Vậy đồ thị hàm số có 3 đường tiệm cận.
	}
\end{bt}

\begin{bt}%[2D1C4-1]
% nguon: Archive/2024/TaiLieuGiangDay/1.HamSo/1. K12 - 4in1 - Hamso/ThubaiGoc/Bai4-Dang2.tex #2441
%[BG-12NEW-4in1, Nguyen Huynh]
	\immini{
		Cho đồ thị hàm số $y=f(x)$ như hình vẽ. Đồ thị của hàm số $y=\log f(x)$ có tất cả bao nhiêu đường tiệm cận đứng?
	}{
		\begin{tikzpicture}[scale=.7, font=\footnotesize, line join=round, line cap=round, >=stealth]
			\draw[->] (-2.5,0)--(4,0)node[above]{$x$};
			\draw[->] (0,-2)--(0,4)node[left]{$y$};
			\draw (-2,-2)--(-1.5,0)node[above left]{$a$}circle(1pt);
			\draw (-1.5,0) .. controls (-.5,4) and (-.2,2.5) .. (.5,.5);
			\draw (.5,.5) .. controls (.7,-.4) and (1.2,-.4) .. (1.5,3);
			\draw (1.5,3)..controls (1.7,3.5) and (2,3.5) .. (2.6,0)--(2.9,-1.5);
			\draw[fill=black] (0,0)node[below left]{$O$}circle(1pt) (0.75,0)node[below]{$b$}circle(1pt) (2.6,0)node[above right]{$c$}circle(1pt);
		\end{tikzpicture}
	}
	\loigiai{
		Dựa vào đồ thị hàm số ta có tập xác định của hàm số là $\mathscr{D}=(a;c)\setminus\{b\}$.\\
		Ta có $\lim\limits_{x\to a^+}\log f(x)=\lim\limits_{x\to 0^+}\log x=-\infty$, suy ra đồ thị có tiệm cận đứng là $x=a$.\\
		Ta có $\lim\limits_{x\to b^+}\log f(x)=\lim\limits_{x\to 0^+}\log x=-\infty$, suy ra đồ thị có tiệm cận đứng là $x=b$.\\
		Ta có $\lim\limits_{x\to c^-}\log f(x)=\lim\limits_{x\to 0^+}\log x=-\infty$, suy ra đồ thị có tiệm cận đứng là $x=c$.\\
		Vậy đồ thị hàm số $y=\log f(x)$ có tất cả ba đường tiệm cận đứng.
	}
\end{bt}

\begin{bt}%[2D1H4-1]
% nguon: Archive/2024/TaiLieuGiangDay/1.HamSo/1. K12 - 4in1 - Hamso/ThubaiGoc/Bai4-Dang3.tex #2452
%[CKP]%[Dự án TL12New-4in1-NCT]
	\immini{
		Cho hàm số $y=f(x)=\dfrac{2-4x}{4-x^2}$ có đồ thị là đường cong như bên. Các đường thẳng $x=2$ và $x=-2$ có phải là các đường tiệm cận đứng của đồ thị hàm số đã cho hay không? Vì sao?
	}{
		\begin{tikzpicture}[>=latex,line join=round, line cap=round, scale=0.7]
			\draw[->] (-6,0)--(6,0) node[below]{$x$};
			\draw[->] (0,-7)--(0,7) node[right]{$y$};
			\draw[red,dashed] (-2,-7)--(-2,7) (2,-7)--(2,7);
			\foreach \x in {-5,-4,-3,-2,-1,1,2,3,4,5} 
			\draw[fill=black] (\x,0) circle (1pt) node[below]{$\x$};
			\foreach \y in {-6,-5,-4,-3,-2,-1,1,2,3,4,5,6} 
			\draw[fill=black] (0,\y) circle (1pt) node[left]{$\y$};
			\clip (-6,-7) rectangle (6,7);
			\draw[samples=200,smooth,blue,line width=1] plot[domain=-6:-2.1] (\x,{(2-4*(\x))/(4-(\x)^2)});
			\draw[samples=200,smooth,blue,line width=1] plot[domain=-1.9:1.9] (\x,{(2-4*(\x))/(4-(\x)^2)});
			\draw[samples=200,smooth,blue,line width=1] plot[domain=2.1:6] (\x,{(2-4*(\x))/(4-(\x)^2)});
			\draw (0,0) node[below left]{$O$};
		\end{tikzpicture}
	}
	\loigiai{
		Ta có $\lim\limits_{x\to 2^{+}} f(x)=\lim\limits_{x\to 2^{+}}\dfrac{2-4x}{4-x^2}=+\infty$ hoặc $\lim\limits_{x\to 2^{-}} f(x)=\lim\limits_{x\to 2^{-}}\dfrac{2-4x}{4-x^2}=-\infty$ nên đường thẳng $x=2$ là đường tiệm cận đứng của đồ thị hàm số $y=f(x)$.\\
		Tương tự, $\lim\limits_{x\to -2^{+}} f(x)=\lim\limits_{x\to -2^{+}}\dfrac{2-4x}{4-x^2}=+\infty$ hoặc $\lim\limits_{x\to -2^{-}}f(x)=\lim\limits_{x\to -2^{-}}\dfrac{2-4x}{4-x^2}=-\infty$ nên đường thẳng $x=-2$ là đường tiệm cận đứng của đồ thị hàm số $y=f(x)$ đã cho.
	}
\end{bt}

\begin{bt}%[2D1H4-1]
% nguon: Archive/2024/TaiLieuGiangDay/1.HamSo/1. K12 - 4in1 - Hamso/ThubaiGoc/Bai4-Dang3.tex #2453
%CKP%[Dự án TL12New-4in1-NCT]
	\immini{
		Cho hàm số $y=f(x)=\dfrac{\sqrt{x^2+1}}{x+1}$ có đồ thị là đường cong như hình bên. Hãy xác định các đường tiệm cận đứng và đường tiệm cận ngang của đồ thị hàm số đã cho.
	}{
		\begin{tikzpicture}[>=latex,line join=round, line cap=round, scale=0.7]
			\draw[->] (-7,0)--(6,0) node[below]{$x$};
			\draw[->] (0,-6)--(0,6) node[right]{$y$};
			\draw[red,dashed] (-1,-6)--(-1,6) (-7,-1)--(6,-1) (-7,1)--(6,1);
			\foreach \x in {-6,-5,-4,-3,-2,-1,1,2,3,4,5} 
			\draw[fill=black] (\x,0) circle (1pt) node[below]{$\x$};
			\foreach \y in {-5,-4,-3,-2,-1,1,2,3,4,5} 
			\draw[fill=black] (0,\y) circle (1pt) node[left]{$\y$};
			\clip (-7,-6) rectangle (6,6);
			\draw[samples=200,smooth,blue,line width=1] plot[domain=-7:-1.1] (\x,{sqrt((\x)^2+1)/(\x+1)});
			\draw[samples=200,smooth,blue,line width=1] plot[domain=-0.9:6] (\x,{sqrt((\x)^2+1)/(\x+1)});
			\draw (0,0) node[above left]{$O$};
		\end{tikzpicture}
	}
	\loigiai{
		Điều kiện xác định: $\mathscr{D}=\mathbb{R}\setminus\{-1\}$.\\
		Ta có $\lim\limits_{x\to -1^{+}}f(x)=\lim\limits_{x\to -1^{+}}\dfrac{\sqrt{x^2+1}}{x+1}=+\infty$ hoặc $\lim\limits_{x\to -1^{-}} f(x)=\lim\limits_{x\to -1^{-}}\dfrac{\sqrt{x^2+1}}{x+1}=-\infty$ nên đường thẳng $x=-1$ là tiệm cận đứng của đồ thị hàm số $y=f(x)$ đã cho.
	}
\end{bt}

\begin{bt}%[2D1H4-1]
% nguon: Archive/2024/TaiLieuGiangDay/1.HamSo/1. K12 - 4in1 - Hamso/ThubaiGoc/Bai4-Dang3.tex #2454
%CKP%[Dự án TL12New-4in1-NCT]
	\immini{
		Đường cong ở hình bên là đồ thị $(C)$ của hàm số $y=f(x)=\dfrac{x^2-5x+7}{x-3}$. Chứng minh rằng đường thẳng $\Delta\colon y=x-2$ là đường tiệm cận xiên của đồ thị hàm số $y=f(x)$. Vẽ $\Delta$ lên cùng hệ trục toạ độ với $(C)$.
	}{
		\begin{tikzpicture}[>=latex,line join=round, line cap=round, scale=0.8]
			\draw[->] (-2.5,0)--(7.5,0) node[below]{$x$};
			\draw[->] (0,-4)--(0,6) node[left]{$y$};
			\draw[red] (3,-4)--(3,6);
			\foreach \x in {-2,-1,1,2,3,4,5,6,7} 
			\draw[fill=black] (\x,0) circle (1pt) node[below]{$\x$};
			\foreach \y in {-3,-2,-1,1,2,3,4,5} 
			\draw[fill=black] (0,\y) circle (1pt) node[left]{$\y$};
			\clip (-2.5,-4) rectangle (7.5,6);
			\draw[samples=200,smooth,blue,line width=1] plot[domain=-2.5:2.9] (\x,{((\x)^2-5*(\x)+7)/(\x-3)});
			\draw[samples=200,smooth,blue,line width=1] plot[domain=3.1:7.5] (\x,{((\x)^2-5*(\x)+7)/(\x-3)});
		\end{tikzpicture}
	}
	\loigiai{
		\begin{center}
			\begin{tikzpicture}[>=latex,line join=round, line cap=round, scale=0.8]
				\draw[->] (-2.5,0)--(7.5,0) node[below]{$x$};
				\draw[->] (0,-4)--(0,6) node[left]{$y$};
				\draw[red] (3,-4)--(3,6);
				\foreach \x in {-2,-1,1,2,3,4,5,6,7} 
				\draw[fill=black] (\x,0) circle (1pt) node[below, fill=white]{$\x$};
				\foreach \y in {-3,-2,-1,1,2,3,4,5} 
				\draw[fill=black] (0,\y) circle (1pt) node[left]{$\y$};
				\clip (-2.5,-4) rectangle (7.5,6);
				\draw[samples=200,smooth,blue,line width=1] plot[domain=-2.5:2.9] (\x,{((\x)^2-5*(\x)+7)/(\x-3)});
				\draw[samples=200,smooth,blue,line width=1] plot[domain=3.1:7.5] (\x,{((\x)^2-5*(\x)+7)/(\x-3)});
				\draw[samples=200,smooth,red] plot[domain=-2.5:7.5] (\x,{\x-2});
				\draw[red] (6,3.5) node[rotate=45]{$y=x-2$};
			\end{tikzpicture}
		\end{center}
		Ta có $\lim\limits_{x\to +\infty}\left[f(x)-(x-2)\right]=\lim\limits_{x\to+\infty}\left[\dfrac{x^2-5x+7}{x-3}-(x-2)\right]=\lim\limits_{x\to+\infty}\left(\dfrac{1}{x-3}\right)=0$.\\
		Suy ra đường thẳng $y=x-2$ là tiệm cận xiên của đồ thị hàm số $y=f(x)$.\\
		Hình ảnh của đường thẳng $\Delta\colon y=x-2$ như hình vẽ. 
	}
\end{bt}

\begin{bt}%[2D1H4-1]
% nguon: Archive/2024/TaiLieuGiangDay/1.HamSo/1. K12 - 4in1 - Hamso/ThubaiGoc/Bai4-Dang3.tex #2455
%CD%[Dự án TL12New-4in1-NCT]
	Đồ thị hàm số ở \texttt{Hình 18a, Hình 18b} đều có đường tiệm cận ngang là đường thẳng màu đỏ. Hỏi đó là đồ thị của hàm số nào trong các hàm số sau?
	
	\begin{tikzpicture}[scale=.7, font=\footnotesize, line join=round, line cap=round, >=stealth]
		\draw[->] (-4,0)--(0,0) node[above right]{$O$}--(5,0) node[below]{$x$};
		\draw[->] (0,-2) --(0,3) node[right]{$y$};
		\draw [domain=-4:5, samples=100] %
		plot (\x, {(2*(\x)^2-2)/((\x)^2+2)});
		\foreach \x/\g in {-1/-100,1/-90} 
		\draw[fill] (\x,0) circle(.5pt)node [shift={(\g:.3)}] {$\x$};
		\foreach \y/\g in {-1/-140,2/135} 
		\draw[fill] (0,\y) circle(.5pt)node [shift={(\g:.3)}] {$\y$};
		%\draw[blue] (1,-4)--(1,5);
		\draw[red!80] (-4,2)--(5,2);
		\draw (2,-2) node{Hình 18a};
	\end{tikzpicture}\hspace{2cm}
	\begin{tikzpicture}[scale=.7, font=\footnotesize, line join=round, line cap=round, >=stealth]
		\draw[->] (-4,0)--(0,0) node[above left]{$O$}--(5,0) node[below]{$x$};
		\draw[->] (0,-2) --(0,2) node[right]{$y$};
		\draw [domain=-4:5, samples=100] %
		plot (\x, {(2*(\x)^2+2*(\x)-1)/(2*(\x)^2+1)});
		\foreach \x/\g in {-2/-90,1/-90} 
		\draw[fill] (\x,0) circle(.5pt)node [shift={(\g:.3)}] {$\x$};
		\foreach \y/\g in {-1/-10,1/135} 
		\draw[fill] (0,\y) circle(.5pt)node [shift={(\g:.3)}] {$\y$};
		%\draw[blue] (1,-4)--(1,5);
		\draw[red!80] (-4,1)--(5,1);
		\draw (2,-2) node{Hình 18b};
	\end{tikzpicture}
	\begin{enumerate}[a)]
		\item$y=\dfrac{x^2+2x-1}{x^2+1}$;
		\item $y=\dfrac{2x^2+x+1}{x-1}$;
		\item $y=\dfrac{2x^2-2}{x^2+2}$.
	\end{enumerate}
	\loigiai{
		\begin{enumerate}[a)]
			\item Ở hình 18a ta thấy đồ thị cắt trục hoành tại $\pm 1$, cắt trục tung tại $-1$ và nhận đường thẳng $y=2$ là tiệm cận ngang và không có tiệm cận đứng. \\
			Vậy trong $3$ hàm số chỉ có hàm $y=\dfrac{2x^2-2}{x^2+1}$ thoả mãn.
			\item Ở hình 18b ta thấy đồ thị có tiệm cận ngang $y=1$ và không có tiệm cận đứng.\\
			Vậy trong $3$ hàm số trên chỉ có hàm số $y=\dfrac{x^2+2x-1}{x^2+1}$ thoả mãn.
		\end{enumerate}
	}
\end{bt}

\begin{bt}%[2D1H4-1]
% nguon: Archive/2024/TaiLieuGiangDay/1.HamSo/1. K12 - 4in1 - Hamso/ThubaiGoc/Bai4-Dang3.tex #2456
%[Dự án TL12New-4in1-NCT]
	Cho bảng biến thiên của hàm số $y=f(x)$ như sau.
	\begin{center}
		\begin{tikzpicture}
			\tikzset{double style/.append style = {draw=\tkzTabDefaultWritingColor,double=\tkzTabDefaultBackgroundColor,double distance=2pt}}
			\tkzTabInit[nocadre=false, lgt=1, espcl=3]{$x$ /1,$y$ /1.7}{$-\infty$,$-\tfrac{1}{2}$,$+\infty$}
			\tkzTabVar{-/ $-\infty $ / , +D+/ $+\infty$ / $+\infty$ , -/ $3$ /}
		\end{tikzpicture}
	\end{center}
	Đồ thị của hàm số đã cho có tổng số bao nhiêu đường tiệm cận đứng và tiệm cận ngang?
	\loigiai{
		Đồ thị hàm số có một đường tiệm cận đứng là đường thẳng $x=-\dfrac{1}{2}$ và có một đường tiệm cận ngang là đường thẳng $y=3$.\\
		Vậy đồ thị hàm số đã cho có tổng số 2 đường tiệm cận đứng và ngang.
	}
\end{bt}

\begin{bt}%[2D1H4-1]
% nguon: Archive/2024/TaiLieuGiangDay/1.HamSo/1. K12 - 4in1 - Hamso/ThubaiGoc/Bai4-Dang3.tex #2457
%[Dự án TL12New-4in1-NCT]
	Cho hàm số $y=f(x)$ xác định trên $\mathbb{R}\backslash \left\{0\right\},$ liên tục trên mỗi khoảng xác định và có bảng biến thiên như hình bên.
	\begin{center}
		\begin{tikzpicture}[>=stealth]
			\tikzset{double style/.append style = {draw=\tkzTabDefaultWritingColor,double=\tkzTabDefaultBackgroundColor,double distance=2pt}}
			\tkzTabInit[nocadre=false,lgt=1,espcl=2.2]{$x$/.6,$y'$/.7,$y$/1.5}{$-\infty$,$0$,$1$,$+\infty$}
			\tkzTabLine{,-, d ,+,z,-,} 
			\tkzTabVar{+/$+\infty$ / , -D- / $-1$ /$-\infty$,+/$2$,-/$-\infty$}
		\end{tikzpicture}
	\end{center} 
Xác định các đường tiệm cận đứng và tiệm cận ngang (nếu có) của đồ thị hàm số.
	\loigiai{
		Do $\lim\limits_{x \to +\infty} y=-\infty$ và $\lim\limits_{x \to -\infty} y=+\infty$  nên đồ thị hàm số không có tiệm cận ngang.\\
		Do $\lim\limits_{x \to 0^+} y=-\infty$ suy ra $x=0$ là tiệm cận đứng của đồ thị hàm số.
	} 
\end{bt}

\begin{bt}%[2D1H4-1]
% nguon: Archive/2024/TaiLieuGiangDay/1.HamSo/1. K12 - 4in1 - Hamso/ThubaiGoc/Bai4-Dang3.tex #2458
%[Dự án TL12New-4in1-NCT]
	Cho hàm số $y=f(x)$ xác định trên $\mathbb{R}\setminus\left\{\pm1\right\}$ liên tục trên mỗi khoảng xác định và có bảng biến thiên như hình vẽ. 
	\begin{center}
		\begin{tikzpicture}
			\tikzset{double style/.append style = {draw=\tkzTabDefaultWritingColor,double=\tkzTabDefaultBackgroundColor,double distance=2pt}}
			\tikzset{double style/.append style = {double distance=0.5pt}} 
			\tkzTabInit[nocadre=false,lgt=1,espcl=2.5]
			{$x$/.7,$y'$ /.7, $y$ /1.6}
			{$-\infty$ ,$-1$,$0$,$1$,$+\infty$}
			\tkzTabLine{,-,d,-,0,+,d,+,}
			\tkzTabVar {+/$-2$,-D+/$-\infty$/$+\infty$,-/$1$,+D-/$+\infty$/$-\infty$,+/$-2$}
		\end{tikzpicture}
	\end{center}
Xác định các đường tiệm cận đứng và tiệm cận ngang của đồ thị hàm số.
	\loigiai{
		Dựa vào bảng biến thiên ta có:\\
		$\lim\limits_{x\to -1^\pm}f(x)=\pm\infty$. 
		$\lim\limits_{x\to 1^\pm}f(x)=\mp\infty$.\\
		Do đó $x=1$ và $x=-1$ là các đường tiệm cận đứng của đồ thị hàm số.\\
		Lại có $\lim\limits_{x\to \pm\infty}f(x)=-2$. Do đó $y=-2$ là tiệm cận ngang của đồ thị hàm số.
	}
\end{bt}

\begin{bt}%[2D1H4-1]
% nguon: Archive/2024/TaiLieuGiangDay/1.HamSo/1. K12 - 4in1 - Hamso/ThubaiGoc/Bai4-Dang3.tex #2459
%[Dự án TL12New-4in1-NCT]
	Cho hàm số $y=f(x)$ có bảng biến thiên như hình vẽ dưới đây. Đồ thị hàm số $y=f(x)$ có bao nhiêu đường tiệm cận đứng và tiệm cận ngang?
	\begin{center}
		\begin{tikzpicture}
			\tikzset{double style/.append style = {draw=\tkzTabDefaultWritingColor,double=\tkzTabDefaultBackgroundColor,double distance=2pt}}
			\tkzTabInit[nocadre=false,lgt=1,espcl=2]
			{$x$ /0.7,$y'$ /0.7,$y$ /1.8}
			{$-\infty$,$-2$,$0$,$2$,$+\infty$}
			\tkzTabLine{,+,d,-,0,+,d,}
			\tkzTabVar{-/$-2$/,+D+/ $3$ / $+\infty$,-/$-2$/,+D/ $+\infty$ / }
			\draw[fill=black!50,opacity=0.8] ($(N43)+(0.1ex,0)$) rectangle (T21);
		\end{tikzpicture}
	\end{center}
	\loigiai
	{
		Dựa vào bảng biến thiên ta có:
		\begin{itemize}
			\item $\displaystyle\lim_{x \to -\infty} y=-2$.
			\item $\displaystyle\lim_{x \to -2^+} y=+\infty$.
			\item $\displaystyle\lim_{x \to 2^-} y=+\infty$.
		\end{itemize}
		Vậy đồ thị hàm số có hai đường tiệm cận đứng và một đường tiệm cận ngang.
	}
\end{bt}

\begin{bt}%[2D1H4-1]
% nguon: Archive/2024/TaiLieuGiangDay/1.HamSo/1. K12 - 4in1 - Hamso/ThubaiGoc/Bai4-Dang3.tex #2465
%[Dự án TL12New-4in1-NCT]
	Hàm số $y=f(x)$ có bảng biến thiên dưới đây. 
	\begin{center}
		\begin{tikzpicture}[>=stealth,scale=.8]
			\tkzTabInit[nocadre=false,lgt=1.2,espcl=3,deltacl=.5]
			{$x$ /1.2, $y'$ /1.2, $y$ /2.5}
			{$-\infty$,$-2$,$0$,$1$,$+\infty$}
			\tkzTabLine{ ,-,d,-, ,+,d,-, }
			\tkzTabVar{+/$-1$,-D+/$-\infty$/$2$,-/$-4$,+/$3$,-/$0$}
		\end{tikzpicture}
	\end{center}
	Xác định các tiệm cận đứng và tiệm cận ngang của đồ thị hàm số $y=f(x)$.
	\loigiai{
		Qua bảng biến thiên ta có $\lim\limits_{x\to-\infty} f(x)=-1$ và $\lim\limits_{x\to+\infty} f(x)=0$ nên đồ thị hàm số có hai đường tiệm cận ngang: $y=-1$ và $y=0$.\\
		Lại có $\lim\limits_{x\to-2^-} f(x)=-\infty$ nên đồ thị hàm số có một tiệm cận đứng $x=-2$.}
\end{bt}

\begin{bt}%[2D1H4-1]
% nguon: Archive/2024/TaiLieuGiangDay/1.HamSo/1. K12 - 4in1 - Hamso/ThubaiGoc/Bai4-Dang3.tex #2466
%[Dự án TL12New-4in1-NCT]
	Tìm các tiệm cận của đồ thị hàm số sau:
	\begin{multicols}{3}
		\begin{tikzpicture}[scale=.8,>=stealth, font=\footnotesize, line join=round, line cap=round]
			\def\xmin{-2} \def\xmax{4}
			\def\ymin{-3} \def\ymax{3}
			%\draw[color=gray!50,dashed] (\xmin,\ymin) grid (\xmax,\ymax);
			\draw[->] (\xmin,0)--(\xmax,0) node [below]{$x$};
			\draw[->] (0,\ymin)--(0,\ymax) node [left]{$y$};
			\fill (0,0) circle (1pt) node[shift={(135:2.5mm)}]{$O$};
			%\node at (current bounding box.south) [below=-2pt] {a) $y=\dfrac{2x-3}{5x^{2}-15x+10}$};
			\clip (\xmin+0.1,\ymin+0.1) rectangle (\xmax-0.1,\ymax-0.1);
			\draw[red,smooth,samples=300,domain=\xmin:0.99] plot(\x,{(2*(\x)-3)/(5*(\x)^2-15*(\x)+10)});
			\draw[red,smooth,samples=300,domain=1.01:1.99] plot(\x,{(2*(\x)-3)/(5*(\x)^2-15*(\x)+10)});
			\draw[red,smooth,samples=300,domain=2.01:\xmax] plot(\x,{(2*(\x)-3)/(5*(\x)^2-15*(\x)+10)});
			\draw[dashed] (1,\ymin)--(1,\ymax);
			\draw[dashed] (2,\ymin)--(2,\ymax);
			\foreach \s/\t in {2/-45,1/-45}
			\fill (\s,0) circle (1pt) node[shift={(\t:3mm)}]{$\s$};
		\end{tikzpicture}
		\begin{tikzpicture}[scale=.6,>=stealth, font=\footnotesize, line join=round, line cap=round]
			\def\xmin{-4} \def\xmax{4}
			\def\ymin{-3} \def\ymax{5}
			%\draw[color=gray!50,dashed] (\xmin,\ymin) grid (\xmax,\ymax);
			\draw[->] (\xmin,0)--(\xmax,0) node [below]{$x$};
			\draw[->] (0,\ymin)--(0,\ymax) node [right]{$y$};
			\fill (0,0) circle (1pt) node[shift={(-135:2.5mm)}]{$O$};
			%\node at (current bounding box.south) [below=-2pt] {b) $y=\dfrac{x^{2}+x-1}{x}$};
			\clip (\xmin+0.1,\ymin+0.1) rectangle (\xmax-0.1,\ymax-0.1);
			\draw[red,smooth,samples=300,domain=\xmin:-0.01] plot(\x,{((\x)^2+(\x)-1)/(\x)});
			\draw[red,smooth,samples=300,domain=0.01:\xmax] plot(\x,{((\x)^2+(\x)-1)/(\x)});
			\draw[dashed,smooth,samples=300,domain=\xmin:\xmax] plot(\x,{(\x)+1});
			\foreach \s/\t in {-1/-90}
			\fill (\s,0) circle (1pt) node[shift={(\t:3mm)}]{$\s$};
			\foreach \p/\r in {1/-20}
			\fill (0,\p) circle (1pt) node[shift={(\r:3mm)}]{$\p$};
		\end{tikzpicture}
		\begin{tikzpicture}[scale=.6,>=stealth, font=\footnotesize, line join=round, line cap=round]
			\def\xmin{-4} \def\xmax{4}
			\def\ymin{-3} \def\ymax{5}
			%\draw[color=gray!50,dashed] (\xmin,\ymin) grid (\xmax,\ymax);
			\draw[->] (\xmin,0)--(\xmax,0) node [below]{$x$};
			\draw[->] (0,\ymin)--(0,\ymax) node [right]{$y$};
			\fill (0,0) circle (1pt) node[shift={(-135:2.5mm)}]{$O$};
			%\node at (current bounding box.south) [below=-2pt] {c) $y=\dfrac{16x^{2}-8x}{16x^{2}+1}$};
			\clip (\xmin+0.1,\ymin+0.1) rectangle (\xmax-0.1,\ymax-0.1);
			\draw[red,smooth,samples=300,domain=\xmin:\xmax] plot(\x,{(16*(\x)^2-8*(\x))/(16*(\x)^2+1)});
			\draw[dashed](\xmin,1)--(\xmax,1);
			\foreach \p/\r in {1/45}
			\fill (0,\p) circle (1pt) node[shift={(\r:3mm)}]{$\p$};
		\end{tikzpicture}
	\end{multicols}
	\begin{multicols}{3}
		\begin{enumerate}
			\item $y=\dfrac{2x-3}{5x^{2}-15x+10}$;
			\item $y=\dfrac{x^{2}+x-1}{x}$;
			\item $y=\dfrac{16x^{2}-8x}{16x^{2}+1}$.
		\end{enumerate}
	\end{multicols}
	\begin{center}
		{\footnotesize Hình 12}
	\end{center}
	\loigiai{
		\begin{enumerate}
			\item Hàm số $y=\dfrac{2x-3}{5x^{2}-15x+10}$ có tập xác định $\mathscr{D}=\mathbb{R} \setminus \left\lbrace 1;2\right\rbrace$.
			\begin{itemize}
				\item 
				Ta có $\lim\limits_{x \rightarrow 1^{-}} \dfrac{2x-3}{5x^{2}-15x+10}=-\infty$; $\lim\limits_{x \rightarrow 1^{+}} \dfrac{2x-3}{5x^{2}-15x+10}=+\infty$.\\
				Ta cũng có $\lim\limits_{x \rightarrow 2^{-}} \dfrac{2x-3}{5x^{2}-15x+10}=-\infty$; $\lim\limits_{x \rightarrow 2^{+}} \dfrac{2x-3}{5x^{2}-15x+10}=+\infty$.\\
				Suy ra đường thẳng $x=1$, $x=2$ là hai tiệm cận đứng của đồ thị hàm số.
				\item Ta có $\lim\limits_{x \rightarrow -\infty} \dfrac{2x-3}{5x^{2}-15x+10}=0$; $\lim\limits_{x \rightarrow+\infty} \dfrac{2x-3}{5x^{2}-15x+10}=0$.\\
				Suy ra đồ thị hàm số có tiệm cận ngang là đường thẳng $y=0$.
			\end{itemize}
			Vậy hàm số $y=\dfrac{2x-3}{5x^{2}-15x+10}$ có hai tiệm cận đứng là $x=1$, $x=2$ và tiệm cận ngang là $y=0$.
			\item Hàm số $y=f(x)=\dfrac{x^{2}+x-1}{x}$ có tập xác định $\mathscr{D}=\mathbb{R} \setminus \left\lbrace 0\right\rbrace$.
			\begin{itemize}
				\item Ta có $\lim\limits_{x \rightarrow 0^{-}} \dfrac{x^{2}+x-1}{x}=+\infty$; $\lim\limits_{x \rightarrow 0^{+}} \dfrac{x^{2}+x-1}{x}=-\infty$.\\
				Suy ra đường thẳng $x=0$ là một tiệm cận đứng của đồ thị hàm số.
				\item Ta có $\begin{aligned}[t]
					a&=\lim\limits_{x \rightarrow+\infty} \dfrac{f(x)}{x}=\lim\limits_{x \rightarrow+\infty} \dfrac{x+1-\dfrac{1}{x}}{x}=1;\\
					b&=\lim\limits_{x \rightarrow+\infty}[f(x)-ax]=\lim\limits_{x \rightarrow+\infty}\left(\dfrac{x^{2}+x-1}{x}-x\right)=\lim\limits_{x \rightarrow+\infty} \dfrac{x-1}{x}=1.
				\end{aligned}$\\
				Ta cũng có $\lim\limits_{x \rightarrow-\infty} \dfrac{f(x)}{x}=1$; $\lim\limits_{x \rightarrow-\infty}[f(x)-x]=1$.\\
				Do đó, đồ thị hàm số có tiệm cận xiên là đường thẳng $y=x+1$.
			\end{itemize}
			Vậy hàm số $y=\dfrac{x^{2}+x-1}{x}$ có tiệm cận đứng là $x=0$ và tiệm cận xiên là $y=x+1$.
			\item Hàm số $y=f(x)=\dfrac{16x^{2}-8x}{16x^{2}+1}$ có tập xác định $\mathscr{D}=\mathbb{R}$.\\
			Ta có $\lim\limits_{x \rightarrow -\infty} \dfrac{16x^{2}-8x}{16x^{2}+1}=1$; $\lim\limits_{x \rightarrow+\infty} \dfrac{16x^{2}-8x}{16x^{2}+1}=1$.\\
			Vậy đồ thị hàm số có tiệm cận ngang là đường thẳng $y=1$.
		\end{enumerate}
	}
\end{bt}

\begin{bt}%[2D1V4-1]
% nguon: Archive/2024/TaiLieuGiangDay/1.HamSo/1. K12 - 4in1 - Hamso/ThubaiGoc/Bai4-Dang3.tex #2467
%[Dự án TL12New-4in1-NCT]
	Cho hàm số bậc ba $y=f(x)$ có bảng biến thiên như hình vẽ
	\begin{center}
		\begin{tikzpicture}
			\tkzTabInit[lgt=1.2,espcl=2]
			{$x$/1, $f’(x)$/1, $f(x)$/2}
			{$-\infty$,$0$,$1$,$+\infty$}
			\tkzTabLine{ ,-,z,+,z,-, }
			\tkzTabVar{+/$+\infty$,-/$-1$,+/$3$,-/$-\infty$}
		\end{tikzpicture}
	\end{center}
	Xác định số đường tiệm cận đứng của đồ thị hàm số $y=\dfrac{x}{f(x)+1}$.
	\loigiai
	{
		Dựa vào bảng biến thiên ta thấy
		$$f(x)+1=0\Leftrightarrow f(x)=-1\Leftrightarrow \hoac{& x^2=0 \\ & x=a>1}\Leftrightarrow\hoac{& x=0\ \text{(nghiệm kép)} \\ & x=a>1.}$$
		Do đó $y=\dfrac{x}{f(x)+1}=\dfrac{x}{kx^2(x-a)}=\dfrac{1}{kx(x-a)}$.\\
		Suy ra $\lim \limits_{x \to 0} y=\infty$, $\lim \limits_{x \to a} y=\infty$ và $x=0$, $x=a$ là các đường tiệm cận đứng.
	}
\end{bt}

\begin{bt}%[2D1V4-1]
% nguon: Archive/2024/TaiLieuGiangDay/1.HamSo/1. K12 - 4in1 - Hamso/ThubaiGoc/Bai4-Dang3.tex #2468
%[Dự án TL12New-4in1-NCT]
	Cho hàm số $y=f(x)$ có bảng biến thiên như hình vẽ sau:
	\begin{center}
		\begin{tikzpicture}
			\tkzTabInit[nocadre,lgt=.8,espcl=3.5,deltacl=0.6]
			{$x$/.6,$y’$/.6,$y$/2}
			{$-\infty$,$-2$,$0$,$+\infty$}
			\tkzTabLine{,,,+,,-,}
			\draw[pattern=north east lines	,draw=none] (T11) rectangle (N23);
			\draw (N21)--(N23);
			\draw ([xshift=2pt]N21)--([xshift=2pt]N23)
			;
			\draw (N31)--(N33);
			\draw ([xshift=2pt]N31)--([xshift=2pt]N33)
			;
			\draw  (T01)--(T21) (T02)--(T22) (T10)--(T13);
			\draw (N23) node[right](A){$-\infty$} 
			(N32) node[below,xshift=-.5cm](B){$+\infty$}
			(N32) node[below right](C){$1$}
			(N43) node[](D){$0$};
			\draw[-stealth] (A)--(B);
			\draw[-stealth] (C)--(D);
		\end{tikzpicture}
	\end{center}
	Hỏi đồ thị hàm số $y=\dfrac{1}{2f(x)-1}$ có tất cả bao nhiêu đường tiệm cận đứng?
	\loigiai{
		Dựa vào bảng biến thiên, ta có $2f(x)-1=0 \Leftrightarrow f(x)=\dfrac{1}{2}\Leftrightarrow\hoac{&x=x_1\in (-2;0)\\&x=x_2\in (0;+\infty).}$\\
		Do $\lim\limits_{x\to x_1}y=\lim\limits_{x\to x_1}\dfrac{1}{2f(x)-1}=\infty$, $\lim\limits_{x\to x_2}y=\lim\limits_{x\to x_2}\dfrac{1}{2f(x)-1}=\infty$ nên $x=x_1$ và $x=x_2$ là các đường tiệm cận đứng của đồ thị hàm số.\\
		Vậy đồ thị hàm số có $2$ đường tiệm cận đứng.	
	}
\end{bt}

\begin{bt}%[2D1C4-1]
% nguon: Archive/2024/TaiLieuGiangDay/1.HamSo/1. K12 - 4in1 - Hamso/ThubaiGoc/Bai4-Dang3.tex #2470
%[Dự án TL12New-4in1-NCT]
	\immini{Cho hàm số $f(x)=ax^3+bx^2+cx+d$ $(a\ne 0)$ có đồ thị như hình vẽ dưới đây.
		Gọi $S$ là tập các giá trị nguyên của $m$ thuộc khoảng $(-2025;2026)$ để đồ thị hàm số $g(x)=\dfrac{(x-1)\sqrt{f(x)}}{\left(f(x)-2\right)(x^2-2mx+m+2)}$ có tất cả $5$ đường tiệm cận (tiệm cận đứng và tiệm cận ngang). Xác định số phần tử của tập $S$.
	}{	\begin{tikzpicture}[scale=1,>=stealth]
			\path
			(1,1) coordinate (A);
			\draw[->](-3,0)--(3,0)node[below]{$x$};
			\draw[->](0,-1)--(0,3)node[left]{$y$};
			\draw[smooth,samples=200,domain=-2.2:2.2]plot(\x,{1/2*(\x)^3-3/2*(\x)+1});
			\draw[dashed](-1,0)--(-1,2)--(0,2) (2,0)--(2,2)--(0,2);
			\draw(0,0)node[below left]{$O$} (1,0)node[below]{$1$} (2,0)node[below]{$2$}
			(-1,0)node[below]{$-1$} (-2,0)node[below left]{$-2$}(0,2)node[above right]{$2$};
	\end{tikzpicture}}
	\loigiai{
		\begin{center}
			\begin{tikzpicture}[scale=1,>=stealth]
				\path
				(1,1) coordinate (A);
				\draw[->](-3,0)--(3,0)node[below]{$x$};
				\draw[->](0,-1)--(0,3)node[left]{$y$};
				\draw[smooth,samples=200,domain=-2.2:2.2]plot(\x,{1/2*(\x)^3-3/2*(\x)+1});
				\draw[dashed](-1,0)--(-1,2)--(0,2) (2,0)--(2,2)--(0,2) (-2.5,2)--(2.5,2);
				\draw(0,0)node[below left]{$O$} (1,0)node[below]{$1$} (2,0)node[below]{$2$}
				(-1,0)node[below]{$-1$} (-2,0)node[below left]{$-2$}(0,2)node[above right]{$2$};
			\end{tikzpicture}
		\end{center}
		Từ đồ thị hàm số ta có $f(x)=0\Leftrightarrow\hoac{&x=-2\\&x=1.}$\\
		Từ đó suy ra $f(x)=a(x+2)(x-1)^2$.\\
		Từ đồ thị hàm số ta có $f(x)-2=0\Leftrightarrow f(x)=2\Leftrightarrow\hoac{&x=-1\\&x=2.}$\\
		Từ đó suy ra $f(x)=a(x+1)^2(x-2)$.\\
		Xét hàm số 
		$\\\begin{aligned}
			g(x)=\dfrac{(x-1)\sqrt{f(x)}}{\left(f(x)-2\right)(x^2-2mx+m+2)}
			&=\dfrac{(x+1)\sqrt{a(x+2)(x-1)^2}}{a (x+1)^2(x-2)(x^2-2mx+m+2)}\\
			&=\dfrac{|x-1|\sqrt{a(x+2)}}{a(x+1)(x-2)(x^2-2mx+m+2)}.
		\end{aligned}$\\
		Hàm số $g(x)$ xác định khi và chỉ khi $\heva{&x+2\geq 0\\&x+1\ne 0\\&x-2\ne 0\\&x^2-2mx+m+2\ne 0}\Leftrightarrow\heva{&x\geq-2\\&x\ne -1\\&x\ne 2\\&x^2-2mx+m+2\ne 0.}$\\
		Ta có $\lim\limits_{x\to +\infty}g(x)=0\Rightarrow y=0$ là tiệm cận ngang duy nhất của đồ thị hàm số $g(x)$.\\
		Vì vậy yêu cầu bài toán tương đương với đồ thị hàm số $g(x)$ có $4$ tiệm cận đứng, điều này tương đương với phương trình $h(x)=x^2-2mx+m+2=0$ có hai nghiệm phân biệt $x_1,x_2$ với $-2\leq x_1<x_2; x_1,x_2\notin\{1;2;-1\}$ khi và chỉ khi
		$$\heva{&\triangle'>0\\&S>-4\\&1.h(-2)\leq 0\\&h(1)\ne 0\\&h(2)\ne 0\\&h(-1)\ne0}\Leftrightarrow\heva{&m^2-m-2>0\\&2m>-4\\&1.(4+4m+m+2)\leq 0\\&1-2m+m+2\ne 0\\& 1+2m+m+2\ne 0}\Leftrightarrow \hoac{&2<m\ne 3\\&-\dfrac{6}{5}\leq m<-1.}$$
		Vậy $m\in\{4;...;2025\}$.}
\end{bt}

\begin{bt}%[2D1H2-4]
% nguon: Archive/2024/TaiLieuGiangDay/1.HamSo/1. K12 - 4in1 - Hamso/ThubaiGoc/Bai2-Dang4.tex #2484
[Sách tham khảo, Mức độ H]%[Dự án Bài giảng 12 - Trung Anh]
	Cho hàm số $y=x^3-3(m+1)x^2+3(7m-3)x$. Tìm tất cả các giá trị nguyên của tham số $m$ để hàm số không có cực trị.
	\loigiai{
		Xét hàm số $y=x^3-3(m+1)x^2+3(7m-3)x (1)\Rightarrow y'=3x^3-6(m+1)x+3(7m-3)$.\\
		Ta có: $y'=0\Leftrightarrow x^2-2(m+1)x+7m-3=0$ \qquad$(2)$\\
		Hàm số đã cho không có cực trị\\
		$\Leftrightarrow$ Phương trình $y'=0$ vô nghiệm hoặc có nghiệm kép\\
		$\Leftrightarrow \triangle'_{(2)}\le0\Leftrightarrow (m+1)^2-1\cdot(7m-3)\le0\Leftrightarrow m^2-5m+4\le0\Leftrightarrow 1\le m\le 4$.\\
		Do $m$ là số nguyên nên $m\in\{1;2;3;4\}$. Vậy tập $S$ có 4 phần tử.
	}
\end{bt}

\begin{bt}%[2D1H2-4]
% nguon: Archive/2024/TaiLieuGiangDay/1.HamSo/1. K12 - 4in1 - Hamso/ThubaiGoc/Bai2-Dang4.tex #2485
[Sách tham khảo, Mức độ H]%[Dự án Bài giảng 12 - Trung Anh]
	Tìm tất cả các giá trị của tham số $m$ để hàm số $y=\dfrac{1}{3}x^3+mx^2+(2m-1)x-1$ có cực trị.
	\loigiai{
		Ta có $y{'}=x^2+2mx+(2m-1)$.\\
		Hàm số đã cho là hàm số bậc ba nên để hàm số có cực trị thì $m^2-2m+1>0 \Leftrightarrow m \ne 1$.	
	}
\end{bt}

\begin{bt}%[2D1H2-4]
% nguon: Archive/2024/TaiLieuGiangDay/1.HamSo/1. K12 - 4in1 - Hamso/ThubaiGoc/Bai2-Dang4.tex #2486
[Sách tham khảo, Mức độ H]%[Dự án Bài giảng 12 - Trung Anh]
	Tìm tất cả tham số thực $m$ để hàm số $y=\dfrac{1}{3}(m+2)x^3+x^2+\dfrac{1}{3}mx-2$ có cực đại, cực tiểu.
	\loigiai{
		Ta có $y'=(m+2)x^2+2x+\dfrac{1}{3}m$.\\
		Điều kiện bài toán tương đương phương trình $y'=0$ có hai nghiệm phân biệt, hay	
		$$\heva{& m+2\ne 0 \\ & 1-\dfrac{1}{3}m(m+2)>0}\Leftrightarrow \heva{& m\ne -2 \\ & -3<m<1}\Leftrightarrow m\in (-3;-2)\cup (-2;1).$$
	}
\end{bt}

\begin{bt}%[2D1H2-4]
% nguon: Archive/2024/TaiLieuGiangDay/1.HamSo/1. K12 - 4in1 - Hamso/ThubaiGoc/Bai2-Dang4.tex #2487
[Sách tham khảo, Mức độ H]%[Dự án Bài giảng 12 - Trung Anh]
	Cho hàm số $y=-\dfrac{1}{3}x^3+2x^2+(m-2)x+m^2.$ Tìm tất cả các giá trị của $m$ để hàm số không có cực trị.
	\loigiai{ Tập xác định $\mathscr{D}=\mathbb{R}.\\ y'=-x^2+4x+m-2$. \\ Hàm số không có cực trị khi và chỉ khi $\Delta' \le 0 \Leftrightarrow 4+m-2 \le 0 \Leftrightarrow m \le -2$.
	}
\end{bt}

\begin{bt}%[2D1H2-4]
% nguon: Archive/2024/TaiLieuGiangDay/1.HamSo/1. K12 - 4in1 - Hamso/ThubaiGoc/Bai2-Dang4.tex #2488
[Sách tham khảo, Mức độ V]%[Dự án Bài giảng 12 - Trung Anh]
	Với giá trị nào của tham số $m$ thì hàm số $y=\dfrac{2}{3}x^3-3x^2+mx-1$ có hai điểm cực trị $x_1$, $x_2$ sao cho $x_2=2x_1$.
	\loigiai{
		Ta có $y'=2x^2-6x+m$. Giả sử hàm số có $2$ cực trị thỏa yêu cầu bài toán. Khi đó $x_1,x_2$ là các nghiệm của phương trình $y'=0$ và thỏa
		$$\heva{&x_1+x_2=3\\ &x_1x_2=\dfrac{m}{2}\\ &x_2=2x_1}\Leftrightarrow \heva{&3x_1=3\\ &2x_1^2=\dfrac{m}{2}\\ &x_2=2x_1}\Leftrightarrow \heva{&x_1=1\\ &m=4\\ &x_2=2.}$$
		Thử lại với $m=4$ thấy thỏa mãn. Vậy $m=4$ là giá trị cần tìm.
	}
\end{bt}

\begin{bt}%[2D1H2-4]
% nguon: Archive/2024/TaiLieuGiangDay/1.HamSo/1. K12 - 4in1 - Hamso/ThubaiGoc/Bai2-Dang4.tex #2492
[Sách tham khảo, Mức H]%[Dự án bài giảng 12, Trung Anh]
	Tìm điều kiện của tham số $m$ để hàm số $y=\dfrac{1}{3}x^3-mx^2+\left(2m+3\right)x+m$ có cực trị.
	\loigiai{
	Tập xác định $\mathscr{D}=\mathbb{R}$.\\
	Ta có $y=\dfrac{1}{3}x^3-mx^2+\left(2m+3\right)x+m$ nên $y'=x^2-2mx+2m+3$.\\
	$y'=0 \Leftrightarrow x^2-2mx+2m+3=0$.\\
	$\Delta'=m^2-2m-3$.\\
	Hàm số $y=\dfrac{1}{3}x^3-mx^2+\left(2m+3\right)x+m$ có cực trị khi và chỉ khi phương trình $y'=0$ có hai nghiệm phân biệt, tức là $\Delta'>0 \Leftrightarrow m^2-2m-3>0 \Leftrightarrow m\in \left(-\infty;-1\right)\cup \left(3;+\infty\right)$.
	}
\end{bt}

\begin{bt}%[2D1H2-4]
% nguon: Archive/2024/TaiLieuGiangDay/1.HamSo/1. K12 - 4in1 - Hamso/ThubaiGoc/Bai2-Dang4.tex #2493
[Sách tham khảo, Mức H]%[Dự án bài giảng 12, Trung Anh]
	Tìm m để hàm số $y=x^3-2x^2+mx-3$ đạt cực đại tại điểm $x=1$.
	\loigiai{
	Ta có $y'=3x^2-4x+m$.\\
	Hàm số có hai cực trị khi và chỉ khi $y'=0$ có hai nghiệm phân biệt.\\
	Tương đương
	\[\triangle>0\Leftrightarrow 16-12m>0\Leftrightarrow m<\dfrac{4}{3}.\]
	Ta có $y'=0\Leftrightarrow\hoac{&x_1=\dfrac{-b+\sqrt{\triangle}}{2a}=\dfrac{4+\sqrt{16-12m}}{6}\\&x_2=\dfrac{-b-\sqrt{\triangle}}{2a}=\dfrac{4-\sqrt{16-12m}}{6}.}$\\
	Bảng biến thiên
	\begin{center}
		\begin{tikzpicture}
			\tikzset{double style/.append style={double distance=2pt}}
			\tkzTabInit[nocadre=true,lgt=1.3,espcl=3.2,deltacl=0.7]
			{$x$/1.2,$f'(x)$/0.6,$f(x)$/2}
			{$-\infty$,$\dfrac{4-\sqrt{16-12m}}{6}$,$\dfrac{4+\sqrt{16-12m}}{6}$,$+\infty$}
			\tkzTabLine{,+,0,-,0,+,}
			\tkzTabVar{-/$-\infty$ ,+/$f(x_2)$,-/$f(x_1)$,+/ $+\infty$}
		\end{tikzpicture}
	\end{center}
	Hàm số đạt cực đại tại $x=1$
	\begin{align*}
		\Leftrightarrow&\, \dfrac{4-\sqrt{16-12m}}{6}=1\\
		\Leftrightarrow&\, \sqrt{16-12m}=-2\, (\text{vô lí}).
	\end{align*}
	Vậy không có giá trị nào của $m$ thỏa mãn.
	}
\end{bt}

\begin{bt}%[2D1H2-4]
% nguon: Archive/2024/TaiLieuGiangDay/1.HamSo/1. K12 - 4in1 - Hamso/ThubaiGoc/Bai2-Dang4.tex #2494
[Sách tham khảo, Mức H]%[Dự án bài giảng 12, Trung Anh]
	Tìm các giá trị của tham số $m$ để hàm số $y=x^3-3x^2+mx-1$ có hai điểm cực trị $x_1$, $x_2$ sao cho $x_1^2+x_2^2-x_1x_2=13$.
	\loigiai{
	Tập xác định: $\mathscr{D}=\mathbb{R}$.\\
	$y'=3x^2-6x+m$, cho $y'=0\Leftrightarrow 3x^2-6x+m=0$.\\
	Hàm số có hai cực trị $x_1$, $x_2$ thì $\Delta'>0\Leftrightarrow 9-3m>0\Leftrightarrow m<3$.\\
	Theo định lí Vi-ét ta có $\heva{&x_1+x_2=2\\&x_1 x_2=\dfrac{m}{3}.}$\\
	Ta có $x_1^2+x_2^2-x_1x_2=13 \Leftrightarrow \left(x_1+x_2\right)^2-3x_1x_2-13=0 \Rightarrow 4-m=13 \Leftrightarrow m=-9 \text{ (thỏa mãn)}.$	
	}
\end{bt}

\begin{bt}%[2D1H2-4]
% nguon: Archive/2024/TaiLieuGiangDay/1.HamSo/1. K12 - 4in1 - Hamso/ThubaiGoc/Bai2-Dang4.tex #2495
[Sách tham khảo, Mức H]%[Dự án bài giảng 12, Trung Anh]
	Tìm tất cả các giá trị của tham số $m$ để hàm số $y=x^3-3mx^2+3(m^2-1)x$ đạt cực tiểu tại $x=1$.
	\loigiai{
		Ta có $y'=3x^2-6mx+3(m^2-1)$.\\
		Điều kiện cần để hàm số đạt cực tiểu tại $x=1$ là 
		\[y'(1)=0\Leftrightarrow 3-6m+3(m^2-1)=0\Leftrightarrow 3m^2-6m=0\Leftrightarrow\hoac{&m=0\\&m=2.}\]
		Với $m=0$, ta có $y'=3x^2-3$. Khi đó $y'=0\Leftrightarrow x=\pm1$.\\
		Bảng biến thiên
		\begin{center}
			\begin{tikzpicture}
				\tikzset{double style/.append style={double distance=2pt}}
				\tkzTabInit[nocadre=true,lgt=1.3,espcl=3.2,deltacl=0.7]
				{$x$/1.2,$f'(x)$/0.6,$f(x)$/2}
				{$-\infty$,$-1$,$1$,$+\infty$}
				\tkzTabLine{,+,0,-,0,+,}
				\tkzTabVar{-/$-\infty$ ,+/$f(-1)$,-/$f(1)$,+/ $+\infty$}
			\end{tikzpicture}
		\end{center}
		Vậy $m=0$ thỏa mãn.\\
		Với $m=2$, ta có $y'=3x^2-12x+9$. Khi đó $y'=0\Leftrightarrow\hoac{&x=1\\&x=3.}$\\  
		Bảng biến thiên
		\begin{center}
			\begin{tikzpicture}
				\tikzset{double style/.append style={double distance=2pt}}
				\tkzTabInit[nocadre=true,lgt=1.3,espcl=3.2,deltacl=0.7]
				{$x$/1.2,$f'(x)$/0.6,$f(x)$/2}
				{$-\infty$,$1$,$3$,$+\infty$}
				\tkzTabLine{,+,0,-,0,+,}
				\tkzTabVar{-/$-\infty$ ,+/$f(1)$,-/$f(3)$,+/ $+\infty$}
			\end{tikzpicture}
		\end{center}
		Vậy $m=2$ không thỏa mãn.
	}
\end{bt}

\begin{bt}%[2D1V2-4]
% nguon: Archive/2024/TaiLieuGiangDay/1.HamSo/1. K12 - 4in1 - Hamso/ThubaiGoc/Bai2-Dang4.tex #2496
[Sách tham khảo, Mức V]%[Dự án bài giảng 12, Trung Anh]
	Có tất cả bao nhiêu số nguyên dương $ m $ để đồ thị hàm số $ y=\dfrac{1}{3}x^3-5x^2+(m^2+m-6)x+3 $ có $ 2 $ điểm cực trị nằm khác phía so với trục tung?
	\loigiai{
		Ta có $ y'=x^2-10x+m^2+m-6 $.\\
		Đồ thị hàm số có hai điểm cực trị nằm khác phía so với trục tung khi và chỉ khi phương trình $y'=0$ có hai nghiệm trái dấu. Tương đương
		$$m^2+m-6 <0 \Leftrightarrow -3<m<2.$$
		Do $ m $ nguyên dương nên $ m=1 $. Vậy có một giá trị nguyên dương $ m $ thỏa mãn yêu cầu bài toán.	
	}
\end{bt}

\begin{bt}%[2D1N4-4]
% nguon: Archive/2024/TaiLieuGiangDay/1.HamSo/1. K12 - 4in1 - Hamso/ThubaiGoc/Bai4-Dang6.tex #2507
[SGK12-CTST, Mức độ 1]%[Bài giảng 12, Phạm Tiến Long]
Theo thuyết tương đối hẹp, khối lượng m (kg) của một hạt phụ thuộc vào tốc độ di chuyển $v(\mathrm{~km} / \mathrm{s})$ của nó trong hệ quy chiếu quán tính theo công thức $m(v)=\dfrac{m_0}{\sqrt{1-\dfrac{v^2}{c^2}}}$, trong đó $m_0$
là khối lượng nghỉ của hạt, $c=300000 \mathrm{~km} / \mathrm{s}$ là tốc độ ánh sáng. 
	Tìm tiệm cận của đồ thị hàm số khối lượng hạt.
	\loigiai{
		Hàm số $y=m(v)=\dfrac{m_{0}}{\sqrt{1-\dfrac{v^{2}}{c^{2}}}}$ có tập xác định $\mathscr{D}=\left[0;c\right)$.\\
		Ta có $\lim\limits_{x \rightarrow c^-} \dfrac{m_{0}}{\sqrt{1-\dfrac{v^{2}}{c^{2}}}}=+\infty$.\\
		Vậy đồ thị hàm số có tiệm cận đứng là đường thẳng $x=c=300\,000$.	
	}
\end{bt}

\begin{bt}%[2D1N4-4]
% nguon: Archive/2024/TaiLieuGiangDay/1.HamSo/1. K12 - 4in1 - Hamso/ThubaiGoc/Bai4-Dang6.tex #2511
[SGK12-CTST, Mức độ 1]%[Bài giảng 12, Phạm Tiến Long]
	Nồng độ oxygen trong hồ theo thời gian $t$ cho bởi công thức $y(t)=5-\dfrac{15t}{9t^{2}+1}$, với $y$ được tính theo mg/l và $t$ được tính theo giờ, $t \geq 0$. Tìm các đường tiệm cận của đồ thị hàm số $y(t)$. Từ đó, có nhận xét gì về nồng độ oxygen trong hồ khi thời gian $t$ trở nên rất lớn?
	\begin{flushright}
		%	(Theo: \href{www.researchgate.net/publication/264903978_Microrespirometric_characterization_of_activated_sludge_inhibition_by_copper_and_zinc}{www.researchgate.net/publication/264903978\_Microrespirometric\_characterization\_of\_activated\_sludge\_inhibition\_by\_copper\_and\_zinc})
		(\textit{Theo}: www.researchgate.net/publication/264903978\_Microrespirometric\_\\
		characterization\_of\_activated\_sludge\_inhibition\_by\_copper\_and\_zinc)
	\end{flushright}
	\loigiai{
		Hàm số $y(t)=5-\dfrac{15t}{9t^{2}+1}$ có tập xác định $\mathscr{D}=\mathbb{R}$.\\
		Ta có $\lim\limits_{x \rightarrow+\infty} \left(5-\dfrac{15t}{9t^{2}+1}\right)=5$.\\
		Vậy đồ thị hàm số có tiệm cận ngang là đường thẳng $y=5$.\\
		Khi thời gian $t$ trở nên rất lớn thì nồng độ oxygen trong hồ sẽ tiến dần về $5$ mg/l.
	}
\end{bt}

\begin{bt}%[2D1N4-4]
% nguon: Archive/2024/TaiLieuGiangDay/1.HamSo/1. K12 - 4in1 - Hamso/ThubaiGoc/Bai4-Dang6.tex #2512
[SGK12-CKP, Mức độ 2]%[Bài giảng 12, Phạm Tiến Long]
	\immini{
		Tại một công ty sản xuất đồ chơi $A$, công ty phải chi $50\,000$ USD để thiết lập dây chuyền sản xuất ban đầu. Sau đó, cứ sản xuất được một sản phẩm đồ chơi $A$, công ty phải trả $5$ USD cho nguyên liệu thô và nhân công. Gọi $x$ ($x\geq 1$) là số đồ chơi $A$ mà công ty đã sản xuất và $T(x)$ (đơn vị USD) là tổng tiền bao gồm cả chi phí ban đầu mà công ty phải chi trả khi sản xuất $x$ đồ chơi $A$. Người ta xác định chi phí trung bình cho mỗi sản phẩm đồ chơi $A$ là $M(x)=\dfrac{T(x)}{x}$.
		\begin{enumerate}
			\item Xem $M(x)$ là hàm số theo $x$ xác định trên nửa khoảng $[1;+\infty)$, tìm tiệm cận ngang của đồ thị hàm số này.
			\item Nêu nhận xét về chi phí trung bình cho mỗi sản phẩm đồ chơi $A$ khi $x$ đủ lớn.
		\end{enumerate}
	}{
		\begin{tikzpicture}[>=latex,line join=round, line cap=round, scale=0.8]
			\draw[->] (-3,0)--(3,0) node[above]{$x$};
			\draw[->] (0,-2)--(0,6) node[right]{$y$};
			\clip (-3,-2) rectangle (3,6);
			\draw[samples=200,smooth,blue,line width=1] plot[domain=-3:-0.51] (\x,{2-2*(\x)+2/(4*(\x)^2-1))});
			\draw[samples=200,smooth,blue,line width=1] plot[domain=-0.49:0.49] (\x,{2-2*(\x)+2/(4*(\x)^2-1))});
			\draw[samples=200,smooth,blue,line width=1] plot[domain=0.51:3] (\x,{2-2*(\x)+2/(4*(\x)^2-1))});
			\draw[samples=200,smooth,red] plot[domain=-3:3] (\x,{2-2*(\x)});
			\draw[red] (-0.5,-2)--(-0.5,6) (0.5,-2)--(0.5,6);
			\draw (0,0) node[above right]{$O$} (-2,0) node[below]{$-2$} (-1,0) node[below]{$-1$} (1,0) node[below]{$1$} (2,0) node[below]{$2$} (0,-1) node[left]{$-1$} (0,1) node[left]{$1$} (0,2) node[left]{$2$} (0,3) node[left]{$3$} (0,4) node[left]{$4$} (0,5) node[left]{$5$};
		\end{tikzpicture}
	}
	\loigiai{
		\begin{enumerate}
			\item Chi phí bỏ ra để sản xuất đồ chơi là $T(x)=50\,000+5x$ (USD).\\
			Chi phí trung bình cho mỗi sản phẩm là $M(x)=\dfrac{T(x)}{x}=\dfrac{50\,000+5x}{x}$ (với $x\geq 1$).\\
			Ta có $\lim\limits_{x\to \pm\infty}M(x)=5$. Suy ra đường thẳng $y=5$ là tiệm cận ngang của đồ thị.
			\item Khi số đồ chơi $A$ được sản xuất với số lượng $x$ đủ lớn thì chi phí trung bình tiệm cận với tổng của giá nguyên liệu thô và nhân công.
		\end{enumerate}
	}
\end{bt}

\begin{bt}%[2D1N2-1]
% nguon: Archive/2024/TaiLieuGiangDay/1.HamSo/1. K12 - 4in1 - Hamso/ThubaiGoc/Bai2-Dang1.tex #2517
[SGK12-CKP, Mức độ 1] 
Tìm cực trị của hàm số $y=f\left(x\right)=\dfrac{1}{3}x^3-x^2-3x+\dfrac{1}{3}$.
	\loigiai{Tập xác định $\mathscr{D}=\mathbb{R}$.\\
		Ta có $y'=x^2-2x-3$;\\
		$y'=0\Leftrightarrow \hoac{x=-1&\\x=3.}$\\
		Bảng biến thiên
		\begin{center}
			\begin{tikzpicture}
				\tkzTabInit[lgt=1.2,espcl=3]
				{$x$ /.8, $f’(x)$ /.8, $f(x)$ /2}
				{$-\infty$,$-1$,$3$,$+\infty$}
				\tkzTabLine{ ,+,z,-,z,+, }
				\tkzTabVar{-/$-\infty$,+/$2$,-/$-\dfrac{26}{3}$,+/$+\infty$}
			\end{tikzpicture}
		\end{center}
		Từ bảng biến thiên, suy ra $x=-1$ là điểm cực đại của hàm số, $f_{\text{CĐ}}=2$; $x=3$ là điểm cực tiểu của hàm số, $f_{\text{CT}}=\dfrac{-26}{3}$.
	}
	
\end{bt}

\begin{bt}%[2D1N2-1]
% nguon: Archive/2024/TaiLieuGiangDay/1.HamSo/1. K12 - 4in1 - Hamso/ThubaiGoc/Bai2-Dang1.tex #2518
[SGK12-CTST, Mức độ 1]
	Tìm cực trị của hàm số $f(x)=x^3-3x^2+3x-4$.
	\loigiai{
		Tập xác định $\mathscr{D}=\mathbb{R}$.\\
		Ta có $f'(x)=3x^2-6x+3$; $f'(x)=0 \Leftrightarrow x=1$.\\
		Bảng biến thiên
		\begin{center}
			\begin{tikzpicture}
				\tkzTabInit[nocadre=false,lgt=1.2,espcl=4,deltacl=0.6]
				{$x$ /0.6,$f'(x)$ /0.6,$f(x)$ /2}
				{$-\infty$,$1$,$+\infty$}
				\tkzTabLine{,+,$0$,+,}
				\tkzTabVar{-/$-\infty$,R,+/$+\infty$}
			\end{tikzpicture}
		\end{center}
		Vậy hàm số không có cực trị.
	}
\end{bt}

\begin{bt}%[2D1H2-1]
% nguon: Archive/2024/TaiLieuGiangDay/1.HamSo/1. K12 - 4in1 - Hamso/ThubaiGoc/Bai2-Dang1.tex #2519
[SGK12-CKP, Mức độ 2]
Tìm cực trị của hàm số $y=f\left(x\right)=\dfrac{3x+1}{x-2}$.
	\loigiai{Hàm số xác định với mọi $x\in \mathbb{R}\backslash \left\{2\right\}$.\\
		Ta có $y'=\dfrac{-7}{\left(x-2\right)^2}<0$ với mọi $x\in \mathbb{R}\backslash \left\{2\right\}$.\\
		Bảng biến thiên
		\begin{center}
			\begin{tikzpicture}
				\tkzTabInit[lgt=1.2,espcl=6]
				{$x$ /.6, $f’(x)$ /.6, $f(x)$ /1.5}
				{$-\infty$,$2$,$+\infty$}
				\tkzTabLine{ ,-,d,-,}
				\tkzTabVar{+/$3$,-D+/$-\infty$/ $+\infty$,-/$3$}
			\end{tikzpicture}
		\end{center}
		Hàm số không có cực trị.
	}
\end{bt}

\begin{bt}%[2D1H2-1]
% nguon: Archive/2024/TaiLieuGiangDay/1.HamSo/1. K12 - 4in1 - Hamso/ThubaiGoc/Bai2-Dang1.tex #2520
[SGK12-CTST, Mức độ 2]
	Tìm cực trị của hàm số $g(x)=\dfrac{x^2+x+4}{x+1}$.
	\loigiai{
		Tập xác định $\mathscr{D}=\mathbb{R}\setminus \{-1\}$.\\
		Ta có $g(x)=x+\dfrac{4}{x+1} \Rightarrow g'(x)=1-\dfrac{4}{(x+1)^2}=\dfrac{x^2+2x-3}{(x+1)^2}$;\\
		$g'(x)=0 \Leftrightarrow x^2+2x-3=0 \Leftrightarrow \hoac{& x=-3\\& x=1.}$\\
		Bảng biến thiên
		\begin{center}
			\begin{tikzpicture}
				\tkzTabInit[nocadre=false,lgt=1.2,espcl=2.5,deltacl=0.6]
				{$x$ /0.6,$g'(x)$ /0.6,$g(x)$ /2}
				{$-\infty$,$-3$,$-1$,$1$,$+\infty$}
				\tkzTabLine{,+,$0$,-,d,-,$0$,+,}
				\tkzTabVar{-/$-\infty$,+/$-5$,-D+/$-\infty$/$+\infty$,-/$3$,+/$+\infty$}
			\end{tikzpicture}
		\end{center}
		Vậy hàm số đạt cực đại tại $x=-3$, $y_{\text{CĐ}}=g(-3)=-5$; và hàm số đạt cực tiểu tại $x=1$, $y_{\text{CT}}=g(1)=3$.
	}
\end{bt}

\begin{bt}%[2D1H2-1]
% nguon: Archive/2024/TaiLieuGiangDay/1.HamSo/1. K12 - 4in1 - Hamso/ThubaiGoc/Bai2-Dang1.tex #2521
[SGK12-KNTT, Mức độ 2]
	Tìm cực trị của hàm số sau $y=x^{4}-3 x^{2}+1$.
	\loigiai{
		$y=x^4-3 x^2+1$; \\
		Tập xác định $\mathscr{D}=\mathbb{R}$.\\
		Ta có $y'=4x^3-6x$. \\
		Khi đó $y'=0\Leftrightarrow \hoac{&x=0\\&x=\pm\dfrac{\sqrt{6}}{2}.}$ \\
		Lập bảng biến thiên của hàm số
		\begin{center}
			\begin{tikzpicture}
				\tkzTabInit[lgt=1.2,espcl=2,deltacl=0.6]
				{$x$ /1, $y'$ /0.6, $y$/1.8}
				{$-\infty$,$-\dfrac{\sqrt{6}}{2}$, $0$, $\dfrac{\sqrt{6}}{2}$,$+\infty$}
				\tkzTabLine{,-,$0$,+,$0$,-,$0$,+,}
				\tkzTabVar{+/$+\infty$,-/$-\dfrac{5}{4}$,+/$1$,-/$-\dfrac{5}{4}$,+/$+\infty$}
			\end{tikzpicture}
		\end{center}
		Từ bảng biến thiên, ta có\\
		Hàm số đạt cực tiểu tại $x=\pm\dfrac{\sqrt{6}}{2}$  và  $y_{\text{CT}}=y\left(\pm\dfrac{\sqrt{6}}{2}\right)=-\dfrac{5}{4}$.\\
		Hàm số đạt cực đại tại $x=0$ và $y_{\text{CĐ}}=y(0)=1$.
	}
\end{bt}

\begin{bt}%[2D1H2-1]
% nguon: Archive/2024/TaiLieuGiangDay/1.HamSo/1. K12 - 4in1 - Hamso/ThubaiGoc/Bai2-Dang1.tex #2522
[Mức độ 2]
	Tìm cực trị của hàm số $y=x\cdot \mathrm{e}^x$. 
	\loigiai{
	Tập xác định: $\mathscr{D}=\mathbb{R}$.\\
	Ta có
	\[y'=\mathrm{e}^x+x \cdot\mathrm{e}^x \Rightarrow y'=0 \Leftrightarrow\mathrm{e}^x+x \cdot\mathrm{e}^x=0 \Leftrightarrow\mathrm{e}^x(1+x)=0 \Leftrightarrow x=-1.\]
%	Giới hạn: $\lim\limits_{x\to-\infty} y=0$ và $\lim\limits_{x\to+\infty} y=+\infty$
	Bảng biến thiên
	\begin{center}
		\begin{tikzpicture}[scale=1, font=\footnotesize]%<DTools>
			\tkzTabInit[nocadre=false, lgt=1.2, espcl=2, deltacl=0.6]
			{$x$/0.8,$f'(x)$/0.6,$f(x)$/2}
			{$-\infty$,$-1$,$+\infty$};
			\tkzTabLine{,-,$0$,+,};
			\tkzTabVar{+/$+\infty$,-/$-\dfrac{1}{\mathrm{e}}$,+/$+\infty$};
		\end{tikzpicture}
	\end{center}
	Từ bảng biến thiên ta có\\
	Hàm số đạt cực tiểu tại $x=-1$ và giá trị cực tiểu $y_\text{CT}=-\dfrac{1}{\mathrm{e}}$.
	}
	
\end{bt}

\begin{bt}%[2D1H2-1]
% nguon: Archive/2024/TaiLieuGiangDay/1.HamSo/1. K12 - 4in1 - Hamso/ThubaiGoc/Bai2-Dang1.tex #2523
[Mức độ 2]
	Tìm cực trị của hàm số $y=x-\ln x$.
	\loigiai{
	Tập xác định $\mathscr{D}=(0,+\infty)$.\\
	Ta có
	\[y'=1-\dfrac{1}{x} \Rightarrow y'=0 \Leftrightarrow x=1.\]
	Bảng biến thiên 
	\begin{center}
		\begin{tikzpicture}[>=stealth]
			%Khai báo thông số các dòng
			\tkzTabInit[nocadre=false,lgt=1.5,espcl=3.5,deltacl=0.5]{$x$/.6,$h'(x)$/.6,$h(x)$/2.5}
			%Dòng 2
			{$-\infty$,$0$,$1$,$+\infty$}
			%Dòng 2 dấu
			\tkzTabLine{ ,  ,d,-,$0$,+, }
			% Các vị trí
			\path
			(N13)node[above](1){}
			(N22)node[below=.3cm](2){$+\infty$}
			(N33)node[above](3){$1$}
			(N42)node[below](4){$+\infty$}
			;
			\fill[pattern = north east lines] (N13) rectangle (N20)
			;
			\draw[double] (N22)--(N23);
			%Vẽ mũi tên
			\foreach \x/\y in {2/3,3/4}
			\draw[->] (\x)--(\y);
		\end{tikzpicture}
	\end{center}
Từ bảng biến thiên ta có\\
Hàm số đạt cực tiểu tại $x=1$ và giá trị cực tiểu là $y_\text{CT}=1$.
	}
	
\end{bt}

\begin{bt}%[2D1H2-1]
% nguon: Archive/2024/TaiLieuGiangDay/1.HamSo/1. K12 - 4in1 - Hamso/ThubaiGoc/Bai2-Dang1.tex #2524
[SGK12-CKP, Mức độ 2]
Tìm cực trị của hàm số sau 
$y=\left|x\right|$.
\loigiai{
\textbf{\textit{Cách 1:	}}
Hàm số $y=\left|x\right|$ xác định với mọi $x\in\mathbb{R}$.\\
Khi $x>0$, ta có $y=x$ và $y'=1$.\\
Khi $x<0$, ta có $y=-x$ và $y'=-1$.\\
Tại $x = 0$, ta có giá trị nhỏ nhất của hàm số là $y = 0$.\\ Vì vậy, ta có cực trị tại điểm $x = 0$, với giá trị cực trị là $y = 0$.\\
\textbf{\textit{Cách 2: Tìm cực trị bằng định nghĩa.}}\\
Hàm số $y=|x|$ xác định là liên tục trên khoảng $\left(-\infty;+\infty\right)$. \\
Xét điểm $x_0=0 \in \left(-\infty;+\infty\right)$.
Ta xét $h=1$ khi đó
$f(x)>f(0)$ với mọi $x\in \left(0-1;0+1\right) \subset \left(-\infty;+\infty\right)$ và $x\ne 0$.\\
Vậy hàm số $f(x)$ đạt cực đại tại $x_0$.

}
\end{bt}

\begin{bt}%[2D1H2-1]
% nguon: Archive/2024/TaiLieuGiangDay/1.HamSo/1. K12 - 4in1 - Hamso/ThubaiGoc/Bai2-Dang1.tex #2525
[Mức độ 2]
	Cho hàm số $y=2x^3-3x^2-12x+1$. Viết phương trình đường thẳng qua hai điểm cực trị của hàm số trên.
	\loigiai{
		Tập xác định $\mathscr{D}=\mathbb{R}$.\\
		Tính $y'=6x^2-6x-12$.\\
		Cho $y'=0\Leftrightarrow\hoac{&x=-1\\&x=2.}$\\
		Bảng biến thiên 
		\begin{center}
			\begin{tikzpicture}[scale=1, font=\footnotesize]%<DTools>
				\tkzTabInit[nocadre=false, lgt=1.2, espcl=2, deltacl=0.6]
				{$x$/0.8,$f'(x)$/0.6,$f(x)$/2}
				{$-\infty$,$-1$,$2$,$+\infty$};
				\tkzTabLine{,+,$0$,-,$0$,+,};
				\tkzTabVar{-/$-\infty$,+/$8$,-/$-19$,+/$+\infty$};
			\end{tikzpicture}
		\end{center}
		Suy ra tọa độ hai điểm cực trị $(-1;8)$, $(2;-19)$.\\
		Vậy phương trình đường thẳng đi qua hai điểm cực trị là $9x+y+1=0$.
	}
\end{bt}

\begin{bt}%[2D1H2-1]
% nguon: Archive/2024/TaiLieuGiangDay/1.HamSo/1. K12 - 4in1 - Hamso/ThubaiGoc/Bai2-Dang1.tex #2529
[SGK12-CD, Mức độ 2]
	Tìm cực trị của hàm số sau $y=x^4-6 x^2+8 x+1$.
	\loigiai{
		$y=x^4-6 x^2+8 x+1$; \\
		Tập xác định $\mathscr{D}=\mathbb{R}$.\\
		Ta có $y'=4x^3-12x+8=4(x-1)^2(x+2)$. \\
		Khi đó $y'=0\Leftrightarrow \hoac{&x=-2\\&x=1.}$ \\
		Lập bảng biến thiên của hàm số
		\begin{center}
			\begin{tikzpicture}[scale=1, font=\footnotesize]%<DTools>
				\tkzTabInit[nocadre=false, lgt=1.2, espcl=2, deltacl=0.6]
				{$x$/0.8,$f'(x)$/0.6,$f(x)$/2}
				{$-\infty$,$-2$,$1$,$+\infty$};
				\tkzTabLine{,-,$0$,+,$0$,+,};
				\tkzTabVar{+/$+\infty$,-/$-23$,R/,+/$+\infty$};
				\tkzTabIma{2}{4}{3}{$4$};
			\end{tikzpicture}
		\end{center}
		Từ bảng biến thiên, ta có hàm số đạt cực tiểu tại $x=-2$  và  $y_{\text{CT}}=y\left(-2\right)=-23$.
	}
\end{bt}

\begin{bt}%[2D1H2-1]
% nguon: Archive/2024/TaiLieuGiangDay/1.HamSo/1. K12 - 4in1 - Hamso/ThubaiGoc/Bai2-Dang1.tex #2530
[SGK12-CKP, Mức độ 2]
	Tìm các điểm cực trị của hàm số sau: $y=3^{x-2 x^2}$.
	\loigiai{
		Hàm số $y=3^{x-2x^{2}}$ xác định với mọi $x\in\mathbb{R}$.\\
		Ta có $y'=3^{x-2x^{2}}\cdot\ln3\cdot\left(1-4x\right)$.\\
		$y'=0\Leftrightarrow x=\dfrac{1}{4}$.\\
		Bảng biến thiên
		\begin{center}
			\begin{tikzpicture}
				\tkzTabInit	
				{$x$/.7,$y'$/.7,$y$/2}
				{$-\infty$,$\tfrac{1}{4}$,$+\infty$}
				\tkzTabLine{,+,0,-,} %
				\tkzTabVar{-/$0$,+/$3^{\frac{1}{8}}$,-/$0$} 
			\end{tikzpicture}
		\end{center}
		Từ bảng biến thiên, suy ra $x=\dfrac{1}{4}$ là điểm cực đại của hàm số, $y_{\text{CĐ}}=3^{\frac{1}{8}}$.
	}
\end{bt}

\begin{bt}%[2D1H2-1]
% nguon: Archive/2024/TaiLieuGiangDay/1.HamSo/1. K12 - 4in1 - Hamso/ThubaiGoc/Bai2-Dang1.tex #2531
[SGK12-CKP, Mức độ 2]
	Tìm các điểm cực trị của hàm số sau $y=\ln \left(x^2+\mathrm{e}\right)$.
	\loigiai{
		Hàm số $y=\ln \left(x^2+\mathrm{e}\right)$ xác định với mọi $x\in\mathbb{R}$.\\
		Ta có $y'=\dfrac{2x}{x^{2}+\mathrm{e}}$.\\
		$y'=0\Leftrightarrow x=0$.\\
		Bảng biến thiên
		\begin{center}
			\begin{tikzpicture}
				\tkzTabInit	
				{$x$/.7,$y'$/.7,$y$/2}
				{$-\infty$,$0$,$+\infty$}
				\tkzTabLine{,-,0,+,} %
				\tkzTabVar{+/$+\infty$,-/$1$,+/$+\infty$} 
			\end{tikzpicture}
		\end{center}
		Từ bảng biến thiên, suy ra $x=0$ là điểm cực tiểu của hàm số, $y_{\text{CT}}=1$.
		}
\end{bt}

\begin{bt}%[2D1H2-1]
% nguon: Archive/2024/TaiLieuGiangDay/1.HamSo/1. K12 - 4in1 - Hamso/ThubaiGoc/Bai2-Dang1.tex #2532
[SGK12-CTST, Mức độ 2]
	Tìm cực trị của hàm số sau $y=\sqrt{-x^2+4}$.
	\loigiai{
		Tập xác định $\mathscr{D}=[-2;2]$.\\
		Ta có $y'=\dfrac{-x}{\sqrt{-x^2+4}}$, $y'=0 \Leftrightarrow x=0$.\\
		Bảng biến thiên
		\begin{center}
			\begin{tikzpicture}
				\tkzTabInit[nocadre=false,lgt=1.2,espcl=4,deltacl=0.6]
				{$x$ /0.6,$y'$ /0.6,$y$ /2}
				{$-2$,$0$,$2$}
				\tkzTabLine{,+,$0$,-,}
				\tkzTabVar{-/$0$,+/$2$,-/$0$}
			\end{tikzpicture}
		\end{center}
		Vậy hàm số đạt cực đại tại $x=0$, $y_{\text{CĐ}}=2$; hàm số không có cực tiểu.
		
	}
\end{bt}

\begin{bt}%[2D1V2-1]
% nguon: Archive/2024/TaiLieuGiangDay/1.HamSo/1. K12 - 4in1 - Hamso/ThubaiGoc/Bai2-Dang1.tex #2533

	[Đề tốt nghiệp 2020 Đợt 2 Mã đề103, Mức độ 2] 
	Cho hàm số $f(x)$ có đạo hàm \[f'(x)=x(x+1)(x-4)^3,\forall x\in \mathbb{R}.\]
	Tìm số điểm cực đại của hàm số đã cho.
	\loigiai{
	$f'(x)=0 \Leftrightarrow $x$ (x+1)(x-4)^3=0 \Leftrightarrow\hoac{&x=0\\&x=-1\\&x=4.}$\\
	Lập bảng biến thiên của hàm số $f(x)$
	\begin{center}
		\begin{tikzpicture}[scale=1, font=\footnotesize]%<DTools>
			\tkzTabInit[nocadre=false, lgt=1.2, espcl=2, deltacl=0.6]
			{$x$/0.8,$f'(x)$/0.6,$f(x)$/2}
			{$-\infty$,$-1$,$0$,$4$,$+\infty$};
			\tkzTabLine{,-,$0$,+,$0$,-,$0$,+,};
			\tkzTabVar{+/$+\infty$,-/$f(-1)$,+/$f(0)$,-/$f(4)$,+/$+\infty$};
		\end{tikzpicture}
	\end{center}
	Vậy hàm số đã cho có một điểm cực đại.
}

	
\end{bt}

\begin{bt}%[2D1N2-3]
% nguon: Archive/2024/TaiLieuGiangDay/1.HamSo/1. K12 - 4in1 - Hamso/ThubaiGoc/Bai2-Dang3.tex #2544
%[BG12, Tran Tony]
	Tìm giá trị của tham số $m$ để hàm số $y=-x^3+x^2+mx$ có hai điểm cực trị.
	\loigiai{
		Ta có $y'=-3x^2+2x+m$.\\
		Hàm số có hai điểm cực trị khi và chỉ khi phương trình $y'=0$ có hai nghiệm phân biệt $$\Leftrightarrow \Delta'=1+3m>0\Leftrightarrow m>-\dfrac{1}{3}.$$
		Vậy $m>-\dfrac{-1}{3}$ thì hàm số đã cho có hai điểm cực trị.
	}
\end{bt}

\begin{bt}%[2D1N2-3]
% nguon: Archive/2024/TaiLieuGiangDay/1.HamSo/1. K12 - 4in1 - Hamso/ThubaiGoc/Bai2-Dang3.tex #2545
%[BG12, Tran Tony]
	Tìm giá trị của tham số $m$ để hàm số $y=2^x+m$ đạt cực đại tại điểm $x=1$.
	\loigiai{
		Ta có $y'=2^2\ln 2>0,\, x\in \mathbb{R}$.\\
		Vậy không có giá trị nào của $m$ để hàm số đạt cực đại tại điểm $x=1$.
	}
\end{bt}

\begin{bt}%[2D1H2-3]
% nguon: Archive/2024/TaiLieuGiangDay/1.HamSo/1. K12 - 4in1 - Hamso/ThubaiGoc/Bai2-Dang3.tex #2546
%[BG12, Tran Tony]
	Tìm giá trị của tham số $m$ để hàm số $y=\dfrac{1}{3}x^3-mx^2+\left(m^2-4\right)x+5$ đạt cực tiểu tại điểm $x=1$.
	\loigiai{
		Ta có $y'=x^2-2mx+m^2-4$.\\
		Hàm số đạt cực tiểu tại $x=1$ thì điều kiện cần là $y'(1)=0\Leftrightarrow 1-2m+m^2-4=0\Leftrightarrow\hoac{& m=3\\ & m=-1.}$\\
		Điều kiện đủ
		\begin{itemize}
			\item Với $m=3$, ta có $y'=x^2-6x+5$, $y'=0\Leftrightarrow\hoac{&x=1\\&x=5.}$\\
			Bảng xét dấu đạo hàm
				\begin{center}
			\begin{tikzpicture}
			\tkzTabInit[lgt=1.2,espcl=3]
			{$x$ /0.6, $y'$ /0.6}
			{$-\infty$,$1$,$5$,$+\infty$}
			\tkzTabLine{ ,+,z,-,z, +, }
			\end{tikzpicture}
			\end{center}
			Từ bảng xét dấu suy ra, hàm số đạt cực đại tại $x=1$. Do đó $m=3$ không thỏa mãn.
			\item Với $m=-1$, ta có $y'=x^2+2x-3$, $y'=0\Leftrightarrow\hoac{&x=1\\&x=-3.}$\\
			Bảng xét dấu đạo hàm
			\begin{center}
				\begin{tikzpicture}
					\tkzTabInit[lgt=1.2,espcl=3]
					{$x$ /0.6, $y'$ /0.6}
					{$-\infty$,$-3$,$1$,$+\infty$}
					\tkzTabLine{ ,+,z,-,z, +, }
				\end{tikzpicture}
			\end{center}
			Từ bảng xét dấu suy ra, hàm số đạt cực tiểu tại $x=1$. Do đó $m=-1$ thỏa mãn.
		\end{itemize}
		Vậy $m=-1$ thì hàm số đã cho đạt cực tiểu tại $x=1$.
	}
\end{bt}

\begin{bt}%[2D1H2-3]
% nguon: Archive/2024/TaiLieuGiangDay/1.HamSo/1. K12 - 4in1 - Hamso/ThubaiGoc/Bai2-Dang3.tex #2547
%[BG12, Tran Tony]
	Có bao nhiêu số nguyên dương $m$ để hàm số $y=(3-m)x^{100}$ đạt cực tiểu tại $x=0$.
	\loigiai
	{Giả sử hàm số $f(x)=(3-m)x^{100}$ đạt cực tiểu tại $x=0$. Khi đó tồn tại khoảng $(a;b)$ chứa $0$ sao cho $$f(x)>f(0),\ \forall x\in (a;b)\setminus\{0\}\Leftrightarrow (3-m)x^{100}>0,\ \forall x\in (a;b)\setminus\{0\}\Leftrightarrow 3-m>0\Leftrightarrow m<3.$$
		Vậy có $2$ số nguyên dương $m$ thỏa mãn bài toán.
	}
\end{bt}

\begin{bt}%[2D1H2-3]
% nguon: Archive/2024/TaiLieuGiangDay/1.HamSo/1. K12 - 4in1 - Hamso/ThubaiGoc/Bai2-Dang3.tex #2548
%[BG12, Tran Tony]
	Cho hàm số $y=f(x)=x^3+ax^2+bx+c$ (với $a$, $b$, $c$ là các số thực) có cực tiểu bằng $-3$ tại điểm $x=1$ và đồ thị hàm số cắt trục tung tại điểm có tung độ là $2$. Tính $H=3a+b-c$.
	\loigiai{
		Đồ thị hàm số cắt trục tung tại điểm có tung độ bằng $2$, ta được $ f(0) = 2 \Rightarrow c = 2 $.\\
		Hàm số đạt cực tiểu bằng $-3$, tại $ x = 1 $, ta được $$\heva{& f'(1) = 0\\ & f(1) = -3} \Rightarrow \heva{& 2a + b + 3 = 0\\ & a + b + 6 = 0} \Leftrightarrow \heva{& a = 3\\ & b = -9.}$$ 
		Khi đó $y'=3x^2+2ax+b\Rightarrow y'(1)=3+2a+b=0$ và $-3=y(1)=1+a+b+2$.\\
		Vậy $H=3\cdot 3-9-2=-2$.
	}
\end{bt}

\begin{bt}%[2D1H2-3]
% nguon: Archive/2024/TaiLieuGiangDay/1.HamSo/1. K12 - 4in1 - Hamso/ThubaiGoc/Bai2-Dang3.tex #2549
%[BG12, Tran Tony]
	Biết rằng đồ thị hàm số $y = x^4 - 2ax^2 + b$ có một điểm cực trị là $(1;2)$. Tính khoảng cách giữa điểm cực đại và điểm cực tiểu của hàm số đã cho.
	\loigiai{
		Ta có $y'= 4x^3 - 4ax$.\\
		Theo đề ta có $\heva{&y'(1) = 0\\&y(1) = 2} \Leftrightarrow \heva{&4-4a=0\\&1-2a+b=2}\Leftrightarrow \heva{&a = 1\\&b=3.}$\\
		Khi đó hàm số trở thành $y=x^4-2x^2 + 3$ có hai điểm cực tiểu là $(1;2)$; $(-1;2)$ và một điểm cực đại là $(0;3)$.\\
		Khoảng cách cần tìm là $\mathrm{d} = \sqrt{(1-0)^2+(3-2)^2} = \sqrt{2}$.
	}
\end{bt}

\begin{bt}%[2D1H2-3]
% nguon: Archive/2024/TaiLieuGiangDay/1.HamSo/1. K12 - 4in1 - Hamso/ThubaiGoc/Bai2-Dang3.tex #2550
%[BG12, Tran Tony]
	Tìm giá trị của tham số $m$ để hàm số $y=\dfrac{x^2-x+m}{x-1}$ đạt cực tiểu tại điểm $x=2$.
	\loigiai{
		Ta có $y=\dfrac{x^2-x+m}{x-1}=x+\dfrac{m}{x-1}\Rightarrow y'=1-\dfrac{m}{(x-1)^2}$.\\
		Hàm số đạt cực tiểu tại $x=2$ thì điều kiện cần là $y'(2)=0\Leftrightarrow 1-\dfrac{m}{(2-1)^2}=0\Leftrightarrow m=1$.\\
		Điều kiện đủ, với $m=1$, ta có $y'=1-\dfrac{1}{(x-1)^2}$, $y'=0\Leftrightarrow\hoac{&x=0\\&x=2.}$\\
		Bảng biến thiên
		\begin{center}
		\begin{tikzpicture}
		\tkzTabInit[lgt=1.2,espcl=3]
		{$x$/0.6,$y'$/0.6,$y$/2.5}
		{$-\infty$,$0$,$1$,$2$,$+\infty$}
		\tkzTabLine{ ,+,z,-,d,-,z,+, }
		\tkzTabVar{-/$-\infty$,+/$y(0)$,-D+/$-\infty$/$+\infty$,-/$y(2)$,+/$+\infty$}
		\end{tikzpicture}
		\end{center}
		Từ bảng biến thiên suy ra, hàm số đạt cực tiểu tại $x=2$. Do đó $m=1$ thỏa mãn.\\
		Vậy $m=1$ thì hàm số đã cho đạt cực tiểu tại $x=2$.
	}
\end{bt}

\begin{bt}%[2D1H2-3]
% nguon: Archive/2024/TaiLieuGiangDay/1.HamSo/1. K12 - 4in1 - Hamso/ThubaiGoc/Bai2-Dang3.tex #2554
%[BG12, Tran Tony]
	Cho hàm số $y=x^3+3mx^2+3\left(m^2-1\right)x+2$. Hãy tìm tất cả các giá trị của $m$ để hàm số có cực tiểu tại $x=2$.
	\loigiai
	{Ta có $y'=3x^2+6mx+3\left(m^2-1\right)$, với  $\forall x\in \mathbb{R}$.\\
		Hàm số đạt cực tiểu tại $x=2$ nên $y'(2)=0 \Leftrightarrow 12+12m+3m^2-3=0 \Leftrightarrow \hoac{& m=-1 \\ & m=-3.}$\\
		Với $m=-1$, ta có $y'=3x^2-6x$.\\
		Bảng xét dấu $y'$
		\begin{center}
			\begin{tikzpicture}
				\tkzTabInit[lgt=1.2,espcl=3]
				{$x$ /0.6, $y'$ /0.6}
				{$-\infty$,$0$,$2$,$+\infty$}
				\tkzTabLine{ ,+,z,-,z, +, }
			\end{tikzpicture}
		\end{center}
		Từ bảng xét dấu suy ra, hàm số đạt cực tiểu tại điểm $x=2$. Do đó $m=-1$ thỏa mãn.\\
		Với $m=-3$, ta có $y'=3x^2-18x+24x$.\\
		Bảng xét dấu $y'$
		\begin{center}
			\begin{tikzpicture}
				\tkzTabInit[lgt=1.2,espcl=3]
				{$x$ /0.6, $y'$ /0.6}
				{$-\infty$,$2$,$4$,$+\infty$}
				\tkzTabLine{ ,+,z,-,z, +, }
			\end{tikzpicture}
		\end{center}
		Từ bảng xét dấu suy ra, hàm số đạt cực đại tại điểm $x=2$. Do đó $m=-3$ không thỏa mãn.\\
		Vậy với $m=-1$ thì hàm số đạt cực tiểu tại $x=2$.
		}
\end{bt}

\begin{bt}%[2D1H2-3]
% nguon: Archive/2024/TaiLieuGiangDay/1.HamSo/1. K12 - 4in1 - Hamso/ThubaiGoc/Bai2-Dang3.tex #2555
%[BG12, Tran Tony]
	Cho hàm số $y=x^3-2x^2+ax+b$, $(a,b\in \mathbb{R})$ có đồ thị $\mathscr{(C)}$. Biết đồ thị $\mathscr{(C)}$ có điểm cực trị là $A(1;3)$. Tính giá trị của $P=4a-b$.
	\loigiai{
		$y=x^3-2x^2+ax+b\Rightarrow y^\prime =3x^2-4x+1$.\\
		Vì đồ thị $\mathscr{(C)}$ có điểm cực trị là $A(1;3)$ $\Rightarrow \heva{&y(1)=3\\&y' (1)=0}\Rightarrow \heva{&a+b-1=3\\&a-1=0}\Rightarrow \heva{a=1\\b=3}\\ \Rightarrow P=4a-b=1$.
	}
	
\end{bt}

\begin{bt}%[2D1H2-3]
% nguon: Archive/2024/TaiLieuGiangDay/1.HamSo/1. K12 - 4in1 - Hamso/ThubaiGoc/Bai2-Dang3.tex #2556
%[BG12, Tran Tony]
	Cho hàm số $y=f(x)=ax^3+bx^2+cx+d$ với $a\ne 0$. Biết đồ thị hàm số có hai điểm cực trị là $A(1;-1)$, $B(-1;3)$. Tính $f(4)$.
	\loigiai{Ta có $y'=3ax^2+2bx+c$. Đồ thị hàm số có hai điểm cực trị là $A(1;-1)$ và $B(-1;3)$ nên
		$$\heva{&y'(1)=0\\&y'(-1)=0\\&y(1)=-1\\&y(-1)=3}\Leftrightarrow\heva{&3a+2b+c=0\\&3a-2b+c=0\\&a+b+c+d=-1\\&-a+b-c+d=3}\Leftrightarrow\heva{&a=1\\&b=0\\&c=-3\\&d=1.}$$
		Suy ra $y=x^3-3x+1$, dễ thấy đồ thị hàm số nhận $A$ và $B$ làm hai điểm cực trị nên thỏa bài toán.\\
		Vậy $f(4)=53$.
	}
\end{bt}

\begin{bt}%[2D1H2-3]
% nguon: Archive/2024/TaiLieuGiangDay/1.HamSo/1. K12 - 4in1 - Hamso/ThubaiGoc/Bai2-Dang3.tex #2557
%[BG12, Tran Tony]
	Tìm tất cả các giá trị thực của tham số $m$ để hàm số $y=x^4+mx^2$ đạt cực tiểu tại $x=0$.
	\loigiai{
		Ta có $y'=4x^3+2mx=2x(2x^2+m)$; $y'=0\Leftrightarrow \hoac{& x=0 \\ & x^2=-\dfrac{m}{2}.}$
		\begin{itemize}
			\item Nếu $m=0$ thì $x=0$ là nghiệm bội ba của phương trình $y'=0$ và khi đó $x=0$ là điểm cực tiểu của hàm số.
			\item Nếu $m>0$ thì $x^2=-\dfrac{m}{2}<0$ (vô lí) do đó $x=0$ là nghiệm đơn của phương trình $y'=0$ và khi đó $x=0$ là điểm cực tiểu của hàm số.
			\item Nếu $m<0$ thì $x^2=-\dfrac{m}{2}>0\Rightarrow x=\pm\sqrt{-\dfrac{m}{2}}$ và  khi đó $x=0$ là điểm cực đại của hàm số (loại).
		\end{itemize}
		Vậy $m\geq 0$.
	}
\end{bt}

\begin{bt}%[2D1V2-3]
% nguon: Archive/2024/TaiLieuGiangDay/1.HamSo/1. K12 - 4in1 - Hamso/ThubaiGoc/Bai2-Dang3.tex #2558
%[BG12, Tran Tony]
	Biết hàm số $f(x)=x^3+ax^2+bx+c$ đạt cực tiểu tại $x=1$ và $f(1)=-3$, đồng thời đồ thị của hàm số cắt trục tung tại điểm có tung độ bằng $2$. Tính giá trị của $f(3)$.
	\loigiai{Ta có $f'(x)=3x^2+2ax+b$.\\
		Hàm số đạt cực tiểu tại $x=1$ suy ra $f'(1)=0\Leftrightarrow 3+2a+b=0\Leftrightarrow 2a+b=-3$.\\
		Hàm số cắt trục tung tại điểm có tung độ bằng $2$ suy ra $c=2$. \\
		Ngoài ra, $f(1)=-3\Leftrightarrow 1+a+b+c=-3\Leftrightarrow a+b=-6$.\\
		Khi đó,
		$$\heva{&2a+b=-3\\&a+b=-6}\Leftrightarrow \heva{&a=3\\&b=-9.}$$
		Suy ra $f(x)=x^3+3x^2-9x+2$.\\
		Vậy $f(3)=29$.
	}
\end{bt}

\begin{bt}%[2D1H4-2]
% nguon: Archive/2024/TaiLieuGiangDay/1.HamSo/1. K12 - 4in1 - Hamso/ThubaiGoc/Bai4-Dang4.tex #2569
%[Dự án TL12New-4in1-NCT]
	Đồ thị hàm số $y=\dfrac{ax+b}{2x+c}$ có tiệm cận ngang $y=2$ và tiệm cận đứng $x=1$. Tính $a+c$. 
	\loigiai{
		Xét hàm số $y=\dfrac{ax+b}{2x+c}$ có tập xác định $\mathscr{D}=\mathbb{R}\setminus\left\{-\dfrac{c}{2}\right\}$. \\
		Tiệm cận ngang $y=\dfrac{a}{2}$ và tiệm cận đứng $x=-\dfrac{c}{2}$.\\
		Theo đề bài, ta có $\heva{&\dfrac{a}{2}=2\\&-\dfrac{c}{2}=1}\Leftrightarrow\heva{&a=4\\&c=-2}\Rightarrow a+c=2$.}
\end{bt}

\begin{bt}%[2D1H4-2]
% nguon: Archive/2024/TaiLieuGiangDay/1.HamSo/1. K12 - 4in1 - Hamso/ThubaiGoc/Bai4-Dang4.tex #2570
%[Dự án TL12New-4in1-NCT]
	Tìm tất cả các giá trị của tham số $m$ để đồ thị hàm số $y=\dfrac{x-m}{x^2-3x+2}$ có đúng hai đường tiệm cận.
	\loigiai{
		Ta có $x^2-3x+2=0\Leftrightarrow\hoac{&x=1\\&x=2.}$\\
		Do $\lim\limits_{x\to\pm\infty} y=0$ nên đồ thị hàm số có một tiệm cận ngang là $y=0$.\\
		Đồ thị hàm số có đúng hai tiệm khi và chỉ khi nó có đúng một tiệm cận đứng, khi đó $m=1 $ hoặc $m=2$.}
\end{bt}

\begin{bt}%[2D1H4-2]
% nguon: Archive/2024/TaiLieuGiangDay/1.HamSo/1. K12 - 4in1 - Hamso/ThubaiGoc/Bai4-Dang4.tex #2571
%[Dự án TL12New-4in1-NCT]
	Cho hàm số $y=\dfrac{x+1}{x^2-2mx+4}$. Tìm tất cả các giá trị thực của tham số $m$ để đồ thị hàm số có ba đường tiệm cận. 
	\loigiai{
		Ta có $\lim\limits_{x\to+\infty} y=0$ suy ra  $y=0$ là tiệm cận ngang của đồ thị hàm số.\\
		Để đồ thị hàm số có $3$ tiệm cận thì phương trình $x^2-2mx+4=0$ có hai nghiệm phân biệt khác $-1$, hay
		\[\heva{&m^2-4>0\\&1+2m+4\neq 0}\Leftrightarrow\heva{&\hoac{&m>2\\&m <-2}\\&m\neq-\dfrac{5}{2}.}\]}
\end{bt}

\begin{bt}%[2D1V4-2]
% nguon: Archive/2024/TaiLieuGiangDay/1.HamSo/1. K12 - 4in1 - Hamso/ThubaiGoc/Bai4-Dang4.tex #2572
%[Dự án TL12New-4in1-NCT]
	Cho đồ thị hai hàm số $f(x)=\dfrac{2x+1}{x+1}$ và $g(x)=\dfrac{ax+1}{x+2}$ với $a\neq\dfrac{1}{2}$. Tìm tất cả các giá trị thực dương của $a$ để các tiệm cận của hai đồ thị hàm số tạo thành một hình chữ nhật có diện tích là $4$. 
	\loigiai{
		Đồ thị hàm số $f(x)=\dfrac{2x+1}{x+1}$ có hai đường tiệm cận là $x=-1$ và $y=2$.\\
		Đồ thị hàm số $g(x)=\dfrac{ax+1}{x+2}$ có hai đường tiệm cận là $x=-2$ và $y=a$.\\
		Hình chữ nhật được tạo thành từ bốn đường tiệm cận của hai đồ thị trên có hai kích thước là $1$ và $|a-2|$.\\
		Theo giả thiết, ta có $|a-2|\cdot 1=4\Leftrightarrow\hoac{&a=6\\&a=-2}$. Vì $a>0$ nên chọn $a=6$.}
\end{bt}

\begin{bt}%[2D1V4-2]
% nguon: Archive/2024/TaiLieuGiangDay/1.HamSo/1. K12 - 4in1 - Hamso/ThubaiGoc/Bai4-Dang4.tex #2577
%[Dự án TL12New-4in1-NCT]
Cho đường cong $(C)$:$y=\dfrac{m x^2+\left(3 m^2-2\right) x-2}{x+3 m}$.
Tìm $\mathrm{m}$ để góc giữa hai đường tiệm cận bằng $45^\circ$.
\loigiai{

Viết lại $(C)$ dưới dạng: $y=m x-2+\dfrac{6 m}{x+3 m}$ $\quad (1)$.\\
Từ (1) suy ra cần xét các khả năng sau:
\begin{itemize}
	\item Nếu $\mathrm{m}=\dfrac{1}{3}$ khi đó từ $(1)$ có $\mathrm{y}=\dfrac{1}{3} \mathrm{x}-2$ với $\mathrm{x} \neq 1$. Trường hợp này $(C)$ không có tiệm cận.
\item Nếu $ \mathrm{m} \neq \dfrac{1}{3}$ khi đó từ $(1)$ suy ra:
\begin{itemize}
\item $(C)$ có tiệm cận đứng $x=-3 m$
\item $(C)$ có tiệm cận xiên $y=m x-2$ (khi $m=0$, thì có tiệm cận ngang $y=-2)$
\end{itemize}
Đoạn thẳng $ x=-3m $ vuông góc với trục hoành, nên $ y=mx-2$ tạo với tiệm cận đứng một góc $45^\circ$ khi và chi khi nó tạo với chiểu dương của trục hoành một góc bằng $45^\circ$ hoặc $135^\circ$, tức là khi $\hoac{& m=\tan 45^\circ\\ & m=\tan 135^\circ}\Leftrightarrow m=\pm 1$.
\end{itemize}

	}
\end{bt}

\begin{bt}%[2D1V4-2]
% nguon: Archive/2024/TaiLieuGiangDay/1.HamSo/1. K12 - 4in1 - Hamso/ThubaiGoc/Bai4-Dang4.tex #2578
%[Dự án TL12New-4in1-NCT]
	Cho đường cong $\left(\mathrm{C}_m\right)$:$y=\dfrac{2 x^2-3 x+m}{x-m}$. Tìm các tiệm cận của $\left(\mathrm{C}_m\right)$ khi $m$ thay đổi.
\loigiai{
	Viết lại $\left(\mathrm{C}_m\right)$ dưới dạng: $y=2x+2 m-3+\dfrac{2 m^2-2 m}{x-m}$ $\quad (1)$.\\
	Ta có $2 m^2+2 m-3=0 \Leftrightarrow m=0$ hoặc $m=1$. Vì thế từ $(1)$ suy ra:
\begin{itemize}
	\item Nếu $m=0$, thì $y=2 x-3$ với $x \neq 0\left(\mathrm{C}_0\right)$ không có tiệm cận.
	\item Nếu $m=1$, thì $y=2 x-1$, với $x \neq 1\left(\mathrm{C}_1\right)$ không có tiệm cận.
	\item Nếu $m \neq 0, m \neq 1$, thì $\left(\mathrm{C}_m\right)$ có tiệm cận đứng $x=m$ và tiệm cận xiên $y=2 x+2 m-3$.
\end{itemize}	
	}
\end{bt}

\begin{bt}%[2D1C4-2]
% nguon: Archive/2024/TaiLieuGiangDay/1.HamSo/1. K12 - 4in1 - Hamso/ThubaiGoc/Bai4-Dang4.tex #2579
%[Dự án TL12New-4in1-NCT]
	Đường tiệm cận đứng và đường tiệm cận ngang của đồ thị hàm số $y=\dfrac{mx+1}{2m+1-x}$ cùng với hai trụ tọa độ tạo thành một hình chữ nhật có diện tích bằng $3$. Tìm $m$. 
	\loigiai{
		Ta có $\lim\limits_{x\to+\infty}\dfrac{mx+1}{2m+1-x}=-m$.\\
		Mặt khác $\lim\limits_{x\to(2m+1)^+}\dfrac{mx+1}{2m+1-x} =\lim\limits_{x\to(2m+1)^+}\dfrac{m(2m+1)+1}{2m+1-x} =\lim\limits_{x\to(2m+1)^+}\dfrac{2m^2+m+1}{2m+1-x}$.\\
		Mà $\heva{&\lim\limits_{x\to(2m+1)^+}\left(2m^2+m+1\right)=2m^2+m+1>0\\&\lim\limits_{x\to(2m+1)^+}(2m+1-x)=0\\&2m+1-x<0, \forall x>2m+1.}$\\
		Suy ra $\lim\limits_{x\to(2m+1)^+}\dfrac{mx+1}{2m+1-x}=-\infty $.\\
		Vậy đồ thị hàm số có hai đường tiệm cận $x=2m+1$ và $y=-m$.\\
		Hai đường tiệm cận tạo với hai trục tọa độ một hình chữ nhật có diện tích bằng $3$ suy ra $$|2m+1|\cdot|m|=3\Leftrightarrow\hoac{&2m^2+m=3\\&2m^2+m=-3(PTVN)}\Leftrightarrow 2m^2+m-3=0\Leftrightarrow\hoac{&m=1\\&m=\dfrac{-3}{2}.}$$}
\end{bt}

\begin{bt}%[2D1C4-2]
% nguon: Archive/2024/TaiLieuGiangDay/1.HamSo/1. K12 - 4in1 - Hamso/ThubaiGoc/Bai4-Dang4.tex #2580
%[Dự án TL12New-4in1-NCT]
Cho $\left(C _{m}\right): y=\dfrac{2x^2+m x-2}{x-1}$. Tìm $m$ để tiệm cận xiên của $\left(C _{m}\right)$ tạo với hệ trục tọa độ một tam giác có diện tích bằng $4$.
\loigiai{
Ta có $\left(C_m\right): y=2 x+m+2+\dfrac{m}{x-1}$ $\quad (1)$.\\
Vậy từ $(1)$ suy ra $\left(C_m\right)$ có tiệm cận xiên $y=2x+m+2$ khi $m \neq 0$. Tiệm cận xiên này cắt trục hoành tại $A\left(-\dfrac{m+2}{2};0\right)$ và trục tung tại $(0;m+2)$.\\
 Vì thế
$$
 S _{OAB}=4\Leftrightarrow\dfrac 12 OA\cdot OB=4\Leftrightarrow\dfrac 12\dfrac{|-m-2|}2\cdot| m+2|=4\Leftrightarrow| m+2|=4\Leftrightarrow\hoac{& 
	 m=2\\ & 
	m=-6 .}
$$
}
\end{bt}

\begin{bt}%[2D1N4-3]
% nguon: Archive/2024/TaiLieuGiangDay/1.HamSo/1. K12 - 4in1 - Hamso/ThubaiGoc/Bai4-Dang5.tex #2592
[Mức độ 1]%[Bài giảng 12, Phạm Tiến Long]
Tìm số giao điểm của tiệm cận ngang của đồ thị hàm số $y=\dfrac{x+2}{2x-1}$ với trục hoành.
\loigiai{
Hàm số đã cho có tập xác định là  $\mathscr{D}=\mathbb{R}\backslash \left\{\dfrac{1}{2}\right\}$.\\
Ta có $\lim\limits_{x\to \pm \infty}\dfrac{x+2}{2x-1}=\dfrac{1}{2}$.\\
Suy ra tiệm cận ngang của đồ thị hàm số là $y=\dfrac{1}{2}$.\\
Số giao điểm của đường thẳng $y=\dfrac{1}{2}$ với trục hoành là $0$.
}
\end{bt}

\begin{bt}%[2D1N4-3]
% nguon: Archive/2024/TaiLieuGiangDay/1.HamSo/1. K12 - 4in1 - Hamso/ThubaiGoc/Bai4-Dang5.tex #2593
[Mức độ 1]%[Bài giảng 12, Phạm Tiến Long]
Tìm tọa độ giao điểm của tiệm cận ngang và tiệm cận đứng của đồ thị hàm số $y=\dfrac{3-x}{x+1}$.
\loigiai{
Hàm số đã cho có tập xác định là  $\mathscr{D}=\mathbb{R}\backslash\left\{-1\right\}$.\\
Ta có 
\begin{itemize}
\item $\lim\limits_{x\to -1^+}\dfrac{3-x}{x+1}=+\infty$;
\item $\lim\limits_{x\to -1^-}\dfrac{3-x}{x+1}=-\infty$;
\item $\lim\limits_{x\to +\infty}\dfrac{3-x}{x+1}=-1$;
\item $\lim\limits_{x\to -\infty}\dfrac{3-x}{x+1}=-1$.
\end{itemize}
Suy ra tiệm cận đứng và tiệm cận ngang của đồ thị hàm số lần lượt có phương trình là $x=-1$, $y=-1$.\\
Vậy giao điểm của tiệm cận ngang và tiệm cận đứng của đồ thị hàm số đã cho có tọa độ là $(-1;-1)$.
}
\end{bt}

\begin{bt}%[2D1N4-3]
% nguon: Archive/2024/TaiLieuGiangDay/1.HamSo/1. K12 - 4in1 - Hamso/ThubaiGoc/Bai4-Dang5.tex #2594
[Mức độ 2]%[Bài giảng 12, Phạm Tiến Long]
Tìm tọa độ giao điểm của tiệm cận đứng của đồ thị hàm số $y=\dfrac{2x}{2-x}$ với đường thẳng $y=2x+3$.
\loigiai{
Hàm số $y=\dfrac{2x}{2-x}$ có tập xác đinh là  $\mathscr{D}=\mathbb{R}\backslash\left\{2\right\}$.\\
Ta có $\lim\limits_{x\to 2^+}\dfrac{2x}{2-x}=-\infty$ và $\lim\limits_{x\to 2^-}\dfrac{2x}{2-x}=+\infty$.\\
Suy ra tiệm cận đứng của đồ thị hàm số là $x=2$.\\
Khi đó tọa độ giao điểm của tiệm cận đứng $x=2$ và đường thẳng $y=2x+3$ là nghiệm của hệ phương trình
$$\heva{&x=2\\&y=2x+3}\Leftrightarrow \heva{&x=2\\&y=7.}$$
Vậy tọa độ giao điểm cần tìm là $(2;7)$.
}
\end{bt}

\begin{bt}%[2D1N4-3]
% nguon: Archive/2024/TaiLieuGiangDay/1.HamSo/1. K12 - 4in1 - Hamso/ThubaiGoc/Bai4-Dang5.tex #2595
[Mức độ 2]%[Bài giảng 12, Phạm Tiến Long]
Tìm tọa độ giao điểm của đường tiệm cận xiên của đồ thị hàm số $f(x)=x-2+\dfrac{3}{x+1}$ và đường thẳng $y=2x-5$.
\loigiai{
Hàm số $f(x)=x-2+\dfrac{3}{x+1}$ có tập xác đinh là $\mathscr{D}=\mathbb{R}\backslash\left\{-1\right\}$.\\
Ta có 
\begin{itemize}
\item $\lim\limits_{x\to +\infty}[f(x)-(x-2)]=\lim\limits_{x\to +\infty}\dfrac{3}{x+1}=x-2$;
\item $\lim\limits_{x\to -\infty}[f(x)-(x-2)]=\lim\limits_{x\to -\infty}\dfrac{3}{x+1}=x-2$.
\end{itemize}
Suy ra tiệm cận xiên của đồ thị hàm số là $y=x-2$.\\
Giao điểm của tiệm cận xiên $y=x-2$ và đường thẳng $y=2x-5$ là nghiệm của hệ phương trình
$$\heva{&y=x-2\\&y=2x-5}\Leftrightarrow \heva{&x-y=2\\&2x-y=5}\Leftrightarrow \heva{&x=3\\&y=1.}$$
Vậy tọa độ giao điểm cần tìm là $(3;1)$.
}
\end{bt}

\begin{bt}%[2D1N4-3]
% nguon: Archive/2024/TaiLieuGiangDay/1.HamSo/1. K12 - 4in1 - Hamso/ThubaiGoc/Bai4-Dang5.tex #2596
[Mức độ 2]%[Bài giảng 12, Phạm Tiến Long]
Biết rằng parabol $y=x^2+3x-4$ cắt tiệm cận ngang của đồ thị hàm số $y=2^x$ tại hai điểm phân biệt $M$ và $N$. Tính độ dài $MN$.
\loigiai{
Hàm số $y=2^x$ có tập xác định là  $\mathscr{D}=\mathbb{R}$.\\
Đồ thị hàm số $y=2^x$ có tiệm cận ngang $y=0$.\\
Ta có $x^2+3x-4=0\Leftrightarrow \hoac{&x=1\\&x=-4.}$\\
Vậy $M(1;0)$, $N(-4;0)$.\\
Suy ra $MN=\sqrt{(-4-1)^2+(0-0)^2}=5$.
}
\end{bt}

\begin{bt}%[2D1N4-3]
% nguon: Archive/2024/TaiLieuGiangDay/1.HamSo/1. K12 - 4in1 - Hamso/ThubaiGoc/Bai4-Dang5.tex #2597
[Mức độ 2]%[Bài giảng 12, Phạm Tiến Long]
	Biết rằng đường thẳng $d\colon x+my+n=0$ đi qua điểm $A(-2;1)$ và song song với tiệm cận đứng của đồ thị hàm số $y=\log x$. Tính $m+n$.
	\loigiai{
Hàm số $y=\log x$ có tập xác định là  $\mathscr{D}=(0;+\infty)$.\\
Tiệm cận đứng của đồ thị hàm số $y=\log x$ là $x=0$.\\
Đường thẳng	$d$ đi qua $A$ và song song với đường thẳng $x=0$ có phương trình là $x+2=0$.\\
Suy ra $m=0$, $n=2$.\\
Vậy $m+n=2$.	
	}
\end{bt}

\begin{bt}%[2D1N4-3]
% nguon: Archive/2024/TaiLieuGiangDay/1.HamSo/1. K12 - 4in1 - Hamso/ThubaiGoc/Bai4-Dang5.tex #2598
[Mức độ 2]%[Bài giảng 12, Phạm Tiến Long]
Tiệm cận xiên của đồ thị hàm số $f(x)=\dfrac{x^2-3x+1}{x-2}$ cắt hai trục tọa độ $Ox$, $Oy$ tại hai điểm $A$ và $B$. Tính diện tích tam giác $OAB$.
\loigiai{
Hàm số đã cho có tập xác định $\mathscr{D}=\mathbb{R}\backslash\{2\}$.\\
	Ta có $\begin{aligned}[t]
	a&=\lim\limits_{x \rightarrow+\infty} \dfrac{f(x)}{x}=\lim\limits_{x \rightarrow+\infty} \dfrac{x^{2}-3 x+1}{x^{2}-2 x}=1;\\
	b&=\lim\limits_{x \rightarrow+\infty}[f(x)-a x]=\lim\limits_{x \rightarrow+\infty}\left(\dfrac{x^{2}-3 x+1}{x-2}-x\right)=\lim\limits_{x \rightarrow+\infty} \dfrac{-x+1}{x-2}=-1.
\end{aligned}$\\
Do đó, đồ thị hàm số có tiệm cận xiên là đường thẳng $y=x-1$.\\
Đường thẳng cắt trục $Ox$, $Oy$ lần lượt tại hai điểm $A(1;0)$ và $B(0;-1)$.\\
Diện tích tam giác $OAB$ là $S=\dfrac{1}{2}\cdot OA\cdot OB=\dfrac{1}{2}\cdot1 \cdot 1=\dfrac{1}{2}$.
}
\end{bt}

\begin{bt}%[2D1N4-3]
% nguon: Archive/2024/TaiLieuGiangDay/1.HamSo/1. K12 - 4in1 - Hamso/ThubaiGoc/Bai4-Dang5.tex #2602
[Mức độ 2]%[Bài giảng 12, Phạm Tiến Long]
Tìm tọa giao điểm của tiệm cận đứng của đồ thị hàm số $y=\dfrac{1}{2-x}$ với trục hoành.
\loigiai{
Hàm số đã cho có tập xác định là $\mathscr{D}=\mathbb{R}\backslash\{2\}$.\\
Ta có $\lim\limits_{x\to 2^-}\dfrac{1}{2-x}=+\infty$ và $\lim\limits_{x\to 2^+}\dfrac{1}{2-x}=-\infty$.\\
Suy ra tiệm cận đứng của đồ thị hàm số là $x=2$.\\
Giao điểm của đường thẳng $x=2$ và trục hoành có tọa độ là $(2;0)$.
}
\end{bt}

\begin{bt}%[2D1H4-3]
% nguon: Archive/2024/TaiLieuGiangDay/1.HamSo/1. K12 - 4in1 - Hamso/ThubaiGoc/Bai4-Dang5.tex #2603
[Mức độ 2]%[Bài giảng 12, Phạm Tiến Long]
	Tìm tọa giao điểm của tiệm cận ngang của đồ thị hàm số $y=\dfrac{3x-4}{x+2}$ với đường thẳng $y=4-x$.
	\loigiai{
		Hàm số đã cho có tập xác định là $\mathscr{D}=\mathbb{R}\backslash\{-2\}$.\\
		Ta có $\lim\limits_{x\to \pm \infty}\dfrac{3x-4}{x+2}=3$.\\
		Suy ra tiệm cận ngang của đồ thị hàm số là $y=3$.\\
		Giao điểm của hai đường thẳng $y=3$ và $y=4-x$ là nghiệm của hệ phương trình
		$$\heva{&y=3\\&y=4-x}\Leftrightarrow \heva{&x=1\\&y=3.}$$
	Vậy tọa độ giao điểm càn tìm là $(1;3)$.
	}
\end{bt}

\begin{bt}%[2D1H4-3]
% nguon: Archive/2024/TaiLieuGiangDay/1.HamSo/1. K12 - 4in1 - Hamso/ThubaiGoc/Bai4-Dang5.tex #2604
[Mức độ 2]%[Bài giảng 12, Phạm Tiến Long]
Tiệm cận xiên của đồ thị hàm số $y=f(x)=\dfrac{x^2-2x+1}{x+2}$ cắt tiệm cận đứng của đồ thị hàm số $y=g(x)=\log_2 x$ tại điểm $M$. Tìm tọa độ điểm $M$.
\loigiai{
Hàm số $f(x)$ có tập xác định là $\mathbb{R}\backslash \{-2\}$.\\
Hàm số $g(x)$ có tập xác định là $(0;+\infty)$.\\
	Ta có $\begin{aligned}[t]
	a&=\lim\limits_{x \rightarrow+\infty} \dfrac{f(x)}{x}=\lim\limits_{x \rightarrow+\infty} \dfrac{x^2-2x+1}{x^2+2x}=1;\\
	b&=\lim\limits_{x \rightarrow+\infty}[f(x)-ax]=\lim\limits_{x \rightarrow+\infty}\left(\dfrac{x^2-2x+1}{x+2}-x\right)=\lim\limits_{x \rightarrow+\infty} \dfrac{-4x+1}{x+2}=-4.
\end{aligned}$\\
Ta cũng có $\lim\limits_{x \rightarrow-\infty} \dfrac{f(x)}{x}=1$; $\lim\limits_{x \rightarrow-\infty}[f(x)-x]=-4$.
\\
Do đó, đồ thị hàm số có tiệm cận xiên là $y=x-4$.\\
Đồ thị hàm số $y=g(x)=\log_2 x$ có tiệm cận đứng là $x=0$.\\
Tọa độ giao điểm cần tìm là nghiệm của hệ phương trình $$\heva{&y=x-4\\&x=0}\Leftrightarrow \heva{&x=0\\&y=-4.}$$
Vậy tọa độ giao điểm cần tìm là $(0;-4)$.
}
\end{bt}

\begin{bt}%[2D1H4-3]
% nguon: Archive/2024/TaiLieuGiangDay/1.HamSo/1. K12 - 4in1 - Hamso/ThubaiGoc/Bai4-Dang5.tex #2605
[Mức độ 2]%[Bài giảng 12, Phạm Tiến Long]
Tiệm cận xiên của đồ thị hàm số $y=f(x)=\dfrac{2x^2-x}{x+3}$ cắt hai trục tọa độ tại hai điểm $A$, $B$. Tính độ dài $AB$.
\loigiai{
Hàm số đã cho có tập xác định là $\mathbb{R}\backslash \{-3\}$.\\
	Ta có $\begin{aligned}[t]
	a&=\lim\limits_{x \rightarrow+\infty} \dfrac{f(x)}{x}=\lim\limits_{x \rightarrow+\infty} \dfrac{2x^2-x}{x^2+3x}=2;\\
	b&=\lim\limits_{x \rightarrow+\infty}[f(x)-ax]=\lim\limits_{x \rightarrow+\infty}\left(\dfrac{2x^2-x}{x+3}-2x\right)=\lim\limits_{x \rightarrow+\infty} \dfrac{-7x}{x+3}=-7.
\end{aligned}$\\
Ta cũng có $\lim\limits_{x \rightarrow-\infty} \dfrac{f(x)}{x}=2$; $\lim\limits_{x \rightarrow-\infty}[f(x)-ax]=-7$.
\\
Do đó, đồ thị hàm số có tiệm cận xiên là đường thẳng $d\colon y=2x-7$.\\
Đường thẳng $d$ cắt hai trục tọa độ tại hai điểm $A\left(0;-7\right)$ và $B\left(\dfrac{7}{2};0\right)$.\\
Ta có $AB=\sqrt{\left(\dfrac{7}{2}-0\right)^2+(0+7)^2}=\dfrac{7\sqrt{5}}{2}$.
}
\end{bt}

\begin{bt}%[2D1V4-3]
% nguon: Archive/2024/TaiLieuGiangDay/1.HamSo/1. K12 - 4in1 - Hamso/ThubaiGoc/Bai4-Dang5.tex #2606
[Mức độ 3]%[Bài giảng 12, Phạm Tiến Long]
Cho hàm số $f(x)=\dfrac{m^2x+3}{x-5}$ và $g(x)=\dfrac{3x^2-4x-2}{x-6}$. Tìm $m$ để tiệm cận ngang của đồ thị hàm số $f(x)$ cắt tiệm cận xiên của đồ thị hàm số $g(x)$ tại điểm có hoành độ bằng $2$.
	\loigiai{
Hàm số $f(x)$ có tập xác định là $\mathbb{R}\backslash\{5\}$.\\
Hàm số $g(x)$ có tập xác định là $\mathbb{R}\backslash\{6\}$.\\
Đồ thị hàm số $f(x)=\dfrac{m^2x+3}{x-5}$ có tiệm cận ngang là đường thẳng $d\colon y=m^2$.\\
Xét hàm số $g(x)=\dfrac{3x^2-4x-2}{x-6}$ ta có\\
 $\begin{aligned}[t]
	a&=\lim\limits_{x \rightarrow+\infty} \dfrac{g(x)}{x}=\lim\limits_{x \rightarrow+\infty} \dfrac{3x^2-4x-2}{x^2-6x}=3;\\
	b&=\lim\limits_{x \rightarrow+\infty}[g(x)-ax]=\lim\limits_{x \rightarrow+\infty}\left(\dfrac{3x^2-4x-2}{x-6}-3x\right)=\lim\limits_{x \rightarrow+\infty} \dfrac{14x-2}{x-6}=14.
\end{aligned}$\\
Ta cũng có $\lim\limits_{x \rightarrow-\infty} \dfrac{g(x)}{x}=3$; $\lim\limits_{x \rightarrow-\infty}[g(x)-ax]=14$.
\\
Do đó, đồ thị hàm số $g(x)$ có tiệm cận xiên là đường thẳng $d'\colon y=3x+14$.\\
Gọi $A\left(x;y\right)=d\cap d'$. Ta có $x=2\Rightarrow y=3\cdot 2+14=20$.\\Vậy $A(3;20)$.\\
Vì $A\in d$ nên ta có $m^2=20\Rightarrow m=\pm2\sqrt{5}$.
	}
\end{bt}

\begin{bt}%[2D1N2-2]
% nguon: Archive/2024/TaiLieuGiangDay/1.HamSo/1. K12 - 4in1 - Hamso/ThubaiGoc/Bai2-Dang2.tex #2617
[Cánh Diều]%[BG12, Tran Tony]
	Dựa vào đồ thị hàm số $y=f(x)=-x^3+3x$ trong hình vẽ, hãy chỉ ra các điểm cực trị của hàm số đó.
	\loigiai{
		\immini{
			\begin{itemize}
				\item Xét khoảng $\left(-\sqrt{3};0\right)$ chứa điểm $x=-1$. Quan sát đồ thị của hàm số $y=f(x)=-x^3+3 x$ trong hình vẽ, ta thấy: $f(x)>f(-1)$ với mọi $x\in\left(-\sqrt{3};0\right)$ và $x \neq-1$.\\
				Vậy $x=-1$ là điểm cực tiểu của hàm số $y=f(x)$.
				\item Xét khoảng $(0;\sqrt{3})$ chứa điểm $x=1$. Quan sát đồ thị của hàm số $y=f(x)=-x^3+3 x$ trong hình vẽ, ta thấy: $f(x)<f(1)$ với mọi $x\in\left(0;\sqrt{3}\right)$ và $x \neq 1$.\\
				Vậy $x=1$ là điểm cực đại của hàm số $y=f(x)$.
			\end{itemize} 
		}{
			\begin{tikzpicture}[scale=0.8, font=\footnotesize, line join=round, line cap=round, >=stealth]
			\path
			(0,0) coordinate (O)
			(-1.73,0) coordinate (A)
			(-1,0) coordinate (B)
			(1,0) coordinate (C)
			(1.73,0) coordinate (D)
			(1,2) coordinate (E)
			(-1,-2) coordinate (F)
			;
				\draw[->] (-2.8,0)--(2.8,0) node[below] {$x$};
				\draw[->] (0,-2.5)--(0,2.8) node[left] {$y$};
				\draw (0,0) node [below left] {$O$};
				\draw[thin] (-1,1pt)--(-1,-1pt) node [above] {$-1$};
				\draw[thin] (1,1pt)--(1,-1pt) node [below] {$1$};
				\draw[thin] (1pt,-2)--(-1pt,-2) node [right] {$-2$};
				\draw[thin] (1pt,2)--(-1pt,2) node [left] {$2$};
				\draw (D) node[above right]{$\sqrt{3}$};
				\draw (A) node[above left]{$-\sqrt{3}$};
				\draw[dashed,thin](-1,0)--(-1,-2)--(0,-2) (1,0)--(1,2)--(0,2);
				\begin{scope}
					\clip (-2.3,-2.6) rectangle (2.3,2.6);
					\draw[samples=200,domain=-2.2:2.2,smooth,variable=\x] plot (\x,{-1*((\x)^3)+3*(\x)});
				\end{scope}
				\foreach \x in {A,B,O,C,D,E,F} \fill(\x) circle (1.2pt);
			\end{tikzpicture}
		}
	}
\end{bt}

\begin{bt}%[2D1N2-2]
% nguon: Archive/2024/TaiLieuGiangDay/1.HamSo/1. K12 - 4in1 - Hamso/ThubaiGoc/Bai2-Dang2.tex #2618
[Chân trời sáng tạo]%[BG12, Tran Tony]
	Tìm cực trị của hàm số $y=f(x)$ có đồ thị được cho trong hình vẽ.
	\begin{center}
		\begin{tikzpicture}[scale=0.5, font=\footnotesize, line join=round, line cap=round, >=stealth]
			\def\xmin{-2} \def\xmax{8}
			\def\ymin{-1} \def\ymax{7} 
			\draw[->] (\xmin,0)--(\xmax,0) node [below]{$x$};
			\draw[->] (0,\ymin)--(0,\ymax) node [left]{$y$};
			\fill (0,0)circle (1pt)node[below right]{$O$};
			\foreach \x/\g in {-1/below,1/below,2/below,3/below,4/below,5/below,6/below,7/below}  \draw[thin] (\x,1pt)--(\x,-1pt) node[\g]{$\x$};
			\foreach \y/\g in {1/left,2/above left,4/left,5/left,6/left} \draw[thin] (1pt,\y)--(-1pt,\y) node[\g]{$\y$};
			\draw[dashed] (-1,0)--(-1,2)--(0,2) (1,0)--(1,5)--(0,5) 
			(2,0)--(2,4)
			(3,0)--(3,2)--(0,2)
			(4,0)--(4,1)--(0,1)
			(5,0)--(5,4)
			(6,0)--(6,6)--(0,6)
			(7,0)--(7,4)--(0,4)
			;
			\draw (0,4) parabola bend (1,5) (2,4) (3,2) parabola bend (4,1) (5,4)
			(5,4) parabola bend (6,6) (7,4) 
			;
			\draw (-1,2)..controls+(64:1) and +(240:1)..(0,4)
			(2,4)..controls+(297:1) and +(122:1)..(3,2)
			;
			\fill (-1,2)circle(1pt)
			(0,4)circle(1pt)
			(1,5)circle(1pt)
			(2,4)circle(1pt)
			(3,2)circle(1pt)
			(4,1)circle(1pt)
			(5,4)circle(1pt)
			(6,6)circle(1pt)
			(7,4)circle(1pt)
			;
		\end{tikzpicture}
	\end{center}
	\loigiai{
		Hàm số $y=f(x)$ có:
		\begin{itemize}
			\item $x=1$ là điểm cực đại vì $f(x)<f(1)$ với mọi $x \in(0;2) \backslash\{1\}$,  $y_{\text{CĐ}}=f(1)=5$;
			\item $x=6$ là điểm cực đại vì $f(x)<f(6)$ với mọi $x \in(5;7) \backslash\{6\}$,  $y_{\text{CĐ}}=f(6)=6$;
			\item $x=4$ là điểm cực tiểu vì $f(x)>f(4)$ với mọi $x \in(3;5) \backslash\{4\}$, $y_{\text{CT}}=f(4)=1$.
		\end{itemize}
	}
\end{bt}

\begin{bt}%[2D1H2-2]
% nguon: Archive/2024/TaiLieuGiangDay/1.HamSo/1. K12 - 4in1 - Hamso/ThubaiGoc/Bai2-Dang2.tex #2619
[Cùng khám phá]%[BG12, Tran Tony]
	Cho hàm số $y=f\left(x\right)$ liên tục trên $\mathbb{R}$ và có bảng biến thiên như sau
	\begin{center}
		\begin{tikzpicture}
			\tkzTabInit[lgt=1.2,espcl=3]
			{$x$ /0.7, $f’(x)$ /0.7, $f(x)$ /2.5}
			{$-\infty$,$0$,$1$,$+\infty$}
			\tkzTabLine{ ,+,d,-,z,+, }
			\tkzTabVar{-/$-\infty$,+/$2$,-/$-3$,+/$+\infty$}
		\end{tikzpicture}
	\end{center}
	Xác định các cực trị của hàm số $f\left(x\right)$.
	\loigiai{ Ta có $f'\left(x\right)>0$ với mọi $x\in \left(-\infty;0\right)$ và $f'\left(x\right)<0$ với mọi $x\in \left(0;1\right)$ nên hàm số đạt cực đại tại $x=0$ và $f_{\text{CĐ}}=f\left(0\right)=2$.\\
		Ta có $f'\left(x\right)<0$ với mọi $x\in \left(0;1\right)$ và $f'\left(x\right)>0$ với mọi $x \in \left(1;+\infty\right)$ nên hàm số đạt cực tiểu tại $x=1$và $f_{CT}=f\left(1\right)=-3$.
	}
\end{bt}

\begin{bt}%[2D1H2-2]
% nguon: Archive/2024/TaiLieuGiangDay/1.HamSo/1. K12 - 4in1 - Hamso/ThubaiGoc/Bai2-Dang2.tex #2620
%[BG12, Tran Tony]
	Cho hàm số $y=f(x)$ liên tục trên $\mathbb{R}$ và có bảng xét dấu $f'(x)$ như sau
	\begin{center}
		\begin{tikzpicture}[scale=1, line join=round, line cap=round, >=stealth]
			\tkzTabInit[nocadre=false, lgt=1.2, espcl=2.5, deltacl=0.6]
			{$x$ /0.6,$f'(x)$ /0.7}
			{$-\infty$,$-1$,$0$,$1$,$2$,$+\infty$}
			\tkzTabLine{,+,$0$,-,$0$,+,||,-,$0$,-,}
		\end{tikzpicture}
	\end{center}
	Tìm số điểm cực đại của hàm số đã cho.
	\loigiai{
		Dựa vào bảng xét dấu đạo hàm, ta thấy $f'(x)$ đổi dấu từ $+$ sang $-$ qua các điểm $x=-1$, $x=1$ và $f(x)$ liên tục trên $\mathbb{R}$ nên $f(x)$ có $2$ điểm cực đại.
	}
\end{bt}

\begin{bt}%[2D1H2-2]
% nguon: Archive/2024/TaiLieuGiangDay/1.HamSo/1. K12 - 4in1 - Hamso/ThubaiGoc/Bai2-Dang2.tex #2621
%[BG12, Tran Tony]
	\immini{Cho hàm số $y=f(x)$ xác định trên $\mathbb{R}$ và có đồ thị hàm số $y=f'(x)$ là đường cong ở hình bên. Hỏi hàm số $y=f(x)$ có bao nhiêu điểm cực trị?
	}
	{
		\begin{tikzpicture}[scale=1, font=\footnotesize, line join=round, line cap=round, >=stealth]
			\clip (-2.5,-.5) rectangle (2.5,2.5);
			\def\xmin{-2.5} \def\xmax{2.5}
			\def\ymin{-.5} \def\ymax{2.5} 
			\draw[->] (\xmin,0)--(\xmax,0) node [above left]{\tiny $x$};
			\draw[->] (0,\ymin)--(0,\ymax) node [below right]{\tiny $y$};
			\node at (0,0) [below left]{\tiny $O$};
			\node at (1.8,1.2) [below]{\tiny $y=f'(x)$};
			\draw (-1.7,2.2)
			..controls +(-80:.3) and +(140:.2) ..(-1.3,-.3)
			..controls +(-40:.1) and +(-120:.3) ..(-1,.7)
			..controls +(60:.2) and +(-150:.2) ..(-.6,0)
			..controls +(-30:.2) and +(-140:.35) ..(.45,2.2)
			..controls +(40:.3) and +(100:.3) ..(1.2,-.4)
			;
		\end{tikzpicture}	
	}
	\loigiai{
		Số cực trị của hàm số $ y=f(x) $ là số điểm cắt của $ y=f'(x) $ với trục hoành.\\
		Dựa vào đồ thị hàm số ta thấy có ba điểm cắt và một điểm đồ thị hàm số tiếp xúc trục hoành.\\
		Vậy hàm số $ y=f(x) $ có $3$ điểm cực trị.		
	}
\end{bt}

\begin{bt}%[2D1H2-2]
% nguon: Archive/2024/TaiLieuGiangDay/1.HamSo/1. K12 - 4in1 - Hamso/ThubaiGoc/Bai2-Dang2.tex #2622
%[BG12, Tran Tony]
	\immini{
		Cho hàm số $y=\dfrac{ax^2+bx+c}{dx+e}$ có đồ thị $(C)$ như hình vẽ bên. Viết phương trình đường thẳng đi qua hai điểm cực trị của đồ thị $(C)$.
	}{
		\begin{tikzpicture}[scale=0.6, font=\footnotesize, line join=round, line cap=round, >=stealth,yscale=0.7]
			\tikzset{every node/.style={scale=0.9}}
			\draw[->] (-3.6,0)--(5.6,0) node[below] {$x$};
			\draw[->] (0,-6.1)--(0,6.1) node[left] {$y$};
			\draw (0,0) node [below left] {$O$};
			\draw[thin] (-1,1pt)--(-1,-1pt) node [above] {$-1$};
			\draw[thin] (3,1pt)--(3,-1pt) node [below] {$3$};
			\draw[thin] (1pt,-4)--(-1pt,-4) node [right] {$-4$};
			\draw[thin] (1pt,4)--(-1pt,4) node [left] {$4$};
			\draw[dashed,thin](-1,0)--(-1,-4)--(0,-4);
			\draw[dashed,thin](3,0)--(3,4)--(0,4);
			\draw[dashed,thin] (1.01,-6)--(1.01,6);
			\begin{scope}
				\clip (-3.5,-6) rectangle (5.5,6);
				\draw[samples=200,domain=-3:0.99,smooth,variable=\x] plot (\x,{(1*((\x)^2)+-2*(\x)+5)/(1*(\x)+-1)});
				\draw[samples=200,domain=1.01:5,smooth,variable=\x] plot (\x,{(1*((\x)^2)+-2*(\x)+5)/(1*(\x)+-1)});
				\draw[dashed,thin] (-3.1,-4.1)--(5.1,4.1);
			\end{scope}
		\end{tikzpicture}
	}
	\loigiai{
		Quan sát đồ thị $(C)$ ta thấy, đồ thị có hai điểm cực trị là $A(-1;-4)$, $B(3;4)$.\\
		Đường thẳng đi qua hai điểm cực trị $A$, $B$ có phương trình là
		$\dfrac{x+1}{3+1}=\dfrac{y+4}{4+4}\Leftrightarrow y=2x-2$.
	}
\end{bt}

\begin{bt}%[2D1H2-2]
% nguon: Archive/2024/TaiLieuGiangDay/1.HamSo/1. K12 - 4in1 - Hamso/ThubaiGoc/Bai2-Dang2.tex #2623
%[BG12, Tran Tony]
	Cho hàm số $y=f(x)$ liên tục trên $\mathbb{R}$ và có bảng biến thiên như sau
	\begin{center}
		\begin{tikzpicture}
			\tkzTabInit[nocadre=false,lgt=1.2,espcl=2.5,deltacl=0.6]
			{$x$ /0.6,$f'(x)$ /0.6,$f(x)$ /2}
			{$-\infty$,$-2$,$0$,$3$,$+\infty$}
			\tkzTabLine{,+,$0$,-,$0$,+,$0$,-,}
			\tkzTabVar{-/$-\infty$, +/$-1$,-/$-3$,+/$2$,-/$-\infty$}
		\end{tikzpicture}
	\end{center}
	Hàm số $y=\ln (f(x))$ có tất cả bao nhiêu điểm cực đại?
	\loigiai{
		Xét $f(x)=0\Leftrightarrow \hoac{& x=a &\text{ với } 0<a<3\\ & x=b &\text{ với } b>3.}$\\
		Hàm số $y=\ln (f(x))$ xác định trên khoảng $(a;b)$.\\
		$y'=\dfrac{f'(x)}{f(x)}$, $y'=0\Leftrightarrow f'(x)=0\Leftrightarrow x=3\in(a;b)$.\\
		\begin{center}
			\begin{tikzpicture}
				\tkzTabInit[nocadre=false,lgt=1,espcl=2.5,deltacl=0.6]
				{$x$ /0.6, $y'$ /0.6, $y$ /2.1}
				{$a$,$3$,$b$}
				\tkzTabLine{,+,$0$,-,}
				\tkzTabVar{-/$-\infty$, +/$\ln2$,-/$-\infty$}
			\end{tikzpicture}
		\end{center}
		Vậy hàm số $y=\ln (f(x))$ có $1$ điểm cực đại.
	}
\end{bt}

\begin{bt}%[2D1H2-2]
% nguon: Archive/2024/TaiLieuGiangDay/1.HamSo/1. K12 - 4in1 - Hamso/ThubaiGoc/Bai2-Dang2.tex #2627
%[BG12, Tran Tony]
	\immini{Cho hàm số $y=f\left(x\right)$ liên tục trên đoạn $\left[-3;3\right]$ và có đồ thị như hình vẽ bên. Hãy xác định các điểm cực trị, giá trị cực trị của hàm số $f\left(x\right)$ trên khoảng $\left(-3;3\right)$.} {\begin{tikzpicture}[>=stealth,scale=0.75,yscale=0.75]
		\draw[->] (-3.5,0) --(0,0) node[above right]{$O$}--(3.5,0)node[below]{$x$};
		\draw[->](0,-3.5)--(0,3.5)node[right]{$y$};
		\draw[dashed] (-3,0)node[above]{$-3$}--(-1,0)node[below]{$-1$}--(1,0)node[above]{$1$}--(2,0)node[below]{$2$}--(3,0)node[above]{$3$};
		\draw (0,-3) node[below left]{$-3$} --(0,1)node[above left ]{$1$}--(0,3)node[left]{$3$};
		\draw[smooth] plot[domain=-3.1:1]
			(\x,{-(\x)^2-2*(\x)});
		\draw[smooth] plot[domain=1:3.05]
			(\x,{-6*(\x)^2+24*(\x)-21});
		\draw[dashed] (-3,0)--(-3,-3)--(0,-3)--(1,-3)--(1,0) ;
		\draw[dashed](1,-3)--(3,-3)--(3,0);
		\draw[dashed] (-1,0)--(-1,1)--(0,1);
		\draw[dashed] (2,0)--(2,3)--(0,3); 
	\end{tikzpicture}}
	\loigiai{
		\begin{itemize}
			\item Với $x \in \left(-3;1\right)\subset \left(-3;3\right)$, hàm số đạt cực đại tại $x=-1$ và giá trị cực đại $f_{\text{CĐ}}=1$.
			\item Với $x \in \left(1;3\right)\subset \left(-3;3\right)$, hàm số đạt cực đại tại $x=2$ và giá trị cực đại $f_{\text{CĐ}}=3$.
		\end{itemize}
	}
\end{bt}

\begin{bt}%[2D1H2-2]
% nguon: Archive/2024/TaiLieuGiangDay/1.HamSo/1. K12 - 4in1 - Hamso/ThubaiGoc/Bai2-Dang2.tex #2628
%[BG12, Tran Tony]
	Cho hàm số $y=f(x)$ liên tục trên các khoảng $\left(-\infty;1\right)$, $\left(1;+\infty\right)$ các có bảng biến thiên như sau
	\begin{center}
		\begin{tikzpicture}[>=stealth]
			\tkzTabInit[nocadre=false,lgt=1.5,espcl=2,deltacl=0.6]{$x$/.6 ,$f'(x)$/.6,$f(x)$/2}
			{$-\infty$ ,$1$ , $2$ , $3$, $4$, $5$ , $+\infty$}
			\tkzTabLine{ , + , d , - , $0$ , + ,$0$ , + ,d, - ,$0$ , + , }
			\tkzTabVar{-/$-\infty$, +D+/$+\infty$/$+\infty$,-/$-2$,R,+/$1$,-/$-1$, +/$+\infty$}
		\end{tikzpicture}
	\end{center}
	Xác định cực trị của hàm số đã cho.		
	\loigiai{	
		Dựa vào bảng biến thiên, ta thấy
		\begin{itemize}
			\item Hàm số đạt cực tiểu tại $x=2$, $y_{\text{CT}}=-2$.
			\item Hàm số đạt cực tiểu tại $x=2$, $y_{\text{CT}}=-1$.
			\item Hàm số đạt cực đại tại $x=4$, $y_{\text{CĐ}}=1$.
		\end{itemize}
	}
\end{bt}

\begin{bt}%[2D1H2-2]
% nguon: Archive/2024/TaiLieuGiangDay/1.HamSo/1. K12 - 4in1 - Hamso/ThubaiGoc/Bai2-Dang2.tex #2629
%[BG12, Tran Tony]
	\immini{Cho hàm số $y=f(x)$ liên tục trên đoạn $[0 ; 3]$ thoả mãn\\ $f'\left(\dfrac{1}{3}\right)=f'(1)=f'\left(\dfrac{5}{2}\right)=0$ và có đồ thị là đường cong như hình bên. Tìm cực trị của hàm số đã cho trên khoảng $(0;3)$.
	}{
	\begin{tikzpicture}[scale=0.9, font=\footnotesize, line join=round, line cap=round, >=stealth]
	\tikzset{declare function={xmin=-2;xmax=4;ymin=-3;ymax=2; f(\x)=32/45*(\x)^4-32/9*(\x)^3+224/45*(\x)^2-32/15*(\x);},smooth,samples=450}
	\draw[->] (xmin,0)--(xmax,0) node[shift={(-100:7pt)}]{$ x $};
	\draw[->] (0,ymin)--(0,ymax) node[shift={(190:7pt)}]{$ y $};
	\fill (0,0) node[shift={(135:10pt)}]{$ O $};
	\draw[dashed](1/3,0)node[shift={(90:8pt)}]{$\tfrac{1}{3}$}|-(0,-1/3)node[shift={(-150:9pt)}]{$-\tfrac{1}{3}$}
			(2.5,0)node[shift={(90:7pt)}]{$\tfrac{5}{2}$}|-(0,-2)node[shift={(180:9pt)}]{$-2$}	(1,0)node[shift={(90:7pt)}]{$1$};
	\clip (xmin,ymin) rectangle (xmax,ymax);
	\draw  plot[domain=xmin:xmax] (\x, {f(\x)});
	\end{tikzpicture}}
	\loigiai{
		\begin{itemize}
			\item Hàm số đạt cực tiểu tại $x=\dfrac{1}{3}$, $y_{\text{CT}}=-\dfrac{1}{3}$.
			\item Hàm số đạt cực tiểu tại $x=\dfrac{5}{2}$, $y_{\text{CT}}=-2$.
			\item Hàm số đạt cực đại tại $x=1$, $y_{\text{CĐ}}=0$.
	\end{itemize}
	}
\end{bt}

\begin{bt}%[2D1H2-2]
% nguon: Archive/2024/TaiLieuGiangDay/1.HamSo/1. K12 - 4in1 - Hamso/ThubaiGoc/Bai2-Dang2.tex #2630
%[BG12, Tran Tony]
	\immini{Cho hàm số $y=f(x)$ có bảng biến thiên như hình vẽ bên. Tìm số điểm cực trị của hàm số $y=\left| f(x)\right|$.
	}{
		\begin{tikzpicture}
			\tkzTabInit[nocadre=false,lgt=1.2,espcl=2.5,deltacl=0.6]
			{$x$ /0.6, $y'$ /0.6, $y$ /2.5}
			{$-\infty$,$-1$,$3$,$+\infty$}
			\tkzTabLine{,+,$0$,-,$0$,+,}
			\tkzTabVar{-/$-\infty$,+/$3$,-/$-1$,+/$+\infty$}
	\end{tikzpicture}
	}	
	\loigiai{Từ bảng biến thiên của $y=f(x)$ ta suy ra bảng biến thiên của $y=\left| f(x)\right| $ như sau
		\begin{center}
			\begin{tikzpicture}
				\tkzTabInit[nocadre=false,lgt=1.2,espcl=2.5,deltacl=0.6]
				{$x$ /0.6, $y$ /2.5}
				{$-\infty$,$x_1$,$-1$,$x_2$,$3$,$x_3$,$+\infty$}
				\tkzTabVar{+/$+\infty$,-/$0$,+/$3$,-/$0$,+/$1$,-/$0$,+/$+\infty$}
			\end{tikzpicture}
		\end{center}	
		(trong đó $x_1$, $x_2$, $x_3$ là nghiệm của phương trình $f(x)=0$).\\
		Vậy số điểm cực trị của hàm số $y=\left| f(x)\right| $ là $5$.	
	}
\end{bt}

\begin{bt}%[2D1V2-2]
% nguon: Archive/2024/TaiLieuGiangDay/1.HamSo/1. K12 - 4in1 - Hamso/ThubaiGoc/Bai2-Dang2.tex #2631
%[BG12, Tran Tony]
	\immini{
		Đạo hàm $f'(x)$ của hàm số $y=f(x)$ có đồ thị như hình vẽ bên. Tìm điểm cực trị của hàm số $y=f(x)$.
	}{
		\begin{tikzpicture}[scale=0.6, font=\footnotesize, line join=round, line cap=round, >=stealth]
			\def\a{0.125} \def\b{-0.625} \def\c{0.25} \def\d{1}
			\def\xt{-3} \def\xp{6} \def\yt{2.5} \def\yd{-3}
			\draw[->] (\xt,0)--(\xp,0) node [below]{$x$};
			\draw[->] (0,\yd)--(0,\yt) node [left]{$y$};
			\node at (0,0) [below right]{$O$};
			\clip (\xt+0.1,\yd-1) rectangle (\xp-0.1,\yt-0.1);
			\draw[smooth,samples=300,domain=-2:5] plot(\x,{\a*(\x)^3+\b*(\x)^2+\c*(\x)+\d}) node [below left]{$y=f'(x)$};
			\draw[dashed] (-2,0) node [below left]{$-2$} -- (-2,-3);
			\draw[dashed] (5,0) node [below]{$5$} -- (5,2.25);
			\fill (-1,0) circle (1.5pt) node[below right]{$-1$};
			\fill (2,0) circle (1.5pt) node[below left]{$2$};
			\fill (4,0) circle (1.5pt) node[below right]{$4$};
		\end{tikzpicture}
	}
	\loigiai{
		Từ đồ thị ta thấy $f'(x)=0 \Leftrightarrow \hoac{& x=-1\\& x=2\\& x=4.}$\\
		Ta có bảng biến thiên như sau
		\begin{center}
			\begin{tikzpicture}
				\tkzTabInit[lgt=1.2,espcl=2.5,deltacl=0.6]
				{$x$ /0.6,$f'(x)$ /0.6,$f(x)$ /2}
				{$-2$,$-1$,$2$,$4$,$5$}
				\tkzTabLine{,-,$0$,+,$0$,-,$0$,+,}
				\tkzTabVar{+/,-/,+/,-/,+/}
			\end{tikzpicture}
		\end{center}
		Hàm số đạt cực đại tại $x=2$ và đạt cực tiểu tại $x=-1$, $x=4$.
	}
\end{bt}

\begin{bt}%[2D1N1-4]
% nguon: Archive/2024/TaiLieuGiangDay/1.HamSo/1. K12 - 4in1 - Hamso/ThubaiGoc/Bai1-Dang4.tex #2643
%%[Mức độ 1]%[Dự án bài giảng new 4in1, Trần Quang Thạnh]
	Dùng phương pháp hàm số, giải phương trình $2^{x^2-x}=2^{x-1}$.
	\loigiai{
	Vì hàm số $y=f(t)=2^t$ đồng biến trên $\mathbb{R}$ nên từ phương trình đã cho ta có 
	$$f(x^2-x)=f(x-1) \Leftrightarrow x^2-x=x-1 \Leftrightarrow x=1.$$
	Vậy phương trình có nghiệm duy nhất $x=1$.
	}
\end{bt}

\begin{bt}%[2D1N1-4]
% nguon: Archive/2024/TaiLieuGiangDay/1.HamSo/1. K12 - 4in1 - Hamso/ThubaiGoc/Bai1-Dang4.tex #2644
%%[Mức độ 1]%[Dự án bài giảng new 4in1, Trần Quang Thạnh]
	Dùng phương pháp hàm số, giải phương trình $\log_{3} (x^2+1) =\log_3 (3-x)$.
	\loigiai{
	Điều kiện $x<3$.\\
	Vì hàm số $y=f(t)=\log_{3} t$ đồng biến trên tập xác định của nó nên từ phương trình đã cho, ta có 
	$$f(x^2+1)=f(3-x) \Leftrightarrow x^2+1 =3-x \Leftrightarrow x^2+x-2=0 \Leftrightarrow \hoac{&x=1\\ &x=-2.}$$
	Đối chiếu với điều kiện, ta có tập nghiệm phương trình là $S=\{1;-2\}$.
	}
\end{bt}

\begin{bt}%[2D1H1-4]
% nguon: Archive/2024/TaiLieuGiangDay/1.HamSo/1. K12 - 4in1 - Hamso/ThubaiGoc/Bai1-Dang4.tex #2645
%%[Mức độ 2]%[Dự án bài giảng new 4in1, Trần Quang Thạnh]
	Giải phương trình $3^x - 3^{x^2-2x} =x^2-3x$.
	\loigiai{
	Phương trình được biến đổi trở thành $3^x + x = 3^{x^2-2x}+x^2-2x$.\\
	Vì hàm số $y=f(t)=3^t$ đồng biến trên $\mathbb{R}$ nên từ phương trình trên ta được
	$$f(x)=f(x^2-2x) \Leftrightarrow x= x^2-2x \Leftrightarrow \hoac{&x=0\\ &x=3.}$$
	Vậy phương trình đã cho có tập nghiệm là $S=\{0;3\}$.
	}
\end{bt}

\begin{bt}%[2D1H1-4]
% nguon: Archive/2024/TaiLieuGiangDay/1.HamSo/1. K12 - 4in1 - Hamso/ThubaiGoc/Bai1-Dang4.tex #2646
%%[Mức độ 2]%[Dự án bài giảng new 4in1, Trần Quang Thạnh]
	Giải bất phương trình $5^x-5^{2x+1}>x+1$.
	\loigiai{
		Phương trình đã cho tương đương với $5^x + x >5^{2x+1}+2x+1$.\\
		Vì hàm số $y=f(t)=5^t+t$ có $y'=5^t\ln5 +1>0$, với mọi $t\in \mathbb{R}$ nên hàm số $y=f(t)$ đồng biến trên $\mathbb{R}$, do đó, từ bất phương trình đã cho ta được $$f(x)>f(2x+1) \Leftrightarrow x>2x+1 \Leftrightarrow x<-1.$$
		Vậy tập nghiệm của bất phương trình đã cho là $S=(-\infty;-1)$.
	}
\end{bt}

\begin{bt}%[2D1C1-4]
% nguon: Archive/2024/TaiLieuGiangDay/1.HamSo/1. K12 - 4in1 - Hamso/ThubaiGoc/Bai1-Dang4.tex #2647
%%[Mức độ 2]%[Dự án bài giảng new 4in1, Trần Quang Thạnh]
	Giải phương trình $\log_2 \dfrac{(x+1)^2}{2x} +x^2+1=0$.
	\loigiai{
		Điều kiện $x>0$.\\
		Ta thực hiện biến đổi như sau
		$$\log_2(x+1)^2 -\log_2 2x +(x+1)^2-2x=0 \Leftrightarrow \log_2(x+1)^2 +(x+1)^2 =\log_2 2x +2x.$$
		Vì hàm số $f(x)=\log_{2} t$ đồng biến trên khoảng $(0;+\infty)$ nên từ phương trình trên ta được
		$$f\left((x+1)^2\right)=f(2x) \Leftrightarrow (x+1)^2=2x \Leftrightarrow x^2+1 =0.$$
		Vậy phương trình đã cho vô nghiệm.
	}
\end{bt}

\begin{bt}%[2D1H1-4]
% nguon: Archive/2024/TaiLieuGiangDay/1.HamSo/1. K12 - 4in1 - Hamso/ThubaiGoc/Bai1-Dang4.tex #2648
%%[Mức độ 2]%[Dự án bài giảng new 4in1, Trần Quang Thạnh]
	Giải phương trình $x+\log (x-3)=4$.
	\loigiai{
		Điều kiện xác định $x>3 $.\\
		Xét hàm số $ f(x)=x+\log (x-3) $ trên khoảng $ (3;+\infty ) $.\\
		Ta có cơ số $ 10>1 $ và hàm số $ y=x $ đồng biến nên $ f(x) $ đồng biến trên khoảng $ (3;+\infty) $.\\
		Mà $f(x)=f(4)=4\Leftrightarrow x=4 $ (thỏa mãn điều kiện xác định).\\
		Vậy $ x=4 $ là nghiệm duy nhất của phương trình.
	}
\end{bt}

\begin{bt}%[2D1H1-4]
% nguon: Archive/2024/TaiLieuGiangDay/1.HamSo/1. K12 - 4in1 - Hamso/ThubaiGoc/Bai1-Dang4.tex #2649
%%[Mức độ 2]%[Dự án bài giảng new 4in1, Trần Quang Thạnh]
	Giải phương trình $ \log_{3}(x+1)+\log_{2}x=2$.
	\loigiai{
		Điều kiện xác định $ x>0$.\\
		Xét hàm số $f(x)=\log_{3}(x+1)+\log_{2}x $ trên khoảng $\left(0;+\infty \right) $.\\
		Với $x>0$, ta có $f'(x)=\dfrac{1}{x\ln3}+\dfrac{1}{x\ln2}>0$, $\forall x>0$ nên hàm số $f(x)$ đồng biến trên $(0;+\infty)$.\\
		Mà $f(2)=2$ nên từ phương trình ta có $f(x)=f(2)\Leftrightarrow x=2 $ (thỏa mãn điều kiện xác định).\\
		Vậy $ x=2 $ là nghiệm của phương trình.
	}
\end{bt}

\begin{bt}%[2D1H1-4]
% nguon: Archive/2024/TaiLieuGiangDay/1.HamSo/1. K12 - 4in1 - Hamso/ThubaiGoc/Bai1-Dang4.tex #2653
[Mức độ 2]%[Dự án bài giảng new 4in1, Trần Quang Thạnh]
	Giải phương trình $5^x-5^{2x+1}=x+1$.
	\loigiai{
	Phương trình đã cho tương đương với $5^x + x =5^{2x+1}+2x+1$.\\
	Vì hàm số $y=f(t)=5^t+t$ có $y'=5^t\ln5 +1>0$, với mọi $t\in \mathbb{R}$ nên hàm số $y=f(t)$ đồng biến trên $\mathbb{R}$, do đó, từ phương trình đã cho ta được $$f(x)=f(2x+1) \Leftrightarrow x=2x+1 \Leftrightarrow x=-1.$$
	Vậy phương trình đã cho có nghiệm duy nhất $x=-1$.
	}
\end{bt}

\begin{bt}%[2D1H1-4]
% nguon: Archive/2024/TaiLieuGiangDay/1.HamSo/1. K12 - 4in1 - Hamso/ThubaiGoc/Bai1-Dang4.tex #2654
[Mức độ 2]%[Dự án bài giảng new 4in1, Trần Quang Thạnh]
	Giải phương trình $3^x - 3^{x^2-2x} >x^2-3x$.
	\loigiai{
		Phương trình được biến đổi trở thành $3^x + x > 3^{x^2-2x}+x^2-2x$.\\
		Vì hàm số $y=f(t)=3^t$ đồng biến trên $\mathbb{R}$ nên từ phương trình trên ta được
		$$f(x)>f(x^2-2x) \Leftrightarrow x> x^2-2x \Leftrightarrow x^2-3x <0 \Leftrightarrow 0<x<3.$$
		Vậy bất phương trình đã cho có tập nghiệm là $S=(0;3)$.
	}
\end{bt}

\begin{bt}%[2D1H1-4]
% nguon: Archive/2024/TaiLieuGiangDay/1.HamSo/1. K12 - 4in1 - Hamso/ThubaiGoc/Bai1-Dang4.tex #2655
[Mức độ 2]%[Dự án bài giảng new 4in1, Trần Quang Thạnh]
	Giải bất phương trình $\log x  > \log (2x-1) + x-1$.
	\loigiai{
	Điều kiện $x>\dfrac{1}{2}$.\\
	Bất phương trình đã cho tương đương với
	$\log x + x > \log(2x-1) +2x-1$.\\
	Vì hàm số $y=f(t)=\log t + t$ có $y'=\dfrac{1}{t\ln10}+1>0$, với mọi $t>0$ nên hàm số $y=f(t)$ đồng biến trên $(0;+\infty)$.\\
	Do đó, từ bất phương trình đã cho ta được
	$$f(x) > f(2x-1) \Leftrightarrow x>2x-1 \Leftrightarrow x<1.$$
	Kết hợp điều kiện, ta có tập nghiệm bất phương trình là $S=\left(\dfrac{1}{2};1\right)$.
	}
\end{bt}

\begin{bt}%[2D1H1-4]
% nguon: Archive/2024/TaiLieuGiangDay/1.HamSo/1. K12 - 4in1 - Hamso/ThubaiGoc/Bai1-Dang4.tex #2656
[Mức độ 2]%[Dự án bài giảng new 4in1, Trần Quang Thạnh]
	Giải phương trình $\log_2 \dfrac{(x+1)^2}{2x} +x^2+1>0$.
	\loigiai{
		Điều kiện $x>0$.\\
		Ta thực hiện biến đổi như sau
		$$\log_2(x+1)^2 -\log_2 2x +(x+1)^2-2x>0 \Leftrightarrow \log_2(x+1)^2 +(x+1)^2 >\log_2 2x +2x.$$
		Vì hàm số $f(x)=\log_{2} t$ đồng biến trên khoảng $(0;+\infty)$ nên từ phương trình trên ta được
		$$f\left((x+1)^2\right)>f(2x) \Leftrightarrow (x+1)^2>2x \Leftrightarrow x^2+1 >0.$$
		Kết hợp điều kiện, ta có tập nghiệm bất phương trình là $S=(0;+\infty)$.
	}
\end{bt}

\begin{bt}%[2D1V1-4]
% nguon: Archive/2024/TaiLieuGiangDay/1.HamSo/1. K12 - 4in1 - Hamso/ThubaiGoc/Bai1-Dang4.tex #2657
[Mức độ 3]%[Dự án bài giảng new 4in1, Trần Quang Thạnh]
	Giải bất phương trình $ \log_2(2^x+1)+\log_3(4^x+2)\le 2$.\hfill $ (*) $
	\loigiai{
		\begin{itemize}
			\item Điều kiện $ \heva{&2^x+1>0\\&4^x+2>0} $ đúng $ \forall x\in \mathbb{R}\Rightarrow  $ Tập xác định $ \mathscr{D}=\mathbb{R} $.
			\item Xét hàm số $ f(x)=\log_2(2^x+1)+\log_3(4^x+2) $ trên $ \mathbb{R} $.\\
			Ta có $ f'(x)=\dfrac{2^x\ln2}{(2^x+1)\ln2}+\dfrac{4^x\ln4}{(4^x+1)\ln3}>0,\forall x\in \mathbb{R}\Rightarrow $ Hàm số $ y=f(x)  $ luôn đồng biến trên $ \mathbb{R} $.\\
			Ta có $ f(0)=2 $ do đó $ \forall x\le 0 $ thì $ f(x)\le f(0)\Leftrightarrow x\le0 $.
			\item Vậy tập nghiệm của bất phương trình là $ x\in (-\infty;0] $.
	\end{itemize}}
\end{bt}

\begin{bt}%[2D1N1-5]
% nguon: Archive/2024/TaiLieuGiangDay/1.HamSo/1. K12 - 4in1 - Hamso/ThubaiGoc/Bai1-Dang5.tex #2668
[SGK12-CTST, Mức độ 1]%[BG12-4in1, Trần Thiên Đức]
	\immini{Trong $8$ phút kể từ khi xuất phát, độ cao $h$ (tính bằng mét) của khinh khí cầu vào tời điểm $t$ phút được cho bởi công thức $h(t)=6t^3-81t^2+324t$. Đồ thị của hàm số $h(t)$ được biểu diễn trong hình bên. Trong các khoảng thời gian nào khinh khí cầu tăng dần độ cao, giảm dần độ cao? Độ cao của khinh khí cầu vào các thời điểm $3$ phút và $6$ phút sau khi xuất phát có gì đặc biệt?
	}{
		\begin{tikzpicture}[scale=0.7, font=\footnotesize, line join=round, line cap=round, >=stealth]
			\def\xmin{-2} \def\xmax{9}
			\def\ymin{-1} \def\ymax{6} 
			\draw[->] (\xmin,0)--(\xmax,0) node [below]{$t$};
			\draw[->] (0,\ymin)--(0,\ymax) node [left]{$h$};
			\clip (\xmin+0.1,\ymin+0.1) rectangle (\xmax-0.1,\ymax-0.1);
			\draw[smooth,samples=300,domain=0:7.3,black] plot(\x,{0.07*(\x)^3-0.89*(\x)^2+3.39*(\x)});
			\fill (0,0)circle (1pt)node[below left]{$O$};
			\fill (0,3.35)circle(1pt )node[left]{$324$} (0,4.05)circle(1pt )node[left]{$405$} (0,4.6)circle(1pt )node[left]{$480$}
			(3,0)circle(1pt )node[below]{$3$}
			(5.6,0)circle(1pt )node[below]{$6$}
			(7.3,0)circle(1pt )node[below]{$8$}
			;
			\draw[dashed] (3,0)--(3,4.05)--(0,4.05) (5.6,0)--(5.6,3.35)--(0,3.35) (7.3,0)--(7.3,4.6)--(0,4.6);
			\fill (3,4.05)circle (1pt)
			(5.6,3.35)circle (1pt)
			(7.3,4.6)circle (1pt)
			;
		\end{tikzpicture}
	}
	\loigiai{
		Dựa vào đồ thị ta thấy khinh khí cầu tăng dần độ cao trong khoảng thời gian $3$ phút đầu và từ phút thứ $6$ đến phút thứ $8$, khinh khí cầu giảm dần độ cao trong khoảng thời gian từ phút thứ $3$ đến phút thứ $6$.\\
		Độ cao của khinh khí cầu tại thời điểm phút thứ $3$ là lớn nhất trong khoảng thời gian $6$ phút đầu, độ cao của khinh khí cầu tại thời điểm phút thứ $6$  là nhỏ nhất trong khoảng thời gian từ phút thứ $3$ đến phút thứ $8$.
		
	}
\end{bt}

\begin{bt}%[2D1N1-5]
% nguon: Archive/2024/TaiLieuGiangDay/1.HamSo/1. K12 - 4in1 - Hamso/ThubaiGoc/Bai1-Dang5.tex #2669
[Mức độ 1]%[BG12-4in1, Trần Thiên Đức]
	Một chất điểm chuyển động theo phương trình $S=-t^3+9t^2+t+10$ trong đó $t$ tính bằng (s) và $S$ tính bằng (m). Xác định khoảng thời gian mà vận tốc của chất điểm tăng.
	\loigiai 
	{
		Ta có $v=S'=-3 t^2+18 t+1$.\\
		$v'=-6t+18$.\\
		$v'=0\Leftrightarrow -6t+18=0 \Leftrightarrow t=3$.\\
		Bảng biến thiên 
		\begin{center}
			\begin{tikzpicture}[scale=1, font=\footnotesize]%<DTools>
				\tkzTabInit[nocadre=false, lgt=1.2, espcl=4, deltacl=0.6]
				{$t$/0.8,$v'(t)$/0.6,$v(t)$/2}
				{$0$,$3$,$+\infty$};
				\tkzTabLine{,+,$0$,-,};
				\tkzTabVar{-/$1$,+/$28$,-/$-\infty$};
			\end{tikzpicture}
		\end{center}
		Vậy trong khoảng thời gian từ $0$ đến $3$ s thì vận tốc của chất điểm tăng.
	}
\end{bt}

\begin{bt}%[2D1H1-5]
% nguon: Archive/2024/TaiLieuGiangDay/1.HamSo/1. K12 - 4in1 - Hamso/ThubaiGoc/Bai1-Dang5.tex #2670
[SGK 12 - Cùng Khám Phá, Mức độ 2]
Thể tích $V$ của $1$ kg nước (tính bằng $\text{cm}^3$) ở nhiệt độ $T$ (đơn vị: $^\circ$C) khi $T$ thay đổi từ $0^\circ$C đến $30^\circ$C được cho xấp xỉ bởi công thức:
	\[V=999{,}87-0{,}06426T+0{,}0085043T^2-0{,}0000769T^3.\]
	(Nguồn: James Stewart, J. (2015). Calculus.Cengage Learning 8th edition, p.284)\\ Tìm nhiệt độ $T_0 \in \left(0;30\right)$ để kể từ nhiệt độ $T_0$ trở lên thì thể tích $V$ tăng (làm tròn kết quả đến hàng đơn vị).
	\loigiai{Xét hàm số $f\left(T\right)=V=999{,}87-0{,}06426T+0{,}0085043T^2-0{,}0000769T^3$.\\ Hàm số xác định trên khoảng $\left(0;30\right)$.\\
		Ta có $f'\left(T\right)=2{,}307 \cdot 10^{-4}T^2+0{,}0170086T-0{,}06426$.\\
		$f'\left(T\right)=0\Leftrightarrow \hoac{&x=3{,}6\\&x=-77{,}33}\Leftrightarrow x=3{,}6$.\\
	Khi đó, ta có bảng biến biên
		\begin{center}
			\begin{tikzpicture}
				\tkzTabInit[lgt=1.5,espcl=3]
				{$T$ /.7, $f’(T)$ /.7, $f(T)$ /2.5}
				{$0$,$3{,}6$,$30$}
				\tkzTabLine{ ,-,z,+,}
				\path
				(N12)node[below](1){$f(0)$}
				(N23)node[above](2){$f(3{,}6)$}
				(N32)node[below](3){$f(30)$}
				;
				\foreach \x/\y in {1/2,2/3}
				\draw[-stealth] (\x)--(\y);
			\end{tikzpicture}
		\end{center}
		Vậy $T_0=4^\circ$C trở lên thì nhiệt độ tăng.
	}
\end{bt}

\begin{bt}%[2D1H1-5]
% nguon: Archive/2024/TaiLieuGiangDay/1.HamSo/1. K12 - 4in1 - Hamso/ThubaiGoc/Bai1-Dang5.tex #2671
[SGK 12 - Cùng Khám Phá, mức độ 2]
	Một công ty tiến hành khai thác $17$ giếng dầu trong
	khu vực được chỉ định. Trung bình mỗi giếng dầu
	chiết xuất được $245$ thùng dầu mỗi ngày. Công ty
	có thể khai thác nhiều hơn $17$ giếng dầu nhưng
	cứ khai thác thêm một giếng thì lượng dầu mỗi
	giếng chiết xuất được hằng ngày sẽ giảm $9$ thùng.
	Để giám đốc công ty có thể quyết định số giếng
	cần thêm cho phù hợp với tài chính, hãy chỉ ra
	số giếng công ty có thể khai thác thêm để sản lượng
	dầu chiết xuất tăng lên.
	\loigiai{
		Gọi số giếng công ty có thể khai thác thêm là $x$ (giếng, $x\in \mathbb{N}$).\\
		Tổng số giếng công ty khai thác hàng ngày là $17+x$ (giếng).\\
		Số thùng chiết xuất được từ mỗi giếng hàng ngày là $245-9x$ (thùng).\\
		Sản lượng dầu chiết xuất hàng ngày là 
		$$
		f(x)=(245-9x)\cdot (17+x)=-9x^2+92x+4165.
		$$
		Ta có $f'(x)=-18x+92$.\\
		$f'(x)=0\Leftrightarrow x=\dfrac{46}{9}\approx 5{,}1$.\\
		Bảng biến thiên
		\begin{center}
			\begin{tikzpicture}[scale=1, font=\footnotesize]%<DTools>
				\tkzTabInit[nocadre=false, lgt=1.2, espcl=2, deltacl=0.6]
				{$x$/0.8,$f'(x)$/0.6,$f(x)$/2}
				{$0$,$\dfrac{46}{9}$,$+\infty$};
				\tkzTabLine{,+,$0$,-,};
				\tkzTabVar{-/$4\,165$,+/$\dfrac{39\,601}{9}$,-/$-\infty$};
			\end{tikzpicture}
		\end{center}
		Ta có $f(5)=4\,400$ và $f(6)=4393$. \\
		Vậy số giếng công ty có thể khai thác thêm để sản lượng dầu chiết suất tăng lên là $5$ giếng.
	}
\end{bt}

\begin{bt}%[2D1H1-5]
% nguon: Archive/2024/TaiLieuGiangDay/1.HamSo/1. K12 - 4in1 - Hamso/ThubaiGoc/Bai1-Dang5.tex #2672
[SGK12 - Cánh Diều, Mức độ 2]
	\immini{
		Máng trượt của một cầu trượt cho trẻ em (\textit{Hình a}) được uốn từ một tấm kim loại có bề rộng $80$ cm, mặt cắt được mô tả ở \textit{Hình b}. Nhà thiết kế khuyến cáo, diện tích mặt cắt càng lớn thì càng đảm bảo an toàn cho trẻ em.
		\begin{enumerate}[a)]
			\item Gọi $S$ là diện tích mặt cắt. Tìm điều kiện của $x$ và viết công thức tính $S$ theo $x$.
			\item Diện tích mặt cắt sẽ tăng hay giảm khi chúng ta tăng giá trị của $x$ từ 0 cm đến 10 cm?
		\end{enumerate}
	}{
		\begin{tikzpicture}[scale=1, font=\footnotesize, line join=round, line cap=round, >=stealth,line width=3.pt]
			\clip(-0.2,-1.2) rectangle (7.5,2.6);
			\fill[line width=3.pt,color=black,fill=black,fill opacity=0.7] (0,0)--(0.41,0.34)--(1,1)--(1.76,1.76)--(2,2)--(2.46,2.08)--(2.83,1.95)--(2.34,1.65)--(1.34,0.65)--(0.75,-0.01)--(0.36,-0.13)-- cycle;
			\draw (0,0)--(0.41,0.34)--(1,1)--(2,2)--(2.46,2.08);
			\draw (0,0)--(0.36,-0.13)--(0.77,0.21)--(0.77,0.21)--(1.36,0.87)--(2.36,1.87)--(2.83,1.95);
			\draw (0.36,-0.13)--(0.75,-0.01)--(1.34,0.65)--(2.34,1.65);
			\draw (2.36,0.68)-- (3.95,0.28);
			\draw (3.42,0.41)--(2.68,2.29)--(2.53,2.37)--(2.16,1.88)--(2.31,1.53);
			\draw (1.63,2.01)--(2,2.5)-- (2.16,2.42)--(2.89,0.55);
			\draw (1.63,2.01)--(1.76,1.76) (2.73,0.98)--(3.26,0.85) (2.52,1.49)--(3.05,1.35);
			\draw (2.34,1.65)--(2.83,1.95)--(2.46,2.08);
			\draw (2,-0.8)node{a)};
			\fill[line width=3.pt,fill=gray,fill opacity=0.3] (5,0.5)--(5,0)--(6.5,0)--(6.5,0.5)--cycle;
			\draw (5,0.5)--(5,0) node [left ,sloped,pos=0.5,rotate=90] {$x$ cm};
			\draw (5,0)--(6.5,0) node [below,sloped,pos=0.5] {$y$ cm};
			\draw (6.5,0)--(6.5,0.5) node [right,sloped,pos=0.5,rotate=-90] {$x$ cm};
			\draw[dashed,thin] (5,0.5)--(6.5,0.5);
			\draw (5.8,-0.8)node{b)};
			%			\draw (4,-1)node{\textit{Hình 5}};
		\end{tikzpicture}
	}
	\loigiai{
		\begin{enumerate}[a)]
			\item Do tấm kim loại có bề rộng $80 \mathrm{~cm}$ nên ta có $2x+y=80 \Leftrightarrow y=80-2x$.\\
			Để có thể thiết kế được máng trượt thì $y>0 \Leftrightarrow 80-2 x>0 \Leftrightarrow x<40$. Suy ra $0<x<40$.\\
			Diện tích của mặt cắt máng trượt là: $S=x y=x(80-2x)=-2x^2+80x$.		
			\item Ta có $S(x)=-2x^2+80x \text{ với } x\in(0; 40)$.
			\allowdisplaybreaks
			\begin{eqnarray*}
				S'(x)&=&-4 x+80 \\
				S'(x)&=&0 \Leftrightarrow-4 x+80=0 \Leftrightarrow x=20.
			\end{eqnarray*}
			Bảng biến thiên của hàm số $S(x)$ như sau
			\begin{center}
				\begin{tikzpicture}[font=\footnotesize,thick,>=stealth]
					\tikzset{double style/.append style={double distance=1.5pt}}
					\tkzTabInit[nocadre=false,lgt=1.2,espcl=2.5,deltacl=0.6,lw=.75pt]
					{$x$ /0.7, $S'(x)$ /0.8, $S(x)$ /2}
					{$0$,$20$,$40$}
					\tkzTabLine{ ,+,,-, }
					\tkzTabVar{-/$0$,+/$800$,-/$0$}
				\end{tikzpicture}
			\end{center}
			Vậy từ khi $x$ tăng từ $0$ cm đến $20$ cm thì diện tích mặt cắt sẽ tăng dần.
		\end{enumerate}
	}
\end{bt}

\begin{bt}%[2D1H1-5]
% nguon: Archive/2024/TaiLieuGiangDay/1.HamSo/1. K12 - 4in1 - Hamso/ThubaiGoc/Bai1-Dang5.tex #2674
[SGK12-CTST, Mức độ 2]
	Kim ngạch xuất khẩu rau quả của Việt Nam trong các năm từ $2010$ đến $2017$ có thể được tính xấp xỉ bằng công thức $f(x)=0{,}01x^3-0{,}04x^2+0{,}25x+0{,}44$ (tỉ USD) với $x$ là số năm tính từ $2010$ đến $2017$ $(0\leq x\leq 7)$.
	\begin{flushright}
		(\textit{Theo:} https://infographics.vn/interactive-xuat-khau-rau-qua-
		du-bao-bung-no-dat-4-ty-usd-trong-nam-2023/116220.vna)
	\end{flushright}
	\begin{enumerate}
		\item Tính đạo hàm của hàm số $y=f(x)$.
		\item Chứng minh rằng kim ngạch xuất khẩu rau quả của Việt Nam tăng liên tục trong các năm từ $2010$ đến $2017$.
	\end{enumerate}
	\loigiai{
		\begin{enumerate}
			\item Ta có $f'(x)=0{,}03x^2-0{,}08x+0{,}25$.
			\item Xét $f'(x)=0 \Leftrightarrow 0{,}03x^2-0{,}08x+0{,}25=0$ (vô nghiệm).\\
			Bảng biến thiên
			\begin{center}
				\begin{tikzpicture}
					\tkzTabInit[nocadre=false,lgt=1.2,espcl=5,deltacl=0.6]
					{$x$ /0.6,$f'(x)$ /0.6,$f(x)$ /2}
					{$0$,$7$}
					\tkzTabLine{,+,}
					\tkzTabVar{-/$0{,}44$,+/$3{,}66$}
				\end{tikzpicture}
			\end{center}
			Từ bảng biến thiên trên, ta thấy $f'(x)>0$, $\forall x\in [0;7]$.\\
			Vậy kim ngạch xuất khẩu rau quả của Việt Nam tăng liên tục trong các năm từ $2010$ đến $2017$.
		\end{enumerate}
	}
\end{bt}

\begin{bt}%[2D1H1-5]
% nguon: Archive/2024/TaiLieuGiangDay/1.HamSo/1. K12 - 4in1 - Hamso/ThubaiGoc/Bai1-Dang5.tex #2678
[Mức độ 2]%
	Định luật Moore được xây dựng bởi Gordon More-một trong những sáng lập viên của tập đoàn sản xuất chip máy tính nổi tiếng Intel. Định luật ban đầu được phát biểu như sau
	\lq\lq Số lượng transistor trên mỗi đơn vị inch vuông sẽ tăng lên gấp đôi sau mỗi $24$ tháng. \rq\rq ($1$ inch vuông xấp xỉ $6{,}45$ cm$^2$).
	Phương trình sau đây mô tả định luật Moore
	$$n_i=n_0\cdot 2^{\tfrac{y_i-y_0}{2}}$$
	trong đó $n_0$ là số lượng transitor trên một đơn vị diện tích ban đầu, $y_0$ là năm ứng với $n_0$, $n_i$ là số lượng transitor tại năm $y_i$ đang xét. \\
	(Nguồn: https://scipython.com/book/chapter-3-simple-plotting-with-pylab/examples/moores-law/)\\
	Nhận xét về tính đơn điệu của hàm số mô tả định luật Moore.
	\loigiai{
		Xét hàm số $n_i=n_0\cdot 2^{\tfrac{y_i-y_0}{2}}$.\\
		Ta có $n_i'=n_0\cdot 2^{\tfrac{y_i-y_0}{2}}\cdot \ln 2>0~\forall y_i$.\\
		Suy ra hàm số đồng biến.
	}
\end{bt}

\begin{bt}%[2D1H1-5]
% nguon: Archive/2024/TaiLieuGiangDay/1.HamSo/1. K12 - 4in1 - Hamso/ThubaiGoc/Bai1-Dang5.tex #2679
[SGK12-CTST, Mức độ 2]
	Xét một chất điểm chuyển động dọc theo trục $Ox$. Toạ độ của chất điểm tại thời điểm $t$ được xác định bởi hàm số $x(t)=t^3-6t^2+9t$ với $t\geq 0$. Khi đó $x'(t)$ là vận tốc của chất điểm tại thời điểm $t$, kí hiệu $v(t)$. Trong khoảng thời gian nào vận tốc của chất điểm tăng, trong khoảng thời gian nào vận tốc của chất điểm giảm?
	\loigiai{
		Ta có $v(t)=x'(t)=3t^2-12t+9$ và $v'(t)=6t-12$.\\
		Xét $v'(t)=0 \Leftrightarrow t=2$.\\
		Bảng biến thiên
		\begin{center}
			\begin{tikzpicture}[scale=1, font=\footnotesize]%<DTools>
				\tkzTabInit[nocadre=false, lgt=1.2, espcl=4, deltacl=0.6]
				{$x$/0.8,$f'(x)$/0.6,$f(x)$/2}
				{$-\infty$,$2$,$+\infty$};
				\tkzTabLine{,-,$0$,+,};
				\tkzTabVar{+/$9$,-/$-3$,+/$+\infty$};
			\end{tikzpicture}
		\end{center}
		Suy ra vận tốc của chất điểm giảm khi $t\in (0;2)$, và tăng khi $t\in (2;+\infty)$.
	}
\end{bt}

\begin{bt}%[2D1H1-5]
% nguon: Archive/2024/TaiLieuGiangDay/1.HamSo/1. K12 - 4in1 - Hamso/ThubaiGoc/Bai1-Dang5.tex #2680
[SGK12-KNTT, Mức độ 2]%[BG12-4in1, Trần Thiên Đức]
	Xét một chất điểm chuyển động trên một trục số nằm ngang, chiều dương từ trái sang phải (Hình vẽ). Giả sử vị trí $s(t)$ (mét) của chất điểm trên trục số đã chọn tại thời điểm $t$ (giây) được cho bởi công thức
	\[s(t)=t^3-9t^2+15t,\,\, t\ge 0.\]
	Hỏi trong khoảng thời gian nào thì chất điểm chuyển động sang phải, trong khoảng thời gian nào  thì chất điểm chuyển động sang trái?
	\begin{center}
		\begin{tikzpicture}[>=stealth,line join=round,line cap=round,font=\footnotesize,scale=1]
			\draw[->] (0,0)coordinate (O)--(5,0)coordinate (B);
			\draw[thick] 
			(1.5,-.05)--(1.5,.05)node[below=1mm,black]{$O$}
			(3,-.05)--(3,.05)node[below=1mm,black]{$s(t)$}
			;
			\path 
			(O)--(B)node[midway,below=6mm,orange]{}
			;
		\end{tikzpicture}
	\end{center}
	\loigiai{
		Theo ý nghĩa cơ học của đạo hàm, vận tốc $v\left(t\right)$ là đạo hàm của $s\left(t\right)$. \\
		Ta có $v\left(t\right) = s'\left(t\right) = \left(t^3 -9t^2 + 15t\right)' = 3t^2 - 18t + 15$.\\
		Tập xác định $\mathscr{D} = \mathbb{R}$.\\
		Ta có $v(t)=0\Leftrightarrow\hoac{&t=1 \\&t=5. }$\\
		Bảng xét dấu
		\begin{center}
			\begin{tikzpicture}[scale=1, font=\footnotesize]
				\tkzTabInit[nocadre=false, lgt=1.2, espcl=2, deltacl=0.6]
				{$t$/0.6,$v(t)$/0.6}
				{$-\infty$,$1$,$5$,$+\infty$};
				\tkzTabLine{,+,$0$,-,$0$,+,};
			\end{tikzpicture}
		\end{center}
		Chất điểm chuyển động theo chiều dương (sang bên phải) khi $v\left(t\right) > 0$, tức là $t \in \left(-\infty; 1\right) \cup \left(5; +\infty\right)$.\\
		Chất điểm chuyển động theo chiều âm (sang bên trái) khi $v\left(t\right) < 0$, tức là $1 < t < 5$.
	}	
\end{bt}

\begin{bt}%[2D1H1-5]
% nguon: Archive/2024/TaiLieuGiangDay/1.HamSo/1. K12 - 4in1 - Hamso/ThubaiGoc/Bai1-Dang5.tex #2681
[Mức độ 2]
	Một vật được phóng thẳng đứng lên trên từ độ cao $2 \mathrm{~m}$ với vận tốc ban đầu là $24{,}5 \mathrm{~m} / \mathrm{s}$. Trong Vật lí, ta biết rằng khi bỏ qua sức cản của không khí thì độ cao $h$ (mét) của vật sau $t$ (giây) được cho bởi công thức
	$$
	h(t)=2+24{,}5 t-4{,}9t^{2}.
	$$
	Hỏi tại từ thời điểm $t=3$ (giây) đến $t=4$ (giây), thì độ cao của vật tăng hay giảm?
	\loigiai{
		Xét hàm số $h(t)=2+24{,}5 t-4{,}9t^{2}$ với $t>0$.\\
		Ta có $h'(t)=24{,}5-9{,}8t$. \\
		$h'(t)=0\Leftrightarrow t=2{,}5$. \\
		Lập bảng biến thiên của $h(t)$: 
		\begin{center}
			\begin{tikzpicture}
				\tkzTabInit[nocadre=false,lgt=1.2,espcl=3.5,deltacl=0.6]
				{$t$ /0.6,$h'(t)$ /0.6,$h(t)$ /1.4}
				{$0$,$2{,}5$,$+\infty$}
				\tkzTabLine{,+,0,-,}
				\tkzTabVar{-/,+/$32{,}625$,-/}
			\end{tikzpicture}
		\end{center}
		Từ bảng biến thiên ta suy ra từ thời điểm $t=3$ (giây) đến $t=4$ (giây), thì độ cao của vật giảm.
	}
\end{bt}

\begin{bt}%[2D1V1-5]
% nguon: Archive/2024/TaiLieuGiangDay/1.HamSo/1. K12 - 4in1 - Hamso/ThubaiGoc/Bai1-Dang5.tex #2682
[Mức độ 3]%
	Giả sử dân số của một thị trấn sau $t$ năm kể từ năm $2000$ được mô tả bởi hàm số $$N\left(t\right)=\frac{25t+10}{t+5}, t\ge0,$$
	trong đó $N\left(t\right)$ được tính bằng nghìn người. Tính đạo hàm $N^\prime\left(t\right)$ và $\lim\limits_{t \to +\infty}N\left(t\right)$. Từ đó giải thích tại sao dân số của thị trấn đó luôn tăng nhưng sẽ không vượt quá một ngưỡng nào đó.
	\loigiai{
		Xét hàm số  $N\left(t\right)=\dfrac{25t+10}{t+5}, t\ge0$.\\
		Ta có $N^\prime\left(t\right)=\dfrac{115}{\left(t+5\right)^2}$.\\
		$\lim\limits_{t \to +\infty}N\left(t\right)=\lim\limits_{t \to +\infty}\dfrac{25t+10}{t+5}=25$.\\
		Bảng biến thiên của hàm số như sau
		\begin{center}
			\begin{tikzpicture}[font=\footnotesize,thick,>=stealth]
				\tikzset{double style/.append style={double distance=1.5pt}}
				\tkzTabInit[nocadre=false,lgt=1.2,espcl=2.5,deltacl=0.6,lw=.75pt]
				{$x$ /0.6,$y'$ /0.6,$y$ /2.5}
				{$-\infty$,$-5$,$+\infty$}
				\tkzTabLine{,+,d,+,}
				\tkzTabVar{-/$25$,+D-/$+\infty$/$-\infty$,+/$25$}
			\end{tikzpicture}
		\end{center}
		Từ bảng biến thiên ta có dân số của thị trấn đó luôn tăng nhưng sẽ không vượt quá $25$ ngàn người.
	}
\end{bt}

\begin{bt}%[2D1V1-1]
% nguon: Archive/2024/TaiLieuGiangDay/1.HamSo/1. K12 - 4in1 - Hamso/ThubaiGoc/Bai1-Dang1.tex #2693
%%[Mức 2]%[BG12new-4in1, Trần Hoà]
		Tìm các khoảng đồng biến, nghịch biến của hàm số $y=\sqrt{x^2-1}$.
		\loigiai{
			Tập xác định: $\mathscr{D}=(-\infty;-1]\cup [1;+\infty)$.\\
			Đạo hàm: $y'=\dfrac{x}{\sqrt{x^2-1}}$ với $x\in (-\infty;-1)\cup (1;+\infty)$.\\
			Cho $y'=0\Leftrightarrow x=0 \notin \mathscr{D}$
			Bảng biến thiên
			\begin{center}
				\begin{tikzpicture}
				\tkzTabInit[lgt=1,espcl=3]
				{$x$/0.7,$y'$/0.7,$y$/2}
				{$-\infty$,$-1$,$1$,$+\infty$}
				\tkzTabLine{,-,d,h,d,+,}
				\tkzTabVar{+/$+\infty$,-H/$0$,-/$0$,+/$+\infty$}
				\end{tikzpicture}
			\end{center}
			Vậy hàm số nghịch biến trên khoảng $(-\infty;-1)$ và đồng biến trên khoảng $(1;+\infty)$.
		}
	\end{bt}

\begin{bt}%[2D1V1-1]
% nguon: Archive/2024/TaiLieuGiangDay/1.HamSo/1. K12 - 4in1 - Hamso/ThubaiGoc/Bai1-Dang1.tex #2694
%%[Mức 2]%[BG12new-4in1, Trần Hoà]
		Tìm các khoảng đồng biến, nghịch biến của hàm số $y=\ln\left(x+\sqrt{9+x^2}\right)$.
		\loigiai{Ta có $x+\sqrt{9+x^2}> x+|x|\ge 0,~\forall x\in\mathbb{R}$.
			Tập xác định: $\mathscr{D}=\mathbb{R}$.\\
			Đạo hàm: $y'=\dfrac{1}{x+\sqrt{9+x^2}}\cdot\left(1+\dfrac{x}{\sqrt{9+x^2}}\right)=\dfrac{1}{\sqrt{9+x^2}}>0,~\forall x\in\mathbb{R}$.\\
			Vậy hàm số nghịch biến trên khoảng $(-\infty;+\infty)$ 
		}
	\end{bt}

\begin{bt}%[2D1H1-1]
% nguon: Archive/2024/TaiLieuGiangDay/1.HamSo/1. K12 - 4in1 - Hamso/ThubaiGoc/Bai1-Dang1.tex #2695
%%[CTST, Mức 2]%[TeX hóa SGK CTST 12]%[Phạm Tiến Long]%[BG12new-4in1, Trần Hoà]
Chứng minh rằng hàm số $g(x)=\dfrac{x}{x-1}$ nghịch biến trên khoảng $(1;+\infty)$.
\loigiai{
Hàm số xác định trên $(1;+\infty)$.\\
Ta có $g'(x)=-\dfrac{1}{(x-1)^2}<0$ với mọi $x\in(1;+\infty)$.\\
Vậy $g(x)$ nghịch biến trên khoảng $(1 ;+\infty)$.
}
\end{bt}

\begin{bt}%[2D1H1-2]
% nguon: Archive/2024/TaiLieuGiangDay/1.HamSo/1. K12 - 4in1 - Hamso/ThubaiGoc/Bai1-Dang1.tex #2696
%%[CTST, Mức 2]%[TeX hóa SGK CTST 12]%[Phạm Tiến Long]%[BG12new-4in1, Trần Hoà]
Chứng minh rằng hàm số $f(x)=3x-\sin x$ đồng biến trên $\mathbb{R}$.
\loigiai{
Tập xác định $\mathscr{D}=\mathbb{R}$.\\
Ta có $f'(x)=3-\cos x$.\\
Vì $-1\le \cos x\le 1$ nên $3-\cos x>0$, $\forall x\in \mathbb{R}$.\\
Suy ra $f'(x)>0$,  $\forall x\in \mathbb{R}$.\\
Vậy hàm số $f(x)=3x-\sin x$ đồng biến trên $\mathbb{R}$.
}
\end{bt}

\begin{bt}%[2D1H1-1]
% nguon: Archive/2024/TaiLieuGiangDay/1.HamSo/1. K12 - 4in1 - Hamso/ThubaiGoc/Bai1-Dang1.tex #2700
%[BG12new-4in1, Trần Hoà]
Tìm các khoảng đơn điệu của hàm số $y=-x^3+6x^2-9x+4$.
\loigiai{
			Hàm số $y=-x^3+6x^2-9x+4$ có tập xác định $\mathscr{D}=\mathbb{R}$.\\
			Ta có $y'=-3x^2+12x-9$. Cho 
				{\allowdisplaybreaks
		\begin{eqnarray*}
		y'&=&0 \Leftrightarrow -3x^2+12x-9=0\\
		&\Leftrightarrow&
		\hoac{&x=1\\&x=3.}
		\end{eqnarray*}}	
Bảng biến thiên
			\begin{center}
				\begin{tikzpicture}
				\tkzTabInit[lgt=1,espcl=3]
				{$x$/0.7,$y'$/0.7,$y$/2}
				{$-\infty$,$1$,$3$,$+\infty$}
				\tkzTabLine{,-,0,+,0,-,}
				\tkzTabVar{+/$+\infty$,-/ $0$/,+/$4$,-/$-\infty$}
				\end{tikzpicture}
			\end{center}
			Vậy hàm số nghịch biến trên mỗi khoảng $(-\infty;1)$, $(3;+\infty)$ và đồng biến trên khoảng $(1;3)$.	
		}
\end{bt}

\begin{bt}%[2D1H1-1]
% nguon: Archive/2024/TaiLieuGiangDay/1.HamSo/1. K12 - 4in1 - Hamso/ThubaiGoc/Bai1-Dang1.tex #2701
%[BG12new-4in1, Trần Hoà]
Tìm các khoảng đơn điệu của hàm số $y=x^4-6x^2+8x+1$.
		\loigiai{
	Hàm số $y=x^4-6x^2+8x+1$ có tập xác định $\mathscr{D}=\mathbb{R}$.\\
	Ta có $y'=4x^3-12x+8=0=4(x-1)^2(x+2)$.
		{\allowdisplaybreaks
		\begin{eqnarray*}
		y'=0 &\Leftrightarrow& 4(x-1)^2(x+2)=0\\ 
		&\Leftrightarrow&
		\hoac{&x=-2\\&x=1.}
		\end{eqnarray*}}			
			Bảng biến thiên
			\begin{center}
				\begin{tikzpicture}
				\tkzTabInit[lgt=1,espcl=3]
				{$x$/0.7,$y'$/0.7,$y$/2}
				{$-\infty$,$-2$,$1$,$+\infty$}
				\tkzTabLine{,-,0,+,0,+,}
				\tkzTabVar{+/$+\infty$,-/ $-23$/,R,+/$+\infty$}
				\end{tikzpicture}
			\end{center}
Vậy hàm số nghịch biến trên khoảng $(-\infty;-2)$ và đồng biến trên khoảng $(-2;+\infty)$.
}
	\end{bt}

\begin{bt}%[2D1H1-1]
% nguon: Archive/2024/TaiLieuGiangDay/1.HamSo/1. K12 - 4in1 - Hamso/ThubaiGoc/Bai1-Dang1.tex #2702
%[BG12new-4in1, Trần Hoà]
		Tìm các khoảng đơn điệu của hàm số $y=\dfrac{3-2x}{x+7}$.
		\loigiai{
			Hàm số $y=\dfrac{3-2x}{x+7}=\dfrac{-2x+3}{x+7}$ có tập xác định $ \mathscr{D}=\mathbb{R}\setminus \{-7\}$.\\
			Ta có $y'=\dfrac{-17}{(x+7)^2}<0,\forall x \neq -7$.\\
			Bảng biến thiên
			\begin{center}
				\begin{tikzpicture}
				\tkzTabInit[lgt=1,espcl=4]
				{$x$/0.7,$y'$/0.7,$y$/2}
				{$-\infty$,$-7$,$+\infty$}
				\tkzTabLine{,-,d,-,}
				\tkzTabVar{+/$-2$,-D+/$-\infty$/$+\infty$,-/$-2$}
				\end{tikzpicture}
			\end{center}
			Hàm số đã cho nghịch biến trên mỗi khoảng $(-\infty;-7)$ và $(-7;+\infty)$.
		}
	\end{bt}

\begin{bt}%[2D1V1-1]
% nguon: Archive/2024/TaiLieuGiangDay/1.HamSo/1. K12 - 4in1 - Hamso/ThubaiGoc/Bai1-Dang1.tex #2703
%[BG12new-4in1, Trần Hoà]
		Tìm các khoảng đơn điệu của hàm số $y=\dfrac{x^2-x+1}{x-1}$.
		\loigiai{
			Hàm số $y=\dfrac{x^2-x+1}{x-1}$ có tập xác định $\mathscr{D}=\mathbb{R}\setminus \{1\}$.\\
			Ta có $y'=\dfrac{x^2-2x}{(x-1)^2},\forall x\in \mathscr{D}$.
		{\allowdisplaybreaks
		\begin{eqnarray*}
		y'=0 &\Leftrightarrow& \dfrac{x^2-2x}{(x-1)^2}=0 \Leftrightarrow x^2-2x=0\\
		& \Leftrightarrow&
		\hoac{&x=0\\&x=-2.}
		\end{eqnarray*}}			
			Bảng biến thiên
			\begin{center}
				\begin{tikzpicture}
				\tkzTabInit[lgt=1,espcl=3]
				{$x$/0.7,$y'$/0.7,$y$/2}
				{$-\infty$,$0$,$1$,$2$,$+\infty$}
				\tkzTabLine{,+,$0$,-,d,-,$0$,+,}
				\tkzTabVar{-/$-\infty$,+/$-1$,-D+/$-\infty$/$+\infty$,-/$3$,+/$+\infty$}
				\end{tikzpicture}
			\end{center}
			Hàm số đồng biến biến trên mỗi khoảng $(-\infty;0)$ và $(2;+\infty)$\\
			Hàm số nghịch biến trên mỗi khoảng $(0;1)$ và $(1;2)$.
		}
	\end{bt}

\begin{bt}%[2D1V1-1]
% nguon: Archive/2024/TaiLieuGiangDay/1.HamSo/1. K12 - 4in1 - Hamso/ThubaiGoc/Bai1-Dang1.tex #2704
%[BG12new-4in1, Trần Hoà]
		Tìm các khoảng đồng biến, nghịch biến của hàm số $y=x+\sqrt{16-x^2}$.
		\loigiai{
			Tập xác định: $\mathscr{D}=[-4;4]$.\\
			Đạo hàm: $y'=1-\dfrac{x}{\sqrt{16-x^2}}=\dfrac{\sqrt{16-x^2}-x}{\sqrt{16-x^2}}$ với $x\in (-4;4)$
				{\allowdisplaybreaks
		\begin{eqnarray*}
		y'=0
		&\Leftrightarrow&
		\heva{&\sqrt{16-x^2}=x\\&x>0}\\
		&\Leftrightarrow&
		\heva{&x>0\\&x^2=8}
		\Leftrightarrow
		x=2\sqrt{2}.
		\end{eqnarray*}}	
			Bảng biến thiên
			\begin{center}
				\begin{tikzpicture}
				\tkzTabInit[lgt=1,espcl=3]
				{$x$/0.7,$y'$/0.7,$y$/2}
				{$-\infty$,$-4$,$2\sqrt{2}$,$4$,$+\infty$}
				\tkzTabLine{,h,d,+,$0$,-,d,h,}
				\tkzTabVar{+H/,-/$-4$,+/$4\sqrt{2}$,-H/$4$,}
				\end{tikzpicture}
			\end{center}
			Vậy hàm số đồng biến trên khoảng $(-4;2\sqrt{2})$ và nghịch biến trên khoảng $(2\sqrt{2};4)$.
		}
	\end{bt}

\begin{bt}%[2D1H2-7]
% nguon: Archive/2024/TaiLieuGiangDay/1.HamSo/1. K12 - 4in1 - Hamso/ThubaiGoc/Bai2-Dang8.tex #2715
[CTST]%[Bài giảng 12-4in1, Nhật Thiện]
	Một phần lát cắt của dãy núi có độ cao tính bằng mét được mô tả bởi hàm số
	$$y=h(x)=-\dfrac{1}{1\,320\,000}x^3+\dfrac{9}{3\,520}x^2-\dfrac{81}{44}x+840, \text { với } 0\leq x\leq 2\,000.$$
	Tìm tọa độ các đỉnh của lát cắt dãy núi trên đoạn $[0; 2\,000]$.
	\begin{center}
		\begin{tikzpicture}[font=\footnotesize, line join=round, line cap=round, >=stealth,scale=0.45,transform shape]
			\def\xmin{-20} \def\xmax{14}
			\def\ymin{-1.25} \def\ymax{7.75}
			\clip (\xmin-.1,\ymin-.1) rectangle (\xmax+.1,\ymax+.1);
			
			\fill[orange!20] (-18,0)..controls +(0:2) and +(178:4) ..(-9.5,1)
			..controls +(2:2) and +(178:4) ..(-4.3,1.5)
			..controls +(45:2) and +(130:2) ..(0,2.3)
			..controls +(-50:.5) and +(160:.5) ..(1.3,1.2)
			..controls +(-20:1.5) and +(150:2) ..(6,4)
			..controls +(-30:1.5) and +(135:2) ..(8.5,1.8)
			..controls +(-40:1.5) and +(155:2) ..(12 ,0)--cycle
			;
			\draw[thick,red] (0,2.3)
			..controls +(-50:.5) and +(160:.5) ..(1.3,1.2)
			..controls +(-20:1.5) and +(150:2) ..(6,4);
			
			\draw[->] (\xmin,0)--(\xmax,0) node [below left]{\Huge $x$};
			\draw[->] (0,\ymin)--(0,\ymax) node [below left]{\Huge $y$};
			\node at (0,0) [below left]{\Huge $O$};
			\node at (4.2,4.3) [thick,red,above]{\Huge $y=h(x)$};
			\draw[dashed](6,4)--(6,0) node[below] {\Huge $2\,000$};
			\foreach \y/\diem in {1.5/500,3/1\,000,4.5/1\,500,6/2\,000}
			\draw[shift={(0,\y)},color=black] (5pt,0pt) -- (-5pt,0pt) node[left]{\Huge $\diem$};
		\end{tikzpicture}
	\end{center}
	\loigiai{
		Trên đoạn $[0; 2\,000]$, ta có $y'=-\dfrac{1}{440\,000}x^2+\dfrac{9}{1\,760}x^2-\dfrac{81}{44}$.\\
		Suy ra $y'=0 \Leftrightarrow \hoac{& x=450\\& x=1800.}$\\
		Bảng biến thiên
		\begin{center}
			\begin{tikzpicture}
				\tkzTabInit[nocadre=true,lgt=1.2,espcl=3,deltacl=0.6]
				{$x$ /0.6,$y'$ /0.6,$y$ /2}
				{$0$,$450$,$1800$,$2000$}
				\tkzTabLine{,-,$0$,+,$0$,-,}
				\tkzTabVar{+/$840$,-/$\dfrac{7365}{16}$,+/$\dfrac{15\,315}{11}$,-/$\dfrac{43720}{33}$}
			\end{tikzpicture}
		\end{center}
		Vậy tọa độ các đỉnh là $\left(450; \dfrac{7365}{16}\right)$, và $\left(1\,800, \dfrac{15\,315}{11}\right)$.
	}
\end{bt}

\begin{bt}%[2D1V2-7]
% nguon: Archive/2024/TaiLieuGiangDay/1.HamSo/1. K12 - 4in1 - Hamso/ThubaiGoc/Bai2-Dang8.tex #2716
[CKP]%[Bài giảng 12-4in1, Nhật Thiện] 
	Thể tích $V$ của $1$ kg nước (tính bằng $\text{cm}^3$) ở nhiệt độ $T$ (đơn vị: $\circ C$) khi $T$ thay đổi từ $0^\circ C$ đến $30^\circ C$ được cho xấp xỉ bởi công thức:
	$$V=999{,}87-0{,}06426T+0,0085043T^2-0{,}0000769T^3.$$
	\textit{(Nguồn: James Stewart, J. (2015). Calculus.Cengage Learning 8th edition, p.284)}\\ Tìm nhiệt độ $T_0 \in \left(0;30\right)$ để thể tích $V$ đạt cực trị (làm tròn kết quả đến hàng đơn vị).
	\loigiai{Xét hàm số $f\left(T\right)=V=999{,}87-0{,}06426T+0,0085043T^2-0{,}0000769T^3$.\\ Hàm số xác định trên khoảng $\left(0;30\right)$\\
		Ta có $f'\left(T\right)=2{,}307 \cdot 10^{-4}T^2+0{,}0170086T-0{,}06426$.\\
		$f'\left(T\right)=0\Leftrightarrow \hoac{&x=3{,}6\\&x=-77{,}33}\Leftrightarrow x=3{,}6$.
		\begin{center}
			\begin{tikzpicture}
				\tkzTabInit[lgt=1.2,espcl=3]
				{$T$ /.7, $f’(T)$ /.7, $f(T)$ /2.5}
				{$0$,$3{,}6$,$30$}
				\tkzTabLine{ ,-,z,+,}
				\tkzTabVar{+/$f\left(0\right)$,-/$f\left(3{,}6\right)$,+/$f\left(30\right)$}
			\end{tikzpicture}
		\end{center}
		Vậy $T_0=3{,}6$ thì thể tích $V$ đạt cực trị.
	}
\end{bt}

\begin{bt}%[2D1V2-7]
% nguon: Archive/2024/TaiLieuGiangDay/1.HamSo/1. K12 - 4in1 - Hamso/ThubaiGoc/Bai2-Dang8.tex #2718
[CD]%[Bài giảng 12-4in1, Nhật Thiện]
	\immini{
		Máng trượt của một cầu trượt cho trẻ em (\textit{Hình 5a}) được uốn từ một tấm kim loại có bề rộng $80$ cm, mặt cắt được mô tả ở \textit{Hình 5b}. Nhà thiết kế khuyến cáo, diện tích mặt cắt càng lớn thì càng đảm bảo an toàn cho trẻ em.
	}{
		\begin{tikzpicture}[scale=1, font=\footnotesize, line join=round, line cap=round, >=stealth,line width=3.pt]
			\clip(-0.2,-1.2) rectangle (7.5,2.6);
			\fill[line width=3.pt,color=red,fill=red,fill opacity=0.7] (0,0)--(0.41,0.34)--(1,1)--(1.76,1.76)--(2,2)--(2.46,2.08)--(2.83,1.95)--(2.34,1.65)--(1.34,0.65)--(0.75,-0.01)--(0.36,-0.13)-- cycle;
			\draw (0,0)--(0.41,0.34)--(1,1)--(2,2)--(2.46,2.08);
			\draw (0,0)--(0.36,-0.13)--(0.77,0.21)--(0.77,0.21)--(1.36,0.87)--(2.36,1.87)--(2.83,1.95);
			\draw (0.36,-0.13)--(0.75,-0.01)--(1.34,0.65)--(2.34,1.65);
			\draw (2.36,0.68)-- (3.95,0.28);
			\draw (3.42,0.41)--(2.68,2.29)--(2.53,2.37)--(2.16,1.88)--(2.31,1.53);
			\draw (1.63,2.01)--(2,2.5)-- (2.16,2.42)--(2.89,0.55);
			\draw (1.63,2.01)--(1.76,1.76) (2.73,0.98)--(3.26,0.85) (2.52,1.49)--(3.05,1.35);
			\draw (2.34,1.65)--(2.83,1.95)--(2.46,2.08);
			\draw (2,-0.6)node{a)};
			\fill[line width=3.pt,fill=gray,fill opacity=0.3] (5,0.5)--(5,0)--(6.5,0)--(6.5,0.5)--cycle;
			\draw (5,0.5)--(5,0) node [left ,sloped,pos=0.5,rotate=90] {$x$ cm};
			\draw (5,0)--(6.5,0) node [below,sloped,pos=0.5] {$y$ cm};
			\draw (6.5,0)--(6.5,0.5) node [right,sloped,pos=0.5,rotate=-90] {$x$ cm};
			\draw[dashed] (5,0.5)--(6.5,0.5);
			\draw (5.8,-0.6)node{b)};
			\draw (4,-1)node{\textit{Hình 5}};
		\end{tikzpicture}
	}
	\begin{enumerate}[a)]
		\item Gọi $S$ là diện tích mặt cắt. Tìm điều kiện của $x$ và viết công thức tính $S$ theo $x$.
		\item Với $x$ đạt giá trị bằng bao nhiêu thì mức độ đảm bảo an của cầu trượt đạt cực đại?
	\end{enumerate}
	\loigiai{
		\begin{enumerate}[a)]
			\item Do tấm kim loại có bề rộng $80 \mathrm{~cm}$ nên ta có: $2x+y=80 \Leftrightarrow y=80-2x$.\\
			Để có thể thiết kế được máng trượt thì $y>0 \Leftrightarrow 80-2 x>0 \Leftrightarrow x<40$. Suy ra $0<x<40$.\\
			Diện tích của mặt cắt máng trượt là: $S=x y=x(80-2x)=-2x^2+80x$.		
			\item Ta có: 
			\allowdisplaybreaks
			\begin{eqnarray*}
				S(x)&=&-2x^2+80x \text{ với } x\in(0; 40);\\
				S'(x)&=&-4 x+80 \\
				S'(x)&=&0 \Leftrightarrow-4 x+80=0 \Leftrightarrow x=20.
			\end{eqnarray*}
			Bảng biến thiên của hàm số $S(x)$ như sau:
			\begin{center}
				\begin{tikzpicture}[font=\footnotesize,thick,>=stealth]
					\tikzset{double style/.append style={double distance=1.5pt}}
					\tkzTabInit[nocadre=false,lgt=1.2,espcl=2.5,deltacl=0.6,lw=.75pt,color,colorL=green!50,colorV=green!50]
					{$x$ /0.7, $S'(x)$ /0.8, $S(x)$ /2}
					{$0$,$20$,$40$}
					\tkzTabLine{ ,+,,-, }
					\tkzTabVar{-/$0$,+/$800$,-/$0$}
				\end{tikzpicture}
			\end{center}
			Do đó, hàm số $S(x)$ đạt cực đại tại $x=20$ và $S_{\text{CĐ}}=800$.\\
			Vậy mức độ đảm bảo an toàn của cầu trượt đạt cực đại thì $x=20$ (cm).
		\end{enumerate}
	}
\end{bt}

\begin{bt}%[2D1V2-7]
% nguon: Archive/2024/TaiLieuGiangDay/1.HamSo/1. K12 - 4in1 - Hamso/ThubaiGoc/Bai2-Dang8.tex #2719
%[Bài giảng 12-4in1, Nhật Thiện]
	Trong các khối trụ có cùng diện tích toàn phần là $6\pi $. 
	\begin{enumerate}
		\item Mô tả thể tích $V$ khối trụ theo bán kính $R$.
		\item Tìm bán kính đáy của khối trụ có thể tích cực đại.
	\end{enumerate}
	
	\loigiai
	{
		\begin{enumerate}
			\item Ta có $S_{tp}=2\pi Rh+2\pi R^2=6\pi \Leftrightarrow Rh+R^2=3\Leftrightarrow h=\dfrac{3-R^2}{R}$. Điều kiện: $0<R<\sqrt{3}$. \\
			$V=\pi R^2h=\pi R^2\cdot \dfrac{3-R^2}{R}=\pi R(3-R^2)=-\pi R^3+3\pi R$.
			\item $V'=-3\pi R^2+3\pi =0\Leftrightarrow \hoac{&R=1 \text{ (nhận)} \\ &R=-1 \text{ (loại)}.}$ \\
			Bảng biến thiên
			\begin{center}
				\begin{tikzpicture}
					\tkzTabInit[lgt=1.5,espcl=4]
					{$R$/1,$V’(R)$/1,$V(R)$/2.5}
					{$0$,$1$,$\sqrt{3}$}
					\tkzTabLine{ ,+,0,-, }
					\tkzTabVar{-/$$,+/$2\pi$,-/$$}
				\end{tikzpicture}
			\end{center}
			Nhìn bảng biến thiên ta thấy $V_{\text{CĐ}}=V(1)=2\pi $ khi $R=1$.
		\end{enumerate} 
		
`	}
\end{bt}

\begin{bt}%[2D1V2-7]
% nguon: Archive/2024/TaiLieuGiangDay/1.HamSo/1. K12 - 4in1 - Hamso/ThubaiGoc/Bai2-Dang8.tex #2720
%[Bài giảng 12-4in1, Nhật Thiện]
	Từ một miếng tôn hình tròn bán kính $2$ m, người ta cắt ra một hình chữ nhật rồi uốn thành mặt xung quanh của một chiếc thùng phi hình trụ như hình vẽ.
	\begin{center}
		\begin{tikzpicture}[scale=1,font=\footnotesize,line join=round, line cap=round,>=stealth]
			\def \r{1.5}
			\path (0,0) coordinate (O);
			\draw (O) circle (\r cm);
		\end{tikzpicture}\qquad
		\begin{tikzpicture}[scale=1,font=\footnotesize,line join=round, line cap=round,>=stealth]
			\def \r{1.5}
			\path
			(0,0) coordinate (O)
			(30:\r) coordinate (A)
			(150:\r) coordinate (B)
			(-150:\r) coordinate (C)
			(-30:\r) coordinate (D);
			\draw (O) circle (\r cm); \draw (A)--(B)--(C)--(D)--cycle;
		\end{tikzpicture}\qquad
		\begin{tikzpicture}[scale=1,font=\footnotesize,line join=round, line cap=round,>=stealth]
			\def \x{1} \def \y{0.3} \def \h{1.5}
			\path
			(0,0) coordinate (A)
			(2*\x,0) coordinate (B)
			($(A)!0.5!(B)$) coordinate (O)
			($(O)+(0,\h)$) coordinate (O')
			($(A)+(0,\h)$) coordinate (A')
			($(B)+(0,\h)$) coordinate (B')
			;
			\draw[dashed] (B) arc(0:180:\x cm and \y cm);
			\draw (B) arc(0:-180:\x cm and \y cm);
			\draw (O') ellipse (\x cm and \y cm);
			\draw (A)--(A') (B)--(B');
		\end{tikzpicture}
	\end{center}
	Tính diện tích phần tôn bị cắt bỏ biết thể tích thùng đạt cực đại (làm tròn đến hàng đơn vị).
	\loigiai{
		\immini{
			Gọi $h$, $x$ lần lượt là chiều rộng và chiều dài của hình chữ nhật. Gọi $r$ là bán kính đáy của hình trụ thu được.\\
			Khi đó $V_{\text{trụ}}=\pi r^2 h$. \quad $(1)$\\
			Vì chu vi đáy của hình trụ là chiều dài của hình chữ nhật nên
			$$C=2\pi r=x \Rightarrow r=\dfrac{x}{2\pi}.\quad (2)$$
		}{
			\begin{tikzpicture}[scale=1,font=\footnotesize,line join=round, line cap=round,>=stealth]
				\def \r{1.5}
				\path
				(0,0) coordinate (O)
				(30:\r) coordinate (A)
				(150:\r) coordinate (B)
				(-150:\r) coordinate (C)
				(-30:\r) coordinate (D);
				\draw (O) circle (\r cm); \draw (A)--node[above]{$x$}(B)--(C)--(D)--node[left]{$h$}(A);
			\end{tikzpicture}\qquad
			\begin{tikzpicture}[scale=1,font=\footnotesize,line join=round, line cap=round,>=stealth]
				\def \x{1} \def \y{0.3} \def \h{1.5}
				\path
				(0,0) coordinate (A)
				(2*\x,0) coordinate (B)
				($(A)!0.5!(B)$) coordinate (O)
				($(O)+(0,\h)$) coordinate (O')
				($(A)+(0,\h)$) coordinate (A')
				($(B)+(0,\h)$) coordinate (B')
				;
				\draw[dashed] (B) arc(0:180:\x cm and \y cm);
				\draw (B) arc(0:-180:\x cm and \y cm);
				\draw (O') ellipse (\x cm and \y cm);
				\draw (A)--(A') (B)--node[right]{$h$}(B'); \draw[dashed] (O)--node[below]{$r$}(B);
		\end{tikzpicture}}
		\noindent Vì miếng tôn hình tròn bán kính $2$ m nên
		$\left(\dfrac{x}{2}\right)^2+\left(\dfrac{h}{2}\right)^2=2^2 \Rightarrow x^2=16-h^2. \quad (3)$\\
		Từ $(1)$, $(2)$, $(3)$ suy ra
		$V_{\text{trụ}}=\pi \cdot \dfrac{16-h^2}{(2\pi)^2}\cdot h =\dfrac{1}{4\pi }\cdot h(16-h^2)=\dfrac{1}{4\pi }\cdot (16h-h^3).$\\
		Xét hàm số $f(h)=16h-h^3$, có $f'(h)=16-3h^2=0 \Rightarrow h=\dfrac{4\sqrt{3}}{3}~\overset{\text{đặt}}{=}~h_0$ và bảng biến thiên
		\begin{center}
			\begin{tikzpicture}
				\tkzTabInit[nocadre=false,lgt=1.2,espcl=2.5,deltacl=0.6]
				{$h$ /0.7,$f'(h)$ /.7,$f(h)$ /2}{$0$,$h_0$,$+\infty$}
				\tkzTabLine{,+,$0$,-,}
				\tkzTabVar{-/$0$,+/$f(h_0)$,-/$-\infty$}
			\end{tikzpicture}
		\end{center}
		Suy ra thể tích thùng đạt cực đại khi $h=\dfrac{4\sqrt{3}}{3}\Rightarrow x=\dfrac{4\sqrt{6}}{3}$.\\
		Vậy diện tích phần tôn bị cắt bỏ là
		$$S=S_{\text{hình tròn}}-S_{\text{hình chữ nhật}}=\pi \cdot 2^2 -x\cdot h=4\pi -\dfrac{4\sqrt{6}}{3} \cdot \dfrac{4\sqrt{3}}{3}\approx 5.$$
	}
\end{bt}

\begin{bt}%[2D1V2-7]
% nguon: Archive/2024/TaiLieuGiangDay/1.HamSo/1. K12 - 4in1 - Hamso/ThubaiGoc/Bai2-Dang8.tex #2721
%[Bài giảng 12-4in1, Nhật Thiện]
	Sau khi tiêu thụ đồ uống có cồn, nồng độ cồn trong máu (\textit{blood alcohol concentration} hoặc BAC) tăng lên khi rượu được hấp thụ, tiếp theo là giảm dần khi rượu được chuyển hóa. Hàm số $$C(t)=1{,}35\cdot t\cdot \mathrm{e}^{-2{,}802t}$$ minh họa BAC trung bình, tính theo đơn vị mg/mL, của một nhóm $8$ người đàn ông sau khi tiêu thụ nhanh $15$ mL ethanol (tương ứng với một thức uống có cồn) trong $t$ giờ. Tính giá trị cực đại của BAC trong $3$ giờ đầu tiên.
	\loigiai{
		Ta có $C'(t)=\mathrm{e}^{-2{,}802 t} (1{,}35 - 3{,}7827 t)=0\Leftrightarrow t=\dfrac{500}{1401}=t_0$.\\
		Bảng biến thiên của hàm số $C(t)$ trên $[0;3]$
		\begin{center}
			\begin{tikzpicture}
				\tkzTabInit[nocadre=false,lgt=1.2,espcl=2.5,deltacl=0.6]
				{$t$ /0.6,$C'(t)$ /0.6,$C(t)$ /2}
				{$0$,$t_0$,$3$}
				\tkzTabLine{,+,0,-,}
				\tkzTabVar{-/,+/$C(t_0)$,-/}
			\end{tikzpicture}
		\end{center}
		Dựa vào bảng biến thiên, hàm số đã cho đạt cực đại bằng $C(t_0)\approx 0{,}18$ mg/mL.
	}
\end{bt}

\begin{bt}%[2D1H2-7]
% nguon: Archive/2024/TaiLieuGiangDay/1.HamSo/1. K12 - 4in1 - Hamso/ThubaiGoc/Bai2-Dang8.tex #2725
%[Bài giảng 12-4in1, Nhật Thiện]
	\immini{Một phần đường ray tàu lượn siêu tốc có dạng đồ thị hàm số bậc ba $y=f(x)=-\dfrac{1}{2\,000}x^3+\dfrac{17}{200}x^2-4x+50$. Trục $Ox$ mô tả quãng đường tàu di chuyển theo chiều ngang (tính bằng centimét), trục $Oy$ mô tả chiều cao của đường ray (tính bằng centimét) tại mỗi vị trí $x$.Tìm tọa độ các cực trị của đường ray trên đoạn $[0;100]$.
}{
	\begin{tikzpicture}[scale=1, font=\footnotesize, line join=round, line cap=round, >=stealth]
		\begin{scope}[scale=.6]
			\draw[->] (-1.1,0)--(12.1,0) node[below left] {$x$};
			\draw[->] (0,-1.6)--(0,6.1) node[below left] {$y$};
			\draw (0,0) node [below left] {$O$};
			\foreach \x/\nx in {1/10,2/20,3/30,4/40,5/50,6/60,7/70,8/80,9/90,10/100,11/110}
			\draw[thin] (\x,1pt)--(\x,-1pt) node [below] {$\nx$};
			\foreach \y/\ny in {1/10,2/20,3/30,4/40,5/50}
			\draw[thin] (1pt,\y)--(-1pt,\y) node [left] {$\ny$};
			\begin{scope}
				\clip (-1,-1.5) rectangle (12,6);
				\draw[samples=200,domain=-1:12,smooth,variable=\x] plot (\x,{(-1/20)*((\x)-2)*((\x)-5)*((\x)-10)});
			\end{scope}
		\end{scope}
	\end{tikzpicture}}
\loigiai{
	Ta có $y'=-\dfrac{3}{2\,000}x^2+\dfrac{17}{100}x-4$.\\
	Xét $y'=0\Leftrightarrow \hoac{&x=\dfrac{100}{3}\\&x=80.}$\\
	Bảng biến thiên
	\begin{center}
		\begin{tikzpicture}
\tkzTabInit[nocadre=false,lgt=1.2,espcl=2.5,deltacl=0.6]
{$x$ /0.6,$y’$ /0.6,$y$ /2}
{$0$,$\tfrac{100}{3}$,$80$,$100$}
\tkzTabLine{,-,0,+,0,-}
\tkzTabVar{+/$50$,-/$-\tfrac{200}{27}$,+/$18$,-/$0$}
\end{tikzpicture}
	\end{center}
Ta có tọa độ các điểm cực trị của đồ thị hàm số $f(x)$ là $A\left(\dfrac{100}{3};-\dfrac{200}{27}\right)$ và $B(80;18)$.
}
\end{bt}

\begin{bt}%[2D1H2-7]
% nguon: Archive/2024/TaiLieuGiangDay/1.HamSo/1. K12 - 4in1 - Hamso/ThubaiGoc/Bai2-Dang8.tex #2726
%[Bài giảng 12-4in1, Nhật Thiện]
	Khi một vật lạ mắc vào khí quản buộc một người phải ho, cơ hoành sẽ đẩy lên trên, làm tăng áp lực trong phổi. Đồng thời, khí quản co lại, khiến vật thể bị tống ra ngoài, không khí di chuyển nhanh hơn và tăng áp lực lên vật thể. Thông qua mô hình toán học, vật tốc (tính bằng cm/s) của luồng khí qua khí quản của một người trưởng thành liên quan đến bán kính $r$ của khí quản (tính bằng cm) theo hàm số 
	$$v(r)=3{,}2(1-r)r^2\, \text{ với }\dfrac{1}{2}\le r\le 1.$$
	\textit{(Nguồn: James Stewart, Precalculus Maththematics for Calculus - 7$^\text{th}$, trang 209)}\\
	Tìm giá trị của $r$ để $v$ đạt cực đại.
	\loigiai{
	Ta có $v'(r)=(6{.}4 - 9{.}6 r) r=0\Leftrightarrow \hoac{&r=0\\&r=\dfrac{2}{3}.}$\\
	Bảng biến thiên của $v(r)$
	\begin{center}
		\begin{tikzpicture}
			\tkzTabInit[nocadre=false,lgt=1.2,espcl=2.5,deltacl=0.6]
			{$r$ /0.6,$v’$ /0.6,$v$ /2}
			{$\frac{1}{2}$,$\frac{2}{3}$,$1$}
			\tkzTabLine{,+,0,-,}
			\tkzTabVar{-/,+/$\dfrac{64}{135}$,-/}
		\end{tikzpicture}
	\end{center}
	Vậy với $r=\dfrac{2}{3}$ thì $v$ đạt cực đại.
}
\end{bt}

\begin{bt}%[2D1H2-7]
% nguon: Archive/2024/TaiLieuGiangDay/1.HamSo/1. K12 - 4in1 - Hamso/ThubaiGoc/Bai2-Dang8.tex #2727
%[Bài giảng 12-4in1, Nhật Thiện]
	\immini{
		Sự phân huỷ của rác thải hữu cơ có trong nước sẽ làm tiêu hao oxygen hoà tan trong nước. Nồng độ oxygen (mg/l) trong một hồ nước sau $t$ giờ $(t \geq 0)$ khi một lượng rác thải hữu cơ bị xả vào hồ được xấp xỉ bởi hàm số (có đồ thị như đường màu đỏ ở hình bên)
		$$
		y(t)=5-\frac{15 t}{9 t^2+1}.
		$$
	}{
		\begin{tikzpicture}[>=stealth,x=1cm,y=0.3cm,scale=2,font=\footnotesize]
			\draw[->] (-0.5,0) -- (4,0) node[below] {$t$};
			\draw[->] (0,-1) -- (0,6) node[left] {$y$};
			\filldraw (0,0) circle (1pt)node[below left]{$O$};
			\draw[domain=0:4,samples=200,red] plot (\x,{5-(15*(\x))/(9*(\x)^2+1)});
			\draw[dashed] (0,5) node [left] {$5$}--(4,5);
			\foreach \x/\g in {1/-90,2/-90,3/-90}
			\draw[thin] (\x,2pt)--(\x,-2pt) + (\g:3mm) node {$\x$};
		\end{tikzpicture}
	}
	\noindent
	Tìm các thời điểm nồng độ oxygen trong nước đạt cực trị?
	\loigiai{
	Ta có $y'(t)=\dfrac{135t^2-15}{(9t^2+1)^2}=0\Leftrightarrow t=\dfrac{1}{3}$.\\
	Bảng biến thiên của hàm $y(t)$
	\begin{center}
		\begin{tikzpicture}[scale=1, font=\footnotesize, line join=round, line cap=round, >=stealth]
\tkzTabInit[nocadre=false,lgt=1.2,espcl=2.5,deltacl=0.6]
{$t$ /0.6,$y’$ /0.6,$y$ /2}
{$0$,$\frac{1}{3}$,$+\infty$}
\tkzTabLine{,-,0,+,}
\tkzTabVar{+/$5$,-/$\dfrac{5}{2}$,+/$5$}
\end{tikzpicture}
	\end{center}
Vậy nồng độ oxygen trong nước đạt cực tiểu tại thời điểm $t=\dfrac{1}{3}$ giờ.
}
\end{bt}

\begin{bt}%[2D1V2-7]
% nguon: Archive/2024/TaiLieuGiangDay/1.HamSo/1. K12 - 4in1 - Hamso/ThubaiGoc/Bai2-Dang8.tex #2728
%[Bài giảng 12-4in1, Nhật Thiện]
	Khối lượng $q$ (kg) của một mặt hàng mà cửa tiệm bán được trong một ngày phụ thuộc vào giá bán $p$ (nghìn đồng/kg) theo công thức $p=15-\dfrac{1}{2} q$. Doanh thu từ việc bán mặt hàng trên của cửa tiệm được tính theo công thức $R=p q$. Tìm giá bán mỗi kilôgam sản phẩm để doanh thu đạt cực trị và xác định doanh thu đó.
	\loigiai{
		 Ta có $p=15-\dfrac{1}{2} q$ $\Rightarrow q=30-2p$.\\
			Lại có $R=pq=p\left(30-2p\right)\Rightarrow R=-2p^2+30p$ với $0<p$.\\
		Xét $R'=-4p+30=0\Leftrightarrow p=7{,}5$.
			Ta có bảng biến thiên như sau
			\begin{center}
				\begin{tikzpicture}[font=\normalsize,t style/.style={style=solid}]
					%dòng khai báo
					\tkzTabInit[lgt=1,espcl=2.5,deltacl=0.6]
					{$x$ /0.75, $R'$/0.75, $R$/2}
					{$ 0$,$ 7{,}5 $,$ +\infty$}
					%dòng xét dấu
					\tkzTabLine{  , +,0 , -,  }  % z, t, d;
					%dòng biến thiên
					\tkzTabVar{-/$0$,+/$112{,}5$,-/$-\infty$} %+ hoac-
				\end{tikzpicture}
			\end{center}
			Từ bảng biến thiên thấy để đạt được doanh thu đạt cực trị thì giá bán mỗi kilôgam sản phẩm $p=7{,}5$ (nghìn đồng/kg) khi đó doanh thu sẽ là $R(7{,}5)=112{,}5$ (nghìn đồng).
		}
\end{bt}

\begin{bt}%[2D1H2-7]
% nguon: Archive/2024/TaiLieuGiangDay/1.HamSo/1. K12 - 4in1 - Hamso/ThubaiGoc/Bai2-Dang8.tex #2734
%[Bài giảng 12-4in1, Nhật Thiện]
	Khi máu di chuyển từ tim qua các động mạch chính rồi đến các mao mạch và quay trở lại qua các tĩnh mạch, huyết áp tâm thu (tức là áp lực của máu lên động mạch khi tim co bóp) liên tục giảm xuống. Giả sử một người có huyết áp tâm thu $P$ (tính bằng mmHg) được cho bởi hàm số
	$$
	P(t)=\dfrac{25 t^2+125}{t^2+1}, 0 \leq t \leq 10,
	$$
	trong đó thời gian $t$ được tính bằng giây. Tìm giá trị cực trị của huyết áp tâm thu $P$.
	\loigiai{
		Ta có tốc độ thay đổi của huyết áp là $P'(t)=\dfrac{-100t}{(t^2+1)^2}$.\\
		Do đó tốc độ thay đổi huyết áp sau $5$ s là $P'(5)=-\dfrac{125}{169}$.
	}
\end{bt}

\begin{bt}%[2D1N1-2]
% nguon: Archive/2024/TaiLieuGiangDay/1.HamSo/1. K12 - 4in1 - Hamso/ThubaiGoc/Bai1-Dang2.tex #2739
[Mức 1]%[BG12new-4in1, Trần Hoà]
    Cho hàm số $y=f(x)$ có bảng biến thiên như hình dưới đây
    \begin{center}
        \begin{tikzpicture}
            \tkzTabInit[nocadre=false,lgt=1.2,espcl=3]
            {$x$ /0.7,$y'$ /0.7,$y$ /2}
            {$-\infty$,$-\tfrac{1}{2}$,$3$,$+\infty$}
            \tkzTabLine{,+,d,+,$0$,-,}
            \tkzTabVar{-/$-\infty$,+D-/$+\infty$/$-\infty$,+/$4$,-/$-\infty$}
        
\end{tikzpicture}
    \end{center}
   Tìm các khoảng đồng biến, nghịch biến của hàm số?
    \loigiai
    {Dựa vào bảng biến thiên, ta thấy hàm số đồng biến trên các khoảng $\left(-\infty;-\dfrac{1}{2}\right)$ và $\left(-\dfrac{1}{2};3\right)$; hàm số nghịch biến trên khoảng $(3;+\infty)$.
    }
\end{bt}

\begin{bt}%[2D1N1-2]
% nguon: Archive/2024/TaiLieuGiangDay/1.HamSo/1. K12 - 4in1 - Hamso/ThubaiGoc/Bai1-Dang2.tex #2740
[Mức 1]%[BG12new-4in1, Trần Hoà]
Cho hàm số $y=f(x)$ có bảng biến thiên như sau.
\begin{center}
\begin{tikzpicture}[>=stealth]
		\tkzTabInit[nocadre=false,lgt=1,espcl=2,deltacl=0.5]{$x$/.7 ,$y'$/.7,$y$/2}
		{$-\infty$ , $-3$ , $4$ , $+\infty$}
		\tkzTabLine{ , - , $0$ , + , $0$ , - , }
		\tkzTabVar{+/$+\infty$ , -/$-1$ , +/$2$ , -/$-\infty$}
\end{tikzpicture}
\end{center}
Tìm các khoảng đồng biến, nghịch biến của hàm số.
	\loigiai{Dựa vào bảng biến thiên ta có:\\
	Hàm số đồng biến trên khoảng $(-3;4)$.\\
	Hàm số nghịch biến trên các khoảng $(-\infty;-3)$ và $(4;+\infty)$.}	
\end{bt}

\begin{bt}%[2D1H1-2]
% nguon: Archive/2024/TaiLieuGiangDay/1.HamSo/1. K12 - 4in1 - Hamso/ThubaiGoc/Bai1-Dang2.tex #2741
[Mức 2]%[BG12new-4in1, Trần Hoà]
Cho hàm số $y=f(x)$ có bảng biến thiên như hình dưới đây
\begin{center}
\begin{tikzpicture}[scale=1, font=\footnotesize, line join=round, line cap=round, >=stealth]
\tkzTabInit[nocadre=false,lgt=1.2,espcl=2.5,deltacl=0.6] %phần bắt buộc
{$x$ /0.6,$y’$ /0.6,$y$ /2}
{$-\infty$,$0$,$2$,$4$,$+\infty$}
\tkzTabLine{ ,+,$0$,-,d,-,$0$,+, }
\tkzTabVar{-/$-\infty$,+/$-2$,-D+/$-\infty$/$+\infty$,-/$6$,+/$+\infty$}
\end{tikzpicture}
\end{center}
Tìm các khoảng đồng biến, nghịch biến của hàm số.
\loigiai{Dựa vào bảng biến thiên ta có:\\
	Hàm số đồng biến trên các khoảng $(-\infty;0)$ và $(4;+\infty)$.\\
	Hàm số nghịch biến trên các khoảng $(0;2)$ và $(2;+4)$.}	
\end{bt}

\begin{bt}%[2D1H1-2]
% nguon: Archive/2024/TaiLieuGiangDay/1.HamSo/1. K12 - 4in1 - Hamso/ThubaiGoc/Bai1-Dang2.tex #2743
[Mức 2]%[BG12new-4in1, Trần Hoà]
	\immini
	{
		Đường cong trong hình bên là đồ thị của hàm số $y=\dfrac{ax+b}{cx+d}$, với $a$, $b$, $c$, $d$ là các số thực. Tìm các khoảng đồng biến, nghịch biến của hàm số.	
	}
	{
		\begin{tikzpicture}[scale=0.8,>=stealth, font=\footnotesize, line join=round, line cap=round]
			\def\a{2} \def\b{1} \def\c{2} \def\d{-1} % Hệ số
			\def\xmin{-3} \def\xmax{4}
			\def\ymin{-2.3} \def\ymax{4}
			\draw[->] (\xmin,0)--(\xmax,0) node [below]{$x$};
			\draw[->] (0,\ymin)--(0,\ymax) node [left]{$y$};
			\node at (0,0) [above left]{$O$};
			\clip (\xmin+0.1,\ymin+0.1) rectangle (\xmax-0.1,\ymax-0.1);
			\draw[smooth,samples=300,domain=\xmin:(-\d/\c-0.1)] plot(\x,{(\a*(\x)+\b)/(\c*(\x)+\d)});
			\draw[smooth,samples=300,domain=(-\d/\c+0.1:\xmax)] plot(\x,{(\a*(\x)+\b)/(\c*(\x)+\d)});
			\draw[dashed] (-\d/\c,\ymin)--(-\d/\c,\ymax);
			\draw[dashed] (\xmin,\a/\c)--(\xmax,\a/\c);
			\node at (0.7,-0.2){2};
			\node at (-0.2,1.2){3};
		\end{tikzpicture}
	}
	\loigiai{
		Nhận thấy hàm số nghịch biến trên từng khoảng $(-\infty; 2)$ và $(2;+\infty)$.
	}
	
\end{bt}

\begin{bt}%[2D1V1-2]
% nguon: Archive/2024/TaiLieuGiangDay/1.HamSo/1. K12 - 4in1 - Hamso/ThubaiGoc/Bai1-Dang2.tex #2744
[Kết nối, mức 2]%[TeX hóa SGK CTST 12]%[Nguyễn Tiến]%[BG12new-4in1, Trần Hoà]%[BG12new-4in1, Trần Hoà]
	\immini{
		Đạo hàm $f'(x)$ của hàm số $y=f(x)$ trên khoảng $(-2;5)$ có đồ thị như hình vẽ bên. Xét tính đơn điệu của hàm số $y=f(x)$ trên khoảng $(-2;5)$.
	}{
		\begin{tikzpicture}[scale=0.7, font=\footnotesize, line join=round, line cap=round, >=stealth]
			\def\a{0.125} \def\b{-0.625} \def\c{0.25} \def\d{1} % Hệ số
			\def\xt{-3} \def\xp{6} \def\yt{2.5} \def\yd{-3} % x_trái, x_phải, y_trên, y_dưới (giới hạn)
			\draw[->] (\xt,0)--(\xp,0) node [below]{$x$};
			\draw[->] (0,\yd)--(0,\yt) node [left]{$y$};
			\node at (0,0) [below right]{$O$};
			\clip (\xt+0.1,\yd-1) rectangle (\xp-0.1,\yt-0.1);
			\draw[smooth,samples=300,domain=-2:5] plot(\x,{\a*(\x)^3+\b*(\x)^2+\c*(\x)+\d}) node [below left]{$y=f'(x)$};
			\draw[dashed] (-2,0) node [below left]{$-2$} -- (-2,-3);
			\draw[dashed] (5,0) node [below]{$5$} -- (5,2.25);
			\fill (-1,0) circle (1.5pt) node[below right]{$-1$};
			\fill (2,0) circle (1.5pt) node[below left]{$2$};
			\fill (4,0) circle (1.5pt) node[below right]{$4$};
			%\node at (1.5,-3.5){\textit{Hình 12}};
		\end{tikzpicture}
	}
	\loigiai{
		Từ đồ thị ta thấy $f'(x)=0 \Leftrightarrow \hoac{& x=-1\\& x=2\\& x=4.}$\\
		Ta có bảng biến thiên như sau
		\begin{center}
			\begin{tikzpicture}
		\tkzTabInit[nocadre=false,lgt=1.2,espcl=2.5,deltacl=0.6]
				{$x$ /0.6,$f'(x)$ /0.6,$f(x)$ /2}
				{$-2$,$-1$,$2$,$4$,$5$}
				\tkzTabLine{,-,$0$,+,$0$,-,$0$,+,}
				\tkzTabVar{+/,-/,+/,-/,+/}
			\end{tikzpicture}
		\end{center}
		Vậy hàm số đồng biến trên $(-1;2)$ và $(4;5)$, nghịch biến trên $(-2;-1)$ và $(2;4)$.
	}
\end{bt}

\begin{bt}%[2D1N1-2]
% nguon: Archive/2024/TaiLieuGiangDay/1.HamSo/1. K12 - 4in1 - Hamso/ThubaiGoc/Bai1-Dang2.tex #2748
%[BG12new-4in1, Trần Hoà]
	Cho hàm số $y=f(x)$ có bảng biến thiên như hình vẽ.
	\begin{center}
		\begin{tikzpicture}[scale=1, font=\footnotesize, line join=round, line cap=round, >=stealth]
			\tkzTabInit[nocadre=false,lgt=1,espcl=3,deltacl=0.5]
			{$x$/.7 ,$y'$/.7,$y$/2}
			{$-\infty$ , $1$ , $2$ , $+\infty$}
			\tkzTabLine{ , + , $0$ , - , d , + , }
			\tkzTabVar{-/$-\infty$ , +/$3$ , -/$0$ , +/$+\infty$}
		\end{tikzpicture}
	\end{center}
	Tìm các khoảng đồng biến, nghịch biến của hàm số?
\loigiai{
		Từ bảng biến thiên ta thấy\\
		Hàm số đồng biến trên các khoảng $(-\inf; 1)$ và $(2;+\infty$.\\
		Hàm số nghịch biến trên khoảng $(1;2)$.
	}
\end{bt}

\begin{bt}%[2D1N1-2]
% nguon: Archive/2024/TaiLieuGiangDay/1.HamSo/1. K12 - 4in1 - Hamso/ThubaiGoc/Bai1-Dang2.tex #2749
%[SGK-CD]%[BG12new-4in1, Trần Hoà]
	Cho hai hàm số $y=f(x)$ xác định trên $\mathbb{R}$ có đồ thị như hình dưới đây. Tìm các khoảng đồng biến, nghịch biến hàm số.
	\begin{center}
		\begin{tikzpicture}[scale=1, font=\footnotesize, line join=round, line cap=round, >=stealth]
		\draw[gray!50,thin,opacity=.5] (-3,-2) grid (4,4);
		\draw[->] (-2.9,0)--(4.1,0) node[below] {$x$};
		\draw[->] (0,-2.1)--(0,4.1) node[left] {$y$};
		\draw (0,0) node [below left] {$O$};
		\foreach \x/\nx in {-1/-1,1/1,2/2}
		\draw[thin] (\x,1pt)--(\x,-1pt) node [below] {$\nx$};
		\foreach \y/\ny in {1/1,2/2}
		\draw[thin] (1pt,\y)--(-1pt,\y) node [left] {$\ny$};
		\begin{scope}
			\clip (-2.8,-2) rectangle (4,4);
			\draw[samples=200,domain=-2.5:-0.32,smooth,variable=\x] plot (\x,{-1.4*(\x)^2-2.92*(\x)-0.51});
			\draw[samples=200,domain=-0.32:0.68,smooth,variable=\x] plot (\x,{2.9*(\x)^2+0.05*(\x)});
			\draw[samples=200,domain=0.68:1.44,smooth,variable=\x] plot (\x,{-5.06*(\x)^2+10.44*(\x)-3.38});
			\draw[samples=200,domain=1.44:3.5,smooth,variable=\x] plot (\x,{4.34*(\x)^2-16.99*(\x)+16.64});
		\end{scope}
		\draw (3.4,2) node{$y=f(x)$};
		\draw (0,-4) node{\text{\textit{Hình }} $6a$};
	\end{tikzpicture}
	\end{center}
	\loigiai{
	Quan sát đồ thị, ta thấy hàm số $y=f(x)$ đồng biến trên các khoảng $(-\infty;-1)$, $(0;1)$, $(2;+\infty)$; nghịch biến trên các khoảng $(-1;0)$, $(1;2)$. 
	}
\end{bt}

\begin{bt}%[2D1N1-2]
% nguon: Archive/2024/TaiLieuGiangDay/1.HamSo/1. K12 - 4in1 - Hamso/ThubaiGoc/Bai1-Dang2.tex #2750
%[BG12new-4in1, Trần Hoà]
	Cho hàm số $y=f(x)$ có bảng biến thiên như hình vẽ dưới đây. Tìm khoảng đồng biến, nghịch biến của hàm số.
	\begin{center}
		\begin{tikzpicture}[line join = round, line cap = round,>=stealth,font=\footnotesize,scale=1]
			\tikzset{double style/.append style={double distance=1.5pt}}
			\tkzTabInit[nocadre=false,lgt=1.2,espcl=2.5,deltacl=0.6]
			{$x$ /0.6,$y'$ /0.6,$y$ /2}
			{$-\infty$,$0$,$1$,$+\infty$}
			\tkzTabLine{,-,d,+,$0$,-,}
			\tkzTabVar{+/$+\infty$,-D-/$-1$/$-\infty$,+/$2$,-/$-\infty$}
		\end{tikzpicture}
	\end{center}
	\loigiai{
		Nhìn vào bảng biến thiên của hàm số ta thấy\\
	Hàm số nghịch biến trên các khoảng $(-\infty;0)$ và $(1;+\infty)$.\\
	Hàm số đồng biến trên khoảng 	$(0;1)$.
	}
\end{bt}

\begin{bt}%[2D1H1-2]
% nguon: Archive/2024/TaiLieuGiangDay/1.HamSo/1. K12 - 4in1 - Hamso/ThubaiGoc/Bai1-Dang2.tex #2751
%[BG12new-4in1, Trần Hoà]
\immini{Cho hàm số $y=f(x)$ xác định trên $\mathbb{R}$ có đồ thị của hàm số $f'(x)$ như hình vẽ. Tìm khoảng đồng biến, nghịch biến của hàm số $y=f(x)$?
}{
\begin{tikzpicture}
[scale=1, font=\footnotesize, line join=round, line cap=round, >=stealth]
\draw[->] (-1,0)--(0,0) node[below left]{$O$}--(3,0) node[below]{$x$};
\draw[->] (0,-2) --(0,2) node[right]{$y$};
\foreach \x in {1,2}{
	\draw[fill] (\x,0) node[below]{$\x$} circle (1pt);
}

\draw [domain=.1:2.5, samples=100] %
plot (\x, {((\x)-2)*((\x)-1)^2});
%\draw [dashed] (-1,0)--(-1,4) --(0,4);%(2,0)--(2,-2)--(0,-2)--(-1,-2)--(-1,0);
\draw[fill] (0,0) circle (1pt);
\end{tikzpicture}
}
\loigiai{
Từ đồ thị hàm số $f'(x)$ ta thấy $f'(x)>0\Leftrightarrow x>2$.\\
Vậy hàm số $y=f(x)$ đồng biến trên $(2;+\infty)$.
}
\end{bt}

\begin{bt}%[2D1V1-2]
% nguon: Archive/2024/TaiLieuGiangDay/1.HamSo/1. K12 - 4in1 - Hamso/ThubaiGoc/Bai1-Dang2.tex #2752
%[BG12new-4in1, Trần Hoà]
Cho hàm số $y=f(x)$ có bảng biến thiên như hình vẽ.
\begin{center}
    \begin{tikzpicture}[scale=0.8, line join=round, line cap=round, >=stealth]
	\tkzTabInit[nocadre=false,lgt=1.3,espcl=2.5,deltacl=0.6]
	{$x$/0.6,$f'(x)$/0.6,$f(x)$/2.5}
	{$-\infty$,$-2$,$0$,$2$,$+\infty$}
	\tkzTabLine{,+,0,-,0,+,0,-,}
	\tkzTabVar{-/$-\infty$,+/$3$,-/$-1$,+/$3$,-/$-\infty$}
\end{tikzpicture}
\end{center}
Tìm khoảng đồng biến của hàm số $y=f(x^2-2)$?

\loigiai{
Ta có $y'=f'(x^2-2)\cdot 2x$;
\[y'=0\Leftrightarrow\hoac{&x=0\\&f'(x^2-2)=0}\Leftrightarrow \hoac{&x=0\\&x^2-2=-2\\&x^2-2=0\\&x^2-2=2}\Leftrightarrow \hoac{&x=0\\&x=\pm\sqrt{2}\\&x=\pm 2.}\]
Bảng biến thiên
\begin{center}
\begin{tikzpicture}[scale=1, line join=round, line cap=round, >=stealth]
	\tkzTabInit[nocadre=false,lgt=1.2,espcl=2,deltacl=0.6]
	{$x$/0.6,$y'$/0.6,$y$/2.5}
	{$-\infty$,$-2$,$-\sqrt{2}$,$0$,$\sqrt{2}$,$2$,$+\infty$}
	\tkzTabLine{,+,0,-,0,+,0,-,0,+,0,-,}
	\tkzTabVar{-/$-\infty$,+/$3$,-/$-1$,+/$3$,-/$-1$,+/$3$,-/$-\infty$}
\end{tikzpicture}
\end{center}
Ta thấy hàm số đồng biến trên các khoảng $(-\infty;-2)$, $(-\sqrt{2};0)$ và $(\sqrt{2};2)$.
}%<MyLT1>
\end{bt}

\begin{bt}%[2D1H1-3]
% nguon: Archive/2024/TaiLieuGiangDay/1.HamSo/1. K12 - 4in1 - Hamso/ThubaiGoc/Bai1-Dang3.tex #2763
%%[Mức độ 1]%[Dự án bài giảng new 4in1, Trần Quang Thạnh]
	Tìm tất cả các giá trị của tham số $m$ để hàm số $y=2^{mx}$ đồng biến trên $\mathbb{R}$.
	\loigiai{
	Ta có $y'=m\cdot 2^{mx}\cdot \ln2$.\\
	Nhận thấy $y'=0$ khi và chỉ khi $m=0$. Khi đó hàm số là hàm hằng.\\
	Do đó, hàm số đồng biến trên $\mathbb{R}$ khi $y' > 0$, với mọi $x\in\mathbb{R}$, tức là $m>0$.
	}
\end{bt}

\begin{bt}%[2D1H1-3]
% nguon: Archive/2024/TaiLieuGiangDay/1.HamSo/1. K12 - 4in1 - Hamso/ThubaiGoc/Bai1-Dang3.tex #2764
%%[Mức độ 1]%[Dự án bài giảng new 4in1, Trần Quang Thạnh]
	Tìm tất cả các giá trị của tham số $m$ để hàm số $y=\left(\dfrac{2}{3}\right)^{mx}$ nghịch biến trên $\mathbb{R}$.
	\loigiai{
	Ta có $y'=m\cdot \left(\dfrac{2}{3}\right)^{mx}\cdot \ln\dfrac{2}{3}$.\\
	Nhận thấy $y'=0$ khi và chỉ khi $m=0$. Khi đó hàm số là hàm hằng.\\
	Hàm số nghịch biến trên $\mathbb{R}$ khi $y'< 0$, với mọi $x\in \mathbb{R}$.\\
	Vì $\left(\dfrac{2}{3}\right)^{mx}\cdot \ln\dfrac{2}{3}<0, \forall x\in \mathbb{R}$ nên $y'\leq 0$, với mọi $x\in \mathbb{R}$ khi $m>0$.\\
	Vậy hàm số đã cho nghịch biến trên $\mathbb{R}$ khi $m>0$.
	}
\end{bt}

\begin{bt}%[2D1H1-3]
% nguon: Archive/2024/TaiLieuGiangDay/1.HamSo/1. K12 - 4in1 - Hamso/ThubaiGoc/Bai1-Dang3.tex #2765
%%[Mức độ 2]%[Dự án bài giảng new 4in1, Trần Quang Thạnh]
	Tìm các giá trị của tham số $m$ để hàm số $y=\dfrac{x^2-m}{x}$ đồng biến trên từng khoảng xác định.
	\loigiai{Tập xác định của hàm số là $(-\infty;0)\cup (0;+\infty)$.\\
	Ta có $y=x-\dfrac{m}{x}$ nên $y'=1+\dfrac{m}{x^2}=\dfrac{x^2+m}{x^2}$, với mọi $x\neq 0$.\\
	Hàm số đồng biến trên từng khoảng xác định khi và chỉ khi $y'\geq 0$, với mọi $x\neq 0$; tức là $x^2+m \geq 0$, với mọi $x\neq 0$; điều này tương đương với $m\geq 0$.
	}
\end{bt}

\begin{bt}%[2D1H1-3]
% nguon: Archive/2024/TaiLieuGiangDay/1.HamSo/1. K12 - 4in1 - Hamso/ThubaiGoc/Bai1-Dang3.tex #2766
%%[Mức độ 2]%[Dự án bài giảng new 4in1, Trần Quang Thạnh]
	Tìm các giá trị của tham số $m$ để hàm số $$y=-x^3+3x^2+3(m^2-1)x-3m^2-1$$ nghịch biến trên $\mathbb{R}$.
	\loigiai{Ta có $y'=-3x^2+6x+3(m^2-1)$.\\
		Vì $a=-3<0$ nên hàm số nghịch biến trên $\mathbb{R}$ khi $\Delta'=9+3\cdot 3(m^2-1)=9m^2\leq  0$ hay $m=0$.\\
		Vậy với $m=0$ thì hàm số nghịch biến trên $\mathbb{R}$.
	}	
\end{bt}

\begin{bt}%[2D1H1-3]
% nguon: Archive/2024/TaiLieuGiangDay/1.HamSo/1. K12 - 4in1 - Hamso/ThubaiGoc/Bai1-Dang3.tex #2767
%%[Mức độ 2]%[Dự án bài giảng new 4in1, Trần Quang Thạnh]
	Có bao nhiêu giá trị nguyên của tham số $m$ để hàm số $y=\dfrac{mx-4}{m-2x}$ nghịch biến trên từng khoảng xác định của nó?
	\loigiai{
		Ta có $y'=\dfrac{m^2-8}{(2x-m)^2}$, với mọi $x\neq \dfrac{m}{2}$.\\
		Hàm số đã cho nghịch biến trên từng khoảng xác định khi $y'<0$, với mọi $x\neq \dfrac{m}{2}$, tức là $m^2-8<0$ hay $-2\sqrt{2}<m<2\sqrt{2}$.\\
		Vậy có $5$ giá trị nguyên của tham số $m$ thỏa mãn yêu cầu bài toán là $-2; -1; 0; 1; 2$.
	}
\end{bt}

\begin{bt}%[2D1H1-3]
% nguon: Archive/2024/TaiLieuGiangDay/1.HamSo/1. K12 - 4in1 - Hamso/ThubaiGoc/Bai1-Dang3.tex #2768
%%[Mức độ 2]%[Dự án bài giảng new 4in1, Trần Quang Thạnh] 
	Cho hàm số $y=\dfrac{x-m^2}{-x+4m}$ ($m$ là tham số thực). Có bao nhiêu giá trị nguyên của $m$ để hàm số đã cho đồng biến trên từng khoảng xác định?
	\loigiai{
		Hàm số xác định khi $x\neq 4m$ và $y'=\dfrac{-m^2+4m}{(-x+4m)^2}$.\\
		Hàm số đồng biến từng khoảng xác định khi và chỉ khi $y'>0$, $\forall x\neq 4m$, tức là $-m^2+4m>0$ hay $0<m<4$.\\
		Do $m$ là các số nguyên nên $m\in\{1;2;3\}$.\\
		Vậy có $3$ giá trị nguyên của tham số $m$ thỏa mãn yêu cầu bài toán.		
	}
\end{bt}

\begin{bt}%[2D1H1-3]
% nguon: Archive/2024/TaiLieuGiangDay/1.HamSo/1. K12 - 4in1 - Hamso/ThubaiGoc/Bai1-Dang3.tex #2769
%%[Mức độ 2]%[Dự án bài giảng new 4in1, Trần Quang Thạnh]
	Có bao nhiêu giá trị nguyên của tham số $m$ để hàm số $y=\dfrac{x+2}{x+5m}$ đồng biến trên khoảng $(-\infty;-10)$?
	\loigiai{
		Ta có tập xác định $\mathscr{D}=\mathbb{R}\setminus \{-5m\}$ và
		$y'=\dfrac{5m-2}{(x+5m)^2}$. \\
		Hàm số đồng biến trên $(-\infty;-10) \Leftrightarrow \left\{\begin{aligned}
			&5m-2>0 \\
			&-5m\geqslant -10
		\end{aligned}\right. \Leftrightarrow \left\{\begin{aligned}
			&m>\dfrac{2}{5} \\
			&m\leqslant 2
		\end{aligned}\right. \Leftrightarrow \dfrac{2}{5}<m\leqslant 2$. \\
		Do $m\in \mathbb{Z}$ nên $m \in \{1;2\}$.\\
		Vậy có $2$ giá trị nguyên của tham số $m$ thỏa mãn yêu cầu bài toán.
	}
\end{bt}

\begin{bt}%[2D1H1-3]
% nguon: Archive/2024/TaiLieuGiangDay/1.HamSo/1. K12 - 4in1 - Hamso/ThubaiGoc/Bai1-Dang3.tex #2774
[Mức độ 2]%[Dự án bài giảng new 4in1, Trần Quang Thạnh]
	Tìm các giá trị của tham số $m$ để hàm số $y=\dfrac{a}{x+1}$ nghịch biến trên $(1;+\infty)$.
	\loigiai{
	Ta có $y'=-\dfrac{a}{(x+1)^2}$.\\
	Nhận thấy rằng, $y'=0$ trên $(1;+\infty)$ khi và chỉ khi $a=0$.\\
	Do đó hàm số nghịch biến trên $(1;+\infty)$ khi $y'<0$, với mọi $x>1$, tức là $a>0$.
	}
\end{bt}

\begin{bt}%[2D1H1-3]
% nguon: Archive/2024/TaiLieuGiangDay/1.HamSo/1. K12 - 4in1 - Hamso/ThubaiGoc/Bai1-Dang3.tex #2775
[Mức độ 2]%[Dự án bài giảng new 4in1, Trần Quang Thạnh]
	Tìm các giá trị của tham số $a$ để hàm số $y=-x^3+ax$ nghịch biến trên $\mathbb{R}$.
	\loigiai{
	Ta có $y'=-3x^2+a$.\\
	Hàm số nghịch biến trên $\mathbb{R}$ khi và chỉ khi $y'\leq 0$, với mọi $x\in \mathbb{R}$, điều này tương đương với $\Delta'_{y'}=3a \leq 0$ hay $a\leq 0$.\\
	Vậy với $a\leq 0$, hàm số đã cho nghịch biến trên $\mathbb{R}$.
	}
\end{bt}

\begin{bt}%[2D1H1-3]
% nguon: Archive/2024/TaiLieuGiangDay/1.HamSo/1. K12 - 4in1 - Hamso/ThubaiGoc/Bai1-Dang3.tex #2776
[Mức độ 2]%[Dự án bài giảng new 4in1, Trần Quang Thạnh]
	Tìm tất cả các giá trị của tham số $m$ để hàm số $y=\dfrac{x^5}{5}+\dfrac{mx^4}{2}+\dfrac{x^3}{3}+m-1$ đồng biến trên $\mathbb{R}$.
	\loigiai{
		Ta có $y'=x^4+2mx^3+x^2=x^2(x^2+2mx+1)$.\\
		Hàm số đồng biến trên $\mathbb{R}$ khi và chỉ khi $y'\geq 0, \forall x\in \mathbb{R}$, hay $x^2+2mx+1\geq 0, \forall x\in \mathbb{R}$; điều này tương đương với $m^2-1 \leq 0$, tức là $-1\leq m\leq 1$.\\
		Vậy hàm số đã cho đồng biến trên $\mathbb{R}$ khi và chỉ khi $-1\leq m\leq 1$.
	}
\end{bt}

\begin{bt}%[2D1H1-3]
% nguon: Archive/2024/TaiLieuGiangDay/1.HamSo/1. K12 - 4in1 - Hamso/ThubaiGoc/Bai1-Dang3.tex #2777
[Mức độ 2]%[Dự án bài giảng new 4in1, Trần Quang Thạnh]
	Tìm các giá trị của tham số $m$ để hàm số $y=\dfrac{mx+m-7}{5x-m+3}$ đồng biến trên mọi khoảng của tập xác định.
	\loigiai{
		Tập xác định $\mathscr{D}=\mathbb{R}\setminus \left\{\dfrac{m-3}{5}\right\}$.\\ Ta có $y'=\dfrac{-m^2-2m+35}{(5x-m+3)^2}$.\\
		Hàm số đồng biến trên mọi khoảng xác định khi và chỉ khi\\
		\centerline{$y'>0,\forall x\ne \dfrac{m-3}{5} \Leftrightarrow -m^2-2m+35>0 \Leftrightarrow m\in (-7;5)$.}
		Vậy, với $m\in (-7;5)$ thì hàm số đồng biến trên mọi khoảng xác định của nó.
	}
\end{bt}

\begin{bt}%[2D1V1-3]
% nguon: Archive/2024/TaiLieuGiangDay/1.HamSo/1. K12 - 4in1 - Hamso/ThubaiGoc/Bai1-Dang3.tex #2778
[Mức độ 3]%[Dự án bài giảng new 4in1, Trần Quang Thạnh]
	Tìm $a$ để hàm số $\displaystyle y = \frac{{{x^2} - 2ax + 4{a^2}}}{{x - 2a}}$ đồng biến trên $\left( {2; + \infty } \right)$.
	\loigiai{
		Tập xác định $\mathscr{D} = \mathbb{R}\backslash \left\{ 2a \right\}$ và $\displaystyle y' = \frac{{{x^2} - 4ax}}{{{{\left( {x - 2a} \right)}^2}}}$.
		
		Hàm số đồng biến trên $\left( {2; + \infty } \right) \Leftrightarrow y' \ge 0,\forall x \in \left( {2; + \infty } \right)$.
		
		Hay $\begin{cases}2a\notin (2;+\infty)\\x^2-4ax\geqslant 0,\forall x \in (2;+\infty)\end{cases}\Leftrightarrow \begin{cases}2a\leq  2\\a\leq  \dfrac{x}{4},\forall x\in (2;+\infty)\end{cases}\Leftrightarrow \begin{cases}
			a\leq  1\\a\leq  \dfrac{1}{2}\end{cases}\Leftrightarrow a\leq  \dfrac{1}{2}$.
		
		Vậy với $a\leq  \dfrac{1}{2}$, hàm số đồng biến trên $\left( {2; + \infty } \right)$.
	}
\end{bt}

\begin{bt}%[2D1N5-5]
% nguon: Archive/2024/TaiLieuGiangDay/1.HamSo/1. K12 - 4in1 - Hamso/ThubaiGoc/Bai5-Dang8.tex #2789

	\immini
	{
		Cho hàm số $ y=f(x) $ xác định và liên tục trên $ [-2;3] $ đồng thời hàm số $ y=f'(x) $ có đồ thị như hình vẽ bên. Xét tính đơn điệu của hàm số đã cho.
	}
	{
		\begin{tikzpicture}[>=stealth,x=1cm,y=1cm,scale=0.75,font=\footnotesize]
			%	Phương trình đường thẳng:
			\path
			(0,0) coordinate (O)
			%	Các điểm mút cho lệnh controls:
			(-2.5,0) coordinate (A1)
			(-1.75,1) coordinate (N)
			(0.2,-2) coordinate (Q1)
			(1.75,1) coordinate (M)
			(3.5,-3) coordinate (Q2)
			;
			%Vẽ đường cong, đường thẳng:
			%	\draw[red] (A1)--(N)--(Q1)--(M)--(Q2);
			\draw[thick,name path=fx] (A1)..controls +(70:0.25) and +(-175:0.25)..(N)..controls +(-5:0.5) and +(170:0.65)..(Q1)..controls +(5:0.5) and +(-170:0.65)..(M)..controls +(-5:0.65) and +(120:0.25)..(Q2)node[black,shift={(-0.7,0.2)}]{$ y=f'(x) $};
			\draw [domain=-3:3, samples=0.1,name path=gx] plot (\x, 0);
			\path[name intersections={of=fx and gx, by={-2,-1,1,2}}] ;
			%Vẽ hệ trục tọa dộ:
			\draw[->] (-3,0)--(0,0) node[above left]{$O$}--(4,0) node[below]{$x$};
			\draw[->] (0,-3.5) --(0,3) node[left]{$y$};
			%-2,-1,1,2,3
			\foreach \p/\r in {-2/-120,-1/-120,1/120,2/60}
			\draw[fill=black] (\p) circle (1.2pt) node[shift={(\r:2.5mm)}]{$\p$};
			\draw[dashed](Q2)--($(A1)!(Q2)!(O)$)node[above]{$ 3 $};
			\draw[fill=black]($(A1)!(Q2)!(O)$)  circle (1.2pt) ;
		\end{tikzpicture}
	}
	\loigiai{
		Dựa vào đồ thị $ y=f'(x) $, hàm số $ y=f(x) $ đồng biến trên các khoảng $ (-2;-1) $ và $ (1;2) $, nghịch biến trên các khoảng $ (-1;1) $ và $ (2;3) $.
	}
\end{bt}

\begin{bt}%[2D1N5-5]
% nguon: Archive/2024/TaiLieuGiangDay/1.HamSo/1. K12 - 4in1 - Hamso/ThubaiGoc/Bai5-Dang8.tex #2790

	\immini
	{
		Cho hàm số $ y=f(x) $ xác định và liên tục trên $ \mathbb{R} $. Biết đồ thị hàm số $ y=f'(x) $ tiếp xúc với trục hoành tại điểm có hoành độ $ x=1 $ và cắt trục hoành tại hai điểm có hoành độ lần lượt là $ x=-1 $, $ x=2 $ như hình vẽ bên. Xét tính đơn điệu của hàm số đã cho.
	}
	{
		\begin{tikzpicture}[>=stealth,x=1cm,y=1cm,scale=0.75,font=\footnotesize]
			%	Phương trình đường thẳng:
			\path
			(0,0) coordinate (O)
			%	Các điểm mút cho lệnh controls:
			(-1.3,1.7) coordinate (N)
			(-0.5,-1.5) coordinate (P)
			(1,0) coordinate (B1)
			(2,-1.5) coordinate (Q)
			(2.7,2.5) coordinate (M)
			;
			%Vẽ đường cong, đường thẳng:
			%	\draw[red] (N)--(P)--(B1)--(Q)--(M);
			\draw[thick,name path=fx] (N)..controls +(-95:0.5) and +(175:0.5)..(P) ..controls +(5:0.25) and +(-170:0.65)..(B1)..controls +(-5:0.5) and +(175:0.5)..(Q)..controls +(5:0.5) and +(100:0.25)..(M)node[black,shift={(0.7,-0.3)}]{$ y=f'(x) $};
			\draw [domain=-2:3, samples=0.1,name path=gx] plot (\x, 0);
			\path[name intersections={of=fx and gx, by={-1,1,2}}] ;
			%Vẽ hệ trục tọa dộ:
			\draw[->] (-2,0)--(0,0) node[above left]{$O$}--(4,0) node[below]{$x$};
			\draw[->] (0,-2) --(0,3) node[left]{$y$};
			%-2,-1,1,2,3
			\foreach \p/\r in {-1/-150,1/90,2/120}
			\draw[fill=black] (\p) circle (1.2pt) node[shift={(\r:2.5mm)}]{$\p$};
		\end{tikzpicture}
	}
	\loigiai{
		Dựa vào đồ thị $ y=f'(x) $, hàm số $ y=f(x) $ đồng biến trên các khoảng $ (-\infty;-1) $ và $ (2;+\infty) $, nghịch biến trên khoảng $ (-1;2) $.
	}
\end{bt}

\begin{bt}%[2D1N5-5]
% nguon: Archive/2024/TaiLieuGiangDay/1.HamSo/1. K12 - 4in1 - Hamso/ThubaiGoc/Bai5-Dang8.tex #2791

	\immini{Cho hàm số $y=f(x)$ có đạo hàm liên tục trên $[-2;2]$ và $y=f'(x)$ có đồ thị như hình bên. Trên khoảng $(-2;2)$, hàm số $y=f(x)$ có bao nhiêu điểm cực trị?}
	{
		\begin{tikzpicture}[smooth, >=latex, line cap =round, line join =round,scale=0.7,font=\scriptsize]
			\draw[->] (-3,0)--(3,0) node[below]{$x$};
			\draw[->] (0,-2.5) -- (0,3) node[below left] {$y$};
			\draw[thick, name path=dcong] (-2,-2)..controls +(90:0.3) and +(180:0.3)..(-0.7,2.7)..controls +(0:0.2) and +(180:0.3)..(0,0.5)..controls +(30:0.2) and +(180:0.3)..(1,2)..controls +(0:0.3) and +(90:0.1).. (2,-2); 
			\draw[thick,dashed] (-2,-2) -- (-2,0) node[above]{$-2$} (-1,0) node[below] {$-1$} --(-1,2) --(1,2) -- (1,0) node[below] {$1$} (2,0) node[above] {$2$}--(2,-2) (0,2) node[above right] {$2$};
			\draw (-1.6,1.8) node[rotate=75] {$y=f'(x)$};
		\end{tikzpicture}
	}
	\loigiai{
			Trên khoảng $(-2;2)$ đồ thị hàm số $y=f'(x)$ cắt trục hoành tại hai điểm phân biệt nên phương trình $f'(x)=0$ có hai nghiệm đơn. Suy ra hàm số $y=f(x)$ có hai điểm cực trị.
				
	}
\end{bt}

\begin{bt}%[2D1N5-5]
% nguon: Archive/2024/TaiLieuGiangDay/1.HamSo/1. K12 - 4in1 - Hamso/ThubaiGoc/Bai5-Dang8.tex #2792

	\immini{
		Cho hàm số $y=f'(x)$ có đồ thị như hình vẽ bên. Đồ thị hàm số $y=f(x)$ cắt trục hoành tại tối đa bao nhiêu điểm?
	}
	{\begin{tikzpicture}[>=stealth,x=1cm,y=1cm,scale=0.5]
			\def\a{0.5} % Hệ số a phải khác 0
			\def\b{-2}
			\def\c{0}
			\def\d{2}
			\draw[->] (-4,0) -- (6,0)node[below]{\scriptsize $x$};
			\draw[->] (0,-4) -- (0,5) node[left] {\scriptsize $y$};
			\draw (0,0)node[below left]{\scriptsize $O$};
			\clip (-4,-4)rectangle(6,5);
			\draw[thick,samples=150,smooth,domain=-5:6] plot(\x,{\a*(\x)^3+(\b)*(\x)^2+(\c)*\x+(\d)});
	\end{tikzpicture}}
	\loigiai{
		Gọi $a, b, c$ $(a<b<c)$ là hoành độ giao điểm của đồ thị hàm số $y=f'(x)$ và trục hoành, dựa vào đồ thị trên ta có bảng biến thiên
		\begin{center}
			% Cần khai báo \usepackage{tkz-tab}
			\begin{tikzpicture}[>=stealth]
				\tkzTabInit[nocadre=false,lgt=1,espcl=2,deltacl=0.5]{$x$/.7 ,$f'(x)$/.7,$f(x)$/2}
				{$-\infty$ , $a$ , $b$ , $c$ , $+\infty$}
				\tkzTabLine{ , - , $0$ , + , $0$ , - , $0$ , + , }
				\tkzTabVar{+/$+\infty$ , -/$f(a)$, +/$f(b)$ , -/$f(c)$ , +/$+\infty$}
			\end{tikzpicture}
		\end{center}
		Dựa vào bảng biến thiên ta thấy đồ thị hàm số $y=f(x)$ cắt trục hoành tại tối đa $4$ điểm.}
\end{bt}

\begin{bt}%[2D1H5-5]
% nguon: Archive/2024/TaiLieuGiangDay/1.HamSo/1. K12 - 4in1 - Hamso/ThubaiGoc/Bai5-Dang8.tex #2793

	\immini
	{
		Cho hàm số $ y=f(x) $ có đạo hàm liên tục trên đoạn $ [-2;-2] $ và đồ thị hàm số $ y=f'(x) $ như hình vẽ. Tìm giá trị nhỏ nhất và  giá trị lớn nhất hàm số $ y=f(x) $ trên đoạn $ [-2;2] $, biết $ f(-1)=4 $, $ f(-2)=-5 $, $ f(2)=-3 $.
	}
	{
		\begin{tikzpicture}[>=stealth,x=1cm,y=1cm,scale=0.9,font=\footnotesize]
			\path
			(2,0) coordinate (B2)
			(-1.7,1.7) coordinate (N)node[shift={(0.8,-0.3)}]{\scriptsize$ y=f'(x)$}
			(-0.7,-0.7) coordinate (P)
			(0,0) coordinate (O)node[above right]{$ O$}
			(1.7,-1.7) coordinate (Q)
			($(B2)!(N)!(O)$) coordinate (A1)node[below]{$ -2 $}
			($(B2)!(Q)!(O)$) coordinate (B1)node[above]{$ 2 $}
			;
			%	\draw[red] (N)--(P)--(O)--(Q);
			\draw[name path=ox] (O)--(-2,0);
			\draw[name path=fx,thick,black]
			(N)..controls +(-70:1) and+(170:0.25)..
			(P)..controls +(9:0.25) and+(-170:0.25)..(O)
			..controls +(-5:0.35) and+(140:0.25)..(Q)
			;
			\draw[dashed](N)--(A1) (Q)--(B1);
			\path[name intersections={of=fx and ox, by={A3,A2}}] ;
			%Vẽ hệ trục tọa dộ:
			\draw[->] (-2,0)--(0,0)--(2,0) node[below]{$x$};
			\draw[->] (0,-2) --(0,2.2) node[left]{$y$};
			%	Node các điểm
			\foreach \p in {N,O,Q}
			\fill (\p) circle (1.5pt) ;
			\fill(A3) circle (1.5pt) node[above right]{$ -1 $};
		\end{tikzpicture}
	}
	
	\loigiai{
		Dựa vào đồ thị ta có  $ f'(x)=0 \Leftrightarrow\hoac{& x=-1 \\ & x=0.}$\\
		Ta có bảng biến thiên của hàm số $ y=f(x) $:
		\begin{center}
			\begin{tikzpicture}
				\tkzTabInit[nocadre=false,lgt=1.2,espcl=2.,deltacl=0.7]
				{$x$ /0.6,$f’(x)$ /0.6,$f(x)$ /2}
				{$-2$,$-1$,$0$,$2$}
				\tkzTabLine{,+,0,-,0,-,}
				\tkzTabVar{-/$ -5 $\  , +/$ 4 $  , R/ , -/$ -3 $  }
				%\tkzTabVar{-/ $f(-2)$  , +/ $f(-1)$ / , R/ , -/ $f(2)$ / }
			\end{tikzpicture}
		\end{center}
		Vậy
		$ \max\limits_{[-2; 2]}f(x)=4, \min\limits_{[-2; 2]}f(x)=-5  $.
	}
\end{bt}

\begin{bt}%[2D1V5-5]
% nguon: Archive/2024/TaiLieuGiangDay/1.HamSo/1. K12 - 4in1 - Hamso/ThubaiGoc/Bai5-Dang8.tex #2794

	\immini
	{
		Cho hàm số $ y=f(x) $ có đạo hàm liên tục trên $ \mathbb{R} $. Biết đồ thị hàm số $ y=f'(x) $ chỉ cắt trục hoành tại ba điểm như hình vẽ bên. Xét tính đơn điệu của hàm số $ g(x)=2 f(x)-x^{2} $ trên khoảng $ (-1;1) $.
	}
	{
		\begin{tikzpicture}[>=stealth,x=1cm,y=1cm,scale=0.8,font=\footnotesize]
			\def\a{1}
			\path
			(0,0) coordinate (O)
			(-1,0) coordinate (A)node[above]{$ -1 $}
			(1,0) coordinate (B)node[below]{$ 1 $}
			(0,-1) coordinate (C)node[right]{$ -1 $}
			(0,1) coordinate (D)node[left]{$ 1 $}
			%	Các điểm mút cho lệnh controls:
			(-1,-1) coordinate (P)
			(-0.3,0.5) coordinate (N)
			(0.3,-0.5) coordinate (Q)
			(1,1) coordinate (M)
			;
			%Vẽ đường cong, đường thẳng:
			%		\draw[red] (P)--(N)--(Q)--(M);
			\draw[thick,name path=fx,shorten >=-1cm, shorten <=-0.5cm] (P)..controls +(75:0.5) and +(185:0.25)..(N)..controls +(-10:0.3) and +(170:0.3)..(Q)..controls +(2:0.35) and +(-105:0.25)..(M)node[black,shift={(0.8,0.3)}]{$ y=f'(x) $};
			%Vẽ hệ trục tọa dộ:
			\draw[->] (-1.5,0)--(0,0) node[below left]{$O$}--(2,0) node[below]{$x$};
			\draw[->] (0,-1.5) --(0,2.35) node[left]{$y$};
			%-2,-1,1,2,3
			\draw[dashed](B)|-(D) (A)|-(C);
			\foreach \p in {M,P,O,A,B,C,D}
			\fill (\p) circle (1pt);
		\end{tikzpicture}
	}
	\loigiai{
		\immini
		{
			Tập xác định: $\mathscr{D} = \mathbb{R}$.\\
			$g'(x)=2\cdot f'(x) -2x$.\\
			Số nghiệm của phương trình $ g'(x)=0 $ chính là số giao điểm của đồ thị hàm số $ y=f'(x) $ với đường thẳng $ y=x $ (như hình vẽ bên).\\
			$ g'(x)=0 \Leftrightarrow f'(x)=x \Leftrightarrow \hoac{& x=-1 \\ & x=0 \\ & x=1.}$\\
		}
		{
			\begin{tikzpicture}[>=stealth,scale=0.85,font=\footnotesize]
				\def\a{1}
				\path
				(0,0) coordinate (O)
				(-1,0) coordinate (A)node[above]{$ -1 $}
				(1,0) coordinate (B)node[below]{$ 1 $}
				(0,-1) coordinate (C)node[right]{$ -1 $}
				(0,1) coordinate (D)node[left]{$ 1 $}
				%	Các điểm mút cho lệnh controls:
				(-1,-1) coordinate (P)
				(-0.3,0.5) coordinate (N)
				(0.3,-0.5) coordinate (Q)
				(1,1) coordinate (M)
				;
				%Vẽ đường cong, đường thẳng:
				%		\draw[red] (P)--(N)--(Q)--(M);
				\draw[thick,name path=fx,shorten >=-1cm, shorten <=-0.5cm] (P)..controls +(75:0.5) and +(185:0.25)..(N)..controls +(-10:0.3) and +(170:0.3)..(Q)..controls +(2:0.35) and +(-105:0.25)..(M)node[black,shift={(1,0.9)}]{$ y=f'(x) $};
				\draw[purple,shorten >=-0.9cm, shorten <=-0.75cm](P)--(M)node[black,shift={(0.9,0.35)}]{$ y=x $};
				%Vẽ hệ trục tọa dộ:
				\draw[->] (-1.5,0)--(0,0) node[above left]{$O$}--(2,0) node[below]{$x$};
				\draw[->] (0,-1.5) --(0,2.35) node[left]{$y$};
				%-2,-1,1,2,3
				\draw[dashed](B)|-(D) (A)|-(C);
				\foreach \p in {M,P,O,A,B,C,D}
				\fill (\p) circle (1pt);
			\end{tikzpicture}
		}
		Lập bảng biến thiên
		\begin{center}
			\begin{tikzpicture}[yscale=0.55,xscale=0.9]
				\def\c{9}%số cột
				\def\d{4}%số dòng
				\foreach \x in {0,...,\c}{
					\foreach \y in {0,...,\d}
					\path[yscale=-1] (\x,\y) node[gray!0,minimum size=1cm] (\x\y) {\x\y};
				}
				
				\draw[shift={(-0.5,0.5)}]
				(0,0) rectangle (\c+1,-\d-1)
				(0,-1)--(\c+1,-1)
				(0,-2)--(\c+1,-2)
				(1,0)--(1,-\d-1);
				
				\path
				(00) node {$x$}(10) node {$-\infty$}(30) node {$-1$}(50) node {$0$}(70) node {$1$}(90) node {$+\infty$}
				(01) node {$g'(x)$}(21) node {$-$}(31) node {$0$}(41) node {$+$}(51) node {$0$}(61) node {$-$}(71) node {$0$}(81) node {$+$}
				(03) node {$g(x)$}
				;
				
				\foreach \a/\b in {12/34,34/52,52/74,74/92}
				\draw[->,>=stealth,shorten >=-0.36cm, shorten <=-0.36cm] (\a)--(\b);
			\end{tikzpicture}
		\end{center}
	}
	Dựa vào bảng biến thiên, hàm số $ g(x)=2 f(x)-x^{2} $ nghịch biến trên các khoảng $ (-\infty;-1) $ và $ (0;1) $; đồng biến trên các khoảng $ (-1;0) $ và $ (1;+\infty) $.
\end{bt}

\begin{bt}%[2D1V5-5]
% nguon: Archive/2024/TaiLieuGiangDay/1.HamSo/1. K12 - 4in1 - Hamso/ThubaiGoc/Bai5-Dang8.tex #2795

	\immini
	{
		Cho hàm số $ y=f(x) $ có đạo hàm liên tục trên đoạn $ [-2;-2] $ và đồ thị hàm số $ y=f'(x) $ như hình vẽ. Tìm giá trị nhỏ nhất và  giá trị lớn nhất của hàm số $ g(x)=f^{2}(x)-4 f(x) $ trên đoạn $ [-2;2] $, biết $ f(-1)=1 $, $ f(-2)=-1 $, $ f(2)=-2 $.
	}
	{
		\begin{tikzpicture}[>=stealth,x=1cm,y=1cm,scale=0.9,font=\footnotesize]
			\path
			(2,0) coordinate (B2)
			(-1.7,1.7) coordinate (N)node[shift={(0.8,-0.3)}]{\scriptsize$ y=f'(x)$}
			(-0.7,-0.7) coordinate (P)
			(0,0) coordinate (O)node[above right]{$ O$}
			(1.7,-1.7) coordinate (Q)
			($(B2)!(N)!(O)$) coordinate (A1)node[below]{$ -2 $}
			($(B2)!(Q)!(O)$) coordinate (B1)node[above]{$ 2 $}
			;
			%	\draw[red] (N)--(P)--(O)--(Q);
			\draw[name path=ox] (O)--(-2,0);
			\draw[name path=fx,thick,black]
			(N)..controls +(-70:1) and+(170:0.25)..
			(P)..controls +(9:0.25) and+(-170:0.25)..(O)
			..controls +(-5:0.35) and+(140:0.25)..(Q)
			;
			\draw[dashed](N)--(A1) (Q)--(B1);
			\path[name intersections={of=fx and ox, by={A3,A2}}] ;
			%Vẽ hệ trục tọa dộ:
			\draw[->] (-2,0)--(0,0)--(2,0) node[below]{$x$};
			\draw[->] (0,-2) --(0,2.2) node[left]{$y$};
			%	Node các điểm
			\foreach \p in {N,O,Q}
			\fill (\p) circle (1.5pt) ;
			\fill(A3) circle (1.5pt) node[above right]{$ -1 $};
		\end{tikzpicture}
	}
	
	\loigiai{
		Ta có bảng biến thiên của hàm số $ y=f(x) $
		\begin{center}
			\begin{tikzpicture}
				\tkzTabInit[nocadre=false,lgt=1.2,espcl=2.,deltacl=0.7]
				{$x$ /0.6,$f’(x)$ /0.6,$f(x)$ /2}
				{$-2$,$-1$,$0$,$2$}
				\tkzTabLine{,+,0,-,0,-,}
				\tkzTabVar{-/$ -1 $\  , +/$ 1 $  , R/ , -/$ -2 $  }
			\end{tikzpicture}
		\end{center}
		Từ bảng biến thiên $ \Rightarrow f(x)\leq 1,\quad \forall x \in [-2;2] $.\\
		Ta có $ g'(x)=2f(x)\cdot f'(x) -4 f'(x)=2f'(x)\cdot\left [f(x)-2\right ] $.\\
		Cho $ g'(x)=0 \Leftrightarrow\hoac{& f'(x)=0&(1) \\ & f(x)=2&(2)} $\\
		Dựa vào đồ thị, $ (1) \Leftrightarrow \hoac{& x=-1 \\ & x=0.}$\\
		Mặt khác ta có $ f(x)\leq 1, \forall x \in [-2;2] \Rightarrow f(x)-2<0,  \forall x \in [-2;2] $.\\
		Vậy phương trình $ (2) $ vô nghiệm.\\
		Do đó, phương trình $ g'(x)=0 $ có hai nghiệm là $ x=-1 $, $ x=0 $.
		\begin{center}
			\begin{tikzpicture}
				\tkzTabInit[nocadre=false,lgt=1.2,espcl=2.,deltacl=0.7]
				{$x$ /0.6,$g’(x)$ /0.6,$g(x)$ /2.5}
				{$-2$,$-1$,$0$,$2$}
				\tkzTabLine{,-,0,+,0,+,}
				\tkzTabVar{+/$ g(-2) $\  , -/$ g(-1) $  , R/ , +/$ g(2) $  }
			\end{tikzpicture}
		\end{center}
		Ta lại có $ g(-1)=-3$, $ g(-2)=5$, $ g(2)=12$.\\
		Vậy $ \max\limits_{[-2; 2]}g(x)=g(2)=12, \min\limits_{[-2; 2]}g(x)=g(-1)=-3  $.
	}
\end{bt}

\begin{bt}%[2D1H5-5]
% nguon: Archive/2024/TaiLieuGiangDay/1.HamSo/1. K12 - 4in1 - Hamso/ThubaiGoc/Bai5-Dang8.tex #2799

	\immini{Cho hàm số $y=f(x)$ xác định, liên tục trên $\mathbb{R}$ và có đạo hàm $f’(x)$. Biết rằng hàm số $f’(x)$ có đồ thị như hình vẽ bên. Hàm số đồng biến trên khoảng nào?
		}
	{\begin{tikzpicture}[>=stealth,line join=round,line cap=round,font=\normalsize,scale=0.6]
			\draw[-stealth] (-4,0)--(0,0)node[above left]{$O$}--(1.5,0)node[below]{$x$};
			\draw[-stealth] (0,-2.5)--(0,2)node[left]{$y$};
			\draw
			(-3.2,-2)..controls +(80:.5)and+(180:.35)..(-2.5,1)
			..controls+(0:.25) and+(105:.25)..(-2,0)
			..controls+(-75:.25) and+(180:.5)..(-1.2,-1.75)
			..controls+(0:.5) and+(180:.25)..(0,0)
			..controls+(0:.25) and+(-90:.25)..(0.5,-2)
			;
			\draw 
			(-3,0) node[below left]{$-3$}
			(-2,0) node[below right]{$-2$}
			;
	\end{tikzpicture}}
	\loigiai{
		Quan sát đồ thị ta thấy $f’(x)>0$ trên khoảng $(-3;-2)$ nên hàm số $y=f(x)$ đồng biến  trên khoảng $(-3;-2)$.}
\end{bt}

\begin{bt}%[2D1H5-5]
% nguon: Archive/2024/TaiLieuGiangDay/1.HamSo/1. K12 - 4in1 - Hamso/ThubaiGoc/Bai5-Dang8.tex #2800

	\immini{Cho hàm số $y=f(x)$ xác định và liên tục trên $\mathbb{R}$ và có đồ thị của đạo hàm $y=f’(x)$ như hình bên.  Hàm số nghịch biến trên khoảng nào? 
	}
	{\begin{tikzpicture}[>=stealth,line join=round,line cap=round,font=\normalsize,scale=0.6]
			\draw[->] (-5,0)--(4,0)node[below]{$x$};
			\foreach \x in {-4,-3,-2,-1,1,2,3}
			\draw[shift={(\x,0)}] (0,2pt)--(0,-2pt) node[below]{\scriptsize $\x$};
			\draw[->] (0,-3)--(0,1)node[right]{$y$};
			\foreach \y in {-2,-1}
			\draw[shift={(0,\y)}] (2pt,0)--(-2pt,0) node[right]{\scriptsize $\y$};
			\draw (0,0)node[below left]{$0$};
			\draw
			(-5,-1)..controls +(80:.25)and+(180:.75)..(-3,0.5)
			..controls+(0:.5) and+(130:.5)..(-2,0)
			..controls+(-50:.5) and+(180:1)..(1,-2.5)
			..controls+(0:1) and+(-90:.25)..(3,0)--(3.25,1)
			;
	\end{tikzpicture}}
	\loigiai{
		Quan sát đồ thị ta thấy $f’(x)<0$ trên khoảng $(-\infty;-4)$ và $(-2;3)$ nên hàm số $y=f(x)$ nghịch biến  trên khoảng $(-\infty;-4)$ và $(-2;3)$.}
\end{bt}

\begin{bt}%[2D1H5-5]
% nguon: Archive/2024/TaiLieuGiangDay/1.HamSo/1. K12 - 4in1 - Hamso/ThubaiGoc/Bai5-Dang8.tex #2801

	\immini{Cho hàm số $y=f(x)$ có đạo hàm liên tục trên $(a,b)$ và $y=f'(x)$ có đồ thị như hình vẽ bên. Trên khoảng $(a,b)$, hàm số $y=f(x)$ có bao nhiêu điểm cực trị?}
	{
		\begin{tikzpicture}[smooth, >=latex, line cap =round, line join =round,scale=0.7,font=\scriptsize]
			\draw[->] (-2.5,0)--(3.5,0) node[below]{$x$};
			\draw[->] (0,-.5) -- (0,4) node[below left] {$y$};
			\draw[thick, name path=dcong] (-1.5,3)..controls +(-90:0.2) and +(180:0.4).. (-0.3,0.3) .. controls + (0:0.4) and +(180:0.2)..(1,2)..controls +(0:0.09) and +(180:0.1)..(1.4,1.5)..controls +(0:0.3) and +(-135:0.1)..(2.5,3);
			\draw[dashed] (-1.5,3) -- (-1.5,0) node[below]{$a$} (2.5,3) node[right]{$y=f'(x)$} -- (2.5,0) node[below] {$b$} (0,0) node[below left] {$O$};
		\end{tikzpicture}
	}
	\loigiai{
		Ta thấy đồ thị $f'(x)$ không cắt trục hoành nên phương trình $f'(x)=0$ vô nghiệm.\\
		Do đó hàm số $y=f(x)$ không có điểm cực trị.
	}
\end{bt}

\begin{bt}%[2D1H5-5]
% nguon: Archive/2024/TaiLieuGiangDay/1.HamSo/1. K12 - 4in1 - Hamso/ThubaiGoc/Bai5-Dang8.tex #2802

	\immini{Cho hàm số $y=f(x)$ có đạo hàm liên tục trên $(a,b)$ và $y=f'(x)$ có đồ thị như hình vẽ bên. Trên khoảng $(a,b)$, hàm số $y=f(x)$ có bao nhiêu điểm cực trị.
	}
	{
		\begin{tikzpicture}[smooth, >=latex, line cap =round, line join =round,scale=0.7,font=\scriptsize]
			\draw[->] (-2.5,1.5)--(3.5,1.5) node[below]{$x$};
			\draw[->] (0,-.5) -- (0,4) node[below left] {$y$};
			\draw[thick, name path=dcong] (-1.5,3)..controls +(-90:0.2) and +(180:0.4).. (-0.3,0) .. controls + (0:0.4) and +(180:0.2)..(1,2)..controls +(0:0.09) and +(180:0.1)..(1.4,1.5)..controls +(0:0.3) and +(-135:0.1)..(2.5,3);
			\draw[dashed] (-1.5,3) -- (-1.5,1.5) node[below]{$a$} (2.5,3) node[right]{$y=f'(x)$} -- (2.5,1.5) node[below] {$b$} (0,1.5) node[above left] {$O$};
		\end{tikzpicture}
	}
	\loigiai{
		Ta thấy đồ thị $f'(x)$ cắt trục hoành tại hai điểm phân biệt và tiếp xúc tại một điểm. Suy ra phương trình $f'(x)=0$ có hai nghiệm đơn.\\
		Do đó hàm số $y=f(x)$ có hai điểm cực trị.}
\end{bt}

\begin{bt}%[2D1V5-5]
% nguon: Archive/2024/TaiLieuGiangDay/1.HamSo/1. K12 - 4in1 - Hamso/ThubaiGoc/Bai5-Dang8.tex #2803

	Cho hàm số $f(x)$, đồ thị của hàm số $y=f'(x)$ là đường cong trong hình vẽ dưới
	\begin{center}
		\begin{tikzpicture}[line join=round, line cap=round,>=stealth,thick,scale=.8]
			\tikzset{label style/.style={font=\footnotesize}}
			\def \xmin{-3}
			\def \xmax{5}
			\def \ymin{-2}
			\def \ymax{7}
			\draw[step=1,gray,very thin,dashed]
			(\xmin,\ymin) grid (\xmax,\ymax);
			\def \hamso{-0.155*(\x)^5+0.83*(\x)^4-0.08*(\x)^3-3.49*(\x)^2+1.65*\x+4}
			\draw[->] (\xmin,0)--(\xmax,0) node[below left] {$x$};
			\draw[->] (0,\ymin)--(0,\ymax) node[below left] {$y$};
			\draw (0,0) node [below left] {$O$};
			\foreach \x in {-1,2,4}
			\draw[thin] (\x,1pt)--(\x,-1pt) node [below left] {$\x$};
			\foreach \y in {-1,1,2,4}
			\draw[thin] (1pt,\y)--(-1pt,\y) node [left] {$\y$};
			\begin{scope}
				\clip (\xmin+0.01,\ymin+0.01) rectangle (\xmax-0.01,\ymax-0.01);
				\draw[samples=200,domain=\xmin+0.01:\xmax-0.01,smooth,variable=\x] plot (\x,{\hamso});
			\end{scope}
			%\draw[dashed] (0,0.9)--(2,0.9)--(2,0) (0,2)--(4,2)--(4,0);
		\end{tikzpicture}
	\end{center}	
	Xét hàm số $g(x)=f(2x+1)-2x^2$ trên đoạn $[-2;5]$. Hãy so sánh các giá trị $g\left(\dfrac{3}{2}\right)$, $g\left(\dfrac{1}{2}\right)$ và $g(-2)$.
	\loigiai{
		\immini
		{
			Ta có $g'(x)=2f'(2x+1)-4x$.\\
			$g'(x)=0 \Leftrightarrow 2f'(2x+1)-4x=0\Leftrightarrow f'(2x+1)=2x.$\hfill$(1)$\\
			Đặt $t=2x+1$ vì $x\in[-2;2]\Rightarrow t\in[-3;5]$.\\
			$(1)\Leftrightarrow f'(t)=t-1\Leftrightarrow\hoac{&t=2\\&t=4}\forall t\in[-3;5].$ 
		}
		{
			\begin{tikzpicture}[line join=round, line cap=round,>=stealth,thick,scale=.7]
				\tikzset{label style/.style={font=\footnotesize}}
				\def \xmin{-3}
				\def \xmax{5}
				\def \ymin{-2}
				\def \ymax{7}
				\draw[step=1,gray,very thin,dashed]
				(\xmin,\ymin) grid (\xmax,\ymax);
				\def \hamso{-0.155*(\x)^5+0.83*(\x)^4-0.08*(\x)^3-3.49*(\x)^2+1.65*\x+4}
				\draw[->] (\xmin,0)--(\xmax,0) node[below left] {$x$};
				\draw[->] (0,\ymin)--(0,\ymax) node[below left] {$y$};
				\draw (0,0) node [below left] {$O$};
				\foreach \x in {-1,2,4}
				\draw[thin] (\x,1pt)--(\x,-1pt) node [below left] {$\x$};
				\foreach \y in {-1,1,2,4}
				\draw[thin] (1pt,\y)--(-1pt,\y) node [left] {$\y$};
				\begin{scope}
					\clip (\xmin+0.01,\ymin+0.01) rectangle (\xmax-0.01,\ymax-0.01);
					\draw[samples=200,domain=\xmin+0.01:\xmax-0.01,smooth,variable=\x] plot (\x,{\hamso});
					\draw[samples=200,domain=\xmin+0.01:\xmax-0.01,smooth,variable=\x] plot (\x,{\x-1});
				\end{scope}
				\draw[dashed] (0,3)--(4,3)--(4,0);
			\end{tikzpicture}
		}
		Bảng biến thiên
		\begin{center}
			\begin{tikzpicture}
				\tkzTabInit[lgt=1.2,espcl=2.5]
				{$x$/1,$t$/1,$g'$/1,$g$/2}
				{$-2$,$\dfrac{1}{2}$,$\dfrac{3}{2}$,$5$}
				\tkzTabLine{$-3$,,$2$,,$4$,,$8$}
				\tkzTabLine{ ,+,0,+,0,-, }
				\tkzTabVar{-/,R,+/,-/}
			\end{tikzpicture}
		\end{center}
		Dựa vào bảng biến thên ta thấy
		$g\left(\dfrac{3}{2}\right)>g\left(\dfrac{1}{2}\right)>g(-2)$.}
\end{bt}

\begin{bt}%[2D1N5-6]
% nguon: Archive/2024/TaiLieuGiangDay/1.HamSo/1. K12 - 4in1 - Hamso/ThubaiGoc/Bai5-Dang9.tex #2815
%[Dự Án Bài Giảng 12 4 in 1, Lê Văn Toàn]
	Viết phương trình tiếp tuyến của đồ thị hàm số $y=x^3-3x^2+1$ tại điểm $A(3;1)$.
	\loigiai{
		Ta có $y'=3x^2-6x$, suy ra $y'(3)=9$.\\
		Tiếp tuyến của đồ thị hàm số tại $A(3;1)$ là
		$$y=y'(3)(x-3)+1=9(x-3)+1=9x-26.$$}
\end{bt}

\begin{bt}%[2D1H5-6]
% nguon: Archive/2024/TaiLieuGiangDay/1.HamSo/1. K12 - 4in1 - Hamso/ThubaiGoc/Bai5-Dang9.tex #2816
%[GHK2-KSCL-Lan1-SGD-BacNinh-2022-2023]
		Viết phương trình tiếp tuyến của đồ thị hàm số $y=x^3-2x+5$ tại điểm có hoành độ $x_0=1$ có phương trình là
	\loigiai{
		Ta có $y'=3x^2-2\Rightarrow{y}'(1)=3-2=1$, và $ y(1)=1^3-2\cdot 1+5=4$.\\
		Phương trình tiếp tuyến với đồ thị hàm số $y=x^3-2x+5$ tại điểm có hoành độ $x_0=1$ là $$ y=y'(1)\left(x-1\right)+y(1)\Leftrightarrow y=1\cdot\left(x-1\right)+4\Leftrightarrow y=x+3.$$
	}
\end{bt}

\begin{bt}%[2D1H5-6]
% nguon: Archive/2024/TaiLieuGiangDay/1.HamSo/1. K12 - 4in1 - Hamso/ThubaiGoc/Bai5-Dang9.tex #2817
%[Dự Án Bài Giảng 12 4 in 1, Lê Văn Toàn]
	Cho hàm số $y=x^3-x-1$ có đồ thị $(C)$. Viết phương trình tiếp tuyến của $(C)$ tại giao điểm của $(C)$ với trục tung.
	\loigiai{
		Gọi tiếp điểm $M(x_0;y_0)$. Vì $M(x_0;y_0)$ là giao điểm của $(C)$ với trục tung nên $M(0;-1)$.\\
		Hệ số góc của tiếp tuyến là $y'(0)=3\cdot 0^2-1=-1$.\\
		Vậy phương trình tiếp tuyến là $y=y'(0)\left( x-x_0\right)+y_0\Leftrightarrow y=-x+1$.}
\end{bt}

\begin{bt}%[2D1H5-6]
% nguon: Archive/2024/TaiLieuGiangDay/1.HamSo/1. K12 - 4in1 - Hamso/ThubaiGoc/Bai5-Dang9.tex #2819
%[Dự Án Bài Giảng 12 4 in 1, Lê Văn Toàn]
	Viết phương trình tiếp tuyến của đường cong $y=x^3+3x^2-2$ tại điểm có hoành độ $x_0=1$.
	\loigiai{
		Ta có $y=x^3+3x^2-2\Rightarrow y'=3x^2+6x$.\\
		Tiếp điểm có hoành độ $x_0=1$ nên tung độ tiếp điểm là $y_0=y(1)=2$.\\
		Hệ số góc của tiếp tuyến là $y'(x_0)=y'(1)=9$.\\
		Vậy phương trình tiếp tuyến tại $(1;2)$ là $y=y'(x_0)(x-x_0)+y_0\Leftrightarrow y=9(x-1)+2\Leftrightarrow y=9x-7$.
	}
\end{bt}

\begin{bt}%[2D1H5-6]
% nguon: Archive/2024/TaiLieuGiangDay/1.HamSo/1. K12 - 4in1 - Hamso/ThubaiGoc/Bai5-Dang9.tex #2821
%[Dự Án Bài Giảng 12 4 in 1, Lê Văn Toàn]
	Cho hàm số $y=\dfrac{x+2}{x+1}$ có đồ thị $(C)$. Viết hương trình tiếp tuyến với đồ thị hàm số tại giao điểm của đồ thị $(C)$ với trục tung.
	\loigiai
	{
		Ta có $y'=f'(x)=\dfrac{-1}{(x+1)^2}$\\
		Gọi $M(a;b)$ là giao điểm của đồ thị $(C)$ với trục tung $\Rightarrow M(0;2)$ và $f'(0)=-1$.\\
		Phương trình tiếp tuyến có dạng
		\[d\colon y=f'(0)\cdot(x-0)+2=-1(x-0)+2=-x+2.\]
	}
\end{bt}

\begin{bt}%[2D1V5-6]
% nguon: Archive/2024/TaiLieuGiangDay/1.HamSo/1. K12 - 4in1 - Hamso/ThubaiGoc/Bai5-Dang9.tex #2822
%[Dự Án Bài Giảng 12 4 in 1, Lê Văn Toàn]
	Tìm điểm $M$ có hoành độ âm trên đồ thị $(C): y=\dfrac{1}{3}x^3-x+\dfrac{2}{3}$ sao cho tiếp tuyến tại $M$ vuông góc với đường thẳng $y=-\dfrac{1}{3}x+\dfrac{2}{3}$.
	\loigiai{
		Gọi $M\left(x_0; \dfrac{1}{3}x_0^3-x_0+\dfrac{2}{3}\right)$. \\
		Do tiếp tuyến tại $M$ vuông góc với đường thẳng $y=-\dfrac{1}{3}x+\dfrac{2}{3}$ nên ta có hệ số góc của tiếp tuyến tại $M$ là $k=3$.\\
		Ta có $y'=x^2-1$. Theo đề bài ta có phương trình $x_0^2-1=3 \Leftrightarrow x_0^2=4 \Leftrightarrow x_0=\pm 2$. \\
		Theo đề bài điểm $M$ có hoành độ âm nên $M(-2; 0)$.}
\end{bt}

\begin{bt}%[2D1V5-6]
% nguon: Archive/2024/TaiLieuGiangDay/1.HamSo/1. K12 - 4in1 - Hamso/ThubaiGoc/Bai5-Dang9.tex #2823
%[Dự Án Bài Giảng 12 4 in 1, Lê Văn Toàn]
	Cho hàm số $y = \dfrac{x+2}{2x+3}$, có đồ thị $\left( H \right)$. Viết phương trình tiếp tuyến của đồ thị hàm số $\left( H \right)$, biết tiếp tuyến đó cắt trục hoành, trục tung lần lượt tại hai điểm phân biệt $A$, $B$ và tam giác $OAB$ cân tại gốc tọa độ $O$.
	\loigiai{
		Tam giác $OAB$ vuông cân tại $O$ nên hệ số góc của tiếp tuyến bằng $ \pm 1$.\\
		Gọi tọa độ tiếp điểm là $(x_0;y_0)$ ta có  $\dfrac {-1}{(2x_0 + 3)^2}=\pm 1 \Leftrightarrow x_0 = - 2$ hoặc $x_0 =- 1$.
		\begin{itemize}
			\item Với $x_0 = - 1$, $y_0 = - 1$, phương trình tiếp tuyến là $y = - x$ (loại).
			\item Với $x_0 =  - 2$, $y_0 = 0$, phương trình tiếp tuyến là $y = - x- 2$ (nhận).
		\end{itemize}
		Vậy phương trình tiếp tuyến của đồ thị hàm số $(H)$ là $y = - x - 2$.}
\end{bt}

\begin{bt}%[2D1H2-1]
% nguon: Archive/2024/TaiLieuGiangDay/1.HamSo/1. K12 - 4in1 - Hamso/Thubai/Bai2-Dang1_converted_converted.tex #2890
%%[SGK12-CKP, Mức độ 2]
Tìm cực trị của hàm số sau
$y=\left|x\right|$.
\loigiai{
Cách 1:
Hàm số $y=\left|x\right|$ xác định với mọi $x\in\mathbb{R}$.\\
Khi $x > 0$, ta có $y=x$ và $y'=1$.\\
Khi $x < 0$, ta có $y=-x$ và $y'=-1$.\\
Tại $x = 0$, ta có giá trị nhỏ nhất của hàm số là $y = 0$.\\ Vì vậy, ta có cực trị tại điểm $x = 0$, với giá trị cực trị là $y = 0$.\\
Cách 2: Tìm cực trị bằng định nghĩa.\\
Hàm số $y=|x|$ xác định là liên tục trên khoảng $\left(-\infty;+\infty\right)$. \\
Xét điểm $x_0=0 \in \left(-\infty;+\infty\right)$.
Ta xét $h=1$ khi đó
$f(x) > f(0)$ với mọi $x\in \left(0-1;0+1\right) \subset \left(-\infty;+\infty\right)$ và $x\ne 0$.\\
Vậy hàm số $f(x)$ đạt cực đại tại $x_0$.
}
\end{bt}

\begin{bt}%[2D1N1-2]
% nguon: Archive/2024/TaiLieuGiangDay/1.HamSo/1. K12 - 4in1 - Hamso/Thubai/Bai1-Dang2_gitlab_gitlab.tex #2906
%%[Mức 1]%[BG12new-4in1, Trần Hoà]
    Cho hàm số $y=f(x)$ có bảng biến thiên như hình dưới đây
    
        \begin{tikzpicture}
            \tkzTabInit[nocadre=false,lgt=1.2,espcl=3]
            {$x$ /0.7,$y'$ /0.7,$y$ /2}
            {$-\infty$,$-\tfrac{1}{2}$,$3$,$+\infty$}
            \tkzTabLine{,+,d,+,$0$,-,}
            \tkzTabVar{-/$-\infty$,+D-/$+\infty$/$-\infty$,+/$4$,-/$-\infty$}
        
\end{tikzpicture}
    
   Tìm các khoảng đồng biến, nghịch biến của hàm số?
    \loigiai
    {Dựa vào bảng biến thiên, ta thấy hàm số đồng biến trên các khoảng $\left(-\infty;-\dfrac{1}{2}\right)$ và $\left(-\dfrac{1}{2};3\right)$; hàm số nghịch biến trên khoảng $(3;+\infty)$.
    }
\end{bt}

\begin{bt}%[2D1H1-4]
% nguon: Archive/2024/TaiLieuGiangDay/1.HamSo/1. K12 - 4in1 - Hamso/Thubai/Bai1-Dang4_converted_converted.tex #2963
%%[Mức độ 2]%[Dự án bài giảng new 4in1, Trần Quang Thạnh]
Giải phương trình $5^x-5^{2x+1}=x+1$.
\loigiai{
Phương trình đã cho tương đương với $5^x + x =5^{2x+1}+2x+1$.\\
Vì hàm số $y=f(t)=5^t+t$ có $y'=5^t\ln5 +1 > 0$, với mọi $t\in \mathbb{R}$ nên hàm số $y=f(t)$ đồng biến trên $\mathbb{R}$, do đó, từ phương trình đã cho ta được $$f(x)=f(2x+1) \Leftrightarrow x=2x+1 \Leftrightarrow x=-1.$$
Vậy phương trình đã cho có nghiệm duy nhất $x=-1$.
}
\end{bt}

\begin{bt}%[2D1H1-4]
% nguon: Archive/2024/TaiLieuGiangDay/1.HamSo/1. K12 - 4in1 - Hamso/Thubai/Bai1-Dang4_converted_converted.tex #2964
%%[Mức độ 2]%[Dự án bài giảng new 4in1, Trần Quang Thạnh]
Giải bất phương trình $\log x > \log (2x-1) + x-1$.
\loigiai{
Điều kiện $x > \dfrac{1}{2}$.\\
Bất phương trình đã cho tương đương với
$\log x + x > \log(2x-1) +2x-1$.\\
Vì hàm số $y=f(t)=\log t + t$ có $y'=\dfrac{1}{t\ln10}+1 > 0$, với mọi $t > 0$ nên hàm số $y=f(t)$ đồng biến trên $(0;+\infty)$.\\
Do đó, từ bất phương trình đã cho ta được
$$f(x) > f(2x-1) \Leftrightarrow x > 2x-1 \Leftrightarrow x < 1.$$
Kết hợp điều kiện, ta có tập nghiệm bất phương trình là $S=\left(\dfrac{1}{2};1\right)$.
}
\end{bt}

\begin{bt}%[2D1V1-4]
% nguon: Archive/2024/TaiLieuGiangDay/1.HamSo/1. K12 - 4in1 - Hamso/Thubai/Bai1-Dang4_converted_converted.tex #2965
%%[Mức độ 3]%[Dự án bài giảng new 4in1, Trần Quang Thạnh]
Giải bất phương trình $ \log_2(2^x+1)+\log_3(4^x+2)\le 2$.\hfill $ (*) $
\loigiai{
\begin{itemize}
\item Điều kiện $ \begin{cases}2^x+1 > 0\\4^x+2 > 0\end{cases} $ đúng $ \forall x\in \mathbb{R}\Rightarrow $ Tập xác định $ \mathscr{D}=\mathbb{R} $.
\item Xét hàm số $ f(x)=\log_2(2^x+1)+\log_3(4^x+2) $ trên $ \mathbb{R} $.\\
Ta có $ f'(x)=\dfrac{2^x\ln2}{(2^x+1)\ln2}+\dfrac{4^x\ln4}{(4^x+1)\ln3} > 0,\forall x\in \mathbb{R}\Rightarrow $ Hàm số $ y=f(x) $ luôn đồng biến trên $ \mathbb{R} $.\\
Ta có $ f(0)=2 $ do đó $ \forall x\le 0 $ thì $ f(x)\le f(0)\Leftrightarrow x\le0 $.
\item Vậy tập nghiệm của bất phương trình là $ x\in (-\infty;0] $.
\end{itemize}}
\end{bt}

\begin{bt}%[2D1N1-5]
% nguon: Archive/2024/TaiLieuGiangDay/1.HamSo/1. K12 - 4in1 - Hamso/Thubai/Bai1-Dang5.tex #3041
%%[SGK12-CTST, Mức độ 1]%[BG12-4in1, Trần Thiên Đức]
	\immini{Trong $8$ phút kể từ khi xuất phát, độ cao $h$ (tính bằng mét) của khinh khí cầu vào tời điểm $t$ phút được cho bởi công thức $h(t)=6t^3-81t^2+324t$. Đồ thị của hàm số $h(t)$ được biểu diễn trong hình bên. Trong các khoảng thời gian nào khinh khí cầu tăng dần độ cao, giảm dần độ cao? Độ cao của khinh khí cầu vào các thời điểm $3$ phút và $6$ phút sau khi xuất phát có gì đặc biệt?
	}{
		\begin{tikzpicture}[scale=0.7, font=\footnotesize, line join=round, line cap=round, >=stealth]
			\def\xmin{-2} \def\xmax{9}
			\def\ymin{-1} \def\ymax{6} 
			\draw[->] (\xmin,0)--(\xmax,0) node [below]{$t$};
			\draw[->] (0,\ymin)--(0,\ymax) node [left]{$h$};
			\clip (\xmin+0.1,\ymin+0.1) rectangle (\xmax-0.1,\ymax-0.1);
			\draw[smooth,samples=300,domain=0:7.3,black] plot(\x,{0.07*(\x)^3-0.89*(\x)^2+3.39*(\x)});
			\fill (0,0)circle (1pt)node[below left]{$O$};
			\fill (0,3.35)circle(1pt )node[left]{$324$} (0,4.05)circle(1pt )node[left]{$405$} (0,4.6)circle(1pt )node[left]{$480$}
			(3,0)circle(1pt )node[below]{$3$}
			(5.6,0)circle(1pt )node[below]{$6$}
			(7.3,0)circle(1pt )node[below]{$8$}
			;
			\draw[dashed] (3,0)--(3,4.05)--(0,4.05) (5.6,0)--(5.6,3.35)--(0,3.35) (7.3,0)--(7.3,4.6)--(0,4.6);
			\fill (3,4.05)circle (1pt)
			(5.6,3.35)circle (1pt)
			(7.3,4.6)circle (1pt)
			;
		\end{tikzpicture}
	}
	\loigiai{
		Dựa vào đồ thị ta thấy khinh khí cầu tăng dần độ cao trong khoảng thời gian $3$ phút đầu và từ phút thứ $6$ đến phút thứ $8$, khinh khí cầu giảm dần độ cao trong khoảng thời gian từ phút thứ $3$ đến phút thứ $6$.\\
		Độ cao của khinh khí cầu tại thời điểm phút thứ $3$ là lớn nhất trong khoảng thời gian $6$ phút đầu, độ cao của khinh khí cầu tại thời điểm phút thứ $6$  là nhỏ nhất trong khoảng thời gian từ phút thứ $3$ đến phút thứ $8$.
		
	}
\end{bt}

\begin{bt}%[2D1H1-5]
% nguon: Archive/2024/TaiLieuGiangDay/1.HamSo/1. K12 - 4in1 - Hamso/Thubai/Bai1-Dang5.tex #3042
%%[SGK12 - Cánh Diều, Mức độ 2]
	\immini{
		Máng trượt của một cầu trượt cho trẻ em (\textit{Hình a}) được uốn từ một tấm kim loại có bề rộng $80$ cm, mặt cắt được mô tả ở \textit{Hình b}. Nhà thiết kế khuyến cáo, diện tích mặt cắt càng lớn thì càng đảm bảo an toàn cho trẻ em.
		\begin{enumerate}[a)]
			\item Gọi $S$ là diện tích mặt cắt. Tìm điều kiện của $x$ và viết công thức tính $S$ theo $x$.
			\item Diện tích mặt cắt sẽ tăng hay giảm khi chúng ta tăng giá trị của $x$ từ 0 cm đến 10 cm?
		\end{enumerate}
	}{
		\begin{tikzpicture}[scale=1, font=\footnotesize, line join=round, line cap=round, >=stealth,line width=3.pt]
			\clip(-0.2,-1.2) rectangle (7.5,2.6);
			\fill[line width=3.pt,color=black,fill=black,fill opacity=0.7] (0,0)--(0.41,0.34)--(1,1)--(1.76,1.76)--(2,2)--(2.46,2.08)--(2.83,1.95)--(2.34,1.65)--(1.34,0.65)--(0.75,-0.01)--(0.36,-0.13)-- cycle;
			\draw (0,0)--(0.41,0.34)--(1,1)--(2,2)--(2.46,2.08);
			\draw (0,0)--(0.36,-0.13)--(0.77,0.21)--(0.77,0.21)--(1.36,0.87)--(2.36,1.87)--(2.83,1.95);
			\draw (0.36,-0.13)--(0.75,-0.01)--(1.34,0.65)--(2.34,1.65);
			\draw (2.36,0.68)-- (3.95,0.28);
			\draw (3.42,0.41)--(2.68,2.29)--(2.53,2.37)--(2.16,1.88)--(2.31,1.53);
			\draw (1.63,2.01)--(2,2.5)-- (2.16,2.42)--(2.89,0.55);
			\draw (1.63,2.01)--(1.76,1.76) (2.73,0.98)--(3.26,0.85) (2.52,1.49)--(3.05,1.35);
			\draw (2.34,1.65)--(2.83,1.95)--(2.46,2.08);
			\draw (2,-0.8)node{a)};
			\fill[line width=3.pt,fill=gray,fill opacity=0.3] (5,0.5)--(5,0)--(6.5,0)--(6.5,0.5)--cycle;
			\draw (5,0.5)--(5,0) node [left ,sloped,pos=0.5,rotate=90] {$x$ cm};
			\draw (5,0)--(6.5,0) node [below,sloped,pos=0.5] {$y$ cm};
			\draw (6.5,0)--(6.5,0.5) node [right,sloped,pos=0.5,rotate=-90] {$x$ cm};
			\draw[dashed,thin] (5,0.5)--(6.5,0.5);
			\draw (5.8,-0.8)node{b)};
			%			\draw (4,-1)node{\textit{Hình 5}};
		\end{tikzpicture}
	}
	\loigiai{
		\begin{enumerate}[a)]
			\item Do tấm kim loại có bề rộng $80 \mathrm{~cm}$ nên ta có $2x+y=80 \Leftrightarrow y=80-2x$.\\
			Để có thể thiết kế được máng trượt thì $y>0 \Leftrightarrow 80-2 x>0 \Leftrightarrow x<40$. Suy ra $0<x<40$.\\
			Diện tích của mặt cắt máng trượt là: $S=x y=x(80-2x)=-2x^2+80x$.		
			\item Ta có $S(x)=-2x^2+80x \text{ với } x\in(0; 40)$.
			\allowdisplaybreaks
			\begin{eqnarray*}
				S'(x)&=&-4 x+80 \\
				S'(x)&=&0 \Leftrightarrow-4 x+80=0 \Leftrightarrow x=20.
			\end{eqnarray*}
			Bảng biến thiên của hàm số $S(x)$ như sau
			\begin{center}
				\begin{tikzpicture}[font=\footnotesize,thick,>=stealth]
					\tikzset{double style/.append style={double distance=1.5pt}}
					\tkzTabInit[nocadre=false,lgt=1.2,espcl=2.5,deltacl=0.6,lw=.75pt]
					{$x$ /0.7, $S'(x)$ /0.8, $S(x)$ /2}
					{$0$,$20$,$40$}
					\tkzTabLine{ ,+,,-, }
					\tkzTabVar{-/$0$,+/$800$,-/$0$}
				\end{tikzpicture}
			\end{center}
			Vậy từ khi $x$ tăng từ $0$ cm đến $20$ cm thì diện tích mặt cắt sẽ tăng dần.
		\end{enumerate}
	}
\end{bt}

\begin{bt}%[2D1H1-5]
% nguon: Archive/2024/TaiLieuGiangDay/1.HamSo/1. K12 - 4in1 - Hamso/Thubai/Bai1-Dang5.tex #3044
%%[SGK12-CTST, Mức độ 2]
	Kim ngạch xuất khẩu rau quả của Việt Nam trong các năm từ $2010$ đến $2017$ có thể được tính xấp xỉ bằng công thức $f(x)=0{,}01x^3-0{,}04x^2+0{,}25x+0{,}44$ (tỉ USD) với $x$ là số năm tính từ $2010$ đến $2017$ $(0\leq x\leq 7)$.
	\begin{flushright}
		(\textit{Theo:} https://infographics.vn/interactive-xuat-khau-rau-qua-
		du-bao-bung-no-dat-4-ty-usd-trong-nam-2023/116220.vna)
	\end{flushright}
	\begin{enumerate}
		\item Tính đạo hàm của hàm số $y=f(x)$.
		\item Chứng minh rằng kim ngạch xuất khẩu rau quả của Việt Nam tăng liên tục trong các năm từ $2010$ đến $2017$.
	\end{enumerate}
	\loigiai{
		\begin{enumerate}
			\item Ta có $f'(x)=0{,}03x^2-0{,}08x+0{,}25$.
			\item Xét $f'(x)=0 \Leftrightarrow 0{,}03x^2-0{,}08x+0{,}25=0$ (vô nghiệm).\\
			Bảng biến thiên
			\begin{center}
				\begin{tikzpicture}
					\tkzTabInit[nocadre=false,lgt=1.2,espcl=5,deltacl=0.6]
					{$x$ /0.6,$f'(x)$ /0.6,$f(x)$ /2}
					{$0$,$7$}
					\tkzTabLine{,+,}
					\tkzTabVar{-/$0{,}44$,+/$3{,}66$}
				\end{tikzpicture}
			\end{center}
			Từ bảng biến thiên trên, ta thấy $f'(x)>0$, $\forall x\in [0;7]$.\\
			Vậy kim ngạch xuất khẩu rau quả của Việt Nam tăng liên tục trong các năm từ $2010$ đến $2017$.
		\end{enumerate}
	}
\end{bt}

\begin{bt}%[2D1H1-5]
% nguon: Archive/2024/TaiLieuGiangDay/1.HamSo/1. K12 - 4in1 - Hamso/Thubai/Bai1-Dang5.tex #3048
%%[Mức độ 2]%
	Định luật Moore được xây dựng bởi Gordon More-một trong những sáng lập viên của tập đoàn sản xuất chip máy tính nổi tiếng Intel. Định luật ban đầu được phát biểu như sau
	\lq\lq Số lượng transistor trên mỗi đơn vị inch vuông sẽ tăng lên gấp đôi sau mỗi $24$ tháng. \rq\rq ($1$ inch vuông xấp xỉ $6{,}45$ cm$^2$).
	Phương trình sau đây mô tả định luật Moore
	$$n_i=n_0\cdot 2^{\tfrac{y_i-y_0}{2}}$$
	trong đó $n_0$ là số lượng transitor trên một đơn vị diện tích ban đầu, $y_0$ là năm ứng với $n_0$, $n_i$ là số lượng transitor tại năm $y_i$ đang xét. \\
	(Nguồn: https://scipython.com/book/chapter-3-simple-plotting-with-pylab/examples/moores-law/)\\
	Nhận xét về tính đơn điệu của hàm số mô tả định luật Moore.
	\loigiai{
		Xét hàm số $n_i=n_0\cdot 2^{\tfrac{y_i-y_0}{2}}$.\\
		Ta có $n_i'=n_0\cdot 2^{\tfrac{y_i-y_0}{2}}\cdot \ln 2>0~\forall y_i$.\\
		Suy ra hàm số đồng biến.
	}
\end{bt}

\begin{bt}%[2D1H1-5]
% nguon: Archive/2024/TaiLieuGiangDay/1.HamSo/1. K12 - 4in1 - Hamso/Thubai/Bai1-Dang5.tex #3049
%%[SGK12-KNTT, Mức độ 2]%[BG12-4in1, Trần Thiên Đức]
	Xét một chất điểm chuyển động trên một trục số nằm ngang, chiều dương từ trái sang phải (Hình vẽ). Giả sử vị trí $s(t)$ (mét) của chất điểm trên trục số đã chọn tại thời điểm $t$ (giây) được cho bởi công thức
	\[s(t)=t^3-9t^2+15t,\,\, t\ge 0.\]
	Hỏi trong khoảng thời gian nào thì chất điểm chuyển động sang phải, trong khoảng thời gian nào  thì chất điểm chuyển động sang trái?
	\begin{center}
		\begin{tikzpicture}[>=stealth,line join=round,line cap=round,font=\footnotesize,scale=1]
			\draw[->] (0,0)coordinate (O)--(5,0)coordinate (B);
			\draw[thick] 
			(1.5,-.05)--(1.5,.05)node[below=1mm,black]{$O$}
			(3,-.05)--(3,.05)node[below=1mm,black]{$s(t)$}
			;
			\path 
			(O)--(B)node[midway,below=6mm,orange]{}
			;
		\end{tikzpicture}
	\end{center}
	\loigiai{
		Theo ý nghĩa cơ học của đạo hàm, vận tốc $v\left(t\right)$ là đạo hàm của $s\left(t\right)$. \\
		Ta có $v\left(t\right) = s'\left(t\right) = \left(t^3 -9t^2 + 15t\right)' = 3t^2 - 18t + 15$.\\
		Tập xác định $\mathscr{D} = \mathbb{R}$.\\
		Ta có $v(t)=0\Leftrightarrow\hoac{&t=1 \\&t=5. }$\\
		Bảng xét dấu
		\begin{center}
			\begin{tikzpicture}[scale=1, font=\footnotesize]
				\tkzTabInit[nocadre=false, lgt=1.2, espcl=2, deltacl=0.6]
				{$t$/0.6,$v(t)$/0.6}
				{$-\infty$,$1$,$5$,$+\infty$};
				\tkzTabLine{,+,$0$,-,$0$,+,};
			\end{tikzpicture}
		\end{center}
		Chất điểm chuyển động theo chiều dương (sang bên phải) khi $v\left(t\right) > 0$, tức là $t \in \left(-\infty; 1\right) \cup \left(5; +\infty\right)$.\\
		Chất điểm chuyển động theo chiều âm (sang bên trái) khi $v\left(t\right) < 0$, tức là $1 < t < 5$.
	}	
\end{bt}

\begin{bt}%[2D1H1-3]
% nguon: Archive/2024/TaiLieuGiangDay/1.HamSo/1. K12 - 4in1 - Hamso/Thubai/Bai1-Dang3 copy.tex #3068
%%[Mức độ 2]%[Dự án bài giảng new 4in1, Trần Quang Thạnh]
Tìm các giá trị của tham số $m$ để hàm số $y=\dfrac{a}{x+1}$ nghịch biến trên $(1;+\infty)$.
\loigiai{
Ta có $y'=-\dfrac{a}{(x+1)^2}$.\\
Nhận thấy rằng, $y'=0$ trên $(1;+\infty)$ khi và chỉ khi $a=0$.\\
Do đó hàm số nghịch biến trên $(1;+\infty)$ khi $y' < 0$, với mọi $x > 1$, tức là $a > 0$.
}
\end{bt}

\begin{bt}%[2D1H1-3]
% nguon: Archive/2024/TaiLieuGiangDay/1.HamSo/1. K12 - 4in1 - Hamso/Thubai/Bai1-Dang3 copy.tex #3069
%%[Mức độ 2]%[Dự án bài giảng new 4in1, Trần Quang Thạnh]
Tìm các giá trị của tham số $a$ để hàm số $y=-x^3+ax$ nghịch biến trên $\mathbb{R}$.
\loigiai{
Ta có $y'=-3x^2+a$.\\
Hàm số nghịch biến trên $\mathbb{R}$ khi và chỉ khi $y'\leq 0$, với mọi $x\in \mathbb{R}$, điều này tương đương với $\Delta'_{y'}=3a \leq 0$ hay $a\leq 0$.\\
Vậy với $a\leq 0$, hàm số đã cho nghịch biến trên $\mathbb{R}$.
}
\end{bt}

\begin{bt}%[2D1H1-3]
% nguon: Archive/2024/TaiLieuGiangDay/1.HamSo/1. K12 - 4in1 - Hamso/Thubai/Bai1-Dang3 copy.tex #3070
%%[Mức độ 2]%[Dự án bài giảng new 4in1, Trần Quang Thạnh]
Tìm tất cả các giá trị của tham số $m$ để hàm số $y=\dfrac{x^5}{5}+\dfrac{mx^4}{2}+\dfrac{x^3}{3}+m-1$ đồng biến trên $\mathbb{R}$.
\loigiai{
Ta có $y'=x^4+2mx^3+x^2=x^2(x^2+2mx+1)$.\\
Hàm số đồng biến trên $\mathbb{R}$ khi và chỉ khi $y'\geq 0, \forall x\in \mathbb{R}$, hay $x^2+2mx+1\geq 0, \forall x\in \mathbb{R}$; điều này tương đương với $m^2-1 \leq 0$, tức là $-1\leq m\leq 1$.\\
Vậy hàm số đã cho đồng biến trên $\mathbb{R}$ khi và chỉ khi $-1\leq m\leq 1$.
}
\end{bt}

\begin{bt}%[2D1H1-3]
% nguon: Archive/2024/TaiLieuGiangDay/1.HamSo/1. K12 - 4in1 - Hamso/Thubai/Bai1-Dang3 copy.tex #3071
%%[Mức độ 2]%[Dự án bài giảng new 4in1, Trần Quang Thạnh]
Tìm các giá trị của tham số $m$ để hàm số $y=\dfrac{mx+m-7}{5x-m+3}$ đồng biến trên mọi khoảng của tập xác định.
\loigiai{
Tập xác định $\mathscr{D}=\mathbb{R}\setminus \left\{\dfrac{m-3}{5}\right\}$.\\ Ta có $y'=\dfrac{-m^2-2m+35}{(5x-m+3)^2}$.\\
Hàm số đồng biến trên mọi khoảng xác định khi và chỉ khi\\
\centerline{$y' > 0,\forall x\ne \dfrac{m-3}{5} \Leftrightarrow -m^2-2m+35 > 0 \Leftrightarrow m\in (-7;5)$.}
Vậy, với $m\in (-7;5)$ thì hàm số đồng biến trên mọi khoảng xác định của nó.
}
\end{bt}

\begin{bt}%[2D1V1-3]
% nguon: Archive/2024/TaiLieuGiangDay/1.HamSo/1. K12 - 4in1 - Hamso/Thubai/Bai1-Dang3 copy.tex #3072
%%[Mức độ 3]%[Dự án bài giảng new 4in1, Trần Quang Thạnh]
Tìm $a$ để hàm số $\displaystyle y = \frac{{{x^2} - 2ax + 4{a^2}}}{{x - 2a}}$ đồng biến trên $\left( {2; + \infty } \right)$.
\loigiai{
Tập xác định $\mathscr{D} = \mathbb{R}\backslash \left\{ 2a \right\}$ và $\displaystyle y' = \frac{{{x^2} - 4ax}}{{{{\left( {x - 2a} \right)}^2}}}$.
Hàm số đồng biến trên $\left( {2; + \infty } \right) \Leftrightarrow y' \ge 0,\forall x \in \left( {2; + \infty } \right)$.
Hay $\begin{cases}2a\notin (2;+\infty)\\x^2-4ax\geqslant 0,\forall x \in (2;+\infty)\end{cases}\Leftrightarrow \begin{cases}2a\leq 2\\a\leq \dfrac{x}{4},\forall x\in (2;+\infty)\end{cases}\Leftrightarrow \begin{cases}
a\leq 1\\a\leq \dfrac{1}{2}\end{cases}\Leftrightarrow a\leq \dfrac{1}{2}$.
Vậy với $a\leq \dfrac{1}{2}$, hàm số đồng biến trên $\left( {2; + \infty } \right)$.
}
\end{bt}

\begin{bt}%[2D1N1-2]
% nguon: Archive/2024/TaiLieuGiangDay/1.HamSo/1. K12 - 4in1 - Hamso/Thubai/Bai1-Dang2.tex #3102
%%[Mức 1]%[BG12new-4in1, Trần Hoà]
Cho hàm số $y=f(x)$ có bảng biến thiên như sau.
\begin{center}
\begin{tikzpicture}[>=stealth]
		\tkzTabInit[nocadre=false,lgt=1,espcl=2,deltacl=0.5]{$x$/.7 ,$y'$/.7,$y$/2}
		{$-\infty$ , $-3$ , $4$ , $+\infty$}
		\tkzTabLine{ , - , $0$ , + , $0$ , - , }
		\tkzTabVar{+/$+\infty$ , -/$-1$ , +/$2$ , -/$-\infty$}
\end{tikzpicture}
\end{center}
Tìm các khoảng đồng biến, nghịch biến của hàm số.
	\loigiai{Dựa vào bảng biến thiên ta có:\\
	Hàm số đồng biến trên khoảng $(-3;4)$.\\
	Hàm số nghịch biến trên các khoảng $(-\infty;-3)$ và $(4;+\infty)$.}	
\end{bt}

\begin{bt}%[2D1H1-2]
% nguon: Archive/2024/TaiLieuGiangDay/1.HamSo/1. K12 - 4in1 - Hamso/Thubai/Bai1-Dang2.tex #3103
%%[Mức 2]%[BG12new-4in1, Trần Hoà]
Cho hàm số $y=f(x)$ có bảng biến thiên như hình dưới đây
\begin{center}
\begin{tikzpicture}[scale=1, font=\footnotesize, line join=round, line cap=round, >=stealth]
\tkzTabInit[nocadre=false,lgt=1.2,espcl=2.5,deltacl=0.6] %phần bắt buộc
{$x$ /0.6,$y’$ /0.6,$y$ /2}
{$-\infty$,$0$,$2$,$4$,$+\infty$}
\tkzTabLine{ ,+,$0$,-,d,-,$0$,+, }
\tkzTabVar{-/$-\infty$,+/$-2$,-D+/$-\infty$/$+\infty$,-/$6$,+/$+\infty$}
\end{tikzpicture}
\end{center}
Tìm các khoảng đồng biến, nghịch biến của hàm số.
\loigiai{Dựa vào bảng biến thiên ta có:\\
	Hàm số đồng biến trên các khoảng $(-\infty;0)$ và $(4;+\infty)$.\\
	Hàm số nghịch biến trên các khoảng $(0;2)$ và $(2;+4)$.}	
\end{bt}

\begin{bt}%[2D1H1-2]
% nguon: Archive/2024/TaiLieuGiangDay/1.HamSo/1. K12 - 4in1 - Hamso/Thubai/Bai1-Dang2.tex #3105
%%[Mức 2]%[BG12new-4in1, Trần Hoà]
	\immini
	{
		Đường cong trong hình bên là đồ thị của hàm số $y=\dfrac{ax+b}{cx+d}$, với $a$, $b$, $c$, $d$ là các số thực. Tìm các khoảng đồng biến, nghịch biến của hàm số.	
	}
	{
		\begin{tikzpicture}[scale=0.8,>=stealth, font=\footnotesize, line join=round, line cap=round]
			\def\a{2} \def\b{1} \def\c{2} \def\d{-1} % Hệ số
			\def\xmin{-3} \def\xmax{4}
			\def\ymin{-2.3} \def\ymax{4}
			\draw[->] (\xmin,0)--(\xmax,0) node [below]{$x$};
			\draw[->] (0,\ymin)--(0,\ymax) node [left]{$y$};
			\node at (0,0) [above left]{$O$};
			\clip (\xmin+0.1,\ymin+0.1) rectangle (\xmax-0.1,\ymax-0.1);
			\draw[smooth,samples=300,domain=\xmin:(-\d/\c-0.1)] plot(\x,{(\a*(\x)+\b)/(\c*(\x)+\d)});
			\draw[smooth,samples=300,domain=(-\d/\c+0.1:\xmax)] plot(\x,{(\a*(\x)+\b)/(\c*(\x)+\d)});
			\draw[dashed] (-\d/\c,\ymin)--(-\d/\c,\ymax);
			\draw[dashed] (\xmin,\a/\c)--(\xmax,\a/\c);
			\node at (0.7,-0.2){2};
			\node at (-0.2,1.2){3};
		\end{tikzpicture}
	}
	\loigiai{
		Nhận thấy hàm số nghịch biến trên từng khoảng $(-\infty; 2)$ và $(2;+\infty)$.
	}
	
\end{bt}

\begin{bt}%[2D1V1-2]
% nguon: Archive/2024/TaiLieuGiangDay/1.HamSo/1. K12 - 4in1 - Hamso/Thubai/Bai1-Dang2.tex #3106
%%[Kết nối, mức 2]%[TeX hóa SGK CTST 12]%[Nguyễn Tiến]%[BG12new-4in1, Trần Hoà]%[BG12new-4in1, Trần Hoà]
	\immini{
		Đạo hàm $f'(x)$ của hàm số $y=f(x)$ trên khoảng $(-2;5)$ có đồ thị như hình vẽ bên. Xét tính đơn điệu của hàm số $y=f(x)$ trên khoảng $(-2;5)$.
	}{
		\begin{tikzpicture}[scale=0.7, font=\footnotesize, line join=round, line cap=round, >=stealth]
			\def\a{0.125} \def\b{-0.625} \def\c{0.25} \def\d{1} % Hệ số
			\def\xt{-3} \def\xp{6} \def\yt{2.5} \def\yd{-3} % x_trái, x_phải, y_trên, y_dưới (giới hạn)
			\draw[->] (\xt,0)--(\xp,0) node [below]{$x$};
			\draw[->] (0,\yd)--(0,\yt) node [left]{$y$};
			\node at (0,0) [below right]{$O$};
			\clip (\xt+0.1,\yd-1) rectangle (\xp-0.1,\yt-0.1);
			\draw[smooth,samples=300,domain=-2:5] plot(\x,{\a*(\x)^3+\b*(\x)^2+\c*(\x)+\d}) node [below left]{$y=f'(x)$};
			\draw[dashed] (-2,0) node [below left]{$-2$} -- (-2,-3);
			\draw[dashed] (5,0) node [below]{$5$} -- (5,2.25);
			\fill (-1,0) circle (1.5pt) node[below right]{$-1$};
			\fill (2,0) circle (1.5pt) node[below left]{$2$};
			\fill (4,0) circle (1.5pt) node[below right]{$4$};
			%\node at (1.5,-3.5){\textit{Hình 12}};
		\end{tikzpicture}
	}
	\loigiai{
		Từ đồ thị ta thấy $f'(x)=0 \Leftrightarrow \hoac{& x=-1\\& x=2\\& x=4.}$\\
		Ta có bảng biến thiên như sau
		\begin{center}
			\begin{tikzpicture}
		\tkzTabInit[nocadre=false,lgt=1.2,espcl=2.5,deltacl=0.6]
				{$x$ /0.6,$f'(x)$ /0.6,$f(x)$ /2}
				{$-2$,$-1$,$2$,$4$,$5$}
				\tkzTabLine{,-,$0$,+,$0$,-,$0$,+,}
				\tkzTabVar{+/,-/,+/,-/,+/}
			\end{tikzpicture}
		\end{center}
		Vậy hàm số đồng biến trên $(-1;2)$ và $(4;5)$, nghịch biến trên $(-2;-1)$ và $(2;4)$.
	}
\end{bt}

\begin{bt}%[2D1H2-6]
% nguon: Archive/2024/TaiLieuGiangDay/5.HH Oxyz/VDC_TDM.tex #3162 
% KIEM_TRA: cần xem lại đáp số/dạng

\immini{
Trong không gian $Oxyz$, cho một điểm sáng $S(-4;2;5)$ và một miếng bìa chắn sáng $(C)$ (không cho ánh sáng đi qua), biết rằng $(C)$ nằm trên mặt phẳng $(\alpha) \equiv (Oxy): z=0$, có tâm $O$ và bán kính $r=2$. Mắt quan sát đặt tại điểm $A(3;0;-2)$. Hãy
\begin{enumerate}
\item Một điểm $M(x;y;z) \in (C)$ thì thỏa mãn điều kiện gì ?
\item Viết phương trình đường thẳng $SA$ ?
\item Hỏi mắt có quan sát được điểm $S$ không ?
\item Hỏi mắt đặt tại $B(6;4;-1)$ có quan sát được điểm $S$ không ?
\end{enumerate}
}{
\begin{tikzpicture}[scale=1, font=\footnotesize, line cap=round, line join=round, >=stealth]
\begin{axis}[axis lines=middle,xlabel={$x$},ylabel={$y$}, zlabel={$z$}, grid=major, xmin=-4.5, xmax=4, ymin=-2, ymax=4, zmin=-3, zmax=5.5, enlargelimits=true, view={120}{30}]
% Vẽ ellipse tại gốc tọa độ
\draw[fill=orange, opacity=0.5] (axis cs:0,0,0) ellipse (2 and 2);
% Vẽ các điểm
\addplot3[only marks, mark=*] coordinates {(0,0,0)} node[anchor=south west] {$O$};
\addplot3[only marks, mark=*] coordinates {(3,0,-2)} node[anchor=north east] {$A$};
\addplot3[only marks, mark=*] coordinates {(-4,2,5)} node[anchor=south] {$S$};
\addplot3 coordinates {(-1,-2.2,0)} node[anchor=south] {$(C)$};
%\addplot3[only marks, mark=*] coordinates {(6,4,-1)} node[anchor=north] {$B$};
%\addplot3[solid, thick, blue] coordinates {(3,0,-2) (-4,2,5)};
%\addplot3[solid, thick, blue] coordinates {(6,4,-1) (-4,2,5)};
\end{axis}
\end{tikzpicture}
}
\loigiai{
\begin{enumerate}
\item Vì $M \in (C)$ và $(C)$ nằm trên mặt phẳng $z=0$, có tâm $O(0;0;0)$ và bán kính $r=2$, nên tọa độ của $M$ phải thỏa mãn $\heva{&x^2 + y^2 \leq 2^2 \\& z = 0.}$
\item Đường thẳng $SA$ đi qua hai điểm $S(-4;2;5)$ và $A(3;0;-2)$ có vectơ chỉ phương $\overrightarrow{AS} = (-7; 2; 7)$.\\
Phương trình tham số của đường thẳng $SA$ là $\heva{&x = 3 - 7t \\&y = 2t \\&z = -2 + 7t.}$
\item Để mắt tại $A$ không quan sát được điểm $S$, đường thẳng $SA$ phải cắt miếng bìa $(C)$. Tức là tồn tại điểm $M$ trên đoạn $SA$ thuộc $(C)$.\\
Giả sử điểm $M$ thuộc đường thẳng $SA$ có tọa độ $M(3-7t; 2t; -2+7t)$. \\
Nếu $M \in (C)$, thì $z=0 \Rightarrow -2+7t = 0 \Rightarrow t = \dfrac{2}{7}$.\\
Với $t = \dfrac{2}{7}$, tọa độ của $M$ là $M\left( 1; \dfrac{4}{7}; 0 \right)$.\\
Ta có $1^2 + \left(\dfrac{4}{7}\right)^2 = 1 + \dfrac{16}{49} = \dfrac{49+16}{49} = \dfrac{65}{49} < 4$.\\
Vì điểm $M\left(1; \dfrac{4}{7}; 0\right)$ thuộc $(C)$ và nằm trên đường thẳng $SA$, nên mắt tại $A$ không quan sát được điểm $S$.
\item Tương tự, xét đường thẳng $SB$ đi qua $S(-4;2;5)$ và $B(6;4;-1)$. \\
Vectơ chỉ phương $\overrightarrow{BS} = (-10; -2; 6)$.\\
Phương trình tham số của đường thẳng $SB$ là $\heva{&x = 6 - 10u \\&y = 4 - 2u \\&z = -1 + 6u.}$\\
Để đường thẳng $SB$ cắt $(C)$, cần có điểm $N$ trên $SB$ có $z=0$.\\
Suy ra $-1 + 6u = 0 \Rightarrow u = \dfrac{1}{6}$. Với $u = \dfrac{1}{6}$, tọa độ của $N$ là $N\left( \dfrac{13}{3}; \dfrac{11}{3}; 0 \right)$.\\
Kiểm tra điều kiện $x^2 + y^2 \leq 4$, ta có $\left(\dfrac{13}{3}\right)^2 + \left(\dfrac{11}{3}\right)^2 = \dfrac{169}{9} + \dfrac{121}{9} = \dfrac{290}{9} > 4$.\\
Vì điểm $N$ không thuộc $(C)$, nên đường thẳng $SB$ không bị chắn bởi $(C)$. Vậy mắt đặt tại $B$ có thể quan sát được điểm $S$.
\end{enumerate}
}
\end{bt}

\begin{bt}%[2D1V3-4]
% nguon: Archive/2024/TaiLieuGiangDay/5.HH Oxyz/VDC_TDM_Dot2.tex #3171
%[Thành Đức Trung]
Trên một mảnh vườn hình chữ nhật $ABCD$ có kích thước $12\text{\;m} \times 10\text{\;m}$, bác Bình đặt một chiếc đèn pin tại vị trí $A$ chiếu chùm sáng phân kì sang phía góc $C$. Biết rằng góc chiếu của chiếc đèn pin bằng $30^\circ$ và không đổi trong quá trình chiếu sáng. Bác Bình quay chiếc đèn pin tại $A$ để điều chỉnh cho hai điểm $M$ và $N$ chạy trên hai cạnh $CD$ và $CB$. Hãy xác định diện tích được chiếu sáng lớn nhất và diện tích được chiếu sáng nhỏ nhất.
\begin{center}
\begin{tikzpicture}[scale=0.5]
\coordinate (A) at (0,0);
\coordinate (B) at (12,0);
\coordinate (C) at (12,10);
\coordinate (D) at (0,10);
\coordinate (M) at (8,10);  
\coordinate (N) at (12,6); 
\draw[thick] (A) -- (B) -- (C) -- (D) -- cycle;
\filldraw[fill=yellow, fill opacity=0.7, draw=black] (A) -- (M) -- (C) -- (N) -- cycle;
\foreach \point/\name in {A/A, B/B, C/C, D/D, M/M, N/N}
\filldraw (\point) circle (2pt) node[above left] {$\name$};
\draw[<->] (13,0) -- (13,10) node[midway, right] {$10\text{\;m}$};
\draw[<->] (0,-1) -- (12,-1) node[midway, below] {$12\text{\;m}$};
\node[below left] at (A) {Đèn pin};
\end{tikzpicture}
\end{center}
\loigiai{
Đặt $\widehat{BAN} = x$ (đơn vị: độ) thì $\widehat{DAM} = 90^\circ - 30^\circ - x = 60^\circ - x$.\\
Vì $\heva{&x\le \widehat{BAC}\approx 39{,}8^\circ\\&60^\circ - x\le \widehat{DAC} \approx 50{,}2^\circ}$ nên $9{,}8^\circ \le x\le 39{,}8^\circ$.
\begin{itemize}
\item Xét $\triangle ABN$ có $BN = AB\cdot \tan x = 12\tan x$, do đó $S_{\triangle ABN} = \dfrac{1}{2}\cdot AB\cdot BN = 72\tan x$.
\item Xét $\triangle DAM$ có $DM = AD\cdot \tan \left(60^\circ - x\right) = 10\tan \left(60^\circ - x\right)$. \\
Do đó $S_{\triangle DAM} = \dfrac{1}{2}\cdot AD\cdot DM = 50\tan \left(60^\circ - x\right)$.
\end{itemize}
Khi đó vùng diện tích được chiếu sáng là
\[S_{AMCN} = S_{ABCD} - S_{\triangle ABN} - S_{\triangle DAM} = 120 - 72\tan x - 50\tan\left(60^\circ - x\right).\]
Đặt $f(x) = 120 - 72\tan x - 50\tan\left(60^\circ - x\right)$, với $x\in\left(9{,}8;39{,}8\right)$. \\
Sử dụng chức năng Table của CASIO, ta được
\begin{itemize}
\item Giá trị lớn nhất của $f(x)$ xấp xỉ $51{,}87$ khi $x\approx21^\circ$.
\item Giá trị nhỏ nhất của $f(x)$ xấp xỉ $41{,}62$ khi $x\approx39{,}8^\circ$.
\end{itemize}
Vậy diện tích được chiếu sáng lớn nhất và diện tích được chiếu sáng nhỏ nhất lần lượt là $51{,}87$ m$^2$ và $41{,}62$ m$^2$.
}
\end{bt}

\begin{bt}%[2D1V4-1]
% nguon: Archive/2024/San pham - To Latex Strong - Tex hoa Lop 12 moi - Dot 3 (Chuong 1 + 2 +3 + 4)/Data/Chuong 1/Canh dieu/To 3 - Dot 16 - Chuong 1 - Bai 3 - CD - De 2.tex #3579
%[Tổ 3 - Đợt 16 - Chương 1 - Bài 3 - Cánh Diều - Đề 2]%[Đình Nguyên]
	Cho hàm số $y=f(x)$ có bảng biến thiên như sau\\
	Biết đồ thị hàm số $g(x)=f\left(\sqrt{x^2+2 x}-x\right)$ có hai đường tiệm cận ngang là $y=a$ và $y=b$, trong đó $a<b$. Tính $S=a-100 b$.
	
	\begin{center}
		\begin{tikzpicture}[scale=1, font=\footnotesize, line join=round, line cap=round, >=stealth]
			\tkzTabInit[nocadre=false,lgt=1,espcl=2,deltacl=0.5]
			{$x$/.7 ,$f'(x)$/.7,$f(x)$/2}
			{$-\infty$ , $0$ , $1$ , $+\infty$}
			\tkzTabLine{ , - , $0$ , + , $0$ , - , }
			\tkzTabVar{+/$4$ , -/$-1$ , +/$3$ , -/$2$}
		\end{tikzpicture}
	\end{center}
	\loigiai{
		$\lim\limits_{x  \to +\infty}\left(\sqrt{x^2+2 x}-x\right)=1$ và $f(1)=3$, suy ra $\lim\limits_{x  \to +\infty} g(x)=\lim\limits_{x  \to +\infty} f\left(\sqrt{x^2+2 x}-x\right)=3$. $\lim\limits_{x  \to -\infty}\left(\sqrt{x^2+2 x}-x\right)=+\infty$ và $\lim\limits_{x  \to +\infty} f(x)=2$, suy ra $\lim\limits_{x  \to -\infty} g(x)=\lim\limits_{x  \to -\infty} f\left(\sqrt{x^2+2 x}-x\right)=2$.\\
		Vậy đồ thị hàm số $g(x)$ có $2$ đường tiệm cận ngang $y=2$ và $y=3$. 
		Suy ra $a=2$, $b=3$.\\
		Suy ra $S=a-100 b=-298$.
	}
\end{bt}

\begin{bt}%[2D1V3-6]
% nguon: Archive/2024/San pham - To Latex Strong - Tex hoa Lop 12 moi - Dot 3 (Chuong 1 + 2 +3 + 4)/Data/Chuong 1/Canh dieu/To 3 - Dot 16 - Chuong 1 - Bai 3 - CD - De 2.tex #3580
%[Tổ 3 - Đợt 16 - Chương 1 - Bài 3 - Cánh Diều - Đề 2]%[Đình Nguyên]
	Một công ty sản xuất đồ gia dụng ước tính chi phí để sản xuất $x$ (sản phẩm) là $C(x)=150 x+900$ (nghìn đồng). Khi sản xuất càng nhiều sản phẩm thì chi phí sản xuất trung bình cho mỗi sản phẩm không vượt quá $t$ (nghìn đồng). Tìm giá trị nhỏ nhất của $t$.
	
	\begin{center}
		\begin{tikzpicture}[scale=1, font=\footnotesize, line join=round, line cap=round, >=stealth]
			\tkzTabInit[nocadre=false,lgt=1,espcl=2,deltacl=0.5]
			{$x$/.7 ,$f'(x)$/.7,$f(x)$/2}
			{$0$ , $+\infty$}
			\tkzTabLine{ , - , }
			\tkzTabVar{+/$+\infty$ , -/$150$}
		\end{tikzpicture}
	\end{center}
	
	\loigiai{
		Chi phí sản xuất trung bình cho mỗi sản phẩm là $f(x)=\dfrac{C(x)}{x}=\dfrac{150 x+900}{x}$.\\
		Suy ra $f'(x)=\dfrac{-900}{x^2}<0 \,\,\forall x>0$. Suy ra
		$\lim\limits_{x\to+\infty}f(x)=\lim\limits_{x  \to+\infty}\dfrac{150x+900}{x}=150$.\\
		Vậy khi sản xuất càng nhiều sản phẩm thì chi phí sản xuất trung bình cho mỗi sản phẩm càng giảm nhưng không dưới $150$ nghìn đồng.
	}
\end{bt}

\begin{bt}%[2D1C4-3]
% nguon: Archive/2024/San pham - To Latex Strong - Tex hoa Lop 12 moi - Dot 3 (Chuong 1 + 2 +3 + 4)/Data/Chuong 1/Canh dieu/To 3 - Dot 16 - Chuong 1 - Bai 3 - CD - De 2.tex #3581
%[Tổ 3 - Đợt 16 - Chương 1 - Bài 3 - Cánh Diều - Đề 2]%[Đình Nguyên]
	Cho hàm số $y=f(x)=a x^3+b x^2+c x+d\quad (a \neq 0)$ có đồ thị như hình vẽ bên dưới\\
	\begin{center}
		
		\begin{tikzpicture}[line join=round, line cap=round,>=stealth,thick]
			\tikzset{every node/.style={scale=0.9}}
			\draw[->] (-3.1,0)--(3.1,0) node[above left] {$x$};
			\draw[->] (0,-3.1)--(0,5.1) node[below left] {$y$};
			\draw (0,0) node [below left] {$O$};
			\foreach \x/\nx in {-3/-3,-2/-2,-1/-1,1/1,2/2,3/3}
			\draw[thin] (\x,1pt)--(\x,-1pt) node [below] {$\nx$};
			\foreach \y/\ny in {-3/-3,-2/-2,-1/-1,1/1,2/2,3/3}
			\draw[thin] (1pt,\y)--(-1pt,\y) node [left] {$\ny$};
			\draw[dashed] (-1,0)--(-1,4)--(0,4);
			\begin{scope}
				\clip (-3,-3) rectangle (4,4);
				\draw[samples=200,domain=-3:3,smooth,variable=\x] plot (\x,{1*((\x)^3)+0*((\x)^2)+-3*(\x)+2});
			\end{scope}
		\end{tikzpicture}
		
	\end{center}
	Tìm số đường tiệm cận đứng của đồ thị hàm số $g(x)=\dfrac{\sqrt{x}}{f\left(x^2-1\right)-4}$.
	\loigiai{
		Điều kiện $x \geq 0$. Từ đồ thị, ta có $a > 0$.\\
		Ta có $f(x)-4=0\Leftrightarrow f(x)=4\Leftrightarrow\hoac{&x=-1\\
			&x=\alpha \quad (\alpha > 1)
			&}$ trong đó $x=-1$ là nghiệm kép.\\
		Suy ra $f(x)-4=a(x+1)^2(x-\alpha)$.\\
		Do đó $f\left(x^2-1\right)-4=a x^4\left(x^2-1-\alpha\right)$. \\
		Suy ra $ g(x)=\dfrac{\sqrt{x}}{f\left(x^2-1\right)-4}=\dfrac{\sqrt{x}}{a x^4\left(x^2-1-\alpha\right)}=\dfrac{1}{a x^3 \sqrt{x}(x-\sqrt{1+\alpha})(x+\sqrt{1+\alpha})}$.\\
		Ta có $\lim\limits_{x \to 0^{+}} g(x)=\lim\limits_{x \to 0^{+}} \dfrac{1}{a x^3\sqrt{x}(x-\sqrt{1+\alpha})(x+\sqrt{1+\alpha})}=-\infty$ nên $x=0$ là một đường tiệm cận đứng của đồ thị $y=g(x)$.\\
		Đường thẳng $x=\sqrt{1+\alpha}$ cũng là một đường tiệm cận đứng của đồ thị $y=g(x)$.\\
		Vậy đồ thị hàm số $g(x)$ có hai đường tiệm cận đứng.
	}
\end{bt}

\begin{bt}%[2D1C4-3]
% nguon: Archive/2024/San pham - To Latex Strong - Tex hoa Lop 12 moi - Dot 3 (Chuong 1 + 2 +3 + 4)/Data/Chuong 1/Canh dieu/To 3 - Dot 16 - Chuong 1 - Bai 3 - CD - De 2.tex #3582
%[Tổ 3 - Đợt 16 - Chương 1 - Bài 3 - Cánh Diều - Đề 2]%[Đình Nguyên]
	Từ một tấm tôn hình chữ nhật có các kích thước là $x$(m), $y$(m) với $x > 2$ và $y > 2$ và diện tích bằng $10$m$^2$, người ta cắt bốn hình vuông bằng nhau ở bốn góc rồi gập thành một cái thùng dạng hình hộp chữ nhật không nắp (như hình vẽ) có chiều cao bằng 1m. Thể tích của thùng là hàm số $V(x)$ trên khoảng $(2;+\infty)$.\\
	Đồ thị hàm số $y=\dfrac{1}{V(x)}$ có bao nhiêu đường tiệm cận đứng?
	\begin{multicols}{2}
		
		\begin{center}
			\begin{tikzpicture}[line join=round, line cap=round,thick]
				\coordinate (A) at (0,3);
				\coordinate (B) at (5,3);
				\coordinate (D) at (0,0);
				\coordinate (C) at ($(B)+(D)-(A)$);
				\draw(A)--(B)--(C)--(D)--cycle;
				\draw (0,0) rectangle (1,1) (A) rectangle (1,2) (B) rectangle (4,2) (4,1) rectangle (C);
				\draw (1,1) rectangle (4,2);
				%		\foreach \i/\g in {A/90,B/90,C/-90,D/-90}{\draw[fill=black](\i) circle (1pt) ($(\i)+(\g:3mm)$) node[scale=1]{$\i$};}
				%				\draw (0,.5) node [left]  {$x$};
				%				\draw (.5,0) node [below]  {$x$};
				%				\draw (0,2.5) node [left]  {$x$};
				%				\draw (0.5,3) node [above]  {$x$};
				%				%%%%%%%%%
				%				\draw (4.5,0) node [below]  {$x$};
				%				\draw (5,0.5) node [right]  {$x$};
				%				\draw (5,2.5) node [right]  {$x$};
				%				\draw (4.5,3) node [above]  {$x$};
				%%%%%%%%
				\draw[<->] (-1,0)--(-1,3);
				\draw[<->] (0,-1)--(5,-1);
				\draw (-1.2,1.5) node[left,rotate=0] {$x$};
				\draw (2.5,-0.8) node[above,rotate=0] {$y$};
			\end{tikzpicture}
			
			
			%%%%%%%%%%%%%%
			\begin{tikzpicture}[scale=1, font=\footnotesize, line join=round, line cap=round, >=stealth]
				\def\bc{4} \def\ba{2} \def\h{2} \def\gocB{35} 
				\coordinate[label=below left:$B$] (B) at (0,0);
				\coordinate[label=above left:$A$] (A) at (\gocB:\ba);
				\coordinate[label=below:$C$] (C) at (\bc,0);
				\coordinate[label=right:$D$] (D) at ($(C)-(B)+(A)$);
				\coordinate[label=above left:$A'$] (A') at ($(A)+(90:\h)$);
				\coordinate[label=left:$B'$] (B') at ($(B)-(A)+(A')$);
				\coordinate[label=below right:$C'$] (C') at ($(C)-(A)+(A')$);
				\coordinate[label=right:$D'$] (D') at ($(D)-(A)+(A')$);
				\draw (B')--(B)--(C)--(D)--(D')--(A')--(B')--(C')--(D') (C)--(C');
				\draw[dashed] (A')--(A)--(D) (A)--(B);
				%	\foreach \diem in {A,B,C,D,A',B',C',D'}	\fill (\diem)circle(1pt);
			\end{tikzpicture}
			
		\end{center}
	\end{multicols}
	%%%%%%%%%%%%%%
	
	\loigiai{
		Do tấm tôn có diện tích bằng $10$m$^2$ nên $x\cdot y=10\Leftrightarrow y=\dfrac{10}{x}$.\\
		Thùng có chiều cao là $1$m và các kích thước còn lại của thùng là $x-2$ và $y-2$.\\
		Thể tích của thùng là $V(x)=1\cdot (x-2)(y-2)=(x-2)\left(\dfrac{10}{x}-2\right)=\dfrac{(x-2)(10-2x)}{x}\cdot$\\
		Suy ra $y=\dfrac{1}{V(x)}=\dfrac{x}{(x-2)(10-2x)}$.\\
		Ta có $\lim\limits_{x \to 2^{+}} \dfrac{1}{V(x)}=\lim\limits_{x \to 2^{+}} \dfrac{x}{(x-2)(10-2x)}=+\infty$ và $\lim\limits_{x \to 2^{-}} \dfrac{1}{V(x)}=\lim\limits_{x \to 2^{-}} \dfrac{x}{(x-2)(10-2x)}=-\infty$.\\
		Suy ra đường thẳng $x=2$ là đường tiệm cận đứng của đồ thị hàm số $y=\dfrac{1}{V(x)}\cdot$\\
		Ta có
		$\lim\limits_{x \to 5^{+}} \dfrac{1}{V(x)}=\lim\limits_{x \to 4^{+}} \dfrac{x}{(x-2)(10-2 x)}=-\infty \text{ và } \lim\limits_{x \to 5^{-}} \dfrac{1}{V(x)}=\lim\limits_{x \to 4^{-}} \dfrac{x}{(x-2)(10-2 x)}=+\infty$.\\
		Suy ra đường thẳng $x=5$ là đường tiệm cận đứng của đồ thị hàm số $y=\dfrac{1}{V(x)}\cdot$\\
		Vậy đồ thị hàm số $y=\dfrac{1}{V(x)}$ có $2$ đường tiệm cận đứng. 
	}
\end{bt}

\begin{bt}%[2D1V4-1]
% nguon: Archive/2024/San pham - To Latex Strong - Tex hoa Lop 12 moi - Dot 3 (Chuong 1 + 2 +3 + 4)/Data/Chuong 1/Canh dieu/To 3 - Dot 16 - Chuong 1 - Bai 3 - CD - De 2.tex #3583
%[Tổ 3 - Đợt 16 - Chương 1 - Bài 3 - Cánh Diều - Đề 2]%[Đình Nguyên]
	Cho hàm số $f(x)=2x-\sqrt{x^2-x}$. Tìm số đường tiệm cận xiên của đồ thị hàm số
	\loigiai{
		Hàm số xác định và liên tục trên $\mathscr{D}=(-\infty; 0] \cup[1;+\infty)$.\\
		Ta có
		\[
		\begin{aligned}
			&a=\lim\limits_{x \to+\infty} \dfrac{f(x)}{x}=\lim\limits_{x \to+\infty} \dfrac{2 x-\sqrt{x^2-x}}{x}=\lim\limits_{x \to+\infty}\left(2-\sqrt{1-\dfrac{1}{x}}\right)=1. \\
			&b=\lim\limits_{x \to+\infty}(f(x)-x)=\lim\limits_{x \to+\infty}\left(x-\sqrt{x^2-x}\right)=\lim\limits_{x \to+\infty} \dfrac{x}{x+\sqrt{x^2-x}}=\lim\limits_{x \to+\infty} \dfrac{1}{1+\sqrt{1-\dfrac{1}{x}}}=\dfrac{1}{2}.
		\end{aligned}
		\]
		Suy ra	$ y=x+\dfrac{1}{2}$ là tiệm cận xiên của đồ thị hàm số khi $x \to+\infty$ 
		\[
		\begin{aligned}	
			&a=\lim\limits_{x \to-\infty} \dfrac{f(x)}{x}=\lim\limits_{x \to-\infty} \dfrac{2 x-\sqrt{x^2-x}}{x}=\lim\limits_{x \to-\infty}\left(2+\sqrt{1-\dfrac{1}{x}}\right)=3. \\
			&b=\lim\limits_{x \to-\infty}(f(x)-3 x)=-\lim\limits_{x \to-\infty}\left(x+\sqrt{x^2-x}\right)=-\lim\limits_{x \to-\infty} \dfrac{x}{x-\sqrt{x^2-x}}=-\lim\limits_{x \to-\infty} \dfrac{1}{1+\sqrt{1-\dfrac{1}{x}}}=-\dfrac{1}{2}.
		\end{aligned}
		\]
		Suy ra $y=3x-\dfrac{1}{2}$ là tiệm cận xiên của đồ thị hàm số khi $x \to-\infty$.\\
		Vậy đồ thị hàm số có $2$ đường tiệm cận xiên.
	}
\end{bt}

\begin{bt}%[2D1C4-4]
% nguon: Archive/2024/San pham - To Latex Strong - Tex hoa Lop 12 moi - Dot 3 (Chuong 1 + 2 +3 + 4)/Data/Chuong 1/Canh dieu/To 3 - Dot 16 - Chuong 1 - Bai 3 - CD - De 2.tex #3584
%[Tổ 3 - Đợt 16 - Chương 1 - Bài 3 - Cánh Diều - Đề 2]%[Đình Nguyên]
	\immini{Người ta muốn làm một cái bể dạng hình hộp chữ nhật không nắp (như hình vẽ) có thể tích bằng $5$m$^3$. Chiều cao của bể là $10$dm, các kích thước khác là $x$m, $y$m với $x > 0$ và $y > 0$. Diện tích toàn phần của bể (không kể nắp) là hàm số $S(x)$ trên khoảng $(0;+\infty)$. Đường tiệm cận xiên của đồ thị hàm số $S(x)$ là đường thẳng $y=a x+b$. Tính giá trị của biểu thức $P=a^2+b^2$}
	{
		\begin{tikzpicture}[scale=1, font=\footnotesize, line join=round, line cap=round, >=stealth]
			\def\bc{4} \def\ba{2} \def\h{2} \def\gocB{35} 
			\coordinate[label=below left:$B$] (B) at (0,0);
			\coordinate[label=above left:$A$] (A) at (\gocB:\ba);
			\coordinate[label=below:$C$] (C) at (\bc,0);
			\coordinate[label=right:$D$] (D) at ($(C)-(B)+(A)$);
			\coordinate[label=above left:$A'$] (A') at ($(A)+(90:\h)$);
			\coordinate[label=left:$B'$] (B') at ($(B)-(A)+(A')$);
			\coordinate[label=below right:$C'$] (C') at ($(C)-(A)+(A')$);
			\coordinate[label=right:$D'$] (D') at ($(D)-(A)+(A')$);
			\draw (B')--(B)--(C)--(D)--(D')--(A')--(B')--(C')--(D') (C)--(C');
			\draw[dashed] (A')--(A)--(D) (A)--(B);
			%	\foreach \diem in {A,B,C,D,A',B',C',D'}	\fill (\diem)circle(1pt);
		\end{tikzpicture}
		
	}
	
	\loigiai{
		Do thể tích của bể là $1$m$^3$ nên $1\cdot x  y \Leftrightarrow xy=5$.\\
		Diện tích toàn phần của bể là $S(x)=xy+2x+2y=5+2x+\dfrac{10}{x}, \quad x>0$.\\
		Ta có $\lim\limits_{x \to+\infty}\left(S(x)-5-2 x\right)=\lim\limits_{x \to+\infty}\left(\dfrac{10}{x}\right)=0$.\\
		Suy ra đồ thị hàm số $S(x)$ có đường tiệm cận xiên là $y=2x+5$.\\
		Suy ra $a=2$; $b=5$.\\
		Vậy $P=a^2+b^2+2^2+5^2=29$.
	}
\end{bt}

\begin{bt}%[2D1C3-1]
% nguon: Archive/SP_KhoLatex/ID6_old/DeThamKhao-08.tex #3607
%[Dự án ĐGNL-ĐGTD-2022 GĐ1]%[Thành Đức Trung]
		Có bao nhiêu số nguyên $m$ để giá trị lớn nhất của hàm số $y=\left|x^{3}-2 x^{2}+(m+2) x+5\right|$ trên đoạn $[-1 ; 2]$ không vượt quá $11$?

\textit{Kết quả:} \tikz\draw (0,0) rectangle (3,0.6);
		\loigiai{
			Ta có
			\begin{eqnarray*}
				\max _{[-1 ; 2]} y \leq 11 & \Leftrightarrow & \left|x^{3}-2 x^{2}+(m+2) x+5\right| \leq 11, \forall x \in[-1 ; 2] \\
				&\Leftrightarrow & -11 \leq x^{3}-2 x^{2}+(m+2) x+5 \leq 11, \forall x \in[-1 ; 2] \\
				&\Leftrightarrow& -x^{3}+2 x^{2}-16 \leq(m+2) x \leq-x^{3}+2 x^{2}+6, \forall x \in[-1 ; 2]\quad (1).
			\end{eqnarray*}
			Đặt $g(x) = -x^{2}+2 x-\dfrac{16}{x}$; $h(x) = -x^{2}+2 x+\dfrac{6}{x}$. 
			Khi đó
			\begin{eqnarray*}
				(1) & \Leftrightarrow  & \heva{&-x^{2}+2 x-\dfrac{16}{x} \leq m+2 \leq -x^{2}+2 x+\dfrac{6}{x}, \forall x \in(0 ; 2] \\ &-x^{2}+2 x+\dfrac{6}{x} \leq m+2 \leq -x^{2}+2 x-\dfrac{16}{x}, \forall x \in[-1 ; 0)} \\
				&\Leftrightarrow & \heva{& \max _{(0 ; 2]} g(x) \leq m+2 \leq \min _{(0 ; 2]} h(x) \\& \max _{[-1 ; 0)} h(x) \leq m+2 \leq \min _{[-1 ; 0)} g(x) } \\
				&\Leftrightarrow  & \heva{& -8 \leq m+2 \leq 3 \\ &-9 \leq m+2 \leq 13} \\
				& \Leftrightarrow & -10 \leq m \leq 1.
			\end{eqnarray*}
		Vậy có 12 giá trị nguyên của tham số $m$ thỏa mãn.
		} 
\end{bt}

\begin{bt}%[2D1H3-1]
% nguon: Archive/SP_KhoLatex/ID6_old/DeThamKhao-08.tex #3608
%[Dự án ĐGNL-ĐGTD-2022 GĐ1]%[Thành Đức Trung]
	Tìm giá trị thực của tham số $m$ để hàm số $f(x)=-x^3-3x^2+m$ có giá trị nhỏ nhất trên đoạn $[-1;1]$ bằng $0$.

	\textit{Kết quả:} \tikz\draw (0,0) rectangle (3,0.6);
	\loigiai{
		Đạo hàm $f'(x)=-3x^2-6x$; $f'(x)=0\Leftrightarrow \hoac{&x=-2\notin [-1;1]\\&x=0\notin [-1;1]}$.\\
		Ta có $\heva{&f(-1)=m-2\\&f(0)=m\\&f(1)=m-4}\Rightarrow \min\limits_{[-1;1]}f(x)=f\left(1\right)=m-4$.\\
		Theo bài ra $\min\limits_{[-1;1]}f(x)=0\Leftrightarrow m-4=0\Leftrightarrow m=4$.\\
		Vậy $m = 4$.
	} 
\end{bt}

\begin{bt}%[2D1H4-1]
% nguon: Archive/SP_KhoLatex/ID6_old/DeThamKhao-08.tex #3609
%[Dự án ĐGNL-ĐGTD-2022 GĐ1]%[Thành Đức Trung]
	Cho hàm số $ y=f(x) $ có bảng biến thiên như sau:
	\begin{center}
		\begin{tikzpicture}[>=stealth,scale=0.85]
			\tkzTabInit[nocadre=false,lgt=1,espcl=2,deltacl=0.5]{$x$/.7 ,$y'$/.7,$y$/2}
			{$-\infty$ , $-1$ , $1$ , $+\infty$}
			\tkzTabLine{ , + , d , h , d , - , }
			\tkzTabVar{-/$+0$ , +DH/$2$, D+/$+\infty$ , -/ $1$}
		\end{tikzpicture}	
	\end{center}
	Tổng số tiệm cận ngang và tiệm cận đứng của đồ thị hàm số đã cho là

	\textit{Kết quả:} \tikz\draw (0,0) rectangle (3,0.6); 
	\loigiai{ Từ bảng biến thiên, ta có
		\begin{enumerate}
			\item	$\lim\limits_{x\to-\infty}f(x)=0\Rightarrow y=0$ là tiệm cận ngang.
			\item $\lim\limits_{x\to+\infty}f(x)=1\Rightarrow y=1$ là tiệm cận ngang.
			\item $\lim\limits_{x\to 1^{+}}f(x)=+\infty\Rightarrow x=1$ là tiệm cận đứng.
			\item $\lim\limits_{x\to(-1)^{-}}f(x)=2\Rightarrow x=-2$ không là tiệm cận đứng.
		\end{enumerate}
		Vậy đồ thị hàm số có tất cả  $ 3 $ đường tiệm cận.
	}
\end{bt}

\begin{bt}%[2D1H2-1]
% nguon: Archive/SP_KhoLatex/ID6_old/DeThamKhao-08.tex #3610
%[Dự án ĐGNL-ĐGTD-2022 GĐ1]%[Thành Đức Trung]
	Hàm số $y=-x^4+2x^2+5$ có bao nhiêu điểm cực trị?

	\textit{Kết quả:} \tikz\draw (0,0) rectangle (3,0.6); 
	\loigiai{
		Ta có $y'=-4x^3+4x$; Giải phương trình $y'=0\Leftrightarrow-4x\left(x^2-1\right)=0\hoac{&x=0\\&x=\pm 1.}$ \\
		Lập bảng biến thiên ta có
		\begin{center}
			\begin{tikzpicture}
				\tkzTabInit[nocadre=false,lgt=1.2,espcl=2.5,deltacl=0.6]
				{$x$ /0.6,$y'$ /0.6,$y$ /2}
				{$-\infty$,$-1$,$0$,$1$,$+\infty$}
				\tkzTabLine{,+,$0$,-,$0$,+,$0$,-,}
				\tkzTabVar{-/$-\infty$, +/$6$,-/$5$,+/$6$,-/$-\infty$}
			\end{tikzpicture}
		\end{center}
		Từ bảng biến thiên ta có hàm số có $3$ điểm cực trị.}
\end{bt}

\begin{bt}%[2D1V1-3]
% nguon: Archive/SP_KhoLatex/ID6_old/DeThamKhao-08.tex #3611
%[Dự án ĐGNL-ĐGTD-2022 GĐ1]%[Thành Đức Trung]
	Hàm số $y=(x+m)^3+(x+n)^3-x^3$ đồng biến trên khoảng $(-\infty;+\infty)$. Giá trị nhỏ nhất của biểu thức $P=4\left(m^2+n^2\right)-m-n$ bằng $-\dfrac{1}{k}$ với $k\in\mathbb{N}$. Giá trị của $b$ bằng

	\textit{Kết quả:} \tikz\draw (0,0) rectangle (3,0.6); 
	\loigiai{
		Ta có $y'=3(x+m)^2+3(x+n)^2-3x^2=3\left[x^2+2(m+n)x+m^2+n^2\right]$.\\
		Hàm số đồng biến trên $(-\infty;+\infty)\Leftrightarrow\heva{&a>0\\&\Delta\leq 0}\Leftrightarrow mn\leq 0$.
		\begin{enumerate}
			\item[TH1:] $mn=0\Leftrightarrow\hoac{&m=0\\&n=0.}$ \\
			Do vai trò của $m,n$ là như nhau nên ta chỉ cần xét trường hợp $m=0$ \\
			$ \Rightarrow P=4n^2-n=\left(2n-\dfrac{1}{4}\right)-\dfrac{1}{16}\geq-\dfrac{1}{16} \quad(1) $.
			\item[TH2:] $m n<0\Leftrightarrow m>0; n<0$.\\
			Ta có $P=\left(2m-\dfrac{1}{4}\right)^2-\dfrac{1}{16}+4n^2+(-n) >-\dfrac{1}{16}\quad (2)$.
		\end{enumerate}
		Từ $(1),(2)$ ta có $P_{\min} =-\dfrac{1}{16}$. Dấu xảy ra khi và chỉ khi $m=\dfrac{1}{8};n=0$ hoặc $m=0;n=\dfrac{1}{8}$.\\
		Vậy $b = 16$.
		}
\end{bt}

\begin{bt}%[2D1V2-6]
% nguon: Archive/SP_KhoLatex/ID6_old/DeThamKhao-08.tex #3612
%[Dự án ĐGNL-ĐGTD-2022 GĐ1]%[Thành Đức Trung]
	\immini{Cho hàm số $f(x)$ với đạo hàm $f'(x)$ có đồ thị như hình vẽ. Hàm số $g(x)=f(x)-\dfrac{x^3}{3}+x^2-x+2$ đạt cực đại tại điểm nào?
		}{\begin{tikzpicture}[scale=1,>=stealth, font=\footnotesize, line join=round, line cap=round]
			\def\a{1} \def\b{-2} \def\c{0} \def\d{1} % Hệ số
			\def\xmin{-2} \def\xmax{3}
			\def\ymin{-3} \def\ymax{3}
			\draw[color=gray!50,dashed] (\xmin,\ymin) grid (\xmax,\ymax);
			\draw[->] (\xmin,0)--(\xmax,0) node [below]{$x$};
			\draw[->] (0,\ymin)--(0,\ymax) node [left]{$y$};
			\node at (0,0) [below left]{$O$};
			\clip (\xmin+0.1,\ymin+0.1) rectangle (\xmax-0.5,\ymax-0.1);
			\draw[smooth,samples=300] plot(\x,{\a*(\x)^3+\b*(\x)^2+\c*(\x)+\d});
			\draw[dashed](-1,0)node[above]{$-1$}--(-1,-2)--(0,-2)node[right]{$-2$} (1,0)node[below]{$1$} (2,0)node[below]{$2$} (0,1)node[above left]{$1$};
	\end{tikzpicture}}

	\textit{Kết quả:} \tikz\draw (0,0) rectangle (3,0.6); 
	\loigiai{
		Ta có $g'(x)=f'(x)-(x-1)^2$.\\
		Điểm cực trị của hàm số $y=g(x)$ là nghiệm của phương trình $g'(x)=0$ tức là nghiệm của phương trình $f'(x)=(x-1)^2$ suy ra điểm cực trị của hàm số $y=g(x)$ cũng là hoành độ giao điểm của các đồ thị hàm số $y=f'(x); y=x^2-2x+1$.\\
		Vẽ đồ thị của các hàm số $y=f'(x); y=x^2-2x+1$ trên cùng một hệ trục tọa độ như hình vẽ sau:\\
		\begin{center}
			\begin{tikzpicture}[scale=1,>=stealth, font=\footnotesize, line join=round, line cap=round]
				\def\a{1} \def\b{-2} \def\c{0} \def\d{1} % Hệ số
				\def\xmin{-2} \def\xmax{3}
				\def\ymin{-3} \def\ymax{3}
				\draw[color=gray!50,dashed] (\xmin,\ymin) grid (\xmax,\ymax);
				\draw[->] (\xmin,0)--(\xmax,0) node [below]{$x$};
				\draw[->] (0,\ymin)--(0,\ymax) node [left]{$y$};
				\node at (0,0) [below left]{$O$};
				\clip (\xmin+0.1,\ymin+0.1) rectangle (\xmax-0.5,\ymax-0.1);
				\draw[red, smooth,samples=300] plot(\x,{\a*(\x)^3+\b*(\x)^2+\c*(\x)+\d});
				\draw[blue, smooth,samples=300] plot(\x,{(\x)^2-2*(\x)+1});
				\draw[dashed](-1,0)node[above]{$-1$}--(-1,-2)--(0,-2)node[right]{$-2$} (1,0)node[below]{$1$} (2,0)node[below]{$2$} (0,1)node[above left]{$1$};
			\end{tikzpicture}
		\end{center}
		Dựa vào đồ thị trên ta có bảng biến thiên của hàm số $y=g(x)$ như sau
		\begin{center}
			\begin{tikzpicture}
				\tkzTabInit[nocadre=false,lgt=1.2,espcl=2.5,deltacl=0.6]
				{$x$ /0.6,$g'(x)$ /0.6,$g(x)$ /2}
				{$-\infty$,$0$,$1$,$2$,$+\infty$}
				\tkzTabLine{,-,$0$,+,$0$,-,$0$,+,}
				\tkzTabVar{+/$ $, -/$ $,+/$ $,-/$ $,+/$ $}
			\end{tikzpicture}
		\end{center}
		Dựa vào BBT ta thấy hàm số $y=g(x)$ có điểm cực đại $x=1$.}
\end{bt}

\begin{bt}%[2D1C2-5]
% nguon: Archive/SP_KhoLatex/ID6_old/DeThamKhao-08.tex #3613
%[Dự án ĐGNL-ĐGTD-2022 GĐ1]%[Thành Đức Trung]
	Gọi $S$ là tập các giá trị của tham số $m$ để đồ thị hàm số $y=x^4-2mx^2+m$ có ba điểm cực trị, đồng thời ba điểm cực trị đó là ba đỉnh của một tam giác có bán kính đường tròn nội tiếp lớn hơn $ 1 $. Biết $S = \left(a;+\infty\right)$, giá trị của $a$ bằng

	\textit{Kết quả:} \tikz\draw (0,0) rectangle (3,0.6); 
	\loigiai{
		Ta có $y'=4x^{3}-4mx=4x(x^{2}-m)$.\\
		Hàm số có 3 điểm cực trị khi $m>0$.\\
		Ba điểm cực trị là $A(0;m),B\left(-\sqrt{m};m-m^{2}\right),C\left(\sqrt{m};m-m^{2}\right)$.\\
		Gọi $I$ là trung điểm của $BC\Rightarrow I\left(0;m-m^{2}\right)$.\\
		$S_{\triangle ABC}=\dfrac{1}{2}AI\cdot BC=m^{2}\sqrt{m}$.\\
		Chu vi của $\triangle ABC$ là $2p=AB+BC+AC=2\left(\sqrt{m+m^{4}}+\sqrt{m}\right)$.\\
		Bán kính đường tròn nội tiếp $\triangle ABC$ là $r=\dfrac{S_{\triangle ABC}}{p}=\dfrac{m^{2}\sqrt{m}}{\sqrt{m+m^{4}}+\sqrt{m}}$.\\
		Theo bài ra 
		\allowdisplaybreaks{\begin{eqnarray*}
				r>1 & \Leftrightarrow & \dfrac{m^{2}\sqrt{m}}{\sqrt{m+m^{4}}+\sqrt{m}}>1\\
				& \Leftrightarrow & \dfrac{m^{2}\sqrt{m}\left(\sqrt{m+m^{4}}-\sqrt{m}\right)}{m^{4}}>1\,(\text{vì }m>0)\\
				& \Leftrightarrow & \sqrt{m}\left(\sqrt{m+m^{4}}-\sqrt{m}\right)>m^{2}\\
				& \Leftrightarrow & \sqrt{m^{2}+m^{5}}>m^{2}+m\\
				& \Leftrightarrow & m^{2}-m-2>0\\
				& \Leftrightarrow & \left[\begin{aligned} & m<-1\\
					& m>2.
				\end{aligned}
				\right.
		\end{eqnarray*}}
		So sánh điều kiện suy ra $m>2$ thỏa mãn. Vậy $a = 2$.
	}
\end{bt}

\begin{bt}%[2D1V4-2]
% nguon: Archive/SP_KhoLatex/ID6_old/DeThamKhao-08.tex #3616
%[Dự án ĐGNL-ĐGTD-2022 GĐ1]%[Thành Đức Trung]
	Có bao nhiêu giá trị thực của tham số $m$ để đồ thị hàm số $y=\dfrac{2x^2-3x+m}{x-m}$ không có tiệm cận đứng?

	\textit{Kết quả:} \tikz\draw (0,0) rectangle (3,0.6); 
	\loigiai{
		TXĐ: $\mathscr{D}=\mathbb{R}\backslash \left\{m\right\}.$\\
		Ta có $y=\dfrac{\left(x-m\right)\left(2x+2m-3\right)+2m\left(m-1\right)}{x-m}=2x+2m-3+\dfrac{2m\left(m-1\right)}{x-m}.$\\
		Để đồ thị hàm số không có tiệm cận đứng thì các giới hạn 
		$$\underset{x\to m^{\pm}}{\mathop{\lim}}y \text{ tồn tại hữu hạn }\Leftrightarrow 2m\left(m-1\right)=0\Leftrightarrow \hoac{
			& m=1 \\ 
			& m=0.}$$
		Vậy có 2 giá trị thực của tham số $m$ thỏa mãn.
	}
\end{bt}

\begin{bt}%[2D1V5-1]
% nguon: Archive/SP_KhoLatex/ID6_old/DeThamKhao-08.tex #3620
%[Dự án ĐGNL-ĐGTD-2022 GĐ1]%[Thành Đức Trung]
	\immini{
		Cho hàm số $y=ax^3+bx^2+cx+d$ có đồ thị là đường cong trong hình vẽ. Có bao nhiêu số dương trong các $a$, $b$, $c$, $d$?
	}{
	\begin{tikzpicture}[scale=.75, font=\footnotesize, line join=round, line cap=round, >=stealth]
			\draw[->] (-4,0)--(1,0) node[below] {$x$};
			\draw[->] (0,-2)--(0,3) node[right] {$y$};
			\draw (0,0) node [below left] {$O$};
			\begin{scope}
				\clip (-4,-2) rectangle (1,3);
				\draw[samples=200,domain=-4:2,thick,variable=\x] plot (\x,{-(2/3)*(\x+3)*(\x+2)*(\x-0.5)});
			\end{scope}
	\end{tikzpicture}}

	\textit{Kết quả:} \tikz\draw (0,0) rectangle (3,0.6); 
	\loigiai{
		Ta có $y'=3ax^2+2bx+c$ . Dựa vào đồ thị ta thấy $a < 0$.\\
		Hàm số có 2 cực trị âm nên 
		\[\heva{&\Delta'_{y'}> 0\\&S<0\\&P>0}\Leftarrow \heva{&b^2-9ac>0\\&-\dfrac{2b}{3a}<0\\&\dfrac{c}{3a}>0}\Rightarrow \heva{&b<0\\&c<0.}\]
		Đồ thị cắt trục $Oy$ tại điểm $\left(0;d\right)$ nên $d > 0$.\\
		Vậy có đúng một số dương trong các số $a$, $b$, $c$, $d$.}
\end{bt}

\begin{bt}%[2D1V3-1]
% nguon: Archive/SP_KhoLatex/ID6_old/DeThamKhao-08.tex #3621
%[Dự án ĐGNL-ĐGTD-2022 GĐ1]%[Thành Đức Trung]
		Gọi $S$ là tập hợp tất cả các giá trị thực của tham số $m$ sao cho giá trị nhỏ nhất của hàm số $f(x)=\dfrac{34}{\sqrt{\left(x^{3}-3 x+2 m\right)^{2}}+1}$ trên đoạn $[0 ; 3]$ bằng $2$. Tổng tất cả các phần tử của $S$ bằng

	\textit{Kết quả:} \tikz\draw (0,0) rectangle (3,0.6); 
		\loigiai{
			Ta có $f(x)=\dfrac{34}{\sqrt{\left(x^{3}-3 x+2 m\right)^{2}}+1}=\dfrac{34}{\left|x^{3}-3 x+2 m\right|+1}=\dfrac{34}{g(x)+1};\  g(x)=\left|x^{3}-3 x+2 m\right|$. \\
			Suy ra  $\displaystyle  \min _{[0 ; 3]} f(x)=\dfrac{34}{ \underset {[0 ; 3]}{\max} ~ g(x)+1}=2 \Leftrightarrow \max _{[0 ; 3]} g(x)=16$. \\
			Xét hàm số $u=x^{3}-3 x+2 m$, dễ thấy $\displaystyle  a=\min _{[0 ; 3]} u=u(1)=2 m-2 ; A=\max _{[0 ; 3]} u=u(3)=2 m+18$. \\
			Vì vậy 
			\begin{align*}
				&  \max _{[0 ; 3]} g(x)  =\max \{|a|,|A|\}=\max \{|2 m-2|,|2 m+18|\} =16  \\
				\Leftrightarrow~& \hoac{& |2 m+18| \leq |2 m-2| = 16 \\& |2 m-2| \leq |2 m+18| = 16 }  \\
				\Leftrightarrow~&  \hoac{&m= -7\\&m=-1.}
			\end{align*}
		Vậy tổng các phần tử của $S$ bằng $-8$.
		}
\end{bt}

\begin{bt}%[2D1C3-6]
% nguon: Archive/SP_KhoLatex/ID6_old/DeThamKhao-08.tex #3623
%[Dự án ĐGNL-ĐGTD-2022 GĐ1]%[Thành Đức Trung]
Gấp góc bên phải của một tờ giấy A$4$ kích thước $297\times210$ mm sao cho sau khi gấp đỉnh ở góc chạm đáy (hình vẽ). Tình độ dài nếp gấp nhỏ nhất.
\begin{center}
\begin{tikzpicture}[scale=1, font=\footnotesize, line join=round, line cap=round, >=stealth]
\path
(0,0)coordinate(D)
(0,2)coordinate(A)
(4,0)coordinate(C)
($(A)+(C)-(D)$)coordinate(B)
($(A)+(6,0)$)coordinate(A')
($(B)+(6,0)$)coordinate(B')
($(C)+(6,0)$)coordinate(C')
($(D)+(6,0)$)coordinate(D')
($(A')+(2,0)$)coordinate(A'')
($(D')+(2,0)$)coordinate(D'')
($(C')+(0,0.5)$)coordinate(E)
($(C')+(-1.414213562,0)$)coordinate(F)
;
\draw(A)--(B) (A)--(D) (B)--(C) (C)--(D) (A')--(A'') (A')--(D') (D')--(C') (C')--(E)--(A'')--(F)--(E);
\draw[dashed](A'')--(B')--(E);
\draw[->](4.5,1)--(5.5,1);
\tkzMarkRightAngle(A'',B',C');
\tkzMarkRightAngle(A'',F,E);
\end{tikzpicture}
\end{center}
\loigiai
{
\begin{center}
\begin{tikzpicture}[scale=1, font=\footnotesize, line join=round, line cap=round, >=stealth]
\path
(0,0)coordinate(D)
(0,2)coordinate(A)
(4,0)coordinate(C)
($(A)+(C)-(D)$)coordinate(B)
($(A)+(6,0)$)coordinate(A')
($(B)+(6,0)$)coordinate(B')node[above]{$B$}
($(C)+(6,0)$)coordinate(C')node[below]{$C$}
($(D)+(6,0)$)coordinate(D')
($(A')+(2,0)$)coordinate(A'')node[above]{$A$}
($(D')+(2,0)$)coordinate(D'')node[below]{$D$}
($(C')+(0,0.5)$)coordinate(E)node[right]{$E$}
($(C')+(-1.414213562,0)$)coordinate(F)node[below]{$F$}
;
\draw(A)--(B) (A)--(D) (B)--(C) (C)--(D) (A')--(A'') (A')--(D') (D')--(C') (C')--(E)--(A'')--(F)--(E)(D'')--(A'');
\draw[dashed](A'')--(B')--(E);
\draw[->](4.5,1)--(5.5,1);
\tkzMarkRightAngle(A'',B',C');
\tkzMarkRightAngle(A'',F,E);
\tkzMarkRightAngle(A'',D'',F);
\end{tikzpicture}
\end{center}
Đặt $EF=BE=x$ mm $\Rightarrow EC=210-x\Rightarrow FC=\sqrt{x^2-(210-x)^2}=\sqrt{420x-210^2}$.\\
Ta có $\triangle ADF\backsim\triangle FCE$ (g.g), suy ra $\dfrac{EF}{AF}=\dfrac{CF}{AD}\Rightarrow AF=\dfrac{EF\cdot AD}{FC}=\dfrac{210x}{\sqrt{420x-210^2}}$.\\
Đặt $y=AE$. Khi đó,
\[y=AE=\sqrt{AF^2+EF^2}=\sqrt{\dfrac{210^2x^2}{420x-210^2}+x^2}=\sqrt{\dfrac{420x^3}{420x-210^2}}.\]
Xét hàm số $f(x)=\dfrac{420x^3}{420x-210^2}$ với $105<x<210$.\\
Ta có $f'(x)=\dfrac{420\cdot3x^2\left(420x-210^2\right)-420\cdot420x^3}{\left(420x-210^2\right)^2}$.\\
Xét
\[\begin{aligned}
f'(x)=0
&\Leftrightarrow \dfrac{420\cdot3x^2\left(420x-210^2\right)-420\cdot420x^3}{\left(420x-210^2\right)^2}=0\\&\Leftrightarrow 3x^2\left(420x-210^2\right)-420x^3=0\\&\Leftrightarrow \hoac{ & x=0\\ & x=157{,}5.}
\end{aligned}\]
Bảng biến thiên của $f(x)$
\begin{center}
\begin{tikzpicture}[scale=1]
\tkzTabInit[nocadre=false, lgt=1.2, espcl=2.5, deltacl=0.6]{$x$/0.6, $f'(x)$/0.6, $f(x)$/2}{$105$, $157{,}5$, $210$}
\tkzTabLine{,-,0,+,}
\tkzTabVar{+/, -/$f\left(157{,}5\right)$, +/}
\end{tikzpicture}
\end{center}	
Suy ra $\min y=\min\limits_{(105;210)} \sqrt{f(x)}=\sqrt{f(157{,}5)}=\dfrac{315\sqrt{3}}{2}$. \\
Vậy độ dài nếp gấp nhỏ nhất là $\dfrac{315\sqrt{3}}{2}$ mm. \\
}
\end{bt}

\begin{bt}%[2D1V5-6]
% nguon: Archive/SP_KhoLatex/ID6_old/bai6.3.tex #3624

Cho hàm số $y=\dfrac{ax+b}{cx+d}$ ($ad-bc \ne 0; c \ne 0$) có đồ thị $(C)$, $M$ là điểm bất kì thuộc $(C)$. 
Tiếp tuyến tại $M$ cắt hai đường tiệm cận của $(C)$ tại $A$ và $B$.
\begin{enumerate}
\item  Gọi $I$ là giao điểm hai đường tiệm cận. Chứng minh rằng: $S_{IAB} = \dfrac{2|ad-bc|}{c^2}$.
\item Chứng minh rằng: Tích các khoảng cách từ $M$ đến hai đường tiệm cận bằng $\dfrac{|ad-bc|}{c^2}$.
\end{enumerate}
\loigiai{
\begin{center}
\begin{tikzpicture}[font=\footnotesize, line join=round, line cap=round,>=stealth, x=0.7cm, y=0.7cm]
		\def \xmin{-3.0};
		\def \xmax{5.5};
		\def \ymin{-3.0};
		\def \ymax{4.5};
		\coordinate (M) at (3,2) ;
		\coordinate (A) at (5,1) ;
		\coordinate (B) at (1,3) ;
		\coordinate (H) at (3,1) ;
		\coordinate (K) at (1,2) ;
		\draw[->] (\xmin, 0.) -- (\xmax,0.) node[anchor=north] {$x$};
		\draw[->] (0.,\ymin) -- (0.,\ymax) node[anchor=west] {$y$};
		\clip(\xmin,\ymin) rectangle (\xmax,\ymax);
		\draw[smooth,samples=100,domain=\xmin:0.99] plot(\x,{(\x+1)/(\x-1)});
		\draw[smooth,samples=100,domain=1.01:\xmax] plot(\x,{(\x+1)/(\x-1)});
		\draw (\xmin,1) -- (\xmax,1);
		\draw (1,\ymin) -- (1,\ymax);
		\draw (A)--(B) (H)--(M)--(K);
		\draw[fill=black] (0,0) circle (1pt) node[above left] {$O$}
		(1,1) circle (1pt) node[above left] {$I$} (3,1) circle (1pt) node[below] {$H$}
		(1,2) circle (1pt) node[left] {$K$};
		\foreach \x/\g in{A/-90,B/180,M/50}
\fill[black](\x)circle(1.1pt) ($(\x)+(\g:3mm)$)node{$\x$}; 
		\end{tikzpicture}
\end{center}
\begin{enumerate}
\item 
Gọi $M(x_0;y_0)$ là điểm thuộc $(C)$. \\
Tiếp tuyến tại $M$ có phương trình
\[
y= \dfrac{ad-bc}{(cx_0+d)^2} (x-x_0) + \dfrac{ax_0+b}{cx_0+d}  ~ (d). 
\]
Giao điểm của TCN và $(d)$ là $A\left ( \dfrac{2cx_0+d}{c} ;\dfrac{a}{c}\right )$. \\
Giao điểm của TCĐ và $(d)$ là $B\left (-\dfrac{d}{c} ; \dfrac{ac^2x_0^2 +2bc^2x_0 +2bcd -ad^2}{c(cx_0+d)^2}\right )$. \\
Suy ra $IA = |x_a - x_I| = \left |\dfrac{2(cx_0+d)}{c}\right |$; $IB =\left |\dfrac{2(bc-ad)}{c(cx_0+d)}\right |$. \\
Vậy $S_{IAB} = \dfrac{1}{2}IA \cdot IB = \dfrac{2|ad-bc|}{c^2}$.
\item Gọi $H$, $K$ là hình chiếu vuông góc của $M$ trên TCN và TCĐ. \\
Ta có $IH = |y_M-y_I| = \left |\dfrac{bc-ad}{c(cx_0+d)}\right |$; $IK = |x_M-x_I| =\left | \dfrac{cx_0+d}{c} \right |$.\\
Suy ra $IH \cdot IK = \dfrac{|ad-bc|}{c^2}$.
\end{enumerate}
}
\end{bt}

\begin{bt}%[2D1H1-1]
% nguon: Archive/SP_KhoLatex/ID6_old/DeThamKhao-01.tex #3631
%[Dự án ĐGNL-ĐGTD-2022 GĐ1]%[Nguyễn Anh Quốc]%[Đề thi HK2, Lương Thế Vinh, Hà Nội 2018]%[Lê Nguyễn Viết Tường]
Tìm các khoảng đồng biến của hàm số $y=x^3-3x+4$.\\
\textit{Kết quả:} \tikz\draw (0,0) rectangle (3,0.6); 
	\loigiai{
		Ta có $y'=3x^2-3$.\\
		$y'=0\Leftrightarrow x=\pm 1$.\\
		Bảng biến thiên
		\begin{center}
			\begin{tikzpicture}
			\tkzTabInit
			[lgt=2,espcl=3] % tùy chọn
			{$x$/1, $y'$/1, $y$/2.5} % cột đầu tiên
			{$-\infty$,$-1$,$1$, $+\infty$} % hàng 1 cột 2
			\tkzTabLine{,+,z,-,z,+,} % hàng 2 cột 2
			\tkzTabVar{-/ $-\infty$,+/ /R,-/ /R,+/ $+\infty$} % hàng 3 cột 2
			
\end{tikzpicture}
		\end{center}
		Vậy hàm số đã cho đồng biến trên $(-3;-2)$.
	}
\end{bt}

\begin{bt}%[2D1V1-3]
% nguon: Archive/SP_KhoLatex/ID6_old/DeThamKhao-01.tex #3632
%[Dự án ĐGNL-ĐGTD-2022 GĐ1]%[Nguyễn Anh Quốc]
	Có bao nhiêu số nguyên dương $m$ để hàm số $y=\dfrac{2}{3}x^3-(2m+9)x^2+2(m^2+9m)x+10$ nghịch biến trên khoảng $(3;6)$?\\
	\textit{Kết quả:} \tikz\draw (0,0) rectangle (3,0.6); 
	\loigiai{
		Ta có $y'=2x^2-2(2m+9)x	+2(m^2+9m), y'=0\Leftrightarrow \hoac{& x=m \\ & x=m+9.}$\\
		Hàm số đã cho liên tục trên đoạn $[3;6]$.\\
		Như vậy để hàm số nghịch biến trên $(3;6)$ thì $y'\le 0\ \forall x\in [3;6]$.
		Ta có $y'\le 0\Leftrightarrow m\le x\le m+9$. Từ đó ta suy ra
		$$\heva{& m\le 3 \\ & m+9\ge 6}\Leftrightarrow -3\le m\le 3.$$
		Vì $m\in\mathbb{N}^*$ nên $m\in \{1;2;3\}$.\\
		Vậy có tất cả $3$ số nguyên dương thỏa mãn.
	}
\end{bt}

\begin{bt}%[2D1H2-1]
% nguon: Archive/SP_KhoLatex/ID6_old/DeThamKhao-01.tex #3633
%[Dự án ĐGNL-ĐGTD-2022 GĐ1]%[Nguyễn Anh Quốc] 
	Cho hàm số $y=\dfrac{x^2+2x+3}{x+1}$ có hai cực trị $y_1$, $y_2$ (với $y_1>y_2$). Tính giá trị biểu thức $P=\dfrac{3\left(y_1^2+y_2^2\right)}{y_1-y_2}$.\\
		\textit{Kết quả:} \tikz\draw (0,0) rectangle (3,0.6); 
	\loigiai{
		Ta có $y'=\dfrac{x^2+2x-1}{(x+1)^2}$, $y'=0\Leftrightarrow \hoac{& x=-1+\sqrt{2}\\ & x=-1-\sqrt{2}}$. Suy ra $y_1=2\sqrt{2}$, $y_2=-2\sqrt{2}$.\\
		Nên $P=\dfrac{3(8+8)}{4\sqrt{2}}=6\sqrt{2}$.
	}
\end{bt}

\begin{bt}%[2D1H2-3]
% nguon: Archive/SP_KhoLatex/ID6_old/DeThamKhao-01.tex #3634
%[Dự án ĐGNL-ĐGTD-2022 GĐ1]%[Nguyễn Anh Quốc]
	Tìm tất cả các giá trị của tham số $m$ để hàm số $y=x^3-3x^2+mx-2$ đạt cực tiểu tại $x=2$ ?\\
	\textit{Kết quả:} \tikz\draw (0,0) rectangle (3,0.6); 
	\loigiai{
		Tập xác định: $\mathscr{D}=\mathbb{R}$. \\
		Ta có $y'=3x^2-6x+m$. \\
		Hàm số đạt cực tiểu tại $x=2 \Rightarrow y'(2)=0 \Leftrightarrow 3 \cdot 2^2-6 \cdot 2+m=0 \Leftrightarrow m=0$. \\
		Ta có $y''=6x-6,y''(2)=6 \cdot 2-6=6>0$. \\
		Vậy hàm số đạt cực tiểu tại $x=2$ khi $m=0$.}
\end{bt}

\begin{bt}%[2D1H3-1]
% nguon: Archive/SP_KhoLatex/ID6_old/DeThamKhao-01.tex #3635
%[Dự án ĐGNL-ĐGTD-2022 GĐ1]%[Nguyễn Anh Quốc]
Tìm giá trị lớn nhất của hàm số $y = \dfrac{x^2 - 3x + 3}{x - 1}$ trên đoạn $\left[-2; \dfrac{1}{2}\right]$.\\
\textit{Kết quả:} \tikz\draw (0,0) rectangle (3,0.6); 
	\loigiai{Ta có $y = x - 2 + \dfrac{1}{x - 1}$. Suy ra $y' = 1 - \dfrac{1}{(x - 1)^2}$. Xét phương trình
		\begin{align*}
		1 - \dfrac{1}{(x - 1)^2} = 0 & \Leftrightarrow (x - 1)^2 = 1\\
		& \Leftrightarrow \hoac{&x - 1 = 1\\&x - 1 = - 1} \\
		& \Leftrightarrow \hoac{&x = 2 & \text{ (loại)}\\&x = 0.}
		\end{align*}
		Xét hàm số $y = \dfrac{x^2 - 3x + 3}{x - 1}$ trên đoạn $\left[-2; \dfrac{1}{2}\right]$, ta có: $y(-2) = -\dfrac{13}{3}, y(0) = -3, y(0{,}5) = -\dfrac{7}{2}$.\\
		Suy ra, $\displaystyle \max_{x \in \left[-2; \tfrac{1}{2}\right]} y = -3$.
		
	}
\end{bt}

\begin{bt}%[2D1V3-1]
% nguon: Archive/SP_KhoLatex/ID6_old/DeThamKhao-01.tex #3636
%[Dự án ĐGNL-ĐGTD-2022 GĐ1]%[Nguyễn Anh Quốc]
	Tìm tất cả các giá trị thực của tham số $m$ để giá trị nhỏ nhất của hàm số $y=\dfrac{x+m^2}{x-1}$ trên đoạn $[2;3]$ bằng $14$.\\
\textit{Kết quả:} \tikz\draw (0,0) rectangle (3,0.6);
	\loigiai{
		Tập xác định $\mathscr{D} = \mathbb{R} \backslash \{1\}$.
		$y'= \dfrac{-1-m^2}{(x-1)^2} < 0$, $ \forall x \neq 1$.\\
		Do đó hàm số nghịch biến trên $ \left( 1; + \infty \right)$ nên hàm số nghịch biến trên $ \left[ 2; 3 \right]$.\\ 
		Suy ra $y(3)$ là giá trị nhỏ nhất.\\
		Theo đề bài $y(3)=14 \Leftrightarrow \dfrac{3+m^2}{2} =14 \Leftrightarrow m= \pm 5$.}
\end{bt}

\begin{bt}%[2D1H4-1]
% nguon: Archive/SP_KhoLatex/ID6_old/DeThamKhao-01.tex #3637
%[Dự án ĐGNL-ĐGTD-2022 GĐ1]%[Nguyễn Anh Quốc]%[Thi GHK1, Trần Nguyên Hãn - Vũng Tàu, 2019]%[Nguyện Ngô, 12EX2-2019]
	Cho hàm số $y=\dfrac{\sqrt {x^2+x+1}}{x-2}$. Số đường tiệm cận của đồ thị hàm số là\\
\textit{Kết quả:} \tikz\draw (0,0) rectangle (3,0.6);
	\loigiai{
		Hàm số có tập xác định $\mathscr{D}=\mathbb{R}\setminus\{2\}$.\\
		Ta có $\lim\limits_{x\to 2^+} \dfrac{\sqrt {x^2+x+1}}{x-2}= +\infty$ nên đồ thị hàm số có đường tiệm cận đứng là $x=2$.\\
		Ta có
		\[\lim\limits_{x\to +\infty} \dfrac{\sqrt {x^2+x+1}}{x-2}
		=\lim\limits_{x\to +\infty} \dfrac{x\sqrt {1+\dfrac{1}{x}+\dfrac{1}{x^2}}}{x\left(1-\dfrac{2}{x}\right)}= 1,\,\,
		\lim\limits_{x\to -\infty} \dfrac{\sqrt {x^2+x+1}}{x-2}
		=\lim\limits_{x\to -\infty} \dfrac{-x\sqrt {1+\dfrac{1}{x}+\dfrac{1}{x^2}}}{x\left(1-\dfrac{2}{x}\right)}= -1.\]
		Do đó đồ thị hàm số có hai đường tiệm cận ngang là $y=1$ và $y=-1$.\\
		Vậy đồ thị hàm số đã cho có $3$ đường tiệm cận.
	}
	%<MyLT>
\end{bt}

\begin{bt}%[2D1V4-1]
% nguon: Archive/SP_KhoLatex/ID6_old/DeThamKhao-01.tex #3638
%[Dự án ĐGNL-ĐGTD-2022 GĐ1]%[Nguyễn Anh Quốc]%[Kiểm tra giữa HKI, Nguyễn Thượng Hiền, HCM, 2018-2019]%[Phạm Đức, 12-EX-2-2019]
	Cho hàm số $y=f(x)$ xác định trên $\mathbb{R}\setminus \{1\}$ và có bảng biến thiên như hình vẽ bên dưới. 
	\begin{center}
		\begin{tikzpicture}[scale=1, font=\footnotesize, line join=round, line cap=round, >=stealth]
		\tkzTab[nocadre=false,lgt=1.2,espcl=2.5,deltacl=0.6]
		{$x$ /0.6,$y{'}$ /0.6,$y$ /2}
		{$-\infty$, $-2$, $1$, $2$, $+\infty$}
		{,-,0,+,d,+,0,-,}
		{+/ $+\infty$, -/ $2$, +D-/ $+\infty$ / $-\infty$, +/ $3$, -/ $-\infty$}
		
\end{tikzpicture}
	\end{center}
	Số đường tiệm cận đứng của đồ thị hàm số $y=\dfrac{1}{2f(x)-5}$ là\\
\textit{Kết quả:} \tikz\draw (0,0) rectangle (3,0.6);
	\loigiai{
		Xét phương trình $2f(x)-5=0\Leftrightarrow f(x)=\dfrac{5}{2}$. Dựa vào bảng biến thiên ta thấy đường thẳng $y=\dfrac{5}{2}$ cắt đồ thị hàm số $y=f(x)$ tại $4$ điểm phân biệt, nên phương trình đang xét có $4$ nghiệm. Mặt khác tử số không có nghiệm, nên đồ thị hàm số đã cho có $4$ đường tiệm cận đứng.}
\end{bt}

\begin{bt}%[2D1H5-4]
% nguon: Archive/SP_KhoLatex/ID6_old/DeThamKhao-01.tex #3639
%[Dự án ĐGNL-ĐGTD-2022 GĐ1]%[Nguyễn Anh Quốc]
	Đường thẳng $y=2x-1$ có bao nhiêu điểm chung với đồ thị của hàm số $y=\dfrac{x^2-x-1}{x+1}$ ?\\
\textit{Kết quả:} \tikz\draw (0,0) rectangle (3,0.6);
	\loigiai{
		Ta có phương trình hoành độ các giao điểm của đường thẳng $y=2x-1$ với đồ thị hàm số $y=\dfrac{x^2-x-1}{x+1}$ là 
		{\allowdisplaybreaks
			\begin{eqnarray*}
				\heva{&\dfrac{x^2-x-1}{x+1}=2x-1 \\& x \ne -1}
				&\Leftrightarrow & \heva{&x^2-x-1=(x+1)(2x-1) \\& x \ne -1}\\
				&\Leftrightarrow & \heva{&x^2+2x=0 \\& x \ne -1} \\
				&\Leftrightarrow &  \hoac{&x=0 \\& x=-2.}
			\end{eqnarray*}
		}
		Vậy đường thẳng $y=2x-1$ cắt đồ thị của hàm số $y=\dfrac{x^2-x-1}{x+1}$ tại hai điểm $A(0;-1)$ và $B(-2;-5)$.
	}
\end{bt}

\begin{bt}%[2D1V5-5]
% nguon: Archive/SP_KhoLatex/ID6_old/DeThamKhao-01.tex #3640
%[Dự án ĐGNL-ĐGTD-2022 GĐ1]%[Nguyễn Anh Quốc]
	\immini
	{Cho hàm số $y=f(x)$ có đạo hàm liên tục trên $\mathbb{R}$ và hàm số $y=f'(x)$ có đồ thị như hình vẽ bên. Số điểm cực đại của đồ thị hàm số $g(x)=f(x^2-3x)$ là bao nhiêu?
	\textit{Kết quả:} \tikz\draw (0,0) rectangle (3,0.6);
}
	{\begin{tikzpicture}[line cap=round,line join=round,font=\footnotesize,>=stealth,scale=0.75]
		\draw[->] (-4,0)--(2,0) node[above] {$x$};
		\draw[->] (0,-5)--(0,1) node[left] {$y$};
		\draw [samples=100, domain=-3.1:1.1] plot (\x,{(\x)^3+3*(\x)^2-4});
		\fill (0,0) circle (1pt) node [below left] {$O$};
		\foreach \x in {-3,-2,-1} \fill (\x,0) circle (1pt) node [above] {$\x$};
		\fill (1,0) node[below right]{$1$} circle(1pt)
		(0,-2) node[shift={(40:2ex)}]{$-2$} circle(1pt)
		(0,-4) node[right]{$-4$} circle(1pt)
		(-3,-4) circle(1pt) (-1,-2) circle(1pt);
		\draw[dashed] (-3,0)|-(0,-4) (-1,0)|-(0,-2);
		
\end{tikzpicture}}
	\loigiai
	{
		Ta có $g'(x)=(2x-3)\cdot f'(x^2-3x)$.\\
		Khi đó $g'(x)=0\Leftrightarrow\hoac{&2x-3=0\\&x^2-3x=-2\\&x^2-3x=1}\Leftrightarrow\hoac{&x=\dfrac{3}{2}\\&x=1\\&x=2\\&x=\dfrac{3-\sqrt{13}}{2}\\&x=\dfrac{3+\sqrt{13}}{2}.}$\\
		Ta có $f'(x^2-3x)>0\Leftrightarrow x^2-3x>1\Leftrightarrow x<\dfrac{3-\sqrt{13}}{2}\text{ hoặc }x>\dfrac{3+\sqrt{13}}{2}$.\\
		Bảng xét dấu của $f'(x)$ như sau
		\begin{center}
			\begin{tikzpicture}
			\tkzTabInit[nocadre=false,lgt=2.3,espcl=2.5,deltacl=0.6]
			{$x$ /1,$2x-3$ /.6,$f'(x^2-3x)$ /.6,$g'(x)$ /0.6}{$-\infty$,$\dfrac{3-\sqrt{13}}{2}$,$1$,$\dfrac{3}{2}$,$2$,$\dfrac{3+\sqrt{13}}{2}$,$+\infty$}
			\tkzTabLine{,-,t,-,t,-,$0$,+,t,+,t,+}
			\tkzTabLine{,+,$0$,-,$0$,-,t,-,$0$,-,$0$,+}
			\tkzTabLine{,-,$0$,+,$0$,+,$0$,-,$0$,-,$0$,+}
			
\end{tikzpicture}
		\end{center}
		Suy ra hàm số $g(x)$ đạt cực đại tại $x=\dfrac{3}{2}$.
	}
\end{bt}

\begin{bt}%[2D1V5-3]
% nguon: Archive/SP_KhoLatex/ID6_old/DeThamKhao-01.tex #3646
%[Dự án ĐGNL-ĐGTD-2022 GĐ1]%[Nguyễn Anh Quốc]
	Cho hàm số $y=x^4-2x^2-1$ có đồ thị $(C)$.
	\begin{enumerate}
		\item Khảo sát sự biến thiên và vẽ đồ thị $(C)$.
		\item Tìm tham số $m$ để phương trình $\big|x^4-2x^2-1\big|=2m$ có $6$ nghiệm phân biệt
	\end{enumerate}
	\loigiai{\begin{enumerate}
			\item 	$\bullet$ Tập xác định: $\mathbb{R}$.
			
			$\bullet$ Sự biến thiên:\\
			$y'=4x^3-4x^2=4x(x^2-1)$; $y'=0\Leftrightarrow \hoac{&x=0\\&x=-1\\&x=1}$.\\
			$y'>0\Leftrightarrow x\in (-1; 0) \cup (0; +\infty)$.\\
			$y'<0\Leftrightarrow x\in (-\infty; -1)\cup (0; 1)$.\\
			Hàm số đồng biến trên các khoảng $(-1; 0)$ và $(1; +\infty)$.\\
			Hàm số nghịch biến trên các khoảng $(-\infty; -1)$ và $(0; 1)$.\\
			+ Cực trị\\
			Hàm số đạt cực đại tại $x=0, y_{\text{CĐ}}=y(0)=-1$.\\
			Hàm số đạt cực tiểu tại $x=- 1$ và $x=1$, $y_{\text{CT}}=y(\pm 1)=-2$.\\
			+ Giới hạn tại vô cực $\lim\limits_{x \to \pm \infty}y=+\infty$.\\
			+ Bảng biến thiên		
			\begin{center}
				\begin{tikzpicture}
					\tkzTabInit[lgt=1,espcl=3]{$x$/1,$y'$ /1, $y$ /2}
					{$-\infty$ , $-1$,  $0$ , $1$, $+\infty$}
					\tkzTabLine{,-,0 ,+, 0,-,0 ,+, }
					\tkzTabVar {+/$+\infty$, -/$-2$,  +/$-1$/, -/$-2$, +/$+\infty$}
				
\end{tikzpicture}
			\end{center}
			
			$\bullet$ Đồ thị:\\
			\begin{center}		
				\begin{tikzpicture}[line cap=round,line join=round,>=stealth,scale=1.2]
					\draw[->,color=black] (-3,0) -- (3,0);
					\foreach \x in {-1,1,2}
					\draw[shift={(\x,0)},color=black] (0pt,1pt) -- (0pt,-1pt) node[above] {\footnotesize $\x$};
					\foreach \x in {-2}
					\draw[shift={(\x,0)},color=black] (0pt,1pt) -- (0pt,-1pt) node[above, xshift=-3pt] {\footnotesize $\x$};
					\draw[color=black] (0pt,0) node[above right] {\footnotesize $O$};
					\draw[->,color=black] (0,-2.5) -- (0,2);
					\foreach \y in {-1}
					\draw[shift={(0,\y)},color=black] (1pt,0pt) -- (-1pt,0pt) node[above left = -2.5pt and -2.5pt] {\footnotesize $\y$};
					\foreach \y in {-2}
					\draw[shift={(0,\y)},color=black] (1pt,0pt) -- (-1pt,0pt) node[below left = -1.5pt and -2.5pt] {\footnotesize $\y$};
					\clip(-3,-2.5) rectangle (3,2);
					\draw[line width=1.0pt, color=black, smooth, samples=100, domain=-2.5:2.5] plot(\x,{(\x)^4 - 2*(\x)^2 - 1});
					\draw [dash pattern=on 2pt off 2pt] (-1,0) -- (-1,-2) -- (1,-2) -- (1,0);
					\draw (2.8,0)  node[above] {$x$}; \draw (0,1.8)  node[right] {$y$};
				
\end{tikzpicture}
			\end{center}
			\item 
			\immini{	Số nghiệm của phương trình (*) bằng số điểm chung của đồ thị$(C'): y=\big|x^4-2x^2-1\big|$ và đường thẳng $(d): y=2m$.\\
				Vậy để phương trình (*) có $6$ nghiệm phân biệt thì $(d)$ phải cắt $(C')$ tại $6$ điểm phân biệt. \\
				$\Leftrightarrow 1<2m<2 \Leftrightarrow \dfrac{1}{2}<m<1$.\\
				Vậy $\dfrac{1}{2}<m<1$ thì phương trình (*) có 4 nghiệm phân biệt}{\begin{tikzpicture}[line cap=round,line join=round,>=stealth,scale=1.2]
					\draw[->,color=black] (-3,0) -- (3,0);
					\foreach \x in {-1,1,2}
					\draw[shift={(\x,0)},color=black] (0pt,1pt) -- (0pt,-1pt) node[below] {\footnotesize $\x$};
					\foreach \x in {-2}
					\draw[shift={(\x,0)},color=black] (0pt,1pt) -- (0pt,-1pt) node[above, xshift=-3pt] {\footnotesize $\x$};
					\draw[color=black] (0pt,0) node[above right] {\footnotesize $O$};
					\draw[->,color=black] (0,-1) -- (0,3);
					\foreach \y in {-1,1,2}
					\draw[shift={(0,\y)},color=black] (1pt,0pt) -- (-1pt,0pt) node[below left = -2.5pt and -2.5pt] {\footnotesize $\y$};
					\clip(-3,-2.5) rectangle (3,3);
					\draw[line width=1.0pt, color=black, smooth, samples=100, domain=-2.5:-1.55] plot(\x,{(\x)^4 - 2*(\x)^2 - 1});
					\draw[line width=1.0pt, color=black, smooth, samples=100, domain=-1.55:1.55] plot(\x,{-(\x)^4 + 2*(\x)^2 + 1});
					\draw[line width=1.0pt, color=black, smooth, samples=100, domain=1.55:2.5] plot(\x,{(\x)^4 - 2*(\x)^2 - 1});
					\draw [dash pattern=on 2pt off 2pt] (-1,0) -- (-1,2) -- (1,2) -- (1,0);
					\draw (2.8,0)  node[above] {$x$}; \draw (0,2.8)  node[right] {$y$};
			
\end{tikzpicture}}
		\end{enumerate}		
	}
\end{bt}

\begin{bt}%[2D1V1-3]
% nguon: Archive/SP_KhoLatex/ID6_old/DeThamKhao-03.tex #3678
%[Dự án ĐGNL-ĐGTD-2022 GĐ1]%[Duong Xuan Loi]
	Có bao nhiêu giá trị nguyên của tham số $m$ để hàm số $y = \dfrac{x + 2}{x + 3m}$ đồng biến trên khoảng $\left(-\infty; -6\right)$?
	
	\textit{Kết quả:} \tikz\draw (0,0) rectangle (3,0.6);
	\loigiai{ 
		Điều kiện $x\neq -3m$.\\
		Có $y'=\dfrac{3m-2}{(x+3m)^3}$. Hàm số đồng biến trên khoảng $(-\infty;-6)$ khi và chỉ khi
		\begin{eqnarray*}
			& \heva{& 3m-2>0\\& -3m\geq -6}
			\Leftrightarrow\heva{& m>\dfrac{2}{3}\\& m\leq 2}
			\Leftrightarrow \dfrac{2}{3}<m\leq 2.
		\end{eqnarray*}
		Suy ra có hai giá trị nguyên của $m$ thỏa mãn.
	}
\end{bt}

\begin{bt}%[2D1V2-6]
% nguon: Archive/SP_KhoLatex/ID6_old/DeThamKhao-03.tex #3679
%[Dự án ĐGNL-ĐGTD-2022 GĐ1]%[Duong Xuan Loi]
	Có bao nhiêu giá trị nguyên của tham số $ m $ để hàm số $ y=|3x^4-4x^3-12x^2+m| $ có $ 7 $ điểm cực trị?
	
	\textit{Kết quả:} \tikz\draw (0,0) rectangle (3,0.6);
	\loigiai{
		Xét hàm số $ y=3x^4-4x^3-12x^2+m $ trên $ \mathbb{R} $.\\ Ta có $ y'=12x^3-12x^2-24x $, $ y'=0\Leftrightarrow 12x^3-12x^2-24x=0\Leftrightarrow\hoac{&x=0\Rightarrow y=m\\&x=-1\Rightarrow y=m-5\\&x=2\Rightarrow y=m-32}. $\\
		Bảng biến thiên
		\begin{center}
			\begin{tikzpicture}[scale=1, font=\footnotesize, line join=round, line cap=round, >=stealth]
				\tkzTabInit[nocadre=false,lgt=1.2,espcl=2.5,deltacl=0.6]
				{$x$ /0.6,$y’$ /0.6,$y$ /2}
				{$-\infty$,$-1$,$0$,$2$,$+\infty$}
				\tkzTabLine{ ,-,0,+,0,-,0,+,}
				\tkzTabVar{+/$+\infty$,-/$m-5$,+/$m$,-/$m-32$,+/$+\infty$}
			\end{tikzpicture}
		\end{center}
		Dựa vào bảng biến thiên của hàm số $ y=3x^4-4x^3-12x^2+m $, suy ra hàm số $ y=|3x^4-4x^3-12x^2+m| $ có $ 7 $ cực trị khi chỉ khi $ \heva{&m>0\\&m-5<0}\Leftrightarrow\heva{&m>0\\&m<5}\Leftrightarrow 0<m<5 $.\\
		Vì $ m\in\mathbb{Z} $ nên $ m$ nhận các giá trị là $1$, $ 2 $, $ 3 $, $4 $.
	}
\end{bt}

\begin{bt}%[2D1V3-5]
% nguon: Archive/SP_KhoLatex/ID6_old/DeThamKhao-03.tex #3680
%[Dự án ĐGNL-ĐGTD-2022 GĐ1]%[Duong Xuan Loi]
	Cho $x$, $y>0$ thỏa mãn $\heva{& x^2-xy+3=0 \\ & 2x+3y-14\le 0}$. Tính tổng giá trị lớn nhất và nhỏ nhất của biểu thức $P=3x^2y-xy^2-2x^3+2x$.
	
	\textit{Kết quả:} \tikz\draw (0,0) rectangle (3,0.6);
	\loigiai{
		Ta có $x^2-xy+3=0$ $\Leftrightarrow y=\dfrac{x^2+3}{x}$.\\
		Khi đó $2x+3y-14\le 0 \Leftrightarrow 2x+\dfrac{3x^2+9}{x}-14\le 0 \Leftrightarrow 5x^2-14x+9\le 0 \Leftrightarrow x\in \left[1;\dfrac{9}{5}\right]$.\\
		Cũng có $P=3x^2\cdot \dfrac{x^2+3}{x}-x\cdot {\left(\dfrac{x^2+3}{x}\right)}^2-2x^3+2x=5x-\dfrac{9}{x}=f(x)$.\\
		Mà $f'(x)=5+\dfrac{9}{x^2}>0$, $\forall x\in \left[1;\dfrac{9}{5}\right]$.\\
		Do đó GTLN của $P$ là $f\left(\dfrac{9}{5}\right)=4$ và GTNN của $P$ là $f(1)=-4$.\\
		Vậy tổng giá trị lớn nhất và nhỏ nhất của biểu thức $P$ bằng $0$.
	}
\end{bt}

\begin{bt}%[2D1V4-2]
% nguon: Archive/SP_KhoLatex/ID6_old/DeThamKhao-03.tex #3681
%[Dự án ĐGNL-ĐGTD-2022 GĐ1]%[Duong Xuan Loi]
	Có bao nhiêu giá trị của tham số $m$ để đồ thị  hàm số $y=\dfrac{mx^2-1}{x^2-3x+2}$ có đúng hai đường tiệm cận?
	
	\textit{Kết quả:} \tikz\draw (0,0) rectangle (3,0.6);
	\loigiai{
		Ta có $\lim \limits_{x \to \pm \infty} y=\lim \limits_{x \to \pm \infty} \dfrac{m-\dfrac{1}{x^2}}{1-\dfrac{3}{x}+\dfrac{2}{x^2}}=m $ $\Rightarrow$ tiệm cận ngang $y=m$.\\
		Để đồ thị hàm số có đúng $2$ đường tiệm cận thì đồ thị hàm số có đúng một tiệm cận đứng.\\
		Vì $x^2-3x+2=0 \Leftrightarrow \hoac{&x=1\\&x=2\\}$,
		suy ra để đồ thị hàm số đúng một đường tiệm cận đứng thì $m x^2-1=0$ có một nghiệm bằng $1$ hoặc bằng $2$.\\
		Khi đó $\hoac{&m-1=0\\&4m-1=0\\} \Leftrightarrow \hoac{&m=1\\&m=\dfrac{1}{4}.\\}$\\
		Với $m=1 \Rightarrow y=\dfrac{x^2-1}{x^2-3x+2}=\dfrac{x+1}{x-2} \Rightarrow \lim \limits_{x \rightarrow 2^+ } y=+\infty \Rightarrow $ đồ thị hàm số có tiệm cận đứng là $x=2$.\\
		Với $m=\dfrac{1}{4} \Rightarrow y=\dfrac{\dfrac{1}{4}x^2-1}{x^2-3x+2}=\dfrac{x+2}{4(x-1)}\Rightarrow \lim \limits_{x \rightarrow 1^+ } y=+\infty \Rightarrow $ đồ thị hàm số có tiệm cận đứng là $x=1$.\\
		Vậy có $2$ giá trị $m$ thỏa mãn yêu cầu bài toán.		
	}
\end{bt}

\begin{bt}%[2D1V1-3]
% nguon: Archive/SP_KhoLatex/ID6_old/DeThamKhao-07.tex #3684
%[Dự án ĐGNL-ĐGTD-2022 GĐ1]%[Dao-V- Thuy]
	Cho hàm số $f(x)=x^3-3(m+1)x^2+3m(m+2)x$. Có bao nhiêu giá trị nguyên của tham số thực $m$ để hàm số nghịch biến trên đoạn $[0;1]$?\\
	\textit{Kết quả:} \tikz\draw (0,0) rectangle (3,0.6);
	\loigiai
	{
		Đạo hàm: $f'(x)=3\left[x^2-2(m+1)x+m(m+2)\right]$; $y'=0\Leftrightarrow\hoac{ & x=m \\ & x=m+2.}$\\
		Bảng biến thiên
\begin{center}
\begin{tikzpicture}
				\tkzTabInit[nocadre=false, lgt=1.2, espcl=2.5, deltacl=0.6]{$x$/0.6, $y'$/0.6, $y$/2}{$-\infty$, $m$, $m+2$, $+\infty$}
				\tkzTabLine{,+,$0$,-,$0$,+,}
				\tkzTabVar{-/$-\infty$, +/$f(m)$, -/$f(m+2)$, +/$+\infty$}
\end{tikzpicture}
\end{center}
		Dựa vào BBT, ta có YCBT $\Leftrightarrow[0;1]\subset[m;m+2]\Leftrightarrow\heva{ & m\leq0 \\ & m+2\geq1}\Leftrightarrow-1\leq m\leq0$.\\
		Vậy có hai giá trị nguyên thỏa mãn bài toán $m\in \{-1; 0\}$.
	}
\end{bt}

\begin{bt}%[2D1H2-2]
% nguon: Archive/SP_KhoLatex/ID6_old/DeThamKhao-07.tex #3685
%[Dự án ĐGNL-ĐGTD-2022 GĐ1]%[Dao-V- Thuy]
	\immini[thm]
	{
		Cho hàm số $f(x)$ liên tục trên $\mathbb{R}$ và có đồ thị như hình vẽ bên. Hỏi hàm số có bao nhiêu điểm cực trị?
			}
	{
\begin{tikzpicture}[>=stealth,line join=round,line cap=round,font=\footnotesize,scale=.71]
			\draw[->] (-2,0)--(4,0) node[below] {$x$};
			\draw[->] (0,-1)--(0,3) node[left] {$y$};
			\draw (0,0) circle (1pt) node[below left] {$O$};
			\clip (-2,-1) rectangle (4,3);
			\draw[smooth, samples=100, domain=0:3] plot (\x, {(\x-1)^2)});
			\draw[smooth, samples=100, domain=-2:0] plot (\x, {.8*\x+1});
			\draw (1,0) circle (1pt) node[below] {$1$}
			(0,1) circle (1pt) node[left] {$1$}
			;
\end{tikzpicture}
	}
\noindent \textit{Kết quả:} \tikz\draw (0,0) rectangle (3,0.6);
	\loigiai
	{
		Dễ nhận thấy hàm số có một điểm cực trị là điểm cực tiểu tại $x=1$.\\
		Xét hàm số $f$ trên khoảng $\left( - \dfrac{1}{2}; \dfrac{1}{2}\right)$.\\
		Ta có $f(x)< f(0)$ với mọi $x\in \left( - \dfrac{1}{2}; 0\right) \cup \left( 0; \dfrac{1}{2}\right)$.\\
		Do đó $x=0$ là điểm cực đại của hàm số.\\
		Vậy hàm số có $2$ điểm cực trị.\\
\textit{Lưu ý.} Tại $x=0$ hàm số không có đạo hàm nhưng vẫn đạt cực trị tại điểm đó.
	}
\end{bt}

\begin{bt}%[2D1H2-1]
% nguon: Archive/SP_KhoLatex/ID6_old/DeThamKhao-07.tex #3686
%[Dự án ĐGNL-ĐGTD-2022 GĐ1]%[Dao-V- Thuy]
	Hàm số $y= x^5+ 2x^3 - 1$ có bao nhiêu điểm cực trị?\\
	\textit{Kết quả:} \tikz\draw (0,0) rectangle (3,0.6);
	\loigiai
	{
		Ta có $y'= 5x^4+6x^2$, $y'=0 \Leftrightarrow x^2 \left(5x^2+6\right)=0 \Leftrightarrow x^2=0$.\\
		Phương trình $y'=0$ không có nghiệm bội lẻ nên hàm số đã cho không có điểm cực trị.
	}
\end{bt}

\begin{bt}%[2D1H4-1]
% nguon: Archive/SP_KhoLatex/ID6_old/DeThamKhao-07.tex #3687
%[Dự án ĐGNL-ĐGTD-2022 GĐ1]%[Dao-V- Thuy]
	Cho hàm số $ y=f(x) $ có bảng biến thiên
\begin{center}
\begin{tikzpicture}
			\tkzTabInit[nocadre=false,lgt=1,espcl=2.5,deltacl=0.5]{$x$/.7 ,$y'$/.7,$y$/2}
			{$-\infty$ , $0$ ,$1$, $+\infty$}
			\tkzTabLine{ , - , d , - ,$0$,+ }
			\tkzTabVar{+/$2$ ,-D+/$-4$/$+\infty$,-/$-2 $ , +/$+\infty$}
\end{tikzpicture}
\end{center}
	Tổng số tiệm cận ngang và tiệm cận đứng của đồ thị hàm số đã cho là\\
	\textit{Kết quả:} \tikz\draw (0,0) rectangle (3,0.6);
	\loigiai{
		Dựa vào bảng biến thiên, ta có:
\begin{enumerate}[•]
			\item	$\lim\limits_{x\to-\infty}f(x)=2\Rightarrow y=2$ là tiệm cận ngang.
			\item $\lim\limits_{x\to 0^{+}}f(x)=+\infty\Rightarrow x=0$ là tiệm cận đứng.
\end{enumerate}
		Vậy đồ thị hàm số có tất cả $ 2 $  đường tiệm cận.
	}
\end{bt}

\begin{bt}%[2D1H5-3]
% nguon: Archive/SP_KhoLatex/ID6_old/DeThamKhao-07.tex #3689
%[Dự án ĐGNL-ĐGTD-2022 GĐ1]%[Dao-V- Thuy]
	\immini[thm]{Cho hàm số $f(x)=ax^4+bx^2+c$ $(a,b,c\in\mathbb R)$. Đồ thị của hàm số $y=f(x)$ như hình vẽ bên. Số nghiệm thực của phương trình $4f(x)-3=0$ là
	}
	{
\begin{tikzpicture}[>=stealth,line join=round,line cap=round,font=\footnotesize,scale=.91]
			\clip(-2,-1) rectangle (2.1,2.1);
			\draw[->] (-2,0) -- (2,0) node[below]{ $x$};
			\draw[->] (0,-1) -- (0,2) node[left]{ $y$};
			\draw[smooth,samples=100,domain=-1.8:1.8] plot(\x,{-(\x)^4+2*(\x)^2});
			\draw [fill=black] (0,0) circle (1pt)node[below left]{\footnotesize $O$}(-1,1) circle (1pt)(1,1) circle (1pt);
			\draw[dashed](-1,0)node[below]{\scriptsize $-1$}--(-1,1)--(0,1)node[above left]{\scriptsize $1$}--(1,1)--(1,0)node[below]{\scriptsize $1$};
\end{tikzpicture}}
\noindent\textit{Kết quả:} \tikz\draw (0,0) rectangle (3,0.6);
	\loigiai{
		\immini{Ta có $ 4f(x)-3=0\Leftrightarrow f(x)=\dfrac{3}{4} $.\\
			Số nghiệm của phương trình là số giao điểm của đồ thị hàm số $ y=f(x) $ và đường thẳng $ y=\dfrac{3}{4} $.\\
			Dựa vào đồ thị hàm số $ y=f(x) $, suy ra số nghiệm phương trình là $ 4 $.}{
\begin{tikzpicture}[scale=1.2,font=\footnotesize,line join= round,line cap=round,>=stealth,x=1cm,y=1cm]
				\clip(-2,-1) rectangle (2.1,2.1);
				\draw[->] (-2,0) -- (2,0) node[below]{ $x$};
				\draw[->] (0,-1) -- (0,2) node[left]{ $y$};
				\draw (-2,0.75) node[above right]{ $y=\frac{3}{4}$} -- (2,0.75);
				\draw[smooth,samples=100,domain=-1.8:1.8] plot(\x,{-(\x)^4+2*(\x)^2});
				\draw [fill=black] (0,0) circle (1pt)node[below left]{\footnotesize $O$}(-1,1) circle (1pt)(1,1) circle (1pt);
				\draw[dashed](-1,0)node[below]{\scriptsize $-1$}--(-1,1)--(0,1)node[above left]{\scriptsize $1$}--(1,1)--(1,0)node[below]{\scriptsize $1$};
\end{tikzpicture}
		}
	}
\end{bt}

\begin{bt}%[2D1H2-1]
% nguon: Archive/SP_KhoLatex/ID6_old/DeThamKhao-06.tex #3699
%[Dự án ĐGNL-ĐGTD-2022 GĐ1]%[Bùi Mạnh Tiến]
	Cho hàm số $y = f(x)$ liên tục trên $\mathbb{R}$ có đạo hàm $f'(x) = x(x^2 - 1)(x^2 - 3x + 2)$. Hỏi hàm số $y = f(x)$ có bao nhiêu điểm cực đại?\\
	\textit{Kết quả:} \tikz\draw (0,0) rectangle (3,0.6);
	\loigiai{
		Tập xác định của hàm số $\mathscr D = \mathbb{R}$.\\
		Ta có $f'(x) = 0 \Leftrightarrow \hoac{&x = 0\\&x^2 - 1 = 0\\&x^2 - 3x + 2 = 0} \Leftrightarrow \hoac{&x = 0\\&x = 1\\&x = -1\\&x = 2.}$\\
		Khi đó ta có bảng biến thiên sau
		\begin{center}
			\begin{tikzpicture}
				\tkzTabInit[nocadre=false,lgt=1.2,espcl=2.5,deltacl=0.6] 
				{$x$ /0.6,$f’(x)$ /0.6,$f(x)$ /2}
				{$-\infty$,$-1$,$0$,$1$,$2$,$+\infty$}
				\tkzTabLine{ ,-,$0$,+,$0$,-,$0$,-,$0$,+, }
				\tkzTabVar{+/$+\infty$,-/$f(-1)$,+/$f(0)$,R,-/$f(2)$,+/$+\infty$}
				\tkzTabIma[draw]{3}{5}{4}{$f(1)$}
			\end{tikzpicture}
		\end{center}
		Dựa vào bảng biến thiên thì hàm số có một điểm cực đại.
	}
\end{bt}

\begin{bt}%[2D1V2-6]
% nguon: Archive/SP_KhoLatex/ID6_old/DeThamKhao-06.tex #3700
%[Dự án ĐGNL-ĐGTD-2022 GĐ1]%[Bùi Mạnh Tiến]
	Cho hàm số $y=f(x)$ liên tục trên $\mathbb{R}$ và có bảng biến thiên như hình vẽ bên dưới. Hàm số $y=\left|f(x-2)+3\right|$ có bao nhiêu điểm cực trị?
	\begin{center}
		\begin{tikzpicture}[>=stealth,scale=1]
			\tkzTabInit[lgt=1.2,espcl=2.5]
			{$x$ /0.6, $f’(x)$ /0.6, $f(x)$ /2}
			{$-\infty$,$-2$,$2$,$+\infty$}
			\tkzTabLine{ ,+,z,-,z,+, }
			\tkzTabVar{-/$-\infty$,+/$3$,-/$-4$,+/$+\infty$}
		\end{tikzpicture}
	\end{center}
	\textit{Kết quả:} \tikz\draw (0,0) rectangle (3,0.6);
	\loigiai{
		Từ bảng biến thiên của hàm số $y=f(x)$ ta có bảng biến thiên của hàm số $y=f(x-2)+3$ như sau
		\begin{center}
			\begin{tikzpicture}[>=stealth,scale=1]
				\tkzTabInit[lgt=1.2,espcl=2.5]
				{$x$ /0.6, $f’(x)$ /0.6, $f(x)$ /2}
				{$-\infty$,$0$,$4$,$+\infty$}
				\tkzTabLine{ ,+,z,-,z,+, }
				\tkzTabVar{-/$-\infty$,+/$6$,-/$-1$,+/$+\infty$}
			\end{tikzpicture}
		\end{center}
		Số cực trị của hàm số $y=f(ax+b)+c$ bằng số cực trị của hàm số $y=f(x)$ và có hai cực trị, hàm số $y=f(x-2)+3$ cắt trục hoành tại $3$ điểm phân biệt nên hàm số $y=|f(x-2)+3|$ có $5$ điểm cực trị.
	}
\end{bt}

\begin{bt}%[2D1H5-3]
% nguon: Archive/SP_KhoLatex/ID6_old/DeThamKhao-06.tex #3701
%[Dự án ĐGNL-ĐGTD-2022 GĐ1]%[Bùi Mạnh Tiến]
	Cho hàm số $y=f(x)$ có đồ thị như hình vẽ bên. Số nghiệm của phương trình $ 2f(x)+1=0$ là
	\begin{center}
		\begin{tikzpicture}[>=stealth,x=1cm,y=1cm,scale=2]
			\def\a{1}
			\def\b{-2}
			\def\c{0}			
			\draw[->] (-2,0) -- (2,0)node[below]{\scriptsize $x$};		
			\draw[->] (0,-2) -- (0,1) node[left] {\scriptsize $y$};
			\draw (0,0)node[above right]{\scriptsize $O$};
			\draw (-1,0)node[above]{\scriptsize $-1$};
			\draw (1,0)node[above]{\scriptsize $1$};
			\draw (0,-1)node[below right]{\scriptsize $1$};
			\draw[dashed] (-1,0)--(-1,-1)--(1,-1)--(1,0);
			\clip (-2,-2)rectangle(2,1);
			\draw[thick,samples=150,smooth,domain=-2:2] plot(\x,{\a*(\x)^4+(\b)*(\x)^2+(\c)});
		\end{tikzpicture}
	\end{center}
	\textit{Kết quả:} \tikz\draw (0,0) rectangle (3,0.6);
	\loigiai{ 
		Phương trình: $f(x) = \dfrac{-1}{2}$. Ta tìm số giao điểm của đồ thị $y=f(x)$ và $y=\dfrac{-1}{2}$.
		Dựa vào đồ thị, ta thấy số giao điểm là $4$. Suy ra số nghiệm phương trình là $4$.
	}
\end{bt}

\begin{bt}%[2D1C2-6]
% nguon: Archive/SP_KhoLatex/ID6_old/DeThamKhao-06.tex #3703
%[Dự án ĐGNL-ĐGTD-2022 GĐ1]%[Bùi Mạnh Tiến]
	\immini{Cho hàm số bậc bốn $y=f(x)$ có đồ thị như hình bên. Số điểm cực trị của hàm số $g(x) = f(x^3-3x)$ là\\
		\textit{Kết quả:} \tikz\draw (0,0) rectangle (3,0.6);
	}{
		\begin{tikzpicture}[scale=1, font=\footnotesize, line join=round, line cap=round, >=stealth]
			\draw[->] (-3.1,0)--(3.1,0) node[below left] {$x$};
			\draw[->] (0,-1.1)--(0,3.1) node[below left] {$y$};
			\draw (0,0) node [below left] {$O$};
			\foreach \x in {-2,2}
			\draw[thin] (\x,1pt)--(\x,-1pt) node [below] {$\x$};
			\begin{scope}
				\clip (-3,-1) rectangle (3,3);
				\draw[samples=200,domain=-3:3,smooth,variable=\x] plot (\x,{1/6*((\x)-1)*((\x)-2)*((\x)+2)*((\x)+2.5)});
			\end{scope}
		\end{tikzpicture}
	}
	\loigiai{
		Ta có $g'(x) = (3x^2-3)f'(x^3-3x)$.\\
		Vì hàm số không đổi dấu khi $x$ đi qua giá trị nghiệm bội chẵn của hàm số đó. Do đó để tìm số cực trị của $g(x)$ ta chỉ cần tìm số nghiệm bội lẻ của phương trình $g'(x)=0$.\\
		Xét phương trình $g'(x) = 0 \Leftrightarrow \hoac{&x=\pm 1\\&f'(x^3-3x)=0} \Leftrightarrow \hoac{&x=\pm 1\\&x^3-3x=a\in(-3;-2)\\&x^3-3x=b\in(-1;0)\\&x^3-3x=c\in(1;2).}$\\
		Ta lại có bảng biến thiên của $y=x^3-3x$ như sau
		\begin{center}
			\begin{tikzpicture}
				\tkzTabInit[nocadre=false,lgt=1.2,espcl=2.5,deltacl=0.6]
				{$x$/0.7,$y'$/0.7,$y$/2.1}
				{$-\infty$,$-1$,$1$,$+\infty$}
				\tkzTabLine{,+,0,-,0,+,}
				\tkzTabVar{-/$-\infty$,+/$2$,-/$-2$,+/$+\infty$}
			\end{tikzpicture}
		\end{center}
		Dựa vào bảng ta thấy 
		\begin{itemize}
			\item $x^3-3x=a$ có một nghiệm $x_1<-1$.
			\item $x^3-3x=b$ có ba nghiệm $x_2<-1<x_3<1<x_4$.
			\item $x^3-3x=b$ có ba nghiệm $x_5<-1<x_6<1<x_7$.
		\end{itemize}
		Và các nghiệm trên đều không trùng nhau (do $a\neq b$, $a\neq c$, $b\neq c$).\\
		Tóm lại $g'(x)=0$ có $9$ nghiệm đơn phân biệt, hay $g(x)$ có $9$ điểm cực trị.
	}
\end{bt}

\begin{bt}%[2D1V1-3]
% nguon: Archive/SP_KhoLatex/ID6_old/DeThamKhao-06.tex #3704
%[Dự án ĐGNL-ĐGTD-2022 GĐ1]%[Bùi Mạnh Tiến]
	Tìm tham số $ m $ để hàm số $y=x^3-3x^2-3mx+1 $ đồng biến trên khoảng $ (0;+\infty) $.
	\loigiai{
		Tập xác định $ \mathscr{D}=\mathbb{R} $.\\
		Yêu cầu bài toán tương đương với 
		\[y'=3x^2-6x-3m\ge 0, \, \forall x\in (0;+\infty) \Leftrightarrow x^2-2x\ge m \, \forall x\in (0;+\infty). \]
		Lập bảng biến thiên hàm số $ y=h(x)=x^2-2x $ trên khoảng $ (0;+\infty) $.
		\begin{center}
			\begin{tikzpicture}[scale=1, font=\footnotesize, line join=round, line cap=round, >=stealth]
				\tkzTabInit[nocadre=false,lgt=1,espcl=2,deltacl=0.5]{$x$/.7 ,$h'$/.7,$h$/2}
				{$0$ , $1$ , $+\infty$}
				\tkzTabLine{ , - , $0$ , + , }
				\tkzTabVar{+/$0$ ,-/$-1$ , +/$+\infty$}
			\end{tikzpicture}
		\end{center}
		Suy ra $ \displaystyle m\le \min_{(0;+\infty)}h(x)=-1 $.
	}
\end{bt}

\begin{bt}%[2D1H2-1]
% nguon: Archive/SP_KhoLatex/ID6_old/DeThamKhao-04.tex #3706
%[ĐGNL-ĐGTD-2022, nguyễn hoàng thanh]
Cho hàm số $f(x)$ có đạo hàm $f'(x)=x^3(x-1)(x-2)$, $\forall x \in \mathbb{R}$. Tìm số điểm cực trị của hàm số $f(x)$.\\
\textit{Kết quả:} \tikz\draw (0,0) rectangle (3,0.6);
\loigiai{
Ta có $f'(x)=0 \Leftrightarrow x^3(x-1)(x-2)=0 \Leftrightarrow \hoac{&x=0 \\&x=1 \\ &x=2.}$\\
Bảng xét dấu đạo hàm
\begin{center}
\begin{tikzpicture}
\tkzTabInit[lgt=1.2,espcl=2]
{$x$ /0.7, $f'(x)$ /0.7}
{$-\infty$, $ 0 $, $ 1 $, $ 2 $, $+\infty$}
\tkzTabLine{ ,-,0,+,0,-,0,+, }
\end{tikzpicture}
\end{center}
Dựa vào bảng trên ta thấy hàm số có $3$ cực trị.
}
\end{bt}

\begin{bt}%[2D1H5-3]
% nguon: Archive/SP_KhoLatex/ID6_old/DeThamKhao-04.tex #3707
%[ĐGNL-ĐGTD-2022, nguyễn hoàng thanh]
\immini[thm]{Cho hàm số bậc ba $ y=f(x) $ có đồ thị là đường cong như hình vẽ bên. Phương trình $ 2f(x)-5=0 $ có bao nhiêu nghiệm thực?\\
\textit{Kết quả:} \tikz\draw (0,0) rectangle (3,0.6);
}
{
\begin{tikzpicture}[scale=.6, font=\footnotesize, line join=round, line cap=round,>=stealth]
\def\a{1/16}
\def\xmin{-3} \def\xmax{3}
\def\ymin{-2} \def\ymax{4}
\draw[->] (\xmin,0)--(\xmax,0) node [below]{$x$};
\draw[->] (0,\ymin)--(0,\ymax) node [left]{$y$};
\draw[fill=black] (0,0) circle (.5pt) node[above left]{\footnotesize $O$};
\draw[fill=black] (1,-1) circle (.5pt);
\draw[fill=black] (-1,3) circle (.5pt);
\draw[fill=black] (-1,0) circle (.5pt) node[below]{\footnotesize $-1$};
\draw[fill=black] (1,0) circle (.5pt) node[above]{\footnotesize $1$};
\draw[fill=black] (0,3) circle (.5pt) node[right]{\footnotesize $3$};
\draw[fill=black] (0,-1) circle (.5pt) node[left]{\footnotesize $-1$};
\clip (\xmin+0.1,\ymin+0.1) rectangle (\xmax-0.1,\ymax-0.1);
\draw[smooth,samples=300,domain=-3:3] plot(\x,{(\x)^3-3*(\x)+1});
\draw[dashed] (-1,0) |- (-1,3) |- (0,3) (1,0) |- (1,-1) |- (0,-1);
\end{tikzpicture}
}
\loigiai{
\immini{
Ta có $ 2f(x)-5=0 \Leftrightarrow f(x)=\dfrac{5}{2}$. $ \quad (1) $\\
Số nghiệm thực của phương trình $ (1) $ là số giao điểm của đồ thị hàm số $ y=f(x) $ và đường thẳng $ y=\dfrac{5}{2} $.\\
Dựa vào hình vẽ ta thấy đồ thị hàm số $ y=f(x) $ và đường thẳng $ y=\dfrac{5}{2} $ có $ 3 $ giao điểm. Vậy phương trình $ (1) $ có $ 3 $ nghiệm thực.
}
{
\begin{tikzpicture}[scale=.65, font=\footnotesize, line join=round, line cap=round,>=stealth]
\def\a{1/16}
\def\xmin{-3} \def\xmax{3.5}
\def\ymin{-2} \def\ymax{4}
\draw[->] (\xmin,0)--(\xmax,0) node [below]{$x$};
\draw[->] (0,\ymin)--(0,\ymax) node [left]{$y$};
\draw[fill=black] (0,0) circle (.5pt) node[above left]{\footnotesize $O$};
\draw[fill=black] (1,-1) circle (.5pt);
\draw[fill=black] (-1,3) circle (.5pt);
\node at (2.5,1.7) [above right]{$y=\frac{5}{2}$};
\draw[fill=black] (-1,0) circle (.5pt) node[below]{\footnotesize $-1$};
\draw[fill=black] (1,0) circle (.5pt) node[above]{\footnotesize $1$};
\draw[fill=black] (0,3) circle (.5pt) node[right]{\footnotesize $3$};
\draw[fill=black] (0,-1) circle (.5pt) node[left]{\footnotesize $-1$};
\clip (\xmin+0.1,\ymin+0.1) rectangle (\xmax-0.1,\ymax-0.1);
\draw[smooth,samples=300,domain=-3:3] plot(\x,{(\x)^3-3*(\x)+1});
\draw[dashed] (-1,0) |- (-1,3) |- (0,3) (1,0) |- (1,-1) |- (0,-1);
\draw (-3,2.5) |- (3,2.5);
\end{tikzpicture}
}
}
\end{bt}

\begin{bt}%[2D1H4-1]
% nguon: Archive/SP_KhoLatex/ID6_old/DeThamKhao-04.tex #3708
%[ĐGNL-ĐGTD-2022, nguyễn hoàng thanh]
Tổng số đường tiệm cận đứng và tiệm cận ngang của đồ thị hàm số $y=\dfrac{\sqrt{x-1}}{2x-3}$ là bao nhiêu?\\
\textit{Kết quả:} \tikz\draw (0,0) rectangle (3,0.6);
\loigiai{
Tập xác định $\mathscr{D}=[1;+\infty)\setminus\left\{\dfrac{3}{2}\right\}$.\\
Ta có
\begin{itemize}
\item $\lim\limits_{x\to +\infty}f(x)=0$ nên đường thẳng $y=0$ là tiệm cận ngang,
\item $\lim\limits_{x\to \tfrac{3}{2}^-} f(x)=-\infty$, $\lim\limits_{x\to \tfrac{3}{2}^+} f(x)=+\infty$ nên đường thẳng $x=\dfrac{3}{2}$ là tiệm cận đứng.
\end{itemize}
Vậy tổng số đường tiệm cận đứng và ngang của đồ thị hàm số đã cho là $2$.
}
\end{bt}

\begin{bt}%[2D1H3-2]
% nguon: Archive/SP_KhoLatex/ID6_old/DeThamKhao-04.tex #3711
%[ĐGNL-ĐGTD-2022, nguyễn hoàng thanh]
Tìm giá trị nhỏ nhất của hàm số $y=x+\dfrac{4}{x}$ trên khoảng $( 0;+\infty )$.\\
\textit{Kết quả:} \tikz\draw (0,0) rectangle (3,0.6);
\loigiai{
Ta có $y'=1-\dfrac{4}{x^2}=\dfrac{x^2-4}{x^2}$ với $x>0$.\\
Suy ra $y'=0\Leftrightarrow x^2-4=0\Leftrightarrow x=2$ vì $x>0$.\\
Bảng biến thiên
\begin{center}
\begin{tikzpicture}
\tkzTabInit[nocadre=false,lgt=1.2,espcl=2.5,deltacl=0.6]
{$x$ /0.7,$y’$ /0.6,$y$ /2}
{$0$,$2$,$+\infty$}
\tkzTabLine{d,-,0,+,}
\tkzTabVar{+/$+\infty$,-/$4$,+/$+\infty$}
\end{tikzpicture}
\end{center}
Vậy giá trị nhỏ nhất của hàm số bằng $4$ tại $x=2.$
}
 \end{bt}

\begin{bt}%[2D1V2-6]
% nguon: Archive/SP_KhoLatex/ID6_old/DeThamKhao-04.tex #3712
%[ĐGNL-ĐGTD-2022, nguyễn hoàng thanh]
Có bao nhiêu giá trị nguyên của tham số $m$ để hàm số $y=\left|x^{3}-3x^{2}+m-4\right|$ có đúng $5$ điểm cực trị?\\
\textit{Kết quả:} \tikz\draw (0,0) rectangle (3,0.6);
\loigiai{
Xét hàm số $f(x) = x^{3}-3x^{2}+m-4 $ trên $\mathbb{R}$.\\
Ta có $f'(x) = 3x^2 - 6x$ nên $f'(x) = 0 \Leftrightarrow \hoac{&x=0\\ & x=2.}$\\
Do đó bảng biến thiên của hàm số $f(x)$ là
\begin{center}
\begin{tikzpicture}[line join=round, line cap = round, >=stealth, scale=.8,font=\footnotesize,transform shape]
\tkzTabInit[nocadre=false,lgt=1.2,espcl=2.5,deltacl=0.6]
{$x$ /0.6, $f'(x)$ /0.6, $f(x)$ /2}
{$-\infty$,$0$,$2$,$+\infty$}
\tkzTabLine{,+,$0$,-,$0$,+,}
\tkzTabVar{-/$-\infty$,+/$m-4$,-/$m-8$,+/$+\infty$}
\end{tikzpicture}
\end{center}
Ta để ý rằng các điểm cực trị của hàm số $y=|f(x)|$ bao gồm các điểm cực trị của $y=f(x)$ tại $x=0$, $x=2$ và các điểm cực trị tại nghiệm của phương trình $f(x) = 0$.\\
Do đó hàm số $y=|f(x)|$ có đúng $5$ điểm cực trị khi và chỉ khi phương trình $f(x)=0$ có đúng $3$ nghiệm khác $0$ và $2$.\\
Nhìn vào bảng biến thiên ta điều này xảy ra khi và chỉ khi $m-8 < 0 < m-4 \Leftrightarrow 4 < m < 8$.\\
Do đó có $3$ giá trị nguyên của $m$ là $5;6;7$.
}
\end{bt}

\begin{bt}%[2D1V5-1]
% nguon: Archive/SP_KhoLatex/ID6_old/DeThamKhao-04.tex #3713
%[ĐGNL-ĐGTD-2022, nguyễn hoàng thanh]
Cho hàm số $ y=f(x) =x^3+mx^2+nx-1 $, với $ m, n$  là các tham số thực thỏa mãn
$ \heva{&m+n>0\\&7+2(2m+n) <0} $. Hãy cho biết số cực trị của hàm số $ y=\left| f(|x|)\right |$.\\
\textit{Kết quả:} \tikz\draw (0,0) rectangle (3,0.6);
\loigiai{
Ta có $ f(x) $ là hàm đa thức liên tục trên $ \mathbb{R} $ và
\begin{itemize}
\item $ f(0)=-1 $ và $f(1)=m+n>0$. Do đó $ f(x) =0$ có ít nhất một nghiệm thuộc khoảng $ (0;1) $.
\item $ f(1)=m+n>0 $ và $ f(2)=7+2(2m+n)<0 $, suy ra $ f(x)=0 $ có ít nhất một nghiệm thuộc khoảng $ (1;2) $.
\item Ta có $ f(2)=7+2(2m+n)<0 $ và	$ \lim\limits_{x \rightarrow +\infty} f(x) =+\infty $ nên tồn tại số $  \alpha $ đủ lớn sao cho $ f(\alpha)>0 $ suy ra $ f(x)=0 $ có ít nhất một nghiệm thuộc khoảng $ (2; \alpha) $.
\end{itemize}
Mặt khác $ f(x) $ là hàm bậc ba nên $ f(x) =0$  có tối đa ba nghiệm. Do đó phương trình $ f(x)=0 $ có 3 nghiệm dương phân biệt với $ m,n $ thỏa mãn giả thiết.\\
Vậy
\begin{itemize}
\item  Hàm số $ y=f(x) $ có hai điểm cực trị dương
\item Suy ra hàm số $ y=f\left (|x|\right ) $ có $ 5 $ điểm cực trị.
\item Hàm số $ y=\left|f\left(|x|\right)\right |   $ có $ 11 $ điểm cực trị.
\end{itemize}
}
\end{bt}

\begin{bt}%[2D1H2-1]
% nguon: Archive/SP_KhoLatex/ID6_old/DeThamKhao-05.tex #3719
%[ĐGNL-ĐGTD-2022]
Cho hàm số $f(x)$ có đạo hàm $f'(x)=-2x^2(x+1)(x+5)^5$, $\forall x \in \mathbb{R}$. Cho biết số điểm cực trị của hàm số $f(x)$ là bao nhiêu?\\
\fbox{\textbf{Đáp số} $\ldots\ldots\ldots$}
\loigiai{
Ta có $f'(x)=0\Leftrightarrow\hoac{&x^2=0\\&x+1=0\\&(x+5)^5=0}\Leftrightarrow\hoac{&x=0\\&x=-1\\&x=-5.}$\\
Phương trình $f'(x)=0$ có $2$ nghiệm bội lẻ nên hàm số đã cho có $2$ điểm cực trị.
}
\end{bt}

\begin{bt}%[2D1H5-3]
% nguon: Archive/SP_KhoLatex/ID6_old/DeThamKhao-05.tex #3720
%[ĐGNL-ĐGTD-2022]
\immini[thm]{$f(x)=ax^4+bx^2+c$ $(a,b,c\in\mathbb R)$. Đồ thị của hàm số $y=f(x)$ như hình vẽ bên. Số nghiệm thực của phương trình $4f(x)-3=0$ là bao nhiêu?\\
\fbox{\textbf{Đáp số} $\ldots\ldots\ldots$}
}
{
\begin{tikzpicture}[scale=.8,font=\footnotesize,line join= round,line cap=round,>=stealth,x=1cm,y=1cm]
	\clip(-2,-1) rectangle (2.1,2.1);
	\draw[->] (-2,0) -- (2,0) node[below]{ $x$};
	\draw[->] (0,-1) -- (0,2) node[left]{ $y$};
	\draw[smooth,samples=100,domain=-1.8:1.8] plot(\x,{-(\x)^4+2*(\x)^2});
	\draw [fill=white] (0,0) circle (1pt)node[below left]{\footnotesize $O$}(-1,1) circle (1pt)(1,1) circle (1pt);
	\draw[dashed](-1,0)node[below]{\scriptsize $-1$}--(-1,1)--(0,1)node[above left]{\scriptsize $1$}--(1,1)--(1,0)node[below]{\scriptsize $1$};
\end{tikzpicture}
}
\loigiai{
\immini{Ta có $ 4f(x)-3=0\Leftrightarrow f(x)=\dfrac{3}{4} $.\\
	Số nghiệm của phương trình là số giao điểm của đồ thị hàm số $ y=f(x) $ và đường thẳng $ y=\dfrac{3}{4} $.\\
	Dựa vào đồ thị hàm số $ y=f(x) $, suy ra số nghiệm phương trình là $ 4 $.}{
	\begin{tikzpicture}[scale=1.2,font=\footnotesize,line join= round,line cap=round,>=stealth,x=1cm,y=1cm]
		\clip(-2,-1) rectangle (2.1,2.1);
		\draw[->] (-2,0) -- (2,0) node[below]{ $x$};
		\draw[->] (0,-1) -- (0,2) node[left]{ $y$};
		\draw (-2,0.75) node[above right]{ $y=\frac{3}{4}$} -- (2,0.75);
		\draw[smooth,samples=100,domain=-1.8:1.8] plot(\x,{-(\x)^4+2*(\x)^2});
		\draw [fill=white] (0,0) circle (1pt)node[below left]{\footnotesize $O$}(-1,1) circle (1pt)(1,1) circle (1pt);
		\draw[dashed](-1,0)node[below]{\scriptsize $-1$}--(-1,1)--(0,1)node[above left]{\scriptsize $1$}--(1,1)--(1,0)node[below]{\scriptsize $1$};
	\end{tikzpicture}
}
}
\end{bt}

\begin{bt}%[2D1H4-1]
% nguon: Archive/SP_KhoLatex/ID6_old/DeThamKhao-05.tex #3721
%[ĐGNL-ĐGTD-2022]
Tìm tất cả các tiệm cận đứng của đồ thị hàm số $y=\dfrac{2x-1-\sqrt{x^2+x+3}}{x^2-5x+6}$.\\
\fbox{\textbf{Đáp số} $\ldots\ldots\ldots$}
\loigiai{
Tập xác định $\mathscr{D}=\mathbb{R}\setminus \{2;3\}.$
Ta có 
	\begin{eqnarray*}
	\lim\limits_{x\to 2^+}\dfrac{2x-1-\sqrt{{x^2+x+3}}}{x^2-5x+6}&=&\lim\limits_{x\to 2^+} \dfrac{(2x-1)^2-(x^2+x+3)}{(x-2)(x-3)\left(2x-1+\sqrt{x^2+x+3}\right)}\\
	&=&\lim\limits_{x\to 2^+}\dfrac{(x-2)(3x+1)}{(x-2)(x-3)\left(2x-1+\sqrt{x^2+x+3}\right)}\\
	&=&\lim\limits_{x\to 2^+}\dfrac{(3x+1)}{(x-3)\left(2x-1+\sqrt{x^2+x+3}\right)}=\dfrac{-7}{6}.
	\end{eqnarray*}
Tương tự ta có 	$\lim\limits_{x\to 2^-}\dfrac{2x-1-\sqrt{{x^2+x+3}}}{x^2-5x+6} =\dfrac{-7}{6}$.\\
Mặt khác $	\lim\limits_{x\to 3^+}\dfrac{2x-1-\sqrt{{x^2+x+3}}}{x^2-5x+6}=	\lim\limits_{x\to 3^+}\dfrac{2x-1-\sqrt{{x^2+x+3}}}{(x-2)(x-3)}= +\infty$.
Suy ra đường thẳng $x=3$ là tiệm cận đứng của đồ thị hàm số đã cho.
}
\end{bt}

\begin{bt}%[2D1H3-2]
% nguon: Archive/SP_KhoLatex/ID6_old/DeThamKhao-05.tex #3724
%[ĐGNL-ĐGTD-2022]
Tính giá trị nhỏ nhất của hàm số $y=3x+\dfrac{4}{x^2}$ trên khoảng $(0;+\infty)$.\\
\fbox{\textbf{Đáp số} $\ldots\ldots\ldots$}
\loigiai{
Ta có $y'=3-\dfrac{8}{x^3}=\dfrac{3x^3-8}{x^3}$;$y'=0\Leftrightarrow 3x^3-8=0 \Leftrightarrow x=\sqrt[3]{\dfrac{8}{3}}=x_0.$\\
Ta có bảng biến thiên
\begin{center}
	\begin{tikzpicture}
		\tkzTabInit[nocadre=false,lgt=1.2,espcl=2.5,deltacl=0.6]
		{$x$ /.7, $y'$ /.7,$y$ /2}
		{$0$ ,$x_0$ , $+\infty$}
		\tkzTabLine{ d,-,0,+, }% thêm d là không xác định
		\tkzTabVar{-D+/ / $+\infty$,  -/$3\sqrt[3]{9}$ , + / $+\infty$ }
	\end{tikzpicture}
\end{center}
Từ bảng biến thiên suy ra:$\min \limits_{(0;+\infty)} y=3\sqrt[3]{9}.$
}
 \end{bt}

\begin{bt}%[2D1V2-6]
% nguon: Archive/SP_KhoLatex/ID6_old/DeThamKhao-05.tex #3725
%[ĐGNL-ĐGTD-2022]
Cho hàm số $y=f(x)$ có đạo hàm $f'(x)=(x-1)^2\left(x^2-7x+12\right)$ với mọi $x \in \mathbb{R}$. Có bao nhiêu giá trị nguyên dương của tham số $m$ để hàm số $g(x)=f\left(x^3-3x+m\right)$ có đúng $6$ điểm cực trị?\\
\fbox{\textbf{Đáp số} $\ldots\ldots\ldots$}
\loigiai{
Ta có
$$
f'(x)=0 \Leftrightarrow(x-1)^2\left(x^2-7x+12\right)=0 \Leftrightarrow\hoac{&x=1 \\
	&x=3 \\
	&x=4.}
$$
trong đó $x=1$ là nghiệm bội $2$. \\
Ta có $g'(x)=\left(3x^2-3\right) f'\left(x^3-3x+m\right)$. Suy ra
\[
g'(x)=0 \Leftrightarrow\left(3x^2-3\right) f'\left(x^3-3x+m\right)=0 \Leftrightarrow\hoac{&x=\pm 1 \\
	&x^3-3x+m=1 \quad (1) \\
	&x^3-3x+m=3 \\
	&x^3-3x+m=4.}
\]
trong đó các nghiệm của $(1)$ là nghiệm bội chẵn. Yêu cầu bài toán trở thành $g'(x)=0$ có $6$ nghiệm bội lẻ. Xét bảng biến thiên của hàm số $y=x^3-3x$.
\begin{center}
	\begin{tikzpicture}[scale=1, font=\footnotesize, line join=round, line cap=round, >=stealth]
		\tkzTabInit[nocadre=false,lgt=1.2,espcl=2.5,deltacl=0.6]
		{$x$ /0.6,$y’$ /0.6,$y$ /2}
		{$-\infty$,$-1$,$1$,$+\infty$}
		\tkzTabLine{,+,0,-,0,+,}
		\tkzTabVar{-/$-\infty$,+/$2$,-/$-2$,+/$+\infty$}
	\end{tikzpicture}
\end{center}
Dựa vào bảng biến thiên ta thấy $g'(x)=0$ có $6$ nghiệm bội lẻ khi và chỉ khi
$$
\hoac{&\heva{&3-m \leq-2 \\
		&-2 < 4-m < 2} \\
	&\heva{&-2<3-m<2 \\
		&2 \leq 4-m}}\Leftrightarrow \hoac{&5 \leq m<6 \\
	&1<m \leq 2.}
$$
Vậy có $2$ giá trị nguyên dương của $m$ thỏa mãn yêu cầu bài toán.
}
\end{bt}

\begin{bt}%[2D1V5-1]
% nguon: Archive/SP_KhoLatex/ID6_old/DeThamKhao-05.tex #3726
%[ĐGNL-ĐGTD-2022]
	Cho hàm số $f(x)=ax^3+bx^2+cx+d$ ($a$, $b$, $c$, $d\in\mathbb{R}$) có bảng biến thiên như sau
\begin{center}
	\begin{tikzpicture}[>=stealth]
		\tkzTabInit[nocadre=false,lgt=1.2,espcl=2.5,deltacl=0.5]{$x$/.7 ,$f'(x)$/.7,$f(x)$/2}
		{$-\infty$ , $-2$ , $0$ , $+\infty$}
		\tkzTabLine{ , + , $0$ , - , $0$ , + , }
		\tkzTabVar{-/$-\infty$ , +/$1$ , -/$-1$ , +/$+\infty$}
	\end{tikzpicture}	
\end{center}
Có bao nhiêu số dương trong các số $a$, $b$, $c$, $d$?\\
\fbox{\textbf{Đáp số} $\ldots\ldots\ldots$}
\loigiai{
	Ta có $f'(x)=3ax^2+2bx+c$.\\
Từ bảng biến thiên ta suy ra
\begin{itemize}
	\item $\lim \limits_{x \to +\infty} f(x)=+\infty\Rightarrow a>0$.
	\item $f(0)=-1\Leftrightarrow d=-1<0$.
	\item Hàm số đạt cực trị tại $x=0$ và $x=-2\Rightarrow \heva{& f'(0)=0 \\ & f'(-2)=0}\Leftrightarrow\heva{& c=0 \\ & 12a-4b=0}\Leftrightarrow\heva{& c=0 \\ & b=3a>0.}$
\end{itemize}
Vậy có $2$ số dương là $a$ và $b$.
}
\end{bt}

\begin{bt}%[2D1V1-3]
% nguon: Archive/SP_KhoLatex/ID6_old/dang7.tex #3734

	Tập hợp tất cả các giá trị thực của tham số $m$ để hàm số $y=\dfrac{x+3}{x+m}$ đồng biến trên khoảng $(-6;-4)$ là
	\loigiai{
		Tập xác định $\mathscr{D}=\mathbb{R}\setminus \{-m\}$.\\
		Đạo hàm $y'=\dfrac{m-3}{(x+m)^2}$.\\
		Hàm số đồng biến trên $(-6;-4)$ khi và chỉ khi
		\[\heva{&m-3>0\\&-m\notin (-6;-4)}\Leftrightarrow \heva{&m>3\\&\hoac{&-m\le -6\\&-m\geq -4}}\Leftrightarrow \hoac{&m\geq 6\\&3<m\leq 4.}\]
		Vậy với $m\in (3;4]\cup [6;+\infty)$ thì hàm số đã cho đồng biến trên khoảng $(-6;-4)$.
	}
\end{bt}

\begin{bt}%[2D1V1-3]
% nguon: Archive/SP_KhoLatex/ID6_old/dang7.tex #3735

	Cho hàm số $ f\left( x \right)=\dfrac{mx-4}{x-m} $ ( $ m $ là tham số thực). Tập hợp tất cả các giá trị thực của tham số $m$ để hàm số nghịch biến trên khoảng $ \left( 0;2\right) $ ? 
	\loigiai{Tập xác định $\mathscr{D}=\mathbb{R}\setminus \{m\}$.\\
		Đạo hàm $ f'\left( x \right)=\dfrac{-{{m}^{2}}+4}{{{\left( x-m \right)}^{2}}} $. \\
		Hàm số nghịch biến trên khoảng $\left( 0;2 \right)$ khi và chỉ khi 
		$$\heva{
			& -{{m}^{2}}+4<0 \\ 
			& m \notin (0;2)  \\ 
		}  \Leftrightarrow \heva{
			& \hoac{&m>2\\&m<-2} \\ 
			& \hoac{&m\geq 2\\&m \leq 0}} \Leftrightarrow \hoac{&m>2\\&m<-2.}$$
		Vậy $m \in (-\infty;-2)\cup (2;+\infty)$ thì hàm số $f(x)$ nghịch biến trên khoảng $(0;2$).
	}
\end{bt}

\begin{bt}%[2D1H1-1]
% nguon: Archive/SP_KhoLatex/ID6_old/dang1.tex #3750

	Tìm các khoảng đơn điệu của hàm số $y=\dfrac{x+2}{\sqrt{x^2-x+3}}$.
	\loigiai{
		Hàm số đã cho xác định khi $x^2-x+3>0$ (đúng với mọi $x\in \mathbb{R}$).\\
		Hàm số đã cho xác định trên $ \mathscr{D}=\mathbb{R}$.\\
		Ta có $y'=\dfrac{\sqrt{x^2-x+3}-\dfrac{(x+2)(2x-1)}{2\sqrt{x^2-x+3}}}{x^2-x+3}=\dfrac{-5x+8}{2(x^2-x+2)\sqrt{x^2-x+2}}$.\\
		Cho $y'=0 \Leftrightarrow \dfrac{-5x+8}{2(x^2-x+2)\sqrt{x^2-x+2}}=0 \Leftrightarrow -5x+8=0 \Leftrightarrow x=\dfrac{8}{5}$.\\
		Bảng biến thiên:
		\begin{center}
			\begin{tikzpicture}
				\tkzTabInit[lgt=1,espcl=3]
				{$x$/0.9,$y'$/0.7,$y$/2}
				{$-\infty$,$\dfrac{8}{5}$,$+\infty$}
				\tkzTabLine{,-,$0$,+,}
				\tkzTabVar{-/$-1$,+/$\dfrac{6}{\sqrt{11}}$,-/$1$}
			\end{tikzpicture}
		\end{center}
		Hàm số đã cho đồng biến trên khoảng $\left(-\infty;\dfrac{8}{5}\right)$ và nghịch biến trên khoảng $\left(\dfrac{8}{5};+\infty \right)$.
	}
\end{bt}

\begin{bt}%[2D1H1-1]
% nguon: Archive/SP_KhoLatex/ID6_old/dang1.tex #3751

	Tìm các khoảng đơn điệu của hàm số $y=(4-3x)\sqrt{6x^2+1}$.
	\loigiai{
		Hàm số đã cho xác định trên $ \mathscr{D}=\mathbb{R}$.\\
		Ta có $y'=-3\sqrt{6x^2+1}+\dfrac{6x(4-3x)}{\sqrt{6x^2+1}}=\dfrac{-36x^2+24x-3}{\sqrt{6x^2+1}}$.\\
		Cho $y'=0 \Leftrightarrow \dfrac{-36x^2+24x-3}{\sqrt{6x^2+1}}=0 \Leftrightarrow -36x^2+24x-3=0 \Leftrightarrow \left[\begin{aligned}
			&x=\dfrac{1}{2} \\
			&x=\dfrac{1}{6}.
		\end{aligned}\right. $\\
		Bảng biến thiên
		\begin{center}
			\begin{tikzpicture}
				\tkzTabInit[lgt=1,espcl=3]
				{$x$/1,$y'$/0.7,$y$/3}
				{$-\infty$,$\dfrac{1}{6}$,$\dfrac{1}{2}$,$+\infty$}
				\tkzTabLine{,+,$0$,-,$0$,+,}
				\tkzTabVar{-/$-\infty$,+/$\dfrac{7\sqrt{42}}{12}$,-/$\dfrac{5\sqrt{10}}{4}$,+/$+\infty$}
			\end{tikzpicture}
		\end{center}
		Hàm số đã cho đồng biến trên các khoảng $\left(-\infty;\dfrac{1}{6}\right)$ và $\left(\dfrac{1}{2};+\infty \right)$.\\
		Hàm số nghịch biến trên khoảng $\left(\dfrac{1}{6};\dfrac{1}{2}\right)$.
	}
\end{bt}

\begin{bt}%[2D1H1-1]
% nguon: Archive/SP_KhoLatex/ID6_old/bai1vd-ghep20dang.tex #3762
[D1B112]
	Tìm các khoảng đơn điệu của hàm số $y=\dfrac{x+2}{\sqrt{x^2-x+3}}$.
	\loigiai{
		Hàm số đã cho xác định khi $x^2-x+3>0$ (đúng với mọi $x\in \mathbb{R}$).\\
		Hàm số đã cho xác định trên $ \mathscr{D}=\mathbb{R}$.\\
		Ta có $y'=\dfrac{\sqrt{x^2-x+3}-\dfrac{(x+2)(2x-1)}{2\sqrt{x^2-x+3}}}{x^2-x+3}=\dfrac{-5x+8}{2(x^2-x+2)\sqrt{x^2-x+2}}$.\\
		Cho $y'=0 \Leftrightarrow \dfrac{-5x+8}{2(x^2-x+2)\sqrt{x^2-x+2}}=0 \Leftrightarrow -5x+8=0 \Leftrightarrow x=\dfrac{8}{5}$.\\
		Bảng biến thiên:
		\begin{center}
			\begin{tikzpicture}
				\tkzTabInit[lgt=1,espcl=3]
				{$x$/0.9,$y'$/0.7,$y$/2}
				{$-\infty$,$\dfrac{8}{5}$,$+\infty$}
				\tkzTabLine{,-,$0$,+,}
				\tkzTabVar{-/$-1$,+/$\dfrac{6}{\sqrt{11}}$,-/$1$}
			\end{tikzpicture}
		\end{center}
		Hàm số đã cho đồng biến trên khoảng $\left(-\infty;\dfrac{8}{5}\right)$ và nghịch biến trên khoảng $\left(\dfrac{8}{5};+\infty \right)$.
	}
\end{bt}

\begin{bt}%[2D1H1-1]
% nguon: Archive/SP_KhoLatex/ID6_old/bai1vd-ghep20dang.tex #3763
[D1B113]
	Tìm các khoảng đơn điệu của hàm số $y=(4-3x)\sqrt{6x^2+1}$.
	\loigiai{
		Hàm số đã cho xác định trên $ \mathscr{D}=\mathbb{R}$.\\
		Ta có $y'=-3\sqrt{6x^2+1}+\dfrac{6x(4-3x)}{\sqrt{6x^2+1}}=\dfrac{-36x^2+24x-3}{\sqrt{6x^2+1}}$.\\
		Cho $y'=0 \Leftrightarrow \dfrac{-36x^2+24x-3}{\sqrt{6x^2+1}}=0 \Leftrightarrow -36x^2+24x-3=0 \Leftrightarrow \left[\begin{aligned}
			&x=\dfrac{1}{2} \\
			&x=\dfrac{1}{6}.
		\end{aligned}\right. $\\
		Bảng biến thiên
		\begin{center}
			\begin{tikzpicture}
				\tkzTabInit[lgt=1,espcl=3]
				{$x$/1,$y'$/0.7,$y$/3}
				{$-\infty$,$\dfrac{1}{6}$,$\dfrac{1}{2}$,$+\infty$}
				\tkzTabLine{,+,$0$,-,$0$,+,}
				\tkzTabVar{-/$-\infty$,+/$\dfrac{7\sqrt{42}}{12}$,-/$\dfrac{5\sqrt{10}}{4}$,+/$+\infty$}
			\end{tikzpicture}
		\end{center}
		Hàm số đã cho đồng biến trên các khoảng $\left(-\infty;\dfrac{1}{6}\right)$ và $\left(\dfrac{1}{2};+\infty \right)$.\\
		Hàm số nghịch biến trên khoảng $\left(\dfrac{1}{6};\dfrac{1}{2}\right)$.
	}
\end{bt}

\begin{bt}%[2D1V1-3]
% nguon: Archive/SP_KhoLatex/ID6_old/bai1vd-ghep20dang.tex #3764
[D1B141]
	Tập hợp tất cả các giá trị thực của tham số $m$ để hàm số $y=\dfrac{x+3}{x+m}$ đồng biến trên khoảng $(-6;-4)$ là
	\loigiai{
		Tập xác định $\mathscr{D}=\mathbb{R}\setminus \{-m\}$.\\
		Đạo hàm $y'=\dfrac{m-3}{(x+m)^2}$.\\
		Hàm số đồng biến trên $(-6;-4)$ khi và chỉ khi
		\[\heva{&m-3>0\\&-m\notin (-6;-4)}\Leftrightarrow \heva{&m>3\\&\hoac{&-m\le -6\\&-m\geq -4}}\Leftrightarrow \hoac{&m\geq 6\\&3<m\leq 4.}\]
		Vậy với $m\in (3;4]\cup [6;+\infty)$ thì hàm số đã cho đồng biến trên khoảng $(-6;-4)$.
	}
\end{bt}

\begin{bt}%[2D1V1-3]
% nguon: Archive/SP_KhoLatex/ID6_old/bai1vd-ghep20dang.tex #3765
[D1B142]
	Cho hàm số $ f\left( x \right)=\dfrac{mx-4}{x-m} $ ( $ m $ là tham số thực). Tập hợp tất cả các giá trị thực của tham số $m$ để hàm số nghịch biến trên khoảng $ \left( 0;2\right) $ ?
	\loigiai{Tập xác định $\mathscr{D}=\mathbb{R}\setminus \{m\}$.\\
		Đạo hàm $ f'\left( x \right)=\dfrac{-{{m}^{2}}+4}{{{\left( x-m \right)}^{2}}} $. \\
		Hàm số nghịch biến trên khoảng $\left( 0;2 \right)$ khi và chỉ khi
		$$\heva{
			& -{{m}^{2}}+4<0 \\
			& m \notin (0;2)  \\
		}  \Leftrightarrow \heva{
			& \hoac{&m>2\\&m<-2} \\
			& \hoac{&m\geq 2\\&m \leq 0}} \Leftrightarrow \hoac{&m>2\\&m<-2.}$$
		Vậy $m \in (-\infty;-2)\cup (2;+\infty)$ thì hàm số $f(x)$ nghịch biến trên khoảng $(0;2$).
	}
\end{bt}

